\documentclass[10pt]{article}
\usepackage[T1]{fontenc}
\usepackage[utf8]{inputenc}
\usepackage[ngerman]{babel}
\usepackage{lmodern}
\usepackage{amsmath,amssymb,array,longtable,graphicx}
\usepackage{geometry}
\usepackage{microtype}
\geometry{paperwidth=6.5in,paperheight=9in,margin=0.52in}
\pagestyle{empty}
\setlength{\parindent}{1em}
\setlength{\parskip}{0.16em}
\setlength{\abovedisplayskip}{0.45em}
\setlength{\belowdisplayskip}{0.45em}
\emergencystretch=2em
\allowdisplaybreaks
\hbadness=10000
\hfuzz=2pt
\newcommand{\err}[4]{#1 & #2 & #3 & #4\\}
\newcommand{\tocentry}[3]{\S\,#1 & #2 & #3\\}
\newcommand{\tocsec}[2]{\multicolumn{3}{c}{\bfseries #1}\\[-0.15em]\multicolumn{3}{c}{\bfseries #2}\\[0.25em]}
\begin{document}
\small

% Source: Clebsch and Gordan, Theorie der Abelschen Functionen, 1866.
% Primary source witness: Google scan PDF theoriederabels00clebgoog.pdf, source PDF pp. 001--016.
% Non-authorial Google/library/blank scan pages 001--008, 010, 012 are included in source witnesses and provenance only.

\begin{center}
{\Large\bfseries THEORIE}\par\vspace{1.2em}
{\small\bfseries DER}\par\vspace{1.5em}
{\Large\bfseries ABELSCHEN FUNCTIONEN.}
\end{center}

\clearpage

\begin{center}
{\Large\bfseries THEORIE}\par\vspace{1.3em}
{\small\bfseries DER}\par\vspace{1.4em}
{\LARGE\bfseries ABELSCHEN FUNCTIONEN}\par\vspace{2.2em}
{\small\bfseries VON}\par\vspace{1.8em}
{\large\bfseries A. CLEBSCH \quad UND \quad P. GORDAN,}\par
{\scriptsize\bfseries PROFESSOREN AN DER UNIVERSITÄT GIESSEN.}\par\vfill
{\Large\bfseries LEIPZIG.}\par
{\bfseries DRUCK UND VERLAG VON B. G. TEUBNER.}\par
{\large 1866.}
\end{center}

\clearpage

\begin{center}
{\Large\bfseries Vorwort.}
\end{center}

Es sind nunmehr fast neun Jahre verflossen, seit der den Wissenschaften so früh entrissene Riemann durch seine unsterblichen Arbeiten die Theorie der Abelschen Functionen geschaffen, und die Grenzen der Analysis unerwartet erweitert hat. Es hat dieser Theorie seither gerechter Weise an Bewunderern nicht gefehlt, während Verständniss und Förderung nicht überall in gleichem Maasse erzielt werden konnte. Erst vor ganz kurzer Zeit hat Herr Neumann in seinem ausgezeichneten Werke über diesen Gegenstand einen grossen Theil von Riemanns Untersuchungen ausführlich dargestellt, ihre Anwendung auf die hyperelliptischen Functionen gegeben, und so dieses Gebiet dem Studium zugänglich gemacht.

Die Schwierigkeiten, auf dem Wege Riemanns in die Theorie der Abelschen Functionen einzudringen, beruhen genau auf denselben Umständen, welche die Bewunderung für den Erfinder am meisten zu erhöhen geeignet sind. Wir meinen einerseits die höchst allgemeinen und doch wieder vorsichtig beschränkten Functionsbegriffen entnommene Grundlage; andrerseits das indirecte Verfahren, durch welches die zu untersuchenden Functionen mehr synthetisch aufgebaut als analytisch entwickelt werden. Es hängt hiermit genau der von Herrn Neumann hervorgehobene Umstand zusammen, dass es bisher nicht gelang, die $\Theta$function als nothwendiges Glied der Theorie naturgemäss einzuführen.

Als die Verfasser sich mit dem Gedanken beschäftigten, diese Theorie auf eine neue Weise zu begründen, schien ihnen namentlich die ursprüngliche Quelle der ganzen Disciplin erneute Aufmerksamkeit zu verdienen, das Jacobische Umkehrproblem, auf welches Herr Weierstrass die Theorie der hyperelliptischen Functionen so glücklich gegründet. Es schien von vornherein möglich, von diesem aus direct zu einer Formel zu gelangen, welche der Jacobischen
\[
\frac{\Theta(u)}{\Theta(0)}=e^{\int_0^u Z(u)\,du}
\]
analog sei, und welche demzufolge, wie in der Theorie des Herrn Weierstrass, die Lösung des Umkehrproblems enthalte. In der That bedurfte es nur einiger geometrischen Betrachtungen und einiger Sätze aus der Theorie der Integration durch das Imaginäre, um diesen directen Weg vollständig durchzuführen. Man vermied auf diese Weise insbesondere alle Betrachtungen über Functionen im Allgemeinen, welche immer misslich sind, weil sie an einen völlig unbestimmten und unbekannte Möglichkeiten enthaltenden Begriff anknüpfen. Die Theorie der Abelschen Functionen muss ihrer Natur nach durch Sätze erledigt werden können, welche sich nur auf die bekannten, insbesondere die algebraischen Functionen beziehen. Wenn wir uns an einigen Stellen auf Sätze berufen, welche in dem Werke der Herren Briot und Bouquet gegeben sind, so verstehen wir diese nur in dem Umfange, wie sie auf algebraische und gewisse aus diesen abgeleitete Functionen Anwendung finden; in diesem Umfange sind sie richtig, bei weiterer Ausdehnung zum Theil mehr als zweifelhaft, wie wir seit Jahren in unsern an hiesiger Universität gehaltenen Vorlesungen auszuführen niemals versäumt haben.

Da es sich zunächst um eine neue Festlegung der allgemeinen Grundzüge einer sehr weiten Theorie handelt, so wird man es uns nicht verdenken, wenn wir auf Besonderheiten nirgends eingegangen sind. Solche Besonderheiten, welche immer an specielle Eigenschaften der algebraischen Fundamentalgleichung geknüpft sind, treten namentlich auch bei den hyperelliptischen Functionen hervor. Wir haben eine besondere Erwähnung der letztern uns umsomehr ersparen zu dürfen geglaubt, einmal weil es sehr leicht ist, ihre Theorie auch auf unserm Wege abzuleiten, zweitens aber, weil ihre so höchst speciellen Eigenschaften die Aufmerksamkeit von den für die allgemeine Theorie wesentlichsten Punkten abzulenken nur zu sehr geeignet sind.

Unter den von Riemann behandelten Theilen der Theorie haben wir die Frage nach der Anzahl der Moduln einer Klasse von Abelschen Functionen ausschliessen zu müssen geglaubt. Diese Frage ist durch die scharfsinnigen Betrachtungen des Herrn Cayley Gegenstand der Controverse geworden; sie ist überhaupt wohl zunächst nur durch tiefe algebraische Untersuchungen endgültig zu entscheiden, für deren Schwierigkeiten die gegenwärtig bekannten Methoden nicht mehr auszureichen scheinen. Und solche Untersuchungen in der That sind es, von denen, unserer Ueberzeugung nach, die Fortentwicklung dieser Disciplin zu erwarten ist, und in welchen die wahren tiefern Schwierigkeiten der Theorie der Abelschen Functionen erblickt werden müssen. Auch diese Disciplin endigt schliesslich in jenen Zweigen der neuern Algebra, welche zum Centrum aller neuern mathematischen Entwicklung bestimmt zu sein scheinen.

\vspace{0.6em}
\noindent Giessen, den 18. August 1866.

\begin{flushright}
{\bfseries Die Verfasser.}
\end{flushright}

\vspace{1.0em}

\begin{center}
{\Large\bfseries Druckfehler.}
\end{center}

\begin{tabular}{@{}r@{\ }l@{\ }l@{\ }l@{}}
p. & 7. & Z. 9. v. o. st. $s_x$ & l. $s_t$\\
--- & & Z. 8. v. u. st. $m+n-3$ & l. $m+n-5$\\
p. & 11. & Z. 2. v. o. st. $s_x-s_x$ & l. $s_x-s_x^\circ$\\
--- & & Z. 3. v. o. st. $\dfrac{\partial F}{\partial s}$ & l. $\dfrac{\partial F}{\partial s_x}$\\
p. & 17. & Z. 10. v. o. st. $\gamma(x-x^\circ)$ & l. $\beta(x-x^\circ)$\\
p. & 28. & Z. 6. v. u. st. $\lim S_{\xi\xi}$ & l. $\lim S_{\xi\eta}$\\
p. & 29. & Z. 2. v. o. st. $\eta$ & l. $\xi$\\
p. & 30. & Z. 14. v. o. bis p. 31. Z. 6. v. o. st. $\eta$ & l. $\xi$\\
p. & 39. & Z. 3. v. u. st. $Z_ns^n$ & l. $Z_ms^m$\\
p. & 40. & Z. 8. v. u. st. $S_nz^n$ & l. $S_mz^m$\\
p. & 74. & Z. 8. v. o. st. $w_{ik}$ & l. $w'_{ik}$\\
p. & 89. & Z. 4. v. u. st. $s_\beta$ & l. $s_\alpha$\\
p. & 93. & Z. 9. v. u. st. links & l. rechts\\
--- & & Z. 8. v. u. st. rechts & l. links\\
p. & 119. & Z. 15. v. o. st. (5) & l. (6)\\
p. & 122. & Z. 15. v. u. st. $s_x, x$ & l. für $s_x, x$\\
p. & 123. & Z. 2. v. u. st. (1) & l. (3)\\
p. & 131. & Z. 4. v. u. st. $Q(q-1)$ & l. $Q(q-2)$\\
p. & 133. & Z. 3. v. u. st. zwei & l. zwei auf einander folgende\\
p. & 157. & Z. 9. v. o. st. $p$ & l. $p-1$\\
p. & 158. & Z. 6. v. u. st. $u-1$ & l. $n-2$\\
\end{tabular}

\vspace{0.7em}
\noindent p. 159. Z. 6. 7. und 9. 10. sind die Ausdrücke
\[
\eta_1\frac{\partial f(\xi)}{\partial \xi_1}
+\eta_2\frac{\partial f(\xi)}{\partial \xi_2}
+\eta_3\frac{\partial f(\xi)}{\partial \xi_3}
\quad\text{und}\quad
\xi_1\frac{\partial f(\eta)}{\partial \eta}
+\xi_2\frac{\partial f(\eta)}{\partial \eta}
+\xi_3\frac{\partial f(\eta)}{\partial \eta}
\]
mit einander zu vertauschen.



% --- Continuation: source PDF pp. 017--038. ---
\clearpage


% Source: Clebsch and Gordan, Theorie der Abelschen Functionen, 1866.
% Primary source witness: Google scan PDF theoriederabels00clebgoog.pdf, source PDF pp. 017--038.
% Source PDF p. 022 is blank/scan whitespace and is represented in the source witness and ledger only.

\begin{center}
{\Large\bfseries Inhaltsverzeichniss.}
\end{center}

{\scriptsize
\begin{longtable}{@{}r@{\ }p{0.74\linewidth}@{\ }r@{}}
\tocsec{Erster Abschnitt.}{Die algebraischen Integrale.}
\tocentry{1.}{Darstellung der zu untersuchenden Differentialausdrücke}{1}
\tocentry{2.}{Classificirung der Integrale}{5}
\tocentry{3.}{Unstetigkeiten. Doppel- und Rückkehrpuncte}{10}
\tocentry{4.}{Integrale erster Gattung}{14}
\tocentry{5.}{Integrale dritter Gattung}{16}
\tocentry{6.}{Bildung der Normalintegrale dritter Gattung}{20}
\tocentry{7.}{Eigenschaften der Normalintegrale dritter Gattung}{22}
\tocentry{8.}{Integrale zweiter Gattung}{27}
\tocsec{Zweiter Abschnitt.}{Das Abel'sche Theorem.}
\tocentry{9.}{Schnittpunctsysteme}{34}
\tocentry{10.}{Transformation der Integralsummen, welche einem Schnittpunctsystem entsprechen}{36}
\tocentry{11.}{Der Jacobi'sche Satz}{39}
\tocentry{12.}{Das Abel'sche Theorem}{44}
\tocentry{13.}{Besondere Fälle}{47}
\tocsec{Dritter Abschnitt.}{Die eindeutigen Transformationen.}
\tocentry{14.}{Transformation der algebraischen Integrale}{50}
\tocentry{15.}{Begriff der eindeutigen Transformationen. Erhaltung der Zahl $p$}{52}
\tocentry{16.}{Directer Beweis für die Erhaltung der Zahl $p$}{54}
\tocentry{17.}{Beweis des Hülfssatzes}{60}
\tocentry{18.}{Erniedrigung der Curvenordnung. Normalcurven}{63}
\tocentry{19.}{Behandlung des Falles $p=0$. Normalcurve: eine Gerade}{67}
\tocentry{20.}{Behandlung des Falles $p=1$. Normalcurve: eine allgemeine Curve dritter Ordnung}{69}
\tocentry{21.}{Behandlung des Falles $p=2$. Normalcurve: eine Curve vierter Ordnung mit einem Doppelpunct}{77}
\tocsec{Vierter Abschnitt.}{Periodicität.}
\tocentry{22.}{Verzweigungspuncte}{80}
\tocentry{23.}{Cyclen und Umgänge}{82}
\tocentry{24.}{Reciprocität zwischen zwei Classen von Umgängen und zwischen zwei Classen von Fundamentalpuncten}{85}
\tocentry{25.}{Integration auf verschiedenen Wegen. Periodicitätsmoduln der Integrale erster Gattung}{88}
\tocentry{26.}{Relationen zwischen den Periodicitätsmoduln der Integrale erster Gattung}{90}
\tocentry{27.}{Form der gefundenen Relation und Eigenschaften ihrer ganzzahligen Coefficienten}{95}
\tocentry{28.}{Beispiel für die Aufstellung der Relation zwischen den Periodicitätsmoduln der Integrale erster Gattung}{99}
\tocentry{29.}{Transformation der Relation}{103}
\tocentry{30.}{Normalform der Integrale erster Gattung}{107}
\tocentry{31.}{Ausführung der Transformation für das in \S\ 28 behandelte Beispiel}{109}
\tocsec{Fünfter Abschnitt.}{Vertauschung von Parameter und Argument.}
\tocentry{32.}{Nähere Bestimmung des Integrals dritter Gattung}{112}
\tocentry{33.}{Vertauschung von Parameter und Argument}{114}
\tocentry{34.}{Endliche Normalform des Integrals dritter Gattung. Einfachste Form des Vertauschungssatzes. Sätze, welche daraus folgen}{116}
\tocentry{35.}{Sätze über Integrale zweiter Gattung, welche aus dem Vertauschungssatze folgen}{119}
\tocentry{36.}{Zurückführung der Differentialquotienten von Integralen dritter Gattung nach einem Parameter auf Integrale zweiter und dritter Gattung}{123}
\tocentry{37.}{Anwendung der oben für Integrale dritter Gattung benutzten Methoden auf algebraische Functionen und Logarithmen von solchen. Das Abel'sche Theorem}{124}
\tocsec{Sechster Abschnitt.}{Das Umkehrproblem.}
\tocentry{38.}{Hülfssatz zur Beurtheilung periodischer Functionen}{129}
\tocentry{39.}{Functionen mit mehreren Perioden}{134}
\tocentry{40.}{Begriff des Jacobi'schen Umkehrproblems}{136}
\tocentry{41.}{Formulirung des Umkehrproblems. Abel'sche Functionen und Transcendenten}{138}
\tocentry{42.}{Reduction des Umkehrproblems auf die Untersuchung der Transcendente $T_{\xi\eta}$}{141}
\tocentry{43.}{Erweitertes Umkehrproblem}{143}
\tocentry{44.}{Perioden der Functionen des erweiterten Umkehrproblems. Modification beim Eintreten von Integralen zweiter Gattung}{148}
\tocsec{Siebenter Abschnitt.}{Die Function $T_{\xi\eta}$.}
\tocentry{45.}{Zerlegung der Function $T_{\xi\eta}$ in die Differenz zweier Functionen, deren jede nur von $p+2$ Puncten abhängt}{152}
\tocentry{46.}{Beweis für die Zerlegbarkeit der Function $\frac12 T_{\xi\eta}(x)$ in zwei von nur $p+1$ Puncten abhängige Functionen}{155}
\tocentry{47.}{Ueber den Character der speciellen Transcendente $T_{\xi\eta}(x)$}{160}
\tocentry{48.}{Theorie der adjungirten Transcendenten}{165}
\tocentry{49.}{Die Function $U$. Die Transcendente $T_{\xi\eta}$ durch sie ausgedrückt}{170}
\tocentry{50.}{Unstetigkeiten der Function $U$. Die Function $V=e^{-U}$}{173}
\tocentry{51.}{Eigenschaften der Function $V=e^{-U}$}{175}
\tocsec{Achter Abschnitt.}{Monodromie.}
\tocentry{52.}{Differential einer Abel'schen Function. Werthsysteme der $v$, für welche das System der $x$ unbestimmt ist}{181}
\tocentry{53.}{Beweis, dass die Function $V$ eine synektische Function der $v$ ist}{187}
\tocentry{54.}{Periodicitätsmoduln von $-U=\log V$}{190}
\tocentry{55.}{Reihenentwickelung}{193}
\tocentry{56.}{Die Argumente der Function $\Theta$}{195}
\tocentry{57.}{Lösung des Umkehrproblems durch $\Theta$functionen}{198}
\tocentry{58.}{Ein Integral dritter Gattung und das Doppelverhältniss $\left(\frac{\varphi}{\psi}\right)_{\xi}:\left(\frac{\varphi}{\psi}\right)_{\eta}$ durch $\Theta$functionen dargestellt}{200}
\tocentry{59.}{Quotienten alternirender Functionen durch $\Theta$functionen dargestellt}{202}
\tocsec{Neunter Abschnitt.}{Die Function $\Theta$.}
\tocentry{60.}{Verschwinden der $\Theta$function. Unbestimmter Fall des Umkehrproblems}{206}
\tocentry{61.}{Besondere Fälle, in welchen die $\Theta$function verschwindet}{211}
\tocentry{62.}{Die Umkehrung des Abel'schen Satzes}{216}
\tocentry{63.}{Producte von $\Theta$functionen auf andere und auf algebraische Functionen zurückgeführt}{218}
\tocentry{64.}{Ueber die Darstellung der im Vorigen enthaltenen algebraischen Ausdrücke}{222}
\tocentry{65.}{Darstellung von Quotienten von $\Theta$producten durch algebraische Functionen}{224}
\tocentry{66.}{Relationen zwischen $\Theta$functionen}{226}
\tocsec{Zehnter Abschnitt.}{Die Theilung.}
\tocentry{67.}{Multiplication}{230}
\tocentry{68.}{Problem der allgemeinen Theilung}{231}
\tocentry{69.}{Das Problem der speciellen Theilung. Zurückführung der allgemeinen Theilung auf die specielle}{235}
\tocentry{70.}{Lösung der allgemeinen Theilungsgleichung}{239}
\tocentry{71.}{Zurückführung der speciellen Theilungsgleichung auf den Fall, wo $m$ Potenz einer Primzahl ist}{241}
\tocentry{72.}{Lösung der speciellen Theilungsgleichung für eine Primzahl und Potenzen einer solchen}{244}
\tocentry{73.}{Algebraisch darstellbare $\Theta$quotienten}{248}
\tocentry{74.}{Algebraische Bestimmung von $\dfrac{\Theta\left(\frac{P}{m}\right)}{\Theta(0)}$}{256}
\tocentry{75.}{Bestimmung von $\dfrac{\Theta\left(\frac{P}{2}\right)}{\Theta(0)}$. Curven $(n-3)^{\mathrm{ter}}$ Ordnung, welche in $p$ Puncten berühren, und durch die Doppel- und Rückkehrpuncte gehen}{259}
\tocentry{76.}{Reduction der Theilungsprobleme für $m=2$}{265}
\tocsec{Elfter Abschnitt.}{Das erweiterte Umkehrproblem.}
\tocentry{77.}{Zurückführung des erweiterten Umkehrproblems auf eine von nur $(p+1)$ Puncten abhängige Function $U^{(q)}$}{270}
\tocentry{78.}{Unstetigkeiten von $U^{(q)}$. Die synektische Function $V^{(q)}=e^{-U^{(q)}}$}{274}
\tocentry{79.}{Reihenentwicklung der Function $V^{(q)}$ aus ihren Periodicitätseigenschaften}{277}
\tocentry{80.}{Bestimmung der Coefficienten $A$ der Function $V^{(q)}$}{281}
\tocentry{81.}{Die Function $\Theta^{(q)}$ ist bei passender Anordnung eine Function ihrer Argumente, welche mit $q$ gerade oder ungerade ist}{287}
\tocentry{82.}{Die Theilung für das erweiterte Umkehrproblem}{290}
\tocentry{83.}{Theilung für $m=2$}{294}
\tocentry{84.}{Einfluss von Rückkehrpuncten}{296}
\tocsec{Zwölfter Abschnitt.}{Die unendlich vielen Formen der Function $\Theta$.}
\tocentry{85.}{Kanonische Determinanten}{300}
\tocentry{86.}{Reduction der kanonischen Determinanten auf einfache}{304}
\tocentry{87.}{Transformation der Integrale erster und dritter Gattung bei Veränderung des Systems von Normalperioden}{308}
\tocentry{88.}{Transformation der $\Theta$function mittelst kanonischer Substitutionen}{311}
\tocentry{89.}{Berührungscurven, welche in den untern Grenzen auftreten}{314}
\tocentry{90.}{Reduction der in der Transformationsformel für $\Theta$ auftretenden Constante auf eine numerische}{318}
\tocentry{91.}{Die numerische Constante bei Zusammensetzung von Transformationen}{321}
\tocentry{92.}{Bestimmung der numerischen Constante für die einfachen Transformationen}{326}
\tocentry{93.}{Berechnung der $a_{ik}$ für die verschiedenen Arten der $\Theta$functionen. Directe Aufstellung der Reihenentwicklungen}{331}
\end{longtable}}

\begin{center}
{\large\bfseries Erster Abschnitt.}\par
{\Large\bfseries Die algebraischen Integrale.}
\end{center}

\begin{center}
{\bfseries \S\ 1. Darstellung der zu untersuchenden Differentialausdrücke.}
\end{center}

Bezeichnen wir durch
\[
(1)\ldots\quad f(x_1,x_2,x_3)=0
\]
die Gleichung einer algebraischen Curve $n^{\mathrm{ter}}$ Ordnung, bezogen auf Dreieckscoordinaten. Wegen der beiden Gleichungen, welche daraus folgen:
\[
\frac{\partial f}{\partial x_1}x_1+\frac{\partial f}{\partial x_2}x_2+\frac{\partial f}{\partial x_3}x_3=nf=0
\]
\[
\frac{\partial f}{\partial x_1}dx_1+\frac{\partial f}{\partial x_2}dx_2+\frac{\partial f}{\partial x_3}dx_3=0,
\]
hat man die Proportion:
\[
x_2dx_3-x_3dx_2:x_3dx_1-x_1dx_3:x_1dx_2-x_2dx_1
=\frac{\partial f}{\partial x_1}:\frac{\partial f}{\partial x_2}:\frac{\partial f}{\partial x_3}.
\]
Daher sind in dem Differentialausdruck\footnote{Diese symmetrische Darstellung solcher Differentiale verdankt man Hrn. Aronhold (vgl. Monatsberichte der Berliner Akademie, Sitzung vom 25. April 1861).}
\[
\frac{\Sigma\pm c_1x_2dx_3}{c_1\frac{\partial f}{\partial x_1}+c_2\frac{\partial f}{\partial x_2}+c_3\frac{\partial f}{\partial x_3}}
\]
die in die einzelnen $c$ in Zähler und Nenner multiplicirten Terme proportional, und der ganze Ausdruck von den $c$ unabhängig.

Multiplicirt man denselben noch mit einer Function $\Phi$, welche rational und homogen in den $x$ ist, und deren Zähler um $n-3$ Ordnungen höher ist als der Nenner, so ist
\[
\frac{\Phi\,\Sigma\pm c_1x_2dx_3}{c_1\frac{\partial f}{\partial x_1}+c_2\frac{\partial f}{\partial x_2}+c_3\frac{\partial f}{\partial x_3}}
\]
ein Differentialausdruck, welcher nur von einer Variabeln abhängt, denn er ist für die $x$ homogen von der nullten Ordnung, ändert sich nicht, wenn man die $x$ gleichzeitig um einen gemeinschaftlichen Factor vermehrt, und hängt also nur von den zwei Verhältnissen der Grössen $x$ ab. Zwischen diesen aber besteht die Gleichung (1); daher hängen diese beiden Verhältnisse algebraisch zusammen, und eines kann als Function des andern aufgefasst werden. So bleibt denn schliesslich nur eine Variable zurück, während der Differentialausdruck wegen Anwendung der Gleichung (1) allerdings irrational wird; und man hat ein irrationales algebraisches Differential vor sich, dessen Irrationalitäten ausschliesslich durch die Gleichung (1) bedingt werden. Indem man dieses Differential von einem Punct der Curve (Werthsystem der $x$) bis zu einem andern integrirt, gelangt man zu dem irrationalen algebraischen Integral
\[
(2)\ldots\quad V=\int\frac{\Phi\cdot\Sigma\pm c_1x_2dx_3}{c_1\frac{\partial f}{\partial x_1}+c_2\frac{\partial f}{\partial x_2}+c_3\frac{\partial f}{\partial x_3}},
\]
welches die Grundlage der folgenden Untersuchungen bildet. Dieses Integral ist so zu nehmen, dass im Verlaufe der Integration die Variabeln $x$ stets durch die Gleichung (1) verbunden sind.

Das Integral (2) ist das allgemeinste Integral einer rationalen Function von $x_1,x_2,x_3$, welches man bilden kann. Soll nämlich irgend ein Differentialausdruck
\[
dV=A_1dx_1+A_2dx_2+A_3dx_3,
\]
in welchem die $A$ homogene Functionen $\mu^{\mathrm{ter}}$ Ordnung bedeuten, mit Hülfe von $f=0$ in einen Ausdruck mit einer Variabeln übergehen, so darf derselbe nur noch von den Verhältnissen der $x$ abhängen, darf sich also nicht ändern, wenn $kx_1,kx_2,kx_3$ an Stelle von $x_1,x_2,x_3$ gesetzt wird. Hierdurch aber erhält man
\[
dV=[(A_1x_1+A_2x_2+A_3x_3)dk+k(A_1dx_1+A_2dx_2+A_3dx_3)]k^\mu.
\]
Daher muss erstlich $k^{\mu+1}=1$, $\mu=-1$ sein, d. h. die $A$ von der $-1^{\mathrm{ten}}$ Dimension; sodann aber
\[
A_1x_1+A_2x_2+A_3x_3=nM\cdot f,
\]
oder
\[
\left(A_1-M\frac{\partial f}{\partial x_1}\right)x_1+
\left(A_2-M\frac{\partial f}{\partial x_2}\right)x_2+
\left(A_3-M\frac{\partial f}{\partial x_3}\right)x_3=0.
\]
Dieser Gleichung wird auf die allgemeinste Weise genügt, wenn man setzt:
\[
A_1=M\frac{\partial f}{\partial x_1}+C_2x_3-C_3x_2,
\]
\[
A_2=M\frac{\partial f}{\partial x_2}+C_3x_1-C_1x_3,
\]
\[
A_3=M\frac{\partial f}{\partial x_3}+C_1x_2-C_2x_1.
\]
Und unser Differentialausdruck geht mit Auslassung des verschwindenden $df$ über in:
\[
dV=\Sigma\pm C_1x_2dx_3,
\]
wo nun die $C$ von der $-2^{\mathrm{ten}}$ Dimension sind. Diese Functionen können einen gemeinsamen Factor $N$ haben. Es sei
\[
C_i=Nc_i
\]
und
\[
\Phi=N\left(c_1\frac{\partial f}{\partial x_1}+c_2\frac{\partial f}{\partial x_2}+c_3\frac{\partial f}{\partial x_3}\right);
\]
sofort geht dann $dV$ in die von uns festgestellte Form über:
\[
V=\int\frac{\Phi\,\Sigma\pm c_1x_2dx_3}{c_1\frac{\partial f}{\partial x_1}+c_2\frac{\partial f}{\partial x_2}+c_3\frac{\partial f}{\partial x_3}},
\]
wo nun $\Phi$ von der $(n-3)^{\mathrm{ten}}$ Dimension ist. Die $c$ haben dabei, soweit sie ausser in $\Phi$ hier vorkommen, keinen Einfluss auf den Werth des Integrals und können daher ausserhalb der Function $\Phi$ durch ganz beliebig gewählte Grössen ersetzt werden. Wir werden sie im Allgemeinen als Constante, d. h. als Coordinaten eines gewissen Punctes $c$ betrachten.

Das Integral $V$ kann auf mannigfache Weise in die Gestalt eines gewöhnlichen Integrals mit einer Variabeln übergeführt werden. Wir wollen dies hier zunächst mit Hülfe einer linearen Transformation bewerkstelligen. Seien $a,b$ zwei beliebig gewählte feste Puncte, welche mit $c$ zusammen ein neues Dreieck bestimmen. Führt man nun drei neue Variable $x,s_x,t$ ein mit Hülfe der Formeln
\[
(3)\ldots\quad
\begin{aligned}
\mu x_1&=a_1t+b_1x+c_1s_x,\\
\mu x_2&=a_2t+b_2x+c_2s_x,\\
\mu x_3&=a_3t+b_3x+c_3s_x,
\end{aligned}
\]
wo $\mu$ eine ganz willkürliche Grösse ist, so sind $t,x,s_x$ die Dreieckscoordinaten des Punctes $x$ bezogen auf das neue Dreieck. Setzt man diese Ausdrücke der $x_i$ in $V$ ein, so hebt sich $\mu$ nach dem Vorigen ganz heraus; es wird nämlich
\[
\mu^n f(x_1,x_2,x_3)=F(t,x,s_x),
\]
\[
\mu^{n-3}\Phi(x_1,x_2,x_3)=\Psi(t,x,s_x),
\]
\[
\begin{aligned}
c_1\frac{\partial f}{\partial x_1}+c_2\frac{\partial f}{\partial x_2}+c_3\frac{\partial f}{\partial x_3}
&=\mu\left(\frac{\partial f}{\partial x_1}\frac{\partial x_1}{\partial s_x}+\frac{\partial f}{\partial x_2}\frac{\partial x_2}{\partial s_x}+\frac{\partial f}{\partial x_3}\frac{\partial x_3}{\partial s_x}\right)\\
&=\mu\frac{\partial f}{\partial s_x}=\frac{1}{\mu^{n-1}}\frac{\partial F}{\partial s_x};
\end{aligned}
\]
endlich
\[
\begin{aligned}
\mu^2\Sigma\pm c_1x_2dx_3&=\left|\begin{array}{ccc}
c_1&a_1t+b_1x+c_1s_x&a_1dt+b_1dx+c_1ds_x\\
c_2&a_2t+b_2x+c_2s_x&a_2dt+b_2dx+c_2ds_x\\
c_3&a_3t+b_3x+c_3s_x&a_3dt+b_3dx+c_3ds_x
\end{array}\right|\\
&=\Sigma\pm a_1b_2c_3\cdot(t\,dx-x\,dt);
\end{aligned}
\]
also, wenn man $\Sigma\pm a_1b_2c_3$ in $\Psi$ eingehen lässt:
\[
V=\int\frac{\Psi(t\,dx-x\,dt)}{\frac{\partial F}{\partial s_x}},
\]
während zwischen $s_x,x,t$ die Gleichung $F=0$ besteht, welche wie die frühere vom $n^{\mathrm{ten}}$ Grade ist.

Auch der neue Ausdruck von $V$ ist homogen von der $0^{\mathrm{ten}}$ Ordnung; daher kann man ohne Beeinträchtigung der Allgemeinheit $t=1$ setzen, und erhält
\[
(4)\ldots\quad V=\int\frac{\Psi\,dx}{\frac{\partial F}{\partial s_x}},
\]
wobei $F=0$ jetzt eine Gleichung zwischen $x,s_x$ ist, welche $s_x$ als Function von $x$ definirt und bis zur $n^{\mathrm{ten}}$ Dimension ansteigt.

Die geometrische Bedeutung von $x,s_x$ für $t=1$ ist sehr einfach. Aus (3) folgen nämlich die Gleichungen:
\[
0=\Sigma\pm a_1b_2x_3+s_x\Sigma\pm c_1b_2x_3,
\]
\[
0=\Sigma\pm a_1c_2x_3+x\Sigma\pm b_1c_2x_3.
\]
Dieses sind die Gleichungen zweier Strahlbüschel, deren ersterer im Schnitt der beiden Linien
\[
\Sigma\pm a_1b_2x_3=0,
\qquad
\Sigma\pm c_1b_2x_3=0,
\]
d. h. in $b$, deren anderer im Schnitt der Linien
\[
\Sigma\pm a_1c_2x_3=0,
\qquad
\Sigma\pm b_1c_2x_3=0,
\]
d. h. in $c$ seinen Scheitel hat, und deren Parameter beziehungsweise $s_x$ und $x$ sind.

Wir werden das Integral $V$ je nach Umständen in der Form (2) oder (4) anwenden. Bemerken wir nur, dass natürlich $\Psi$ jede beliebige rationale Function von $s_x$ und $x$ sein kann, und dass also auch (4) alle algebraischen Integrale rationaler Functionen von $s_x$ und $x$ umfasst, bei denen $s_x,x$ durch die Gleichung
\[
F(s_x,x)=0
\]
verbunden sind.

\begin{center}
{\bfseries \S\ 2. Classificirung der Integrale.}
\end{center}

Das Integral $V$ enthält die gebrochene Function $\Psi$, deren Nenner von der $m^{\mathrm{ten}}$ Ordnung sein mag, und deren Zähler dann bis zu der $(m+n-3)^{\mathrm{ten}}$ Ordnung ansteigt\footnote{Die Differenz zwischen den Ordnungen in Zähler und Nenner kann scheinbar grösser werden, indem der Nenner eine Potenz von $t$ als Factor enthält, welche gleich $1$ gesetzt wird. Dies lässt sich aber offenbar durch veränderte Wahl der Puncte $c,b$, deren Verbindungslinie $t=0$ ist, immer vermeiden.}. Setzen wir
\[
\Psi=\frac{M}{N},
\]
so können wir auf diesen Bruch eine Art von Partialbruchzerlegung anwenden, und das Integral $V$ dadurch in eine Anzahl einfacherer Integrale zerlegen.

Die Curven $F=0$, $N=0$ schneiden sich in $mn$ Puncten. Nach diesen Puncten lassen wir von dem Puncte $b$ aus $mn$ Strahlen gehen, denen ebensoviel Werthe von $x$ entsprechen. Das Product der Gleichungen aller dieser Strahlen ist eine Gleichung
\[
X=0
\]
vom Grade $mn$. Sie entsteht, wenn man die Werthe von $x$ aufsucht, die mit entsprechenden Werthen von $s_x$ zusammen die Gleichungen $F=0$, $N=0$ befriedigen; d. h. wenn man $s_x$ aus den Gleichungen $F=0$, $N=0$ eliminirt. Sei, nach $s_x$ geordnet:
\[
F=f_0s_x^n+f_1s_x^{n-1}+\cdots+f_n,
\]
\[
N=n_0s_x^m+n_1s_x^{m-1}+\cdots+n_m,
\]
so ist $X$ die Determinante von $m+n$ Reihen:
\[
X=\left|\begin{array}{r|cccc}
1)&f_0&f_1&f_2&\cdots\\
2)&0&f_0&f_1&\cdots\\
3)&0&0&f_0&\cdots\\
&\cdots&\cdots&\cdots&\\
m+1)&n_0&n_1&n_2&\cdots\\
m+2)&0&n_0&n_1&\cdots\\
m+3)&0&0&n_0&\cdots\\
&\cdots&\cdots&\cdots&
\end{array}\right|.
\]
Denkt man sich die Vertikalreihen von $X$ der Reihe nach mit
\[
s_x^{m+n-1},\quad s_x^{m+n-2},\ldots,1
\]
multiplicirt, und addirt man alle so multiplicirten Reihen zur letzten, so werden die Glieder der letzten Vertikalreihe folgende:
\[
s_x^{m-1}F,\quad s_x^{m-2}F,\ldots,F,\quad s_x^{n-1}N,\quad s_x^{n-2}N,\ldots,N.
\]
Indem man daher nach dieser Reihe ordnet, erhält man $X$ in der Form
\[
(1)\ldots\quad X=AF+BN,
\]
wo $A$ und $B$ Determinanten sind, welche mit $X$ ganz übereinstimmen, nur dass die letzte Vertikalreihe bei $A$
\[
s_x^{m-1},\ s_x^{m-2},\ldots,1,\quad 0,\quad 0,\ldots,0,
\]
bei $B$ aber
\[
0,\quad 0,\ldots,0,\quad s_x^{n-1},\ s_x^{n-2},\ldots,1
\]
ist. Man kann daher die Curven $n^{\mathrm{ter}}$ und $m^{\mathrm{ter}}$ Ordnung $F=0$, $N=0$, mit den linken Theilen der Gleichungen zweier andern Curven $A=0$, $B=0$ so multipliciren, dass die Summe der Producte die Gleichung $X=0$ als Resultat ergiebt. Die geometrische Bedeutung hievon ist evident. Die Geraden des Büschels $X=0$ schneiden $F=0$ noch in $mn(n-1)$ Puncten, welche nicht auf $N=0$ liegen, und diese liegen auf einer Curve $B=0$ der Ordnung $m(n-1)$. Ebenso schneiden jene Geraden die Curve $N=0$ noch in $mn(m-1)$ Puncten, welche nicht auf $F=0$ liegen, und diese liegen auf einer Curve $A=0$ der Ordnung $n(m-1)$.

Die Curven $A=0$, $B=0$ sind durch die Gleichung (1) nicht vollständig bestimmt, denn jene Gleichung bleibt unverändert, wenn man $A$ und $B$ beziehungsweise durch
\[
A+CN,\qquad B-CF
\]
ersetzt, unter $C$ eine beliebige Curve $(mn-m-n)^{\mathrm{ter}}$ Ordnung verstanden. Aber es ist sehr wichtig zu bemerken, dass die oben erhaltenen speciellen Curven $A=0$, $B=0$ besondere wichtige Eigenschaften haben. Es enthält nämlich $A$ die Grösse $s_x$ nur zur $(m-1)^{\mathrm{ten}}$, $B$ dieselbe nur zur $(n-1)^{\mathrm{ten}}$ Potenz. Denkt man sich also $A$, $B$ durch Einführung der oben gleich $1$ gesetzten Variabeln $t$ wieder homogen gemacht, so sind alle Glieder in $A$ und $B$ für $s_x$ und $x$ wenigstens von der Ordnung $(n-1)(m-1)$; die beiden Curven $A$ und $B$ sind also hier so bestimmt, dass beide noch in dem Puncte $c$ einen $(n-1)(m-1)$ fachen Punct haben.

Multiplicirt man also den Bruch
\[
\Psi=\frac{M}{N}
\]
in Zähler und Nenner mit $B$, so kann man für den Nenner, der Gleichung $F=0$ wegen, auch $X$ setzen, und erhält also
\[
\Psi=\frac{MB}{X},
\]
wo nun der Nenner eine Function $mn^{\mathrm{ter}}$ Ordnung von $x$ allein ist, während der Zähler eine Function der $(mn+n-3)^{\mathrm{ten}}$ Ordnung ist, welche aber wegen der Eigenschaften von $B$ die Grösse $s_x$ nur bis zur $(m+2n-4)^{\mathrm{ten}}$ Potenz enthält. Es ist also $MB$ von der Form:
\[
MB=s_x^{m+2n-4}X_{(m-1)(n-1)}+s_x^{m+2n-5}X_{(m-1)(n-1)+1}+\cdots+X_{mn+n-3},
\]
wo die $X$ Functionen von $x$ allein sind, deren Ordnung durch den angehängten Index angegeben wird. Ist nun in gleicher Weise $C$ eine Function der $(mn-3)^{\mathrm{ten}}$ Ordnung, welche $s_x$ nur zur $(m+n-4)^{\mathrm{ten}}$ Potenz enthält:
\[
C=s_x^{m+n-4}\xi_{(m-1)(n-1)}+s_x^{m+n-3}\xi_{(m-1)(n-1)+1}+\cdots+\xi_{mn-3},
\]
so kann man die Functionen $\xi$ so bestimmen, dass
\[
MB+CF
\]
nur noch Potenzen von $s_x$ enthält, welche unterhalb der $n^{\mathrm{ten}}$ bleiben. Ist nämlich nur $f_0$ nicht gleich Null, was immer stattfindet, wenn der Punct $c$ ($t=0$, $x=0$) nicht auf der Curve selbst liegt, so müssen hiezu die Gleichungen bestehen:
\[
0=X_{(m-1)(n-1)}+\xi_{(m-1)(n-1)}f_0,
\]
\[
0=X_{(m-1)(n-1)+1}+\xi_{(m-1)(n-1)+1}f_0+\xi_{(m-1)(n-1)}f_1,
\]
\[
0=X_{(m-1)(n-1)+2}+\xi_{(m-1)(n-1)+2}f_0+\xi_{(m-1)(n-1)+1}f_1+\xi_{(m-1)(n-1)}f_2,
\]
\[
\cdots
\]
\[
0=X_{mn-3}+\xi_{mn-3}f_0+\xi_{mn-4}f_1+\cdots+\xi_{mn-n-3}f_n,
\]
wodurch die Functionen $\xi$, da $f_0$ eine von Null verschiedene Constante ist, vollkommen und eindeutig bestimmt sind.

Nach dieser Bestimmung der $\xi$ wird die höchste Potenz von $s_x$ in $MB+CF$ die $(n-1)^{\mathrm{te}}$, und es ist also
\[
MB+CF=X'_{mn-2}s_x^{n-1}+X'_{mn-1}s_x^{n-2}+\cdots+X'_{mn+n-3},
\]
wo die $X'$ wieder Functionen von $x$ allein sind, deren Grad der Index angiebt. Da ferner
\[
\frac{\partial F}{\partial s_x}=nf_0s_x^{n-1}+(n-1)f_1s_x^{n-2}+\cdots,
\]
so kann man den ersten Term der obigen Reihe durch Abziehen von
\[
\frac{1}{nf_0}\cdot X'_{mn-2}\frac{\partial F}{\partial s_x}=D\frac{\partial F}{\partial s_x}
\]
zerstören, und erhält einen Ausdruck der Form:
\[
MB+CF-D\frac{\partial F}{\partial s_x}=X''_{mn-1}s_x^{n-2}+X''_{mn}s_x^{n-3}+\cdots+X''_{mn+n-3}.
\]
Mit Rücksicht auf die Gleichung $F=0$ wird also jetzt
\[
\frac{M}{N}=\frac{MB-D\frac{\partial F}{\partial s_x}}{X}+\frac{D}{X}\frac{\partial F}{\partial s_x}
\]
\[
=\frac{1}{X}\left\{X''_{mn-1}s_x^{n-2}+X''_{mn}s_x^{n-3}+\cdots+X''_{mn+n-3}\right\}+\frac{D}{X}\frac{\partial F}{\partial s_x}.
\]
Unser Integral
\[
V=\int\frac{M}{N}\frac{dx}{F'_{s_x}}
\]
zerfällt hienach in zwei Theile. Der erste, welcher von dem letzten Gliede in dem obigen Ausdruck von $\frac{M}{N}$ herrührt, ist
\[
\int\frac{D}{X}\,dx,
\]
also das Integral einer rationalen Function, und kann daher durch Logarithmen und rationale Functionen integrirt werden.

Es genügt $V$ auf den Theil zu beschränken, welcher von den übrigen Gliedern des Ausdrucks für $\frac{M}{N}$ herrührt.

Den ersten Theil von $\frac{M}{N}$ können wir zunächst durch Division in eine ganze Function zerlegen, welche in Bezug auf $s_x$ und $x$ bis zur $(n-3)^{\mathrm{ten}}$ Dimension ansteigt; und in einen Rest, welcher in Bezug auf $x$ von der $(mn-1)^{\mathrm{ten}}$ Ordnung wird, und mit Partialbruchzerlegung zu behandeln ist. Bei diesem letzten Prozesse, bei welchem wir die Wurzeln von $X=0$ zunächst verschieden voraussetzen, entstehen $mn$ verschiedene Glieder, deren Zähler ganze Functionen $(n-2)^{\mathrm{ten}}$ Grades von $s_x$ sind, während der Nenner immer die Form $x-a$ hat.

Wir können also jedes Integral der hier betrachteten Form auf Aggregate folgender Arten von Integralen zurückführen:
\begin{enumerate}
\item Integrale rationaler Functionen von $x$.
\item Integrale der Form
\[
\int\frac{Q\,dx}{F'_{s_x}},
\]
wo $Q$ eine ganze Function ist, welche für $s_x$ und $x$ bis zur $(n-3)^{\mathrm{ten}}$ Dimension aufsteigt.
\item Integrale der Form
\[
\int\frac{Q\,dx}{(x-a)F'_{s_x}},
\]
wo $Q$ bis zur $(n-2)^{\mathrm{ten}}$ Dimension aufsteigt.
\end{enumerate}
Endlich können, wenn im Obigen die Gleichung $X=0$ gleiche Wurzeln hatte, noch Integrale auftreten, welche höhere Potenzen von $x-a$ im Nenner haben. Diese aber kann man durch einen Grenzübergang aus den unter 3). angeführten entstehen lassen. Nach der Theorie der Partialbrüche sind die entsprechenden Glieder von $\frac{M}{N}$ dann, wenn $\hat P$ den nach Absonderung der gleichen Factoren $(x-a)$ von $X$ übrigbleibenden Nenner bedeutet, von der Form
\[
\frac{\left(\frac{R}{P}\right)_{x=a}}{(x-a)^\mu}
+\frac{1}{1}\frac{\left(\frac{\partial}{\partial x}\frac{R}{P}\right)_{x=a}}{(x-a)^{\mu-1}}
+\frac{1}{1\cdot2}\frac{\left(\frac{\partial^2}{\partial x^2}\frac{R}{P}\right)_{x=a}}{(x-a)^{\mu-2}}+\cdots.
\]
Die hieraus entstehenden Integrale erhält man also, wenn man gewisse in der Form 3). enthaltene Integrale nach dem Parameter $a$ differenzirt\footnote{Diese Art der Zurückführung auf das Integral 3). ist hier nur zur vorläufigen Orientirung gegeben. Es wird sich später (vergl. Abschnitt 5) für die Darstellung solcher Integrale ein ganz anderer Weg zeigen. Ebenso werden wir die logarithmischen und rationalen Functionen durch Integrale der Formen 1)--3) ausdrücken lernen.}.

Die unter den Formen 2). und 3). enthaltenen Integrale sind nun genauer zu untersuchen; und zwar insbesondere bezüglich der Unstetigkeiten, welche in denselben auftreten können, und welche ihren Charakter bestimmen.

\begin{center}
{\bfseries \S\ 3. Unstetigkeiten. Doppel- und Rückkehrpuncte.}
\end{center}

Die Integrale der Form 2). sind in $x$ und $s_x$ ausgedrückt:
\[
V=\int\frac{Q\,dx}{F'_{s_x}},
\]
wo $Q$ zur $(n-3)^{\mathrm{ten}}$ Dimension ansteigt. In den $x_i$ haben sie daher die Form
\[
V=\int\frac{Q\cdot\Sigma\pm c_1x_2dx_3}{c_1\frac{\partial f}{\partial x_1}+c_2\frac{\partial f}{\partial x_2}+c_3\frac{\partial f}{\partial x_3}},
\]
wo $Q$ eine homogene Function $(n-3)^{\mathrm{ter}}$ Ordnung ist. Die zu integrirende Function kann hier nur dann unstetig werden, wenn der Nenner verschwindet, wenn also $F'(s_x)=0$ oder
\[
c_1\frac{\partial f}{\partial x_1}+c_2\frac{\partial f}{\partial x_2}+c_3\frac{\partial f}{\partial x_3}=0.
\]
Dies ist die Gleichung der Polare des Punctes $c$ in Bezug auf die Curve $f=0$; beide Curven schneiden sich in $n(n-1)$ Puncten, in welchen die von dem willkürlich gewählten Puncte $c$ an die Curve $f=0$ zu ziehenden Tangenten dieselbe berühren. In einem dieser Puncte seien $s_x^\circ,x^\circ$ die Werthe von $s_x$ und $x$. Man hat dann
\[
F(x^\circ,s_x^\circ)=0,\qquad \frac{\partial F(x^\circ,s_x^\circ)}{\partial s_x^\circ}=0;
\]
daher ist in der Nähe eines solchen Puncts, wenn $s_x-s_x^\circ$, $x-x^\circ$ sehr kleine Grössen sind, und wenn der Kürze wegen $F^0$ für $F(s_x^\circ,x^\circ)$ geschrieben wird:
\[
F(s_x,x)=0=\frac{\partial F^0}{\partial x^0}(x-x^0)+\frac12\frac{\partial^2F^0}{\partial s_x^{0\,2}}(s_x-s_x^0)^2+\cdots
\]
\[
\frac{\partial F}{\partial s_x}=\frac{\partial^2F^0}{\partial s_x^{0\,2}}(s_x-s_x^0)+\cdots,
\]
so dass der ersten Gleichung zufolge $s_x-s_x^0$ mit $\sqrt{x-x^0}$ von gleicher Ordnung wird, und also auch $\frac{\partial F}{\partial s}$, der Nenner von $V$, mit $\sqrt{x-x^0}$ proportional. Das Integral $V$ wird also, wenn sich die Integration in der Nähe dieser Puncte bewegt, proportional mit
\[
\int\frac{dx}{\sqrt{x-x^0}}=2\sqrt{x-x^0},
\]
d. h. es bleibt endlich, obgleich die zu integrirende Function unendlich wird.

Für den Berührungspunct einer von $c$ an die Curve gezogenen Tangente wird also das Integral $V$ nicht unstetig.

Aber wenn die Curve $f=0$ einige der gewöhnlichen Singularitäten (und auf diese werden wir uns beschränken)\footnote{Es mag bemerkt werden, dass ein $r$ facher Punct in allen Beziehungen immer durch $\frac{r\cdot r-1}{2}$ Doppelpuncte ersetzt gedacht werden kann, sobald nur seine Tangenten sämmtlich verschieden sind. Auf diesen Fall lässt sich z. B. die Function $f$ bei den hyperelliptischen Functionen immer zurückführen, so dass diese keine weiteren Schwierigkeiten machen. Aber besondere Betrachtungen werden erfordert, wenn mehrere der Tangenten eines vielfachen Puncts zusammenfallen.}, wenn sie Doppel- und Rückkehrpuncte besitzt, so erfordern diese Betrachtungen einige Modificationen. Bekanntlich vermindert jeder Doppelpunct die Anzahl der von einem Puncte an die Curve zu ziehenden Tangenten um zwei, jeder Rückkehrpunct um drei. Nennen wir $d$ die Anzahl der Doppelpuncte, $r$ die der Rückkehrpuncte, so ist die Anzahl der eigentlichen Tangenten, welche man von $c$ an die Curve ziehen kann, oder die Anzahl der beweglichen Schnittpuncte von Curve und Polare
\[
n(n-1)-2d-3r;
\]
die Polare jedes Punctes, also auch die von $c$, geht durch jeden Doppel- oder Rückkehrpunct der Curve, und berührt die Tangente jedes der letztern, so dass von den Schnittpuncten der Polare mit der Curve in jedem Doppelpunct zwei, in jedem Rückkehrpunct drei vereinigt und fest liegen.

Ein Doppelpunct ist bekanntlich dadurch charakterisirt, dass in ihm
\[
\frac{\partial f}{\partial x_1}=0,\qquad \frac{\partial f}{\partial x_2}=0,\qquad \frac{\partial f}{\partial x_3}=0,
\]
oder, wenn man sich der Gleichung $F(s_x,x)=0$ bedient, dass zugleich
\[
F=0,\qquad \frac{\partial F}{\partial s_x}=0,\qquad \frac{\partial F}{\partial x}=0.
\]
Entwickelt man also wie oben in der Nähe eines Doppelpuncts $s_x^\circ,x^\circ$ nach Potenzen von $s_x-s_x^\circ$, $x-x^\circ$, so hat man:
\[
\begin{aligned}
F(s_x,x)=0={}&\frac12\left\{\frac{\partial^2F^\circ}{\partial s_x^{\circ\,2}}(s_x-s_x^\circ)^2
+2\frac{\partial^2F^\circ}{\partial s_x^\circ\partial x^\circ}(s_x-s_x^\circ)(x-x^\circ)
+\frac{\partial^2F^\circ}{\partial x^{\circ\,2}}(x-x^\circ)^2\right\}+\cdots,
\end{aligned}
\]
\[
\frac{\partial F}{\partial s_x}=\frac{\partial^2F^\circ}{\partial s_x^{\circ\,2}}(s_x-s_x^\circ)+\frac{\partial^2F^\circ}{\partial x^\circ\partial s_x^\circ}(x-x^\circ)+\cdots.
\]
Aus der ersten dieser Gleichungen erkennt man, dass hier wieder $s_x-s_x^0$ in der Grenze mit $x-x^0$ selbst, nicht wie früher mit $\sqrt{x-x^0}$, proportional wird, und zwar ergeben sich für $\frac{s_x-s_x^0}{x-x^0}$ zwei Werthe aus der quadratischen Gleichung
\[
\left(\frac{s_x-s_x^0}{x-x^0}\right)^2\frac{\partial^2F^0}{\partial s_x^{0\,2}}
+2\frac{s_x-s_x^0}{x-x^0}\frac{\partial^2F^0}{\partial s_x^0\partial x^0}
+\frac{\partial^2F^0}{\partial x^{0\,2}}=0.
\]
Daher wird auch $\frac{\partial F}{\partial s_x}$ in der Nähe des Doppelpuncts mit $x-x^\circ$ proportional, und das Integral $V$ wird, wenn die Integration sich in der Nähe von $s_x^\circ,x^\circ$ bewegt, proportional mit
\[
\int\frac{dx}{x-x^0}=\log(x-x^0),
\]
d. h. das Integral wird logarithmisch unendlich. Ausgenommen hiervon ist der Fall, wo der Zähler $Q$ selbst für $s_x=s_x^\circ$, $x=x^\circ$ verschwindet. Ist dieses der Fall, so nähert sich $Q$ in der Nähe dieser Werthe nicht wie sonst einem endlichen constanten Werth, sondern man hat
\[
Q=\frac{\partial Q^\circ}{\partial s_x}(s_x-s_x^\circ)+\frac{\partial Q^\circ}{\partial x}(x-x^\circ)+\cdots;
\]
also auch $Q$ mit $x-x^\circ$ proportional. In diesem Falle nähert sich sonach $\dfrac{Q}{\frac{\partial F}{\partial s_x}}$ einer endlichen Grenze, und das Integral $V$ bleibt in der Nähe des Doppelpuncts vollkommen endlich.

Ist $s_x^\circ,x^\circ$ ein Rückkehrpunct, so tritt zu den Bedingungen des Doppelpuncts noch die hinzu, dass die obige quadratische Gleichung zwei gleiche Wurzeln habe, d. h. dass
\[
\frac{\partial^2F^0}{\partial s_x^{0\,2}}\cdot\frac{\partial^2F^0}{\partial x^{0\,2}}-
\left(\frac{\partial^2F^0}{\partial s_x^0\partial x^0}\right)^2=0.
\]
Man kann dann zwei Grössen $a,b$ so bestimmen, dass
\[
\frac{\partial^2F^0}{\partial s_x^{0\,2}}=a^2,
\qquad
\frac{\partial^2F^0}{\partial x^{0\,2}}=b^2,
\qquad
\frac{\partial^2F^0}{\partial s_x^0\partial x^0}=ab.
\]
Indem man bei der Entwicklung von $F(s_x,x)$ noch die Glieder dritter Ordnung berücksichtigt, hat man
\[
\begin{aligned}
F(s_x,x)=0={}&\frac12\left[a(s_x-s_x^0)+b(x-x^0)\right]^2\\
&+\frac16\left\{\frac{\partial^3F^0}{\partial s_x^{0\,3}}(s_x-s_x^0)^3
+3\frac{\partial^3F^0}{\partial s_x^{0\,2}\partial x^0}(s_x-s_x^0)^2(x-x^0)\right.\\
&\left.\quad+3\frac{\partial^3F^0}{\partial s_x^0\partial x^{0\,2}}(s_x-s_x^0)(x-x^0)^2
+\frac{\partial^3F^0}{\partial x^{0\,3}}(x-x^0)^3\right\}+\cdots,
\end{aligned}
\]
während
\[
\frac{\partial F}{\partial s_x}=a\left[a(s_x-s_x^0)+b(x-x^0)\right]+\cdots.
\]
Die erste Gleichung zeigt, dass die Combination
\[
a(s_x-s_x^0)+b(x-x^0)
\]
von höherer als der ersten, nämlich von der $\frac32^{\mathrm{ten}}$ Ordnung ist; daher können wir in den Gliedern $3^{\mathrm{ter}}$ Ordnung von $F(s_x,x)=0$ die Gleichung $a(s_x-s_x^0)+b(x-x^0)=0$ anwenden, oder
\[
s_x-s_x^0=-\frac{b}{a}(x-x^0)
\]
setzen, wodurch der Ausdruck
\[
a(s_x-s_x^0)+b(x-x^0)
\]
der $\frac32^{\mathrm{ten}}$ Potenz von $x-x^0$ proportional wird. Es ist also auch $\frac{\partial F}{\partial s_x}$, welches sich von diesem Ausdrucke nur durch den Factor $a$ unterscheidet, mit $(x-x^0)^{\frac32}$ proportional, und das Integral $V$ wird, so lange die Integration sich in der Nähe des Rückkehrpuncts bewegt, indem der Coefficient von $\frac{1}{x-x^0}$ unter dem Integralzeichen verschwindet, von der Form:
\[
\int\left(\frac{A}{\sqrt{(x-x^0)^3}}+\frac{B}{\sqrt{x-x^0}}\right)dx=
\frac{2}{3}\frac{A}{\sqrt{x-x^0}}+2B\sqrt{x-x^0},
\]
also algebraisch unendlich, und zwar der $-\frac12^{\mathrm{ten}}$ Potenz von $x-x^0$ proportional. Ausgenommen ist auch hier wieder der Fall, wo der unter dem Integralzeichen befindliche Zähler $Q$ selbst für den Rückkehrpunct verschwindet, wo dann $V$ mit
\[
\int\frac{dx}{\sqrt{x-x^0}}=2\sqrt{x-x^0}
\]
proportional wird, also endlich bleibt. Man sieht, dass sich der Rückkehrpunct verhält wie die Vereinigung eines Doppelpuncts mit einem Berührungspuncte, so dass nach Aufhebung des Doppelpuncts noch die Eigenschaften des Berührungspunctes übrig bleiben.

\begin{center}
{\bfseries \S\ 4. Integrale erster Gattung.}
\end{center}

Aus dem Gesagten erkennt man, dass es Integrale geben kann, welche überhaupt nie unendlich werden. Dieselben müssen die Eigenschaft haben, dass ihr Zähler $Q$ oder $\Theta$ für alle Doppel- und Rückkehrpuncte verschwindet. Dies ist, wenn $n>2$, in der That immer möglich; für $n=1$ oder $n=2$ existiren Integrale der hier betrachteten Art überhaupt nicht, da es keine ganze Function $Q$ $(n-3)^{\mathrm{ter}}$ Ordnung mehr giebt.

Die höchste Anzahl von Doppel- und Rückkehrpuncten, welche eine eigentliche Curve $n^{\mathrm{ter}}$ Ordnung besitzen kann, ist
\[
\frac{n-1\cdot n-2}{2}.
\]
Wäre nämlich ein Doppel- oder Rückkehrpunct mehr vorhanden, so könnte man durch die
\[
\frac{n-1\cdot n-2}{2}+1
\]
Puncte dieser Art und durch noch $2n-3$ andere auf der Curve gewählte Puncte, also im Ganzen durch
\[
\frac{(n-1)(n+2)}{2}
\]
Puncte, eine Curve $(n-1)^{\mathrm{ter}}$ Ordnung legen, welche dadurch gerade bestimmt ist. Diese Curve schneidet dann die gegebene, da jeder Doppel- oder Rückkehrpunct als zweifacher Schnittpunct zu rechnen ist, in
\[
(n-1)(n-2)+2+2n-3=n(n-1)+1
\]
Puncten, also in einem Schnittpuncte mehr, als einer Curve $n^{\mathrm{ter}}$ und einer Curve $(n-1)^{\mathrm{ter}}$ Ordnung überhaupt zukommt. Daher müsste ein Theil beider Curven zusammenfallen, mithin die Curve $n^{\mathrm{ter}}$ Ordnung selbst in Theile zerfallen und aufhören eine eigentliche Curve $n^{\mathrm{ter}}$ Ordnung zu sein.

Somit ist $d+r$ höchstens gleich
\[
\frac{n-1\cdot n-2}{2}.
\]
Die letztere Zahl aber stimmt genau überein mit der Zahl der Coefficienten in der Curve $(n-3)^{\mathrm{ter}}$ Ordnung $\Theta=0$, welche den Zähler unter dem Integralzeichen von $V$ bildet. Die Curve $\Theta=0$ kann also so bestimmt werden, dass sie durch die $d+r$ Doppel- und Rückkehrpuncte hindurchgeht, d. h. dass sie den $d+r$ für ihre Coefficienten linearen Gleichungen genügt, welche aussagen, dass $\Theta$ für die Coordinaten eines jener Puncte verschwinde. Für die höchste Zahl von Doppel- und Rückkehrpuncten erhält man hierbei freilich so viel homogene Gleichungen als Unbekannte, und diese müssen sämmtlich verschwinden; für
\[
d+r=\frac{n-1\cdot n-2}{2}
\]
ist also die Anzahl der überall endlichen Integrale gleich Null. Im Allgemeinen aber ist die Zahl der in $\Theta$ noch unbestimmt bleibenden Coefficienten
\[
p=\frac{n-1\cdot n-2}{2}-d-r,
\]
und das allgemeinste überall endliche Integral setzt sich daher aus ebenso vielen speciellen Integralen zusammen; oder wie man sich ausdrücken kann, es giebt
\[
p=\frac{n-1\cdot n-2}{2}-d-r
\]
verschiedene Integrale erster Gattung; wenn man mit letzterem Ausdruck überall endlich bleibende Integrale bezeichnet. Solche Integrale werden wir im Folgenden im Allgemeinen Integrale $I$ nennen.

Alle Integrale von der Form
\[
V=\int\frac{\Theta\Sigma\pm c_1x_2dx_3}{c_1\frac{\partial f}{\partial x_1}+c_2\frac{\partial f}{\partial x_2}+c_3\frac{\partial f}{\partial x_3}}
=\int\frac{Q\,dx}{\frac{\partial F}{\partial s_x}}
\]
kann man aus $\frac{n-1\cdot n-2}{2}$ speciellen Integralen zusammensetzen. Als solche specielle Integrale kann man ausser den eben behandelten überall endlichen Integralen solche wählen, welche nur in je einem der $d+r$ Ausnahmspuncte unendlich werden. In der That, sei $H=0$ eine beliebige Curve $(n-3)^{\mathrm{ter}}$ Ordnung, seien $\Theta_1=0$, $\Theta_2=0\ldots\Theta_{d+r}=0$ Curven $(n-3)^{\mathrm{ter}}$ Ordnung, welche durch die $d+r$ Ausnahmspuncte, immer einen ausgenommen, hindurchgehen, so kann man Constanten $\alpha_1,\alpha_2\ldots\alpha_{d+r}$ immer so bestimmen, dass die Curve
\[
(1)\ldots\quad H-\alpha_1\Theta_1-\alpha_2\Theta_2\cdots-\alpha_{d+r}\Theta_{d+r}=\Theta=0
\]
durch alle $d+r$ Ausnahmspuncte geht, und also
\[
\begin{aligned}
\int\frac{H\Sigma\pm c_1x_2dx_3}{c_1\frac{\partial f}{\partial x_1}+c_2\frac{\partial f}{\partial x_2}+c_3\frac{\partial f}{\partial x_3}}
={}&\int\frac{\Theta\Sigma\pm c_1x_2dx_3}{c_1\frac{\partial f}{\partial x_1}+c_2\frac{\partial f}{\partial x_2}+c_3\frac{\partial f}{\partial x_3}}\\
&+\Sigma_i\alpha_i\int\frac{\Theta_i\Sigma\pm c_1x_2dx_3}{c_1\frac{\partial f}{\partial x_1}+c_2\frac{\partial f}{\partial x_2}+c_3\frac{\partial f}{\partial x_3}}
\end{aligned}
\]
sich auf die angegebene Weise zusammensetzt. Bezeichnet man nämlich durch obere Indices die Werthe, welche die Functionen $H,\Theta_1,\Theta_2\ldots$ für die Ausnahmspuncte annehmen, so hat man, indem man die Bedingungen aufstellt, dass die Curve (1) durch alle Ausnahmspuncte hindurchgehe:
\[
H^{(1)}-\alpha_1\Theta_1^{(1)}=0,
\]
\[
H^{(2)}-\alpha_2\Theta_2^{(2)}=0,
\]
\[
\cdots
\]
\[
H^{(d+r)}-\alpha_{d+r}\Theta_{d+r}^{(d+r)}=0.
\]
Alle übrigen $\Theta_i^{(k)}$ verschwinden der Annahme nach; die vorstehenden Gleichungen aber bestimmen die $\alpha$ vollkommen und eindeutig. Die in Rede stehende Darstellung ist daher auch nur auf eine Weise möglich, wenn die Curven $\Theta_i$ einmal gewählt sind. Auf diese werden wir an einer späteren Stelle wieder zurückkommen.

\begin{center}
{\bfseries \S\ 5. Integrale dritter Gattung.}
\end{center}

Statt der in \S\ 2. unter 3). aufgeführten Integrale
\[
V=\int\frac{Q\,dx}{(x-a)\frac{\partial F}{\partial s_x}},
\]
wo $Q$ eine Function $(n-2)^{\mathrm{ten}}$ Grades ist, wollen wir die etwas allgemeinere Form betrachten
\[
V=\int\frac{Q\,dx}{(\alpha s_x+\beta x+\gamma)\frac{\partial F}{\partial s_x}}
=\int\frac{\Omega\Sigma\pm c_1x_2dx_3}{(\alpha_1x_1+\alpha_2x_2+\alpha_3x_3)\left(c_1\frac{\partial f}{\partial x_1}+c_2\frac{\partial f}{\partial x_2}+c_3\frac{\partial f}{\partial x_3}\right)}.
\]

% Continuation package Abel03: source pp. 039--055, with Toronto backfill for missing Google printed p. 30.
% --- Continuation: source PDF pp. 039--055. ---
% Scan-authorial text: printed pp. 17--33, end of §5 through end of Erster Abschnitt.

bei welcher im Nenner eine beliebige lineare Function erscheint. Das vorliegende Integral verhält sich bezüglich der Berührungspuncte der von $c$ an $f=0$ zu ziehenden Tangenten genau wie das vorige; auch wird es im Allgemeinen für die Doppel- und Rückkehrpuncte unstetig. Aber es wird auch noch für diejenigen $n$ Puncte unstetig, in denen die Curve $F=0$ von der Geraden
\[
\alpha s_x+\beta x+\gamma=0
\]
geschnitten wird. Ist $s_x^\circ$, $x^\circ$ einer dieser Puncte, so ist in der Nähe desselben
\[
\begin{aligned}
\alpha s_x+\beta x+\gamma&=\alpha(s_x-s_x^\circ)+\gamma(x-x^\circ),\\
F(s_x,x)=0&=\frac{\partial F^\circ}{\partial s_x^\circ}(s_x-s_x^\circ)+\frac{\partial F^\circ}{\partial x^\circ}(x-x^\circ)+\cdots .
\end{aligned}
\]
Aus der letzten Gleichung geht hervor, dass $s_x-s_x^\circ$ mit $x-x^\circ$ proportional ist; die erste zeigt sodann, dass der Nenner unter dem Integralzeichen von $V$ ebenfalls mit $x-x^\circ$ proportional, also $V$ logarithmisch unendlich wird. Die Schnittpuncte der Graden mit $F=0$ verhalten sich also in dieser Beziehung wie die Doppelpuncte der Curve.

Da die Function $Q$ oder $\Omega$ von der $(n-2)^{\mathrm{ten}}$ Ordnung ist, so enthält sie
\[
\frac{n(n-1)}{2}
\]
willkürliche Coefficienten, und man kann sie als Combination von ebenso vielen speciellen Functionen desselben Grades darstellen, oder, was dasselbe ist, das Integral $V$ als Aggregat von
\[
\frac{n(n-1)}{2}
\]
speciellen Integralen derselben Art. Nun kann man offenbar specielle Curven $\Omega=0$ so legen, dass sie durch $n-2$ Schnittpuncte der Geraden mit $f=0$ gehen. Mehr ist nicht möglich, wenn nicht $\Omega$ in die Gerade und eine Curve $(n-3)^{\mathrm{ter}}$ Ordnung zerfallen, $V$ also ein Integral erster Gattung werden soll; denn sollte eine Curve $\Omega=0$ der $(n-2)^{\mathrm{ten}}$ Ordnung durch $n-1$ dieser Schnittpuncte gehen, so müsste sie mit der Geraden $n-1$ Puncte gemein haben und also die Gerade selbst enthalten. Aber man kann eine solche specielle Curve ausserdem noch durch die $d+r$ Ausnahmspuncte von $f=0$ gehen lassen; denn da deren Zahl höchstens
\[
\frac{(n-1)(n-2)}{2}
\]
beträgt, so wird hienach die Zahl der für $\Omega=0$ gegebenen Puncte höchstens
\[
\frac{(n-1)(n-2)}{2}+n-2=\frac{(n+1)(n-2)}{2},
\]
welche Zahl gerade hinreicht, um eine Curve $(n-2)^{\mathrm{ter}}$ Ordnung zu bestimmen.

Seien also $\Omega_2=0$, $\Omega_3=0\ldots\Omega_n=0$ specielle Curven, welche sämmtlich die Ausnahmspuncte, und welche ausserdem alle Schnittpuncte der Geraden mit $f=0$ enthalten, beziehungsweise ausgenommen bei $\Omega_2=0$ den ersten und zweiten, bei $\Omega_3=0$ den ersten und dritten $\ldots$, bei $\Omega_n=0$ den ersten und $n^{\mathrm{ten}}$. Dann kann man in der Gleichung
\[
\Omega-k_2\Omega_2-k_3\Omega_3\cdots-k_n\Omega_n=0
\]
die Constanten $k$ so bestimmen, dass die hiedurch dargestellte Curve durch $n-1$ Schnittpuncte der Geraden mit $f=0$ geht, also die Gerade ganz enthält. Man hat nämlich, wenn man durch obere Indices bezeichnet, dass die Coordinaten des entsprechenden Schnittpuncts eingesetzt werden sollen, nur die Gleichungen
\[
\Omega^{(2)}-k_2\Omega_2^{(2)}=0,
\qquad
\Omega^{(3)}-k_3\Omega_3^{(3)}=0,
\]
\[
\cdots
\qquad
\Omega^{(n)}-k_n\Omega_n^{(n)}=0
\]
zu erfüllen, damit der $2^{\mathrm{te}}$, $3^{\mathrm{te}}\ldots$, $n^{\mathrm{te}}$ Schnittpunct auf der obigen Curve liege. Bestimmt man hieraus die $k$, so zerfällt die obige Curve in die Gerade und in eine Curve $(n-3)^{\mathrm{ter}}$ Ordnung $\Theta=0$; man hat also
\[
\Omega=k_2\Omega_2+k_3\Omega_3\cdots+k_n\Omega_n+(\alpha_1x_1+\alpha_2x_2+\alpha_3x_3)\Theta.
\]
Und somit erhält man für $V$ folgende Darstellung:
\[
\begin{aligned}
V={}&\sum_{i=2}^{i=n}k_i\int
\frac{\Omega_i\Sigma\pm c_1x_2dx_3}{(\alpha_1x_1+\alpha_2x_2+\alpha_3x_3)\left(c_1\frac{\partial f}{\partial x_1}+c_2\frac{\partial f}{\partial x_2}+c_3\frac{\partial f}{\partial x_3}\right)}\\
&+\int\frac{\Theta\Sigma\pm c_1x_2dx_3}{c_1\frac{\partial f}{\partial x_1}+c_2\frac{\partial f}{\partial x_2}+c_3\frac{\partial f}{\partial x_3}}.
\end{aligned}
\]
Diese Darstellung zeigt, dass das Integral $V$ sich auf die früher behandelten Integrale zurückführen lässt, und auf solche, welche nur in zwei Puncten logarithmisch unendlich werden. Denn in der That wird ein Integral
\[
\int\frac{\Omega_i\Sigma\pm c_1x_2dx_3}{(\alpha_1x_1+\alpha_2x_2+\alpha_3x_3)\left(c_1\frac{\partial f}{\partial x_1}+c_2\frac{\partial f}{\partial x_2}+c_3\frac{\partial f}{\partial x_3}\right)}
\]
nur noch in den beiden Schnittpuncten der Geraden mit der Curve $f=0$ unstetig, durch welche die Curve $\Omega_i=0$ nicht geht. Solche nur in zwei Puncten unstetige Integrale heissen Integrale dritter Gattung.

Auf diese Integrale lassen sich nicht nur, nebst Integralen der im vorigen \S\ behandelten Form, alle unter \S\ 3 angeführten Integrale zurückführen, sondern auch diejenigen im Vorigen behandelten Integrale
\[
\int\frac{\Theta_i\Sigma\pm c_1x_2dx_3}{c_1\frac{\partial f}{\partial x_1}+c_2\frac{\partial f}{\partial x_2}+c_3\frac{\partial f}{\partial x_3}}
\]
selbst, welche nur in einem Doppelpuncte unendlich wurden. Denn lässt man in dem Integral
\[
\int\frac{\Omega\Sigma\pm c_1x_2dx_3}{(\alpha_1x_1+\alpha_2x_2+\alpha_3x_3)\left(c_1\frac{\partial f}{\partial x_1}+c_2\frac{\partial f}{\partial x_2}+c_3\frac{\partial f}{\partial x_3}\right)}
\]
die Gerade $\alpha$ durch einen Doppelpunct gehen, und bestimmt $\Omega$, wie im Frühern gelehrt wurde, so dass es durch alle Doppelpuncte (und Rückkehrpuncte) und ausserdem durch die übrigen $n-2$ Puncte geht, welche die Gerade $\alpha$ mit der Curve $f=0$ gemein hat, so hat man offenbar ein Integral von dem Character der oben behandelten vor sich. Aber $\Omega$ hat dann mit der Geraden $\alpha$ $n-1$ Puncte gemein und enthält diese demnach ganz; daher geht das obige Integral, indem
\[
\Omega=(\alpha_1x_1+\alpha_2x_2+\alpha_3x_3)\Theta
\]
wird, in die Form
\[
\int\frac{\Theta\Sigma\pm c_1x_2dx_3}{c_1\frac{\partial f}{\partial x_1}+c_2\frac{\partial f}{\partial x_2}+c_3\frac{\partial f}{\partial x_3}}
\]
über, wo nun $\Theta$ durch alle Ausnahmspuncte geht, bis auf den auf $\alpha$ liegenden, so dass dieses Integral zugleich das Schema der ad 2) \S\ 2 angeführten nicht überall endlichen Integrale liefert, soweit dieselben in Doppelpuncten unstetig werden. Auch diese also rechnen wir zu den Integralen dritter Gattung. Und da diese immer in zwei Puncten unstetig sein müssen, so kann man sagen, jene seien unstetig in den beiden Puncten, welche in dem betrachteten Doppelpuncte vereinigt liegen.

Wie diese Betrachtungen zu modificiren sind, wenn der Unstetigkeitspunct ein Rückkehrpunct wird, ist weiter unten zu erörtern.

Da sowohl die Integrale dritter Gattung, als auch deren Differentialquotienten nach Constanten, als endlich die oben abgesonderten algebraischen oder logarithmischen Functionen, überhaupt also alle hier auftretenden Integrale in irgend welchen Puncten unstetig werden, die Integrale erster Gattung allein ausgenommen, so folgt als Corollar der Satz:

Jedes Integral der vorliegenden Form, welches für keinen Punct von $f=0$ unstetig wird, ist ein Integral erster Gattung.

\begin{center}
{\bfseries \S\ 6. Bildung der Normalintegrale dritter Gattung.}
\end{center}

Ueber die Bildung eines Integrals $3^{\mathrm{ter}}$ Gattung können wir nun Folgendes genauer ausführen.

Es seien $\xi$, $\eta$ die beiden Unstetigkeitspuncte des Integrals, so dass, da beide auf $f=0$ liegen, die Gleichungen stattfinden:
\[
f(\xi_1,\xi_2,\xi_3)=0,
\qquad
f(\eta_1,\eta_2,\eta_3)=0.
\]
Dann wird die Gleichung der Geraden $\alpha$ folgende:
\[
\Sigma\pm x_1\xi_2\eta_3=0,
\]
und das Integral wird:
\[
V=\int\frac{\Omega\Sigma\pm c_1x_2dx_3}{\Sigma\pm x_1\xi_2\eta_3\left(c_1\frac{\partial f}{\partial x_1}+c_2\frac{\partial f}{\partial x_2}+c_3\frac{\partial f}{\partial x_3}\right)}.
\]
Die Curve $(n-2)^{\mathrm{ter}}$ Ordnung $\Omega=0$ ist nun so zu bestimmen, dass sie durch die sämmtlichen Ausnahmspuncte von $f=0$, und durch die $n-2$ Puncte geht, in denen die Gerade $\xi\eta$ die Curve $f=0$ noch schneidet. Die Anzahl aller für $\Omega=0$ festgelegten Puncte ist also $d+r+n-2$, so dass man noch
\[
\frac{(n-1)(n-2)}{2}-d-r=p
\]
Puncte beliebig hinzufügen kann, durch welche $\Omega$ zu gehen genöthigt sein soll.

Die Coordinaten eines Punctes der Geraden $\xi$, $\eta$ sind bekanntlich
\[
\xi_1+\lambda\eta_1,
\qquad
\xi_2+\lambda\eta_2,
\qquad
\xi_3+\lambda\eta_3,
\]
wo $\lambda$ ein variabler Parameter ist.

Diejenigen Werthe nun von $\lambda$, welche den $n-2$ andern Schnittpuncten der Geraden $\xi\eta$ mit $f=0$ entsprechen, erhält man, indem man diese Ausdrücke der Coordinaten in $f=0$ einsetzt. Diese Gleichung nimmt dann, nach Potenzen von $\lambda$ geordnet, die Form an:
\[
u_0+\lambda u_1+\lambda^2u_2\cdots+\lambda^nu_n=0,
\]
und giebt also eine Gleichung $n^{\mathrm{ten}}$ Grades für $\lambda$ zur Bestimmung der $n$ Schnittpuncte unserer Geraden mit $f=0$. Aber zwei Wurzeln kennen wir bereits: $\lambda=0$, welche den Punct $\xi$, und $\lambda=\infty$, welche den Punct $\eta$ liefert. Man muss daher in der obigen Gleichung $u_0=0$, $u_n=0$ haben; und in der That ist ja
\[
u_0=f(\xi_1,\xi_2,\xi_3),
\qquad
u_n=f(\eta_1,\eta_2,\eta_3).
\]
Daher geht die obige Gleichung nach Auslassung des Factors $\lambda$ in die Gleichung $(n-2)^{\mathrm{ten}}$ Grades
\[
(1)\ldots\quad u_1+\lambda u_2+\lambda^2u_3\cdots+\lambda^{n-2}u_{n-1}=0
\]
über, welche zur Bestimmung der übrigen $n-2$ Schnittpuncte dient.

Diese Schnittpuncte sollen zugleich auf der Curve $\Omega=0$ liegen. Dieses heisst, dass, wenn man in $\Omega$ statt der Coordinaten die Grössen $\xi_i+\lambda\eta_i$ einsetzt, für $\lambda$ dieselbe Gleichung resultiren muss, welche wir soeben erhalten haben. Setzt man also
\[
\Omega(\xi_1+\lambda\eta_1,\xi_2+\lambda\eta_2,\xi_3+\lambda\eta_3)
=v_1+\lambda v_2+\lambda^2v_3\cdots+\lambda^{n-2}v_{n-1},
\]
so können sich die $v$ von den entsprechenden Grössen $u$ nur durch einen gemeinschaftlichen Factor unterscheiden. Man hat also die Gleichungen:
\[
(2)\ldots\quad
\begin{cases}
 v_1=\mu u_1,\\
 v_2=\mu u_2,\\
 \cdots\\
 v_{n-1}=\mu u_{n-1},
\end{cases}
\]
wobei $\mu$ diesen unbestimmten gemeinschaftlichen Factor bedeutet.

Die Gleichungen (2) sind $n-1$ an der Zahl, linear für die in den $v$ auftretenden Coefficienten von $\Omega$, und für die Grösse $\mu$, welche übrigens unbeschadet der Allgemeinheit gleich 1 gesetzt werden kann, was im Folgenden geschehen soll. Fügen wir die
\[
\frac{(n-1)(n-2)}{2}
\]
andern Gleichungen hinzu, welche aussagen, dass $\Omega=0$ durch ebensoviel gegebene Puncte geht, so hat man
\[
\frac{n(n-1)}{2}
\]
lineare Gleichungen mit eben so vielen Unbekannten, den Coefficienten von $\Omega$, welche durch dieselben vollständig bestimmt sind. Trägt man die erhaltenen Werthe ein, so wird $\Omega$ der negative Quotient zweier Determinanten. Die Determinante des Zählers enthält in der ersten Vertikalreihe die Elemente
\[
0, u_1, u_2\ldots u_{n-1}, 0,0\ldots0;
\]
in den folgenden Vertikalreihen aber die in die Coefficienten von $\Omega$ multiplicirten Glieder der Functionen $\Omega$, $v_1$, $v_2\ldots v_{n-1}$, so wie derjenigen, welche die linken Theile der übrigen linearen Gleichungen bilden. Die Determinante des Nenners ist die Determinante der linearen Gleichungen selbst, oder die Unterdeterminante der vorigen genommen nach ihrem ersten Element\footnote{Dieser Nenner hat übrigens eine sehr bemerkenswerthe Eigenschaft. Er zerfällt in das Product von $\frac{(n-1)(n-2)}{2}$ dreigliedrigen Determinanten mit einer grösseren Determinante. Eine der ersten gleich Null gesetzt sagt aus, dass einer der $\frac{(n-1)(n-2)}{2}$ festen Puncte auf der Geraden $\xi\eta$ liege; das Verschwinden der letzten giebt an, dass jene Puncte selbst auf einer Curve $(n-3)^{\mathrm{ter}}$ Ordnung liegen, ein Fall, welcher immer zu vermeiden ist.}.

Diese Darstellungsweise der Function $\Omega$ wollen wir im Folgenden zu Grunde legen, und das so gebildete Integral dritter Gattung durch $S$ oder $S_{\xi\eta}$ oder $S_{\xi\eta}(x)$ bezeichnen. Dasselbe hat einige wichtige Eigenschaften, welche hier entwickelt werden sollen.

\begin{center}
{\bfseries \S.\ 7. Eigenschaften der Normalintegrale dritter Gattung.}
\end{center}

Betrachten wir zunächst das Verhalten von $\Omega$ und $S$ in der Nähe der Unstetigkeitspuncte. In der Nähe des Punctes $\xi$ kann man die $x$ überall da, wo es sich nicht um verschwindend kleine Factoren handelt, durch die $\xi$ ersetzen. Nun sind die ersten beiden Horizontalreihen des Zählers von $\Omega$, die erste von $\Omega$ selbst, die zweite von der ersten Gleichung (2) herrührend, folgende:
\[
\begin{array}{cccccc}
0 & x_1^{n-2} & nx_1^{n-3}x_2 & \cdots & x_3^{n-2},\\
u_1 & \xi_1^{n-2} & n\xi_1^{n-3}\xi_2 & \cdots & \xi_3^{n-2}.
\end{array}
\]
Wenn $x$ in die Nähe von $\xi$ rückt, so unterscheiden sich diese Reihen nur noch durch das erste Glied; zieht man die zweite Reihe von der ersten ab, so wird das erste Glied der ersten Reihe $-u_1$, die übrigen verschwinden; der Zähler von $\Omega$ reducirt sich also auf $-u_1$ multiplicirt mit der Unterdeterminante des ersten Gliedes, und da diese genau dem Nenner von $\Omega$ gleich ist, so erhält man
\[
\Omega=u_1,
\]
oder wenn man die Entstehung von $u_1$ als des Coefficienten von $\lambda$ in der Entwicklung von
\[
f(\xi_1+\lambda\eta_1,\xi_2+\lambda\eta_2,\xi_3+\lambda\eta_3)
\]
berücksichtigt:
\[
\Omega=\eta_1\frac{\partial f(\xi)}{\partial \xi_1}+\eta_2\frac{\partial f(\xi)}{\partial \xi_2}+\zeta_3\frac{\partial f(\xi)}{\partial \xi_3}.
\]
So lange also die Integration sich in der Nähe des Punctes $\xi$ bewegt, kann man für $S$ setzen:
\[
S=\frac{\eta_1\frac{\partial f(\xi)}{\partial \xi_1}+\eta_2\frac{\partial f(\xi)}{\partial \xi_2}+\zeta_3\frac{\partial f(\xi)}{\partial \xi_3}}{c_1\frac{\partial f(\xi)}{\partial \xi_1}+c_2\frac{\partial f(\xi)}{\partial \xi_2}+c_3\frac{\partial f(\xi)}{\partial \xi_3}}
\cdot
\int\frac{\Sigma\pm c_1\xi_2dx_3}{\Sigma\pm x_1\xi_2\eta_3}.
\]
Setzen wir hier für die ganz willkürlichen Grössen $c$ die $\eta$ ein, so findet sich sofort
\[
S=-\int\frac{\Sigma\pm dx_1\xi_2\eta_3}{\Sigma\pm x_1\xi_2\eta_3}=-\log\Sigma\pm x_1\xi_2\eta_3.
\]
Das Normalintegral $S$ verhält sich also in der Nähe des Unstetigkeitspuncts $\xi$ wie $-\log\Sigma\pm x_1\xi_2\eta_3$; mithin in der Nähe des Unstetigkeitspuncts $\eta$ wie $+\log\Sigma\pm x_1\xi_2\eta_3$\footnote{Dies ändert sich nicht, auch wenn ein Unstetigkeitspunct in den Berührungspunct einer der von $c$ an die Curve zu ziehenden Tangenten fällt. Ueberhaupt ist ein solches Zusammentreffen immer als vollkommen gleichgültig zu betrachten.}.

Nehmen wir ferner drei verschiedene Puncte auf der Curve $f=0$ an, die wir durch $\xi$, $\eta$, $\zeta$ bezeichnen, und bilden die Summe
\[
S_{\xi\eta}+S_{\eta\zeta}+S_{\zeta\xi}.
\]
Unter dem Integralzeichen befindet sich dann, von einem gemeinschaftlichen Factor abgesehen, eine Summe von der Form:
\[
\frac{\Omega}{\Sigma\pm x_1\xi_2\eta_3}
+\frac{\Omega'}{\Sigma\pm x_1\eta_2\zeta_3}
+\frac{\Omega''}{\Sigma\pm x_1\zeta_2\xi_3},
\]
wobei die $\Omega'$, $\Omega''$ ähnliche Bildungen für die Punctepaare $\eta\zeta$, $\zeta\xi$ bedeuten, wie $\Omega$ für das Punctepaar $\xi\eta$. Bringen wir diesen Ausdruck auf gleichen Nenner, so ist der Zähler die Function $n^{\mathrm{ter}}$ Ordnung:
\[
\begin{aligned}
\varphi={}&\Omega\,\Sigma\pm x_1\eta_2\zeta_3\cdot\Sigma\pm x_1\zeta_2\xi_3\\
&+\Omega'\,\Sigma\pm x_1\zeta_2\xi_3\cdot\Sigma\pm x_1\xi_2\eta_3\\
&+\Omega''\,\Sigma\pm x_1\xi_2\eta_3\cdot\Sigma\pm x_1\eta_2\zeta_3.
\end{aligned}
\]
Dies gleich Null gesetzt, ist die Gleichung einer Curve $n^{\mathrm{ter}}$ Ordnung. Dieselbe geht durch alle Schnittpuncte von $f=0$ mit den Geraden $\xi\eta$, $\eta\zeta$, $\zeta\xi$. Denn was z. B. die Gerade $\xi\eta$ angeht, so ist der linke Theil ihrer Gleichung in zweien der obigen Terme geradezu enthalten. In dem ersten aber verschwindet $\Omega$ für die $n-2$ Schnittpuncte, welche die Gerade ausser $\xi$ und $\eta$ mit $f=0$ gemein hat, während $\Sigma\pm x\eta\zeta$ für $\eta$, $\Sigma\pm x\zeta\xi$ für $\xi$ verschwindet. Diese Eigenschaft wird auch nicht verändert, wenn man statt $\varphi$ den Ausdruck $\varphi+kf$ betrachtet, wo $k$ eine Constante ist. Nun kann man diese Constante immer so wählen, dass die Curve $\varphi+kf=0$ die drei Puncte $\xi$, $\eta$, $\zeta$ zu Doppelpuncten hat. Stellen wir für den Punct $\xi$ die hierin enthaltenen Bedingungen auf. Sie bestehen darin, dass die Differentialquotienten von $\varphi+kf$ nach den $x$ für $x=\xi$ verschwinden. Nun erhält man aber für $x=\xi$, indem man die Glieder auslässt, deren Verschwinden für $x=\xi$ man unmittelbar einsieht:
\[
\begin{aligned}
\frac{\partial(\varphi+kf)}{\partial x_1}
={}&\bigl[\Omega(\xi_2\zeta_3-\zeta_2\xi_3)+\Omega''(\xi_2\eta_3-\eta_2\xi_3)\bigr]\Sigma\pm \xi_1\eta_2\zeta_3+k\frac{\partial f(\xi)}{\partial \xi_1},\\
\frac{\partial(\varphi+kf)}{\partial x_2}
={}&\bigl[\Omega(\zeta_3\xi_1-\xi_3\zeta_1)+\Omega''(\eta_3\xi_1-\xi_3\eta_1)\bigr]\Sigma\pm \xi_1\eta_2\zeta_3+k\frac{\partial f(\xi)}{\partial \xi_2},\\
\frac{\partial(\varphi+kf)}{\partial x_3}
={}&\bigl[\Omega(\xi_1\zeta_2-\zeta_1\xi_2)+\Omega''(\xi_1\eta_2-\eta_1\xi_2)\bigr]\Sigma\pm \xi_1\eta_2\zeta_3+k\frac{\partial f(\xi)}{\partial \xi_3}.
\end{aligned}
\]
In diesen Ausdrücken wird nach dem Vorigen
\[
\Omega=-\left(\eta_1\frac{\partial f(\xi)}{\partial \xi_1}+\eta_2\frac{\partial f(\xi)}{\partial \xi_2}+\zeta_3\frac{\partial f(\xi)}{\partial \xi_3}\right),
\]
\[
\Omega''=\left(\zeta_1\frac{\partial f(\xi)}{\partial \xi_1}+\zeta_2\frac{\partial f(\xi)}{\partial \xi_2}+\zeta_3\frac{\partial f(\xi)}{\partial \xi_3}\right).
\]
Bemerkt man noch, dass
\[
0=\xi_1\frac{\partial f(\xi)}{\partial \xi_1}+\xi_2\frac{\partial f(\xi)}{\partial \xi_2}+\xi_3\frac{\partial f(\xi)}{\partial \xi_3}=nf(\xi),
\]
so kann man den in der ersten der obigen Gleichungen eingeklammerten Ausdruck auch so schreiben:
\[
\left|
\begin{array}{ccc}
\xi_1\frac{\partial f(\xi)}{\partial \xi_1}+\xi_2\frac{\partial f(\xi)}{\partial \xi_2}+\xi_3\frac{\partial f(\xi)}{\partial \xi_3} & \xi_2 & \xi_3\\
\eta_1\frac{\partial f(\xi)}{\partial \xi_1}+\eta_2\frac{\partial f(\xi)}{\partial \xi_2}+\zeta_3\frac{\partial f(\xi)}{\partial \xi_3} & \eta_2 & \eta_3\\
\zeta_1\frac{\partial f(\xi)}{\partial \xi_1}+\zeta_2\frac{\partial f(\xi)}{\partial \xi_2}+\zeta_3\frac{\partial f(\xi)}{\partial \xi_3} & \zeta_2 & \zeta_3
\end{array}
\right|
=\frac{\partial f(\xi)}{\partial \xi_1}\Sigma\pm \xi_1\eta_2\zeta_3.
\]
Indem man die andern Ausdrücke ähnlich behandelt, findet man endlich
\[
\frac{\partial(\varphi+kf)}{\partial x_i}=\frac{\partial f(\xi)}{\partial \xi_i}\left[k+(\Sigma\pm\xi_1\eta_2\zeta_3)^2\right].
\]
Es wird daher $\xi$ ein Doppelpunct der Curve $\varphi+kf=0$, wenn
\[
k=-(\Sigma\pm\xi_1\eta_2\zeta_3)^2
\]
gesetzt wird. Aber wenn man in gleicher Weise $k$ so bestimmt, dass $\eta$ oder $\zeta$ Doppelpunct wird, so erhält man offenbar stets den nämlichen Werth von $k$, da dieser durch Vertauschung von $\xi$, $\eta$, $\zeta$ sich nicht ändert. Daher hat die Curve $\varphi+kf=0$ alle drei Puncte zu Doppelpuncten. Sie hat deswegen mit jeder der Linien $\xi\eta$, $\eta\zeta$, $\zeta\xi$ nicht blos die $n-2$ übrigen Schnittpuncte der Geraden mit $f=0$, sondern auch noch zwei Doppelpuncte, also im Ganzen $n+2$ Puncte gemein, muss also diese Geraden ganz enthalten. Daher ist
\[
\varphi+kf=\Sigma\pm x_1\xi_2\eta_3\cdot\Sigma\pm x_1\eta_2\zeta_3\cdot\Sigma\pm x_1\zeta_2\xi_3\cdot\Theta,
\]
wo $\Theta$ ein Ausdruck $(n-3)^{\mathrm{ter}}$ Ordnung ist.

Es wird demnach, mit Berücksichtigung der während der Integration stattfindenden Gleichung $f=0$:
\[
\frac{\Omega}{\Sigma\pm x_1\xi_2\eta_3}
+\frac{\Omega'}{\Sigma\pm x_1\eta_2\zeta_3}
+\frac{\Omega''}{\Sigma\pm x_1\zeta_2\xi_3}=\Theta,
\]
und demnach
\[
S_{\xi\eta}+S_{\eta\zeta}+S_{\zeta\xi}
=\int\frac{\Theta\Sigma\pm c_1x_2dx_3}{c_1\frac{\partial f}{\partial x_1}+c_2\frac{\partial f}{\partial x_2}+c_3\frac{\partial f}{\partial x_3}}.
\]
Da der Voraussetzung nach die $\Omega$ für die Doppel- und Rückkehrpuncte verschwinden, so gilt dasselbe für $\Theta$; die rechte Seite dieser Gleichung ist also ein Integral erster Gattung, und man hat den Satz:

Die Summe dreier Normalintegrale dritter Gattung, gebildet für je zwei von drei Puncten der Curve $f=0$ als Unstetigkeitspuncte, ist ein Integral erster Gattung.

Diese Gleichung, welche wir kürzer in der Form
\[
(1)\ldots\quad S_{\xi\eta}+S_{\eta\zeta}+S_{\zeta\xi}=I
\]
schreiben können, ist für die Theorie der Integrale dritter Gattung von grosser Wichtigkeit. Setzen wir nämlich für $\zeta$ irgend einen beliebig gewählten constanten Punct der Curve $f=0$, und bemerken, dass wegen der Bildungsweise der Function $S$
\[
S_{\eta\xi}=-S_{\xi\eta},
\qquad
S_{\xi\zeta}=-S_{\zeta\xi},
\]
so hat man aus (1):
\[
S_{\xi\eta}=S_{\xi\zeta}-S_{\eta\zeta}+I,
\]
was man in folgendem Satze aussprechen kann:

Das Normalintegral dritter Gattung $S_{\xi\eta}$ ist, bis auf ein Glied, welches einem Integrale erster Gattung gleich wird, die Differenz zweier ganz gleich gebildeter Functionen, deren eine nur von den $\xi$, deren andre nur von den $\eta$ abhängt, und welche sich nur durch diese Parameter unterscheiden.

Auf diese Weise sind die von zwei Puncten $\xi$, $\eta$ abhängigen Integrale $S$ auf Functionen zurückgeführt, welche nur von einem der Unstetigkeitspuncte abhängen.

Ein weiterer Satz bezieht sich auf den Umstand, dass $\Omega$ und daher auch $S$ nicht vollkommen bestimmt waren, sondern dass bei der Bildung von $\Omega$ eine Reihe von ganz willkürlichen Puncten benutzt war. Indem wir sie durch andre ersetzen, erhalten wir eine Function $\Omega'$ und ein Integral $S'$, welche in allen wesentlichen Puncten mit $\Omega$ und $S$ übereinstimmen. Die Curve $(n-2)^{\mathrm{ter}}$ Ordnung
\[
\Omega-\Omega'=0
\]
geht durch die $n-2$ übrigen Schnittpuncte der Geraden $\xi\eta$ mit $f=0$, da für diese Puncte $\Omega$ und $\Omega'$ einzeln verschwinden. Aber da $\Omega$ sich für $x=\xi$ nach dem Vorigen auf einen ganz bestimmten Werth reducirt, der von den willkürlich zu wählenden Puncten vollkommen unabhängig ist, so nimmt $\Omega'$ für $x=\xi$ denselben Werth an, und $\Omega-\Omega'$ verschwindet für $x=\xi$, und ebenso für $x=\eta$. Die Curve $\Omega-\Omega'=0$ hat also mit der Geraden $\xi\eta$ $n$ Puncte gemein und enthält demnach dieselbe ganz, so dass
\[
\Omega-\Omega'=\Theta\cdot\Sigma\pm x_1\xi_2\eta_3,
\]
wo $\Theta=0$ eine Curve $(n-3)^{\mathrm{ter}}$ Ordnung ist, welche wie $\Omega=0$, $\Omega'=0$ durch die Doppel- und Rückkehrpuncte von $f=0$ hindurchgeht. Daher ist
\[
S-S'=\int\frac{\Theta\Sigma\pm c_1x_2dx_3}{c_1\frac{\partial f}{\partial x_1}+c_2\frac{\partial f}{\partial x_2}+c_3\frac{\partial f}{\partial x_3}}=I,
\]
d. h.:

Wenn man das Normalintegral $S$ auf zwei verschiedene Arten für dieselben Unstetigkeitspuncte $\xi$, $\eta$ bildet, so unterscheiden sich beide Bildungen nur um ein Integral erster Gattung.

Für diejenige Form der Integrale dritter Gattung, in welchen $x$, $s_x$ eingeführt, und $s_x$ als Function von $x$ betrachtet wird, ersetzen wir ebenso die Puncte $\xi$, $\eta$ durch Werthpaare $s_\xi$, $\xi$, $s_\eta$, $\eta$, welche den Gleichungen 3. \S\ 1. gemäss mit den Coordinaten jener Puncte zusammenhängen. Es nähert sich dann das Integral dritter Gattung dem Obigen zufolge in der Nähe von $\xi$ dem Ausdruck $-\log\Sigma\pm x_1\xi_2\eta_3$, was hier zunächst in
\[
-\log\left|
\begin{array}{ccc}
1&x&s_x\\
1&\xi&s_\xi\\
1&\eta&s_\eta
\end{array}
\right|
=-\log\bigl[(x-\xi)(s_\xi-s_\eta)-(s_x-s_\xi)(\xi-\eta)\bigr]
\]
übergeht. Wenn man nun bemerkt, dass in der Nähe von $\xi$
\[
0=(x-\xi)\frac{\partial F(\xi,s_\xi)}{\partial \xi}+(s_x-s_\xi)\frac{\partial F(\xi,s_\xi)}{\partial s_\xi},
\]
so wird obiger Ausdruck bis auf eine additive Constante zu $-\log(x-\xi)$, und es verhält sich also das Integral in der Nähe dieses Unstetigkeitspuncts wie der Logarithmus der Differenz der Variabeln und des Parameters.

\begin{center}
{\bfseries \S\ 8. Integrale zweiter Gattung.}
\end{center}

Ein Integral zweiter Gattung entsteht, wenn die beiden Unstetigkeitspuncte eines Integrals dritter Gattung unendlich nahe zusammenrücken, und also die Gerade $\xi\eta$ zur Tangente der Curve $f=0$ wird. Man kann ein solches Integral direct bilden in der Form
\[
\int
\frac{\Omega\Sigma\pm c_1x_2dx_3}{\left(x_1\frac{\partial f(\xi)}{\partial \xi_1}+x_2\frac{\partial f(\xi)}{\partial \xi_2}+x_3\frac{\partial f(\xi)}{\partial \xi_3}\right)
\left(c_1\frac{\partial f}{\partial x_1}+c_2\frac{\partial f}{\partial x_2}+c_3\frac{\partial f}{\partial x_3}\right)},
\]
wo nun im Nenner ein Ausdruck erscheint, welcher gleich Null gesetzt die Gleichung der Tangente des Punctes $\xi$ liefert. Man hat dann die Function $(n-2)^{\mathrm{ter}}$ Ordnung $\Omega$ so zu bestimmen, dass $\Omega=0$ durch die Doppel- und Rückkehrpuncte von $f=0$ und durch die $n-2$ Puncte geht, in denen die Tangente des Punctes $\xi$ die Curve $f=0$ noch schneidet.

Man erkennt hieraus sofort, dass ein nur in einem Rückkehrpunct unendliches Integral ein Integral zweiter Gattung ist, denn jede durch den Rückkehrpunct gelegte Gerade verbindet in diesem zwei unendlich nahe Puncte, ungleich einer durch einen Doppelpunct gelegten Geraden, welche als Verbindungslinie zweier ganz verschiedener Puncte zu betrachten ist. Dieselbe Betrachtung, welche oben das in einem Doppelpunct unendliche Integral als Integral dritter Gattung erkennen lehrte, kennzeichnet also das in einem Rückkehrpunct unendliche, welches übrigens ebenfalls die Form eines Integrals erster Gattung hat, als in zwei unendlich nahen Puncten unstetig, d. h. als Integral zweiter Gattung.

Aber die oben angedeutete Darstellung eines Integrals zweiter Gattung führt im Allgemeinen auf erhebliche algebraische Schwierigkeiten. Um diese zu vermeiden, kann man auch folgenden Weg einschlagen. Wenn man in der Gleichung
\[
S_{\xi\eta}=S_{\xi\zeta}-S_{\eta\zeta}+I
\]
die Puncte $\xi$, $\eta$ einander sich nähern lässt, wenn man also etwa setzt:
\[
\eta_1=\xi_1-\varepsilon\alpha_1,
\qquad
\eta_2=\xi_2-\varepsilon\alpha_2,
\qquad
\eta_3=\xi_3-\varepsilon\alpha_3,
\]
und $\varepsilon$ sich der Null nähert, so erhält man:
\[
\lim S_{\xi\eta}=\varepsilon\left(\alpha_1\frac{\partial S_{\xi\zeta}}{\partial \xi_1}+\alpha_2\frac{\partial S_{\xi\zeta}}{\partial \xi_2}+\alpha_3\frac{\partial S_{\xi\zeta}}{\partial \xi_3}\right)+\lim I.
\]
Da nun $I$ die $\xi$, $\eta$, $\zeta$ nur in den Coefficienten enthält, so bleibt die Grenze, der sich dasselbe nähert, immer noch ein Integral erster Gattung; indem wir also davon abstrahiren, können wir als Integral zweiter Gattung den Ausdruck
\[
E_\xi=\alpha_1\frac{\partial S_{\xi\zeta}}{\partial \xi_1}+\alpha_2\frac{\partial S_{\xi\zeta}}{\partial \xi_2}+\alpha_3\frac{\partial S_{\xi\zeta}}{\partial \xi_3}
\]
zu Grunde legen. Weil aber auch $S_{\xi\eta}$ gleich der Differenz einer Function von $\xi$ und derselben Function von $\eta$ ist, abgesehen von einem Integral erster Gattung, so kann auch obiger Ausdruck $\zeta$ nur noch in einem Gliede enthalten, welches sich aus Integralen erster Gattung zusammensetzt, oder man kann ein Integral $I$ so bestimmen, dass
\[
\alpha_1\frac{\partial S_{\xi\zeta}}{\partial \xi_1}+\alpha_2\frac{\partial S_{\xi\zeta}}{\partial \xi_2}+\alpha_3\frac{\partial S_{\xi\zeta}}{\partial \xi_3}-I
\]
von $\zeta$ unabhängig wird.

Die Grössen $\alpha$ sind nicht vollkommen willkürlich, sondern müssen so gewählt werden, dass $\eta$ auf der Curve bleibt, oder dass
\[
0=f(\eta)=f(\xi)-\varepsilon\left(\alpha_1\frac{\partial f}{\partial \xi_1}+\alpha_2\frac{\partial f}{\partial \xi_2}+\alpha_3\frac{\partial f}{\partial \xi_3}\right)+\cdots .
\]
Da nun das erste Glied ohnedies verschwindet, so bleibt die Bedingung
\[
\alpha_1\frac{\partial f}{\partial \xi_1}+\alpha_2\frac{\partial f}{\partial \xi_2}+\alpha_3\frac{\partial f}{\partial \xi_3}=0.
\]
Man kann dieselbe auf die allgemeinste Weise erfüllen, indem man setzt:
\[
\alpha_1=\beta_3\frac{\partial f}{\partial \xi_2}-\beta_2\frac{\partial f}{\partial \xi_3},
\quad
\alpha_2=\beta_1\frac{\partial f}{\partial \xi_3}-\beta_3\frac{\partial f}{\partial \xi_1},
\quad
\alpha_3=\beta_2\frac{\partial f}{\partial \xi_1}-\beta_1\frac{\partial f}{\partial \xi_2},
\]
wo denn die $\beta$ ganz beliebig sind.

Setzt man dies ein, so wird
\[
(1)\ldots\quad
\alpha_1\frac{\partial S_{\xi\zeta}}{\partial \xi_1}+\alpha_2\frac{\partial S_{\xi\zeta}}{\partial \xi_2}+\alpha_3\frac{\partial S_{\xi\zeta}}{\partial \xi_3}
=\left|
\begin{array}{ccc}
\frac{\partial S_{\xi\zeta}}{\partial \xi_1}&\frac{\partial f}{\partial \xi_1}&\beta_1\\
\frac{\partial S_{\xi\zeta}}{\partial \xi_2}&\frac{\partial f}{\partial \xi_2}&\beta_2\\
\frac{\partial S_{\xi\zeta}}{\partial \xi_3}&\frac{\partial f}{\partial \xi_3}&\beta_3
\end{array}
\right|.
\]
Es ist aber $S_{\xi\zeta}$ eine Function, welche nur von den Verhältnissen der $\xi$ abhängt, wie eine Abzählung lehrt; daher hat man identisch:
\[
\xi_1\frac{\partial S_{\xi\zeta}}{\partial \xi_1}+\xi_2\frac{\partial S_{\xi\zeta}}{\partial \xi_2}+\xi_3\frac{\partial S_{\xi\zeta}}{\partial \xi_3}=0.
\]
Ausserdem ist:
\[
\xi_1\frac{\partial f}{\partial \xi_1}+\xi_2\frac{\partial f}{\partial \xi_2}+\xi_3\frac{\partial f}{\partial \xi_3}=0.
\]
Aus beiden Gleichungen geht hervor:
\[
\begin{aligned}
\frac{\partial S_{\xi\zeta}}{\partial \xi_2}\frac{\partial f}{\partial \xi_3}
-\frac{\partial S_{\xi\zeta}}{\partial \xi_3}\frac{\partial f}{\partial \xi_2}&=M\cdot\xi_1,\\
\frac{\partial S_{\xi\zeta}}{\partial \xi_3}\frac{\partial f}{\partial \xi_1}
-\frac{\partial S_{\xi\zeta}}{\partial \xi_1}\frac{\partial f}{\partial \xi_3}&=M\cdot\xi_2,\\
\frac{\partial S_{\xi\zeta}}{\partial \xi_1}\frac{\partial f}{\partial \xi_2}
-\frac{\partial S_{\xi\zeta}}{\partial \xi_2}\frac{\partial f}{\partial \xi_1}&=M\cdot\xi_3,
\end{aligned}
\]
also auch
\[
\alpha_1\frac{\partial S_{\xi\zeta}}{\partial \xi_1}
+\alpha_2\frac{\partial S_{\xi\zeta}}{\partial \xi_2}
+\alpha_3\frac{\partial S_{\xi\zeta}}{\partial \xi_3}
=M(\beta_1\xi_1+\beta_2\xi_2+\beta_3\xi_3).
\]
Die Grösse $M$ ist von den zufälligen Werthen der Constanten $\beta$ ganz frei geworden.

Man kann dieselbe auch leicht darstellen. Denn man erhält
\[
\frac{E_\xi}{\beta_1\xi_1+\beta_2\xi_2+\beta_3\xi_3}
=M=\int\frac{N\cdot\Sigma\pm c_1x_2dx_3}{c_1\frac{\partial f}{\partial x_1}+c_2\frac{\partial f}{\partial x_2}+c_3\frac{\partial f}{\partial x_3}},
\]
wo
\[
\begin{aligned}
(\beta_1\xi_1+\beta_2\xi_2+\beta_3\xi_3)N
={}&\alpha_1\frac{\partial}{\partial \xi_1}\frac{\Omega}{\Sigma\pm x_1\xi_2\eta_3}
+\alpha_2\frac{\partial}{\partial \xi_2}\frac{\Omega}{\Sigma\pm x_1\xi_2\eta_3}\\
&+\alpha_3\frac{\partial}{\partial \xi_3}\frac{\Omega}{\Sigma\pm x_1\xi_2\eta_3}.
\end{aligned}
\]
Dies giebt
\[
\begin{aligned}
(\beta_1\xi_1+\beta_2\xi_2+\beta_3\xi_3)N
={}&\frac{1}{(\Sigma\pm x_1\xi_2\eta_3)^2}
\left\{
\left|\begin{array}{ccc}
\frac{\partial f}{\partial \xi_1}&\beta_1&\frac{\partial\Omega}{\partial \xi_1}\\
\frac{\partial f}{\partial \xi_2}&\beta_2&\frac{\partial\Omega}{\partial \xi_2}\\
\frac{\partial f}{\partial \xi_3}&\beta_3&\frac{\partial\Omega}{\partial \xi_3}
\end{array}\right|
\left|\begin{array}{ccc}
x_1&\xi_1&\eta_1\\
x_2&\xi_2&\eta_2\\
x_3&\xi_3&\eta_3
\end{array}\right|
\right.\\
&\left.
-\Omega
\left|\begin{array}{ccc}
\frac{\partial f}{\partial \xi_1}&\beta_1&\eta_2x_3-\eta_3x_2\\
\frac{\partial f}{\partial \xi_2}&\beta_2&\eta_3x_1-\eta_1x_3\\
\frac{\partial f}{\partial \xi_3}&\beta_3&\eta_1x_2-\eta_2x_1
\end{array}\right|
\right\}.
\end{aligned}
\]
oder wenn man das erste Determinantenproduct ausführt, und die zweite Determinante in bekannter Weise transformirt:
\[
\begin{aligned}
(\beta_1\xi_1+\beta_2\xi_2+\beta_3\xi_3)N
={}&\frac{1}{(\Sigma\pm x_1\xi_2\eta_3)^2}
\left\{
\left|\begin{array}{ccc}
\sum x_i\frac{\partial f}{\partial \xi_i}&0&\sum \eta_i\frac{\partial f}{\partial \xi_i}\\
\sum x_i\beta_i&\sum \xi_i\beta_i&\sum \eta_i\beta_i\\
\sum x_i\frac{\partial\Omega}{\partial \xi_i}&\Omega&\sum \eta_i\frac{\partial\Omega}{\partial \xi_i}
\end{array}\right|
\right.\\
&\left.
-\Omega
\left|\begin{array}{cc}
\sum \eta_i\frac{\partial f}{\partial \xi_i}&\sum x_i\frac{\partial f}{\partial \xi_i}\\
\sum \eta_i\beta_i&\sum x_i\beta_i
\end{array}\right|
\right\}\\
={}&(\sum\beta_i\xi_i)
\frac{\sum x_i\frac{\partial f}{\partial \xi_i}\sum \eta_i\frac{\partial\Omega}{\partial \xi_i}
-\sum x_i\frac{\partial\Omega}{\partial \xi_i}\sum \eta_i\frac{\partial f}{\partial \xi_i}}{(\Sigma\pm x_1\xi_2\eta_3)^2}.
\end{aligned}
\]
Man hat also
\[
N=-\frac{\sum x_i\frac{\partial f}{\partial \xi_i}\sum \eta_i\frac{\partial\Omega}{\partial \xi_i}
-\sum x_i\frac{\partial\Omega}{\partial \xi_i}\sum \eta_i\frac{\partial f}{\partial \xi_i}}{(\Sigma\pm x_1\xi_2\eta_3)^2},
\]
und daher für das Integral zweiter Gattung die von den willkürlichen Grössen $\alpha$, $\beta$ befreite Form:
\[
\frac{E_\xi}{\beta_1\xi_1+\beta_2\xi_2+\beta_3\xi_3}=\int
\frac{\sum x_i\frac{\partial f}{\partial \xi_i}\sum \eta_i\frac{\partial\Omega}{\partial \xi_i}
-\sum x_i\frac{\partial\Omega}{\partial \xi_i}\sum \eta_i\frac{\partial f}{\partial \xi_i}}{(\Sigma\pm x_1\xi_2\eta_3)^2}
\cdot
\frac{\Sigma\pm c_1x_2dx_3}{c_1\frac{\partial f}{\partial x_1}+c_2\frac{\partial f}{\partial x_2}+c_3\frac{\partial f}{\partial x_3}}.
\]
Da $S_{\xi\zeta}$ bis auf Integrale erster Gattung die Differenz zweier Functionen ist, deren eine nur von $\xi$, deren andre nur von $\zeta$ abhängt, so sieht man, dass $E_\xi$ von $\zeta$ unabhängig ist, bis auf Glieder, welche Integrale erster Gattung sind. Der Punct $\zeta$ ist kein Unstetigkeitspunct von $E_\xi$ mehr; denn in der Nähe dieses Punctes geht $S_{\xi\zeta}$ in
\[
\log\Sigma\pm x_1\zeta_2\xi_3
\]
über, welcher in $E_\xi$ den Term giebt:
\[
\frac{\Sigma\pm x_1\zeta_2\alpha_3}{\Sigma\pm x_1\zeta_2\xi_3}.
\]
Setzt man hierin
\[
x_i=\zeta_i+\varrho\xi_i+\tau\alpha_i,
\]
wo nun $\varrho$, $\tau$ sehr kleine Grössen sind, so erhält man $-\frac{\varrho}{\tau}$. Aber zugleich ist
\[
0=f(x)=\varrho\left(\xi_1\frac{\partial f(\zeta)}{\partial \zeta_1}+\xi_2\frac{\partial f(\zeta)}{\partial \zeta_2}+\xi_3\frac{\partial f(\zeta)}{\partial \zeta_3}\right)
+\tau\left(\alpha_1\frac{\partial f(\zeta)}{\partial \zeta_1}+\alpha_2\frac{\partial f(\zeta)}{\partial \zeta_2}+\alpha_3\frac{\partial f(\zeta)}{\partial \zeta_3}\right),
\]
und also
\[
\lim\frac{\Sigma\pm x_1\zeta_2\alpha_3}{\Sigma\pm x_1\zeta_2\xi_3}
=-\frac{\varrho}{\tau}
=\frac{\alpha_1\frac{\partial f(\zeta)}{\partial \zeta_1}+\alpha_2\frac{\partial f(\zeta)}{\partial \zeta_2}+\alpha_3\frac{\partial f(\zeta)}{\partial \zeta_3}}
{\xi_1\frac{\partial f(\zeta)}{\partial \zeta_1}+\xi_2\frac{\partial f(\zeta)}{\partial \zeta_2}+\xi_3\frac{\partial f(\zeta)}{\partial \zeta_3}},
\]
also nicht unendlich gross.

Dagegen ist in der Nähe von $\xi$
\[
\lim S_{\xi\zeta}=-\log\Sigma\pm x_1\xi_2\zeta_3,
\]
also
\[
\lim E_\xi=-\frac{\Sigma\pm x_1\alpha_2\zeta_3}{\Sigma\pm x_1\xi_2\zeta_3},
\]
und wenn man nun
\[
x_i=\xi_i+\varrho\zeta_i+\tau\alpha_i
\]
setzt, wo wieder $\varrho$, $\tau$ sehr kleine Grössen sind:
\[
\lim E_\xi=\frac{1}{\tau}.
\]
Es wird also $E_\xi$ in der Nähe des Punctes $\xi$ algebraisch unendlich, und zwar von der ersten Ordnung. Bemerkt man, dass wieder
\[
0=f(x)=\varrho\left(\zeta_1\frac{\partial f(\xi)}{\partial \xi_1}+\zeta_2\frac{\partial f(\xi)}{\partial \xi_2}+\zeta_3\frac{\partial f(\xi)}{\partial \xi_3}\right)
+\tau\left(\alpha_1\frac{\partial f(\xi)}{\partial \xi_1}+\alpha_2\frac{\partial f(\xi)}{\partial \xi_2}+\alpha_3\frac{\partial f(\xi)}{\partial \xi_3}\right),
\]
und dass hier der Coefficient von $\tau$ wegen der den $\alpha$ auferlegten Bedingung verschwindet, so sieht man, dass $\varrho$ von höherer Ordnung wird, und dass also kürzer
\[
x_i=\xi_i+\tau\alpha_i
\]
gesetzt werden kann. Es ergeben sich daraus folgende, nur die Verhältnisse der $x$ enthaltende Ausdrücke:
\[
\lim E_\xi=\frac{1}{\tau}
=-\frac{x_2\alpha_3-x_3\alpha_2}{x_2\xi_3-x_3\xi_2}
=-\frac{x_3\alpha_1-x_1\alpha_3}{x_3\xi_1-x_1\xi_3}
=-\frac{x_1\alpha_2-x_2\alpha_1}{x_1\xi_2-x_2\xi_1},
\]
welche, wie es sein muss, von $\zeta$ unabhängig sind.

Führen wir mit Hülfe der Gleichungen
\[
\nu\xi_i=a_i+b_i\xi+c_is_\xi,
\qquad
\nu_1\zeta_i=a_i+b_i\zeta+c_is_\zeta
\]
in $S_{\xi\zeta}$ die Grössen $\xi$, $s_\xi$, $\zeta$, $s_\zeta$ ebenso wie an Stelle der $x$ die $s_x$, $x$ ein, wo denn
\[
F(\xi,s_\xi)=0,
\qquad
F(\zeta,s_\zeta)=0,
\]
und setzen wir zugleich
\[
\alpha_i=\frac{1}{\nu}\left(b_i+c_i\frac{ds_\xi}{d\xi}\right),
\]
so erhalten wir für $E_\xi$ die einfache Form
\[
E_\xi=\frac{\partial S_{\xi\zeta}}{\partial \xi}+\frac{\partial S_{\xi\zeta}}{\partial s_\xi}\frac{ds_\xi}{d\xi}
=\frac{dS_{\xi\zeta}}{d\xi}.
\]
Man kann also bei dieser Darstellung das Integral zweiter Gattung ersetzen durch den Differentialquotienten des Integrals dritter Gattung nach dem Parameter $\xi$.

Das oben den $\alpha$ beigelegte Werthsystem ist in der That möglich, denn die Gleichung
\[
\alpha_1\frac{\partial f(\xi)}{\partial \xi_1}+\alpha_2\frac{\partial f(\xi)}{\partial \xi_2}+\alpha_3\frac{\partial f(\xi)}{\partial \xi_3}=0
\]
geht über in
\[
\frac{\partial F(s_\xi,\xi)}{\partial \xi}
+\frac{\partial F(s_\xi,\xi)}{\partial s_\xi}\frac{ds_\xi}{d\xi}=0,
\]
was eine identische Gleichung ist.

Da in der Nähe von $\xi$ das Integral $S$ in $-\log(x-\xi)$ überging, so wird daselbst
\[
\lim E_\xi=\frac{1}{x-\xi};
\]
und zugleich kann man bemerken, dass weil $E$ durch Differentiation einer Function entsteht, welche nur in zwei Puncten logarithmisch unendlich wird und an allen andern Stellen den Character einer algebraischen Function besitzt, in $E$ keine logarithmisch unendlich werdenden Terme vorkommen können\footnote{Ueber den einfachsten der in diesem Abschnitte behandelten Fälle $(p=0,n=2)$ vergleiche man die ausgezeichnete Arbeit des Hrn. Aronhold, Crelles Journal, Bd. 61. p. 95 folgg.}.

% --- Continuation: source PDF pp. 056--071. ---
\clearpage

\begin{center}
{\large\bfseries Zweiter Abschnitt.}\par\vspace{0.8em}
{\Large\bfseries Das Abelsche Theorem.}
\end{center}

\begin{center}
{\bfseries \S\ 9. Schnittpunctsysteme.}
\end{center}

Die gegebene Curve $n^{\text{ter}}$ Ordnung $f=0$ wird von einer andern Curve $m^{\text{ter}}$ Ordnung $\varphi=0$ in $mn$ Puncten geschnitten. Diese $mn$ Puncte sind, wenn $\varphi=0$ gegeben vorliegt, vollständig bestimmt, und zwar mittelst einer Gleichung $mn^{\text{ten}}$ Grades. Aber wenn die Curve $\varphi=0$ nicht selbst gegeben ist, so genügt eine gewisse Anzahl von Schnittpuncten, um die übrigen zu bestimmen\footnote{Jacobi, de relationibus etc. Crelles Journal, Bd. 15.}. Ist $m<n$, so sind
\[
\frac{m\cdot(m+3)}{2}
\]
Schnittpuncte als gegeben erforderlich; durch diese ist die Curve $m^{\text{ter}}$ Ordnung $\varphi=0$ vollkommen bestimmt, mithin auch die übrigen Schnittpuncte. Ist dagegen $m\geq n$, so kann man, wenn es sich nur um die Schnittpuncte beider Curven handelt, an Stelle der Gleichung $\varphi=0$ die Gleichung
\[
\varphi'=\varphi+\mu f=0
\]
setzen, wo $\mu$ ein Ausdruck $(m-n)^{\text{ter}}$ Ordnung ist; und man kann die
\[
\frac{(m-n+1)(m-n+2)}{2}
\]
Coefficienten von $\mu$ so wählen, dass ebensoviel Coefficienten der reducirten Function $\varphi'$ verschwinden. Man kann also, ohne das Schnittpunctsystem zu verändern, statt einer allgemeinen Curve $\varphi=0$ eine specielle
\[
\varphi'=\varphi+\mu f=0
\]
setzen, welche nur noch von
\[
\frac{(m+1)(m+2)}{2}-\frac{(m-n+1)(m-n+2)}{2}-1
=mn-\frac{(n-1)(n-2)}{2}
\]
Coefficienten abhängt. Diese Curve ist dann durch
\[
mn-\frac{(n-1)(n-2)}{2}
\]
Puncte bestimmt; sind also so viel Schnittpuncte von $f=0$, $\varphi=0$ gegeben, so sind die übrigen
\[
\frac{(n-1)(n-2)}{2}
\]
dadurch bestimmt.

Die beiden Zahlen
\[
mn-\frac{m(m+3)}{2}
\quad\text{und}\quad
\frac{(n-1)(n-2)}{2},
\]
welche angeben, wie viel Schnittpuncte des Systems durch die übrigen bestimmt sind, stimmen noch überein für $m=n-1$ und $m=n-2$. Für $m>n-3$ kann man also immer die Zahl $\frac{(n-1)(n-2)}{2}$ gelten lassen, welche, wie man sieht, von $m$ ganz unabhängig ist. Für kleinere Werthe von $m$ aber ist die Zahl der Puncte, welche durch die übrigen gegeben werden, eine geringere.

Diese Betrachtungen verlangen eine kleine Modification, wenn $\varphi$ durch einige Doppel- oder Rückkehrpuncte von $f$ geht. Diese Puncte rechnen sich dann als Bestimmungsstücke von $\varphi=0$ einfach, dagegen als Schnittpuncte doppelt; so dass die Anzahl der durch die übrigen bestimmten Schnittpuncte um die Anzahl derjenigen Puncte vermindert wird, die in diese Ausnahmspuncte hineinfallen und also von vorn herein bekannt sind. Man hat also den Satz:

Von den Schnittpuncten einer gegebenen Curve $n^{\text{ter}}$ Ordnung $f=0$ mit einer Curve $m^{\text{ter}}$ Ordnung, welche durch $\delta$ Ausnahmspuncte von $f=0$ gehen soll, sind
\[
\begin{array}{ll}
1) & \text{wenn }m\geq n-2,\quad \dfrac{(n-1)(n-2)}{2}-\delta\text{ durch die übrigen bestimmt;}\\[0.4em]
2) & \text{wenn }m<n-2,\quad mn-\dfrac{m(m+3)}{2}-\delta\text{ durch die übrigen bestimmt.}
\end{array}
\]
Dieser Satz erleidet nur dann eine Ausnahme, wenn die gegebenen Schnittpuncte selbst nicht von einander unabhängig sind. Zerfällt nämlich die Curve $\varphi=0$ in mehrere Curven niedern Grades, so zerfällt auch das ganze Schnittpunctsystem in mehrere, und in jedem dieser Schnittpunctsysteme dürfen nur so viel Puncte sich unter den gegebenen befinden, als für eine solche Curve niedern Grades hinreichen, vermittelst des obigen Theorems die übrigen zu bestimmen.

Es ist leicht, die Gleichung zu finden, von der diese übrigen Schnittpuncte abhängen. Eliminirt man aus $\varphi=0$, $f=0$ eine der Variabeln, so bleibt eine Gleichung $mn^{\text{ten}}$ Grades; dieselbe wird durch diejenigen Werthe der Unbekannten erfüllt, welche den gegebenen Schnittpuncten entsprechen. Wenn man also etwa mit $z$ die Unbekannte bezeichnet, mit $z_1,z_2,\ldots,z_{\varrho}$ die den gegebenen Schnittpuncten entsprechenden Werthe derselben, durch $Z=0$ die erwähnte Resultante, so ist
\[
\frac{Z}{(z-z_1)(z-z_2)\cdots(z-z_{\varrho})}=0
\]
die Gleichung, welche die übrigen Schnittpuncte liefert.

\begin{center}
{\bfseries \S\ 10. Transformation der Integralsummen, welche einem Schnittpunctsystem entsprechen.}
\end{center}

Das Abel'sche Theorem lehrt transcendente Beziehungen kennen, welche aus den zwischen den Coordinaten eines Schnittpunctsystems bestehenden Gleichungen entspringen. Es beschäftigt sich nämlich mit der Bestimmung einer Summe von $mn$ gleichartigen Integralen einer Function, welche keine andern Irrationalitäten enthält als die durch $f=0$ bedingten; wobei die untern Grenzen die Schnittpuncte einer Curve $m^{\text{ter}}$ Ordnung $\varphi=0$ mit $f=0$ sind, die obern Grenzen die Schnittpuncte einer andern Curve $\psi=0$ mit $f=0$. Nach unsern frühern Betrachtungen handelt es sich also um Bestimmung einer Summe von Integralen der Form:
\[
\Sigma\int F(x_1,x_2,x_3)\,\Sigma\pm c_1x_2\,dx_3,
\]
wo $F$ eine homogene rationale Function $(-2)^{\text{ter}}$ Ordnung ist, und wo die Summe sich auf die $mn$ Punctepaare erstreckt, welche, dem einen und dem andern Schnittpunctsystem angehörig, untere und obere Grenzen der Integrale werden sollen.

Aber da wir gelernt haben, jede Function $F$ in solche Theile zu zerlegen, dass ihr Integral eines unserer Normalintegrale oder doch aus Differentialquotienten solcher Integrale nach Constanten zusammengesetzt wird (abgesehen von gewissen sofort integrirbaren Theilen), so wird es hinreichen, unter $F$ die Form zu verstehen, welche unter dem Zeichen unseres Normalintegrals dritter Gattung sich befindet, d. h. die Summe zu betrachten:
\[
V=\Sigma\int
\frac{\Omega_{n-2}\cdot \Sigma\pm c_1x_2\,dx_3}
{\Sigma\pm x_1\xi_2\eta_3\cdot(c_1 f'_{x_1}+c_2 f'_{x_2}+c_3 f'_{x_3})}.
\]
Die verschiedenen Integrale entstehen, indem die einzelnen Schnittpuncte sich von ihren Anfangslagen in ihre Endlagen hinbewegen. Man kann dieses dadurch prägnanter hervortreten lassen, dass man einen Curvenbüschel $\varphi+\lambda\psi=0$ betrachtet, und zwischen den Schnittpunctsystemen, welche $f=0$ mit $\varphi=0$ $(\lambda=0)$ und $\psi=0$ $(\lambda=\infty)$ gemein hat, diejenigen einschaltet, welche $f=0$ mit $\varphi+\lambda\psi=0$ (für zwischen $0$ und $\infty$ liegende Werthe von $\lambda$) gemein hat. Diese Zwischensysteme verbinden die Grenzsysteme der Art, dass während $\lambda$ sich von $0$ bis $\infty$ bewegt, die Punctsysteme sich von den untern Grenzen der Integrale zu den obern hinbewegen. Indem man also in jedem Augenblick die $mn$ Schnittpuncte von $f=0$ mit einer Curve des Systems $\varphi+\lambda\psi=0$ betrachtet, werden die Coordinaten dieser Puncte Functionen von $\lambda$ allein, und man erhält
\[
V=\int_0^\infty d\lambda\cdot \Sigma
\frac{\Omega_{n-2}\,\Sigma\pm c_1x_2\dfrac{\partial x_3}{\partial\lambda}}
{\Sigma\pm x_1\xi_2\eta_3\cdot(c_1 f'_{x_1}+c_2 f'_{x_2}+c_3 f'_{x_3})},
\]
wo sich nun die Summe nur noch auf die algebraische Function unter dem Integralzeichen bezieht, und also eine symmetrische Function der Coordinaten der $mn$ Puncte bildet, welche durch den Schnitt von $f=0$ mit $\varphi+\lambda\psi=0$ bestimmt sind.

Transformiren wir zunächst den Ausdruck unter dem Summenzeichen, indem wir Zähler und Nenner mit der Functionaldeterminante
\[
\begin{vmatrix}
\dfrac{\partial f}{\partial x_1}&\dfrac{\partial f}{\partial x_2}&\dfrac{\partial f}{\partial x_3}\\
\dfrac{\partial\varphi}{\partial x_1}&\dfrac{\partial\varphi}{\partial x_2}&\dfrac{\partial\varphi}{\partial x_3}\\
\dfrac{\partial\psi}{\partial x_1}&\dfrac{\partial\psi}{\partial x_2}&\dfrac{\partial\psi}{\partial x_3}
\end{vmatrix}
\]
multipliciren. Diese, mit
\[
\begin{vmatrix}
c_1&c_2&c_3\\
x_1&x_2&x_3\\
\dfrac{\partial x_1}{\partial\lambda}&\dfrac{\partial x_2}{\partial\lambda}&\dfrac{\partial x_3}{\partial\lambda}
\end{vmatrix}
\]
multiplicirt, giebt nach bekannten Sätzen über die homogenen Functionen:
\[
\begin{vmatrix}
c_1\dfrac{\partial f}{\partial x_1}+c_2\dfrac{\partial f}{\partial x_2}+c_3\dfrac{\partial f}{\partial x_3} & nf & \dfrac{\partial f}{\partial\lambda}\\
c_1\dfrac{\partial\varphi}{\partial x_1}+c_2\dfrac{\partial\varphi}{\partial x_2}+c_3\dfrac{\partial\varphi}{\partial x_3} & m\varphi & \dfrac{\partial\varphi}{\partial\lambda}\\
c_1\dfrac{\partial\psi}{\partial x_1}+c_2\dfrac{\partial\psi}{\partial x_2}+c_3\dfrac{\partial\psi}{\partial x_3} & m\psi & \dfrac{\partial\psi}{\partial\lambda}
\end{vmatrix}.
\]
Nun ist für alle in dem Integral vorkommenden Werthsysteme der $x$ immer $f=0$, also auch $\frac{\partial f}{\partial\lambda}=0$; ferner $\varphi+\lambda\psi=0$, also auch
\[
\frac{\partial\varphi}{\partial\lambda}+\lambda\frac{\partial\psi}{\partial\lambda}+\psi=0,
\]
oder wenn man mit $\psi$ multiplicirt, und $-\varphi$ für $\lambda\psi$ setzt:
\[
\psi\frac{\partial\varphi}{\partial\lambda}-\varphi\frac{\partial\psi}{\partial\lambda}+\psi^2=0.
\]
Daher wird unser im Zähler stehendes Determinantenproduct
\[
m\left(c_1\frac{\partial f}{\partial x_1}+c_2\frac{\partial f}{\partial x_2}+c_3\frac{\partial f}{\partial x_3}\right)\psi^2.
\]
Im Nenner ist die Functionaldeterminante zu multipliciren mit
\[
\begin{vmatrix}
x_1&x_2&x_3\\
\xi_1&\xi_2&\xi_3\\
\eta_1&\eta_2&\eta_3
\end{vmatrix};
\]
das Product wird:
\[
\begin{vmatrix}
nf&\xi_1\dfrac{\partial f}{\partial x_1}+\xi_2\dfrac{\partial f}{\partial x_2}+\xi_3\dfrac{\partial f}{\partial x_3}&\eta_1\dfrac{\partial f}{\partial x_1}+\eta_2\dfrac{\partial f}{\partial x_2}+\eta_3\dfrac{\partial f}{\partial x_3}\\
m\varphi&\xi_1\dfrac{\partial\varphi}{\partial x_1}+\xi_2\dfrac{\partial\varphi}{\partial x_2}+\xi_3\dfrac{\partial\varphi}{\partial x_3}&\eta_1\dfrac{\partial\varphi}{\partial x_1}+\eta_2\dfrac{\partial\varphi}{\partial x_2}+\eta_3\dfrac{\partial\varphi}{\partial x_3}\\
m\psi&\xi_1\dfrac{\partial\psi}{\partial x_1}+\xi_2\dfrac{\partial\psi}{\partial x_2}+\xi_3\dfrac{\partial\psi}{\partial x_3}&\eta_1\dfrac{\partial\psi}{\partial x_1}+\eta_2\dfrac{\partial\psi}{\partial x_2}+\eta_3\dfrac{\partial\psi}{\partial x_3}
\end{vmatrix}
\]
\[
\begin{aligned}
&=m\psi\Bigl[\Bigl(\xi_1\frac{\partial f}{\partial x_1}+\xi_2\frac{\partial f}{\partial x_2}+\xi_3\frac{\partial f}{\partial x_3}\Bigr)\Bigl(\eta_1\frac{\partial(\varphi+\lambda\psi)}{\partial x_1}+\eta_2\frac{\partial(\varphi+\lambda\psi)}{\partial x_2}+\eta_3\frac{\partial(\varphi+\lambda\psi)}{\partial x_3}\Bigr)\\
&\quad-\Bigl(\eta_1\frac{\partial f}{\partial x_1}+\eta_2\frac{\partial f}{\partial x_2}+\eta_3\frac{\partial f}{\partial x_3}\Bigr)\Bigl(\xi_1\frac{\partial(\varphi+\lambda\psi)}{\partial x_1}+\xi_2\frac{\partial(\varphi+\lambda\psi)}{\partial x_2}+\xi_3\frac{\partial(\varphi+\lambda\psi)}{\partial x_3}\Bigr)\Bigr].
\end{aligned}
\]
Und wenn man die letzte Klammer für den Augenblick mit $R$ bezeichnet, hat man nach Einführung der erhaltenen Resultate für $V$ den transformirten Ausdruck:
\[
V=\int_0^\infty d\lambda\cdot\Sigma\frac{\Omega\psi}{R}.
\]

Es liegt nun nahe, statt $x_1,x_2,x_3$ neue Variabeln $z,s,t$ einzuführen, welche mit den $x_i$ durch die linearen Beziehungen verbunden sind:
\[
x_1=\xi_1s+\eta_1z+\zeta_1t,\quad
x_2=\xi_2s+\eta_2z+\zeta_2t,\quad
x_3=\xi_3s+\eta_3z+\zeta_3t.
\]
Dabei sind dann $s,z,t$ Dreieckscoordinaten; und zwar bezogen auf ein Coordinatendreieck, dessen Ecken die Unstetigkeitspuncte $\xi$, $\eta$ des Integrals und ein beliebiger dritter Punct $\zeta$ sind. In diesen neuen Variabeln ist für jede Function $f$:
\[
\frac{\partial f}{\partial s}=\xi_1\frac{\partial f}{\partial x_1}+\xi_2\frac{\partial f}{\partial x_2}+\xi_3\frac{\partial f}{\partial x_3},\qquad
\frac{\partial f}{\partial z}=\eta_1\frac{\partial f}{\partial x_1}+\eta_2\frac{\partial f}{\partial x_2}+\eta_3\frac{\partial f}{\partial x_3},
\]
und $R$ nimmt daher den einfachen Ausdruck an:
\[
R=\frac{\partial f}{\partial s}\frac{\partial(\varphi+\lambda\psi)}{\partial z}
-\frac{\partial f}{\partial z}\frac{\partial(\varphi+\lambda\psi)}{\partial s}.
\]
Um die Summe $\frac{\Omega\psi}{R}$ nun auszuführen, welche über alle Werthsysteme auszudehnen ist, die gleichzeitig den Gleichungen $f=0$ und $\varphi+\lambda\psi=0$ genügen, ist es nöthig, auf einen algebraischen Satz zurückzugehen, welcher von Jacobi herrührt\footnote{Jacobi, theoremata nova algebraica, Crelles Journal Bd. 14.}.

\begin{center}
{\bfseries \S\ 11. Der Jacobi'sche Satz.}
\end{center}

Wenn zwei Gleichungen (Curven), $f=0$, $F=0$ $n^{\text{ter}}$ und $m^{\text{ter}}$ Ordnung gegeben sind, welche homogen sind in Bezug auf die drei Variabeln $s,z,t$, so kann man $s$ oder $z$ aus beiden Gleichungen eliminiren und findet dann zwei Gleichungen $Z=0$, $S=0$, welche vom $mn^{\text{ten}}$ Grade bezüglich für $z,t$ und $s,t$ sind. Ist, nach Potenzen von $s$ geordnet:
\[
f=z_0+z_1s+z_2s^2\cdots+z_ns^n,
\qquad
F=Z_0+Z_1s+Z_2s^2\cdots+Z_ns^n,
\]
so wird
\[
Z=\begin{vmatrix}
z_0&z_1&z_2&\cdots\\
0&z_0&z_1&\cdots\\
0&0&z_0&\cdots\\
\cdots&\cdots&\cdots&\cdots\\
Z_0&Z_1&Z_2&\cdots\\
0&Z_0&Z_1&\cdots\\
\cdots&\cdots&\cdots&\cdots
\end{vmatrix};
\]
und wenn man die zweite Verticalreihe mit $s$, die dritte mit $s^2$ u. s. w. multiplicirt zur ersten Reihe addirt, so findet sich
\[
Z=\begin{vmatrix}
f&z_1&z_2&\cdots\\
sf&z_0&z_1&\cdots\\
s^2f&0&z_0&\cdots\\
\cdots&\cdots&\cdots&\cdots\\
F&Z_1&Z_2&\cdots\\
sF&Z_0&Z_1&\cdots\\
\cdots&\cdots&\cdots&\cdots
\end{vmatrix}=Af+BF,
\]
wo $A$, $B$ die Determinanten sind:
\[
A=\begin{vmatrix}
1&z_1&z_2&\cdots\\
s&z_0&z_1&\cdots\\
s^2&0&z_0&\cdots\\
\cdots&\cdots&\cdots&\cdots\\
0&Z_1&Z_2&\cdots\\
0&Z_0&Z_1&\cdots\\
\cdots&\cdots&\cdots&\cdots
\end{vmatrix},\quad
B=\begin{vmatrix}
0&z_1&z_2&\cdots\\
0&z_0&z_1&\cdots\\
0&0&z_0&\cdots\\
\cdots&\cdots&\cdots&\cdots\\
1&Z_1&Z_2&\cdots\\
s&Z_0&Z_1&\cdots\\
\cdots&\cdots&\cdots&\cdots
\end{vmatrix}.
\]
Ganz ebenso hat man, nach $z$ geordnet:
\[
f=s_0+s_1z+s_2z^2\cdots+s_nz^n,
\qquad
F=S_0+S_1z+S_2z^2\cdots+S_nz^n,
\]
und
\[
S=\begin{vmatrix}
s_0&s_1&s_2&\cdots\\
0&s_0&s_1&\cdots\\
0&0&s_0&\cdots\\
\cdots&\cdots&\cdots&\cdots\\
S_0&S_1&S_2&\cdots\\
0&S_0&S_1&\cdots\\
\cdots&\cdots&\cdots&\cdots
\end{vmatrix}=Cf+DF,
\]
wo
\[
C=\begin{vmatrix}
1&s_1&s_2&\cdots\\
z&s_0&s_1&\cdots\\
z^2&0&s_0&\cdots\\
\cdots&\cdots&\cdots&\cdots\\
0&S_1&S_2&\cdots\\
0&S_0&S_1&\cdots\\
\cdots&\cdots&\cdots&\cdots
\end{vmatrix},\quad
D=\begin{vmatrix}
0&s_1&s_2&\cdots\\
0&s_0&s_1&\cdots\\
0&0&s_0&\cdots\\
\cdots&\cdots&\cdots&\cdots\\
1&S_1&S_2&\cdots\\
z&S_0&S_1&\cdots\\
\cdots&\cdots&\cdots&\cdots
\end{vmatrix}.
\]
Die Grade, bis zu denen die Functionen $A,B,C,D$ in Bezug auf die Variabeln $s,z,t$ ansteigen, sind folgende:
\[
\begin{array}{c|ccc}
& s&z&t\\ \hline
A&m-1&mn-n&mn-n\\
B&n-1&mn-m&mn-m\\
C&mn-n&m-1&mn-n\\
D&mn-m&n-1&mn-m
\end{array}
\]
Die Gleichung $Z=0$ liefert $mn$ Werthe von $\frac{z}{t}$, welche den $mn$ Schnittpuncten von $f=0$, $F=0$ entsprechen; ebenso liefert $S=0$ $mn$ Werthe von $\frac{s}{t}$; diese Werthe gehören paarweise zusammen. Denken wir uns die Wurzeln von $Z=0$, sowie die von $S=0$ geordnet, und die zusammengehörigen Wurzeln mit gleichen Indices bezeichnet.

Setzt man in den Gleichungen
\[
Z=Af+BF,
\qquad
S=Cf+DF
\]
für $\frac{z}{t}$, $\frac{s}{t}$ zwei nicht zusammengehörige Werthepaare, so verschwinden $S$ und $Z$, aber nicht $f$ und $F$. Für diese Werthe muss also die Determinante
\[
\Delta=AD-BC
\]
verschwinden.

Setzt man aber in die obigen Gleichungen ein zusammengehöriges Werthepaar, so verschwinden auch $f$, $F$; man erhält durch Differentiation nach $s$ und $z$, indem man die in $f$, $F$ multiplicirten Terme als verschwindend auslässt:
\[
\frac{\partial Z}{\partial z}=A\frac{\partial f}{\partial z}+B\frac{\partial F}{\partial z},\qquad
0=C\frac{\partial f}{\partial z}+D\frac{\partial F}{\partial z},
\]
\[
0=A\frac{\partial f}{\partial s}+B\frac{\partial F}{\partial s},\qquad
\frac{\partial S}{\partial s}=C\frac{\partial f}{\partial s}+D\frac{\partial F}{\partial s},
\]
und es ergiebt sich durch Determinantenbildung:
\[
\begin{vmatrix}
\dfrac{\partial Z}{\partial z}&0\\[0.2em]
0&\dfrac{\partial S}{\partial s}
\end{vmatrix}
=
\begin{vmatrix}A&B\\C&D\end{vmatrix}
\cdot
\begin{vmatrix}
\dfrac{\partial f}{\partial z}&\dfrac{\partial F}{\partial z}\\[0.2em]
\dfrac{\partial f}{\partial s}&\dfrac{\partial F}{\partial s}
\end{vmatrix}.
\]
Die letzte Determinante sei dem früheren analog durch $-R$ bezeichnet; man findet dann als Werth von $\Delta$ für ein zusammengehöriges Werthepaar aus dieser Gleichung
\[
\Delta_i=-\left(\frac{\dfrac{\partial Z}{\partial z}\dfrac{\partial S}{\partial s}}{R}\right)_i,
\]
wo der beigesetzte Index der des fraglichen Werthepaares ist.

Betrachten wir $\Delta$ genauer. Wegen der obigen Tabelle ist $\Delta$ im Ganzen vom Grade $2mn-m-n$; aber es enthält
\[
s\text{ nur bis zum Grade }mn-1,
\qquad
z\text{ nur bis zum Grade }mn-1,
\]
oder, geometrisch ausgedrückt, die Curve $\Delta=0$ hat zwei $(m-1)(n-1)$fache Puncte; den einen bei $s=0$, $t=0$, den andern bei $z=0$, $t=0$. Wegen dieser Eigenschaft von $\Delta$ kann der homogene Ausdruck nullter Dimension
\[
\frac{\Delta t^{m+n}}{SZ},
\]
der nur noch die Quotienten $\frac{s}{t}$, $\frac{z}{t}$ enthält, in Partialbrüche nach $\frac{s}{t}$ und $\frac{z}{t}$ zerlegt werden, und zwar ohne Entstehung einer ganzen Function. Die Zerlegung nach $\frac{s}{t}$ giebt im Zähler den Werth von $\Delta$, darin für $\frac{s}{t}$ die betreffenden Wurzeln von $S=0$ gesetzt. Und wenn man dann auch nach $\frac{z}{t}$ die Zerlegung ausführt, so werden die Zähler eben jene Werthe, für $\frac{z}{t}$ auch noch die betreffenden Wurzeln von $Z=0$ gesetzt. Daher verschwinden von vorn herein alle Glieder, in deren Nenner nicht zusammengehörige Wurzeln von $S$, $Z$ benutzt sind, und es bleibt also nur die einfache Summe:
\[
\frac{\Delta t^{m+n}}{SZ}
=t^2\sum_{i=1}^{i=mn}
\frac{\Delta_i\,t_i^{m+n}}
{\left(\dfrac{\partial S}{\partial s}\right)_i\left(\dfrac{\partial Z}{\partial z}\right)_i(st_i-ts_i)(zt_i-tz_i)},
\]
in welcher nur die im Nenner befindlichen negativen Potenzen von $t$, $t_i$ wieder heraufmultiplicirt sind; oder man hat, wenn man für $\Delta$ seinen Werth setzt,
\[
\frac{\Delta t^{m+n}}{SZ}
=-t^2\sum_{i=1}^{i=mn}
\frac{t_i^{m+n}}{R_i(st_i-ts_i)(zt_i-tz_i)}.
\]
Wenn man bei dieser Partialbruchzerlegung im Zähler statt $t^{m+n}$ eine beliebige Function $U$ $(m+n-2)^{\text{ter}}$ Ordnung mit $t^2$ multiplicirt gehabt hätte, so wären allerdings noch Terme hinzugekommen, welche in Bezug auf $z$ oder $s$, oder auf beide Grössen den Character einer ganzen Function gehabt hätten. Aber abgesehen von diesen Gliedern (was der Strich durch das Gleichheitszeichen andeuten soll), würde man erhalten haben:
\[
\frac{\Delta U}{SZ}\not=
\sum_{i=1}^{i=mn}
\frac{U_i t_i^2}{R_i(st_i-ts_i)(zt_i-tz_i)},
\]
wo nur auf beiden Seiten der Factor $t^2$ ausgelassen ist.

Entwickelt man auf beiden Seiten dieser Gleichung nach aufsteigenden Potenzen von $t$, so müssen die Coefficienten gleicher Potenzen auf beiden Seiten einander gleich sein, solche Glieder ausgenommen, welche auch bei der Entwicklung der übergangenen Terme auftreten würden. Betrachten wir insbesondere die mit dem Factor $\frac{t^0}{sz}$ behafteten Glieder. Da die ausgelassenen Terme nur entweder $z$ oder $s$ im Nenner enthalten, so kann bei der Entwicklung derselben kein Term erscheinen, der gleichzeitig $z$ und $s$ im Nenner enthielte; man braucht daher auf diese Terme keine Rücksicht zu nehmen. Aber statt die betreffenden Coefficienten aus der obigen Gleichung nun durch Entwicklung zu bestimmen, kann man einfach auf beiden Seiten $t=0$ setzen, und dann erst links den Coefficienten von $\frac{1}{sz}$ aufsuchen. Hiedurch reduciren sich $S$, $Z$ auf ihre höchsten Terme, $Z$ auf $\alpha z^{mn}$, $S$ auf $\beta s^{mn}$; und $\alpha$, $\beta$ sind nun nichts anderes als die linken Theile der Gleichungen, welche herauskommen, wenn man in $f=0$, $F=0$ von vorn herein $t=0$ setzt und sodann einmal $s$, einmal $z$ eliminirt. Aber hiebei geht die andre Variable immer von selbst mit heraus; man sieht daher, dass $\beta=\pm\alpha$ sein muss, und dass, wenn für $t=0$
\[
f=a_0s^n+a_1s^{n-1}z+\cdots+a_nz^n=f_0,
\qquad
F=b_0s^m+b_1s^{m-1}z+\cdots+b_mz^m=F_0,
\]
wir $\alpha$ nichts anderes ist als die Determinante:
\[
\alpha=\begin{vmatrix}
a_0&a_1&a_2&\cdots\\
0&a_0&a_1&\cdots\\
0&0&a_0&\cdots\\
\cdots&\cdots&\cdots&\cdots\\
b_0&b_1&b_2&\cdots\\
0&b_0&b_1&\cdots\\
0&0&b_0&\cdots\\
\cdots&\cdots&\cdots&\cdots
\end{vmatrix}.
\]
Schreibt man, um $\beta$ zu bilden, die Coefficienten $a$ und die Coefficienten $b$ in umgekehrter Folge, so entsteht eine Determinante
\[
\beta=(-1)^{mn}\alpha.
\]
Und so hat man den von Jacobi aufgestellten Satz:

Sind $f$, $F$ zwei homogene Functionen $n^{\text{ter}}$ und $m^{\text{ter}}$ Ordnung von $s,z,t$ und erstreckt sich die Summe über die $mn$ Werthsysteme, welche den Gleichungen $f=0$, $F=0$ gemeinsam sind; sind ferner
\[
Z=Af+BF=0,
\qquad
S=Cf+DF=0
\]
die Gleichungen, welche aus $f=0$, $F=0$ durch Elimination von $s$ oder $z$ entstehen, $\alpha$ der gemeinsame Coefficient von $z^{mn}$ in $Z$ und von $s^{mn}$ in $S$, $\Delta=AD-BC$, endlich $U$ eine beliebige ganze Function $m+n-2^{\text{ter}}$ Ordnung, und $U_0$, $\Delta_0$ die Werthe von $U$, $\Delta$ für $t=0$, so ist
\[
\sum_{i=1}^{i=mn}
\frac{U_i}{\left(\dfrac{\partial f}{\partial z}\dfrac{\partial F}{\partial s}-\dfrac{\partial f}{\partial s}\dfrac{\partial F}{\partial z}\right)_i}
=\frac{(-1)^{mn}}{\alpha^2}\left[\frac{\Delta_0U_0}{(sz)^{mn}}\right]_{s^{-1}z^{-1}}.
\]

\begin{center}
{\bfseries \S\ 12. Das Abel'sche Theorem.}
\end{center}

Wenden wir diesen Satz jetzt auf unsern Fall an, indem wir $F=\varphi+\lambda\psi$ setzen, so ist zunächst zu bemerken, dass in $V$ statt $U$ das Product $\Omega\psi$ auftritt. Nun ist $t=0$ die Verbindungslinie der Unstetigkeitspuncte $\xi\eta$ unseres Integrals dritter Gattung. Die Curve $\Omega=0$ wurde aber so bestimmt, dass sie (vgl. \S\ 6) für die Puncte dieser Linie in $\frac{f}{sz}$ übergeht, welches für $t=0$ eine ganze Function $(n-2)^{\text{ter}}$ Ordnung wird, da die Puncte $s=0$, $t=0$ und $z=0$, $t=0$ auf $f=0$ liegen. Mithin kann man
\[
U_0=\frac{f_0}{sz}\cdot\psi_0
\]
setzen.

Ausserdem findet sich aus
\[
Z=Af+BF,
\qquad
S=Cf+DF
\]
durch Combination
\[
DZ-BS=\Delta f,
\]
oder
\[
\frac{\Delta f}{ZS}=\frac{D}{S}-\frac{B}{Z}.
\]
Der Entwicklungscoefficient
\[
\left[\frac{\Delta U}{SZ}\right]_{s^{-1}z^{-1}}
\]
geht also über in
\[
\left[\left(\frac{D}{S}-\frac{B}{Z}\right)\frac{\psi}{sz}\right]_{s^{-1}z^{-1}}
\]
für $t=0$, oder in
\[
\frac{(-1)^{mn}}{\alpha}\left[\frac{D_0\psi_0}{s^{mn+1}}\right]_{s^{-1}}
-\frac{1}{\alpha}\left[\frac{B_0\psi_0}{z^{mn+1}}\right]_{z^{-1}},
\]
und man sieht, dass man in dem ersten Entwicklungscoefficienten nicht nur $t$, sondern auch $z$ gleich Null setzen darf, in dem zweiten nicht nur $t$, sondern auch $s$. Sehen wir, welche Werthe $D_0$, $B_0$ dann annehmen.

Da $Z=\alpha z^{mn}$, so können wir $B_0$ dadurch bilden, dass wir $z^{mn}$ in $z^{(m-1)(n-1)}$ und $z^{m+n-1}$ zerlegen, mit dem letzteren die letzte Verticalreihe von $\alpha$ multipliciren, und die vorhergehenden Verticalreihen mit $sz^{m+n-2}$, $s^2z^{m+n-3}\ldots s^{m+n-1}$ multiplicirt hinzufügen. Die letzte Verticalreihe in $\alpha$ wird dann aus folgenden Elementen bestehen:
\[
s^{m-1}f_0,\\ s^{m-2}zf_0,\ldots,z^{m-1}f_0,
\quad s^{n-1}F_0,
\quad s^{n-2}zF_0,\ldots,z^{n-1}F_0.
\]
Lassen wir nun die Terme mit $f_0$ aus, und heben aus den übrigen den Factor $F_0$ heraus, so ist die übrigbleibende Determinante, mit $z^{(m-1)(n-1)}$ multiplicirt, $B_0$. Aber wir brauchen $B_0$ nur, sofern darin auch $s=0$ gesetzt ist, d. h. indem wir die letzte Verticalreihe auf ihr letztes Glied reduciren. Es ist also $B_0$ gleich zu setzen $z^{m(n-1)}$ multiplicirt mit der Unterdeterminante von $\alpha$, genommen nach dem letzten Element der letzten Reihe.

Ebenso findet man $D_0$. Da $S=\beta s^{mn}$, so theilen wir $s^{mn}$ in $s^{(m-1)(n-1)}$ und in $s^{m+n-1}$, multipliciren mit diesem die letzte Verticalreihe von $\alpha$, fügen die übrigen mit
\[
s^{m+n-2}\cdot z,
\quad s^{m+n-3}\cdot z^2,
\ldots,
 z^{m+n-1}
\]
multiplicirt hinzu, und finden dann als letzte Verticalreihe
\[
s^{m-1}f_0,
\quad s^{m-2}zf_0,\ldots,z^{m-1}f_0,
\quad s^{n-1}F_0,
\quad s^{n-2}zF_0,\ldots,z^{n-1}F_0.
\]
Lassen wir hier die Terme mit $f_0$ aus, heben aus dem übrigen den Factor $F_0$ hervor, und multipliciren die entstandene Determinante mit $s^{(m-1)(n-1)}$, so haben wir $D_0$. Aber wir brauchen $D_0$ nur für $z=0$. Daher können wir die letzte Verticalreihe auf das Glied $s^{n-1}F_0$ beschränken, und können also für $D_0$ setzen $s^{m(n-1)}$ multiplicirt mit der Unterdeterminante von $\beta$ nach dem letzten Element $(b_0)$ der letzten Verticalreihe, oder was dasselbe ist, $(-1)^{mn}\,s^{m(n-1)}$ multiplicirt mit der Unterdeterminante von $\alpha$ nach dem $(m+1)^{\text{ten}}$ Element $b_0$ der ersten Verticalreihe.

Inzwischen haben wir noch nicht davon Gebrauch gemacht, dass in $f_0$ selbst das erste und letzte Glied ($a_0$ und $a_n$) verschwindet, da ja $\xi$ ($z=0$, $t=0$) und $\eta$ ($s=0$, $t=0$) auf der Curve $f=0$ liegen. Bezeichnen wir also mit $\gamma$ die zweite Unterdeterminante von $\alpha$, genommen nach dem Element $b_0$ der ersten, und nach dem letzten Element $b_n$ der letzten Verticalreihe, welches in diesen Reihen die einzigen nicht verschwindenden Elemente sind, so ist
\[
\alpha=b_0b_n\gamma;
\]
und für $B_0$ haben wir zu setzen
\[
z^{m(n-1)}\,b_0\gamma,
\]
für $D_0$
\[
(-1)^{mn}s^{m(n-1)}\,b_n\gamma.
\]
Und so wird unser Entwicklungscoefficient, indem $\gamma$ sich forthebt:
\[
\frac{1}{b_0}\left[\frac{\psi_0}{s^{m+1}}\right]_{s^{-1}}
-\frac{1}{b_n}\left[\frac{\psi_0}{z^{m+1}}\right]_{z^{-1}}.
\]
In der ersten Klammer müssen wir für $\psi_0$ seinen höchsten Term in $s$ setzen. Aber $\psi_0$ war $\psi$, für die $x$ darin, weil $t=0$, gesetzt
\[
x_i=s\xi_i+z\eta_i.
\]
Daher ist der höchste Term in $s$ offenbar
\[
s^m\psi(\xi).
\]
Ebenso wird der höchste Term in $z$, welcher in der zweiten Klammer zu nehmen ist:
\[
z^m\psi(\eta).
\]
Endlich ist zu bemerken, dass $b_0$, $b_n$ ganz ebenso die Coefficienten der höchsten Terme in $s$ und $z$ sind, welche entstehen, wenn man in $F=\varphi+\lambda\psi$ dieselbe Substitution macht, dass also
\[
b_0=\varphi(\xi)+\lambda\psi(\xi),
\qquad
b_n=\varphi(\eta)+\lambda\psi(\eta).
\]
Unser Entwicklungscoefficient ist also
\[
\frac{\psi(\xi)}{\varphi(\xi)+\lambda\psi(\xi)}
-\frac{\psi(\eta)}{\varphi(\eta)+\lambda\psi(\eta)}.
\]
Und so wird unser Integral:
\[
V=-\int_0^\infty \frac{\psi(\xi)\,d\lambda}{\varphi(\xi)+\lambda\psi(\xi)}
+\int_0^\infty \frac{\psi(\eta)\,d\lambda}{\varphi(\eta)+\lambda\psi(\eta)}
\]
\[
=\left[\log\frac{\varphi(\eta)+\lambda\psi(\eta)}{\varphi(\xi)+\lambda\psi(\xi)}\right]_0^\infty
=\log\frac{\psi(\eta)\cdot\varphi(\xi)}{\varphi(\eta)\cdot\psi(\xi)}.
\]
Und wir haben den Satz:

Die Summe von $mn$ gleichartigen Normalintegralen dritter Gattung mit den Unstetigkeitspuncten $\xi$, $\eta$, deren untere Grenzen die $mn$ Werthsysteme sind, welche den $mn$ Schnittpuncten von $f=0$ mit einer Curve $m^{\text{ter}}$ Ordnung $\varphi=0$ entsprechen, und deren obere Grenzen die $mn$ Schnittpuncte von $f=0$ mit einer Curve $m^{\text{ter}}$ Ordnung $\psi=0$ angeben, hat den Werth
\[
\log\frac{\psi(\eta)\cdot\varphi(\xi)}{\varphi(\eta)\cdot\psi(\xi)}.
\]

\begin{center}
{\bfseries \S\ 13. Besondere Fälle.}
\end{center}

Diesem allgemeinen Satze haben wir nun die Erörterung einiger besondern Fälle hinzuzufügen.

1) Wenn das gegebene Integral in ein Integral erster Gattung übergehen soll, oder doch in ein solches Integral dritter Gattung, welches nur in einem Ausnahmspuncte von $f=0$ unstetig wird, so kann man dies dadurch ausdrücken, dass $\Omega$ den Factor $\Sigma\pm x_1\xi_2\eta_3$ enthalten soll. Dieser Factor verschwindet mit $t$, denn die Dreiecksseite $t=0$ ist mit der Geraden $\Sigma\pm x_1\xi_2\eta_3=0$ vollkommen identisch. Daher verschwindet in derjenigen Stelle unserer Entwicklung, wo $t=0$ gesetzt wurde, der Entwicklungscoefficient, und somit $V$. Man hat also den Satz:

Die Summe von gleichartigen Integralen erster Gattung, oder von solchen Integralen dritter Gattung, die nur in einem Ausnahmspuncte von $f=0$ unstetig werden, von dem Schnittpunctsystem der Curve $f=0$ mit einer Curve $m^{\text{ter}}$ Ordnung $\varphi=0$ ausgedehnt bis zu dem Schnittpunctsystem von $f=0$ mit einer Curve $m^{\text{ter}}$ Ordnung $\psi=0$, ist immer gleich Null.

2) Wenn die Curven $\varphi=0$, $\psi=0$ Puncte mit der Curve $f=0$ gemein haben, so fallen die ihnen entsprechenden Integrale einfach aus, weil ihre obern und untern Grenzen zusammenfallen.

3) Wenn der Punct $\xi$ selbst einer der festliegenden Schnittpuncte der Curven $\varphi$, $\psi$ ist, so wird die Gleichung des Theorems illusorisch. Während also im Allgemeinen $\frac{(n-1)(n-2)}{2}$ Integralsummen verschwinden, entsprechend den $\frac{(n-1)(n-2)}{2}$ Coefficienten einer Function $(n-3)^{\text{ter}}$ Ordnung, so reducirt sich diese Zahl um $\delta$, wenn die Curven $\varphi$, $\psi$ durch $\delta$ Doppel- oder Rückkehrpuncte von $f=0$ gehen.

4) Im Allgemeinen existiren also $\frac{(n-1)(n-2)}{2}-\delta$ unabhängige transscendente Bedingungen, welche ausdrücken, dass ebenso viel Schnittpuncte einer Curve $m^{\text{ter}}$ Ordnung mit $f=0$ gegeben sind, wenn diese Curve durch $\delta$ Doppel- oder Rückkehrpuncte von $f=0$ gehen soll. Diese Zahl stimmt genau mit der Anzahl von algebraischen Bedingungen überein, welche wir oben aufgestellt haben, und es sind die transscendenten Bedingungen des Abel'schen Theorems nur eine andre Form dieser letztern.

Wenn aber $m\leq n-3$, so ist die Anzahl der algebraischen Bedingungen kleiner, nämlich nur $mn-\frac{m(m+3)}{2}-\delta$. Daher können dann die transscendenten Bedingungen nicht von einander unabhängig sein, sondern müssen sich auf ebenso viel von einander unabhängige Bedingungen reduciren. So erhält man für $m=n-3$ von transscendenten Bedingungen $\frac{n(n-3)}{2}-\delta+1$, während nur $\frac{n(n-3)}{2}-\delta$ algebraische Bedingungen vorhanden sind.

5) Schreibt man die Gleichung des Theorems für die Integrale dritter Gattung in der Form
\[
(1)\quad \Sigma S_{\xi\eta}(x)-\Sigma S_{\xi\eta}(y)
=\log\frac{\varphi(\xi)}{\psi(\xi)}-\log\frac{\varphi(\eta)}{\psi(\eta)},
\]
wo die $x$ das eine, die $y$ das andre Schnittpunctsystem darstellen, so erhält man daraus die Gleichung für die Integrale zweiter Gattung, indem man $\eta_i=\xi_i-\varepsilon\alpha_i$ setzt (vgl. \S\ 8) und $\varepsilon$ gegen Null convergiren lässt. Es tritt dann an die Stelle von $S_{\xi\eta}$ das Integral $E_\xi$, vermehrt um Integrale erster Gattung, die aber aus obiger Summe wegen der soeben bewiesenen Sätze\footnote{Der Satz 1) lässt sich auch aus der Formel (1) ableiten. Bildet man nämlich $S_{\xi\eta}$ auf doppelte Weise (vgl. p. 27.), so gilt jedesmal die Gleichung (1). Aber die Differenz entsprechender $S$ ist nach dem a. a. O. angeführten Satz ein Integral erster Gattung $I(x)$, und man hat daher, wenn man die beiden Formen der Gleichung (1) von einander abzieht, $\Sigma I(x)-\Sigma I(y)=0$, was wieder der Satz 1) ist.} verschwinden, und es bleibt:
\[
\Sigma E_\xi(x)-\Sigma E_\xi(y)
=\alpha_1\frac{\partial\log\frac{\varphi(\xi)}{\psi(\xi)}}{\partial\xi_1}
+\alpha_2\frac{\partial\log\frac{\varphi(\xi)}{\psi(\xi)}}{\partial\xi_2}
+\alpha_3\frac{\partial\log\frac{\varphi(\xi)}{\psi(\xi)}}{\partial\xi_3},
\]
welche wiederum illusorisch wird, wenn der Punct $\xi$ beiden Curven $\varphi=0$, $\psi=0$ angehört.

6) Es treten Fälle ein, wo in dem allgemeinen Satze über die Normalintegrale dritter Gattung die logarithmische Function verschwindet. Wenn nämlich $\xi$ und $\eta$ mit den Grundpuncten des Büschels $\varphi+\lambda\psi=0$ auf einer Curve $m^{\text{ter}}$ Ordnung liegen, so ist es dieselbe Curve dieses Büschels, welche durch $\xi$ und durch $\eta$ geht. Ist $\varrho$ ihr Parameter, so hat man
\[
\varphi(\xi)+\varrho\psi(\xi)=0,
\qquad
\varphi(\eta)+\varrho\psi(\eta)=0,
\]
und daher
\[
\frac{\varphi(\xi)}{\psi(\xi)}\cdot\frac{\psi(\eta)}{\varphi(\eta)}=1,
\]
sodass die Gleichung des Abel'schen Theorems in
\[
\Sigma S_{\xi\eta}(x)-\Sigma S_{\xi\eta}(y)=0
\]
übergeht.



% --- Continuation source pp. 072--088 ---

\begin{center}
{\large Dritter Abschnitt.}\\[0.35em]
{\Large Die eindeutigen Transformationen.}
\end{center}

\subsection*{\S\,14. Transformation der algebraischen Integrale.}

Das Integral, welches, wie die von uns betrachteten, sich auf eine Curve
\[
f(x_1,x_2,x_3)=0
\]
bezieht, und keine andern Irrationalitäten enthält, als diejenigen, welche durch diese Gleichung selbst begründet werden, kann auf unendlich viele Arten in andre Integrale ähnlicher Natur übergeführt werden, welche sich auf eine andre Curve
\[
F(y_1,y_2,y_3)=0
\]
beziehen. Setzen wir nämlich an Stelle der \(x\) in unserm Integral irgend rationale Functionen der \(y\), so geht einerseits das Integral in ein ganz ähnliches Integral in den \(y\) über, andrerseits geht \(f=0\) in eine andre Gleichung \(F=0\) über. Um das erste zu beweisen, muss man nur zeigen, dass der Ausdruck
\[
\frac{\Omega\cdot \Sigma\pm c_1x_2\,dx_3}
 {c_1f'_{x_1}+c_2f'_{x_2}+c_3f'_{x_3}},
\]
wo \(\Omega\) eine rationale (nicht nothwendig ganze) Function \((n-3)^{\text{ter}}\) Dimension bedeutet, in einen ähnlichen übergeführt wird. Seien die Transformationsformeln
\[
\tag{1}
\begin{aligned}
x_1&=\varphi_1(y_1,y_2,y_3),\\
x_2&=\varphi_2(y_1,y_2,y_3),\\
x_3&=\varphi_3(y_1,y_2,y_3),
\end{aligned}
\]
(wo links auch, wenn man will, noch ein unbestimmter Factor hinzugefügt werden kann), so wird \(f=0\) in \(MF=0\) übergehen, wo \(M\) ein Factor ist, welcher sich möglicherweise absondert, der aber \(1\) sein kann. Da
\[
f=MF,
\]
so ist auch
\[
\frac{\partial f}{\partial x_i}
=M\left(
\frac{\partial F}{\partial y_1}\frac{\partial y_1}{\partial x_i}
+\frac{\partial F}{\partial y_2}\frac{\partial y_2}{\partial x_i}
+\frac{\partial F}{\partial y_3}\frac{\partial y_3}{\partial x_i}
\right)
\]
für alle Werthsysteme, für welche, wie hier angenommen werden soll, \(F\) verschwindet; also:
\[
c_1f'_{x_1}+c_2f'_{x_2}+c_3f'_{x_3}
=M\left[k_1F'_{y_1}+k_2F'_{y_2}+k_3F'_{y_3}\right],
\]
wo
\[
k_h=c_1\frac{\partial y_h}{\partial x_1}
+c_2\frac{\partial y_h}{\partial x_2}
+c_3\frac{\partial y_h}{\partial x_3}.
\]
Multipliciren wir ferner Zähler und Nenner des Differentialausdrucks mit der Determinante
\[
\Delta=\Sigma\pm
\frac{\partial y_1}{\partial x_1}
\frac{\partial y_2}{\partial x_2}
\frac{\partial y_3}{\partial x_3},
\]
so wird
\[
\Delta\cdot\Sigma\pm c_1x_2\,dx_3
=r\begin{vmatrix}
k_1&k_2&k_3\\
y_1&y_2&y_3\\
dy_1&dy_2&dy_3
\end{vmatrix},
\]
wo \(r\) die Ordnung der Functionen \(\varphi\) angiebt. Man findet also endlich, wenn man noch für \(\frac1\Delta\) die Determinante
\[
D=\Sigma\pm
\frac{\partial x_1}{\partial y_1}
\frac{\partial x_2}{\partial y_2}
\frac{\partial x_3}{\partial y_3},
\]
eine ganze Function der \(y\), setzt:
\[
\frac{\Omega\cdot\Sigma\pm c_1x_2\,dx_3}
 {c_1f'_{x_1}+c_2f'_{x_2}+c_3f'_{x_3}}
=
\frac{r\Omega D\cdot\Sigma\pm k_1y_2\,dy_3}
 {M\cdot (k_1F'_{y_1}+k_2F'_{y_2}+k_3F'_{y_3})}.
\]
Hier ist wirklich der Ausdruck rechts ganz ähnlicher Natur, wie der Ausdruck links, nur dass an Stelle der Function \(\Omega\) die Function \(\frac{r\Omega D}{M}\) getreten ist, und dass an Stelle der ganz willkürlichen Grössen \(c\) die ebenso willkürlichen Grössen \(k\) erscheinen.

Aber es ist sehr wichtig, dass, wenn das ursprüngliche Integral ein Integral erster Gattung, also \(\Omega\) eine ganze Function \((n-3)^{\text{ter}}\) Ordnung war, auch das transformirte Integral wieder ein Integral erster Gattung wird; indem sich der Zähler \(\Omega D\) dann immer mit Hülfe von \(F=0\) so transformiren lässt, dass er durch \(M\) theilbar wird, und dass der Rest für die Doppel- und Rückkehrpuncte von \(F=0\) verschwindet.

Um dies zu beweisen, betrachten wir die Gleichungen:
\[
\begin{aligned}
M\frac{\partial F}{\partial y_1}=0
&=\frac{\partial f}{\partial x_1}\frac{\partial x_1}{\partial y_1}
+\frac{\partial f}{\partial x_2}\frac{\partial x_2}{\partial y_1}
+\frac{\partial f}{\partial x_3}\frac{\partial x_3}{\partial y_1},\\
M\frac{\partial F}{\partial y_2}=0
&=\frac{\partial f}{\partial x_1}\frac{\partial x_1}{\partial y_2}
+\frac{\partial f}{\partial x_2}\frac{\partial x_2}{\partial y_2}
+\frac{\partial f}{\partial x_3}\frac{\partial x_3}{\partial y_2},\\
M\frac{\partial F}{\partial y_3}=0
&=\frac{\partial f}{\partial x_1}\frac{\partial x_1}{\partial y_3}
+\frac{\partial f}{\partial x_2}\frac{\partial x_2}{\partial y_3}
+\frac{\partial f}{\partial x_3}\frac{\partial x_3}{\partial y_3}.
\end{aligned}
\]
Diese Gleichungen müssen zusammenbestehen für alle etwaigen Unstetigkeitspuncte des gegebenen Integrals. Es ist zu zeigen, dass, wenn die hier gleich Null gesetzten Ausdrücke verschwinden, der Zähler \(\Omega D\) immer in ebenso hoher Ordnung Null wird. Dann folgt, dass das neue Integral nie unstetig wird, und nach \S\,5. ein Integral erster Gattung sein muss. Nun können aber jene Gleichungen entweder dadurch neben einander bestehen, dass die Determinante der rechten Theile, also \(D\) verschwindet; oder dadurch, dass
\[
\frac{\partial f}{\partial x_1}=0,\qquad
\frac{\partial f}{\partial x_2}=0,\qquad
\frac{\partial f}{\partial x_3}=0.
\]
Im letzten Falle hat entweder \(f\) im Puncte \(x\) einen Doppelpunct, und dann verschwindet \(\Omega\), weil das gegebene Integral erster Gattung sein sollte; oder es verschwinden sämmtliche \(x\), indem für das fragliche Werthsystem der \(y\) die drei Functionen \(\varphi\) gleichzeitig verschwinden. In diesem Falle werden die Ausdrücke rechts Null von der \((n-1)^{\text{ten}}\) Ordnung, da die \(\frac{\partial f}{\partial x}\) von der \((n-1)^{\text{ten}}\) Ordnung Null werden, sobald alle \(x\) verschwinden; aber zugleich wird \(D\) (vgl. unten) von der ersten Ordnung Null und \(\Omega\) als Function \((n-2)^{\text{ter}}\) Ordnung der \(x\) wird von der \((n-2)^{\text{ten}}\) Ordnung Null; das Product \(\Omega D\) also Null von der \((n-1)^{\text{ten}}\) Ordnung, wie der Nenner.

Man sieht also wirklich, dass, sobald der Nenner verschwindet, der Zähler in gleicher Weise verschwindet, und dass man also ein Integral erster Gattung vor sich hat.

\subsection*{\S\,15. Begriff der eindeutigen Transformationen. Erhaltung der Zahl \(p\).}

Unter den verschiedenen Transformationen sind nun vorläufig diejenigen von besonderem Interesse, in denen sich die \(x\) ebensowohl rational durch die \(y\), als umgekehrt die \(y\) rational durch die \(x\) ausdrücken lassen. Dies ist insofern denkbar, als man, indem man die Gleichungen (1) mit \(F=0\) verbindet, durch Elimination der \(y\) wieder zu der Gleichung \(f=0\) zurückkehren kann, und im Verlauf des Eliminationsprocesses die \(y\) mit ganzen rationalen Functionen der \(x\) proportional findet. Aber allerdings wird dies nur eintreten, wenn sich ein oben angenommener Factor \(M\) wirklich absondert, und also \(F\) nicht eine Function von \(\varphi_1,\varphi_2,\varphi_3\) ist. In diesem Falle freilich würde der Eliminationsprocess die \(y\) als Functionen der \(x\) nicht anders geben, als durch Auflösung der Gleichungen (1) selbst, mithin im Allgemeinen mit Irrationalitäten behaftet.

Solche Transformationen, in welchen, eine gewisse Gleichung \(f=0\) oder \(F=0\) vorausgesetzt, die \(x\) sich durch die \(y\), und diese (oder doch ihre Verhältnisse, worauf es allein ankommt) sich durch jene rational ausdrücken lassen, nennen wir eindeutige Transformationen. Die beiden Curven \(f=0\), \(F=0\) entsprechen sich nämlich in diesem Fall Punct für Punct eindeutig; jedem Puncte \(y\) der zweiten Curve entspricht im Allgemeinen aus (1) ein einziger Punct \(x\) der ersten und umgekehrt.

In einem solchen Fall liefert also jedes für die \(x\) zu bildende Integral erster Gattung ein solches in den \(y\), und jedes in den \(y\) zu bildende Integral erster Gattung ein solches für die \(x\). Hieraus folgt, dass die Zahl der Integrale erster Gattung in beiden Fällen die gleiche sein muss; oder wie wir es auch ausdrücken können:

\begin{quote}
Für Curven, welche einander Punct für Punct eindeutig entsprechen, hat die Zahl \(p\) denselben Werth.\footnote{Vgl. Riemann, Crelles Journal Bd. 54. p. 133.}
\end{quote}

Sind also \(n,\nu\) die Ordnungen, \(d+r,\delta+\varrho\) die Zahlen der Doppel- und Rückkehrpuncte für beide Curven, so ist immer
\[
\frac{(n-1)(n-2)}{2}-d-r
=
\frac{(\nu-1)(\nu-2)}{2}-\delta-\varrho.\footnote{Eine interessante Anwendung gestattet dieser Satz auf die Theorie der reciproken Curven, oder auf die Vergleichung der Singularitäten, welche eine algebraische Curve in Punct- und Liniencoordinaten aufweist. Die Gleichungen, welche die Coordinaten \(y\) der Tangente mit den Coordinaten \(x\) des Berührungspuncts verbinden, gehören unstreitig in die hier behandelte Classe von Transformationsgleichungen, da im Allgemeinen jedem Punct nur eine Tangente, jeder Tangente nur ein Berührungspunct entspricht. Bei dem Uebergange von Punct- zu Liniencoordinaten entspricht aber der Ordnung \(n\) die Classe \(k\), der Anzahl der Doppelpuncte \(d\) die Anzahl der Doppeltangenten \(t\), der Anzahl der Rückkehrpuncte \(r\) die Anzahl der Wendetangenten \(w\). Die Gleichsetzung der beiden Werthe von \(p\) giebt also
\[
\frac{(n-1)(n-2)}{2}-d-r
=
\frac{(k-1)(k-2)}{2}-t-w;
\]
eine Gleichung, welche man mit Hülfe der bekannten Plücker'schen Formeln verificiren kann, welche aber zugleich, da die Ausdrücke von \(k,w\) leicht zu bilden sind, die directeste Herleitung für die Anzahl der Doppeltangenten einer Curve bildet.}
\]

\subsection*{\S\,16. Directer Beweis für die Erhaltung der Zahl \(p\).}

Dieser Satz scheint wichtig genug, um einen andern, direct algebraischen Beweis wünschen zu lassen. Einen solchen kann man folgendermassen führen.

Gehen wir von der Gleichung \(f=0\) aus, und von den umgekehrten Formeln (1):
\[
\tag{2}
\begin{aligned}
\mu y_1&=\Phi_1(x_1,x_2,x_3),\\
\mu y_2&=\Phi_2(x_1,x_2,x_3),\\
\mu y_3&=\Phi_3(x_1,x_2,x_3),
\end{aligned}
\]
wo \(\mu\) einen unbestimmten Factor bedeutet.

Die Functionen \(\Phi\) seien vom Grade \(s\); die Gleichungen \(\Phi_1=0\), \(\Phi_2=0\), \(\Phi_3=0\) stellen drei Curven \(s^{\text{ter}}\) Ordnung dar; nehmen wir der Allgemeinheit wegen an, dass eine Reihe von Puncten existirt, welche diesen drei Curven und zugleich der Curve \(f=0\) gemeinsam seien, und bezeichnen wir die Anzahl dieser gemeinsamen Puncte mit \(\sigma+\tau\), wo \(\sigma\) die Anzahl solcher Puncte sind, welche einfache Puncte auf \(f=0\), \(\tau\) die Anzahl derjenigen, welche etwa Doppel- oder Rückkehrpuncte von \(f=0\) sind.

Wenn man aus den Gleichungen (2) und aus \(f=0\) die \(x\) eliminirt, so gelangt man zu einer Gleichung \(F=0\) \(\nu^{\text{ten}}\) Grades in den \(y\); es ist die Gleichung der transformirten Curve. Bestimmen wir zunächst die Ordnung dieser Curve.

Schneiden wir \(F=0\) mit einer beliebigen Geraden
\[
\alpha_1y_1+\alpha_2y_2+\alpha_3y_3=0,
\]
so erhalten wir die Ordnung von \(F\) in der Anzahl von Lösungen, welche dieses Problem zulässt. Aber setzen wir statt der uns noch unbekannten Gleichung \(F=0\) die Gleichung \(f=0\) in Verbindung mit den Gleichungen (2). Indem wir dann die Werthe der \(y\) in die Gleichung der Geraden eintragen, bekommen wir zur Bestimmung der Puncte \(x\), welche den Schnittpuncten von \(f=0\) mit der Geraden entsprechen, die Gleichungen:
\[
\tag{3}
\begin{cases}
f=0,\\
\alpha_1\Phi_1+\alpha_2\Phi_2+\alpha_3\Phi_3=0.
\end{cases}
\]
Jede Lösung dieser Gleichungen liefert im Allgemeinen einen Punct \(y\), sodass im Ganzen \(ns\) Schnittpuncte existiren müssten. Aber hievon sind solche Lösungen auszuschliessen, welche etwa die Gleichungen (3) unabhängig von den speciellen Werthen der \(\alpha\) erfüllen, welche also mit dem vorgelegten Problem nichts gemein haben. Für diese müssen \(f\) und die \(\Phi\) gleichzeitig verschwinden; es existiren also solche fremde Lösungen, und zwar führen sie auf die \(\sigma+\tau\) Puncte \(x\), von denen oben die Rede war, und von denen die \(\sigma\) ersten als einfache Schnittpuncte der beiden Curven (3), die \(\tau\) andern aber als zwiefache erscheinen, da sie Doppel- oder Rückkehrpuncte der ersten Curve sind. Man sieht also, dass die Zahl \(ns\) der Lösungen von (3) um \(\sigma+2\tau\) zu vermindern ist, und man findet daher die Ordnung von \(F\) aus der Gleichung
\[
\nu=ns-\sigma-2\tau.
\]

Bestimmen wir jetzt die Zahl der Doppel- und Rückkehrpuncte für die transformirte Curve. Diese können zum Theil den Doppel- und Rückkehrpuncten der gegebenen Curve entsprechen; andrerseits aber können, so wie Doppelpuncte der ersten Curve sich in Punctepaare der zweiten aufgelöst haben können, oder Rückkehrpuncte in zwei unendlich nahe Puncte, so auch Punctepaare der ersten Curve sich in Doppelpuncte der zweiten vereinigt haben, und dadurch auf der transformirten Curve neue Doppelpuncte oder, wenn jene Puncte unendlich nahe lagen, Rückkehrpuncte entstanden sein.

Bezeichnen wir durch \(\delta',\varrho'\) die Anzahl der Doppel- und Rückkehrpuncte von \(F\), denen Doppel- und Rückkehrpuncte von \(f\) entsprechen, durch \(\delta'',\varrho''\) die Anzahl derjenigen, denen Punctepaare auf \(f\) entsprechen, so dass
\[
\delta+\varrho=\delta'+\varrho'+\delta''+\varrho''.
\]

Differenzirt man die Gleichungen (2), so drücken sich die Differentiale der \(y\) durch die Differentiale der \(x\) aus mit Hülfe der Formeln:
\[
\tag{4}
\begin{aligned}
\mu\,dy_1+y_1\,d\mu
&=\frac{\partial\Phi_1}{\partial x_1}dx_1+
  \frac{\partial\Phi_1}{\partial x_2}dx_2+
  \frac{\partial\Phi_1}{\partial x_3}dx_3,\\
\mu\,dy_2+y_2\,d\mu
&=\frac{\partial\Phi_2}{\partial x_1}dx_1+
  \frac{\partial\Phi_2}{\partial x_2}dx_2+
  \frac{\partial\Phi_2}{\partial x_3}dx_3,\\
\mu\,dy_3+y_3\,d\mu
&=\frac{\partial\Phi_3}{\partial x_1}dx_1+
  \frac{\partial\Phi_3}{\partial x_2}dx_2+
  \frac{\partial\Phi_3}{\partial x_3}dx_3.
\end{aligned}
\]
Aus diesen Gleichungen folgt, dass wenn \(x\) ein solcher Doppel- oder Rückkehrpunct von \(f=0\) ist, für welchen nicht alle \(\Phi\) gleichzeitig, also auch \(\mu\), verschwinden, jeder Fortschreitungsrichtung \(dx\) auf der ersten Curve auch eine Fortschreitungsrichtung \(dy\) auf der zweiten entspricht, also dem Doppel- oder Rückkehrpunct \(x\) ein Doppel- oder Rückkehrpunct \(y\). Ist dagegen \(x\) einer derjenigen Doppel- oder Rückkehrpuncte, für welche die \(\Phi\) verschwinden, so verschwindet auch \(\mu\). Von den \(dx\) kann man eines beliebig Null setzen; zwischen den übrigen besteht dann die Gleichung \(2^{\text{ten}}\) Grades:
\[
\frac{\partial^2 f}{\partial x_1^2}dx_1^2
+2\frac{\partial^2 f}{\partial x_1\partial x_2}dx_1dx_2+\cdots=0.
\]
Indem man aus dieser Gleichung die Verhältnisse der \(dx\) bestimmt, geben die Gleichungen (4), aus denen die \(dy\) verschwunden sind, zwei Werthsysteme der \(y\), also zwei Puncte von \(F=0\), welche einem solchen Doppelpunct von \(f=0\) entsprechen.\footnote{Dies wird nicht aufgehoben durch den Umstand, dass die beiden Puncte \(y\) einander unendlich nahe rücken können, wenn der Doppelpunct auf \(f=0\) in einen Rückkehrpunct übergeht.} Die \(\tau\) Doppel- oder Rückkehrpuncte von \(f=0\), welche auch auf den Curven \(\Phi=0\) liegen, geben also keine Doppel- oder Rückkehrpuncte von \(F=0\). Mithin ist
\[
\delta'+\varrho'=d+r-\tau.
\]

Die übrigen Doppel- und Rückkehrpuncte von \(F=0\) entstehen nun dadurch, dass zwei Puncten \(x\) derselbe Punct \(y\) entspricht; dass also Werthsysteme der \(x,x'\) bestehen, so dass
\[
\tag{5}
\begin{gathered}
\Phi_1(x)=\Phi_1(x'),\qquad
\Phi_2(x)=\Phi_2(x'),\qquad
\Phi_3(x)=\Phi_3(x'),\\
f(x)=0,\qquad f(x')=0;
\end{gathered}
\]
wobei der Kürze wegen nur immer ein Argument statt dreier geschrieben ist. Die Anzahl von Lösungen, welche diese Gleichungen zulassen, bestimmt die Zahl \(\delta''+\varrho''\).

Nun ist aber in diesen Gleichungen die selbstverständliche und unbrauchbare Lösung enthalten:
\[
x_1=x_1',\qquad x_2=x_2',\qquad x_3=x_3'.
\]
Um diese zu vermeiden, setzen wir an Stelle der Gleichungen (5) zunächst diese:
\[
\tag{6}
\begin{gathered}
\Phi_1(x)+\varepsilon\Psi_1(x)=\Phi_1(x'),\\
\Phi_2(x)+\varepsilon\Psi_2(x)=\Phi_2(x'),\qquad f(x)=0,\quad f(x')=0,\\
\Phi_3(x)+\varepsilon\Psi_3(x)=\Phi_3(x'),
\end{gathered}
\]
wo die \(\Psi\) ganz beliebige Functionen \(s^{\text{ter}}\) Ordnung sind, und eliminiren die \(x'\) aus diesen Gleichungen. Indem wir im Resultat \(\varepsilon\) gegen Null convergiren lassen, erhalten wir die gesuchte Gleichung, welche mit \(f=0\) zusammen die den neuentstehenden Doppel- und Rückkehrpuncten entsprechenden Punctenpaare liefert.

Da inzwischen der Voraussetzung nach die Elimination der \(x\) aus den Gleichungen
\[
\Phi_1(x')=\mu y_1,\qquad
\Phi_2(x')=\mu y_2,\qquad
\Phi_3(x')=\mu y_3,\qquad
f(x')=0
\]
die Gleichung
\[
F(y)=0
\]
liefert, mithin auch für alle beliebigen Werthe der \(x'\):
\[
F[\Phi_1(x'),\Phi_2(x'),\Phi_3(x')]=0,
\]
so giebt die Elimination der \(x'\) aus den Gleichungen (6):
\[
F[\Phi_1(x)+\varepsilon\Psi_1(x),\Phi_2(x)+\varepsilon\Psi_2(x),
\Phi_3(x)+\varepsilon\Psi_3(x)]=0
\]
\[
=F(\Phi)+\varepsilon\left[
\Psi_1(x)F'_{\Phi_1}+\Psi_2(x)F'_{\Phi_2}+\Psi_3(x)F'_{\Phi_3}
\right]+\cdots,
\]
und das erste Glied dieser Entwicklung verschwindet, so dass nach Division mit \(\varepsilon\) die gesuchte Resultante in der Form bleibt:
\[
\tag{7}
R=\Psi_1(x)F'_{\Phi_1}+\Psi_2(x)F'_{\Phi_2}+\Psi_3(x)F'_{\Phi_3}=0.
\]
Aber diese Gleichung hängt noch von den ganz zufälligen Functionen \(\Psi\) ab; es kommt darauf an, diese in einem unwesentlichen Factor abzutrennen.

Bemerken wir nun, dass die Gleichungen (6) für unendlich kleine Werthe von \(\varepsilon\) noch befriedigt werden können, sobald die Grössen \(x\) ausser der Gleichung \(f(x)=0\) noch einer zweiten von den \(\Psi\) abhängigen Gleichung genügen, und zugleich die \(x'\) sich unendlich wenig von ihnen unterscheiden, wenn diese unendlich kleinen Unterschiede passend gewählt sind. Setzt man nämlich
\[
x_i'=x_i+\varepsilon\xi_i,
\]
und entwickelt nach dem Taylor'schen Lehrsatz, so gehen die Gleichungen (6) mit Weglassung des Factors \(\varepsilon\) über in:
\[
\begin{aligned}
\Psi_1(x)&=\xi_1\Phi'_1(x_1)+\xi_2\Phi'_1(x_2)+\xi_3\Phi'_1(x_3),\\
\Psi_2(x)&=\xi_1\Phi'_2(x_1)+\xi_2\Phi'_2(x_2)+\xi_3\Phi'_2(x_3),\\
\Psi_3(x)&=\xi_1\Phi'_3(x_1)+\xi_2\Phi'_3(x_2)+\xi_3\Phi'_3(x_3),\\
0&=\xi_1f'(x_1)+\xi_2f'(x_2)+\xi_3f'(x_3).
\end{aligned}
\]
Nach Elimination der \(\xi\) giebt dies für die \(x\) die einzige Bedingung:
\[
S=0=
\begin{vmatrix}
\Psi_1(x)&\Phi'_1(x_1)&\Phi'_1(x_2)&\Phi'_1(x_3)\\
\Psi_2(x)&\Phi'_2(x_1)&\Phi'_2(x_2)&\Phi'_2(x_3)\\
\Psi_3(x)&\Phi'_3(x_1)&\Phi'_3(x_2)&\Phi'_3(x_3)\\
0&f'(x_1)&f'(x_2)&f'(x_3)
\end{vmatrix}.
\]
Diejenigen Werthsysteme der \(x\), welche die Gleichungen \(f=0\), \(S=0\) befriedigen, genügen auch der Gleichung \(R=0\); aber die erstern entsprechen keinen Doppel- oder Rückkehrpuncten, sondern Doppel- oder Rückkehrpuncte liefern nur solche \(x\), welche der Gleichung \(\frac{R}{S}=0\) genügen. Die Division links auszuführen ist im Allgemeinen selbst mit Hülfe von \(f=0\) nicht möglich.

Da die Functionen \(\Psi\) ganz beliebig sind, so muss man haben
\[
\frac{R}{S}
=\frac{F'_{\Phi_1}}{\Sigma\pm\Phi'_2(x_1)\Phi'_3(x_2)f'(x_3)}
=\frac{F'_{\Phi_2}}{\Sigma\pm\Phi'_3(x_1)\Phi'_1(x_2)f'(x_3)}
=\frac{F'_{\Phi_3}}{\Sigma\pm\Phi'_1(x_1)\Phi'_2(x_2)f'(x_3)}.
\]
Es ist leicht, einen vierten Werth dieses Verhältnisses anzugeben, welcher ein symmetrischer ist. Ist nämlich
\[
F(\Phi_1,\Phi_2,\Phi_3)=M\cdot f,
\]
so hat man wegen \(f=0\) auch
\[
F'_{\Phi_1}\frac{\partial\Phi_1}{\partial x_i}
+F'_{\Phi_2}\frac{\partial\Phi_2}{\partial x_i}
+F'_{\Phi_3}\frac{\partial\Phi_3}{\partial x_i}
-M\frac{\partial f}{\partial x_i}=0.
\]
Wenn man in den drei Gleichungen, welche hiedurch dargestellt sind, die Grössen \(F'_{\Phi_1},F'_{\Phi_2},F'_{\Phi_3},-M\) als die Unbekannten ansieht und ihre Verhältnisse durch Unterdeterminanten ausdrückt, so findet man die Gleichheit der oben gleich \(\frac{R}{S}\) gefundenen Verhältnisse und ausserdem den vierten Werth
\[
-\frac{M}{\Sigma\pm
\frac{\partial\Phi_1}{\partial x_1}
\frac{\partial\Phi_2}{\partial x_2}
\frac{\partial\Phi_3}{\partial x_3}}.
\]

Die Curve \(R=0\) ist vom Grade \(\nu s\), die Curve \(S=0\) vom Grade \(n+3s-3\). Die Anzahl der Schnittpuncte von \(R=0\), \(f=0\), welche ausser denen von \(S=0\), \(f=0\) noch existiren, ist also
\[
n[\nu s-n-3s+3].
\]
Indessen auch diese Schnittpuncte liefern nicht sämmtlich solche Puncte \(x\), welche Doppel- oder Rückkehrpuncten entsprechen, sondern man hat noch diejenigen Werthsysteme abzusondern, welche mit den Coordinaten der gemeinsamen Schnittpuncte der Curven \(\Phi=0\) mit \(f=0\) übereinstimmen, da für diese die Gleichungen (5) überhaupt nichts aussagen. Nun wird für die Puncte \(\sigma,\tau\) offenbar \(R\) von der \((\nu-1)^{\text{ten}}\) Ordnung Null; man kann ferner zeigen, dass \(S\) in den Puncten \(\sigma\) von der zweiten, in den Puncten \(\tau\) von der dritten Ordnung Null wird. Man sieht dann, dass \(\frac{R}{S}\) für die erstern in der \((\nu-3)^{\text{ten}}\), für die andern in der \((\nu-4)^{\text{ten}}\) Ordnung verschwindet. Nimmt man hinzu, dass die erstern einfache, die letztern Doppelpuncte von \(f=0\) sind, so sieht man, dass jeder der erstern \((\nu-3)\), jeder der letztern \(2(\nu-4)\) Schnittpuncte absorbirt. Man hat also von der oben gegebenen Zahl noch
\[
(\nu-3)\sigma+2(\nu-4)\tau
\]
abzuziehen, erhält also
\[
n[\nu s-n-3s+3]-(\nu-3)\sigma-2(\nu-4)\tau
=(\nu-1)(\nu-2)-(n-1)(n-2)+2\tau
\]
brauchbare Schnittpuncte. Von diesen entsprechen zusammen immer zwei einem Doppel- oder Rückkehrpunct von \(F=0\); daher hat man
\[
\delta''+\varrho''
=
\frac{(\nu-1)(\nu-2)}{2}
-\frac{(n-1)(n-2)}{2}
+\tau,
\]
\[
\delta+\varrho
=\delta'+\varrho'+\delta''+\varrho''
=
\frac{(\nu-1)(\nu-2)}{2}
-\frac{(n-1)(n-2)}{2}
+d+r,
\]
oder endlich
\[
\frac{(\nu-1)(\nu-2)}{2}-\delta-\varrho
=
\frac{(n-1)(n-2)}{2}-d-r,
\]
was zu beweisen war.

\subsection*{\S\,17. Beweis des Hülfssatzes.}

Es bleibt übrig, die Sätze nachzuweisen, welche hier benutzt sind, und welche auf folgenden hinauskommen:

\begin{quote}
Wenn man drei Functionen \(\varphi,\psi,f\) hat, deren erste beide von gleicher Ordnung sind, so verschwindet ihre Functionaldeterminante in der zweiten Ordnung für jeden etwa den Curven \(\varphi=0\), \(\psi=0\), \(f=0\) gemeinsamen Punct; und in der dritten, wenn dieser Punct ein Doppel- oder Rückkehrpunct von \(f=0\) ist.
\end{quote}

Dieser Satz kann geometrisch so ausgesprochen werden, wenn man die Curve
\[
\Sigma\pm
\frac{\partial\varphi}{\partial x_1}
\frac{\partial\psi}{\partial x_2}
\frac{\partial f}{\partial x_3}
\]
dabei als Jacobi'sche Curve\footnote{Cremona, Introd. ad una teor. geom. delle curve piane.} bezeichnet:

\begin{quote}
Gemeinsame Schnittpuncte von \(\varphi=0\), \(\psi=0\), \(f=0\) liegen auch auf der Jacobi'schen Curve; diese berührt in denselben die Curve \(f=0\), sobald \(\varphi,\psi\) von gleichem Grade sind, und hat, wenn \(f\) einen solchen Punct zum Doppelpunct hat, diesen gleichfalls zum Doppelpunct, so dass ihre Tangenten mit denen von \(f=0\) übereinstimmen.
\end{quote}

Ist nämlich \(\Delta=0\) die Gleichung der Jacobi'schen Curve, \(x\) einer der betrachteten Puncte, und setzt man in der Nähe dieses Punctes für \(x_i\) die Ausdrücke \(x_i+\lambda\xi_i\) ein, wo denn \(\lambda\) sehr klein; so erhält man statt \(\Delta=0\) die Gleichung
\[
\Delta+\lambda\sum \xi_i\frac{\partial\Delta}{\partial x_i}
+\frac{\lambda^2}{1\cdot2}\sum\sum \xi_i\xi_k
\frac{\partial^2\Delta}{\partial x_i\partial x_k}
+\cdots=0,
\]
und an Stelle von \(f=0\) die Gleichung
\[
f+\lambda\sum \xi_i\frac{\partial f}{\partial x_i}
+\frac{\lambda^2}{1\cdot2}\sum\sum \xi_i\xi_k
\frac{\partial^2f}{\partial x_i\partial x_k}
+\cdots=0.
\]
Ist nun \(\Delta=0\), \(f=0\), so haben beide Gleichungen den Schnittpunct \(x\) gemein. Stimmen aber auch noch die Glieder erster Ordnung überein, so dass
\[
\frac{\partial\Delta}{\partial x_i}
=m\frac{\partial f}{\partial x_i},
\]
so kann man mit Hülfe der zweiten Gleichung die Glieder erster Ordnung in der ersten zerstören, die erste Gleichung beginnt mit \(\lambda^2\), und indicirt also zwei in \(x\) zusammenfallende Schnittpuncte. Haben endlich beide Curven in \(x\) Doppelpuncte, so sind
\[
\frac{\partial\Delta}{\partial x_i},\qquad
\frac{\partial f}{\partial x_i}
\]
gleich Null; hat man dann auch
\[
\frac{\partial^2\Delta}{\partial x_i\partial x_k}
=m\frac{\partial^2f}{\partial x_i\partial x_k},
\]
fallen also die Richtungen der Tangenten des Doppelpuncts für beide Curven zusammen, so kann man die erste Gleichung auf die Terme dritter Ordnung reduciren, während die zweite mit Termen zweiter Ordnung beginnt, und der Doppelpunct absorbirt also 6 Schnittpuncte.

Seien, um den obigen Satz zu beweisen, \(\varphi,\psi\) von der \(\varrho^{\text{ten}}\), \(f\) von der \(n^{\text{ten}}\) Ordnung. Ferner seien \(c_1,c_2,c_3\) beliebige Constanten. Dann ist
\[
\tag{1}
\begin{aligned}
&\Delta\,(c_1x_1+c_2x_2+c_3x_3)\\
&=
\begin{vmatrix}
\frac{\partial\varphi}{\partial x_1}&\frac{\partial\varphi}{\partial x_2}&\frac{\partial\varphi}{\partial x_3}&0\\
\frac{\partial\psi}{\partial x_1}&\frac{\partial\psi}{\partial x_2}&\frac{\partial\psi}{\partial x_3}&0\\
\frac{\partial f}{\partial x_1}&\frac{\partial f}{\partial x_2}&\frac{\partial f}{\partial x_3}&0\\
c_1&c_2&c_3&c_1x_1+c_2x_2+c_3x_3
\end{vmatrix}\\
&=
\begin{vmatrix}
\frac{\partial\varphi}{\partial x_1}&\frac{\partial\varphi}{\partial x_2}&\frac{\partial\varphi}{\partial x_3}&-\varrho\varphi\\
\frac{\partial\psi}{\partial x_1}&\frac{\partial\psi}{\partial x_2}&\frac{\partial\psi}{\partial x_3}&-\varrho\psi\\
\frac{\partial f}{\partial x_1}&\frac{\partial f}{\partial x_2}&\frac{\partial f}{\partial x_3}&-nf\\
c_1&c_2&c_3&0
\end{vmatrix}.
\end{aligned}
\]
Hieraus sieht man, dass \(\Delta\) mit \(\varphi,\psi,f\) zusammen verschwindet. Benutzt man dies, und differenzirt nach \(x_i\), so hat man für diejenigen Werthe der \(x\), in denen \(\varphi,\psi,f,\Delta\) gleich Null sind:
\[
\begin{aligned}
&\frac{\partial\Delta}{\partial x_i}(c_1x_1+c_2x_2+c_3x_3)\\
&=
\begin{vmatrix}
\frac{\partial\varphi}{\partial x_1}&\frac{\partial\varphi}{\partial x_2}&\frac{\partial\varphi}{\partial x_3}&-\varrho\frac{\partial\varphi}{\partial x_i}\\
\frac{\partial\psi}{\partial x_1}&\frac{\partial\psi}{\partial x_2}&\frac{\partial\psi}{\partial x_3}&-\varrho\frac{\partial\psi}{\partial x_i}\\
\frac{\partial f}{\partial x_1}&\frac{\partial f}{\partial x_2}&\frac{\partial f}{\partial x_3}&-n\frac{\partial f}{\partial x_i}\\
c_1&c_2&c_3&0
\end{vmatrix}\\
&=
\begin{vmatrix}
\frac{\partial\varphi}{\partial x_1}&\frac{\partial\varphi}{\partial x_2}&\frac{\partial\varphi}{\partial x_3}&0\\
\frac{\partial\psi}{\partial x_1}&\frac{\partial\psi}{\partial x_2}&\frac{\partial\psi}{\partial x_3}&0\\
\frac{\partial f}{\partial x_1}&\frac{\partial f}{\partial x_2}&\frac{\partial f}{\partial x_3}&(\varrho-n)\frac{\partial f}{\partial x_i}\\
c_1&c_2&c_3&\varrho c_i
\end{vmatrix}\\
&=-(\varrho-n)\frac{\partial f}{\partial x_i}\,
 \Sigma\pm
 \frac{\partial\varphi}{\partial x_1}
 \frac{\partial\psi}{\partial x_2}c_3,
\end{aligned}
\]
d. h. \(\Delta=0\) und \(f=0\) berühren sich.

Ist aber \(x\) ein Doppelpunct von \(f=0\), so verschwindet die rechte Seite, also auch die linke, d. h. \(\Delta=0\) hat den Punct \(x\) gleichfalls zum Doppelpunct. Differenzirt man also jetzt (1) zweimal, nach \(x_i\) und nach \(x_k\), so kann man zur Vereinfachung bemerken, dass, weil hier gleichzeitig
\[
(n-1)\frac{\partial f}{\partial x_i}
=x_1\frac{\partial^2 f}{\partial x_1\partial x_i}
+x_2\frac{\partial^2 f}{\partial x_2\partial x_i}
+x_3\frac{\partial^2 f}{\partial x_3\partial x_i}=0,
\]
\[
\varrho\varphi
=x_1\frac{\partial\varphi}{\partial x_1}
+x_2\frac{\partial\varphi}{\partial x_2}
+x_3\frac{\partial\varphi}{\partial x_3}=0,
\]
\[
\varrho\psi
=x_1\frac{\partial\psi}{\partial x_1}
+x_2\frac{\partial\psi}{\partial x_2}
+x_3\frac{\partial\psi}{\partial x_3}=0,
\]
auch die Determinante
\[
\begin{vmatrix}
\frac{\partial^2 f}{\partial x_1\partial x_i}&
\frac{\partial^2 f}{\partial x_2\partial x_i}&
\frac{\partial^2 f}{\partial x_3\partial x_i}\\
\frac{\partial\varphi}{\partial x_1}&
\frac{\partial\varphi}{\partial x_2}&
\frac{\partial\varphi}{\partial x_3}\\
\frac{\partial\psi}{\partial x_1}&
\frac{\partial\psi}{\partial x_2}&
\frac{\partial\psi}{\partial x_3}
\end{vmatrix}
\]
verschwindet. Nun braucht man in der transformirten Determinante (1) diejenigen Glieder gar nicht zu berücksichtigen, in denen die dritte Horizontalreihe nicht differenzirt wird, weil diese dann ganz verschwindet. Differenzirt man die dritte Reihe einmal, und die erste oder zweite auch einmal, so kommen Determinanten heraus, in denen die letzte Verticalreihe nur ein von Null verschiedenes Glied enthält. Zerstört man dasselbe mit Hülfe der übrigen Verticalreihen, so erhält man nur das letzte Glied der letzten Verticalreihe von Null verschieden, und zwar ist es mit einem der Ausdrücke multiplicirt, auf deren Verschwinden wir soeben aufmerksam gemacht haben. Es bleibt also, indem wir nur die dritte Horizontalreihe zweimal differenziren:
\[
\frac{\partial^2\Delta}{\partial x_i\partial x_k}
(c_1x_1+c_2x_2+c_3x_3)
=-(\varrho-n)\frac{\partial^2f}{\partial x_i\partial x_k}
\,\Sigma\pm
\frac{\partial\varphi}{\partial x_1}
\frac{\partial\psi}{\partial x_2}c_3,
\]
d. h. die Tangenten des Doppelpuncts von \(\Delta\) fallen mit denen des Doppelpuncts von \(f\) zusammen.

\subsection*{\S\,18. Erniedrigung der Curvenordnung. Normalcurven.}

Wir wollen jetzt die Frage behandeln, welches der niedrigste Grad \(\nu\) einer Curve \(F=0\) ist, welche wir durch Transformation aus einer gegebenen \(f=0\) erhalten können. Nehmen wir zunächst an, die Ordnung \(s\) der Ausdrücke \(\Phi\), welche bei der Transformation den neuen Variabeln gleich gesetzt werden, sei grösser als \(n-3\). In diesem Falle sind (vgl. \S\,9.) die Puncte, in denen eine Curve \(\Phi=0\) die Curve \(f=0\) schneidet, und auf welche es, da \(\Phi\) nur mit \(f=0\) verbunden auftritt, allein ankommt, gegeben durch
\[
ns-\frac{(n-1)(n-2)}{2}=ns-p-d-r
\]
unter ihnen. Die Zahl der allen \(\Phi\) mit \(f=0\) gemeinschaftlichen Puncte \(\sigma+\tau\) muss daher wenigstens um 2 kleiner sein als diese Zahl. Wäre sie nämlich nur um 1 kleiner, so könnte man zunächst mit Hülfe von \(f=0\) in einem \(\Phi\) alle Coefficienten bis auf \(ns-p-d-r+1\) zerstören, und die übrigen hätten dann den \(ns-p-d-r-1\) linearen Gleichungen zu genügen, welche die gegebenen Puncte mit sich führten. Die Function \(\Phi\) wäre dadurch bis auf einen einzigen Coefficienten \(\lambda\) bestimmt, die drei \(\Phi\) hätten alle Formen \(u+\lambda_1v\), \(u+\lambda_2v\), \(u+\lambda_3v\), und es bestände also eine lineare Relation zwischen ihnen, was nicht sein darf. Man muss also haben:
\[
\sigma+\tau\le ns-p-d-r-2.
\]
Nehmen wir die Gleichung hinzu:
\[
\nu=ns-\sigma-2\tau,
\]
so folgt
\[
\nu\ge p+d+r+2-\tau.
\]

Der kleinste Werth also, den die Ordnung \(\nu\) der transformirten Curve erhalten kann, entsteht, wenn \(\tau=d+r\) und \(\nu=p+2\). Dieser Werth ist, wie man sieht, von der Ordnung der Transformationsgleichungen ganz unabhängig; sämmtliche Curven \(\Phi=0\) gehen durch sämmtliche Doppelpuncte der Curve \(f=0\), und ausserdem durch
\[
\sigma=ns-\nu-2\tau=ns-p-2-2(d+r)
\]
\[
=ns-\frac{(n-1)(n-2)}{2}-2-d-r
\]
feste gemeinsame Puncte. Es ist zur Möglichkeit der Transformation noch erforderlich, dass \(\sigma\) Null oder positiv sei. Nun ist \(s>n-3\), daher auch
\[
\sigma>\frac{n(n-3)}{2}-3-d-r=p-4.
\]
Daher ist \(\sigma\) in der That Null oder positiv, sobald \(p>2\). Nur die Fälle \(p=0\), \(p=1\), \(p=2\) bedürfen einer besonderen Behandlung.

Wir sehen, dass wir mit Hülfe von Substitutionen, deren Ordnung die von \(n-3\) übersteigt, niedrigstens Curven der \((p+2)^{\text{ten}}\) Ordnung erhalten. Untersuchen wir nun Substitutionen von niedrigerer Ordnung.

Ist die Ordnung der Substitutionen
\[
s=n-3-k,\qquad (k=0,1\ldots,n-4),
\]
so ist erstlich wieder
\[
\nu=ns-\sigma-2\tau=n(n-3-k)-\sigma-2\tau.
\]
Ferner muss die Anzahl \(\sigma+\tau\) der gemeinschaftlichen Puncte wenigstens um 2 kleiner sein, als diejenige Zahl von Puncten, welche eine Curve \(s^{\text{ter}}\) Ordnung bestimmt, d. h. man muss haben
\[
\sigma+\tau\le \frac{(n-3-k)(n-k)}{2}-2.
\]
Dies mit der Gleichung
\[
\sigma+2\tau=n(n-3-k)-\nu
\]
verbunden, giebt die Ungleichung:
\[
\tau\ge \frac{(n-3-k)(n+k)}{2}-\nu+2,
\]
oder
\[
\tau\ge p+1+d+r-\frac{k\cdot k+3}{2}-\nu,
\]
oder endlich
\[
\nu\ge p+1+(d+r-\tau)-\frac{k\cdot k+3}{2}.
\]
Hieraus sieht man, dass für ein gegebenes \(k\) oder \(s\) der günstigste Fall eintritt, wenn
\[
\nu=p+1+(d+r-\tau)-\frac{k\cdot k+3}{2},
\]
und es wird dann zugleich
\[
\sigma=p-3-kn+(d+r-\tau)+\frac{k\cdot k+3}{2}.
\]

Es ist also immer möglich, die Curve \(f=0\) auf eine Curve \((p+1)^{\text{ter}}\) Ordnung zurückzuführen; und zwar wird dies geleistet durch eine Transformation \((n-3)^{\text{ter}}\) Ordnung \((k=0)\), bei welcher alle Curven \(\Phi\) durch sämmtliche Doppel- und Rückkehrpuncte gehen \((\tau=d+r)\) und ausserdem durch \(p-3\) beliebig auf \(f=0\) gewählte Puncte \(\sigma\).

Auch hier sind natürlich wie oben die Fälle \(p=0,1,2\) auszuschliessen und besonders zu behandeln.

Diese Transformation ist allgemein anwendbar, indem sie eintritt, wie gross auch die Anzahl der Doppelpuncte, welche die Curve enthält.

Aus den auf p. 64 aufgestellten Ungleichungen geht hervor, dass immer
\[
\tau\le \frac{(n-3-k)(n-k)}{2}-2,
\]
\[
\tau\ge \frac{(n-3-k)(n+k)}{2}+2-\nu.
\]
Man schliesst hieraus
\[
\frac{(n-3-k)(n-k)}{2}-2
\ge
\frac{(n-3-k)(n+k)}{2}+2-\nu
\]
oder
\[
\nu\ge 4+k(n-k-3),
\]
oder endlich
\[
\nu-n\ge (k-1)(n-k-4).
\]
Da nun \(k\) nur die Werthe \(0,1,\ldots,n-4\) annehmen kann, so ist die rechte Seite dieser Ungleichung Null oder positiv, also \(\nu\ge n\), ausser wenn \(k=0\). Man sieht daher, dass eine Erniedrigung der Curvenordnung durch eine Transformation im Allgemeinen (d. h. ohne besondere Voraussetzungen über die Lage der Doppel- und Rückkehrpuncte) nur entstehen kann bei \(k=0\).

Wenn also überhaupt der Grad der transformirten Curve niedriger als \(n\) sein soll, so muss dies durch die oben erwähnte allgemeine Transformation geschehen, welche von der \((n-3)^{\text{ten}}\) Ordnung ist, und auf den Grad \(p+1\) führt. Man muss dann haben
\[
p+1<n,
\]
also für die Zahl der Doppelpuncte die Bestimmung
\[
d+r>\frac{(n-1)(n-4)}{2}.
\]
Da die transformirte Curve von der \((p+1)^{\text{ten}}\) Ordnung wird, wenn man eine Transformation \((n-3)^{\text{ter}}\) Ordnung in der angegebenen Weise anwendet, und da die Zahl \(p\) erhalten bleibt, so hat die neue Curve
\[
\frac{p(p-1)}{2}-p=\frac{p(p-3)}{2}
\]
Doppelpuncte. Die Curven \((p+1)^{\text{ter}}\) Ordnung mit \(\frac{p(p-3)}{2}\) Doppelpuncten bilden also gewissermassen die Normalform für alle Curven, welche einer gewissen Zahl \(p\) zugehören, indem alle solche sich in diese transformiren lassen. Dagegen sind umgekehrt nach dem Vorigen diese Curven im Allgemeinen in keine niedern transformirbar, und alle Curven niederer Ordnung, denen dasselbe \(p\) zukommt, können nur aus speciellen Curven \((p+1)^{\text{ter}}\) Ordnung mit \(\frac{p(p-3)}{2}\) Doppelpuncten hervorgehn, welche wegen ihrer besondern Natur die Erniedrigung der Ordnung ausnahmsweise gestatten. Unsere Normalcurven sind genauer aufgezählt von \(p=3\) ab folgende:
\[
\begin{array}{ll}
p=3:& \text{Curven }4^{\text{ter}}\text{ Ordnung mit }0\text{ Doppelpuncten},\\
p=4:& \text{Curven }5^{\text{ter}}\text{ Ordnung mit }2\text{ Doppelpuncten},\\
p=5:& \text{Curven }6^{\text{ter}}\text{ Ordnung mit }5\text{ Doppelpuncten},\\
p=6:& \text{Curven }7^{\text{ter}}\text{ Ordnung mit }9\text{ Doppelpuncten},\\
p=7:& \text{Curven }8^{\text{ter}}\text{ Ordnung mit }14\text{ Doppelpuncten},\\
&\text{u. s. w.}
\end{array}
\]
Die allgemeine Curve \(5^{\text{ter}}\) Ordnung erscheint dabei als Transformation einer speciellen Curve \(7^{\text{ter}}\) Ordnung mit 9 Doppelpuncten, die Curve \(5^{\text{ter}}\) Ordnung mit einem Doppelpunct Transformation einer speciellen Curve \(6^{\text{ter}}\) Ordnung mit 5 Doppelpuncten etc.

% Continuation: source PDF pp. 089--106.

\subsection*{\S\,19. Behandlung des Falles \(p=0\). Normalcurve: eine Gerade.}

Die Fälle, welche nun noch eine besondre Untersuchung erheischen, sind diejenigen, wo \(p=0,1,2\). Sie sollen der Reihe nach behandelt werden. Die einfachsten Curven dieser Art sind beziehungsweise die Gerade, die Curve \(3^{\text{ter}}\) Ordnung ohne Doppelpunct, und die Curve \(4^{\text{ter}}\) Ordnung mit einem Doppelpunct. Es soll gezeigt werden, wie die allgemeinsten Curven, für welche \(p=0,1,2\) sich beziehungsweise in diese Curven transformiren lassen.

Wenn für eine Curve \(n^{\text{ter}}\) Ordnung \(p=0\), so besitzt dieselbe
\[
\frac{(n-1)(n-2)}{2}
\]
Doppel- und Rückkehrpuncte. Betrachten wir diese zusammen mit \(2n-3\) andern auf der Curve gewählten Puncten als Grundpuncte eines Büschels von Curven \((n-1)^{\text{ter}}\) Ordnung. Derselbe ist dadurch vollständig bestimmt; denn durch die
\[
\frac{(n-1)(n+2)}{2}-1
\]
festen Puncte sind die Constanten einer Curve \((n-1)^{\text{ter}}\) Ordnung bis auf eine bestimmt. Sei diese eine, willkürlich gebliebene, \(\lambda\), so ist die Gleichung des Büschels von der Form
\[
(1)\ldots\quad u+\lambda v=0,
\]
wo \(u\) und \(v\) Functionen \((n-1)^{\text{ter}}\) Ordnung sind, welche gleich Null gesetzt irgend zwei Curven des Büschels darstellen.

Verbinden wir die Gleichung (1) mit der Gleichung \(f=0\) der gegebenen Curve, so erhalten wir die \(n\cdot(n-1)\) Puncte, in welchen eine Curve des Büschels die gegebene schneidet. Aber diese bestehen aus den \(\frac{(n-1)(n-2)}{2}\) doppelt zu rechnenden Doppel- und Rückkehrpuncten, aus den \(2n-3\) willkürlich gewählten festen Puncten, und aus
\[
n(n-1)-(n-1)(n-2)-(2n-3)=1
\]
beweglichen, d.\ h. von \(\lambda\) abhängigen Schnittpuncten. Verbinden wir die Gleichungen
\[
f=0,\qquad u+\lambda v=0
\]
mit der Gleichung
\[
\alpha=\alpha_1x_1+\alpha_2x_2+\alpha_3x_3=0,
\]
in welcher die \(\alpha\) Coordinaten einer beliebigen Geraden bedeuten, und eliminiren wir aus diesen Gleichungen die \(x\), so entsteht eine Gleichung für die \(\alpha\), welche aussagt, dass die Gerade \(\alpha\) durch einen Schnittpunct der beiden Curven gehe, d.\ h. es entsteht das Product der \(n\cdot(n-1)\) Gleichungen der einzelnen Schnittpuncte. Die Endgleichung der Elimination ist, da die Grade der gegebenen Gleichungen in den Variabeln durch die Tabelle
\[
\begin{array}{c|ccc}
 & x & \alpha & \lambda\\ \hline
f=0 & n & 0 & 0\\
u+\lambda v=0 & n-1 & 0 & 1\\
\alpha=0 & 1 & 1 & 0
\end{array}
\]
dargestellt sind, in der That vom Grade \(n\cdot(n-1)\) in den \(\alpha\) und ausserdem vom Grade \(n\) in \(\lambda\). Diese Endgleichung muss sich in \(n\cdot(n-1)\) lineare Factoren auflösen, von denen nur einer \(\lambda\) enthält, der nämlich, welcher die Gleichung des beweglichen Schnittpuncts angiebt. Es muss sich also ein Factor vom \([n(n-1)-1]^{\text{ten}}\) Grade in den \(\alpha\) absondern lassen, welcher \(\lambda\) nicht enthält. Nach Absonderung desselben hat die Eliminationsgleichung die Form
\[
A_1\alpha_1+A_2\alpha_2+A_3\alpha_3=0,
\]
wo die \(A\) Functionen \(n^{\text{ten}}\) Grades von \(\lambda\) sind. Dieses ist die Gleichung des beweglichen Schnittpuncts; seine Coordinaten verhalten sich also zu einander wie die Functionen \(A\), und man hat also die Coordinaten \(x\) des beweglichen Schnittpuncts als rationale Functionen \(n^{\text{ter}}\) Ordnung eines Parameters \(\lambda\) dargestellt mit Hülfe der Gleichungen:
\[
(1)\ldots\quad
\begin{gathered}
\mu x_1=A_1,\\
\mu x_2=A_2,\\
\mu x_3=A_3,
\end{gathered}
\]
wobei \(\mu\) ein willkürlicher Factor ist.

Setzen wir \(\lambda=\frac{y_1}{y_2}\), so haben wir die \(x\) als homogene Functionen \(n^{\text{ter}}\) Ordnung von \(y_1,y_2\) dargestellt. Sehen wir nun \(y_1,y_2,y_3\) als Coordinaten eines Puncts der transformirten Curve an, so können wir die \(x\) als Functionen der \(y\) betrachten, bei welchen
\[
y_3=0;
\]
es ist also dieses die Gleichung der transformirten Curve, d.\ h. dieselbe ist eine Gerade, wie es sein sollte.

Wir können so jede Curve \(n^{\text{ter}}\) Ordnung mit \(\frac{(n-1)(n-2)}{2}\) Doppel- und Rückkehrpuncten in eine Gerade überführen mit Hülfe der Transformation (1), welche von der \(n^{\text{ten}}\) Ordnung ist.

\subsection*{\S\,20. Behandlung des Falles \(p=1\). Normalcurve: eine allgemeine Curve dritter Ordnung.}

Gehen wir zu den Curven \(n^{\text{ter}}\) Ordnung mit \(\frac{n(n-3)}{2}\) Doppel- und Rückkehrpuncten über, für welche \(p=1\). Man betrachte die Doppel- und Rückkehrpuncte nebst \(2n-2\) andern auf der Curve beliebig gewählten festen Puncten als die
\[
\frac{(n-1)(n+2)}{2}-1
\]
Grundpuncte eines Curvenbüschels \(u+\lambda v=0\) der \((n-1)^{\text{ten}}\) Ordnung. Verfährt man wieder ganz wie im Vorigen, so hat man die Eliminationsresultante aus den Gleichungen
\[
f=0,\qquad u+\lambda v=0,\qquad \alpha=0
\]
zu bilden, welche wie im Vorigen für die \(\alpha\) vom \(n(n-1)^{\text{ten}}\), für \(\lambda\) vom \(n^{\text{ten}}\) Grade ist, und sich als das Product von \(n(n-1)\) in den \(\alpha\) linearen Factoren, den Gleichungen der \(n(n-1)\) Schnittpuncte von \(f=0\) mit \(u+\lambda v=0\) darstellt. Von diesen fallen jetzt \(n(n-3)\) in die Doppel- und Rückkehrpuncte, \(2n-2\) in die beliebig gewählten festen Puncte hinein, es bleiben also
\[
n(n-1)-n(n-3)-(2n-2)=2
\]
bewegliche Schnittpuncte übrig. Die Resultante geht daher nach Absonderung des Factors \([n(n-1)-2]^{\text{ten}}\) Grades, der \(\lambda\) nicht enthält, in die Form über
\[
(1)\ldots\quad
0=w_{11}\alpha_1^2+w_{22}\alpha_2^2+w_{33}\alpha_3^2
+2w_{23}\alpha_2\alpha_3+2w_{31}\alpha_3\alpha_1+2w_{12}\alpha_1\alpha_2,
\]
wo die \(w\) Functionen \(n^{\text{ten}}\) Grades in \(\lambda\) sind. Da diese Gleichung das Product der Gleichungen der beiden beweglichen Schnittpuncte darstellt, so muss sie sich in zwei lineare Factoren auflösen lassen, d.\ h. es muss unabhängig von dem Werthe von \(\lambda\)
\[
(2)\ldots\quad
W=\begin{vmatrix}
w_{11}&w_{12}&w_{13}\\
w_{21}&w_{22}&w_{23}\\
w_{31}&w_{32}&w_{33}
\end{vmatrix}=0
\]
sein. Sind ferner \(x_1,x_2,x_3\) und \(\xi_1,\xi_2,\xi_3\) die Coordinaten der beweglichen Schnittpuncte, so muss die obige Gleichung mit
\[
(3)\ldots\quad
0=(x_1\alpha_1+x_2\alpha_2+x_3\alpha_3)(\xi_1\alpha_1+\xi_2\alpha_2+\xi_3\alpha_3)
\]
gleichbedeutend sein, d.\ h. man hat zur Bestimmung der \(x,\xi\) die Gleichungen:
\[
(4)\ldots\quad
\begin{array}{lll}
x_1\xi_1=w_{11},&
x_2\xi_3+x_3\xi_2=2w_{23},\\
x_2\xi_2=w_{22},&
x_3\xi_1+x_1\xi_3=2w_{31},\\
x_3\xi_3=w_{33},&
x_1\xi_2+x_2\xi_1=2w_{12}.
\end{array}
\]
Bezeichnet man sodann die Coordinaten der Verbindungslinie der Puncte \(x,\xi\) mit \(p_1,p_2,p_3\), so dass
\[
(5)\ldots\quad
\begin{gathered}
p_1=x_2\xi_3-x_3\xi_2,\\
p_2=x_3\xi_1-x_1\xi_3,\\
p_3=x_1\xi_2-x_2\xi_1,
\end{gathered}
\]
so findet man aus (4) sofort
\[
p_ip_k=-4W_{ik},
\]
wo die \(W_{ik}\) die Unterdeterminanten von \(W\) sind. Die aus den \(W_{ik}\) zu bildenden Unterdeterminanten sind also sämmtlich identisch Null, so dass für jeden Werth von \(\lambda\):
\[
\begin{array}{lll}
W_{22}W_{33}=W_{23}^{2},& W_{12}W_{13}=W_{11}W_{23},\\
W_{33}W_{11}=W_{31}^{2},& W_{23}W_{21}=W_{22}W_{31},\\
W_{11}W_{22}=W_{12}^{2},& W_{31}W_{32}=W_{33}W_{12}.
\end{array}
\]
Dies ist nicht anders möglich, als wenn die \(W_{ik}\) bis auf einen allen gemeinschaftlichen Factor, den wir durch \(-Q\) bezeichnen wollen, den Quadraten und Producten rationaler Functionen \(l_1,l_2,l_3\) gleich werden, wenn also
\[
W_{ik}=-l_il_kQ.
\]
Es folgt dann
\[
p_1=2l_1\sqrt Q,\qquad p_2=2l_2\sqrt Q,\qquad p_3=2l_3\sqrt Q,
\]
und ferner aus den Gleichungen (4), (5) zusammen:
\[
(6)\ldots\quad
\begin{array}{lll}
x_1\xi_1=w_{11},&x_2\xi_1=w_{21}-l_3\sqrt Q,&x_3\xi_1=w_{31}+l_2\sqrt Q,\\
x_1\xi_2=w_{12}+l_3\sqrt Q,&x_2\xi_2=w_{22},&x_3\xi_2=w_{32}-l_1\sqrt Q,\\
x_1\xi_3=w_{13}-l_2\sqrt Q,&x_2\xi_3=w_{23}+l_1\sqrt Q,&x_3\xi_3=w_{33}.
\end{array}
\]
Statt dieser Gleichungen kann man auch folgende setzen, welche entstehen, wenn man diese Reihen, mit drei willkürlichen Grössen \(\gamma_1,\gamma_2,\gamma_3\) multiplicirt, addirt, und \(\varrho=\xi_1\gamma_1+\xi_2\gamma_2+\xi_3\gamma_3\) als willkürlichen Factor betrachtet:
\[
(7)\ldots\quad
\begin{aligned}
\varrho x_1&=w_{11}\gamma_1+w_{12}\gamma_2+w_{13}\gamma_3+(\gamma_2l_3-\gamma_3l_2)\sqrt Q,\\
\varrho x_2&=w_{21}\gamma_1+w_{22}\gamma_2+w_{23}\gamma_3+(\gamma_3l_1-\gamma_1l_3)\sqrt Q,\\
\varrho x_3&=w_{31}\gamma_1+w_{32}\gamma_2+w_{33}\gamma_3+(\gamma_1l_2-\gamma_2l_1)\sqrt Q.
\end{aligned}
\]
Diese Ausdrücke sind von den Grössen \(\gamma\) nur formell abhängig, wie aus den Gleichungen (6) hervorgeht. Lässt man das Zeichen von \(\sqrt Q\) unbestimmt, so geben diese Gleichungen ebensowohl die \(x\) als die \(\xi\), indem man der Wurzel nur einmal das positive, einmal das negative Zeichen zu geben braucht.

Es kommt alles darauf an, den Grad der Function \(Q\) zu bestimmen. Dies geschieht durch folgende Betrachtung. Es giebt Curven des Büschels \(u+\lambda v=0\), welche die Curve \(f=0\) berühren, indem die beiden beweglichen Schnittpuncte zusammenfallen. Für diese muss offenbar \(Q\) verschwinden. Untersuchen wir also, wie viel solcher Curven es giebt. Da \(f=0\) von einer Curve \(u+\lambda v=0\) berührt werden soll, so hat man für den Berührungspunct:
\[
\varrho\frac{\partial f}{\partial x_1}=\frac{\partial u}{\partial x_1}+\lambda\frac{\partial v}{\partial x_1},\qquad
\varrho\frac{\partial f}{\partial x_2}=\frac{\partial u}{\partial x_2}+\lambda\frac{\partial v}{\partial x_2},\qquad
\varrho\frac{\partial f}{\partial x_3}=\frac{\partial u}{\partial x_3}+\lambda\frac{\partial v}{\partial x_3}.
\]
Daraus folgt
\[
0=\Sigma\pm \frac{\partial f}{\partial x_1}\frac{\partial u}{\partial x_2}\frac{\partial v}{\partial x_3};
\]
die Berührungspuncte sind also Schnittpuncte der Jacobi'schen Curve mit \(f=0\). Aber nicht alle diese Schnittpuncte sind wesentlich. Denn nach dem Frühern fallen in jeden einfachen Punct, welchen \(u,v\) mit \(f\) gemein haben, 2 Schnittpuncte der Jacobi'schen Curve mit \(f=0\), und sogar 6, wenn dieser Punct ein Doppelpunct von \(f\) war. Nun ist der Grad der Jacobi'schen Curve \(3n-5\), \(u,v\) haben \(\frac{n(n-3)}{2}\) Doppelpuncte, und \(2n-2\) einfache Puncte mit \(f\) gemein; also wird die Anzahl der wesentlichen Schnittpuncte, d.\ h. der Grad von \(Q\):
\[
n(3n-5)-2(2n-2)-3n(n-3)=4.
\]
Die Function \(Q\) ist also vom vierten Grade, die \(l\) werden vom \((n-2)^{\text{ten}}\) in \(\lambda\).

Schreiben wir nun \(Q\) in die Form
\[
Q=(\lambda-a_1)(\lambda-a_2)(\lambda-a_3)(\lambda-a_4)
\]
und setzen
\[
y_1=(\lambda-a_1)(a_3-a_2),\qquad
y_2=(\lambda-a_2)(a_3-a_1),
\]
\[
(8)\ldots\quad y_1y_3=\frac{\sqrt Q}{m},\qquad
m^2=\frac{a_4-a_2}{(a_3-a_1)(a_1-a_2)^2(a_3-a_2)^2},
\]
so lassen die \(w,l\) sich als Functionen \(w',l'\) von \(y_1,y_2\) darstellen, deren Grade den Ordnungen der \(w,l\) beziehungsweise gleich sind, \(\sqrt Q\) wird bis auf eine Constante, welche in die \(l\) eingehen mag, gleich \(y_1y_3\), und man findet die \(x\) durch die Formeln ausgedrückt:
\[
(9)\ldots\quad
\begin{aligned}
\varrho x_1&=\gamma_1w'_{11}+\gamma_2w'_{12}+\gamma_3w'_{13}
+m(\gamma_2l'_3-\gamma_3l'_2)y_1y_3,\\
\varrho x_2&=\gamma_1w'_{21}+\gamma_2w'_{22}+\gamma_3w'_{23}
+m(\gamma_3l'_1-\gamma_1l'_3)y_1y_3,\\
\varrho x_3&=\gamma_1w'_{31}+\gamma_2w'_{32}+\gamma_3w'_{33}
+m(\gamma_1l'_2-\gamma_2l'_1)y_1y_3.
\end{aligned}
\]
Die \(x\) sind also als ganze Functionen \(n^{\text{ter}}\) Ordnung der \(y\) dargestellt; und auch umgekehrt kann man die \(y\) durch die \(x\) rational ausdrücken, da \(\lambda=-\frac{u}{v}\) eine rationale Function der \(x\) ist, wodurch also \(y_1,y_2\) sofort so dargestellt sind, \(y_3\) aber in den Gleichungen (9) nur linear vorkommt. Zwischen den \(y\) endlich besteht aus (8) die Gleichung
\[
(10)\ldots\quad
F=y_3^2y_1-y_2(y_1-y_2)(y_1-k^2y_2)=0,\qquad
k^2=\frac{\frac{a_3-a_2}{a_3-a_1}}{\frac{a_4-a_2}{a_4-a_1}}.
\]
Dies ist die Gleichung der transformirten Curve. Sie ist von der dritten Ordnung, und das Coordinatensystem der \(y\) hat in Bezug auf dieselbe eine sehr bemerkenswerthe Lage. Für \(y_1=0\) geht die Gleichung der Curve in \(y_2^3=0\) über; die Axe \(y_1=0\) schneidet also die Curve in drei unendlich nahen Puncten, d.\ h. in \(y_1=0,\ y_2=0\) ist ein Wendepunct und \(y_1=0\) seine Tangente. Durch ihn gehen die drei Geraden
\[
y_2=0,\qquad y_1-y_2=0,\qquad y_1-k^2y_2=0,
\]
für welche die Gleichung der Curve immer in
\[
y_3^2y_1=0
\]
übergeht. Jene drei Geraden sind also die Tangenten, welche von dem Wendepunct an die Curve gelegt werden können, und \(y_3=0\) ist die Gerade, auf welcher die Berührungspuncte der drei Tangenten sich befinden, die zu dem Wendepunct gehörige harmonische Gerade. Es ist zugleich \(k^2\) das Doppelverhältniss der vom Wendepunct an die Curve gezogenen Tangenten, die Wendetangente selbst inbegriffen.

Die Curven \(n^{\text{ter}}\) Ordnung, welche entstehen, indem man die \(x\) gleich Null setzt, müssen mit der Curve (10) eine gewisse Anzahl \(\sigma\) fester Puncte gemein haben. Die letztere hat keine Doppelpuncte, daher muss der Grad der Curve \(f=0\), welcher, wie wir wissen, \(n\) ist, und welcher aus (10) durch die Substitutionen \(n^{\text{ter}}\) Ordnung (9) wieder hervorgeht, gleich \(3n-\sigma\) sein, oder \(\sigma=2n\). Diese \(\sigma\) Puncte sind leicht nachzuweisen. Nehmen wir irgend eine lineare Verbindung der rechten Theile von (9) und bilden also die Gleichung:
\[
(11)\ldots\quad
\Sigma\gamma_i\beta_kw'_{ik}+y_1y_3\Sigma\pm\gamma_1\beta_2l'_3=0.
\]
Die durch sie dargestellte Curve muss ebenfalls jene \(2n\) constanten Schnittpuncte mit (10) haben. Aber in der Nähe von \(y_1=0,\ y_2=0\) geht (11) in das Büschel von \(n-1\) Geraden
\[
y_1\cdot \Sigma\pm\gamma_1\beta_2l'_3=0
\]
über. Die Curve (11) hat also im Puncte \(y_1=0,\ y_2=0\) einen \((n-1)\)fachen Punct, dessen eine Tangente, \(y_1=0\), die Curve (10) berührt. Daher fallen in diesen Punct \(n\) Schnittpuncte der beiden Curven.

Ferner giebt die Gleichung (11)
\[
(12)\ldots\quad
(\Sigma\gamma_i\beta_kw'_{ik})^2
=y_1^2y_3^2(\Sigma\pm\gamma_1\beta_2l'_3)^2,
\]
wo nun \(y_3^2\) mit Hülfe von (10) eliminirt werden kann. Aber dann geht \(l'_il'_ky_1^2y_3^2\) nach dem Vorigen in die Unterdeterminante \(W'_{ik}\) der \(w'_{ik}\), mit \(-1\) multiplicirt, über, und die rechte Seite von (12) in die Determinante:
\[
-\begin{vmatrix}
W'_{11}&W'_{12}&W'_{13}&\gamma_1&\beta_1\\
W'_{21}&W'_{22}&W'_{23}&\gamma_2&\beta_2\\
W'_{31}&W'_{32}&W'_{33}&\gamma_3&\beta_3\\
\gamma_1&\gamma_2&\gamma_3&0&0\\
\beta_1&\beta_2&\beta_3&0&0
\end{vmatrix}
=(\Sigma w'_{ik}\gamma_i\beta_k)^2
-(\Sigma w'_{ik}\gamma_i\gamma_k)(\Sigma w'_{ik}\beta_i\beta_k).
\]
Die Gleichung (12) zerfällt daher in die beiden Gleichungen
\[
\Sigma w'_{ik}\gamma_i\gamma_k=0,\qquad
\Sigma w'_{ik}\beta_i\beta_k=0.
\]
Nur die zweite von diesen ist von den speciellen Constanten \(\beta\) der Curve (11) abhängig; die erste liefert \(n\) constante Verhältnisse \(y_1:y_2\), und jedem entspricht aus (11) ein Werth von \(y_3\), welcher gleichfalls von den \(\beta\) unabhängig ist, so dass man aus
\[
(13)\ldots\quad \Sigma\gamma_i\gamma_kw'_{ik}=0
\]
in Verbindung mit (10) \(n\) weitere constante Schnittpuncte findet. Was nämlich den Werth von \(y_3\) betrifft, so ist, wenn wirklich der aus (11) gefundene Werth von den \(\beta\) unabhängig sein soll, nothwendig
\[
y_3=
\frac{w'_{11}\gamma_1+w'_{12}\gamma_2+w'_{13}\gamma_3}{\gamma_3l'_2-\gamma_2l'_3}
=
\frac{w'_{21}\gamma_1+w'_{22}\gamma_2+w'_{23}\gamma_3}{\gamma_1l'_3-\gamma_3l'_1}
=
\frac{w'_{31}\gamma_1+w'_{32}\gamma_2+w'_{33}\gamma_3}{\gamma_2l'_1-\gamma_1l'_2}.
\]
Setzt man aber die Quotienten paarweise einander gleich, so hat man die drei Bedingungen
\[
\begin{aligned}
\gamma_1\Sigma w'_{ik}\gamma_il'_k-l'_1\Sigma w'_{ik}\gamma_i\gamma_k&=0,\\
\gamma_2\Sigma w'_{ik}\gamma_il'_k-l'_2\Sigma w'_{ik}\gamma_i\gamma_k&=0,\\
\gamma_3\Sigma w'_{ik}\gamma_il'_k-l'_3\Sigma w'_{ik}\gamma_i\gamma_k&=0,
\end{aligned}
\]
welche wegen (13) sich auf die eine Gleichung
\[
\Sigma w'_{ik}\gamma_il'_k=0
\]
reduciren. Diese aber geht, mit \(y_1y_3l'_h\) multiplicirt, in
\[
\sum_{ik}w'_{ik}\gamma_iW'_{kh}=W\cdot\gamma_h=0
\]
über, und ist also identisch erfüllt.

Es entsteht nun die Frage, ob der Modul \(k^2\) eine wesentliche Constante ist, oder ob man ihren Werth abändern kann, indem man irgend eine andere Transformation anwendet.

Zur Beantwortung dieser Frage\footnote{Wir verdanken diesen Beweis einer schriftlichen Mittheilung des Hrn. Cayley. Man vergleiche über denselben, so wie über verschiedenes auf die hier behandelten Gegenstände Bezügliches die Abhandlung des Hrn. Cayley, Verhandlungen der London Mathematical Society, Oct. 16, 1865.} denken wir uns die Curve \(f=0\) durch eine beliebige andre Transformation in die Curve dritter Ordnung
\[
(14)\ldots\quad z_3^2z_1=z_2(z_1-z_2)(z_1-\lambda^2z_2)
\]
verwandelt. Da die \(x\) den \(y\) wie den \(z\) eindeutig entsprechen, so entsprechen auch diese einander gegenseitig eindeutig, und man kann also die neue Transformation zusammensetzen aus unserer ursprünglichen, und aus einer Transformation, welche die Curve dritter Ordnung (10) in die Curve dritter Ordnung (14) verwandelt. Es ist die Frage, ob der Werth von \(\lambda^2\) dabei beliebig gross gemacht werden kann.

Die Grösse \(\lambda^2\) stellt das Doppelverhältniss der vom Puncte \(z_1=0,\ z_2=0\) an die Curve (14) zu legenden Tangenten dar, wie \(k^2\) das entsprechende bei der Curve \(F=0\) darstellt. Nach einem bekannten Satze aus der Theorie der Curven dritter Ordnung weiss man, dass dieses Doppelverhältniss ungeändert bleibt, von welchem Puncte der Curve aus man das Tangentenbüschel auch lege.

Wenn die von (10) zu (14) führende Transformation
\[
z_i=\varphi_i(y_1,y_2,y_3)
\]
von der \(\varrho^{\text{ten}}\) Ordnung ist, so müssen die Curven \(\varphi=0\) durch \(3\varrho-3\) feste Puncte der Curve (10) gehen, damit das Endresultat von der dritten Ordnung werde. Einer Geraden
\[
(15)\ldots\quad \alpha_1z_1+\alpha_2z_2+\alpha_3z_3=0
\]
entspricht bei dieser Transformation eine Curve \(\varrho^{\text{ter}}\) Ordnung
\[
(16)\ldots\quad \alpha_1\varphi_1+\alpha_2\varphi_2+\alpha_3\varphi_3=0,
\]
dergestalt, dass die beweglichen Schnittpuncte der einen Curve mit (14) denen der andern mit (10) entsprechen. Lässt man die Gerade (15) sich um einen Punct drehen, fixirt also einen der beweglichen Schnittpuncte, so erhält man das Linienbüschel
\[
(17)\ldots\quad Z+\lambda Z'=0,
\]
und ihm entspricht aus (16) ein Büschel von Curven \(\varrho^{\text{ter}}\) Ordnung
\[
(18)\ldots\quad Y+\lambda Y'=0,
\]
welcher entsteht, indem der entsprechende bewegliche Schnittpunct der Curve (16) fixirt wird. Die Gebilde (17) und (18) sind projectivisch, insofern dasselbe reihende Element \(\lambda\) in beiden auftritt und die entsprechenden Elemente verbindet.

Da den beiden übrigen Schnittpuncten einer Geraden des Büschels (17) immer die beiden übrigen Schnittpuncte einer Curve des Büschels (18) entsprechen, so fallen beide Punctenpaare gleichzeitig zusammen. Den 4 von dem fixirten Puncte auf (14) an diese Curve zu ziehenden Tangenten entsprechen also 4 Curven des Büschels (18), welche die Curve \(F=0\) berühren, und da die Grössen \(\lambda\) der entsprechenden Curven und Geraden dieselben sind, so ist das Doppelverhältniss der vier Berührungscurven gleich dem der vier Tangenten. Nun blieb letzteres ungeändert, wenn man den auf (14) fixirten Punct über die Curve wandern lässt. Ebenso also bleibt das Doppelverhältniss der vier Berührungscurven ungeändert, wenn man den vierten Punct über die Curve \(F=0\) wandern lässt. Und was von diesem Puncte gilt, muss ebenso von den andern Puncten des Büschels gelten. Man hat also den Satz:

\emph{Wenn man durch irgend \(3s-2\) Puncte einer Curve dritter Ordnung einen Büschel von Curven \(s^{\text{ter}}\) Ordnung legt}\footnote{Der Büschel ist hiedurch im Allgemeinen nicht vollständig bestimmt, wohl aber die Schaaren von Durchschnittspuncten, welche seine Curven mit der Curve dritter Ordnung gemein haben.}\emph{, so giebt es in demselben 4 Curven, welche die Curve dritter Ordnung berühren, und ihr Doppelverhältniss ist von der Lage jener \(3s-2\) Puncte unabhängig.}

Lässt man nun \(3s-3\) jener Puncte so liegen, dass sie den vollständigen Durchschnitt einer Curve \((s-1)^{\text{ter}}\) Ordnung mit der Curve \(3^{\text{ter}}\) Ordnung bilden, so geht der Curvenbüschel \(s^{\text{ter}}\) Ordnung in eine feste Curve \((s-1)^{\text{ter}}\) Ordnung (resp. deren Schnittpunctsystem mit der Curve \(3^{\text{ter}}\) Ordnung) über, und in einen Büschel von Geraden, welcher den letzten festen Punct zum Scheitel hat. Hieraus folgt ferner:

\emph{Jenes Doppelverhältniss ist auch von der Zahl \(s\) unabhängig, und ist gleich dem Doppelverhältniss der Tangenten, welche von einem Punct der Curve \(3^{\text{ter}}\) Ordnung an dieselbe gelegt werden können.}

Somit hat also \(k^2\) wirklich einen absolut unveränderlichen Werth, und \(\lambda^2\) muss entweder gleich \(k^2\) sein, oder doch gleich einem der fünf andern Doppelverhältnisse
\[
\left(\frac1{k^2},\ 1-k^2,\ \frac1{1-k^2},\ \frac{k^2-1}{k^2},\ \frac{k^2}{k^2-1}\right),
\]
welche sich aus denselben Elementen bilden lassen.

Diejenigen Curven, für welche \(p=1\), besitzen also eine absolute Invariante.

\subsection*{\S\,21. Behandlung des Falls \(p=2\). Normalcurve: eine Curve vierter Ordnung mit einem Doppelpunct.}

Die Reduction der Curven mit \(\frac{(n-1)(n-2)}{2}-2\) Doppel- und Rückkehrpuncten \((p=2)\) geschieht auf ganz ähnliche Weise wie bei \(p=1\). Nehmen wir einen Curvenbüschel \((n-1)^{\text{ter}}\) Ordnung, welches durch die Doppel- und Rückkehrpuncte und durch \(2n\) andere auf der Curve \(f=0\) beliebig gewählte Puncte geht, und bezeichnen denselben durch
\[
u+\lambda v=0,
\]
so haben die Curven dieses Büschels wieder mit \(f=0\) nur zwei bewegliche Schnittpuncte gemein. Eliminirt man also die \(x\) aus den Gleichungen
\[
f=0,\qquad u+\lambda v=0,\qquad \alpha_1x_1+\alpha_2x_2+\alpha_3x_3=0,
\]
so erhält man wieder mit Uebergehung der von \(\lambda\) unabhängigen Factoren das Product der Gleichungen der beiden beweglichen Schnittpuncte in der Form
\[
\Sigma w_{ik}\alpha_i\alpha_k=0,
\]
wo die \(w\) ganze Functionen \(n^{\text{ter}}\) Ordnung von \(\lambda\) sind, deren Determinante verschwindet. Man hat daher auch wie im Vorigen für die \(x\) hier in \(\lambda\) die Ausdrücke:
\[
\begin{aligned}
\varrho x_1&=w_{11}\gamma_1+w_{12}\gamma_2+w_{13}\gamma_3+(\gamma_2l_3-\gamma_3l_2)\sqrt Q,\\
\varrho x_2&=w_{21}\gamma_1+w_{22}\gamma_2+w_{23}\gamma_3+(\gamma_3l_1-\gamma_1l_3)\sqrt Q,\\
\varrho x_3&=w_{31}\gamma_1+w_{32}\gamma_2+w_{33}\gamma_3+(\gamma_1l_2-\gamma_2l_1)\sqrt Q,
\end{aligned}
\]
wo die \(\gamma\) willkürliche Grössen sind, und \(Q,l_1,l_2,l_3\) rationale Functionen von \(\lambda\), so dass
\[
l_il_kQ=-W_{ik}.
\]
Aber der Grad der Function \(Q\) ist hier ein anderer als im Vorigen. Er wird erhalten, wenn man die Anzahl der Schnittpuncte von \(f=0\) mit der Jacobi'schen Curve
\[
\Sigma\pm\frac{\partial f}{\partial x_1}\frac{\partial u}{\partial x_2}\frac{\partial v}{\partial x_3}=0
\]
vermindert um die doppelte Anzahl der festen einfachen Puncte, welche \(u,v\) mit \(f\) gemein haben, und um die 6fache Anzahl der Doppelpuncte; der Grad von \(Q\) ist also:
\[
n(3n-5)-4n-6\left(\frac{(n-1)(n-2)}{2}-2\right)=6.
\]
Setzt man also
\[
Q=(\lambda-a_1)(\lambda-a_2)(\lambda-a_3)(\lambda-a_4)(\lambda-a_5)(\lambda-a_6)
\]
und sodann
\[
y_1=(\lambda-a_1)(a_3-a_2),\qquad y_2=(\lambda-a_2)(a_3-a_1),
\]
\[
y_1y_2y_3=\frac{1}{m}\sqrt Q,\qquad
m^2=-\frac{(a_4-a_2)(a_5-a_2)(a_6-a_2)}
{(a_3-a_1)(a_3-a_1)^4(a_1-a_2)^4},
\]
so kann man die Functionen \(w,l\) als homogene Functionen \(w',l'\) von \(y_1\) und \(y_2\) darstellen, und erhält für die \(x\) die in den \(y\) rationalen Ausdrücke:
\[
(1)\ldots\quad
\begin{aligned}
\varrho x_1&=\gamma_1w'_{11}+\gamma_2w'_{12}+\gamma_3w'_{13}
+m(\gamma_2l'_3-\gamma_3l'_2)y_1y_2y_3,\\
\varrho x_2&=\gamma_1w'_{21}+\gamma_2w'_{22}+\gamma_3w'_{23}
+m(\gamma_3l'_1-\gamma_1l'_3)y_1y_2y_3,\\
\varrho x_3&=\gamma_1w'_{31}+\gamma_2w'_{32}+\gamma_3w'_{33}
+m(\gamma_1l'_2-\gamma_2l'_1)y_1y_2y_3,
\end{aligned}
\]
während zwischen den \(y\) selbst die Gleichung besteht:
\[
(2)\ldots\quad
y_1y_2y_3^2=(y_1-y_2)(y_1-k^2y_2)(y_1-\lambda^2y_2)(y_1-\mu^2y_2),
\]
wo
\[
k^2=\frac{\frac{a_3-a_2}{a_3-a_1}}{\frac{a_4-a_2}{a_4-a_1}},
\qquad
\lambda^2=\frac{\frac{a_3-a_2}{a_3-a_1}}{\frac{a_5-a_2}{a_5-a_1}},
\qquad
\mu^2=\frac{\frac{a_3-a_2}{a_3-a_1}}{\frac{a_6-a_2}{a_6-a_1}}.
\]
Die Gleichung (2) ist die Gleichung der transformirten Curve; und zwar ist es die Gleichung einer Curve \(4^{\text{ter}}\) Ordnung mit dem Doppelpunct \(y_1=0,\ y_2=0\); die Axen \(y_1=0,\ y_2=0\) selbst sind die Tangenten des Doppelpuncts, die 4 Geraden
\[
y_1-y_2=0,\qquad y_1-k^2y_2=0,\qquad
y_1-\lambda^2y_2=0,\qquad y_1-\mu^2y_2=0
\]
sind die vier Tangenten, die man von dem Doppelpunct an die Curve ziehen kann, und ihre Berührungspuncte liegen auf der Geraden \(y_3=0\).

Man kann statt des Curvensystems \((n-1)^{\text{ter}}\) Ordnung auch ein System \((n-2)^{\text{ter}}\) Ordnung benutzen, welches durch alle Doppel- und Rückkehrpuncte, und durch \(n\) feste Puncte von \(f=0\) geht.

Dieses Curvensystem entsteht aus dem Vorigen, wenn \(n\) der festen Puncte in gerader Linie liegen, so dass das System \((n-1)^{\text{ter}}\) Ordnung in diese feste Gerade selbst und in ein System \((n-2)^{\text{ter}}\) Ordnung übergeht. Endlich kann man ein System \((n-3)^{\text{ter}}\) Ordnung benutzen, welches durch alle Doppel- und Rückkehrpuncte, einen ausgenommen, hindurchgeht. Dieses entsteht, wenn in dem System \((n-1)^{\text{ter}}\) Ordnung \(2n\) der festen Puncte auf einem Kegelschnitt liegen, wodurch das System \((n-1)^{\text{ter}}\) Ordnung in diesen und in ein System \((n-3)^{\text{ter}}\) Ordnung übergeht, welches die sämmtlichen Doppel- und Rückkehrpuncte zu Grundpuncten hat.

\begin{center}
{\large Vierter Abschnitt.}\\[0.35em]
{\Large Periodicität.}
\end{center}

\subsection*{\S\,22. Verzweigungspuncte.}

Stellen wir einen Werth der unabhängigen Variabeln \(x\) durch einen Punct der \(XY\)-Ebene dar, dessen Coordinaten der reelle und imaginäre Theil der Grösse \(x=X+Y\sqrt{-1}\) sind, so entsprechen jedem Puncte der Ebene \(n\) Werthe von \(s_x\), welche aus der Gleichung \(F(s_x,x)=0\) sich ergeben. Diese Werthe sind alle verschieden, sobald nicht \(x\) einer der Werthe ist, für welche zugleich
\[
F=0,\qquad \frac{\partial F}{\partial s_x}=0.
\]
Die solchen Werthen entsprechenden Puncte heissen Verzweigungspuncte\footnote{Riemann in Crelles Journal Bd. 54.}; in denselben fallen im Allgemeinen zwei Werthe von \(s_x\) zusammen; wir werden uns auf die Betrachtung dieses Falles beschränken, indem alle andern als Grenzfälle daraus folgen\footnote{Eine genaue Untersuchung aller möglichen Fälle findet man bei Puiseux, Recherches sur les fonctions algébriques, in Liouvilles Journal, Bd. 15.}.

Wenn man an einer beliebigen Stelle der \(XY\)-Ebene einen der Werthe von \(s_x\) wählt, und diesen sich auf einer geschlossenen Curve continuirlich ändern lässt, so kommt man, wie Puiseux gezeigt hat, immer zu dem ursprünglichen Werthe von \(s_x\) zurück, sobald die geschlossene Curve keinen Verzweigungspunct einschliesst; sind aber Verzweigungspuncte eingeschlossen, so kommt man unter Umständen zu einem andern Werthe von \(s_x\). Und dieser wieder bleibt ungeändert, wenn man die geschlossene Curve so variirt, dass bei der Gestaltsänderung der Curve nie ein Verzweigungspunct ein- oder austritt.

Ist nur ein einziger Verzweigungspunct eingeschlossen, so kann man zu einer andern Wurzel \(s_x\) gelangen. Wir können die geschlossene Curve dann zusammenziehen in einen Weg \(A\), welcher nach dem Verzweigungspunct hinführt, in einen unendlich kleinen Kreis um denselben, und in den Weg \(A\) in umgekehrter Richtung. Beginnen wir nun mit dem Werthe \(s'_x\), und mag dieser bei Durchlaufung des Weges \(A\) nach dem Verzweigungspunct hin sich in \(s_x^\circ\) verwandelt haben, welcher Werth dem Verzweigungspuncte als zugehörige Doppelwurzel entsprechen soll. Brechen wir den Weg \(A\) in der Entfernung \(\varepsilon\) von dem Verzweigungspuncte ab, um den kleinen Kreis zu durchlaufen, so wird auf demselben \(s_x\) von \(s_x^\circ\), \(x\) von dem Verzweigungswerthe \(x^\circ\) sehr wenig verschieden sein, und wenn man durch \(\varphi\) den Winkel bezeichnet, welchen ein Radius des Kreises gegen die \(X\)-Axe macht, so hat man
\[
x-x^\circ=\varepsilon e^{i\varphi}.
\]
Ferner geht nach dem Frühern \(F(s_x,x)=0\) über in:
\[
\frac{\partial F^\circ}{\partial x^\circ}(x-x^\circ)
+\frac{1}{2}\frac{\partial^2F^\circ}{\partial (s_x^\circ)^2}(s_x-s_x^\circ)^2=0,
\]
so dass
\[
s_x-s_x^\circ=\operatorname{Const.}\sqrt{x-x^\circ}
=\operatorname{Const.}\sqrt\varepsilon\, e^{\frac{i\varphi}{2}}.
\]
Beim Umlauf um den Kreis wächst nun \(\varphi\) um \(2i\pi\), und man erhält, wenn man wieder auf den Weg \(A\) zurückgekommen ist:
\[
s_x-s_x^\circ=-\operatorname{Const.}\sqrt\varepsilon\, e^{\frac{i\varphi}{2}}.
\]
Der Werth von \(s_x\) hat sich also beim Umlauf um den Kreis geändert; und man kann also auch, indem man den Weg \(A\) zurückverfolgt, nicht zu dem ursprünglichen Werthe \(s'_x\) gelangen, sondern muss zu einem andern Werth \(s''_x\) kommen. Und zwar sieht man auch, dass man umgekehrt von \(s''_x\) ausgehend zu \(s'_x\) zurückgelangt.

Fixiren wir also in der Ebene \(XY\) irgend einen Punct \(O\), und ziehen von hier gerade oder krumme, einander nicht schneidende Linien \(A\) nach den Verzweigungspuncten, so ist jeder Verzweigungspunct zweien der Werthe von \(s_x\), die dem Punct \(O\) angehören, und welche durch \(s^{(1)},s^{(2)},\ldots,s^{(n)}\) bezeichnet werden sollen, zugeordnet, denjenigen nämlich, welche durch Beschreibung eines Umgangs um den Verzweigungspunct in einander übergehen, während, sobald man einen Umgang um einen Verzweigungspunct mit einem nicht zugeordneten Werthe von \(s_x\) beginnt, der Verzweigungspunct sich vollkommen neutral verhält, und man also zu demselben Werthe von \(s_x\) zurückgelangt. Mit einer Abänderung der Curven \(A\) ändert sich die Art der Zugehörigkeit, sobald irgend eine der Curven während der Abänderung einen Verzweigungspunct überschreitet.

Dieses Verhalten der Verzweigungspuncte rührt wesentlich von dem Umstande her, dass \(s_x-s_x^\circ\) in der Nähe eines solchen Puncts mit \(\sqrt{x-x^\circ}\) proportional wird. Nimmt man statt eines Verzweigungspuncts, welcher einem Berührungspuncte entspricht (\S\,3), einen Punct, welcher einem Doppelpuncte der Curve \(F(s_x,x)=0\) entspricht, so wird in der Nähe eines solchen \(s_x-s_x^\circ\) mit \(x-x^\circ\) selbst proportional, und die soeben dargelegte charakteristische Eigenschaft des Verzweigungspuncts geht verloren. Daher werden wir solche Puncte nicht zu den Verzweigungspuncten rechnen; wohl aber die den Rückkehrpuncten entsprechenden, in denen \(s_x-s_x^\circ\) mit \(\sqrt{x-x^\circ}^{\,3}\) proportional wird. Da die Anzahl der Berührungspuncte \(n(n-1)-2d-3r\) ist, so wird also die Zahl der Verzweigungspuncte
\[
\nu=n(n-1)-2d-2r=2p+2(n-1).
\]

Der Weg nach einem Verzweigungspuncte hin, in kleinem Kreise um diesen herum und wieder zurück, soll der Kürze wegen eine Schleife genannt werden. Jede Schleife verbindet zwei bestimmte, dem Puncte \(O\) angehörige Wurzeln mit einander.

\subsection*{\S\,23. Cyclen und Umgänge.}

Die Gleichung \(F(s_x,x)=0\) ist als Gleichung für \(s_x\) nothwendig irreductibel, wenn nicht die ganze Curve in Theile zerfallen soll. Daher kann man den Satz aussprechen:

\emph{Die Lage des ganz willkürlichen Punctes \(O\) darf man immer so gewählt voraussetzen, dass man durch eine passende Reihe von Schleifen, von jedem beliebig vorausbestimmten, in \(O\) stattfindenden Werthe von \(s_x\) zu jedem beliebigen andern vorausbestimmten gelangen kann.}

Wäre dies nämlich nicht der Fall, so müsste es einen Cyclus von Wurzeln \(s_x^{(\alpha)},s_x^{(\beta)},\ldots\) geben, welche nur immer in einander übergingen, welche Schleifen man auch durchläuft. Eine jede symmetrische Function von \(s_x^{(\alpha)},s_x^{(\beta)},\ldots\) würde also, indem man sie über einen beliebigen Weg hin, oder über die Schleifen, in welche er zusammengezogen werden kann, continuirlich sich ändern liesse, in jedem beliebig gewählten Puncte \(O\) immer denselben Werth haben, würde also überall in der Ebene \(XY\) nur einen einzigen Werth besitzen, und demnach eine rationale Function von \(x\) sein\footnote{Briot et Bouquet, Théorie des fonct. doubl. périod.}. Man könnte also eine Gleichung aufstellen, deren Coefficienten rationale Functionen von \(x\) wären, und deren Wurzeln nur gewisse der \(n\) Wurzeln \(s_x\) wären; diese Gleichung müsste also ein Factor von \(F(s_x,x)\) sein, was ausgeschlossen ist.

Heben wir also irgend einen Verzweigungspunct heraus, welcher etwa den Wurzeln \(s^{(1)},s^{(2)}\) zugeordnet ist, so wird es immer eine Wurzel \(s^{(3)}\) geben, welche entweder mit \(s^{(1)}\) oder mit \(s^{(2)}\) durch eine Schleife zusammenhängt, sodann eine Wurzel \(s^{(4)}\), welche mit \(s^{(1)},s^{(2)}\) oder \(s^{(3)}\) zusammenhängt u. s. w. Wir wählen \(n-1\) Verzweigungspuncte als Fundamentalpuncte erster Art, welche die Eigenschaft haben, dass einer etwa von \(s^{(1)}\) zu \(s^{(2)}\), der zweite von \(s^{(1)}\) oder \(s^{(2)}\) zu \(s^{(3)}\), der dritte von \(s^{(1)},s^{(2)}\) oder \(s^{(3)}\) zu \(s^{(4)}\) führt u. s. w. Jeder andere Verzweigungspunct bildet mit einigen von diesen einen Cyclus, insofern die Schleife jedes andern Verzweigungspuncts, welche etwa von \(s^{(i)}\) nach \(s^{(k)}\) führt, mit passend gewählten Schleifen der Fundamentalpuncte, und zwar nur auf eine Art, so verbunden werden kann, dass man den ganzen so gebildeten Weg durchlaufend zu dem Ausgangswerthe von \(s\) zurückkommt. Um dies deutlich zu machen, combiniren wir die Fundamentalschleifen zunächst zu solchen Wegen \(W_2,W_3,\ldots,W_n\), dass \(W_2\) von \(s^{(1)}\) nach \(s^{(2)}\), \(W_3\) von \(s^{(1)}\) nach \(s^{(3)}\) führt u. s. w.; die von \(s^{(i)}\) nach \(s^{(k)}\) führende Schleife bildet dann mit \(W_i\) und \(W_k\) einen Cyclus, wobei der etwa dem \(s^{(1)}\) entsprechende Weg \(W_1\) auszulassen ist.

Aus den Fundamentalpuncten erster Art selbst lässt sich offenbar nie ein Cyclus zusammensetzen. Jeder beliebige Cyclus aber lässt sich in solche Cyclen auflösen, welche nur einen andern Verzweigungspunct enthalten. Denn indem man jede Schleife des gegebenen Cyclus, Fundamentalschleifen ausgenommen, durch Zufügung von Fundamentalschleifen zu solchen einfachen Cyclen ergänzt, erhält man den gegebenen Cyclus gleich der Summe der neuen, vermehrt oder vermindert um gewisse Fundamentalschleifen. Das Aggregat dieser müsste also auch einen Cyclus bilden, was nie geschehen kann, und daher müssen diese Fundamentalpuncte geradezu wegfallen, indem die Schleife jedes derselben in positivem und negativem Sinn auftritt; so dass der gegebene Cyclus gleich einer Summe einfacher Cyclen wird.

Unter den Cyclen giebt es einige vorzugsweise ausgezeichnete. Es sind die, welche aus vollständigen Umgängen entstehn. Unter einem vollständigen Umgange verstehen wir einen Weg, welcher alle Schleifen der Reihe nach umfasst; und zwar denken wir uns die Schleifen dabei so geordnet, wie die Tangenten der vom Puncte \(O\) gezogenen Wege in \(O\) selbst auf einander folgen, wenn man sich in einem kleinen Kreise um \(O\) bewegt. Der ganze Weg eines Umgangs lässt sich dann durchs Unendliche ohne Ueberschreitung eines Verzweigungspuncts auf den Punct \(O\) zusammenziehen, woraus folgt, dass jeder Umgang zu dem Ausgangswerthe von \(s_x\) zurückführt. Es giebt also \(n\) getrennte Umgänge, deren jeder mit einem der in \(O\) stattfindenden Werthe von \(s_x\) beginnt und mit demselben endigt. In einem solchen Umgange sind einige Verzweigungspuncte wesentlich, d.\ h. sie führen von einer Wurzel zu einer andern, andere unwesentlich. In zwei verschiedenen Umgängen kann dieselbe Schleife nicht mit denselben Wurzelwerthen durchlaufen werden, weil sonst alle vorhergehenden und folgenden Schleifen, und also auch die ganzen Umgänge, übereinstimmen müssten. Daher kann jeder Verzweigungspunct nur in zwei Umgängen wesentlich sein, indem er in dem einen von einem seiner zugeordneten Wurzeln \(s^{(i)}\) nach der andern \(s^{(k)}\) führt, in dem andern Umgange aber umgekehrt von \(s^{(k)}\) nach \(s^{(i)}\). Die Cyclen der \(n\) Umgänge bestehen also aus den Schleifen der wesentlichen Verzweigungspuncte derselben, und man kann jeden Umgang durch eine Summe von einfachen Cyclen ersetzen. Da sich der ganze Umgang durch


% --- Continuation: source PDF pp. 107--121. ---
% Primary source witness: Google scan PDF theoriederabels00clebgoog.pdf, source PDF pp. 107--121.
% This tranche begins mid-sentence from §23 and continues through the opening of §28.

stetige Aenderung auf einen Punct reduciren lässt, und damit sein Effect identisch Null wird, so kann man auch sagen, dass einer dieser einfachen Cyclen durch die übrigen desselben Umgangs ausgedrückt werden kann. So könnte es scheinen, als könne man \(n\) Cyclen, die in verschiedenen Umgängen vorkommen, durch die übrigen ausdrücken; in der That ist dies nur mit \(n-1\) derselben der Fall. Nämlich der Weg, welcher aus allen Umgängen besteht, enthält jede Schleife zweimal wesentlich, und zwar einmal von \(s^{(i)}\) zu \(s^{(k)}\), das andre Mal von \(s^{(k)}\) zu \(s^{(i)}\) führend. Die verschiedenen Schleifen dieses Gesammtumgangs zerstören einander daher geradezu, ohne dass man die Verschiebung des Weges zu Hülfe zu nehmen braucht. Daher ist der letzte Umgang nichts anderes, als eine Combination der frühern, und kann daher auch zu keiner neuen Relation zwischen den Cyclen Veranlassung geben. Da wir nun \(2p+n-1\) einfache Cyclen hatten, und zwischen denselben \(n-1\) Relationen bestehen, so haben wir den Satz:

\begin{center}
\emph{Es giebt \(2p\) von einander unabhängige einfache Cyclen.}
\end{center}

\subsection*{\S\,24. Reciprocität zwischen zwei Classen von Umgängen und zwischen zwei Classen von Fundamentalpuncten.}

Zwischen den \(n\) Umgängen einerseits und den \(n\) im Puncte \(O\) stattfindenden Werthen von \(s_x\) andrerseits findet nun eine merkwürdige Reciprocität statt, welche man, wie folgt, darstellen kann.

Verbinden wir die Verzweigungspuncte in der oben angegebenen Reihenfolge durch eine gebrochene Linie, welche bei dem ersten beginnt und bei dem letzten aufhört. Die nach den verschiedenen Verzweigungspuncten gezogenen Schleifen \(a_1,a_2\ldots a_{\nu}\) liegen dann sämmtlich auf derselben Seite dieser gebrochenen Linie. Beschreiben wir nun ferner solche Wege \(A_1,A_2\ldots A_{\nu}\), welche von \(O\) ausgehend immer nur den einen, durch den Index angezeigten Verzweigungspunct ausschliessen, alle übrigen aber umfassen. Ein solcher Weg ist immer zwei Umgängen zugeordnet, den beiden nämlich, in denen der ausgeschlossene Verzweigungspunct als wesentlicher auftritt. Verbindet nämlich die dem Verzweigungspunct angehörende Schleife die Wurzeln \(s^{(i)}\) und \(s^{(k)}\), und ist der Punct ein wesentlicher für die mit \(s^{(\alpha)}\) und \(s^{(\beta)}\) beginnenden Umgänge; so führt zunächst der betreffende Weg \(A\), wenn er nicht mit einer der Wurzeln \(s^{(\alpha)},s^{(\beta)}\) begonnen wird, immer zu der Ausgangswurzel zurück, indem der ausgeschlossene Verzweigungspunct hier gleichgültig wird, und der Weg \(A\) also durch einen vollständigen Umgang ersetzt werden kann. Beginnt man aber den Weg \(A\) mit der Wurzel \(s^{(\alpha)}\), so stimmt bis zu der ausgeschlossenen Schleife alles genau mit dem \(\alpha^{\text{ten}}\) Umgange überein, bis zu der Stelle also, bei welcher durch jene Schleife \(s^{(i)}\) in \(s^{(k)}\) übergeführt werden müsste. Statt dessen geht man mit \(s^{(i)}\) zur nächsten Schleife weiter, verhält sich also so, wie im \(\beta^{\text{ten}}\) Umgange, bei welchem umgekehrt die ausgeschlossene Schleife mit \(s^{(i)}\) geendigt hätte. Der Weg \(A\) besteht also, wenn mit \(s^{(\alpha)}\) begonnen wird, aus dem Anfange des \(\alpha^{\text{ten}}\) und dem Ende des \(\beta^{\text{ten}}\) Umgangs; und ebenso, wenn mit \(s^{(\beta)}\) begonnen wird, mit dem Anfange des \(\beta^{\text{ten}}\) und dem Ende des \(\alpha^{\text{ten}}\) Umgangs.

Dieser Weg \(A\) ist also den Umgängen \(\alpha\) und \(\beta\) zugeordnet, wie die dem Verzweigungspunct entsprechende Schleife den Wurzeln \(i\) und \(k\). Nun kann man aber den Weg \(A\) durch stetige Veränderung in eine Schleife überführen, welche sich ähnlich wie die Schleife \(a\) verhält, aber auf der andern Seite der gebrochenen Linie liegt. \emph{Das System der neuen Schleifen \(A\) ist also dem System der ursprünglichen Schleifen \(a\) ganz analog; und zwar sind beide einander in Bezug auf die gebrochene Linie conjugirt, in der Weise, dass die Schleife des einen Systems die Wurzeln \(s^{(\alpha)},s^{(\beta)}\) verbindet, in deren entsprechenden Umgängen die entsprechende Schleife des andern Systems wesentlich ist, und umgekehrt.}

Dem oben ausgesprochenen Satze, dass man durch die Schleifen erster Art von jeder Wurzel \(s^{(i)}\) zu jeder andern \(s^{(k)}\) gelangen könne, entspricht also der Satz, dass man durch die Schleifen zweiter Art aus jedem Umgange in jeden andern gelangen könne, also auch durch Schleifen erster Art, da jede Schleife einer Art sich durch Schleifen der andern Art darstellen lässt. Den Fundamentalpuncten erster Art entsprechen also auch \(n-1\) Fundamentalpuncte zweiter Art, welche zu dem zweiten Schleifensystem in derselben Beziehung stehen, wie jene zum ersten und welche also zu diesem die Beziehung haben, dass die ihnen angehörigen Verzweigungspuncte alle Umgänge des ersten Systems unter einander verbinden.

Man kann aber ferner den merkwürdigen Satz beweisen, \emph{dass die beiden Reihen von Fundamentalpuncten immer so gewählt werden können, dass kein Fundamentalpunct einer Art zugleich der andern Art angehört.} Dies zeigt man auf folgende Weise.

In einem Umgange, in welchem der erste Fundamentalpunct erster Art vorkommt, ist, da die Puncte des Umgangs einen Cyclus bilden, also nicht lauter Fundamentalpuncte erster Art sein können, wenigstens noch ein anderer wesentlicher Verzweigungspunct vorhanden. Wir wählen ihn zum ersten Fundamentalpunct der zweiten Art. Dieser Punct verbindet den ersten Umgang mit einem zweiten. Nun beweist man zweitens, dass es im ersten und zweiten Umgange zusammen wenigstens einen Punct geben muss, welcher kein Fundamentalpunct erster Art ist, und welcher nur in einem der beiden Umgänge vorkommt. Dieser Punct als zweiter Fundamentalpunct zweiter Art verbindet die ersten beiden Umgänge mit einem dritten. Sodann zeigt man drittens, dass in allen drei Umgängen zusammen wenigstens ein Punct vorkommt, welcher nur einmal auftritt und kein Fundamentalpunct erster Art ist. Dieser Punct als dritter Fundamentalpunct zweiter Art verbindet die drei ersten Umgänge mit einem vierten. Indem man so fortfährt, gelangt man zu einem vollständigen System von Fundamentalpuncten zweiter Art, welches mit dem ersten System keinen Punct gemein hat. Und zwar ist dabei nur folgender Hülfssatz zu beweisen:

\begin{center}
\emph{Wenn \(i\) Umgänge durch Fundamentalpuncte zweiter Art bereits verbunden sind, so giebt es in ihnen ausser Fundamentalpuncten erster Art wenigstens noch einen Punct, der nur in einem dieser Umgänge vorkommt.}
\end{center}

Der Beweis dieses Satzes ist sehr einfach. Nehmen wir eine symmetrische Function der Wurzeln \(s^{(1)},s^{(2)}\ldots s^{(i)}\), mit denen diese \(i\) Umgänge beginnen, und lassen wir die \(s\) sich stetig ändern, während wir um alle Verzweigungspuncte gehen. Wir gelangen nach Beschreibung des Umgangs zu den ursprünglichen Werthen der \(s\), also auch zu dem Functionswerth zurück. Für diesen Umgang aber und diese symmetrische Function sind alle in zweien der \(i\) Umgänge als wesentlich vorkommenden Verzweigungspuncte unwesentlich, da sie nur gleichzeitig \(s^{(k)}\) in \(s^{(h)}\), \(s^{(h)}\) in \(s^{(k)}\) überführen, also die symmetrische Function nicht ändern. Der für die symmetrische Function gebildete Umgang muss aber wesentliche Verzweigungspuncte besitzen, weil die Function sonst rational wäre, was dem Frühern widerspricht. Es muss also in den \(i\) Umgängen einzeln auftretende Verzweigungspuncte geben. Aber dieselben bilden nothwendig eine Reihe von Cyclen\footnote{Man kann nämlich folgende Sätze aussprechen: \emph{Zwei cyclische Wege, in denen ein Verzweigungspunct gemeinsam ist, dessen Schleife in beiden entgegengesetzt durchlaufen wird, lassen sich zu einem cyclischen Wege vereinigen, in dem jener doppelte Verzweigungspunct fortgefallen ist.} \emph{Ein cyclischer Weg, in dem ein Verzweigungspunct doppelt vorkommt, so dass seine Schleife beide Male entgegengesetzt durchlaufen wird, lässt sich in zwei cyclische Wege theilen, in denen jene Schleife nicht mehr vorkommt.} Mit Hülfe dieser beiden Sätze kann man offenbar immer eine beliebige Zahl von Umgängen auf Cyclen reduciren, aus denen die den Umgängen gemeinsamen Verzweigungspuncte ausgefallen sind. Den ersten Satz beweist man folgendermassen. Der erste Weg führe von \(s^{(\alpha)}\) irgendwie zu \(s^{(i)}\), womit die gemeinsame Schleife beginnt, und nach Durchlaufung derselben von \(s^{(k)}\) nach \(s^{(\alpha)}\) zurück. Der andre führe von \(s^{(\beta)}\) nach \(s^{(k)}\), und nach Durchlaufung der gemeinsamen Schleife von \(s^{(i)}\) nach \(s^{(\beta)}\) zurück. Der vereinigte Weg, bei welchem jene Schleife ausfällt, geht dann wie der erste von \(s^{(k)}\) über \(s^{(\alpha)}\) nach \(s^{(i)}\), und von dort an den zweiten anknüpfend von \(s^{(i)}\) über \(s^{(\beta)}\) nach \(s^{(k)}\) zurück. Ebenso führe im zweiten Satze der Weg von \(s^{(\alpha)}\) irgendwie nach \(s^{(i)}\), nach Durchlaufung der Schleife von \(s^{(k)}\) über \(s^{(\beta)}\) nach \(s^{(k)}\), und endlich nach nochmaliger Durchlaufung der Schleife von \(s^{(i)}\) nach \(s^{(\beta)}\) zurück. Der eine cyclische Weg geht dann von \(s^{(k)}\) über \(s^{(\beta)}\) nach \(s^{(k)}\) zurück, der andre von \(s^{(i)}\) über \(s^{(\alpha)}\) nach \(s^{(i)}\) zurück, wodurch die geforderte Zerlegung geleistet ist.}. Da nun Fundamentalpuncte erster Art keinen Cyclus bilden können, so muss es wenigstens noch einen andern Punct unter diesem Reste geben, was zu beweisen war.

\subsection*{\S\,25. Integration auf verschiedenen Wegen. Periodicitätsmoduln der Integrale erster Gattung.}

Wir wollen jetzt die oben bezüglich der continuirlichen Aenderung der Function \(s_x\) betrachteten Wege als Integrationswege eines Integrals erster Gattung untersuchen. Ein solches ist vollständig bestimmt, wenn die Anfangs- und Endwerthe von \(s_x\) und \(x\) nebst dem Wege sich angegeben finden, auf welchem man von dem einen Werthsysteme zu dem andern gelangt. Als Anfangswerth von \(x\) können wir immer denjenigen betrachten, welcher dem Puncte \(O\) zukommt, und als Anfangswerth von \(s_x\) einen der zugehörigen Werthe dieser Grösse. Zwei Integrationswege, die durch das Endliche oder Unendliche ohne Ueberschreitung eines Verzweigungspuncts in einander überführbar sind, liefern immer dasselbe Resultat als Werth des Integrals. Wenn sie aber Verzweigungspuncte einschliessen, und zwar weder alle noch keinen, so dass eine directe Ueberführung eines Integrationsweges in den andern nicht möglich ist, so hat man im Allgemeinen zwei verschiedene Werthe des Integrals vor sich, obgleich die Werthpaare \(x,s_x\) für die obern und untern Grenzen dieselben sind. Man kann dann statt des zweiten Integrationsweges den ersten setzen, vermehrt um eine geschlossene Curve, welche den zweiten Integrationsweg in directer, den ersten in entgegengesetzter Richtung enthält; man kann also den zweiten Integrationsweg aus dem ersten durch Hinzufügung eines cyclischen Weges entstanden denken. Nun können, wie oben gezeigt, alle cyclischen Wege aus \(2p\) derselben zusammengesetzt werden, welche den \(2p\) unabhängigen einfachen Cyclen entsprechen. Der allgemeinste Integrationsweg, welcher neue Werthe des Integrals liefern kann, ist also aus irgend einem entstanden zu denken, indem man die einfachen cyclischen Wege beliebig oft positiv oder negativ hinzugefügt hat. Nennen wir also den Werth des Integrals, über einen cyclischen Weg ausgedehnt, einen \emph{Periodicitätsmodul} des Integrals, so zeigt sich erstlich, \emph{dass jeder Periodicitätsmodul durch \(2p\) einfache Periodicitätsmoduln linear mit ganzzahligen Coefficienten ausdrückbar ist, oder dass ein Integral erster Gattung \(2p\) unabhängige Periodicitätsmoduln hat.} Nennen wir ferner \(I\) das Integral auf beliebigem Wege genommen, \(I',I^{(2)}\ldots I^{(2p)}\) seine Periodicitätsmoduln, so ist
\[
I+m'I'+m^{(2)}I^{(2)}\cdots+m^{(2p)}I^{(2p)}
\]
der allgemeinste Werth, welchen das Integral durch Abänderung des Integrationsweges erhalten kann, wo \(m',m^{(2)},\ldots,m^{(2p)}\) positive oder negative ganze Zahlen sind.

Die \(p\) von einander unabhängigen Integrale erster Gattung wollen wir, auf irgend einem gemeinsamen Wege zwischen denselben Grenzen genommen, durch \(I_1,I_2,\ldots,I_p\) bezeichnen. Aendern wir dann den Integrationsweg, so ändern sich gleichzeitig die Werthe aller Integrale; indem die Zahlen \(m\), welche nur von der Natur des neuen Weges abhängen, für alle Integrale dieselben bleiben, während die Periodicitätsmoduln für jedes Integral individuelle Werthe haben. Das allgemeine Werthsystem der Integrale \(I\) auf irgend einem Wege ist also durch das Schema dargestellt:
\[
\begin{gathered}
I_1+m'I_1'+m^{(2)}I_1^{(2)}\cdots+m^{(2p)}I_1^{(2p)},\\
I_2+m'I_2'+m^{(2)}I_2^{(2)}\cdots+m^{(2p)}I_2^{(2p)},\\
\vdots\\
I_p+m'I_p'+m^{(2)}I_p^{(2)}\cdots+m^{(2p)}I_p^{(2p)}.
\end{gathered}
\]
Die Periodicitätsmoduln bilden ein System von \(2p^2\) Constanten, bei welchem die \(2p\) Grössen derselben Horizontale sich auf dasselbe Integral und verschiedene cyclische Wege, die \(p\) Grössen derselben Verticale sich auf denselben cyclischen Weg und auf verschiedene Integrale beziehen.

Zwischen den Grössen \(I_h^{(k)}\) bestehen algebraische Relationen; und zwar insbesondere zwischen den Grössen \(I_h^{(k)}\) zweier Horizontalreihen bestehen Gleichungen, welche linear für die \(I_h^{(k)}\) jeder der beiden Reihen sind. Diese sollen jetzt entwickelt werden.

\subsection*{\S\,26. Relationen zwischen den Periodicitätsmoduln der Integrale erster Gattung.\footnotemark}\footnotetext{Die hier aufzustellenden Relationen umfassen keineswegs sämmtliche existirende, sondern enthalten nur diejenigen, welche für den Verlauf der Theorie nothwendig sind.}

Bezeichnen wir durch \(I\) und \(H\) irgend zwei der \(p\) Integrale erster Gattung, und führen wir das Integral \(\int I\,dH\) über alle \(n\) vollständigen Umgänge. Die zu integrirende Function \(I\dfrac{\partial H}{\partial x}\) wird nirgends unstetig und hat keine Verzweigungspuncte als die von \(F(s_x,x)=0\). Daher ist das Integral, über einen Umgang genommen, gleich Null, also auch das Integral, welches wir betrachten. Es kommt nun darauf an, die linke Seite der Gleichung
\[
\int I\,dH=0
\]
anders darzustellen. Dieses Integral zerfällt zunächst in \(n\) Theile, den \(n\) Umgängen entsprechend; jeder Theil wieder in \(\nu\) Theile, die Integrale über die einzelnen Schleifen. Zunächst sieht man nun leicht, dass jede Schleife als ihren Antheil des Integrals den Werth Null giebt, sobald die Schleife in dem betreffenden Umgange unwesentlich war. Denn betrachten wir den Hin- und Rückweg, welchen wir bei einer solchen Schleife zwischen \(O\) und dem Verzweigungspuncte machen müssen, so sehen wir, dass an derselben Stelle \(I\) beidemal den gleichen, \(dH\) aber den entgegengesetzten Werth besitzt. Die Elemente des Hin- und Rückwegs heben einander also genau auf. Es bleiben daher nur die wesentlichen Schleifen der Umgänge zu betrachten, so dass jede Schleife nur zweimal zu berücksichtigen ist. Und zwar werden beide Male dieselben Werthe von \(s_x\) und \(x\) nur in entgegengesetzter Reihenfolge durchlaufen, da die Schleife einmal von \(s^{(i)}\) nach \(s^{(k)}\), das andere Mal von \(s^{(k)}\) nach \(s^{(i)}\) führt. Vergleichen wir also wieder zwei Elemente solcher derselben Schleife entsprechenden Integrale, für welche \(s_x\) und \(x\) dieselben Werthe besitzen, so hat \(dH\) für beide gleiche und entgegengesetzte Werthe; bezeichnen wir also durch \(I\) und \(I'\) die beiden Werthe von \(I\), so ist das Element unseres Integrals
\[
(I-I')\,dH.
\]
Nun ist \(I\) das Integral, genommen vom Anfang des ersten Umgangs, in welchem der betrachtete Verzweigungspunct wesentlich ist, und zwar zunächst bis zu dem Werthe \(s^{(i)}\), mit welchem die Schleife des Verzweigungspuncts beginnt, sodann aber von dort bis zu dem variabeln Puncte der Schleife. Ebenso ist \(I'\) das Integral, genommen von dem Anfang des zweiten Umgangs, in welchem der Verzweigungspunct wesentlich ist, und zwar zunächst bis \(s^{(k)}\), womit die Schleife in diesem Umgange beginnt, sodann aber weiter bis zu dem variabeln Punct. Statt \(-I'\) kann man also den Rest dieses Umgangs setzen, d. h. zunächst das Integral von dem variabeln Puncte bis zum Werthe \(s^{(i)}\), mit dem die Schleife in diesem Umgange endigt, sodann aber von dort weiter bis zum Ende des Umgangs. Es wird also \(I-I'\), indem sich die Integrale von dem variabeln Puncte bis zu \(s^{(i)}\) aufheben, ein Integral über einen Weg, welcher alle Verzweigungspuncte bis auf den betrachteten einschliesst; ein Weg, den wir im vorigen \S{} in eine Schleife zweiter Art übergeführt haben.

Bezeichnen wir also die Werthe des Integrals \(I\), geführt über die Schleifen, so dass man dabei von einer Wurzel \(s\) mit kleinerem Index zu einer Wurzel \(s\) mit grösserm Index gelangt, durch \(a_1,a_2\ldots a_{\nu}\), die entsprechenden Werthe über die Schleifen zweiter Art geführt durch \(A_1,A_2\ldots A_{\nu}\), so dass, wenn man \(A_i\), wie oben geschah, durch die \(a\) ausdrückt, der Weg von \(A\) mit demjenigen Umgange beginnt, in welchem die Schleife \(+a_i\) vorkommt. Dann wird das oben untersuchte Element der \(i^{\text{ten}}\) Schleife
\[
(I-I')\,dH=A_i\,dH;
\]
also der von dieser Schleife herrührende Antheil des Integrals wird \(A_i\int dH\), und wenn man die für \(H\) genommenen Integrale, welche den \(a_i,A_i\) entsprechen, durch \(b_i,B_i\) bezeichnet, wird dieser Antheil \(A_i b_i\). Die obige Gleichung geht also in folgende über:
\[
(1)\quad A_1b_1+A_2b_2+\cdots+A_{\nu}b_{\nu}=0.
\]
Bemerken wir dabei, dass wir diese Gleichung auch in der Form schreiben können:
\[
(2)\quad B_1a_1+B_2a_2+\cdots+B_{\nu}a_{\nu}=0.
\]
Denn addiren wir beide Gleichungen und gehen auf ihre Ableitung zurück, so haben wir
\[
\Sigma(A_i b_i+B_i a_i)=\int(I\,dH+H\,dI).
\]
Die rechte Seite dieses Ausdrucks unbestimmt integrirt giebt \(H.I\); und da \(H\) sowohl als \(I\) am Anfange und am Ende jedes Umgangs denselben Werth besitzen, so verschwindet das Integral rechts über jeden einzelnen Umgang, also auch über alle Umgänge zusammen genommen. Man sieht also, dass man die Ausdrücke (1), (2) in einander muss überführen können, mit alleiniger Anwendung der linearen zwischen den \(a\), den \(b\), den \(A\), und den \(B\) entstehenden Relationen, welche aussagen, dass das Integral \(I\) oder \(H\), über alle Schleifen desselben Umgangs ausgedehnt, den Werth Null giebt.

Die Gleichungen (1), (2) bestehen zwischen den über die einzelnen Schleifen ausgedehnten Integralen. Man kann sie so verwandeln, dass sie nur Integrale, welche über cyclische Wege ausgedehnt sind, also Periodicitätsmoduln enthalten. Dies gelingt durch drei auf einander folgende Transformationen.

Statt der Grössen \(b_1,b_2\ldots\) wollen wir zunächst die den betreffenden Verzweigungspuncten entsprechenden, über die einfachen cyclischen Wege ausgedehnten Integrale \((b_1),(b_2)\ldots\) einführen. Bilden wir als Combination der Integrale über Fundamentalschleifen erster Art uns zunächst \(n-1\) Grössen \(f_2,f_3\ldots f_n\), deren Wege beziehungsweise von \(s^{(1)}\) nach \(s^{(2)}\), von \(s^{(1)}\) nach \(s^{(3)}\), \(\ldots\) von \(s^{(1)}\) nach \(s^{(n-1)}\) führen (vgl. \S\,23), so ist, wenn \(b_i\) von \(s^{(\alpha)}\) zu \(s^{(\beta)}\) führt, das über den entsprechenden einfachen cyclischen Weg genommene Integral (Periodicitätsmodul):
\[
(b_i)=b_i+f_{\alpha}-f_{\beta};
\]
wobei immer \(f_1\) durch Null zu ersetzen ist. Setzen wir nun also in (1)
\[
b_i=(b_i)-f_{\alpha}+f_{\beta},
\]
so wird auf der linken Seite von (1) jede Grösse \(f_{\alpha}\) multiplicirt mit allen denjenigen \(A\), für deren Verzweigungspunct \(s^{(\alpha)}\) eine der zugeordneten Wurzeln ist, und zwar ist das betreffende \(A\) negativ, wenn die positive Schleife erster Art des Verzweigungspuncts mit \(s^{(\alpha)}\) beginnt, positiv im entgegengesetzten Falle. Es bilden also die diesen \(A\) zugehörigen Schleifen zweiter Art einen vollständigen Umgang zweiter Art, oder ihre Summe verschwindet. Somit erhält man die Transformation
\[
(3)\quad \Sigma A_i b_i=\Sigma A_i(b_i).
\]
Die Summe links enthält \(n-1\) Glieder weniger als die Summe rechts; denn alle diejenigen \((b_i)\) verschwinden identisch, welche den Fundamentalpuncten erster Art selbst zugehören.

Ganz ähnlich besitzen die zwischen den \(a\) oder \(b\) bestehenden linearen Gleichungen, welche aussagen, dass ein Integral, über einen Umgang geführt, verschwindet, die Eigenschaft, dass man darin die \(a\) oder \(b\) durch die \((a)\), \((b)\) ersetzen kann. Dies folgt schon aus dem in \S\,23 gegebenen Beweis, dass man in einem Umgange die Fundamentalpuncte erster Art auslassen, die andern Puncte aber durch die entsprechenden einfachen Cyclen erster Art ersetzen kann.

Denken wir uns nun, um die zweite Transformation auszuführen, die Grössen \(A\) als lineare Functionen der \(a\) ausgedrückt, wie die Wege der \(A\) ursprünglich aus Schleifen erster Art zusammengesetzt wurden. Setzen wir in diesen Ausdrücken statt der \(a\) die entsprechenden Perioden \((a)\), so mögen die \(A\) in Ausdrücke \((A)\) übergehen. Bezeichnen wir die Combinationen der \(a\), welche den oben durch \(f\) bezeichneten Combinationen der \(b\) entsprechen, durch \(g\), und beginnt der Umgang \(A_i\) mit \(s^{(\alpha)}\) und endigt mit \(s^{(\beta)}\), so ist
\[
(A_i)=A_i+g_{\alpha}-g_{\beta}.
\]
Setzen wir also in der rechten Seite der Gleichung (3)
\[
A_i=(A_i)-g_{\alpha}+g_{\beta},
\]
so wird darin \(g_{\alpha}\) mit allen denjenigen \((b_i)\) multiplicirt, deren Verzweigungspuncte dem \(\alpha^{\text{ten}}\) Umgange angehören; und zwar negativ, wenn die betreffende Schleife erster Art in demselben positiv auftritt, und positiv, wenn das Gegentheil stattfindet. Daher bilden alle mit \(g_{\alpha}\) multiplicirten \((b_i)\) einen vollständigen Umgang; ihre algebraische Summe verschwindet nach dem Vorigen, und man hat:
\[
(4)\quad \Sigma A_i b_i=\Sigma(A_i)(b_i).
\]
Wir können nun drittens den Ausdruck (4) dadurch transformiren, dass wir mit Hülfe der zwischen den \((b_i)\) stattfindenden Bedingungsgleichungen \(n-1\) dieser Grössen eliminiren. Nach der Theorie, welche wir oben über die Fundamentalpuncte zweiter Art aufgestellt haben, lassen sich die diesen entsprechenden Grössen \((b_i)\) aus den linearen Bedingungsgleichungen immer bestimmen. Denn in jedem folgenden Umgange war wenigstens ein neuer Fundamentalpunct zweiter Art enthalten, also auch eine neue Grösse \((b_i)\), welche den Coefficienten \(\pm1\) besitzt, und sich also als lineare Function der übrigen \((b_i)\) mit ganzzahligen Coefficienten darstellt. Denken wir uns diese Bestimmungen ausgeführt, so hat man endlich:
\[
(5)\quad \Sigma A_i b_i=\Sigma[A_i](b_i),
\]
wo die Summe sich nur noch über \(2p\) Verzweigungspuncte erstreckt, welche weder Fundamentalpuncte erster noch zweiter Art sind. Die \(2p\) Grössen \((b_i)\) sind jetzt keiner linearen Relation mehr unterworfen, da die Bedingungsgleichungen zur Elimination der \(n-1\) übrigen \((b_i)\) verwandt wurden. Die Grössen \([A_i]\) sind genauer zu untersuchen und zu characterisiren.

\subsection*{\S\,27. Form der gefundenen Relation und Eigenschaften ihrer ganzzahligen Coefficienten.}

Bilden wir uns aus den \(n-1\) Perioden \((A_i)\), welche den \(n-1\) Fundamentalpuncten zweiter Art entsprechen, solche Combinationen \(h_2,h_3\ldots h_n\), dass die entsprechend aus den \(A_i\) zusammengesetzten Wege aus dem ersten Umgange beziehungsweise in den zweiten, dritten etc. führen. Gehört dann der \(i^{\text{te}}\) Verzweigungspunct dem \(\alpha^{\text{ten}}\) Umgang mit positiv durchlaufener Schleife, dem \(\beta^{\text{ten}}\) mit negativ durchlaufener an, so ist
\[
{}[A_i]=(A_i)+h_{\alpha}-h_{\beta},
\]
wobei \(h_1\) immer durch Null zu ersetzen ist. In der That, nämlich, wenn man hieraus
\[
(A_i)=[A_i]-h_{\alpha}+h_{\beta}
\]
in die rechte Seite der Gleichung (4) einsetzt, geschieht Folgendes. Die den Fundamentalpuncten erster Art zugehörigen \((A_i)\) kommen in jenem Ausdrucke überhaupt nicht vor. Die \([A_i]\), welche den Fundamentalpuncten zweiter Art entsprechen, verschwinden identisch, weil die Combination \(h_{\alpha}-h_{\beta}\) sich dann nothwendig auf \(-(A_i)\) reducirt, da Fundamentalpuncte zweiter Art nie einen Cyclus zweiter Art erzeugen können. Endlich verschwindet auch alles, was mit den \(h_{\alpha},h_{\beta}\) multiplicirt ist. Denn es wird \(h_{\alpha}\) mit allen denjenigen \((b_i)\) multiplicirt, deren Verzweigungspuncte auf dem \(\alpha^{\text{ten}}\) Umgange liegen, und zwar negativ, wenn auf diesem Umgange die Schleife des Puncts positiv vorkommt, positiv, wenn das Entgegengesetzte eintritt. Der Coefficient von \(h_{\alpha}\) ist also genau einer der Ausdrücke, welche wegen der linearen Bedingungsgleichungen verschwinden. Man kommt also auf die Form (5) zurück, und die neue Definition der \([A_i]\) muss also mit der dort gegebenen übereinstimmen, da zwischen den übrig gebliebenen \((b_i)\) keine Relationen mehr auftreten.

Die Grössen \([A_i]\) sind Periodicitätsmoduln von solcher Art, dass die bei ihrer Bildung benutzten Verzweigungspuncte nicht blos Cyclen erster Art, sondern auch Cyclen zweiter Art bilden, d. h. solche, welche in Rücksicht auf die Schleifen zweiter Art gebildet sind, oder welche in stetiger Folge immer einen Umgang erster Art mit dem andern verbinden, zuletzt aber zu dem anfangs benutzten Umgange zurückkehren, so wie die Cyclen erster Art stetig von einer Wurzel zur andern und endlich zu der anfänglich benutzten zurückführen.

Die Grössen \(A_i\) erschienen bei ihrer Bildung als lineare Functionen der \(a_i\) mit ganzzahligen Coefficienten. Da nun die den Fundamentalpuncten zweiter Art entsprechenden \((a_i)\) durch die \(2p\) übrigen in derselben Weise ausgedrückt werden können, so hat man schliesslich die \(A_i\), und also auch die aus diesen zusammengesetzten \([A_i]\) als lineare Functionen der letztern \(2p\) Grössen \((a_i)\). Setzt man also
\[
{}[A_i]=c_{i1}a_1+c_{i2}a_2\cdots+c_{i,2p}a_{2p}
\]
indem wir jetzt der Bequemlichkeit wegen die Indices \(1,2\ldots2p\) für die übriggebliebenen Verzweigungspuncte anwenden), so ist die zwischen den \((a)\), \((b)\) bestehende Bedingung:
\[
(6)\quad \sum_{i,k=1}^{2p} c_{ik}b_i a_k=0.
\]
Von den ganzzahligen Coefficienten \(c_{ik}\) dieser für das Folgende fundamentalen Gleichung, werden wir jetzt zwei Haupteigenschaften beweisen.

Da die Gleichung
\[
\Sigma A_i b_i+\Sigma B_i a_i=0
\]
mit Hülfe der zwischen den \(a\) allein und zwischen den \(b\) allein bestehenden Gleichungen zu einer identischen wurde, so muss die aus der Anwendung der obigen Transformation beider Theile hervorgehende Gleichung
\[
\Sigma[A_i]b_i+\Sigma[B_k]a_k=0,
\]
in welcher zwischen den \((a_i)\), \((b_i)\) keine solchen Relationen mehr gelten, eine vollkommen identische sein. Man hat nun auch
\[
[B_k]=c_{k1}b_1+c_{k2}b_2\cdots+c_{k,2p}b_{2p},
\]
und die obige Identität geht also über in:
\[
\sum c_{ik}b_i a_k+\sum c_{ki}b_i a_k=0,
\]
was sich in die \(\dfrac{2p\cdot(2p+1)}{2}\) Gleichungen
\[
c_{ik}+c_{ki}=0
\]
auflöst. Alle Producte \((a_i)(b_i)\) verschwinden daher aus dem Ausdruck (6), und die Coefficienten von \((a_i)(b_k)\) und \((b_i)(a_k)\) sind gleich und entgegengesetzt. \emph{Die Coefficienten \(c_{ik}\) bilden also eine überschlagene Determinante.}

Schon hieraus folgt, dass die Determinante der \(c\) das Quadrat einer ganzen Zahl sein muss, da jede überschlagene Determinante gerader Ordnung das Quadrat eines aus ihren Elementen rational gebildeten Ausdrucks ist. Aber es lässt sich weiter der wichtige Satz beweisen:

\begin{center}
\emph{Die Determinante der \(c\) ist gleich 1.}
\end{center}

Dieser Satz, welcher gewissermassen der Angelpunct dieser Theorie ist, lässt sich folgendermassen beweisen.

Die \([A_i]\) lassen sich als lineare Functionen der \(2p\) übrig gebliebenen \((a_i)\) mit ganzzahligen Coefficienten darstellen. Der Werth der ganzzahligen Determinante der \(c\) sei \(C\). Kann man nun auch analog die \((a_i)\) als ganzzahlige lineare Functionen der \([A_i]\) darstellen, und ist die Determinante der letztern Ausdrücke \(\Gamma\), so sind \(C,\Gamma\) adjungirte Determinanten, also \(C\cdot\Gamma=1\), daher \(C=\Gamma=\pm1\). Und da \(C\) nothwendig positiv ist, so wird \(C=1\), was zu beweisen war.

Um nun die \(2p\) Grössen \((a_i)\), welche keinen Fundamentalpuncten entsprechen, durch die \(2p\) zugehörigen \([A]\) auszudrücken, betrachten wir folgende Reihenfolge von Gleichungen:

\begin{center}
\scriptsize
\begin{tabular}{p{0.43\linewidth}|p{0.43\linewidth}}
\centering I. & \centering II.\tabularnewline
Die \(A\) werden durch die \(a\) ausgedrückt, indem man Umgänge erster Art bildet, welche nur je einen Verzweigungspunct ausschliessen. &
Die \(a\) werden durch die \(A\) ausgedrückt, indem man Umgänge zweiter Art bildet, welche nur je einen Verzweigungspunct ausschliessen.\tabularnewline[0.6em]
\centering III. & \centering IV.\tabularnewline
Durch Beschreibung vollständiger Umgänge erster Art bildet man Gleichungen zwischen den \(a\). &
Durch Beschreibung vollständiger Umgänge zweiter Art bildet man Gleichungen zwischen den \(A\).\tabularnewline[0.6em]
\centering V. & \centering VI.\tabularnewline
Aus den \(a\) bildet man die \((a)\) durch Hinzufügung von \(a\), welche den Fundamentalpuncten erster Art entsprechen. Dabei verschwinden die den letztern zugehörigen \(a\). &
Man bildet aus den \(A\) die \((A)\), indem man die \(A\) aus I. entnimmt und in diesen Ausdrücken die \(a\) zu \((a)\) ergänzt. Von den \(A\) verschwindet im Allgemeinen keines.\tabularnewline[0.6em]
\centering VII. & \centering VIII.\tabularnewline
Die Gleichungen I. haben also fort zu bestehen, wenn man die \(a,A\) darin durch die \((a),A\) ersetzt. &
Die Gleichungen II. bestehen ebenfalls noch, wenn die \(a,A\) durch die \((a),(A)\) ersetzt werden. Denn bliebe von den \(a\), welche den Fundamentalpuncten erster Art entsprechen, eine Spur zurück, so müssten Integrale über nicht cyclische Wege allgemein auf cyclische zurückkommen, was nicht möglich ist.\tabularnewline[0.6em]
\centering IX. & \centering X.\tabularnewline
\multicolumn{2}{c}{Dasselbe gilt von den Bedingungen III. IV.}\tabularnewline[0.6em]
\centering XI. & \centering XII.\tabularnewline
Aus den \((A)\) bildet man die \([A]\), indem man \(A\) hinzufügt, welche den Fundamentalpuncten zweiter Art entsprechen. &
Aus den \((a)\) bildet man die \([a]\), indem man letztere dadurch definirt, dass die Gleichungen VIII. auch für die \([a],[A]\) zu bestehen fortfahren sollen.
\end{tabular}
\end{center}

Die Grösse \([a_i]\) drückt sich also durch die \((a_i)\) aus, vermehrt um diejenigen Combinationen \(h_{\alpha},h_{\beta}\) der den Fundamentalpuncten zweiter Art entsprechenden \((A)\), welche, aus den \(A\) gebildet, den unvollständigen zur Bildung der \(a\) beschriebenen Umgang der \(A\) schliessen würden (\(h_1\) durch Null ersetzt). Daher sind \(\alpha,\beta\) die Indices der vollständigen Umgänge zweiter Art, in denen der \(i^{\text{te}}\) Verzweigungspunct wesentlich ist, oder wegen der zwischen den beiden Schleifensystemen bestehenden Reciprocität die Indices derjenigen Wurzeln, welche die \(i^{\text{te}}\) Schleife erster Art mit einander verbindet. Nun sind aber einige der \((a_i)\) Null, nämlich diejenigen, welche den Fundamentalpuncten erster Art entsprechen. Daher werden die ihnen zugehörigen \([a_i]\) geradezu gleich den Differenzen der betreffenden \(h\). Und da die entsprechenden \(n-1\) Grössen \(a_i\), als Fundamentalpuncten erster Art angehörig, so combinirt werden können, dass sie von \(s_1\) zu jedem beliebigen andern \(s\) führen, so kann man die \(n-1\) \([a_i]\) ebenso combiniren, und erhält dann direct eine der Grössen \(h_{\alpha}\) durch diese \([a_i]\) ausgedrückt. Man kann also auch die Grössen \((a_i)\) jetzt durch die \([a_i]\), d. h. linear und ganzzahlig durch die \([A_i]\) ausdrücken, deren Anzahl gleich \(2p+n-1\) ist, weil die den Fundamentalpuncten zweiter Art entsprechenden \([A_i]\) verschwinden.

Nun bleiben aber die Gleichungen IV. X. offenbar noch bestehen, wenn man die \(A,(A)\) durch die \([A]\) ersetzt. Denn um von den \((A)\) zu den \([A]\) überzugehen, hat man nur immer die Differenzen der \(k\) hinzuzufügen, welche aus den \(A\) gebildet und den \(A\) hinzugefügt einen einfachen Cyclus zweiter Art hervorrufen würden. Die Indices zweier in einer solchen Differenz auftretenden \(k\) geben die Wurzeln an, welche eine Schleife zweiter Art verbindet. Da also jede folgende Differenz mit dem \(k\) beginnt, mit welchem die vorhergehende aufhörte, und die letzte mit dem \(k\) schliesst, mit dem die erste begann, so heben sich wirklich alle \(k\) auf und es bleiben nur die Gleichungen IV. X. mit den \([A]\) geschrieben zurück.

Wie wir aber früher bewiesen, dass sich aus den Gleichungen I. die den ersten Fundamentalpuncten entsprechenden \(a\) durch die übrigen ganzzahlig und linear ausdrücken lassen, so gilt reciprok der Satz, dass aus den Gleichungen II. sich die den zweiten Fundamentalpuncten entsprechenden \(A\) durch die übrigen ganzzahlig und linear ausdrücken. Setzen wir also in jenen Gleichungen statt der \(A\) die \([A]\), so erhalten wir die \(n-1\) den Fundamentalpuncten zweiter Art entsprechenden \([A]\) durch die \(2p\) übrigen ausgedrückt. Wir können daher auch die \((a_i)\), welche oben durch sämmtliche \([A_i]\) ausgedrückt waren, durch die \(2p\) letztern ganzzahlig und linear ausdrücken, was zu beweisen war.

\subsection*{\S\,28. Beispiel für die Aufstellung der Relation zwischen den Periodicitätsmoduln der Integrale erster Gattung.}

Wir schalten hier zur Erläuterung für den Fall \(n=4\) und für 12 Verzweigungspuncte folgendes Schema ein, in welchem die oben angedeuteten Operationen ausgeführt sind.



% Source PDF pp. 122--133 / printed pp. 100--111.
% Continuation of § 28 after the introductory sentence already present in the cumulative text.

\begin{center}
\scriptsize
\begin{tabular}{c@{\qquad}c@{\qquad}c}
\textbf{Index der}\\[-0.25em]
\textbf{Verzwei-}\\[-0.25em]
\textbf{gungspuncte.}
&
\begin{tabular}{c}
\textbf{Ihre Schleifen \(1^{\mathrm{ter}}\) Art ver-}\\
\textbf{binden die Wurzeln, oder sie}\\
\textbf{liegen auf den Umgängen}\\
\textbf{\(2^{\mathrm{ter}}\) Art.}
\end{tabular}
&
\begin{tabular}{c}
\textbf{Ihre Schleifen \(2^{\mathrm{ter}}\) Art}\\
\textbf{verbinden die Wurzeln, oder}\\
\textbf{liegen auf den Umgängen}\\
\textbf{\(1^{\mathrm{ter}}\) Art.}
\end{tabular}
\\[0.6em]\hline\\[-0.6em]
1. & \(1-2\) & \(1-2\)\\
2. & \(2-3\) & \(1-3\)\\
3. & \(1-4\) & \(2-4\)\\
4. & \(2-4\) & \(3-2\)\\
5. & \(3-4\) & \(1-3\)\\
6. & \(1-3\) & \(4-3\)\\
7. & \(2-4\) & \(2-1\)\\
8. & \(1-2\) & \(3-1\)\\
9. & \(2-3\) & \(3-4\)\\
10. & \(2-4\) & \(4-2\)\\
11. & \(3-4\) & \(3-4\)\\
12. & \(3-4\) & \(4-3\)
\end{tabular}
\end{center}

\noindent Relationen zwischen den \(a\) aus Umgängen \(1^{\mathrm{ter}}\) Art:
\[
\mathrm{I.}\quad
\begin{array}{r@{\quad}l}
1&0=a_1+a_2+a_3-a_7-a_8,\\
2&0=-a_1+a_3-a_4+a_7-a_{10},\\
3&0=-a_2+a_4-a_5-a_6+a_8+a_9+a_{11}-a_{12},\\
4&0=-a_3+a_5-a_9+a_{10}-a_{11}+a_{12}.
\end{array}
\]

\noindent Relationen zwischen den \(A\) aus Umgängen \(2^{\mathrm{ter}}\) Art:
\[
\mathrm{II.}\quad
\begin{array}{r@{\quad}l}
1&0=A_1+A_3+A_6+A_8,\\
2&0=-A_1+A_2+A_4+A_7-A_8+A_9+A_{10},\\
3&0=-A_2+A_5-A_6-A_9+A_{11}+A_{12},\\
4&0=-A_3-A_4-A_5-A_7-A_{10}-A_{11}-A_{12}.
\end{array}
\]

\noindent Ausdrücke der \(A\) durch die \(a\):
\[
\mathrm{III.}\quad
\begin{array}{rcl}
A_1&=&a_3-a_4+a_7-a_{10}\\
A_2&=&a_1+a_4-a_5-a_6+a_8+a_9+a_{11}-a_{12}\\
A_3&=&-a_1+a_6-a_9+a_{10}-a_{11}+a_{12}\\
A_4&=&-a_2+a_7-a_{10}\\
A_5&=&a_1+a_2-a_6+a_8+a_9+a_{11}-a_{12}\\
A_6&=&-a_3+a_8+a_9+a_{11}-a_{12}\\
A_7&=&-a_1+a_3-a_4-a_5\\
A_8&=&-a_2+a_4-a_5-a_6\\
A_9&=&-a_2+a_4-a_5-a_6+a_8+a_{10}-a_{11}+a_{12}\\
A_{10}&=&-a_3+a_6-a_9\\
A_{11}&=&-a_2+a_4-a_5-a_6+a_8+a_9+a_{12}\\
A_{12}&=&-a_3+a_6-a_9+a_{10}-a_{11}.
\end{array}
\]

\noindent Ausdrücke der \(a\) durch die \(A\):
\[
\mathrm{IV.}\quad
\begin{array}{rcl}
a_1&=&A_2+A_4+A_7-A_8+A_9+A_{10}\\
a_2&=&-A_1+A_5-A_6-A_9+A_{11}+A_{12}\\
a_3&=&A_1-A_4-A_5-A_7-A_{10}-A_{11}-A_{12}\\
a_4&=&-A_1+A_2-A_5-A_7-A_{10}-A_{11}-A_{12}\\
a_5&=&-A_2-A_7-A_{10}-A_{11}-A_{12}\\
a_6&=&A_1+A_3-A_9+A_{11}+A_{12}\\
a_7&=&-A_1+A_2+A_4-A_{10}-A_{11}-A_{12}\\
a_8&=&A_1+A_3+A_6+A_9+A_{10}\\
a_9&=&-A_1+A_2+A_4+A_7-A_8+A_{11}+A_{12}\\
a_{10}&=&-A_1+A_2+A_4+A_7-A_8+A_9-A_{11}-A_{12}\\
a_{11}&=&-A_2+A_5-A_6-A_9-A_{12}\\
a_{12}&=&-A_2+A_5-A_6-A_9+A_{11}.
\end{array}
\]

\noindent Fundamentalpuncte \(1^{\mathrm{ter}}\) Art:\quad \(1.\quad2.\quad3.\)
\hfill Fundamentalpuncte \(2^{\mathrm{ter}}\) Art:\quad \(4.\quad5.\quad6.\)

\noindent Combinationen, welche im ersten System von der Wurzel \(1\) zu \(2,3,4\) führen:
\[
g_2=a_1,\qquad g_3=a_1+a_2,\qquad g_4=a_3.
\]
Combinationen, welche im zweiten Systeme von der Wurzel \(1\) zu \(2,3,4\) führen:
\[
k_2=A_4+A_5,\qquad k_3=A_5,\qquad k_4=A_5-A_6.
\]

\noindent Bildung der rund eingeklammerten Grössen:
\[
\begin{array}{rcl@{\qquad\qquad}rcl}
(a_1)&=&0&(A_1)&=&A_1-a_1\\
(a_2)&=&0&(A_2)&=&A_2-a_1-a_2\\
(a_3)&=&0&(A_3)&=&A_3-a_3+a_1\\
(a_4)&=&a_4-a_3+a_1&(A_4)&=&A_4+a_2\\
(a_5)&=&a_5-a_3+a_1+a_2&(A_5)&=&A_5-a_1-a_2\\
(a_6)&=&a_6-a_1-a_2&(A_6)&=&A_6+a_3-a_1-a_2\\
(a_7)&=&a_7-a_3+a_1&(A_7)&=&A_7+a_1\\
(a_8)&=&a_8-a_1&(A_8)&=&A_8+a_1+a_2\\
(a_9)&=&a_9-a_2&(A_9)&=&A_9-a_3+a_1+a_2\\
(a_{10})&=&a_{10}-a_3+a_1&(A_{10})&=&A_{10}+a_3-a_1\\
(a_{11})&=&a_{11}-a_3+a_1+a_2&(A_{11})&=&A_{11}-a_3+a_1+a_2\\
(a_{12})&=&a_{12}-a_3+a_1+a_2&(A_{12})&=&A_{12}+a_3-a_1-a_2.
\end{array}
\]

\noindent Bildung der eckig eingeklammerten Grössen:
\[
\begin{array}{rcl@{\qquad\qquad}rcl}
{}[A_1]&=&(A_1)-(A_4)-(A_5)&[a_1]&=&-(A_4)-(A_5)\\
{}[A_2]&=&(A_2)-(A_5)&[a_2]&=&(A_4)\\
{}[A_3]&=&(A_3)+(A_4)+(A_6)&[a_3]&=&-(A_5)+(A_6)\\
{}[A_4]&=&0&[a_4]&=&(a_4)+(A_4)+(A_6)\\
{}[A_5]&=&0&[a_5]&=&(a_5)+(A_6)\\
{}[A_6]&=&0&[a_6]&=&(a_6)-(A_5)\\
{}[A_7]&=&(A_7)+(A_4)+(A_5)&[a_7]&=&(a_7)+(A_4)+(A_6)\\
{}[A_8]&=&(A_8)+(A_5)&[a_8]&=&(a_8)-(A_4)-(A_5)\\
{}[A_9]&=&(A_9)+(A_6)&[a_9]&=&(a_9)+(A_4)\\
{}[A_{10}]&=&(A_{10})-(A_4)-(A_6)&[a_{10}]&=&(a_{10})+(A_4)+(A_6)\\
{}[A_{11}]&=&(A_{11})+(A_6)&[a_{11}]&=&(a_{11})+(A_6)\\
{}[A_{12}]&=&(A_{12})-(A_6)&[a_{12}]&=&(a_{12})+(A_6).
\end{array}
\]

\noindent Ausdrücke der \([A]\) durch die \((a)\) und umgekehrt:
\[
\begin{array}{rcl}
{}[A_7]&=&-(a_4)-(a_6)+(a_7)+(a_9)-(a_{10})+(a_{11})-(a_{12})\\
{}[A_8]&=&(a_4)-(a_5)-2(a_6)+(a_8)+(a_9)+(a_{11})-(a_{12})\\
{}[A_9]&=&(a_4)-(a_5)-(a_6)+2(a_8)+(a_9)+(a_{10})\\
{}[A_{10}]&=&(a_6)-(a_7)-(a_8)-2(a_9)+(a_{10})-(a_{11})+(a_{12})\\
{}[A_{11}]&=&(a_4)-(a_5)-(a_6)+2(a_8)+2(a_9)+(a_{11})\\
{}[A_{12}]&=&(a_6)-(a_8)-2(a_9)+(a_{10})-2(a_{11})+(a_{12})\\[0.3em]
(a_7)&=&-2[A_1]+2[A_2]+2[A_7]-[A_8]+[A_9]+[A_{10}]\\
(a_8)&=&[A_1]-[A_2]+[A_3]-[A_7]+[A_8]\\
(a_9)&=&[A_2]+[A_7]-[A_8]+[A_9]\\
(a_{10})&=&-2[A_1]+2[A_2]+3[A_7]-2[A_8]+2[A_9]+2[A_{10}]\\
(a_{11})&=&-2[A_1]+2[A_7]-[A_8]-[A_9]+2[A_{10}]+2[A_{11}]+[A_{12}]\\
(a_{12})&=&-2[A_1]+2[A_7]-[A_8]-[A_9]+2[A_{10}]+3[A_{11}]+2[A_{12}].
\end{array}
\]

\noindent Ausdrücke von \((a_4),(a_5),(a_6),[A_1],[A_2],[A_3]\) aus den Relationen I. II.:
\[
\begin{array}{rcl}
(a_4)&=&(a_7)-(a_{10})\\
(a_5)&=&(a_7)+(a_8)\\
(a_6)&=&(a_9)-(a_{10})+(a_{11})-(a_{12})\\
{}[A_1]&=&[A_7]-[A_8]+[A_{10}]+[A_{11}]+[A_{12}]\\
{}[A_2]&=&-[A_9]+[A_{11}]+[A_{12}]\\
{}[A_3]&=&-[A_7]-[A_{10}]-[A_{11}]-[A_{12}].
\end{array}
\]

\noindent Definitive Ausdrücke der Grössen \((a_i),[A_i]\), welche keinen Fundamentalpuncten entsprechen:
\[
\begin{array}{rcl@{\qquad}rcl}
{}[A_7]&=&(a_{10})&(a_7)&=&[A_8]-[A_9]-[A_{10}]\\
{}[A_8]&=&-(a_9)+(a_{10})-(a_{11})+(a_{12})&(a_8)&=&-[A_7]+[A_9]-[A_{11}]-[A_{12}]\\
{}[A_9]&=&(a_8)+(a_{10})-(a_{11})+(a_{12})&(a_9)&=&[A_7]-[A_8]+[A_{11}]+[A_{12}]\\
{}[A_{10}]&=&-(a_7)-(a_8)-(a_9)&(a_{10})&=&[A_7]\\
{}[A_{11}]&=&(a_8)+(a_9)+(a_{12})&(a_{11})&=&[A_8]-[A_9]-[A_{12}]\\
{}[A_{12}]&=&-(a_8)-(a_9)-(a_{11})&(a_{12})&=&[A_8]-[A_9]+[A_{11}].
\end{array}
\]
Die überschlagene Determinante der \(c_{ik}\) ist daher
\[
\left|\begin{array}{rrrrrr}
 0& 0& 0& 1& 0& 0\\
 0& 0&-1& 1&-1& 1\\
 0& 1& 0& 1&-1& 1\\
-1&-1&-1& 0& 0& 0\\
 0& 1& 1& 0& 0& 1\\
 0&-1&-1& 0&-1& 0
\end{array}\right|.
\]

\subsection*{\S\,29. Transformation der Relation.}

Die Grössen \((a_i)\), welche keinen Fundamentalpuncten entsprechen, können als fundamentale Periodicitätsmoduln bezeichnet werden, durch welche alle andern Periodicitätsmoduln desselben Integrals linear und ganzzahlig sich ausdrücken. Zwischen den fundamentalen Periodicitätsmoduln je zweier verschiedener Integrale erster Gattung besteht die Gleichung
\[
\tag{1}\sum c_{ik}(a_i)(b_k)=0,
\]
wo die \(c\) jedesmal wieder dieselben Zahlen bedeuten.

Statt der Periodicitätsmoduln \((a_i)\) kann man aber, um alle möglichen Periodicitätsmoduln linear und ganzzahlig auszudrücken, irgend welche \(2p\) ganzzahlige lineare Combinationen derselben anwenden, deren Determinante \(1\) ist. Wir können also statt der \((a_i)\) die \(2p\) neuen Periodicitätsmoduln einführen:
\[
\tag{2}M^{(h)}=\alpha_1^{(h)}(a_1)+\alpha_2^{(h)}(a_2)+\cdots+\alpha_{2p}^{(h)}(a_{2p}),\quad(h=1,2..2p)
\]
und also gleichzeitig statt der \((b)\) die \(2p\) neuen Periodicitätsmoduln:
\[
\tag{3}N^{(h)}=\alpha_1^{(h)}(b_1)+\alpha_2^{(h)}(b_2)+\cdots+\alpha_{2p}^{(h)}(b_{2p}),\quad(h=1,2..2p)
\]
und man kann, sobald die Determinante der ganzen Zahlen \(\alpha\) gleich \(1\) ist, alle Periodicitätsmoduln des ersten Integrals durch die \(M\), alle des zweiten durch die \(N\) ausdrücken. Sind dann die Auflösungen der Systeme (2), (3):
\[
\tag{4}
\left\{\begin{array}{rcl}
(a_i)&=&\beta_i^{(1)}M^{(1)}+\beta_i^{(2)}M^{(2)}+\cdots+\beta_i^{(2p)}M^{(2p)}\\
(b_i)&=&\beta_i^{(1)}N^{(1)}+\beta_i^{(2)}N^{(2)}+\cdots+\beta_i^{(2p)}N^{(2p)},
\end{array}\right.
\]
so sind die \(\beta\) wiederum ganze Zahlen, deren Determinante \(1\) ist, und die Gleichung (1) verwandelt sich in
\[
\tag{5}\sum C^{(h,m)}M^{(h)}N^{(m)}=0,
\]
wo
\[
\tag{6}C^{(h,m)}=\sum c_{ik}\beta_i^{(h)}\beta_k^{(m)}.
\]
Hieraus sieht man, dass erstlich immer wieder
\[
C^{(h,m)}=-C^{(m,h)},
\]
denn wenn man \(h\) mit \(m\), und gleichzeitig (was irrelevant ist) \(i\) mit \(k\) vertauscht, so ändert die rechte Seite von (6) nur ihr Zeichen. Zweitens ist nach bekannten Determinanten-Sätzen die Determinante der \(C\) gleich der Determinante der \(c\) multiplicirt mit dem Quadrat der Determinante der \(\beta\), also gleich \(1\). Daher hat die Determinante der \(C\) genau die Eigenschaften, welche der Determinante der \(c\) zukommen.

Hieraus folgt weiter, dass auch durch eine Reihenfolge von linearen Substitutionen, wo an Stelle der \((a),(b)\) successive immer andere fundamentale Periodicitätsmoduln auftreten, diese Eigenschaften der Coefficienten fortwährend erhalten bleiben.

Unter den linearen Substitutionen giebt es nun solche, welche auf die Determinante der \(c\) dieselbe Wirkung haben, als wenn in derselben eine Verticalreihe, mit einem Factor multiplicirt, einer andern hinzugefügt wird, und zugleich dasselbe mit den entsprechenden Horizontalreihen geschieht. Setzen wir nämlich
\[
(a_i)=M^{(i)},
\]
und nur für einen bestimmten Werth \(i=\rho\) davon abweichend
\[
(a_\rho)=M^{(\rho)}+mM^{(\sigma)}\quad(\sigma\ \text{von}\ \rho\ \text{verschieden}).
\]
Die Determinante der \(\beta\) enthält dann in der Diagonale lauter \(1\), übrigens Nullen, bis auf eine einzige Stelle, in welcher das Element \(m\) erscheint. Daher hat die Determinante der \(\beta\) wirklich den Werth \(1\). Die linke Seite der Gleichung (1) geht über in
\[
\sum c_{ik}M^{(i)}N^{(k)}+
m\left[M^{(\sigma)}\sum_k c_{\rho k}N^{(k)}+
N^{(\sigma)}\sum_i c_{i\rho}M^{(i)}\right].
\]
Man hat also im Allgemeinen
\[
C^{(ik)}=c_{ik},
\]
insbesondere aber
\[
C^{(i\sigma)}=c_{i\sigma}+mc_{i\rho},\qquad
C^{(\sigma i)}=c_{\sigma i}+mc_{\rho i},
\]
und
\[
C^{(\sigma\sigma)}=c_{\sigma\sigma}+m(c_{\sigma\rho}+c_{\rho\sigma})=0.
\]
Die Determinante der \(C\) entsteht also aus der Determinante der \(c\), wenn man die \(\rho^{\mathrm{te}}\) Verticalreihe mit \(m\) multiplicirt zu der \(\sigma^{\mathrm{ten}}\) addirt, und mit den Horizontalreihen dasselbe ausführt. Diese Operation ist äquivalent mit der linearen Transformation
\[
\begin{array}{rcl@{\qquad}rcl}
(a_\rho)&=&M^{(\rho)}+mM^{(\sigma)},&(a_i)&=&M^{(i)}\\
M^{(\rho)}&=&(a_\rho)-m(a_\sigma),&M^{(i)}&=&(a_i)
\end{array}
\quad(i=1,2..\rho-1,\rho+1..2p).
\]

Sehen wir nun, wie wir die Determinante der \(c\) durch eine Reihenfolge solcher Operationen auf eine möglichst einfache Form bringen können.

Da die Determinante den Werth \(1\) hat, so haben die Elemente einer Reihe, z. B. der ersten Horizontalreihe, keinen gemeinschaftlichen Factor. Daher können wir durch passendes Zufügen von Verticalreihen es dahin bringen, dass eine Zahl der ersten Horizontalreihe \(1\) wird; und sofort können wir dann die übrigen Glieder derselben Horizontal- und Verticalreihe mit Hülfe dieser Glieder zerstören. Wir können auch die \(1\) an jede beliebige Stelle der ersten Verticalreihe bringen und dann die \(1\) an ihrer frühern Stelle zerstören. Daher kann man es immer so einrichten, dass die zweite Stelle der ersten Horizontalreihe \(1\) ist, alle andern derselben Horizontalreihe und der zweiten Verticalreihe aber aus Nullen bestehen; und da wir immer die entsprechenden Operationen mit den Verticalreihen ausführen, so wird dann die erste Verticalreihe an zweiter Stelle \(-1\), an allen übrigen \(0\) enthalten, und die zweite Horizontalreihe an erster Stelle \(-1\), an allen übrigen Stellen aber \(0\). Man kann nun weiter mit der dritten und vierten Reihe genau so verfahren, wie soeben mit der ersten und zweiten, sodann ebenso mit der fünften und sechsten u. s. w., so dass endlich die Determinante der \(c\) die Gestalt annimmt:
\[
\left|\begin{array}{rrrrrrr}
0&1&0&0&\cdots&0&0\\
-1&0&0&0&\cdots&0&0\\
0&0&0&1&\cdots&0&0\\
0&0&-1&0&\cdots&0&0\\
\multicolumn{7}{c}{\cdots\qquad\qquad\cdots}\\
0&0&0&0&\cdots&0&1\\
0&0&0&0&\cdots&-1&0
\end{array}\right|.
\]
Es ist also endlich durch successive Transformation eine solche Transformation entstanden, für welche in der transformirten Gleichung (1) alle \(c_{ik}\) verschwinden bis auf
\[
c_{12},\ c_{34},\ c_{56}\ldots c_{2p-1,2p},
\]
welche gleich \(1\) sind. Oder mit andern Worten: Man kann immer solche ganzzahlige lineare Transformationsformeln
\[
\begin{array}{rcl}
M^{(h)}&=&\alpha_1^{(h)}(a_1)+\alpha_2^{(h)}(a_2)+\cdots+\alpha_{2p}^{(h)}(a_{2p})\\
N^{(h)}&=&\alpha_1^{(h)}(b_1)+\alpha_2^{(h)}(b_2)+\cdots+\alpha_{2p}^{(h)}(b_{2p})
\end{array}
\]
aufstellen, durch welche identisch
\[
\begin{split}
\sum c_{ik}(a_i)(b_k)
&=(M^{(1)}N^{(2)}-N^{(1)}M^{(2)})+(M^{(3)}N^{(4)}-N^{(3)}M^{(4)})\\
&\quad+\cdots+(M^{(2p-1)}N^{(2p)}-N^{(2p-1)}M^{(2p)})
\end{split}
\]
wird.

Diese Combinationen der \((a),(b)\), auf welche hier die Transformation der Gleichung (1) geführt hat, wollen wir normale Periodicitätsmoduln nennen; es giebt unendlich viele Systeme von normalen Periodicitätsmoduln. In einem solchen System sind die normalen Periodicitätsmoduln eines Integrals paarweise conjugirt (insofern die mit den Indices \(1\) und \(2\), \(3\) und \(4,\ldots,2p-1\) und \(2p\) behafteten Grössen in besondern Beziehungen zu einander stehen), und zwischen den normalen Periodicitätsmoduln zweier Integrale erster Gattung besteht die Relation:
\[
\tag{7}\sum_{i=1}^{i=p}(M^{(2i-1)}N^{(2i)}-N^{(2i-1)}M^{(2i)})=0.\footnotemark
\]
\footnotetext{In dieser Form haben die zwischen den Periodicitätsmoduln auftretenden Relationen Rosenhain und Weierstrass für die hyperelliptischen, Riemann für die Abel'schen Functionen gegeben.}

\subsection*{\S\,30. Normalform der Integrale erster Gattung.}

Diese Betrachtungen endlich führen auf eine bestimmte Normalform der Integrale erster Gattung.

Bisher hatten wir unter den \(p\) Integralen erster Gattung Ausdrücke verstanden, welche aus irgend \(p\) Integralen erster Gattung linear zusammengesetzt waren. Jedes solche Integral enthält also noch \(p\) willkürliche Constanten. Sind \(H_1,H_2\ldots H_p\) irgend welche von einander unabhängige Integrale erster Gattung (d. h. solche, zwischen denen keine Relation mit constanten Coefficienten besteht), so sind die allgemeinsten \(p\) Integrale erster Gattung, welche wir zu Grunde legen können:
\[
\begin{array}{rcl}
I_1&=&m_{11}H_1+m_{12}H_2+\cdots+m_{1p}H_p\\
I_2&=&m_{21}H_1+m_{22}H_2+\cdots+m_{2p}H_p\\
&\vdots&\\
I_p&=&m_{p1}H_1+m_{p2}H_2+\cdots+m_{pp}H_p,
\end{array}
\]
wo die Determinante der willkürlichen Constanten \(m\), damit auch die \(I\) von einander unabhängig seien, nicht verschwinden darf.

Bilden wir für diese \(I\) ein System von normalen Periodicitätsmoduln, so verwandeln die \(H\) sich in Integrale zwischen bestimmten Grenzen, also in gegebene Constante. Wir können daher die \(m\) so bestimmen, dass \(p\) der normalen Periodicitätsmoduln jedes Integrals bestimmte Werthe annehmen. Bezeichnen wir die normalen Periodicitätsmoduln von \(I_k\) durch \(I_k^{(1)}, I_k^{(2)}\ldots I_k^{(2p)}\), die entsprechenden der \(H_k\) durch \(H_k^{(1)}, H_k^{(2)}\ldots H_k^{(2p)}\), so wollen wir die \(m_{ik}\) so bestimmen, dass alle \(I_k^{(h)}\) mit oberem ungeraden Index verschwinden, mit Ausnahme der \(I_k^{(2k-1)}\), welche den gemeinsamen Werth \(2i\pi\) annehmen sollen. Wir haben dann zur Bestimmung von \(m_{i1},m_{i2}\ldots m_{ip}\) folgendes System linearer Gleichungen:
\[
\begin{array}{rcl}
m_{i1}H_1^{(1)}+m_{i2}H_2^{(1)}+\cdots+m_{ip}H_p^{(1)}&=&0\\
m_{i1}H_1^{(3)}+m_{i2}H_2^{(3)}+\cdots+m_{ip}H_p^{(3)}&=&0\\
\multicolumn{3}{c}{\cdots}\\
m_{i1}H_1^{(2i-1)}+m_{i2}H_2^{(2i-1)}+\cdots+m_{ip}H_p^{(2i-1)}&=&2i\pi\\
\multicolumn{3}{c}{\cdots}\\
m_{i1}H_1^{(2p-1)}+m_{i2}H_2^{(2p-1)}+\cdots+m_{ip}H_p^{(2p-1)}&=&0,
\end{array}
\]
und also, wenn \(\Delta\) die Determinante der \(H\) mit ungeradem obern Index, \(\Delta_k^{(2m-1)}\) eine ihrer Unterdeterminanten bezeichnet:
\[
m_{ik}=\frac{2\pi\sqrt{-1}}{\Delta}\Delta_k^{(2i-1)}.
\]
Die Normalintegrale erster Gattung, welche auf diese Weise aus \(I_1,I_2\ldots I_p\) entstehen, wollen wir fortan durch \(u_1,u_2\ldots u_p\) bezeichnen. Die erste Hälfte des Systems von normalen Periodicitätsmoduln der Normalintegrale ist gegeben durch das Schema:
\[
\begin{array}{c@{\quad}cccc}
u_1&2i\pi&0&\cdots&0\\
u_2&0&2i\pi&\cdots&0\\
\vdots&&\ddots&&\vdots\\
u_p&0&0&\cdots&2i\pi.
\end{array}
\]
Gehen wir nun auf die Gleichung (7) zurück, um zu untersuchen, welche Relationen zwischen den übrigen Normalperiodicitätsmoduln eintreten, so können wir zunächst jene Gleichung, auf zwei Integrale \(I_h,I_k\) angewandt, in der Form schreiben:
\[
\sum_{i=1}^{i=p}\left(I_h^{(2i-1)}I_k^{(2i)}-I_k^{(2i-1)}I_h^{(2i)}\right)=0.
\]
Da aber \(I_h^{(2i-1)}\) und \(I_k^{(2i-1)}\) verschwinden, bis auf \(I_h^{(2h-1)}\) und \(I_k^{(2k-1)}\), welche gleich \(2i\pi\) werden, so geht diese Gleichung über in die einfache Form:
\[
I_k^{(2h)}=I_h^{(2k)}.
\]
Setzen wir also
\[
I_k^{(2h)}=I_h^{(2k)}=a_{hk},
\]
wo \(a_{hk}\) durch Vertauschung der Indices sich nicht ändern soll, so sind die aus (7) fliessenden Gleichungen erfüllt, und man hat den Satz:

Die zweite Hälfte der normalen Periodicitätsmoduln von Normalintegralen erster Gattung bildet ein System von symmetrischer Determinante:
\[
\begin{array}{c@{\quad}cccc}
u_1&a_{11}&a_{12}&\cdots&a_{1p}\\
u_2&a_{21}&a_{22}&\cdots&a_{2p}\\
\vdots&&\cdots&&\\
u_p&a_{p1}&a_{p2}&\cdots&a_{pp}.
\end{array}
\]
Und die allgemeinsten Werthe, welche die Normalintegrale \(u_1,u_2\ldots u_p\) durch Abänderung des Integrationsweges annehmen können, sind:
\[
\begin{array}{rcl}
u_1'&=&u_1+2i\pi m_1+a_{11}q_1+a_{12}q_2+\cdots+a_{1p}q_p\\
u_2'&=&u_2+2i\pi m_2+a_{21}q_1+a_{22}q_2+\cdots+a_{2p}q_p\\
&\vdots&\\
u_p'&=&u_p+2i\pi m_p+a_{p1}q_1+a_{p2}q_2+\cdots+a_{pp}q_p,
\end{array}
\]
wo die \(m,q\) beliebige positive oder negative ganze Zahlen bedeuten (vgl. Riemann).

\subsection*{\S\,31. Ausführung der Transformation für das in \S\,28 behandelte Beispiel.}

Wir wollen an dem oben behandelten Beispiel die Aufsuchung der normalen Periodicitätsmoduln durchführen. Unsere Determinante war:
\[
\left|\begin{array}{rrrrrr}
0&0&0&1&0&0\\
0&0&-1&1&-1&1\\
0&1&0&1&-1&1\\
-1&-1&-1&0&0&0\\
0&1&1&0&0&1\\
0&-1&-1&0&-1&0
\end{array}\right|.
\]
Die dritte Horizontal-, resp. Verticalreihe von der zweiten abgezogen, liefert:
\[
\mathrm{I.}\quad
\left|\begin{array}{rrrrrr}
0&0&0&1&0&0\\
0&0&-1&0&0&0\\
0&1&0&1&-1&1\\
-1&0&-1&0&0&0\\
0&0&1&0&0&1\\
0&0&-1&0&-1&0
\end{array}\right|.
\]
Die erste Reihe von der dritten abgezogen:
\[
\mathrm{II.}\quad
\left|\begin{array}{rrrrrr}
0&0&0&1&0&0\\
0&0&-1&0&0&0\\
0&1&0&0&-1&1\\
-1&0&0&0&0&0\\
0&0&1&0&0&1\\
0&0&-1&0&-1&0
\end{array}\right|.
\]
Die zweite Reihe zu der fünften addirt:
\[
\mathrm{III.}\quad
\left|\begin{array}{rrrrrr}
0&0&0&1&0&0\\
0&0&-1&0&0&0\\
0&1&0&0&0&1\\
-1&0&0&0&0&0\\
0&0&0&0&0&1\\
0&0&-1&0&-1&0
\end{array}\right|.
\]
Die zweite Reihe von der sechsten abgezogen:
\[
\mathrm{IV.}\quad
\left|\begin{array}{rrrrrr}
0&0&0&1&0&0\\
0&0&-1&0&0&0\\
0&1&0&0&0&0\\
-1&0&0&0&0&0\\
0&0&0&0&0&1\\
0&0&0&0&-1&0
\end{array}\right|.
\]
Diese Determinante hat, abgesehen von der Ordnung der Reihen, auf die es schliesslich nicht ankommt, die gesuchte Form.

Die \(4\) Umformungen der Determinante geben folgende Substitutionen:
\[
\begin{array}{c|c|c|c}
\mathrm{I.}&\mathrm{II.}&\mathrm{III.}&\mathrm{IV.}\\
(a_7)=a_7'&a_7'=a_7''-a_9''&a_7''=a_7'''&a_7'''=M_7\\
(a_8)=a_8'&a_8'=a_8''&a_8''=a_8'''+a_{11}'''&a_8'''=M_8-M_{12}\\
(a_9)=a_9'-a_8'&a_9'=a_9''&a_9''=a_9'''&a_9'''=M_9\\
(a_{10})=a_{10}'&a_{10}'=a_{10}''&a_{10}''=a_{10}'''&a_{10}'''=M_{10}\\
(a_{11})=a_{11}'&a_{11}'=a_{11}''&a_{11}''=a_{11}'''&a_{11}'''=M_{11}\\
(a_{12})=a_{12}'&a_{12}'=a_{12}''&a_{12}''=a_{12}'''&a_{12}'''=M_{12}
\end{array}
\]
Zwischen den ersten und letzten Variabeln bestehen also die Gleichungen:
\[
\begin{array}{rcl@{\qquad}rcl}
(a_7)&=&M_7-M_9&M_7&=&(a_7)+(a_9)+(a_8)\\
(a_8)&=&M_8+M_{11}-M_{12}&M_8&=&(a_8)-(a_{11})+(a_{12})\\
(a_9)&=&-M_8+M_9-M_{11}+M_{12}&M_9&=&(a_8)+(a_9)\\
(a_{10})&=&M_{10}&M_{10}&=&(a_{10})\\
(a_{11})&=&M_{11}&M_{11}&=&(a_{11})\\
(a_{12})&=&M_{12}&M_{12}&=&(a_{12}).
\end{array}
\]
und setzt man ebenso
\[
\begin{array}{rcl@{\qquad}rcl}
(b_7)&=&N_7-N_9&N_7&=&(b_7)+(b_9)+(b_8)\\
(b_8)&=&N_8+N_{11}-N_{12}&N_8&=&(b_8)-(b_{11})+(b_{12})\\
(b_9)&=&-N_8+N_9-N_{11}+N_{12}&N_9&=&(b_8)+(b_9)\\
(b_{10})&=&N_{10}&N_{10}&=&(b_{10})\\
(b_{11})&=&N_{11}&N_{11}&=&(b_{11})\\
(b_{12})&=&N_{12}&N_{12}&=&(b_{12}),
\end{array}
\]
so geht die zwischen den \((a),(b)\) bestehende Gleichung über in:
\[
0=(M_{10}N_7-M_7N_{10})+(M_8N_9-M_9N_8)+(M_{12}N_{11}-M_{11}N_{12}).
\]
Bezeichnet man also die den \(3\) verschiedenen Integralen erster Gattung entsprechenden Periodicitätsmoduln durch \(M,N,P\), so hat man die drei Gleichungen:
\[
\begin{array}{l}
(M_{10}N_7-M_7N_{10})+(M_8N_9-M_9N_8)+(M_{12}N_{11}-M_{11}N_{12})=0\\
(N_{10}P_7-N_7P_{10})+(N_8P_9-N_9P_8)+(N_{12}P_{11}-N_{11}P_{12})=0\\
(P_{10}M_7-P_7M_{10})+(P_8M_9-P_9M_8)+(P_{12}M_{11}-P_{11}M_{12})=0.
\end{array}
\]
Bestimmt man also die Constanten der Integrale so, dass:
\[
\begin{array}{lll}
M_7=2i\pi,&N_7=0,&P_7=0,\\
M_9=0,&N_9=2i\pi,&P_9=0,\\
M_{11}=0,&N_{11}=0,&P_{11}=2i\pi,
\end{array}
\]
so gehen die obigen Gleichungen über in:
\[
M_8=N_{10},\qquad N_{12}=P_8,\qquad P_{10}=M_{12},
\]
und das vollständige System der normalen Periodicitätsmoduln wird:
\[
\begin{array}{llllll}
M_7=2i\pi,&M_9=0,&M_{11}=0,&M_{10}=a_{11},&M_8=a_{12},&M_{12}=a_{13},\\
N_7=0,&N_9=2i\pi,&N_{11}=0,&N_{10}=a_{12},&N_8=a_{22},&N_{12}=a_{23},\\
P_7=0,&P_9=0,&P_{11}=2i\pi,&P_{10}=a_{13},&P_8=a_{23},&P_{12}=a_{33}.
\end{array}
\]


% Source PDF pp. 134--149 / printed pp. 112--127.
\section*{Fünfter Abschnitt.}
\subsection*{Vertauschung von Parameter und Argument.}

\subsubsection*{\S\,32. Nähere Bestimmung des Integrals dritter Gattung.}

Ein Integral dritter Gattung verhält sich bezüglich der Verzweigungspuncte ganz ähnlich wie ein Integral erster Gattung; es ändert im Allgemeinen seinen Werth, wenn der Weg, über den das Integral geführt wird, bei seiner Verschiebung einen Verzweigungspunct überschreitet. Aber dazu kommt, dass der Werth eines Integrals dritter Gattung sich auch dann ändert, wenn der Weg bei seiner Verschiebung einen der beiden Unstetigkeitspuncte überschreitet. Man kann daher den allgemeinsten Werth des Integrals aus einem derselben ableiten, indem man erstlich cyclische, um Verzweigungspuncte führende Wege, sodann aber einen weitern Weg hinzufügt, welcher einen oder den andern der beiden Unstetigkeitspuncte einschliesst.

Im ersten Abschnitt \S\,7 wurde gezeigt, dass unser Normalintegral dritter Gattung sich in der Nähe des ersten Unstetigkeitspuncts \(\xi\) wie \(-\log(x-\xi)\), in der Nähe des zweiten aber wie \(\log(x-\eta)\) verhält. Bei einem Umlauf in kleinem Kreise, im Sinne der Bewegung eines Uhrzeigers, ändert sich also, wenn der kleine Kreis den ersten Unstetigkeitspunct zum Mittelpunct hat, das Integral dritter Gattung um \(2\pi\sqrt{-1}\), dagegen um \(-2\pi\sqrt{-1}\), wenn der zweite Unstetigkeitspunct der Mittelpunct ist. Daher wird durch einen Umlauf um die beiden Unstetigkeitspuncte zugleich das Integral gar nicht, durch Umläufe um einen derselben das Integral nur um \(2m\pi\sqrt{-1}\) geändert, wo \(m\) eine positive oder negative ganze Zahl oder Null ist.

Führen wir nun auch hier wie bei den Integralen erster Gattung den Begriff des Periodicitätsmoduls so ein, dass derselbe das Integral über einen cyclischen Weg bedeutet, und sind \(S_{\xi\eta}^{(1)}\), \(S_{\xi\eta}^{(2)}\), \(\ldots\), \(S_{\xi\eta}^{(2p)}\) von einander unabhängige Periodicitätsmoduln des Integrals \(S_{\xi\eta}\), so ist der allgemeinste Werth, welchen dies Integral durch Abänderung des Integrationsweges anzunehmen im Stande ist:
\[
\Sigma_{\xi\eta}=S_{\xi\eta}+2m\pi\sqrt{-1}+m^{(1)}S_{\xi\eta}^{(1)}+m^{(2)}S_{\xi\eta}^{(2)}+\cdots+m^{(2p)}S_{\xi\eta}^{(2p)},
\]
wo die \(m\) ganze Zahlen sind.

Wir wollen nun im Folgenden den Werth von \(S_{\xi\eta}\) insofern beschränken, als wir die Unstetigkeitspuncte \(\xi,\eta\) durch eine Curve verbinden wollen, welche von \(\xi\) und dem zugehörigen Werthe \(s_{\xi}\) zu \(\eta\) und dem zugehörigen Werthe \(s_{\eta}\) führt; welche keine der zu den Verzweigungspuncten führenden Schleifen so schneidet, dass im Schnittpuncte auch die Werthe von \(s_x\) übereinstimmen; diese Curve setzen wir uns vor, niemals zu überschreiten. Die ganzen Integrale \(S_{\xi\eta}^{(1)}\), \(S_{\xi\eta}^{(2)}\), \(\ldots\), welche wir uns zunächst über eine Anzahl von Schleifen geführt dachten, kann man dann zwar über äquivalente Wege führen, indem man die Schleifensysteme zu geschlossenen Curven erweitert, aber auch diese dürfen nie über einen der Puncte \(\xi,\eta\) geschoben werden, so dass jene Integrale wirklich vollkommen bestimmte Werthe erhalten. In der obigen Formel fällt zugleich das Glied \(2m\pi\sqrt{-1}\) aus, wenn die \(S_{\xi\eta}^{(1)}\), \(S_{\xi\eta}^{(2)}\) etc. auf die angegebene Weise verstanden werden.

Wir wollen ferner die Curve, welche beide Unstetigkeitspuncte verbindet, zu einer Schleife erweitern, welche von \(0\) zu einem Puncte jener Curve, auf ihr bis zu dem Unstetigkeitspuncte \(\xi\), in kleinem Kreise, und zwar im Sinne des Zeigers einer Uhr, um diesen herum, längs der Curve zu dem andern Unstetigkeitspuncte, um diesen herum, wieder eine Strecke längs der Curve, und auf demselben von \(0\) nach der Curve laufenden Wege nach \(0\) wieder zurück geht. Diese Schleife, welche keine der nach den Verzweigungspuncten geführten so schneiden darf, dass für den Schnittpunct auch die Werthe von \(s_x\) übereinstimmen, soll dem frühern Schleifensystem hinzugefügt werden. Sie ist einem bestimmten der in \(0\) stattfindenden Wurzelwerthe zugeordnet; denn nur von einem aus wird man bei den Unstetigkeitspuncten auch auf die betreffenden Werthe von \(s_x\) geführt, während für andere Anfangswerthe jene Schleife nicht existirt.

Aber das Integral, über diese neue Schleife geführt, giebt identisch Null; denn nach bekannten Sätzen kann man die Schleife (welche für die Werthe von \(s_x\) ein geschlossener Weg ist) auf zwei um \(\xi\) und \(\eta\) in gleichem Sinne beschriebene Kreise zurückführen; und da der Werth des Integrals über beide Kreise genommen gleich und entgegengesetzt ist \((2\pi\sqrt{-1}\) und \(-2\pi\sqrt{-1})\), so ist die Summe der entstehenden Werthe Null.

Bildet man also das Integral \(S_{\xi\eta}\), geführt um einen der geschlossenen Umgänge, von welchen im vorigen Abschnitte die Rede war, und welche jetzt nur um die neue Schleife erweitert sind, so erhält man Null gleich einer Summe von Integralen über einzelne Schleifen, welche man, wie im Vorigen, zu Periodicitätsmoduln ergänzen kann; aber die neue Schleife hat an diesen Gleichungen keinen Antheil. Die linearen Relationen zwischen ganzen Integralen dritter Gattung sind also genau dieselben, wie die Relationen zwischen den entsprechenden Periodicitätsmoduln der Integrale erster Gattung.

\subsubsection*{\S\,33. Vertauschung von Parameter und Argument.}

Betrachten wir nun zwei Integrale dritter Gattung, \(S_{\xi\eta}\), \(S_{\alpha\beta}\), deren Unstetigkeitspuncte durch die angefügten Indices bezeichnet werden. Bilden wir für jedes dieser Integrale die seine Unstetigkeitspuncte umkreisende Schleife, und fügen diese den Umgängen ein; wobei wir die Lage der Schleifen so einrichten wollen, dass die letztern weder einander noch die nach den Verzweigungspuncten gerichteten Schleifen in der oben bezeichneten verbotenen Weise schneiden.

Das Integral \(\int S_{\xi\eta}\,dS_{\alpha\beta}\), über einen vollständigen Umgang geführt, giebt dann Null, weil man es auf einen Punct zusammenziehn kann. Dasselbe Integral, über alle \(n\) Umgänge nach einander ausgeführt, giebt daher Veranlassung zu einer Gleichung, welche der im vorigen Abschnitte entwickelten ganz ähnlich ist. Da nämlich der Werth von \(S_{\xi\eta}\) oder von \(S_{\alpha\beta}\) nach Durchlaufung einer um die Unstetigkeitspuncte gelegten Schleife sich in keiner Weise ändert, so haben diese Schleifen auf die Functionswerthe der folgenden gar keinen Einfluss. So weit also das Resultat von diesen übrigen Schleifen herrührt, hat es genau die Form des im vorigen Abschnitte für Integrale erster Gattung gewonnenen, zumal auch die zwischen den ganzen Integralen bestehenden Bedingungen genau dieselben wie früher sind. Bezeichnen wir also die von einander unabhängigen Periodicitätsmoduln der beiden Integrale \(S_{\xi\eta}\) und \(S_{\alpha\beta}\) durch \(A_i,B_i\) \((i=1,2,\ldots,2p)\), so wird der von den Schleifen der Verzweigungspuncte herrührende Theil unsers Integrals
\[
\sum c_{ik}B_iA_k,
\]
wo die ganzen Zahlen \(c\) die denselben im vorigen Abschnitte beigelegten Werthe haben.

Was nun den Antheil des Integrals angeht, welcher von den Schleifen der Unstetigkeitspuncte herrührt, so braucht man dabei offenbar nur auf die Theile derselben einzugehen, welche \(\xi\) mit \(\eta\), resp. \(\alpha\) mit \(\beta\) rückwärts und vorwärts verbinden. Der zwischen \(\xi\) und \(\eta\) liegende Weg wird einmal in der Richtung von \(\xi\) nach \(\eta\), das andre Mal entgegengesetzt durchlaufen; an entsprechenden Stellen hat also \(dS_{\alpha\beta}\) den gleichen aber entgegengesetzten Werth. An den entsprechenden Stellen beider Wege ist aber der Werth des Integrals \(S_{\xi\eta}\) um \(2\pi\sqrt{-1}\) verschieden, weil inzwischen einer der Unstetigkeitspuncte, \(\eta\), umkreist wurde. Der Werth des Elements \(S_{\xi\eta}dS_{\alpha\beta}\) ist also das erste Mal \(S_{\xi\eta}dS_{\alpha\beta}\), das zweite Mal \(-(S_{\xi\eta}-2\pi\sqrt{-1})dS_{\alpha\beta}\), sodass beide zusammen den Werth \(2\pi\sqrt{-1}\,dS_{\alpha\beta}\) liefern. Der von der ersten Schleife herrührende Bestandtheil unsers Integrals wird also
\[
2\pi\sqrt{-1}\int_{\xi}^{\eta}dS_{\alpha\beta}.\footnotemark
\]
\footnotetext{Die Bezeichnung der Grenzen verstehen wir so, dass \(\xi,\eta\) nicht nur die Variabeln, sondern die Puncte \(\xi,\eta\) der Curve \(f=0\) bedeuten sollen, wobei denn nicht nur eine Coordinate, sondern beide Coordinaten derselben als gegeben vorausgesetzt werden.}

Was nun die um \(\alpha,\beta\) führende Schleife anbetrifft, so können wir sie auf zwei um \(\alpha,\beta\) beschriebene kleine Kreise zusammenziehen, wo denn \(S_{\xi\eta}\) beziehungsweise die constanten Werthe erhält, welche ihm in den Puncten \(\alpha,\beta\) zukommen. Da die Schleife sich im Sinne eines Uhrzeigers um die Puncte \(\alpha,\beta\) bewegen sollte, so giebt die Integration von \(dS_{\alpha\beta}\) bei \(\alpha\) den Werth \(2\pi\sqrt{-1}\), bei \(\beta\) den Werth \(-2\pi\sqrt{-1}\), und der von dieser Schleife herrührende Bestandtheil ist also die Summe von \(2\pi\sqrt{-1}\) multiplicirt mit dem Werth von \(S_{\xi\eta}\) in \(\alpha\), und von \(-2\pi\sqrt{-1}\) multiplicirt mit dem Werthe von \(S_{\xi\eta}\) in \(\beta\); oder er kann dargestellt werden durch
\[
2\pi\sqrt{-1}\int_{\beta}^{\alpha}dS_{\xi\eta}.
\]

Wir haben also endlich die Gleichung erhalten:
\[
\tag{1}
2\pi\sqrt{-1}\left[\int_{\alpha}^{\beta}dS_{\xi\eta}-\int_{\xi}^{\eta}dS_{\alpha\beta}\right]=\sum c_{ik}B_iA_k.
\]

\noindent\textbf{I.} Bildet man also zwei Integrale dritter Gattung, bei denen die Grenzen des einen die Parameter des andern sind, und umgekehrt, so ist ihre Differenz durch die Periodicitätsmoduln beider Integrale darstellbar.

\subsubsection*{\S\,34. Endliche Normalform der Integrale dritter Gattung. Einfachste Form des Vertauschungssatzes. Sätze, welche daraus folgen.}

Dieser wichtige Satz lässt sich präciser aussprechen, wenn wir die Normalform unserer Integrale dritter Gattung auf passende Weise noch genauer feststellen. Wir haben schon im ersten Abschnitte gesehen, dass die Integrale dritter Gattung, wie wir sie aufstellten, nur bis auf additive Integrale erster Gattung bestimmt waren, d. h. bis auf lineare Ausdrücke mit \(p\) Constanten. Statt eines Integrals \(S_{\xi\eta}\) können wir also, ohne irgend etwas zu ändern, den Ausdruck
\[
S_{\xi\eta}+c_1u_1+c_2u_2+\cdots+c_pu_p
\]
setzen, wobei wir absichtlich die lineare Function durch unsere Normalintegrale erster Gattung ausgedrückt haben. Bilden wir nun die erste Hälfte derjenigen Combinationen der Periodicitätsmoduln, welche wir oben für Integrale erster Gattung als normale bezeichnet haben, und welche wir hier für Integrale dritter Gattung ebenso bezeichnen wollen, so erhalten wir die Werthe:
\[
S_{\xi\eta}^{(1)}+2\pi\sqrt{-1}\,c_1,\quad
S_{\xi\eta}^{(2)}+2\pi\sqrt{-1}\,c_2,\quad\ldots,\quad
S_{\xi\eta}^{(p)}+2\pi\sqrt{-1}\,c_p,
\]
wo \(S_{\xi\eta}^{(1)}\), \(S_{\xi\eta}^{(2)}\) \(\ldots\) die von dem irgendwie fixirten Integral \(S_{\xi\eta}\) herrührenden Theile bezeichnen. Nun können wir offenbar die Constanten \(c\) so bestimmen, dass diese Ausdrücke sämmtlich verschwinden. Alsdann ist der an Stelle von \(S_{\xi\eta}\) zu setzende Ausdruck
\[
\tag{2}
\Pi_{\xi\eta}=S_{\xi\eta}-\frac{1}{2\pi\sqrt{-1}}\left\{S_{\xi\eta}^{(1)}u_1+S_{\xi\eta}^{(2)}u_2+\cdots+S_{\xi\eta}^{(p)}u_p\right\}
\]
vollkommen bestimmt; und wir werden fortan diesen Ausdruck ausschliesslich als Normalintegral dritter Gattung bezeichnen. Das Normalintegral dritter Gattung ist also so bestimmt, dass die erste Hälfte seiner normalen Periodicitätsmoduln verschwindet.

Aber im vorigen Capitel wurde gezeigt, dass der Ausdruck \(\sum c_{ik}a_ib_k\) so transformirt werden kann, dass jeder Term das Product eines Periodicitätsmoduls der einen Hälfte mit einem der andern Hälfte ist. Wenden wir also dieselbe Transformation auf die rechte Seite der Gleichung (1) an, so muss dieselbe verschwinden, da wir die Integrale dritter Gattung jetzt so normirt haben, dass die ganze eine Hälfte ihrer Periodicitätsmoduln verschwindet. Die Gleichung (1) geht also über in
\[
\tag{3}
\int_{\xi}^{\eta}d\Pi_{\alpha\beta}=\int_{\alpha}^{\beta}d\Pi_{\xi\eta},
\]
und man hat folgenden Satz für die Vertauschung von Parameter und Argument:

\noindent\textbf{II.} Das Normalintegral dritter Gattung \(\Pi\) ändert sich nicht, wenn man die Grenzen mit den Unstetigkeitspuncten vertauscht.

Dieser Satz ist höchst zahlreicher Anwendungen fähig. Erinnern wir uns z. B. des im ersten Abschnitte aufgestellten Satzes, dass die Summe dreier Integrale dritter Gattung, deren zwei immer einen Unstetigkeitspunct gemein haben, ein Integral erster Gattung ist. Wir können diese Summe so schreiben:
\[
\int_{\alpha}^{\beta}\left(d\Pi_{\xi\eta}+d\Pi_{\eta\zeta}+d\Pi_{\zeta\xi}\right),
\]
und sie mit Hülfe unseres Satzes in
\[
\int_{\xi}^{\eta}d\Pi_{\alpha\beta}+\int_{\eta}^{\zeta}d\Pi_{\alpha\beta}+\int_{\zeta}^{\xi}d\Pi_{\alpha\beta}
\]
verwandeln, was identisch Null ist, wenn die aus den Wegen \(\xi\eta\), \(\eta\zeta\), \(\zeta\xi\) zusammengesetzte geschlossene Curve sich auf einen Punct zusammenziehen lässt. Vorausgesetzt, dass also die vier Integrale mit den Unstetigkeitspuncten \(\xi\eta\), \(\eta\zeta\), \(\zeta\xi\), \(\alpha\beta\) die in der geometrischen Darstellung oben niedergelegten Bedingungen erfüllen, dass die zu ihrer Definition angewandten, um die Unstetigkeitspuncte geführten Schleifen einander und die Schleifen der Verzweigungspuncte nicht schneiden, hat man den Satz:

\noindent\textbf{III.} Die Summe dreier Normalintegrale dritter Gattung, von denen je zwei einen Unstetigkeitspunct gemein haben, ist identisch Null.

Nehmen wir für \(\zeta\) irgend einen numerischen Werth, so ergiebt sich weiter der Satz:

\noindent\textbf{IV.} Das Normalintegral dritter Gattung ist die Differenz zweier gleichartigen Functionen, deren jede nur von einem seiner Parameter abhängt.

Bezeichnet man solche Functionen durch \(\log \Phi_\xi\), \(\log \Phi_\eta\), wobei \(\Phi\) eine keineswegs vollkommen bestimmte Function ist, so hat man also
\[
\tag{4}
\Pi_{\xi\eta}=\log\frac{\Phi_\eta}{\Phi_\xi}.
\]

Das Theorem III. dient auch dazu, das Integral erster Gattung zu bestimmen, welches der Summe der betreffenden \(S\) gleich wird. Denn wenn man für \(\Pi\) aus (2) seinen Werth in \(S\) setzt, so erhält man:
\[
\tag{5}
\int_{\alpha}^{\beta}\left(dS_{\xi\eta}+dS_{\eta\zeta}+dS_{\zeta\xi}\right)
=\frac{1}{2\pi\sqrt{-1}}\sum_{h=1}^{h=p}
\left(S_{\xi\eta}^{(h)}+S_{\eta\zeta}^{(h)}+S_{\zeta\xi}^{(h)}\right)
\int_{\alpha}^{\beta}du_h.
\]

Es wurde oben bestimmt, dass die erste Hälfte der Periodicitätsmoduln für das Integral dritter Gattung \(\Pi_{\xi\eta}\) verschwinden solle. Es bleibt übrig, für dieses Integral die andre Hälfte der Periodicitätsmoduln zu bestimmen. Wir erhalten dieselben, indem wir in der Gleichung (1) das eine Integral \(S_{\alpha\beta}\) in ein Integral erster Gattung \(u_h\) übergehn lassen. Dann verschwindet \(\int_{\alpha}^{\beta}dS_{\xi\eta}\), und wenn wir \(S\) in \(\Pi\) verwandeln, verschwindet rechts die Hälfte der Glieder, weil sie in die verschwindenden Periodicitätsmoduln von \(\Pi_{\xi\eta}\) multiplicirt werden. Die entsprechenden Periodicitätsmoduln von \(u\) aber verschwinden nicht sämmtlich, sondern es bleibt der \(h^{\mathrm{te}}\) bestehn, welcher gleich \(2\pi\sqrt{-1}\) wird. Dieser ist nun in dem transformirten Ausdruck der rechten Seite in den negativen Periodicitätsmodul von \(S_{\xi\eta}\) multiplicirt, welcher in der zweiten Hälfte der Periodicitätsmoduln mit dem Index \(h\) behaftet ist. Bezeichnen wir ihn durch \(\Pi_{\xi\eta}^{(h)}\), so haben wir also die Gleichung
\[
\tag{6}
\Pi_{\xi\eta}^{(h)}=-\int_{\xi}^{\eta}du_h.
\]

\noindent\textbf{V.} Die zweite Hälfte der Periodicitätsmoduln des Integrals \(\Pi\) drückt sich also durch die zwischen den Unstetigkeitspuncten als Grenzen genommenen Integrale erster Gattung aus.

Setzt man in (5) für \(\Pi\) seinen Ausdruck in \(S\), so erhält man die zweite Reihe \(\Sigma^{(h)}\) der Periodicitätsmoduln von \(S\) ausgedrückt durch die Formel:
\[
\tag{7}
\Sigma_{\xi\eta}^{(h)}=-\int_{\xi}^{\eta}du_h+\frac{1}{2\pi\sqrt{-1}}\sum_{k=1}^{k=p}S_{\xi\eta}^{(k)}a_{kh}.
\]

\subsubsection*{\S\,35. Sätze über Integrale zweiter Gattung, welche aus dem Vertauschungssatze folgen.}

Wir gelangen zu einer weitern Reihe von Sätzen, wenn wir im Vorigen das eine oder beide Integrale dritter Gattung in Integrale zweiter Gattung übergehn lassen, oder was dasselbe ist, wenn wir nach einem Parameter differenziren.

Wir haben oben (p. 32)
\[
E_\xi=\frac{\partial S_{\xi\zeta}}{\partial \xi}
\]
gesetzt. Dieses Integral \(E_\xi\) ist von dem Parameter \(\xi\) nicht unabhängig, aber es enthält \(\xi\) nur in einem Theil, welcher ein Aggregat von Integralen erster Gattung ist. In der That, setzen wir für \(S\) seinen Ausdruck durch \(\Pi\) und durch die ersten \(p\) Periodicitätsmoduln von \(S\), so haben wir
\[
\tag{8}
E_\xi=\frac{\partial \Pi_{\xi\zeta}}{\partial \xi}
+\frac{1}{2\pi\sqrt{-1}}\sum_{h=1}^{h=p}
\frac{\partial S_{\xi\zeta}^{(h)}}{\partial \xi}\,u_h.
\]

Betrachten wir nun zunächst den nur von dem einen Parameter \(\xi\) abhängigen Ausdruck
\[
\tag{9}
Z_\xi=\frac{\partial \Pi_{\xi\zeta}}{\partial \xi}=\frac{\partial \log\Phi_\xi}{\partial \xi},
\]
und schliessen denselben zwischen die Grenzen \(\alpha,\beta\) ein, so finden wir nach Theorem II.:
\[
\tag{10}
\int_{\alpha}^{\beta}dZ_\xi
=\frac{\partial}{\partial \xi}\int_{\xi}^{\zeta}d\Pi_{\alpha\beta}
=-\left(\frac{\partial\Pi_{\alpha\beta}}{\partial x}\right)_{x=\xi}
=-\left(\frac{\partial\Pi_{\alpha\beta}}{\partial x}\right)_\xi.
\]
Die rechte Seite aber ist eine algebraische Function der dem Puncte \(\xi\) zugehörigen Werthe von \(x,s_x\). Man hat also den Satz:

\noindent\textbf{VI.} Das zwischen bestimmten Grenzen genommene, aus Differentiation von \(\Pi\) nach einem Parameter entstandene Integral zweiter Gattung \(Z\) ist eine algebraische Function des Parameters, nach welchem differenzirt wird.

Setzen wir für \(\Pi\) seinen Werth in \(S\), so erhalten wir die Formel
\[
\tag{11}
\int_{\alpha}^{\beta}dZ_\xi
=-\left(\frac{\partial S_{\alpha\beta}}{\partial x}\right)_\xi
+\frac{1}{2\pi\sqrt{-1}}\sum_{h=1}^{h=p}S_{\alpha\beta}^{(h)}
\left(\frac{\partial u_h}{\partial x}\right)_\xi,
\]
und daher den Satz:

\noindent\textbf{VII.} Das zwischen bestimmten Grenzen genommene Integral zweiter Gattung \(Z\) drückt sich durch eine algebraische Function, und durch eine linear auftretende Reihe von transcendenten Functionen seiner Grenzen aus, welche die erste Reihe von Periodicitätsmoduln desjenigen Integrals dritter Gattung sind, das jene Grenzen zu Unstetigkeitspuncten hat.

Endlich erhält man für \(E\) aus (8) die Formel:
\[
\tag{12}
\begin{aligned}
\int_{\alpha}^{\beta}dE_\xi
&=-\left(\frac{\partial S_{\alpha\beta}}{\partial x}\right)_\xi\\
&\quad+\frac{1}{2\pi\sqrt{-1}}\sum_{h=1}^{h=p}\left[
S_{\alpha\beta}^{(h)}\left(\frac{\partial u_h}{\partial x}\right)_\xi
+\frac{\partial S_{\xi\zeta}^{(h)}}{\partial \xi}\int_{\alpha}^{\beta}du_h\right],
\end{aligned}
\]
woraus sich der Satz ergiebt:

\noindent\textbf{VIII.} Das zwischen bestimmten Grenzen genommene Integral zweiter Gattung \(E\) drückt sich durch algebraische Functionen und durch Integrale erster Gattung aus, sowie ausserdem linear

1) durch die erste Reihe der Periodicitätsmoduln desjenigen Integrals dritter Gattung \(S\), dessen Unstetigkeitspuncte die Grenzen sind,

2) durch die Periodicitätsmoduln erster Reihe
\[
\tag{13}
E_\xi^{(h)}=\frac{\partial S_{\xi\zeta}^{(h)}}{\partial \xi}
\]
des Integrals \(E\) selbst.

Aus den Gleichungen (8), (9), (13) finden wir ferner die Combination:
\[
\tag{14}
Z_\xi=E_\xi-\frac{1}{2\pi\sqrt{-1}}\sum_{h=1}^{h=p}E_\xi^{(h)}u_h,
\]
welche der Gleichung (2) zwischen \(\Pi\) und \(S\) ganz analog ist. Hieraus folgt sofort:

\noindent\textbf{IX.} Die erste Reihe der Periodicitätsmoduln von \(Z_\xi\) verschwindet.

Indem man sodann die Gleichung (6) nach \(\xi\) differenzirt, erhält man die zweite Reihe von Periodicitätsmoduln von \(Z\):
\[
\tag{15}
Z_\xi^{(h)}=\left(\frac{\partial u_h}{\partial x}\right)_\xi.
\]
Also:

\noindent\textbf{X.} Die zweite Reihe von Periodicitätsmoduln der Function \(Z\) besteht aus algebraischen Functionen des Parameters.

Und die zweite Reihe der Periodicitätsmoduln von \(E_\xi\) setzt sich, wie sich aus (14) oder durch Differentiation von (7) ergiebt, aus der ersten Reihe und aus den Periodicitätsmoduln von \(S_{\xi\zeta}^{(h)}\) zusammen mit Hülfe der Formel
\[
\tag{16}
H_\xi^{(h)}=\left(\frac{\partial u_h}{\partial x}\right)_\xi
+\frac{1}{2\pi\sqrt{-1}}\sum_{k=1}^{k=p}E_\xi^{(k)}a_{kh}.
\]

Aus den obigen Sätzen erkennt man, dass die Integrale zweiter Gattung \(E_\xi\) sich immer ausdrücken lassen durch algebraische Functionen und Integrale erster Gattung, durch deren Periodicitätsmoduln, welche bei der Bildung der Normalintegrale erster Gattung benutzt wurden; durch die erste Reihe von Periodicitätsmoduln der Integrale dritter Gattung, und endlich durch deren nach einem Parameter genommenen Differentialquotienten, welche die erste Reihe der Periodicitätsmoduln des Integrals zweiter Gattung bilden.

Für die letzteren Grössen kann man nun noch Relationen aufstellen, welche den Satz nachweisen:

\noindent\textbf{XI.} Die erste Reihe der Periodicitätsmoduln der Integrale \(E_\xi\) mit beliebigem Unstetigkeitspuncte \(\xi\) lässt sich immer durch die entsprechenden Periodicitätsmoduln solcher \(p\) Integrale \(E_\alpha,E_\beta\ldots E_\pi\) ausdrücken, welche in bestimmten, beliebig gewählten Puncten \(\alpha,\beta\ldots\pi\) unstetig werden.

Zum Beweise dieses Satzes führt eine weitere Eigenschaft der Function \(Z\). Man hat nämlich, wenn der beigefügte Index anzeigt, dass in der Function für \(x,s_x\) die dem durch den Index angedeuteten Punct entsprechende Werthe gesetzt werden sollen:
\[
\left(\frac{\partial Z_\xi}{\partial x}\right)_\alpha
=-\frac{\partial}{\partial\alpha}\int_{\alpha}^{\beta}dZ_\xi
=-\frac{\partial}{\partial\alpha}\int_{\alpha}^{\beta}d\frac{\partial\Pi_{\xi\zeta}}{\partial\xi}
=-\frac{\partial^2}{\partial\alpha\partial\xi}\int_{\alpha}^{\beta}d\Pi_{\xi\zeta}.
\]
Da nun dieser Ausdruck sich nicht ändert, wenn man \(\xi\) mit \(\alpha\), \(\zeta\) mit \(\beta\) vertauscht, so hat man den Satz:

\noindent\textbf{XII.} Wenn man \(Z_\xi\) nach dem Argument differenzirt und dann für \(s_x,x\) die Werthe setzt, welche diese Variabeln für einen beliebigen Punct \(\alpha\) erhalten, so ist das Resultat dasselbe, wie wenn im Differentialquotienten von \(Z_\alpha\) nach dem Argument \(s_x,x\) die dem Puncte \(\xi\) entsprechenden Werthe gesetzt werden.

Diesen Satz kann man auch in folgender Form aussprechen:

\noindent\textbf{XIII.} Hat man \(q\) Puncte und bildet die Summen \(M_1,M_2,\ldots,M_q\) so, dass jedes \(M\) die Summe von \(q-1\) Functionen \(Z\) ist, deren Argumente \(q-1\) der obigen Puncte entsprechen, während der Parameter dem \(q^{\mathrm{ten}}\) Puncte entspricht; und sind \(x_1,x_2\ldots x_q\) die Werthe von \(x\) für diese \(q\) Puncte, so ist immer
\[
M_1dx_1+M_2dx_2+\cdots+M_qdx_q
\]
ein vollständiges Differential.

Die Gleichung
\[
\tag{17}
\left(\frac{\partial Z_\xi}{\partial x}\right)_\alpha=
\left(\frac{\partial Z_\alpha}{\partial x}\right)_\xi,
\]
giebt aber, wenn man die Functionen \(Z\) durch die \(E\) ausdrückt:
\[
\tag{18}
\begin{aligned}
&\left(\frac{\partial E_\xi}{\partial x}\right)_\alpha
-\frac{1}{2\pi\sqrt{-1}}\sum_{h=1}^{h=p}E_\xi^{(h)}\left(\frac{\partial u_h}{\partial x}\right)_\alpha\\
&\qquad=\left(\frac{\partial E_\alpha}{\partial x}\right)_\xi
-\frac{1}{2\pi\sqrt{-1}}\sum_{h=1}^{h=p}E_\alpha^{(h)}\left(\frac{\partial u_h}{\partial x}\right)_\xi.
\end{aligned}
\]

Bildet man diese Gleichung nun, indem man statt \(\alpha\) der Reihe nach \(p\) beliebig gewählte constante Puncte \(\alpha,\beta,\ldots\pi\) benutzt, und bemerkt man, dass die eingeklammerten Grössen rein algebraische Ausdrücke sind, so hat man ein System linearer Gleichungen vor sich, vermöge dessen man die \(E_\xi^{(h)}\) durch die \(E_\alpha^{(h)}\), \(E_\beta^{(h)}\ldots E_\pi^{(h)}\) und durch algebraische Functionen von \(\xi\) ausdrücken kann, wie das Theorem XI es verlangt.

Bildet man ebenso aus den Gleichungen (12), indem man für \(\xi\) darin der Reihe nach \(a,b,c\ldots\) einsetzt, \(p\) verschiedene Gleichungen, so findet man die ganzen Integrale \(S_{\alpha\beta}^{(h)}\) vermittelst eines Systemes linearer Gleichungen durch rein algebraische Integrale \(2^{\mathrm{ter}}\) Gattung, durch algebraische Functionen und die Grössen \(E_a^{(h)}\), \(E_b^{(h)}\) etc. ausgedrückt. Man hat also noch den Satz:

\noindent\textbf{XIV.} Die ganzen Integrale \(S_{\alpha\beta}^{(h)}\) lassen sich durch rein algebraische Integrale erster und zweiter Gattung und durch algebraische Functionen ausdrücken, wobei nur ganze Integrale zweiter Gattung mit absolut festen Unstetigkeitspuncten in den Coefficienten erscheinen.

\subsubsection*{\S\,36. Zurückführung der Differentialquotienten von Integralen dritter Gattung nach einem Parameter auf Integrale zweiter und dritter Gattung.}

Im ersten Abschnitte wurde bewiesen, dass jedes Integral einer rationalen Function von \(x\) und \(s_x\) sich ausdrücken lässt durch algebraische und logarithmische Functionen, durch Integrale erster und dritter Gattung und durch die nach den Unstetigkeitspuncten genommenen Differentialquotienten der letzteren. Wir wollen jetzt zeigen, wie auch alle höheren Differentialquotienten der Integrale dritter Gattung sich zurückführen lassen auf algebraische Functionen, auf Integrale erster Gattung und deren Periodicitätsmoduln, auf die Integrale dritter Gattung und deren Periodicitätsmoduln, und auf die erste Reihe der Periodicitätsmoduln von \(p\) Integralen zweiter Gattung, deren Unstetigkeitspuncte beliebig gewählt sind.

Differenziren wir nämlich die Gleichung (1) \(\mu\) mal nach \(\xi\), \(\nu\) mal nach \(\eta\). Man erhält dann
\[
\frac{\partial^{\mu+\nu}}{\partial\xi^\mu\partial\eta^\nu}\int_{\alpha}^{\beta}d\Pi_{\xi\eta}
=
\frac{\partial^{\mu+\nu}}{\partial\xi^\mu\partial\eta^\nu}\int_{\xi}^{\eta}d\Pi_{\alpha\beta}.
\]
Drückt man nun die \(\Pi\) durch die Integrale \(S\) aus, so wird daraus:
\[
\tag{19}
\begin{aligned}
&\frac{\partial^{\mu+\nu}}{\partial\xi^\mu\partial\eta^\nu}\int_{\alpha}^{\beta}dS_{\xi\eta}
-\frac{1}{2\pi\sqrt{-1}}\sum_{h=1}^{h=p}
\frac{\partial^{\mu+\nu}S_{\xi\eta}^{(h)}}{\partial\xi^\mu\partial\eta^\nu}\int_{\alpha}^{\beta}du_h\\
&\quad=\frac{\partial^{\mu+\nu}}{\partial\xi^\mu\partial\eta^\nu}\int_{\xi}^{\eta}dS_{\alpha\beta}
-\frac{1}{2\pi\sqrt{-1}}\sum_{h=1}^{h=p}S_{\alpha\beta}^{(h)}
\frac{\partial^{\mu+\nu}}{\partial\xi^\mu\partial\eta^\nu}\int_{\xi}^{\eta}du_h.
\end{aligned}
\]
Diese Gleichung drückt in der That schon das verlangte Integral auf die geforderte Weise aus, nur dass in dieser Formel noch die Grössen
\[
\frac{\partial^{\mu+\nu}S_{\xi\eta}^{(h)}}{\partial\xi^\mu\partial\eta^\nu},
\]
die höheren Differentialquotienten der ersten Reihe von Periodicitätsmoduln des Integrals \(S_{\xi\eta}\) auftreten. Wenn wir aber jetzt in (18) \(\xi,\eta\) an Stelle von \(\alpha,\beta\), und für \(\xi\) der Reihe nach \(p\) beliebig gewählte Puncte setzen, und wenn wir sodann \(\mu\) mal nach \(\xi\) und \(\nu\) mal nach \(\eta\) differenziren, so haben wir \(p\) Gleichungen vor uns, welche für die gesuchten Grössen
\[
\frac{\partial^{\mu+\nu}S_{\xi\eta}^{(h)}}{\partial\xi^\mu\partial\eta^\nu}
\]
linear sind, und welche ausserdem nur noch algebraische Ausdrücke und die erste Reihe der jenen \(p\) Puncten entsprechenden Periodicitätsmoduln von Integralen zweiter Gattung enthalten. Durch diese also und algebraische Ausdrücke können wir die gesuchten Grössen darstellen, und indem wir dies in (19) eintragen, erhalten wir die geforderte Formel.

\subsubsection*{\S\,37. Anwendung der oben für Integrale dritter Gattung benutzten Methoden auf algebraische Functionen und Logarithmen von solchen. Das Abel'sche Theorem.}

Wir haben soweit alle Resultate abgeleitet, die sich aus der Untersuchung des Integrals \(\int I\,dH\) ergeben, wenn \(I,H\) Integrale erster, zweiter oder dritter Gattung sind. Wir wollen jetzt noch die Resultate untersuchen, welche sich ergeben, wenn wir für eine der Functionen \(I,H\) eine algebraische Function oder den Logarithmus einer solchen einführen.

Um das Integral \(\int \log\frac{\varphi}{\psi}\,dS_{\alpha\beta}\) zu behandeln, in welchem \(\varphi,\psi\) beliebige ganze Functionen \(m^{\mathrm{ter}}\) Ordnung von \(s_x,x\) sind, construiren wir uns zunächst die Schleife für \(S_{\alpha\beta}\) wie im Früheren. Die Functionen \(\varphi,\psi\) verschwinden jede in \(mn\) Puncten, den Schnittpuncten der Curve \(\varphi=0\), \(f=0\) einerseits, \(\psi=0\), \(f=0\) andrerseits entsprechend. Fallen einige der ersten mit einigen der letzten zusammen, so bleibt \(\log\frac{\varphi}{\psi}\) für sie endlich, während für die übrigen \(\log\frac{\varphi}{\psi}\) unendlich gross wird. Wir wollen nun immer einen Punct \(s_x,x\), für welchen \(\varphi\) verschwindet, mit einem Punct \(s_x,x\) verbinden, für welchen \(\psi\) verschwindet; aus diesen Curven construiren wir Schleifen, welche um je zwei solche Puncte und zum Puncte \(0\) führen, wie die früher den Integralen dritter Gattung entsprechenden gezogen wurden, und legen sie so, dass niemals zwei verschiedene Schleifen ein Werthepaar \(s_x,x\) gemein haben. Wie man die Puncte, welche von \(\varphi=0\) herrühren, mit den von \(\psi=0\) herrührenden verbindet, ist bei den ersten Punctpaaren gleichgültig; die letzten aber combiniren wir so, dass, wenn die Curve \(\varphi=0\) durch stetige Veränderung in \(\psi=0\) übergeführt wird, und wenn dies so geschieht, dass von den zuerst combinirten Puncten dabei immer der eine in den zugehörigen übergeht, auch von den letzteren Paaren immer ein Punct in den andern übergeführt wird.

Die Function \(\log\frac{\varphi}{\psi}\) ändert sich beim Umgange um eine solche Schleife nicht; denn sie ändert sich beim Umgange um den einen Unstetigkeitspunct immer um \(-2\pi\sqrt{-1}\), beim Umgange um den andern um \(+2\pi\sqrt{-1}\), bleibt also bei gleichzeitigem Umgange um beide ungeändert. Bilden wir nun \(\int d\log\frac{\varphi}{\psi}\) um einen geschlossenen Umgang, so haben auf den Werth dieses Integrals die eingeschalteten Schleifen keinen Einfluss, und man erhält zwischen den Periodicitätsmoduln von \(\log\frac{\varphi}{\psi}\) dieselben Relationen, wie sie für die der Integrale erster Gattung stattfinden. Uebrigens haben diese Periodicitätsmoduln sämmtlich den Werth \(2m\pi\sqrt{-1}\), wo die \(m\) irgend welche ganze Zahlen sein können; denn dies ist der einzige Zuwachs, den \(\log\frac{\varphi}{\psi}\) beim Umgange um einen cyclischen Weg überhaupt erhalten kann.

Das Integral \(\int\log\frac{\varphi}{\psi}\,dS_{\alpha\beta}\) um alle Umgänge genommen besteht nun ganz wie früher aus einem Theil, welcher von den um die Verzweigungspuncte geführten Integralen herrührt, und welcher sich aus den Periodicitätsmoduln von \(S_{\alpha\beta}\) und \(\log\frac{\varphi}{\psi}\) zusammensetzt; aus einem zweiten Theile, der von der um \(\alpha,\beta\) gelegten Schleife herrührt, und auf die oben auseinandergesetzte Weise den Werth
\[
2\pi\sqrt{-1}\int_{\beta}^{\alpha}d\log\frac{\varphi}{\psi}
\]
giebt; und endlich aus den Gliedern, welche von den Unstetigkeitspuncten der Function \(\log\frac{\varphi}{\psi}\) herrühren. Diese Glieder kann man genau so behandeln, wie oben diejenigen, welche von \(S_{\xi\eta}\) herrührten, und findet
\[
-2\pi\sqrt{-1}\sum_{\xi}\int_{\xi}^{\eta}dS_{\alpha\beta},
\]
wo die \(\xi\) sich auf die Puncte beziehen, in denen \(\varphi=0\), die \(\eta\) aber auf diejenigen, in welchen \(\psi=0\).

Dividirt man nun durch \(2\pi\sqrt{-1}\), so erhält man, da die Periodicitätsmoduln von \(\log\frac{\varphi}{\psi}\) sich dadurch auf ganze Zahlen reduciren:
\[
\sum_{\xi}\int_{\xi}^{\eta}dS_{\alpha\beta}-\int_{\beta}^{\alpha}d\log\frac{\varphi}{\psi}=K,
\]
wo \(K\) eine ganzzahlige Combination der Perioden von \(S_{\alpha\beta}\) ist, die von der Lage der um die Puncte \(\xi,\eta\) geführten Schleifen gegen die Verzweigungsschleifen zwar abhängen kann, nicht aber von der speciellen Lage der Puncte \(\xi,\eta\). Daher kann man zur Bestimmung von \(K\) die Curven \(\varphi,\psi\) einander unendlich nahe rücken lassen, und da der Voraussetzung nach dann auch die paarweise einander zugeordneten Puncte \(\xi,\eta\) zusammenfallen sollen, so hat man \(K=0\), und daher
\[
\log\left(\frac{\varphi}{\psi}\right)_\alpha-\log\left(\frac{\varphi}{\psi}\right)_\beta
=\sum_{\xi}\int_{\xi}^{\eta}dS_{\alpha\beta}.
\]
Diese Gleichung ist das Abel'sche Theorem, und zwar in einer Weise, bei welcher sowohl die Integrale rechts als die Logarithmen vollkommen bestimmt sind. Man kann darin der Ableitung der Formeln nach \(\Pi\) an Stelle von \(S\) setzen; und dies stimmt mit dem Satze überein, dass \(\sum_{\xi}\int_{\xi}^{\eta}du_h\) verschwinden muss, welcher einen Theil des Abel'schen Theorems bildet. Bei dieser Darstellung des Abel'schen Theorems ist vorausgesetzt, dass erstlich der Weg, auf dem man von \(\log(\frac{\varphi}{\psi})_\beta\) zu \(\log(\frac{\varphi}{\psi})_\alpha\) gelangt, von keinem der Integrationswege rechts geschnitten wird, und dass zweitens die obern und untern Grenzen dieser Wege einander so zugeordnet sind, dass sie unendlich klein werden, wenn man die Curven \(\varphi=0\), \(\psi=0\) unendlich nahe zusammenrücken lässt.

Setzt man aber \(\Pi\) an Stelle von \(S\), und wendet man den Satz über Vertauschung von Parameter und Argument an, so findet man:
\[
\mathrm{I.}\quad
\log\left(\frac{\varphi}{\psi}\right)_\alpha-\log\left(\frac{\varphi}{\psi}\right)_\beta
=\sum_{\xi}\int_{\alpha}^{\beta}d\Pi_{\xi\eta}.
\]
Diese Formel enthält die Zurückführung des Logarithmus einer algebraischen Function auf eine Summe von Normalintegralen dritter Gattung. Sie hört nicht auf zu gelten, wenn man einige der Puncte \(\xi\) oder der Puncte \(\eta\) einander sehr nahe rücken lässt, wodurch nur einzelne Theile der um die Unstetigkeitspuncte gelegten Schleifen einander unendlich nahe rücken; oder auch ganze Schleifen, sobald zwei Paare zugeordneter Puncte einander sehr nahe rücken, d. h. wenn entsprechende Puncte der Curven \(\varphi=0\), \(\psi=0\) Berührungspuncte dieser Curven mit \(f=0\) werden. —

Wenn wir nun zur Untersuchung des Integrals \(\int \frac{\varphi}{\psi}\,dS_{\alpha\beta}\) übergehen, so legen wir zunächst ein Schleifensystem, welches wegen der Unstetigkeitspuncte von \(\frac{\varphi}{\psi}\) nothwendig wird, folgendermassen. Jede Schleife geht vom Puncte \(0\) aus nach einem der Unstetigkeitspuncte von \(\frac{\varphi}{\psi}\), um diesen der Bewegung eines Uhrzeigers entsprechend in kleinem Kreise herum und nach \(0\) zurück. Dieselbe muss in \(0\) mit einem bestimmten Werthe von \(s_x\) beginnen, um sich in der Nähe des Unstetigkeitswerthes von \(x\) zugleich dem Unstetigkeitswerthe von \(s_x\) zu nähern, und wird demgemäss einem bestimmten Umgange eingereiht. Wieder dürfen diese Schleifen einander und die Schleifen der Verzweigungspuncte nicht in der oben verbotenen Art schneiden.

Da die Periodicitätsmoduln von \(\frac{\varphi}{\psi}\) sämmtlich Null sind, weil die Function nach Durchlaufung eines cyclischen Weges immer ihren ursprünglichen Werth wiedererlangt, so hat man in dem Integral \(\int \frac{\varphi}{\psi}\,dS_{\alpha\beta}\) nur diejenigen Theile zu untersuchen, welche von den um die Unstetigkeitspuncte gelegten Schleifen herrühren. Die zu \(\alpha,\beta\) gehörige Schleife giebt, ähnlich wie oben
\[
2\pi\sqrt{-1}\int_{\beta}^{\alpha} d\frac{\varphi}{\psi}
=2\pi\sqrt{-1}\left\{\left(\frac{\varphi}{\psi}\right)_{\alpha}-\left(\frac{\varphi}{\psi}\right)_{\beta}\right\}.
\]
Ist ferner \(\xi\) einer der Unstetigkeitspuncte von \(\frac{\varphi}{\psi}\), und zieht man die ihm entsprechende Schleife in einen unendlich kleinen, um den Unstetigkeitspunct beschriebenen Kreis zusammen, so wird der von dieser Schleife herrührende Theil des Integrals:
\[
-2\pi\sqrt{-1}\left(\frac{\varphi}{\frac{\partial\psi}{\partial x}}\right)_{\xi}\cdot
\left(\frac{\partial S_{\alpha\beta}}{\partial x}\right)_{\xi},
\]
so dass man die Gleichung erhält:
\[
\left(\frac{\varphi}{\psi}\right)_{\alpha}-\left(\frac{\varphi}{\psi}\right)_{\beta}
=\sum\left(\frac{\varphi}{\frac{\partial\psi}{\partial x}}\right)_{\xi}
\left(\frac{\partial S_{\alpha\beta}}{\partial x}\right)_{\xi},
\]
oder
\[
\left(\frac{\varphi}{\psi}\right)_{\alpha}-\left(\frac{\varphi}{\psi}\right)_{\beta}
=-\sum\left(\frac{\varphi}{\frac{\partial\psi}{\partial x}}\right)_{\xi}
\frac{\partial}{\partial\xi}\int_{\xi}^{\zeta}dS_{\alpha\beta}.
\]
Hier kann man nun für \(S\) die Function \(\Pi\) setzen, und findet dann mit Anwendung des Satzes über Vertauschung von Parameter und Argument:
\[
\mathrm{II}\ldots\quad
\left(\frac{\varphi}{\psi}\right)_{\alpha}-\left(\frac{\varphi}{\psi}\right)_{\beta}
=-\sum\left(\frac{\varphi}{\frac{\partial\psi}{\partial x}}\right)_{\xi}
\int_{\alpha}^{\beta}dZ_{\xi}.
\]
Diese Formel enthält, wenn man \(\alpha\) als die Variable betrachtet, die Darstellung einer beliebigen rationalen Function von \(\alpha\) und \(s_{\alpha}\) als Aggregat von Integralen zweiter Gattung. Wenn mehrere Puncte \(\xi\) einander unendlich nahe rücken, so treten Differentialquotienten von Integralen \(Z\) nach dem Parameter auf.

\section*{Sechster Abschnitt.}
\subsection*{Das Umkehrproblem.}

\subsubsection*{\S\,38. Hülfssatz zur Beurtheilung periodischer Functionen.}

Die naheliegende Frage, ob man ebenso, wie aus dem Integral
\[u=\int_0^x\frac{dx}{\sqrt{1-x^2}}
\]
die obere Grenze als Function \(x=\sin u\) des Integrals, oder aus \(u=\int\frac{dx}{x}\) die obere Grenze als Function \(x=e^u\) des Integrals erscheint, auch bei einem allgemeinen algebraischen Integral erster Gattung die obere Grenze als Function des Integrals behandeln könne, hat \textsc{Jacobi} in einer berühmten Abhandlung, welcher wir den Begriff der Abel'schen Functionen verdanken, behandelt und den Charakter der dadurch entstehenden Functionen festgelegt.\footnote{Crelle's Journal Bd. XIII.} Obgleich die Untersuchung Jacobi's sich zunächst nur auf den Fall \(p=2\) bezieht, so liegt doch die Verallgemeinerung auf der Hand. Man braucht zunächst einen Hülfssatz, welchen Jacobi für jenen Fall aufstellt, und welcher verallgemeinert sich folgendermassen ausspricht:

Es seien \(q-1\) Reihen von je \(q\) reellen Grössen gegeben:
\[
(1)\ldots\quad
\begin{array}{cccc}
a_1^{(1)},&a_2^{(1)},&\ldots&a_q^{(1)}\\
a_1^{(2)},&a_2^{(2)},&\ldots&a_q^{(2)}\\[-0.15em]
\multicolumn{4}{c}{\cdots\cdots\cdots}\\
a_1^{(q-1)},&a_2^{(q-1)},&\ldots&a_q^{(q-1)}.
\end{array}
\]
Dann kann man immer \(q\) positive oder negative ganze Zahlen \(m_1,m_2\ldots m_q\) so bestimmen, dass die \(q-1\) Ausdrücke
\[
(2)\ldots\quad
\begin{array}{rcl}
Q^{(1)}&=&m_1a_1^{(1)}+m_2a_2^{(1)}+\cdots+m_qa_q^{(1)}\\
Q^{(2)}&=&m_1a_1^{(2)}+m_2a_2^{(2)}+\cdots+m_qa_q^{(2)}\\[-0.15em]
\multicolumn{3}{c}{\cdots\cdots\cdots}\\
Q^{(q-1)}&=&m_1a_1^{(q-1)}+m_2a_2^{(q-1)}+\cdots+m_qa_q^{(q-1)}
\end{array}
\]
beziehungsweise kleiner werden als \(q-1\) beliebig gewählte noch so kleine Grössen \(\varepsilon^{(1)},\varepsilon^{(2)}\ldots\varepsilon^{(q-1)}\).

Nehmen wir an, der Satz wäre bewiesen, wenn man für \(q\) die Zahlen \(1,2\ldots q-1\) setzt. Wir wollen dann zeigen, dass er für \(q\) selbst gilt.

Unter den Grössen \(a_q^{(1)},a_q^{(2)}\ldots a_q^{(q-1)}\) muss wenigstens eine sein, welche von Null verschieden ist. Dieselbe sei etwa \(a_q^{(q-1)}\). Bilden wir nun die Ausdrücke
\[
(3)\ldots\quad
\begin{array}{c}
Q^{(1)}-Q^{(q-1)}\dfrac{a_q^{(1)}}{a_q^{(q-1)}}\\[0.45em]
Q^{(2)}-Q^{(q-1)}\dfrac{a_q^{(2)}}{a_q^{(q-1)}}\\[-0.1em]
\cdots\cdots\cdots\\[0.25em]
Q^{(q-2)}-Q^{(q-1)}\dfrac{a_q^{(q-2)}}{a_q^{(q-1)}}
\end{array}
\]
Die Anzahl dieser Gleichungen ist \(q-2\), die Anzahl der in denselben enthaltenen unbekannten ganzen Zahlen ist \(q-1\), da \(m_q\) eliminirt ist. Wir können daher der Voraussetzung nach die in diesen Ausdrücken enthaltenen Zahlen \(m_1,m_2\ldots m_{q-1}\) so bestimmen, dass diese \(q-2\) Ausdrücke absolut kleiner werden, als beliebig kleine vorherbestimmte Grössen. Sodann aber kann man \(m_q\) so bestimmen, dass \(\frac{Q^{(q-1)}}{a_q^{(q-1)}}\) absolut kleiner wird als \(\frac12\). Die absoluten Werthe der Ausdrücke (3) waren kleiner als beliebig kleine gegebene Zahlen, und daher müssen auch die Ausdrücke \(Q^{(1)},Q^{(2)}\ldots Q^{(q-1)}\) selbst sich den Werthen \(\frac12a_q^{(1)},\frac12a_q^{(2)}\ldots\frac12a_q^{(q-2)}\) beliebig nähern. Ist insbesondere der absolute Werth von \(\frac{Q^{(q-1)}}{a_q^{(q-1)}}\) gleich \(\frac12-\varepsilon\), und ist abgesehen vom Zeichen
\[
Q^{(i)}-Q^{(q-1)}\frac{a_q^{(i)}}{a_q^{(q-1)}}=\varepsilon^{(i)},
\]
so ist der absolute Werth von \(Q^{(i)}\) nicht grösser als der von \(\varepsilon_i\) vermehrt um den von \(a_q^{(i)}\left(\frac12-\varepsilon\right)\), oder nicht grösser als \(\varepsilon^{(i)}+\frac12\sqrt{\bigl(a_q^{(i)}\bigr)^2}\). Setzt man also \(\varepsilon^{(i)}=\frac{\sqrt{\bigl(a_q^{(i)}\bigr)^2}}{m}\), wo \(m\) eine beliebig grosse Zahl ist, so hat man schliesslich \(Q^{(i)}\) gleich oder kleiner als der absolute Werth von \(\left(\frac12+\frac1m\right)a_q^{(i)}\); und man kann daher für alle Werthe von \(i\) schreiben (abgesehen vom Zeichen)
\[
Q^{(i)}<\left(\frac12+\frac1m\right)a_q^{(i)},
\]
wobei nur, wenn eine der Grössen \(a_q^{(1)},a_q^{(2)}\ldots a_q^{(q-2)}\) verschwindend klein wird, das betreffende \(Q^{(i)}\) kleiner als eine beliebig vorausbestimmte Grenze zu setzen ist. Bezeichnen wir die auf diese Weise gebildeten ganzzahligen Verbindungen \(Q^{(1)},Q^{(2)}\ldots\) durch \(a_{q+1}^{(1)},a_{q+1}^{(2)}\ldots a_{q+1}^{(q-1)}\), so dass also abgesehen vom Zeichen
\[
a_{q+1}^{(i)}<\frac12a_q^{(i)}\qquad(i=1,2\ldots q-1).
\]
Wenn nun nur eine der Grössen \(a_{q+1}^{(i)}\) nicht verschwindend klein wird, so kann man die Ausdrücke
\[
\begin{array}{rcl}
R^{(1)}&=&n_2a_2^{(1)}+n_3a_3^{(1)}+\cdots+n_{q+1}a_{q+1}^{(1)}\\
R^{(2)}&=&n_2a_2^{(2)}+n_3a_3^{(2)}+\cdots+n_{q+1}a_{q+1}^{(2)}\\[-0.15em]
\multicolumn{3}{c}{\cdots\cdots\cdots}\\
R^{(q-1)}&=&n_2a_2^{(q-1)}+n_3a_3^{(q-1)}+\cdots+n_{q+1}a_{q+1}^{(q-1)}
\end{array}
\]
ganz ebenso behandeln, wie oben die Ausdrücke \(Q\); d. h. man kann die ganzen Zahlen \(n\) so bestimmen, dass
\[
R^{(i)}<\frac12a_{q+1}^{(i)}\qquad(i=1,2\ldots q-1),
\]
oder, wenn man diese linearen ganzzahligen Verbindungen durch \(a_{q+2}^{(i)}\) bezeichnet, dass
\[
a_{q+2}^{(i)}<\frac12a_{q+1}^{(i)}\qquad(i=1,2\ldots q-1),
\]
wobei nur, sobald eine der Grössen \(a_{q+1}^{(i)}\) verschwindend klein ist, \(a_{q+2}^{(i)}\) kleiner als eine beliebige Grenze zu setzen ist.

Indem man auf diese Weise fortfährt, bildet man sich Reihen von Grössen,
\[
\begin{array}{cccc}
a_{q+1}^{(1)},&a_{q+2}^{(1)},&\ldots&a_{q+k}^{(1)}\\
a_{q+1}^{(2)},&a_{q+2}^{(2)},&\ldots&a_{q+k}^{(2)}\\[-0.15em]
\multicolumn{4}{c}{\cdots\cdots\cdots}\\
a_{q+1}^{(q-1)},&a_{q+2}^{(q-1)},&\ldots&a_{q+k}^{(q-1)}
\end{array}
\]
deren jede entweder verschwindend klein oder kleiner als die Hälfte der unmittelbar links neben ihr stehenden Grösse ist. Man kann diesen Process also so lange fortsetzen, bis alle \(a\) der letzten Verticalreihe kleiner sind als beliebig vorherbestimmte Grössen.

Nun ist aber jedes \(a\) eine lineare ganzzahlige Function der \(q\) Grössen \(a\), welche links neben ihm sich befinden (wobei man sich das obige Schema durch das System der gegebenen \(a\) nach links hin zu ergänzen hat), und die Coefficienten sind bei sämmtlichen \(a\) der gleichen Verticalreihe dieselben. Daher kann man auch die \(a\) der letzten Verticalreihe durch die gegebenen \(a\) linear und ganzzahlig ausdrücken, so dass
\[
\begin{array}{rcl}
a_{q+k}^{(1)}&=&\mu_1a_1^{(1)}+\mu_2a_2^{(1)}+\cdots+\mu_qa_q^{(1)}\\
a_{q+k}^{(2)}&=&\mu_1a_1^{(2)}+\mu_2a_2^{(2)}+\cdots+\mu_qa_q^{(2)}\\[-0.15em]
\multicolumn{3}{c}{\cdots\cdots\cdots}\\
a_{q+k}^{(q-1)}&=&\mu_1a_1^{(q-1)}+\mu_2a_2^{(q-1)}+\cdots+\mu_qa_q^{(q-1)}.
\end{array}
\]
Und diese Ausdrücke liefern also, wie im Theorem verlangt wurde, ganzzahlige lineare Verbindungen der gegebenen \(a\), welche beliebig klein gemacht worden sind.

Der Algorithmus bricht ab in dem oben ausgenommenen Falle, dass alle Grössen \(a\) einer Verticalreihe verschwindend klein werden. Man hat dann Verbindungen der verlangten Natur vor sich, und braucht die Operationen also nicht weiter fortzusetzen.

Wir haben sonach bewiesen, dass unser Satz für eine beliebige Zahl \(q\) richtig ist, wenn er für \(q-1\) besteht. Er ist also nur noch für \(q=2\) nachzuweisen, wo er so lautet: Sind \(a_1,a_2\) reelle Grössen, so kann man immer zwei Zahlen \(m_1,m_2\) so bestimmen, dass \(m_1a_1+m_2a_2\) kleiner wird als eine beliebig vorherbestimmte noch so kleine Grösse.

Wie man aber in diesem Falle die Zahlen \(m_1,m_2\) findet, lehrt die Theorie der Kettenbrüche. Ist \(\left(\frac{a_1}{a_2}\right)\) der absolute Werth von \(\frac{a_1}{a_2}\), und sind \(\frac{\mu}{\nu},\frac{\mu'}{\nu'}\) zwei Näherungswerthe der Kettenbruchsentwicklung von \(\left(\frac{a_1}{a_2}\right)\), so ist bekanntlich der absolute Werth von \(\left(\frac{a_1}{a_2}\right)-\frac{\mu}{\nu}\) kleiner als \(\frac{1}{\nu\nu'}\), mithin auch der absolute Werth von \(\nu(a_1)-\mu(a_2)\) kleiner als \(\frac{(a_2)}{\nu'}\). Man braucht daher die Kettenbruchentwicklung nur so weit fortzusetzen, bis \(\frac{(a_2)}{\nu'}\) gehörig klein ist, um die geforderte Bestimmung eines aus \((a_1),(a_2)\) gebildeten linearen Ausdrucks zu leisten, welcher kleiner als eine vorausbestimmte Grösse ist. Somit ist der im Eingange gegebene Satz allgemein bewiesen.

Bemerken wir noch wegen des Folgenden, dass im Allgemeinen zwischen den Grössen \(a\) mit gleichem obern Index keine ganzzahlige lineare Relation besteht. Deswegen wird auch im Allgemeinen keine der linearen Combinationen, welche kleiner als gegebene Grenzen sind, absolut Null werden, was das Bestehen einer solchen linearen Relation voraussetzen würde.

\subsubsection*{\S\,39. Functionen mit mehreren Perioden.}

Betrachten wir nun ein Integral erster Gattung, und den allgemeinsten Werth, welchen dasselbe für bestimmte gegebene Grenzen annehmen kann. Sind die Periodicitätsmoduln des Integrals
\[
a_1+b_1\sqrt{-1},\quad a_2+b_2\sqrt{-1},\quad\ldots,\quad a_{2p}+b_{2p}\sqrt{-1},
\]
wo die \(a,b\) nun reelle Zahlen bedeuten, und ist \(I\) ein Werth des Integrals, so ist der allgemeinste, welchen dasselbe erhalten kann:
\[
I+m_1(a_1+b_1\sqrt{-1})+m_2(a_2+b_2\sqrt{-1})\cdots+m_{2p}(a_{2p}+b_{2p}\sqrt{-1}).
\]
Ist nun \(p>1\), so kann man die Zahlen \(m_1,m_2\ldots m_{2p}\) nach dem obigen Satze immer so bestimmen, dass die Ausdrücke
\[
\begin{array}{c}
m_1a_1+m_2a_2+\cdots+m_{2p}a_{2p}\\
m_1b_1+m_2b_2+\cdots+m_{2p}b_{2p}
\end{array}
\]
beliebig klein werden. Man kann also durch passende Veränderung des Weges auch den Werth von \(I\) beliebig ändern, oder man hat den Satz:

Ein Integral erster Gattung kann für \(p>1\) bei beliebig gegebenen Grenzen durch passende Wahl des Weges jeden beliebigen Werth erhalten.

Dieser Satz könnte nur eine Ausnahme finden, wenn die Aenderung des Werthes von \(I\) nicht unendlich klein, sondern absolut Null würde, d. h. wenn es solche ganze Zahlen \(m\) gäbe, dass
\[
\begin{array}{rcl}
m_1a_1+m_2a_2+\cdots+m_{2p}a_{2p}&=&0,\\
m_1b_1+m_2b_2+\cdots+m_{2p}b_{2p}&=&0,
\end{array}
\]
oder dass zwischen den Periodicitätsmoduln selbst die Relation bestünde:
\[
m_1(a_1+b_1\sqrt{-1})+m_2(a_2+b_2\sqrt{-1})\cdots+m_{2p}(a_{2p}+b_{2p}\sqrt{-1})=0.
\]
In diesem Falle würde man diese \(2p\) Periodicitätsmoduln durch \(2p-1\) andere ganzzahlig ausdrücken können, so dass also das Integral in Wahrheit nicht \(2p\), sondern nur \(2p-1\) von einander unabhängige Periodicitätsmoduln besässe.

Setzen wir, um dies einzusehen, der Kürze wegen
\[
a_i+b_i\sqrt{-1}=P_i,
\]
so dass die Gleichung
\[
(1)\ldots\quad m_1P_1+m_2P_2+\cdots+m_{2p}P_{2p}=0
\]
vorausgesetzt wird. Ist nun \(m_2'\) der grösste gemeinschaftliche Theiler von \(m_1,m_2\), so kann man zwei Zahlen \(n_1,n_2\) so bestimmen, dass
\[
\frac{m_1}{m_2'}n_1-\frac{m_2}{m_2'}n_2=1,
\]
und indem man
\[
\begin{array}{rcl}
n_1P_1+n_2P_2&=&P_1',\\
\dfrac{m_2}{m_2'}P_1+\dfrac{m_1}{m_2'}P_2&=&P_2'
\end{array}
\]
setzt, werden \(P_1',P_2'\) zwei neue Periodicitätsmoduln, durch welche man \(P_1,P_2\) ganzzahlig ausdrücken kann. Die Gleichung (1) geht nun in die Gleichung
\[
m_2'P_2'+m_3P_3+m_4P_4+\cdots+m_{2p}P_{2p}=0
\]
über, welche einen Modul weniger enthält.

Ist ferner \(m_3''\) der grösste gemeinschaftliche Factor von \(m_2'\) und \(m_3\), und bestimmt man \(n_2',n_3'\) so, dass
\[
\frac{m_2'}{m_3''}n_2'-\frac{m_3}{m_3''}n_3'=1,
\]
so kann man wieder an Stelle von \(P_2',P_3\) die Periodicitätsmoduln
\[
\begin{array}{rcl}
n_2'P_2'+n_3'P_3&=&P_2'',\\
\dfrac{m_3}{m_3''}P_2'+\dfrac{m_2'}{m_3''}P_3&=&P_3''
\end{array}
\]
einführen, durch welche sich \(P_2',P_3\) ganzzahlig ausdrücken lassen. Man hat also \(P_1,P_2,P_3\) durch \(P_1',P_2'',P_3''\) linear und ganzzahlig ausgedrückt; die Gleichung (1) aber geht über in die Gleichung
\[
m_3''P_3''+m_4P_4+\cdots+m_{2p}P_{2p}=0,
\]
welche nur \(2p-2\) Moduln enthält.

Indem man auf diese Weise fortfährt, gelangt man dazu, \(P_1,P_2\ldots P_{2p}\) durch neue Moduln \(P_1',P_2'',P_3''\ldots P_{2p-1}^{(2p-1)}\) ganzzahlig auszudrücken, während die Gleichung (1) in
\[
P_{2p}^{(2p-1)}=0
\]
übergeht. Man kann also wirklich die \(2p\) Periodicitätsmoduln als lineare ganzzahlige Verbindungen von nur \(2p-1\) Moduln darstellen, und hat also ein Integral mit nur \(2p-1\) unabhängigen Periodicitätsmoduln vor sich.

Da das Integral erster Gattung für jede obere Grenze jeden Werth annehmen kann, so hat auch umgekehrt die obere Grenze als Function des Integrals betrachtet für jeden beliebigen Werth des letztern jeden beliebigen Werth. Man hat also eine unendlich vieldeutige Function vor sich, bei welcher es nicht sowohl auf den Werth der Argumente, als auf die Entstehungsweise derselben ankommt, d. h. auf den Weg, welchen das Integral hat durchlaufen müssen, um einen gewissen Werth zu erlangen; da lediglich durch diesen Weg die obere Grenze ihre Bestimmung findet.

\subsubsection*{\S\,40. Begriff des Jacobi'schen Umkehrproblems.}

Die hieraus erwachsenden Schwierigkeiten hat Jacobi dadurch beseitigen gelehrt, dass er statt eines einzigen Integrals die Summe von \(p\) gleichartigen Integralen mit verschiedenen obern Grenzen gleich einer neuen Variabeln \(v\) setzt. Bilden wir nun \(p\) Gleichungen, in deren jeder eine Summe von \(p\) gleichartigen Integralen einer neuen Variabel gleich gesetzt wird, so dass in jeder Gleichung die nämlichen obern Grenzen wiederkehren, jede Gleichung aber durch Benutzung eines besondern Integrals erster Gattung charakterisirt ist. Diese Gleichungen verbinden \(2\) Systeme von Variabeln: das System der obern Grenzen (oder vielmehr der ihnen entsprechenden Werthepaare \(x,s_x\)) und die \(p\) Grössen \(v\). Die Aufgabe, aus diesen Gleichungen die obern Grenzen als Functionen der Grössen \(v\) aufzustellen, ist das Jacobi'sche Umkehrproblem.

Sehen wir zunächst, wie die oben auseinandergesetzte Schwierigkeit hier zu bestehen aufhört. Die allgemeinsten Werthe der \(v\), welche durch Veränderung der Integrationswege entstehen (wobei wir natürlich voraussetzen, dass sämmtliche mit derselben obern Grenze behafteten Integrale auch auf demselben Wege zu nehmen sind), gehen aus einem Werthsysteme hervor, indem wir lineare und ganzzahlige Combinationen der Periodicitätsmoduln demselben hinzufügen; und zwar sind die entsprechenden Periodicitätsmoduln der verschiedenen Integrale immer mit derselben unbestimmten ganzen Zahl zu behaften. Theilen wir nun die \(p\) Zuwächse, welche die verschiedenen Functionen \(v\) so erhalten haben, in ihre \(2p\) reellen und imaginären Theile, so können wir durch passende Bestimmung der ganzen Zahlen zwar \(2p-1\) dieser Grössen beliebig klein machen, nicht aber die \((2p)^{\mathrm{te}}\), vielmehr wird diese letzte Grösse im Allgemeinen um so grösser, je kleiner man die erstern macht, je grösser also die Werthe der ganzen Zahlen durch fortgesetzte Näherung gemacht werden. In zwei Werthsystemen der \(v\) also, welche denselben obern Grenzen entsprechen, sind die beiden Werthe wenigstens für eine der Variabeln um eine endliche Grösse unterschieden. Man kann daher nicht mehr sagen, dass jedem System der obern Grenzen jedes System der \(v\) entspräche. Aber dies würde noch stattfinden, sobald man weniger als \(p\) Integrale in einer Gleichung und weniger als \(p\) Gleichungen einführen wollte, um auf sie ein Umkehrproblem zu begründen.

Man kann das Jacobi'sche Umkehrproblem auch, wie Jacobi selbst gethan hat,\footnote{Crelle's Journal Bd. IX.} mit dem Abel'schen Theorem in Verbindung setzen. Dieses Theorem lehrt eine Summe von beliebig vielen gleichartigen Integralen erster Gattung durch die Summe von \(p\) Integralen derselben Art ersetzen, wobei die oberen Grenzen algebraische Functionen der frühern werden. Man braucht nämlich nur durch die Puncte der Curve \(F(x,s_x)=0\) [oder \(f(x_1,x_2,x_3)=0\)], welche den obern Grenzen der gegebenen Summe entsprechen, und durch eine passende Zahl constanter Puncte der Curve eine algebraische Curve zu legen, deren Schnittpunctsystem dadurch vollständig bestimmt ist. Die \(p\) Puncte, welche dann das Schnittpunctsystem ergänzen, haben die Eigenschaft, dass die Summe der ihnen entsprechenden Integrale erster Gattung der negativen Summe der den gegebenen Puncten entsprechenden Integrale bis auf eine Constante gleich wird. Wenden wir diese Betrachtung auf alle \(p\) von einander unabhängigen Integrale erster Gattung an, so haben wir einen speciellen Fall des Umkehrproblems vor uns, in welchem nämlich die Grössen \(v\) als negative Integralsummen gegeben sind, und in welchem das Umkehrproblem sich durch eine algebraische Gleichung absolvirt.

\subsubsection*{\S\,41. Formulirung des Umkehrproblems. Abel'sche Functionen und Transcendenten.}

Bezeichnen wir durch \(x^{(1)},x^{(2)}\ldots x^{(p)}\) die den obern Grenzen, durch \(c^{(1)},c^{(2)}\ldots c^{(p)}\) die den untern Grenzen entsprechenden Puncte und Variabeln. Nehmen wir ferner die Integrale erster Gattung in der Normalform an, dann ist das Umkehrungsproblem durch folgende Gleichungen gegeben:
\[
(1)\ldots\quad
\left\{\begin{array}{rcl}
\displaystyle\sum_{i=1}^{i=p}\int_{c^{(i)}}^{x^{(i)}}du_1&=&v_1,\\[0.35em]
\displaystyle\sum_{i=1}^{i=p}\int_{c^{(i)}}^{x^{(i)}}du_2&=&v_2,\\[-0.1em]
\multicolumn{3}{c}{\cdots\cdots\cdots}\\[-0.1em]
\displaystyle\sum_{i=1}^{i=p}\int_{c^{(i)}}^{x^{(i)}}du_p&=&v_p.
\end{array}\right.
\]
Ist nun \(\frac{\varphi}{\psi}\) irgend eine homogene Function von \(x,s_x\), oder von \(x_1,x_2,x_3\), deren Zähler und Nenner von gleich hoher Ordnung sind, so werden die Werthe, welche diese Function für die verschiedenen \(p\) Puncte annimmt, als Wurzeln einer algebraischen Gleichung
\[
(2)\ldots\quad
\left(\frac{\varphi}{\psi}\right)^p+M_1\left(\frac{\varphi}{\psi}\right)^{p-1}+M_2\left(\frac{\varphi}{\psi}\right)^{p-2}\cdots+M_p=0
\]
betrachtet werden können, deren Coefficienten Functionen der \(v\) sind. Diese Functionen, oder allgemeiner, symmetrische Functionen der \(p\) Werthe, welche eine homogene Function nullter Dimension von \(x,s_x\), oder von \(x_1,x_2,x_3\) für \(p\) verschiedene Puncte annimmt, sollen Abel'sche Functionen der zugehörigen Grössen \(v\) genannt werden.

Die Beschränkung, welche darin liegt, dass wir hier nur symmetrische Functionen der \(p\) Werthe einer Function \(\frac{\varphi}{\psi}\) betrachten, ist nur eine scheinbare. Denn jede rationale symmetrische (d. h. durch Vertauschung zweier Puncte nicht veränderliche) Function der Coordinatenpaare \(x^{(1)},s_{x^{(1)}};x^{(2)},s_{x^{(2)}};\ldots x^{(p)},s_{x^{(p)}}\) ist bekanntlich aus Functionen der von uns betrachteten Art zusammensetzbar, und damit haben alle über diese speciellen Functionen anzustellenden Betrachtungen auch für jene allgemeinern Gültigkeit, und wir können die Definition der Abel'schen Functionen auf letztere ausdehnen, ohne den aufgestellten Begriff zu erweitern oder zu verengern.

Die allgemeinsten Werthe, welche die \(v\) durch Veränderung der Integrationswege bei unveränderten obern Grenzen erhalten können, sind:
\[
(3)\ldots\quad
\begin{array}{rcl}
v_1'&=&v_1+2m_1\pi\sqrt{-1}+a_{11}q_1+a_{12}q_2\cdots+a_{1p}q_p\\
v_2'&=&v_2+2m_2\pi\sqrt{-1}+a_{21}q_1+a_{22}q_2\cdots+a_{2p}q_p\\[-0.15em]
\multicolumn{3}{c}{\cdots\cdots\cdots}\\
v_p'&=&v_p+2m_p\pi\sqrt{-1}+a_{p1}q_1+a_{p2}q_2\cdots+a_{pp}q_p,
\end{array}
\]
wo die \(m,q\) beliebige ganze Zahlen bedeuten. Die Abel'schen Functionen bleiben offenbar ungeändert, wenn man die \(v\) durch die \(v'\) ersetzt; sie bleiben also ungeändert, wenn man die \(v\) gleichzeitig um eines der folgenden \(2p\) Systeme vermehrt:
\[
(4)\ldots\quad
\begin{array}{rcccccc}
v_1)&2\pi\sqrt{-1}&0&\cdots&0&a_{11}&a_{12}\cdots a_{1p}\\
v_2)&0&2\pi\sqrt{-1}&\cdots&0&a_{21}&a_{22}\cdots a_{2p}\\[-0.15em]
\multicolumn{7}{c}{\cdots\cdots\cdots}\\
v_p)&0&0&\cdots&2\pi\sqrt{-1}&a_{p1}&a_{p2}\cdots a_{pp}.
\end{array}
\]
Man nennt daher diese Grössensysteme die Systeme der Perioden der Abel'schen Functionen, und diese selbst, da es \(2p\) solcher Systeme giebt, \(2p\)fach periodische Functionen.

Betrachten wir ferner die Summe von \(p\) Integralen dritter Gattung \(\Pi_{\xi\eta}\), welche zwischen denselben Puncten \(c,x\) geführt sind und zwei gegebene Unstetigkeitspuncte \(\xi,\eta\) haben. Diese Summe
\[
(5)\ldots\quad
\sum_{i=1}^{i=p}\int_{c^{(i)}}^{x^{(i)}}d\Pi_{\xi\eta}
=T_{\xi\eta}\begin{pmatrix}x^{(1)},x^{(2)}\ldots x^{(p)}\\ c^{(1)},c^{(2)}\ldots c^{(p)}\end{pmatrix}
\]
können wir, da sie als symmetrische Function der obern Grenzen


% --- Continuation: source pp. 162--173 ---

\small

% Source: Clebsch--Gordan, Theorie der Abelschen Functionen, Google source pp. 162--173.
% Scan-first continuation. Source page 162 continues the sentence begun on source p. 161.

auftritt, die symmetrischen Functionen der letzten aber schon als Functionen der \(v\) betrachtet wurden, als Function der \(v\) selbst untersuchen. In diesem Sinne soll sie den Namen einer Abel'schen Transcendente führen, ebenso wie die durch Differentiation nach \(\xi\) entstandene Function
\[
(6)\ldots\quad
\sum_{i=1}^{i=p}\int_{c^{(i)}}^{x^{(i)}}dZ_\xi
=T_\xi\begin{pmatrix}x^{(1)}x^{(2)}\cdots x^{(p)}\\ c^{(1)}c^{(2)}\cdots c^{(p)}\end{pmatrix}.
\]
Beide Transcendenten werden wir, wenn alle Argumente darin gleichartig sind, kürzer durch
\[
T_{\xi\eta}\begin{pmatrix}x\\c\end{pmatrix},\qquad
T_\xi\begin{pmatrix}x\\c\end{pmatrix}
\]
bezeichnen.

Als Eigenschaften der Function \(T_{\xi\eta}\) kennen wir zunächst die schon dem Integral \(\Pi\) selbst zukommende, dass \(T_{\xi\eta}\) die Differenz zweier gleichartiger Functionen ist, deren eine nur von \(\xi\), deren andere ebenso nur von \(\eta\) abhängt. Und mit Rücksicht auf die Grenzen kann man also \(T_{\xi\eta}\) aus vier gleichartigen Functionen mit je \(p+1\) Argumenten zusammensetzen, dergestalt, dass
\[
\begin{aligned}
T_{\xi\eta}={}&\varphi(x^{(1)},x^{(2)}\ldots x^{(p)},\xi)
-\varphi(x^{(1)},x^{(2)}\ldots x^{(p)},\eta)\\
&-\varphi(c^{(1)},c^{(2)}\ldots c^{(p)},\xi)
+\varphi(c^{(1)},c^{(2)}\ldots c^{(p)},\eta);
\end{aligned}
\]
eine Darstellung, über welche im folgenden Abschnitt ausführlicher gehandelt werden soll.

Setzen wir ferner in den Gleichungen (5), (6) für die Integrationswege diejenigen, welche die Integralsummen \(v\) in die Integralsummen \(v'\) (3) übergehen lassen, so ändern sich die \(\Pi,Z\) um die betreffenden Periodicitätsmoduln; und wenn man diese den Sätzen der vorigen Abschnitte entnimmt, so findet man folgende Bestimmung der Periodicitätsmoduln von \(T,T\):

Wenn die Integrationswege sich so ändern, dass die Integrale erster Gattung bezüglich um
\[
2m_h\pi\sqrt{-1}+a_{1h}q_1+a_{2h}q_2+\cdots+a_{ph}q_p
\]
wachsen, so wächst
\[
T_{\xi\eta}\begin{pmatrix}x\\c\end{pmatrix}\text{ um }
q_1\int_\eta^\xi du_1+q_2\int_\eta^\xi du_2+\cdots+q_p\int_\eta^\xi du_p,
\]
\[
T_\xi\begin{pmatrix}x\\c\end{pmatrix}\text{ um }
q_1\frac{\partial u_1(\xi)}{\partial\xi}+q_2\frac{\partial u_2(\xi)}{\partial\xi}+\cdots+q_p\frac{\partial u_p(\xi)}{\partial\xi}.
\]

\subsubsection*{\S\,42. Reduction des Umkehrproblems auf die Untersuchung der Transcendente \(T_{\xi\eta}\).}

Die Bestimmung der Coefficienten in der Gleichung (2) als Functionen der \(v\), und damit das Umkehrproblem überhaupt, kann auf die Untersuchung der Transcendente \(T\) als Function der \(v\) zurückgeführt werden, indem man jene Coefficienten zunächst durch Functionen \(T\) ausdrückt.

Setzen wir nämlich in der Gleichung des Abel'schen Theorems (Gl. I. p. 127) für \(\alpha\) der Reihe nach \(x^{(1)},x^{(2)}\ldots x^{(p)}\), für \(\beta\) aber zugleich der Reihe nach \(c^{(1)},c^{(2)}\ldots c^{(p)}\). Setzen wir ferner an Stelle der Curve \(\varphi=0\) irgend eine Curve \(\varphi-\lambda\psi=0\) des aus \(\varphi=0\), \(\psi=0\) entspringenden Büschels; und bezeichnen wir durch \(\eta\) wie dort die Schnittpuncte von \(f=0\) mit \(\psi=0\), welche nicht zugleich auf \(\varphi=0\) liegen, durch \(\xi\) aber die Schnittpuncte von \(\varphi-\lambda\psi=0\) mit \(f=0\), welche nicht auf \(\varphi\) oder \(\psi\) liegen. Dann erhalten wir aus dem Abel'schen Theorem die folgende Reihe von Gleichungen:
\[
\begin{aligned}
\log\!\biggl[\Bigl(\frac{\varphi}{\psi}\Bigr)_{x^{(1)}}-\lambda\biggr]
-\log\!\biggl[\Bigl(\frac{\varphi}{\psi}\Bigr)_{c^{(1)}}-\lambda\biggr]
&=-\sum\int_{c^{(1)}}^{x^{(1)}}d\Pi_{\xi\eta},\\
(7)\ldots\quad
\log\!\biggl[\Bigl(\frac{\varphi}{\psi}\Bigr)_{x^{(2)}}-\lambda\biggr]
-\log\!\biggl[\Bigl(\frac{\varphi}{\psi}\Bigr)_{c^{(2)}}-\lambda\biggr]
&=-\sum\int_{c^{(2)}}^{x^{(2)}}d\Pi_{\xi\eta},\\[-0.3em]
&\hspace{-1.0em}\cdots\cdots\cdots\\[-0.3em]
\log\!\biggl[\Bigl(\frac{\varphi}{\psi}\Bigr)_{x^{(p)}}-\lambda\biggr]
-\log\!\biggl[\Bigl(\frac{\varphi}{\psi}\Bigr)_{c^{(p)}}-\lambda\biggr]
&=-\sum\int_{c^{(p)}}^{x^{(p)}}d\Pi_{\xi\eta}.
\end{aligned}
\]
Die Paare \(\xi\eta\) sind einander so zugeordnet wie das Abel'sche Theorem a. a. O. es vorschreibt, und die Summen rechts beziehen sich auf die so gebildeten Schnittpunctepaare. Aber die Integrationswege sind hier ebenfalls durch das Abel'sche Theorem bedingt; sie sind daher möglicherweise nicht dieselben, welche in den Integralsummen \(v\) erscheinen, sondern würden Integralsummen \(v'\) liefern, welche durch Periodicitätsmoduln von denselben unterschieden sind. Indessen hat in Wahrheit eine Aenderung in den Integrationswegen der \(x\) auf die rechten Theile der Summen (7) keinen Einfluss; denn wenn z. B. \(x^{(1)}\) einen Periodenweg durchläuft, so ändert sich die rechte Seite der ersten Gleichung (7) nur um einen Ausdruck der Form
\[
\sum_\eta\int_\eta^\xi dI,
\]
wo \(I\) ein Integral erster Gattung bedeutet, und die Summe sich auf alle Punctepaare \(\xi,\eta\) bezieht. Diese Summe aber verschwindet nach dem Abel'schen Theorem, da die \(\xi,\eta\) Schnittpunctsysteme sind, welche in einander übergeführt werden können. Daher braucht man hierauf gar keine Rücksicht zu nehmen, und man findet also, indem man die Gleichungen (7) addirt, und zugleich von den Logarithmen zu den Zahlen übergeht:
\[
\prod_{i=1}^{i=p}
\frac{\bigl(\frac{\varphi}{\psi}\bigr)_{x^{(i)}}-\lambda}
{\bigl(\frac{\varphi}{\psi}\bigr)_{c^{(i)}}-\lambda}
=e^{-\sum T_{\xi\eta}\begin{smallmatrix}x^{(1)}x^{(2)}\cdots x^{(p)}\\c^{(1)}c^{(2)}\cdots c^{(p)}\end{smallmatrix}}.
\]
Setzt man
\[
M=\biggl[\lambda-\Bigl(\frac{\varphi}{\psi}\Bigr)_{c^{(1)}}\biggr]
\biggl[\lambda-\Bigl(\frac{\varphi}{\psi}\Bigr)_{c^{(2)}}\biggr]
\cdots e^{-\sum T_{\xi\eta}\binom{x}{c}},
\]
wo also \(M\) ausser Constanten nur Functionen \(T\) enthält, so kann man mit Rücksicht auf die in \S\,41 (2) durch \(M_1,M_2\ldots\) bezeichneten Coefficienten diese Gleichung auch in der Form schreiben:
\[
\lambda^p+M_1\lambda^{p-1}+\cdots+M_p=N.
\]
Man braucht nur diese Gleichung \(p\) mal hinter einander für ebensoviel verschiedene Werthe von \(\lambda\) zu bilden, um ebensoviel Gleichungen als Unbekannte \(M\) vor sich zu haben, welche man dann mittelst dieses Systems linearer Gleichungen durch Functionen \(T\) ausdrücken kann. Man hat also den Satz:

\begin{quote}
Die Coefficienten der Gleichung \(p^{\mathrm{ten}}\) Grades, deren Wurzeln die \(p\) Functionen \(\bigl(\frac{\varphi}{\psi}\bigr)_{x^{(i)}}\) sind, werden rationale Verbindungen verschiedener Transcendenten \(e^T\).
\end{quote}

Bezüglich der Zurückführung des Umkehrungsproblems auf eine solche Gleichung \(p^{\mathrm{ten}}\) Grades mag noch bemerkt werden, dass die Auflösung einer einzigen solchen Gleichung hinreicht, um alle möglichen Functionen \(\frac{\varphi}{\psi}\) auszudrücken. Denn kennen wir die Wurzeln der einen Gleichung
\[
\Bigl(\frac{\varphi}{\psi}\Bigr)^p+M_1\Bigl(\frac{\varphi}{\psi}\Bigr)^{p-1}+\cdots+M_p=0,
\]
so kann die Bestimmung der entsprechenden \(p\) Werthe einer andern Function \(\frac{\Phi}{\Psi}\) immer auf lineare Gleichungen zurückgeführt werden. Bilden wir nämlich der Reihe nach die Gleichungen \(p^{\mathrm{ten}}\) Grades, deren Wurzeln die \(p\) Werthe der Functionen
\[
\frac{\Phi}{\Psi},\quad \frac{\Phi}{\Psi}\frac{\varphi}{\psi},\quad
\frac{\Phi}{\Psi}\left(\frac{\varphi}{\psi}\right)^2,
\ldots,\quad
\frac{\Phi}{\Psi}\left(\frac{\varphi}{\psi}\right)^{p-1}
\]
sind, und entnehmen wir jeder dieser Gleichungen den Ausdruck für die Summen der Wurzeln, so haben wir Gleichungen der Form
\[
\begin{aligned}
\Bigl(\frac{\Phi}{\Psi}\Bigr)_{x^{(1)}}+\Bigl(\frac{\Phi}{\Psi}\Bigr)_{x^{(2)}}+\cdots+\Bigl(\frac{\Phi}{\Psi}\Bigr)_{x^{(p)}}&=A,\\
\Bigl(\frac{\Phi}{\Psi}\Bigr)_{x^{(1)}}\Bigl(\frac{\varphi}{\psi}\Bigr)_{x^{(1)}}+\Bigl(\frac{\Phi}{\Psi}\Bigr)_{x^{(2)}}\Bigl(\frac{\varphi}{\psi}\Bigr)_{x^{(2)}}+\cdots+\Bigl(\frac{\Phi}{\Psi}\Bigr)_{x^{(p)}}\Bigl(\frac{\varphi}{\psi}\Bigr)_{x^{(p)}}&=A',\\
\Bigl(\frac{\Phi}{\Psi}\Bigr)_{x^{(1)}}\Bigl(\frac{\varphi}{\psi}\Bigr)_{x^{(1)}}^{2}+\Bigl(\frac{\Phi}{\Psi}\Bigr)_{x^{(2)}}\Bigl(\frac{\varphi}{\psi}\Bigr)_{x^{(2)}}^{2}+\cdots+\Bigl(\frac{\Phi}{\Psi}\Bigr)_{x^{(p)}}\Bigl(\frac{\varphi}{\psi}\Bigr)_{x^{(p)}}^{2}&=A'',\\[-0.3em]
\cdots\cdots\cdots&\\[-0.3em]
\sum_{i=1}^{p}\Bigl(\frac{\Phi}{\Psi}\Bigr)_{x^{(i)}}\Bigl(\frac{\varphi}{\psi}\Bigr)_{x^{(i)}}^{p-1}&=A^{(p-1)},
\end{aligned}
\]
wo links rationale Verbindungen verschiedener Functionen \(e^T\) stehen. Da die Coefficienten der Gleichungen als bekannt vorausgesetzt werden, so genügen dieselben, um die Functionen \(\bigl(\frac{\Phi}{\Psi}\bigr)_{x^{(i)}}\) linear aus ihnen zu bestimmen.

\subsubsection*{\S\,43. Erweitertes Umkehrproblem.}

Man kann das Jacobi'sche Umkehrproblem in folgender Weise erweitern. Nehmen wir \(p\) Summen von je \(p+q\) gleichartigen Integralen erster Gattung mit den obern Grenzpuncten \(x^{(1)},x^{(2)}\ldots x^{(p+q)}\) und den untern \(c^{(1)},c^{(2)}\ldots c^{(p+q)}\) gleich gegebenen Grössen \(v_1,v_2\ldots v_p\) an; ausserdem aber \(q\) Summen von je \(p+q\) gleichartigen Integralen dritter Gattung mit den nämlichen Grenzen gleich gegebenen Grössen \(w_1,w_2\ldots w_q\). Die Unstetigkeitspuncte der Integrale dritter Gattung seien \(\xi^{(1)},\eta^{(1)};\xi^{(2)},\eta^{(2)};\ldots;\xi^{(q)},\eta^{(q)}\). Dann sind die Gleichungen des erweiterten Umkehrproblems
\[
(1)\ldots\quad
\left\{\begin{array}{rcl}
\displaystyle\sum_{h=1}^{h=p+q}\int_{c^{(h)}}^{x^{(h)}}du_1&=&v_1,\\
\displaystyle\sum_{h=1}^{h=p+q}\int_{c^{(h)}}^{x^{(h)}}du_2&=&v_2,\\[-0.2em]
\multicolumn{3}{c}{\cdots\cdots\cdots}\\[-0.2em]
\displaystyle\sum_{h=1}^{h=p+q}\int_{c^{(h)}}^{x^{(h)}}du_p&=&v_p,\\
\displaystyle\sum_{h=1}^{h=p+q}\int_{c^{(h)}}^{x^{(h)}}d\Pi_{\xi^{(1)}\eta^{(1)}}&=&w_1,\\
\displaystyle\sum_{h=1}^{h=p+q}\int_{c^{(h)}}^{x^{(h)}}d\Pi_{\xi^{(2)}\eta^{(2)}}&=&w_2,\\[-0.2em]
\multicolumn{3}{c}{\cdots\cdots\cdots}\\[-0.2em]
\displaystyle\sum_{h=1}^{h=p+q}\int_{c^{(h)}}^{x^{(h)}}d\Pi_{\xi^{(q)}\eta^{(q)}}&=&w_q.
\end{array}\right.
\]
Die Aufgabe ist hier, die Werthepaare \(x,s_x\) oder die Coordinaten der obern Grenzpuncte als Functionen der Grössen \(v,w\) darzustellen, oder als solche die Coefficienten einer Gleichung
\[
\left(\frac{\varphi}{\psi}\right)^{p+q}+N_1\left(\frac{\varphi}{\psi}\right)^{p+q-1}+N_2\left(\frac{\varphi}{\psi}\right)^{p+q-2}\cdots+N_{p+q}=0
\]
darzustellen, deren Wurzeln die den obern Grenzen entsprechenden Werthe einer homogenen Function \(\frac{\varphi}{\psi}\) sind, deren Zähler und Nenner gleich hohe Ordnung haben.

Man kann dieses Umkehrproblem mit Hülfe des Abel'schen Theorems auf das Jacobi'sche zurückführen. Bestimmen wir zunächst bei beliebig gewählten untern Grenzpuncten \(c^{(p+q+1)},c^{(p+q+2)}\ldots c^{(2p+q)}\) obere Grenzpuncte \(x^{(p+q+1)},x^{(p+q+2)}\ldots x^{(2p+q)}\) so, dass für die entsprechenden Integrale erster Gattung die Gleichungen gelten:
\[
(2)\ldots\quad
\left\{\begin{array}{rcl}
\displaystyle\sum_{i=1}^{i=p}\int_{c^{(p+q+i)}}^{x^{(p+q+i)}}du_1&=&-v_1,\\
\displaystyle\sum_{i=1}^{i=p}\int_{c^{(p+q+i)}}^{x^{(p+q+i)}}du_2&=&-v_2,\\[-0.2em]
\multicolumn{3}{c}{\cdots\cdots\cdots}\\[-0.2em]
\displaystyle\sum_{i=1}^{i=p}\int_{c^{(p+q+i)}}^{x^{(p+q+i)}}du_p&=&-v_p,
\end{array}\right.
\]
was mit Hülfe des Jacobi'schen Umkehrproblems geschieht, so gehen die \(p\) ersten Gleichungen des Systems (1) in
\[
\sum_{i=1}^{i=2p+q}\int_{c^{(i)}}^{x^{(i)}}du_k=0,\qquad(k=1,2\ldots p)
\]
über; Gleichungen, welche nach dem Abel'schen Theorem aussagen, dass, sobald die untern Grenzpuncte mit festen Puncten der Curve \(f=0\) auf einer Curve \(m^{\mathrm{ter}}\) Ordnung liegen, auch die obern Grenzpuncte mit denselben sich auf einer solchen Curve befinden. Nehmen wir aber solche feste Puncte in genügender Anzahl an, und legen durch sie und die \(2p+q\) Puncte \(c\) eine Curve \(\Psi=0\), durch sie und die schon jetzt bekannten Puncte \(x^{(p+q+1)},x^{(p+q+2)}\ldots x^{(2p+q)}\) eine Curve \(\Phi=0\), die auch die gesuchten Puncte \(x^{(1)},x^{(2)}\ldots x^{(p+q)}\) enthalten soll, und in welcher man ausserdem so viel Coefficienten gleich Null annehmen kann, als durch Anwendung von \(f=0\) zerstörbar sind. Dann erhalten wir, indem wir auch zu den übrigen Gleichungen des Systems (1) die den Puncten \(x^{(p+q+1)},x^{(p+q+2)}\ldots x^{(2p+q)}\) entsprechenden Integrale auf beiden Seiten addiren, und dann links das Abel'sche Theorem anwenden:
\[
\log\left(\frac{\Phi}{\Psi}\right)_{\eta^{(h)}}-
\log\left(\frac{\Phi}{\Psi}\right)_{\xi^{(h)}}
=w_h+T_{\xi^{(h)}\eta^{(h)}}\begin{pmatrix}x^{(p+q+1)}\cdots x^{(2p+q)}\\c^{(p+q+1)}\cdots c^{(2p+q)}\end{pmatrix},
\]
\[
h=1,2\ldots q.
\]
Setzt man nun der Kürze wegen
\[
(3)\ldots\quad
l_h=\frac{\Psi_{\eta^{(h)}}}{\Psi_{\xi^{(h)}}}
\exp\!\left(w_h+T_{\xi^{(h)}\eta^{(h)}}\begin{pmatrix}x^{(p+q+1)}\cdots\\c^{(p+q+1)}\cdots\end{pmatrix}\right),
\]
so dass \(l_1,l_2\ldots l_q\) vollkommen gegebene Grössen sind, so hat man also für die Bestimmung der Function \(\Phi\) folgende Gleichungen:
\[
(4)\ldots\quad
\begin{array}{rcl}
\Phi_{\eta^{(1)}}-l_1\Phi_{\xi^{(1)}}&=&0,\\
\Phi_{\eta^{(2)}}-l_2\Phi_{\xi^{(2)}}&=&0,\\[-0.3em]
\multicolumn{3}{c}{\cdots\cdots\cdots}\\[-0.3em]
\Phi_{\eta^{(q)}}-l_q\Phi_{\xi^{(q)}}&=&0.
\end{array}
\]
Diese Gleichungen, welche für die Coefficienten von \(\Phi\) linear sind, genügen vollkommen, um dieselben zu bestimmen. Denn es waren in \(\Phi\) zunächst so viel Coefficienten zu Null gemacht, als durch Hinzufügung von \(f=0\) zum Verschwinden gebracht werden konnten, und sodann soviel Schnittpuncte von \(\Phi=0\) mit \(f=0\) gegeben, dass nur \(p+q\) derselben übrig blieben. Da nun \(p\) Schnittpuncte immer durch die übrigen bestimmt sind, so müssen \(q\) Gleichungen hinreichen, um die Function \(\Phi\) vollends zu bestimmen, was durch die Gleichungen (4) erreicht wird.

Wir bestimmen also auf diese Weise die Puncte \(x^{(1)},x^{(2)}\ldots x^{(p+q)}\), welche wir suchen, als die noch unbekannten \(p+q\) Schnittpuncte der Curven \(f=0\), \(\Phi=0\), deren vollständige Bestimmung im Vorigen gezeigt ist.

Nun könnte es freilich scheinen, als ob zur vollständigen Bestimmung der Function \(\Phi\) auch die vollständige Auflösung derjenigen Gleichung \(p^{\mathrm{ten}}\) Grades gefordert würde, von welcher die Bestimmung der Hülfs puncte \(x^{(p+q+1)},x^{(p+q+2)}\ldots x^{(2p+q)}\) abhängt. Aber in Wirklichkeit genügt es schon, die symmetrischen Functionen der Coordinaten dieser Puncte zu kennen, welche sich, wie oben gezeigt, als rationale Combinationen verschiedener Functionen \(e^{T_{\xi\eta}}\) darstellen. Bilden wir nämlich die Gleichungen
\[
\Phi_{x^{(p+q+1)}}=0,\quad
\Phi_{x^{(p+q+2)}}=0,\ldots,
\Phi_{x^{(2p+q)}}=0,
\]
welche aussagen, dass jene Puncte auf der Curve \(\Phi=0\) liegen sollen, so können wir doch an Stelle dieser Gleichungen die \(p\) Gleichungen setzen
\[
(5)\ldots\quad
\begin{array}{rcl}
\displaystyle\sum_{i=1}^{i=p}(\Theta_1\cdot\Phi)_{x^{(p+q+i)}}&=&0,\\
\displaystyle\sum_{i=1}^{i=p}(\Theta_2\cdot\Phi)_{x^{(p+q+i)}}&=&0,\\[-0.3em]
\multicolumn{3}{c}{\cdots\cdots\cdots}\\[-0.3em]
\displaystyle\sum_{i=1}^{i=p}(\Theta_p\cdot\Phi)_{x^{(p+q+i)}}&=&0,
\end{array}
\]
wo \(\Theta_1,\Theta_2\ldots\Theta_p\) von einander unabhängige gebrochene rationale Functionen sind, welche man, wenn homogene Coordinaten angewandt werden, so schreiben muss, dass Zähler und Nenner der Ausdrücke \(\Theta\cdot\Phi\) von gleich hoher Ordnung werden. Diese Gleichungen sagen genau dasselbe aus, wie die früheren, da ihre Determinante der Voraussetzung nach nicht verschwindet, und sie enthalten nur noch symmetrische Functionen der \(p\) Puncte, welche im Vorigen durch die Functionen \(e^{T_{\xi\eta}}\) ausgedrückt wurden. Wendet man die Gleichungen (4), (5) zur Bestimmung der Function \(\Phi\) an und fügt etwa noch die Bedingungen hinzu, welche aussagen, dass \(\Phi=0\) durch die Doppel- und Rückkehrpuncte gehen soll, so hat man die genügende Anzahl von Gleichungen vor sich, um die Coefficienten von \(\Phi\) linear zu bestimmen. Und zwar drücken alle Coefficienten von \(\Phi\) sich hiernach rational durch die Grössen \(e^w\) und durch verschiedene Transcendenten \(e^{T_{\xi\eta}}\) aus. Das erweiterte Umkehrproblem hat also die letztern mit dem ursprünglichen gemein.

Man kann nun auch, wie es dort geschah, die Werthe, welche eine gegebene Function \(\frac{\varphi}{\psi}\) für die gesuchten Puncte \(x^{(1)},x^{(2)}\ldots x^{(p)}\) annimmt, als Wurzeln einer Gleichung darstellen, deren Coefficienten rationale Functionen der \(e^w\) und verschiedener Transcendenten \(e^{T_{\xi\eta}}\) werden. Wir brauchen nämlich nur die oben angewandten Ausgangsformeln hier für alle \(p+q\) gesuchten Punctepaare \(c,x\) hinzuschreiben:
\begingroup\footnotesize
\[
\begin{aligned}
\log\!\biggl[\Bigl(\frac{\varphi}{\psi}\Bigr)_{x^{(1)}}-\lambda\biggr]
-\log\!\biggl[\Bigl(\frac{\varphi}{\psi}\Bigr)_{c^{(1)}}-\lambda\biggr]
&=-\sum\int_{c^{(1)}}^{x^{(1)}}d\Pi_{\xi\eta},\\
(6)\ldots\quad
\log\!\biggl[\Bigl(\frac{\varphi}{\psi}\Bigr)_{x^{(2)}}-\lambda\biggr]
-\log\!\biggl[\Bigl(\frac{\varphi}{\psi}\Bigr)_{c^{(2)}}-\lambda\biggr]
&=-\sum\int_{c^{(2)}}^{x^{(2)}}d\Pi_{\xi\eta},\\[-0.4em]
&\hspace{-1.0em}\cdots\cdots\cdots\\[-0.4em]
\log\!\biggl[\Bigl(\frac{\varphi}{\psi}\Bigr)_{x^{(p+q)}}-\lambda\biggr]
-\log\!\biggl[\Bigl(\frac{\varphi}{\psi}\Bigr)_{c^{(p+q)}}-\lambda\biggr]
&=-\sum\int_{c^{(p+q)}}^{x^{(p+q)}}d\Pi_{\xi\eta}.
\end{aligned}
\]
\endgroup
wo wieder die \(\xi\) die Schnittpuncte von \(f=0\) mit \(\varphi-\lambda\psi=0\), die \(\eta\) die Schnittpuncte von \(f=0\) mit \(\psi=0\) bedeuten. Indem man nun wieder summirt, erhält man
\begingroup\footnotesize
\[
\begin{aligned}
\log\Biggl\{&
\frac{\bigl[(\frac{\varphi}{\psi})_{x^{(1)}}-\lambda\bigr]
\bigl[(\frac{\varphi}{\psi})_{x^{(2)}}-\lambda\bigr]\cdots}
{\bigl[(\frac{\varphi}{\psi})_{c^{(1)}}-\lambda\bigr]
\bigl[(\frac{\varphi}{\psi})_{c^{(2)}}-\lambda\bigr]\cdots}
\cdot
\frac{\bigl[(\frac{\varphi}{\psi})_{x^{(p+q)}}-\lambda\bigr]}
{\bigl[(\frac{\varphi}{\psi})_{c^{(p+q)}}-\lambda\bigr]}
\Biggr\}\\
&={} -\sum_{\xi\eta}\left\{\int_{c^{(1)}}^{x^{(1)}}d\Pi_{\xi\eta}+
\int_{c^{(2)}}^{x^{(2)}}d\Pi_{\xi\eta}+\cdots+
\int_{c^{(p+q)}}^{x^{(p+q)}}d\Pi_{\xi\eta}\right\}.
\end{aligned}
\]
\endgroup
Die rechte Seite dieser Gleichung, vermehrt um die Summe
\[
-T_{\xi\eta}\begin{pmatrix}x^{(p+q+1)}\cdots x^{(2p+q)}\\c^{(p+q+1)}\cdots c^{(2p+q)}\end{pmatrix},
\]
ist nach dem Abel'schen Theorem gleich
\[
\log\left[\left(\frac{\Phi_\xi}{\Phi_\eta}\right)\cdot\left(\frac{\Psi_\eta}{\Psi_\xi}\right)\right],
\]
wo \(\Phi,\Psi\) die oben definirten Ausdrücke bedeuten, so dass jene Seite diesem Logarithmus vermehrt um
\[
-T_{\xi\eta}\begin{pmatrix}x^{(p+q+1)}\cdots x^{(2p+q)}\\c^{(p+q+1)}\cdots c^{(2p+q)}\end{pmatrix}
\]
gleich wird. Trägt man dies ein und geht von den Logarithmen zu den Zahlen über, so hat man
\begingroup\footnotesize
\[
\begin{aligned}
(7)\ldots\quad
&\frac{\bigl[(\frac{\varphi}{\psi})_{x^{(1)}}-\lambda\bigr]
\bigl[(\frac{\varphi}{\psi})_{x^{(2)}}-\lambda\bigr]\cdots
\bigl[(\frac{\varphi}{\psi})_{x^{(p+q)}}-\lambda\bigr]}
{\bigl[(\frac{\varphi}{\psi})_{c^{(1)}}-\lambda\bigr]
\bigl[(\frac{\varphi}{\psi})_{c^{(2)}}-\lambda\bigr]\cdots
\bigl[(\frac{\varphi}{\psi})_{c^{(p+q)}}-\lambda\bigr]}\\
&\qquad=\frac{\Phi_{\xi^{(1)}}\Psi_{\eta^{(1)}}}{\Phi_{\eta^{(1)}}\Psi_{\xi^{(1)}}}\cdot
\frac{\Phi_{\xi^{(2)}}\Psi_{\eta^{(2)}}}{\Phi_{\eta^{(2)}}\Psi_{\xi^{(2)}}}\cdots
 e^{-\sum_{\xi\eta}T_{\xi\eta}\begin{smallmatrix}x^{(p+q+1)}\cdots\\c^{(p+q+1)}\cdots\end{smallmatrix}}.
\end{aligned}
\]
\endgroup
Setzen wir in (7) für \(\lambda\) wieder der Reihe nach \(\lambda_1,\lambda_2\ldots\lambda_{p+q}\), so haben wir lineare Gleichungen, um die symmetrischen Functionen der Grössen \(\bigl(\frac{\varphi}{\psi}\bigr)_x\) zu bilden, und also die Coefficienten der Gleichung darzustellen, deren Wurzeln dieselben sind. Auch diese Ausdrücke sind, wie man sieht, rationale Functionen der \(e^w\) und verschiedener Transcendenten \(e^{T_{\xi\eta}}\), wie dies bei den früheren Problemen der Fall war.

Die Bemerkung des \S\,42, dass die Kenntniss eines Functionensystems \(\frac{\varphi}{\psi}\) genügt, um alle andern daraus linear zu bestimmen, bleibt hier natürlich vollkommen in Gültigkeit.

\subsubsection*{\S\,44. Perioden der Functionen des erweiterten Umkehrproblems. Modification beim Eintreten von Integralen zweiter Gattung.}

Bezeichnen wir auch diese symmetrischen Functionen als Abel'sche Functionen, so sehen wir, dass dieselben ebenfalls periodische Functionen ihrer Argumente \(v_1,v_2\ldots v_p,w_1,w_2\ldots w_q\) sind. Nun ist das allgemeinste Werthsystem, welches aus einem gegebenen System \(v_1,v_2\ldots v_p,w_1,w_2\ldots w_q\) durch Abänderung der Integrationswege entstehen kann, das folgende:
\[
\begin{array}{rcl}
v_1'&=&v_1+2m_1\pi\sqrt{-1}+\displaystyle\sum_{h=1}^{h=p}a_{1h}q_h,\\
v_2'&=&v_2+2m_2\pi\sqrt{-1}+\displaystyle\sum_{h=1}^{h=p}a_{2h}q_h,\\[-0.3em]
\multicolumn{3}{c}{\cdots\cdots\cdots}\\[-0.3em]
v_p'&=&v_p+2m_p\pi\sqrt{-1}+\displaystyle\sum_{h=1}^{h=p}a_{ph}q_h,\\
w_1'&=&w_1+2n_1\pi\sqrt{-1}+\displaystyle\sum_{h=1}^{h=p}q_h\int_{\eta^{(1)}}^{\xi^{(1)}}du_h,\\
w_2'&=&w_2+2n_2\pi\sqrt{-1}+\displaystyle\sum_{h=1}^{h=p}q_h\int_{\eta^{(2)}}^{\xi^{(2)}}du_h,\\[-0.3em]
\multicolumn{3}{c}{\cdots\cdots\cdots}\\[-0.3em]
w_q'&=&w_q+2n_q\pi\sqrt{-1}+\displaystyle\sum_{h=1}^{h=p}q_h\int_{\eta^{(q)}}^{\xi^{(q)}}du_h.
\end{array}
\]
Diese Functionen sind also \(2p+q\)fach periodisch; und zwar erscheinen hier \(p+q\) Perioden \(2\pi\sqrt{-1}\), welche den verschiedenen Argumenten unabhängig von einander zukommen.

Wir können auch Transcendenten diesem erweiterten Umkehrproblem entsprechend bilden. Aber wie oben schon gezeigt, unterscheidet sich eine Summe von \(p+q\) gleichartigen Integralen dritter Gattung von der negativ genommenen Transcendente
\[
T_{\xi\eta}\begin{pmatrix}x^{(p+q+1)}\cdots\\c^{(p+q+1)}\cdots\end{pmatrix}
\]
nur um den Logarithmus einer Function, welche, wie oben gezeigt, aus den \(e^w\) und verschiedenen Ausdrücken \(e^{T_{\xi\eta}}\) sich rational zusammensetzt. Bezeichnen wir nämlich durch \(T^{(q)}_{\xi\eta}\) die Transcendente des erweiterten Umkehrproblems
\[
T^{(q)}_{\xi\eta}\begin{pmatrix}x^{(1)},x^{(2)}\cdots x^{(p+q)}\\c^{(1)},c^{(2)}\cdots c^{(p+q)}\end{pmatrix}
=\sum_{i=1}^{i=p+q}\int_{c^{(i)}}^{x^{(i)}}d\Pi_{\xi\eta},
\]
wenn darin die Integrationswege mit denen der Gleichungen (10) übereinstimmen, so haben wir bei der oben bereits getroffenen Wahl der Integrationswege in \(T\):
\begingroup\footnotesize
\[
T^{(q)}_{\xi\eta}\begin{pmatrix}x^{(1)},x^{(2)}\cdots x^{(p+q)}\\c^{(1)},c^{(2)}\cdots c^{(p+q)}\end{pmatrix}
=-T_{\xi\eta}\begin{pmatrix}x^{(p+q+1)}\cdots x^{(2p+q)}\\c^{(p+q+1)}\cdots c^{(2p+q)}\end{pmatrix}
-\log\frac{\Phi_\xi\Psi_\eta}{\Phi_\eta\Psi_\xi}.
\]
\endgroup
An diese Transcendente knüpft sich eine zweite Behandlung des erweiterten Umkehrproblems, bei welcher man die symmetrischen Functionen der Werthe von \(\frac{\varphi}{\psi}\) ganz ebenso auf die Transcendente \(T_{\xi\eta}^{(q)}\) zurückführt, wie oben bei dem Jacobi'schen Umkehrproblem die entsprechenden Functionen auf \(T_{\xi\eta}\) zurückgeführt wurden.

Aus der obigen Gleichung für \(T_{\xi\eta}^{(q)}\) erhält man durch Differentiation eine andre, welche die Transcendente
\[
T_{\xi}^{(q)}\begin{pmatrix}x^{(1)},x^{(2)}\cdots x^{(p+q)}\\c^{(1)},c^{(2)}\cdots c^{(p+q)}\end{pmatrix}
=\sum_{i=1}^{i=p+q}\int_{c^{(i)}}^{x^{(i)}}dZ_\xi
\]
auf \(T_\xi\) zurückführt. Denn man hat bei gleicher Wahl der Integrationswege in \(T_\xi\):
\[
T_\xi^{(q)}\begin{pmatrix}x^{(1)},x^{(2)}\cdots x^{(p+q)}\\c^{(1)},c^{(2)}\cdots c^{(p+q)}\end{pmatrix}
=-T_\xi\begin{pmatrix}x^{(p+q+1)}\cdots x^{(2p+q)}\\c^{(p+q+1)}\cdots c^{(2p+q)}\end{pmatrix}
-\frac{\partial}{\partial\xi}\log\frac{\Phi_\xi}{\Psi_\xi}.
\]
Von Wichtigkeit ist ferner die Bemerkung, dass von den in dem erweiterten Umkehrproblem gegebenen Integralen dritter Gattung einige durch das Zusammenrücken der betreffenden Unstetigkeitspuncte \(\xi,\eta\) in Integrale zweiter Gattung übergehen können. Bei denjenigen Grössen \(w\), für welche dies eintritt, verschwindet dann die von den Integralen dritter Gattung herrührende logarithmische Periode \(2n\pi\sqrt{-1}\); wenn also \(r\) Integralsummen dritter Gattung in Integralsummen zweiter Gattung übergehen, so sind die entsprechenden Abel'schen Functionen nur \((2p+q-r)\)fach periodisch. In der Behandlung des Umkehrproblems aber würde hierbei eine Modification eintreten.

Eine hieher gehörige der Gleichungen (1) lautet nämlich nicht mehr:
\[
\sum_{h=1}^{h=p+q}\int_{c^{(h)}}^{x^{(h)}}d\Pi_{\xi^{(i)}\eta^{(i)}}=w_i,
\]
sondern
\[
\sum_{h=1}^{h=p+q}\int_{c^{(h)}}^{x^{(h)}}dZ_{\xi^{(i)}}=t_i.
\]
Indem man also jetzt rechts und links die den \(p\) Hülfsgrössen \(x\) entsprechenden Integrale hinzufügt, findet man:
\[
\sum_{h=1}^{h=2p+q}\int_{c^{(h)}}^{x^{(h)}}dZ_{\xi^{(i)}}=t_i+T_{\xi^{(i)}}\begin{pmatrix}x^{(p+q+1)}\cdots x^{(2p+q)}\\c^{(p+q+1)}\cdots c^{(2p+q)}\end{pmatrix},
\]
und daher mit Anwendung des Abel'schen Theorems:
\[
-\frac{\partial}{\partial\xi^{(i)}}\log\frac{\Phi_{\xi^{(i)}}}{\Psi_{\xi^{(i)}}}
=t_i+T_{\xi^{(i)}}\begin{pmatrix}x^{(p+q+1)}\cdots x^{(2p+q)}\\c^{(p+q+1)}\cdots c^{(2p+q)}\end{pmatrix},
\]
oder
\begingroup\footnotesize
\[
\Phi_{\xi^{(i)}}\frac{\partial\Psi_{\xi^{(i)}}}{\partial\xi^{(i)}}
-\Psi_{\xi^{(i)}}\frac{\partial\Phi_{\xi^{(i)}}}{\partial\xi^{(i)}}
=\Phi_{\xi^{(i)}}\Psi_{\xi^{(i)}}
\left[t_i+T_{\xi^{(i)}}\begin{pmatrix}x^{(p+q+1)}\cdots x^{(2p+q)}\\c^{(p+q+1)}\cdots c^{(2p+q)}\end{pmatrix}\right].
\]
\endgroup
Dieses ist ebenfalls eine lineare Gleichung für die Coefficienten von \(\Phi\) und sie muss an Stelle der entsprechenden Gleichung (4) gesetzt werden. Die Coefficienten von \(\Phi\) enthalten also die \(t\) in rationaler Weise, und ausser verschiedenen Functionen \(e^{T_{\xi\eta}}\) auch noch Functionen \(T_\xi=\frac{\partial T_{\xi\eta}}{\partial\xi}\), welche den Integralsummen zweiter Gattung entsprechen.

Aber hierin besteht auch in der That die ganze Veränderung. Denn wenn man übrigens die Behandlungsweise des erweiterten Umkehrproblems beibehält, ergiebt sich nur eine Veränderung in den Ausdrücken der \(T_{\xi\eta}\), insofern diese von \(\Phi\) abhängen und also durch die veränderte Bildungsweise dieser Function modificirt werden.

Von besonderm Interesse ist der Fall des erweiterten Umkehrproblems, in welchem die Integrale dritter und zweiter Gattung solche sind, welche in den Doppel- und Rückkehrpuncten von \(f=0\) unstetig werden.

% --- continuation source pp. 174--181 ---

% Source: Clebsch--Gordan, Theorie der Abelschen Functionen, Google source pp. 174--181.
% Continuation slice only; scan render and equation zooms used as authority.

\begin{center}
{\large\bfseries Siebenter Abschnitt.}\par\vspace{0.8em}
{\Large\bfseries Die Function \(T_{\xi\eta}\).}
\end{center}

\subsubsection*{\S\,45. Zerlegung der Function \(T_{\xi\eta}\) in die Differenz zweier Functionen, deren jede nur von \(p+2\) Puncten abhängt.}

Die Transcendente
\[
T_{\xi\eta}\begin{pmatrix}x^{(1)},x^{(2)}\ldots x^{(p)}\\ \alpha^{(1)},\alpha^{(2)}\ldots \alpha^{(p)}\end{pmatrix}
=\int_{\alpha^{(1)}}^{x^{(1)}}d\Pi_{\xi\eta}
+\int_{\alpha^{(2)}}^{x^{(2)}}d\Pi_{\xi\eta}
+\cdots+\int_{\alpha^{(p)}}^{x^{(p)}}d\Pi_{\xi\eta}
\]
hängt im Ganzen noch von \(2p+2\) Puncten ab, nämlich von \(\xi\), von \(\eta\) und von den \(2p\) Grenzen der benutzten Integrale. Aber wir haben schon oben (\S\,41) darauf aufmerksam gemacht, dass diese Transcendente sich aus vier Functionen zusammensetzen lässt, welche nur von je \(p+1\) Argumenten abhängen, dergestalt, dass
\[
\begin{aligned}
T_{\xi\eta}\binom{x}{\alpha}={}&\varphi(x^{(1)},x^{(2)}\ldots;\xi)
-\varphi(x^{(1)},x^{(2)}\ldots;\eta)\\
&-\varphi(\alpha^{(1)},\alpha^{(2)}\ldots;\xi)
+\varphi(\alpha^{(1)},\alpha^{(2)}\ldots;\eta).
\end{aligned}
\]
Die Function \(\varphi\) ist keineswegs bestimmt; vielmehr kann \(\varphi(x^{(1)},x^{(2)}\ldots;\xi)\) um eine beliebige Function der \(x\) und ausserdem um eine beliebige Function von \(\xi\) geändert werden; wenn sodann auch die andern Functionen die entsprechenden Aenderungen erfahren, so bleibt der Ausdruck für \(T_{\xi\eta}\) völlig ungeändert.

Man kann nun eine Theilung von \(T_{\xi\eta}\) sehr leicht auf die Art bewerkstelligen, dass man \(T_{\xi\eta}\) in die Differenz zweier Integralsummen mit beliebig gewählten untern Grenzen zerlegt, und dann mit Einführung eines neuen constanten Puncts \(\zeta\) und mit Benutzung der Formel
\[
d\Pi_{\xi\eta}=d\Pi_{\xi\zeta}-d\Pi_{\eta\zeta}
\]
jede dieser Summen wieder als Differenz darstellt. Aber der hierdurch gemachte Fortschritt ist nur ein scheinbarer, da \(p+1\) neue Grössen eingeführt worden sind, und also jeder der vier Theile wieder von \(2p+2\) Grössen abhängt. Will man daher die Function \(T_{\xi\eta}\) auf genuine Weise in Functionen zerlegen, welche von einer geringern Anzahl von Grössen abhängen, so muss man einen andern Weg einschlagen.

Bemerken wir zunächst Folgendes. Wenn man zwei Curven \((n-2)^{\mathrm{ter}}\) Ordnung legt, welche durch die übrigen auf der Geraden \(\xi,\eta\) liegenden Puncte von \(f=0\), und durch die Doppel- und Rückkehrpuncte von \(f=0\) gehen, und von denen sodann die eine noch durch die Puncte \(x\), die andre durch die Puncte \(\alpha\) geht, so sind beide Curven genau bestimmt, und jede derselben schneidet \(f=0\) noch in \(p\) andern Puncten. Die mit den \(x\) auf einer Curve gelegenen seien \(y^{(1)},y^{(2)}\ldots y^{(p)}\), die mit den \(\alpha\) auf einer Curve gelegenen \(\beta^{(1)},\beta^{(2)}\ldots\beta^{(p)}\). Es hängen dann die \(y\) nur von den \(x\) und von \(\xi,\eta\), die \(\beta\) nur von den \(\alpha\) und \(\xi,\eta\) ab. Aber die beiden Curven stehen in der besondern Beziehung zu einander, dass sie sich mit \(f=0\) in \(n-2\) Puncten schneiden, welche auf einer Geraden liegen. Daher tritt hier der Fall des Satzes (\S\,13. 6) ein, wonach die Summe der entsprechenden Integrale dritter Gattung bei Bildung des Abel'schen Satzes nicht einer logarithmischen Function gleich wird, sondern verschwindet. Daher hat man
\[
\sum_{h=1}^{h=p}\int_{\alpha^{(h)}}^{x^{(h)}}d\Pi_{\xi\eta}
+\sum_{h=1}^{h=p}\int_{\beta^{(h)}}^{y^{(h)}}d\Pi_{\xi\eta}=0;
\]
wobei die Integrationswege der zweiten Summe so zu wählen sind, dass, wenn man in der ersten die \(x\) den \(\alpha\) nähert, indem man die Integrationswege verschwindend klein macht, sich auch die \(y\) den \(\beta\) so nähern, dass die Integrationswege verschwindend klein werden; und wo natürlich ein und derselbe Weg von \(\xi\) nach \(\eta\) von allen Integrationswegen nicht geschnitten wird.

Diese Gleichung kann man nun auch in die Form kleiden:
\[
2\sum_{h=1}^{h=p}\int_{\alpha^{(h)}}^{x^{(h)}}d\Pi_{\xi\eta}
=\sum_{h=1}^{h=p}\int_{y^{(h)}}^{x^{(h)}}d\Pi_{\xi\eta}
-\sum_{h=1}^{h=p}\int_{\beta^{(h)}}^{\alpha^{(h)}}d\Pi_{\xi\eta};
\]
dabei müssen nun auch die Integrationswege so genommen werden, dass diese Formel mit der vorigen übereinkommt. Es muss also der Integrationsweg von \(y^{(h)}\) bis \(x^{(h)}\) sich aus dem frühern von \(\alpha^{(h)}\) nach \(x^{(h)}\) gelegten Wege und einem von \(y^{(h)}\) nach \(\alpha^{(h)}\) gelegten zusammensetzen, während der letztere mit dem frühern Wege von \(\beta^{(h)}\) nach \(y^{(h)}\) den neuen Weg von \(\beta^{(h)}\) nach \(\alpha^{(h)}\) liefert. Wenn wir nun die Functionen \(T\) einführen und bei jeder derselben die hier stattfindenden obern und untern Grenzwerthe bemerken, so haben wir:
\[
\begin{aligned}
(1)\ldots\quad
T_{\xi\eta}\begin{pmatrix}x^{(1)},x^{(2)}\ldots x^{(p)}\\ \alpha^{(1)},\alpha^{(2)}\ldots \alpha^{(p)}\end{pmatrix}
&=\frac12 T_{\xi\eta}\begin{pmatrix}x^{(1)},x^{(2)}\ldots x^{(p)}\\ y^{(1)},y^{(2)}\ldots y^{(p)}\end{pmatrix}\\
&\quad-\frac12 T_{\xi\eta}\begin{pmatrix}\alpha^{(1)},\alpha^{(2)}\ldots\alpha^{(p)}\\ \beta^{(1)},\beta^{(2)}\ldots\beta^{(p)}\end{pmatrix};
\end{aligned}
\]
eine Gleichung, welche wir, da die untern Grenzen rechts durch die obern bestimmt sind, der Einfachheit wegen in der Form schreiben wollen
\[
T_{\xi\eta}\binom{x}{\alpha}=\frac12T_{\xi\eta}(x)-\frac12T_{\xi\eta}(\alpha),
\]
in welcher rechts die untern Grenzen nicht besonders angegeben sind. Wir sehen hier ohne Einführung neuer Constanten die allgemeine Function \(T_{\xi\eta}\) in die halbe Differenz zweier andern Functionen derselben Art umgeformt, deren eine keine Spur der \(\alpha\), deren andre keine Spur der \(x\) mehr enthält, deren jede also nur noch von \(p+2\) Grössen abhängig ist. Wir werden also fortan die einfachere Function
\[
\frac12T_{\xi\eta}\begin{pmatrix}x^{(1)},x^{(2)}\ldots x^{(p)}\\ y^{(1)},y^{(2)}\ldots y^{(p)}\end{pmatrix}=\frac12T_{\xi\eta}(x)
\]
zu untersuchen haben.

Die Function \(\frac12T_{\xi\eta}(x)\) ist die Differenz zweier gleichartigen Functionen, deren eine \(\xi\), nicht aber \(\eta\) enthält, während die andere sich von ersterer dadurch unterscheidet, dass \(\eta\) an die Stelle von \(\xi\) getreten ist. Diese Eigenschaft erscheint in dem obigen Ausdruck verdeckt, da auch die untern Grenzen \(y\) noch von \(\xi,\eta\) abhängen; es ist deswegen nothwendig, sie zu beweisen.

\subsubsection*{\S\,46. Beweis für die Zerlegbarkeit der Function \(\frac12T_{\xi\eta}(x)\) in zwei von nur \(p+1\) Puncten abhängige Functionen.}

Bemerken wir zu diesem Zwecke, dass nach der Formel (1) unser Ausdruck sich durch
\[
\sum\int_{\alpha^{(h)}}^{x^{(h)}}d\Pi_{\xi\eta}
-\frac12\sum\int_{\beta^{(h)}}^{\alpha^{(h)}}d\Pi_{\xi\eta}
\]
darstellen lässt, wo die \(\alpha\) irgend \(p\) constante Puncte bedeuten. Von dem ersten Theile dieses Ausdrucks wissen wir schon, dass er die erforderliche Eigenschaft besitzt, welche constanten Puncte \(\alpha\) wir auch wählen; es ist daher nur der zweite Ausdruck noch zu betrachten, in welchem die \(\beta\) noch von \(\xi,\eta\) abhängig sind. Nähern wir nun die \(p\) constanten Puncte \(\alpha\) einer solchen Lage, bei welcher sie sich mit den Doppel- und Rückkehrpuncten von \(f=0\) auf einer Curve \((n-3)^{\mathrm{ter}}\) Ordnung befinden. Alsdann geht die Curve \((n-2)^{\mathrm{ter}}\) Ordnung, auf welcher die \(\alpha\) und \(\beta\) mit jenen Ausnahmspuncten und mit den \(n-2\) übrigen Puncten der Linie \(\xi\eta\) liegen sollen, in die Verbindung dieser Curve \((n-3)^{\mathrm{ter}}\) Ordnung mit der Linie \(\xi\eta\) selbst über. Von den Puncten \(\beta\) also fallen zwei in \(\xi,\eta\) selbst, während die \(p-2\) übrigen in die übrigen Schnittpuncte der Curve \((n-3)^{\mathrm{ter}}\) Ordnung mit der Curve \(f=0\) fallen. Von den \(p\) Gliedern der Summe
\[
\sum\int_{\beta^{(h)}}^{\alpha^{(h)}}d\Pi_{\xi\eta}
\]
erhalten also \(p-2\) constante Grenzen, und besitzen demnach ebenfalls die Eigenschaft, in die Differenz von zwei Functionen zu zerfallen, deren eine nicht mehr von \(\eta\), deren andere nicht mehr von \(\xi\) abhängt. Daher bleibt nur die Grenze zu untersuchen, welcher sich die beiden übrigen Integrale nähern, deren untere Grenzen wir vorläufig als von \(\xi,\eta\) wenig verschieden annehmen, und somit durch \(\xi+\delta x,\eta+\delta y\) bezeichnen wollen. Die aus diesem Grenzübergang hervorgehende Integralsumme bedarf einer genaueren Untersuchung, weil in derselben die unteren Grenzen mit Unstetigkeitspuncten zusammenfallen. Statt der Integrale
\[
\int_{\xi+\delta x}^{\alpha^{(1)}}d\Pi_{\xi\eta}
+\int_{\eta+\delta y}^{\alpha^{(2)}}d\Pi_{\xi\eta}
\]
können wir aber zunächst nach einem oft angewandten Theorem setzen:
\[
\begin{aligned}
&\int_{\xi+\delta x}^{\alpha^{(1)}}d\Pi_{\xi\zeta}
-\int_{\xi+\delta x}^{\alpha^{(1)}}d\Pi_{\eta\zeta}
+\int_{\eta+\delta y}^{\alpha^{(2)}}d\Pi_{\xi\zeta}
-\int_{\eta+\delta y}^{\alpha^{(2)}}d\Pi_{\eta\zeta}.
\end{aligned}
\]
Im zweiten und dritten dieser Integrale kann man den Grenzübergang sofort ausführen, und erhält
\[
\int_{\eta}^{\alpha^{(2)}}d\Pi_{\xi\zeta}-\int_{\xi}^{\alpha^{(1)}}d\Pi_{\eta\zeta},
\]
in welcher Differenz wegen des Satzes über Vertauschung von Parameter und Argument die von \(\xi,\eta\) zugleich abhängenden Theile sich zerstören, denn man hat sofort
\[
\begin{aligned}
\int_{\eta}^{\alpha^{(2)}}d\Pi_{\xi\zeta}-\int_{\xi}^{\alpha^{(1)}}d\Pi_{\eta\zeta}
={}&\int_{\eta}^{\delta}d\Pi_{\xi\varepsilon}
+\int_{\eta}^{\delta}d\Pi_{\varepsilon\zeta}
+\int_{\delta}^{\alpha^{(2)}}d\Pi_{\xi\zeta}\\
&-\int_{\xi}^{\varepsilon}d\Pi_{\eta\delta}
-\int_{\xi}^{\delta}d\Pi_{\delta\zeta}
-\int_{\varepsilon}^{\alpha^{(1)}}d\Pi_{\eta\zeta},
\end{aligned}
\]
wo in der That das erste positive Glied gegen das erste negative sich aufhebt. Es bleibt also endlich der Ausdruck zu betrachten:
\[
\lim\left\{\int_{\xi+\delta x}^{\alpha^{(1)}}d\Pi_{\xi\zeta}
-\int_{\eta+\delta y}^{\alpha^{(2)}}d\Pi_{\eta\zeta}\right\}.
\]
Addirt man nun unter dem ersten Integralzeichen \(d\log\Sigma\pm x_1\xi_2\zeta_3\), so entsteht das Integral
\[
\int_{\xi+\delta x}^{\alpha^{(1)}}d\bigl(\Pi_{\xi\zeta}(x)+\log\Sigma\pm x_1\xi_2\zeta_3\bigr),
\]
welches im Puncte \(\xi\) nicht mehr unstetig ist und daher unmittelbar in das Integral von \(\xi\) bis \(\alpha^{(1)}\) verwandelt werden kann, welches nur noch von \(\xi\) abhängt. Ebenso wird das letzte der obigen Integrale durch Zufügung von \(d\log\Sigma\pm y_1\eta_2\zeta_3\) in
\[
\int_{\eta+\delta y}^{\alpha^{(2)}}d\bigl(\Pi_{\eta\zeta}(y)+\log\Sigma\pm y_1\eta_2\zeta_3\bigr)
\]
verwandelt, welches in \(\eta\) stetig ist, daher durch das Integral von \(\eta\) bis \(\alpha^{(2)}\) ersetzt werden kann, und also nur von \(\eta\) abhängt.

So ist die Betrachtung unseres Ausdrucks zurückgeführt auf die Untersuchung der Summe der oben zugefügten Integrale, welche wieder abgezogen werden müssen, und so die Terme liefern:
\[
-\int_{\alpha^{(1)}}^{\xi+\delta x}d\log\Sigma\pm x_1\xi_2\zeta_3
+\int_{\alpha^{(2)}}^{\eta+\delta y}d\log\Sigma\pm y_1\eta_2\zeta_3
\]
\[
=-\log\left(\frac{\Sigma\pm\delta x_1\xi_2\zeta_3}{\Sigma\pm\alpha_1^{(1)}\xi_2\zeta_3}\cdot
\frac{\Sigma\pm\alpha_1^{(2)}\eta_2\zeta_3}{\Sigma\pm\delta y_1\eta_2\zeta_3}\right),
\]
bei welchem wir übrigens nur auf den Ausdruck
\[
\log\frac{\Sigma\pm\delta x_1\xi_2\zeta_3}{\Sigma\pm\delta y_1\eta_2\zeta_3}
\]
zu achten haben, da die andern Theile die geforderte Eigenschaft bereits besitzen.

Die Differentiale \(\delta x,\delta y\) sind dadurch an einander gebunden, dass der Voraussetzung nach die Puncte \(\xi+\delta x,\eta+\delta y\) immer mit den \(n-2\) andern Puncten der Geraden \(\xi\eta\), mit den Doppel- und Rückkehrpuncten und \(p\) constanten Puncten \(\alpha\) auf einer Curve \((n-2)^{\mathrm{ter}}\) Ordnung liegen müssen. Die hierdurch ausgesprochene Bedingung ist genau so zu bilden, wie die Function \((n-2)^{\mathrm{ter}}\) Ordnung \(\Omega\), welche unter dem Integralzeichen des Integrals dritter Gattung im Zähler steht, und deren Bildung wir im ersten Abschnitt ausführlich gelehrt haben. Wir erhalten eine Determinante von \(\frac{n(n-1)}{2}\) Reihen, welche das Resultat einer Elimination ist. Bezeichnen wir durch \(w=0\) die Gleichung einer Curve \((n-2)^{\mathrm{ter}}\) Ordnung, und wird, wenn darin \(\xi_1+\lambda\eta_1,\xi_2+\lambda\eta_2,\xi_3+\lambda\eta_3\) an Stelle der Variabeln treten,
\[
w(\xi+\lambda\eta)=v+\lambda v_1+\lambda^2v_2+\cdots+\lambda^{n-2}v_{n-2};
\]
wird ferner
\[
\frac{f(\xi+\lambda\eta)}{\lambda}=u_1+\lambda u_2+\lambda^2u_3+\cdots+\lambda^{n-2}u_{n-1},
\]
so schreibe man zunächst die Gleichung \(w=0\), geschrieben mit den Variabeln \(\xi+\delta x\), sodann mit \(\eta+\delta y\), ferner geschrieben mit den constanten Argumenten, welche den Doppel- und Rückkehrpuncten und den Puncten \(\alpha\) entsprechen. Diesen Gleichungen füge man das System hinzu:
\[
\begin{array}{rcl}
v&=&\mu u_1,\\
v_1&=&\mu u_2,\\[-0.3em]
\multicolumn{3}{c}{\cdots\cdots\cdots}\\[-0.3em]
v_{n-2}&=&\mu u_{n-1}.
\end{array}
\]
Man eliminire die Coefficienten von \(w\), welche auch in den \(v\) linear vorkommen, und \(\mu\); das Resultat ist nach \S\,6 die gesuchte Gleichung.

Aber an Stelle der Gleichungen
\[
w(\xi+\delta x)=0,\qquad w(\eta+\delta y)=0
\]
kann man die Gleichungen setzen:
\[
\begin{aligned}
w(\xi+\delta x)-v&=-\mu u_1
=\delta x_1\frac{\partial w(\xi)}{\partial\xi_1}
+\delta x_2\frac{\partial w(\xi)}{\partial\xi_2}
+\delta x_3\frac{\partial w(\xi)}{\partial\xi_3},\\
w(\eta+\delta y)-v_{n-2}&=-\mu u_{n-1}
=\delta y_1\frac{\partial w(\eta)}{\partial\eta_1}
+\delta y_2\frac{\partial w(\eta)}{\partial\eta_2}
+\delta y_3\frac{\partial w(\eta)}{\partial\eta_3}.
\end{aligned}
\]
Indem diese Gleichungen an Stelle der beiden ersten treten, enthalten die ersten Reihen der Eliminationsdeterminante nur noch unendlich kleine Terme, bis auf die letzten Glieder, welche bezüglich \(-u_1\) und \(-u_{n-1}\) sind. Um also nur Terme erster Ordnung zu erhalten, müssen wir die Determinante auf die mit diesen Elementen \(u_1\) und \(u_{n-1}\) behafteten Glieder beschränken. Die entwickelte Determinante hat also die Form
\[
Au_{n-1}-Bu_1=0,
\]
wo \(A\) und \(B\) Determinanten sind, welche sich nur durch eine Reihe unterscheiden. Es entsteht nämlich \(A=0\) aus der Elimination der Coefficienten von \(w\) aus den Gleichungen
\[
\begin{array}{c}
w(\xi+\delta x)=0,\\
w(\alpha^{(1)})=0,\\
w(\alpha^{(2)})=0,\\[-0.3em]
\cdots\cdots\cdots\\[-0.3em]
v=0,\\
v_1=0,\\[-0.3em]
\cdots\cdots\cdots\\[-0.3em]
v_{n-2}=0,
\end{array}
\]
während bei \(B\) nur die erste Gleichung durch
\[
w(\eta+\delta y)=0
\]
zu ersetzen ist. Daher ist \(A=0\), wenn darin die Grössen \(\xi+\delta x\) als die Variabeln betrachtet werden, eine Curve \((n-2)^{\mathrm{ter}}\) Ordnung, welche durch die Coordinaten \(\xi_i+\lambda\eta_i\), für welche
\[
w(\xi+\lambda\eta)=v+\lambda v_1+\cdots+\lambda^{n-2}v_{n-2},
\]
identisch befriedigt wird, welche also die Gerade \(\xi\eta\) als Factor enthält, und welche daher ausserdem nur noch aus der Curve \((n-3)^{\mathrm{ter}}\) Ordnung bestehen kann, welche durch die Puncte \(\alpha\) und durch die Doppel- und Rückkehrpuncte bestimmt ist. Sei die Gleichung der letztern \(\Phi(\xi+\delta x)=0\), so ist also
\[
A=M\cdot\Phi(\xi+\delta x)\cdot\Sigma\pm(\xi+\delta x)_1\xi_2\eta_3,
\]
wo \(M\) von den \(\xi+\delta x\) nicht mehr abhängt. Daher auch
\[
B=M\cdot\Phi(\eta+\delta y)\cdot\Sigma\pm(\eta_1+\delta y_1)\xi_2\eta_3.
\]
Mit Hinweglassung des Factors \(M\) und indem wir \(u_{n-1},u_1\) durch ihre Werthe ersetzen, wird also unsere Bedingungsgleichung
\[
\begin{aligned}
&\Phi(\xi+\delta x)\,[\Sigma\pm(\xi_1+\delta x_1)\xi_2\eta_3]
\left(\eta_1\frac{\partial f(\xi)}{\partial\xi_1}+\eta_2\frac{\partial f(\xi)}{\partial\xi_2}+\eta_3\frac{\partial f(\xi)}{\partial\xi_3}\right)\\
&=\Phi(\eta+\delta y)\,[\Sigma\pm(\eta_1+\delta y_1)\xi_2\eta_3]
\left(\xi_1\frac{\partial f(\eta)}{\partial\eta_1}+\xi_2\frac{\partial f(\eta)}{\partial\eta_2}+\xi_3\frac{\partial f(\eta)}{\partial\eta_3}\right),
\end{aligned}
\]
oder wenn wir jetzt den Grenzübergang ausführen:
\[
\begin{aligned}
&\Phi(\xi)\Sigma\pm\delta x_1\xi_2\eta_3
\left(\eta_1\frac{\partial f(\xi)}{\partial\xi_1}+\eta_2\frac{\partial f(\xi)}{\partial\xi_2}+\eta_3\frac{\partial f(\xi)}{\partial\xi_3}\right)\\
&=\Phi(\eta)\Sigma\pm\delta y_1\xi_2\eta_3
\left(\xi_1\frac{\partial f(\eta)}{\partial\eta_1}+\xi_2\frac{\partial f(\eta)}{\partial\eta_2}+\xi_3\frac{\partial f(\eta)}{\partial\eta_3}\right).
\end{aligned}
\]
Da nun \(\xi+\delta x\) auf der Curve \(f=0\) liegen muss, hat man
\[
\frac{\partial f(\xi)}{\partial\xi_1}\delta x_1
+\frac{\partial f(\xi)}{\partial\xi_2}\delta x_2
+\frac{\partial f(\xi)}{\partial\xi_3}\delta x_3=0;
\]
und die identische Gleichung
\[
\left|
\begin{array}{cccc}
\delta x_1&\delta x_2&\delta x_3&\displaystyle\sum\frac{\partial f(\xi)}{\partial\xi_i}\delta x_i\\
\xi_1&\xi_2&\xi_3&\displaystyle\sum\frac{\partial f(\xi)}{\partial\xi_i}\xi_i\\
\eta_1&\eta_2&\eta_3&\displaystyle\sum\frac{\partial f(\xi)}{\partial\xi_i}\eta_i\\
\zeta_1&\zeta_2&\zeta_3&\displaystyle\sum\frac{\partial f(\xi)}{\partial\xi_i}\zeta_i
\end{array}
\right|=0
\]
liefert also
\[
\Sigma\pm\delta x_1\xi_2\eta_3=
\Sigma\pm\delta x_1\xi_2\zeta_3\cdot
\frac{\displaystyle\sum\frac{\partial f(\xi)}{\partial\xi_i}\eta_i}
{\displaystyle\sum\frac{\partial f(\xi)}{\partial\xi_i}\zeta_i}.
\]
Ebenso hat man
\[
\Sigma\pm\delta y_1\xi_2\eta_3=
-\Sigma\pm\delta y_1\eta_2\zeta_3\cdot
\frac{\displaystyle\sum\frac{\partial f(\eta)}{\partial\eta_i}\xi_i}
{\displaystyle\sum\frac{\partial f(\eta)}{\partial\eta_i}\zeta_i}.
\]

% --- continuation source pp. 182--193 ---

Setzt man diese Werthe in die obige Bedingungsgleichung ein und entfernt überflüssige Factoren, so hat man
\[
0=\frac{\Phi'\xi}{\sum_i \frac{\partial f(\xi)}{\partial \xi_i}\zeta_i}
\sum\pm\delta x_1\xi_2\zeta_3
+\frac{\Phi'\eta}{\sum_i \frac{\partial f(\eta)}{\partial \eta_i}\zeta_i}
\sum\pm\delta y_1\eta_2\zeta_3 .
\]
Der zu untersuchende Logarithmus wird also
\[
\log\left(-\frac{\Phi'\eta}{\Phi'\xi}\cdot
\frac{\sum_i \frac{\partial f(\xi)}{\partial \xi_i}\zeta_i}
{\sum_i \frac{\partial f(\eta)}{\partial \eta_i}\zeta_i}\right),
\]
welcher offenbar die Differenz einer Function von \(\xi\) und einer Function von \(\eta\) ist, wie oben verlangt wurde.

So haben wir dargethan, dass unsere von nur \(p+2\) Grössen abhängige Transcendente
\[
\frac12 T_{\xi\eta}\begin{pmatrix}x^{(1)}&x^{(2)}&\cdots\\ y^{(1)}&y^{(2)}&\cdots\end{pmatrix}
\]
die wesentliche Eigenschaft noch besitzt, dass sie in einen nur von \(\eta\) und in einen nur von \(\xi\) abhängigen Theil zerfällt, wie dieses der ursprünglichen Transcendente \(T_{\xi\eta}\) mit \(2p+2\) Argumenten zukam.

Es folgt hieraus, dass, wenn man drei verschiedene Transcendenten \(T_{\xi\eta}\), \(T_{\eta\zeta}\), \(T_{\zeta\xi}\) bildet, deren obere Grenzen immer die nämlichen Puncte \(x\) sind, deren untere Grenzen aber nach der oben angegebenen Construction entstehen und bezüglich durch \(y,z,t\) bezeichnet werden mögen, die Summe aller drei bei passend gelegten Wegen verschwindet, so dass
\[
T_{\xi\eta}(x)+T_{\eta\zeta}(x)+T_{\zeta\xi}(x)=0.
\]

\subsubsection*{\S\,47. Ueber den Character der speciellen Transcendente \(T_{\xi\eta}(x)\).}

Um den Character und die Art der Bestimmtheit der neu eingeführten Transcendente völlig hervortreten zu lassen, fügen wir noch folgende Erläuterungen hinzu.

Gehen wir zunächst von den \(p\) festen Puncten \(c\) aus, welche die untern Grenzpuncte des Umkehrproblems sind, und welche nur nicht mit den Doppel- und Rückkehrpuncten auf einer Curve \((n-3)^{\mathrm{ter}}\) Ordnung liegen sollen.\footnote{Diese Voraussetzung ist hier der Kürze wegen gemacht. Es hat keine wesentliche Schwierigkeit, auch die entgegengesetzte zu behandeln.} Diese Puncte seien die obern Grenzen einer solchen Transcendente, zwei feste Puncte \(\lambda,\mu\) ihre Parameter. Sie bestimmen \(p\) andere Puncte \(b\) als untere Grenzen; ordnen wir diese den \(c\) in beliebiger Weise zu, und verbinden die zugeordneten Paare durch bestimmte Integrationswege, welche einen beliebig vorgeschriebenen Weg von \(\lambda\) nach \(\mu\) nicht schneiden, so erhalten wir eine vollkommen bestimmte Function \(T_{\lambda\mu}(c)\), welche eine Constante ist.

Wir können übrigens in dieser Constante die Zuordnung der \(c,b\) verändern, wenn wir folgende Vorsichtsmassregeln beobachten. Verbinden wir alle Puncte \(c\) in der imaginären Ebene durch bestimmte, die Parametercurve \(\lambda,\mu\) nicht schneidende Wege mit einem gegebenen Puncte \(0\), welcher ebenfalls einem bestimmten Puncte von \(f=0\) entspricht. Setzen wir dann fest, dass bei veränderter Zuordnung der \(b\) zu den \(c\) die neuen Integrationswege sich aus den ursprünglichen und aus diesen neuen Linien zusammensetzen sollen, so bleibt der Werth von \(T_{\lambda\mu}(c)\) dadurch völlig ungeändert.

Lassen wir nun den einen Parameter sich von \(\lambda\) nach \(\xi\) hin bewegen, wobei die Puncte \(c,\mu\) vollkommen fest bleiben sollen. Hierbei bewegen sich zugleich die Puncte \(b\) und gehen schliesslich in gewisse andre Puncte \(b'\) über. Das Differential der Transcendente \(T_{\xi\mu}\) ist in jeder Lage völlig bestimmt (vgl. auch \S\,51.); aber es kann vorkommen, dass der Werth von \(T_{\xi\mu}\) unendlich gross wird. Nach den Eigenschaften des Integrals dritter Gattung könnte dies dadurch eintreten, dass einer der Puncte \(\xi,\mu\) mit einem der Puncte \(c,b'\) zusammenfällt; also, da \(\mu\) und die \(c\) als constant gelten, in folgenden Fällen:
\[
\begin{array}{ll}
1) & \text{wenn }\xi\text{ mit einem der Puncte }c\text{ zusammenfällt,}\quad\hfill\\
2) & \text{wenn }\xi\text{ mit einem der Puncte }b'\text{ zusammenfällt,}\quad\hfill\\
3) & \text{wenn }\mu\text{ mit einem der Puncte }b'\text{ zusammenfällt.}
\end{array}
\]
Aber in der That begründet nur der erste Fall ein Unendlichwerden, indem die andern beiden Fälle gar nicht eintreten können. Wenn nämlich einer der Parameter, etwa \(\xi\), mit einem der \(b'\) zusammenfallen sollte, so hätte die Curve \((n-2)^{\mathrm{ter}}\) Ordnung, welche die \(b'\) bestimmt, \(n-1\) Puncte mit der Geraden \(\xi\mu\) gemein und enthielte daher diese ganz; daher müsste ein zweiter der Puncte \(b'\) in \(\mu\) fallen, und die \(p-2\) übrigen Puncte \(b'\) lägen mit den Puncten \(c\) auf einer durch die Doppel- und Rückkehrpuncte gehenden Curve \((n-3)^{\mathrm{ter}}\) Ordnung, was der über die \(c\) gemachten Voraussetzung nach unmöglich ist.

Indem man also den Punct \(\xi\) sich so bewegen lässt, dass er niemals mit einem der Puncte \(c\) zusammenfällt, bleibt die Function \(T_{\xi\mu}(c)\) jederzeit endlich und bestimmt. Da bei der ganzen Bewegung keiner der Parameter mit einem der Grenzpuncte zusammenfällt, so macht auch die in der Definition des Integrals dritter Gattung liegende Bedingung keine Schwierigkeit, dass eine gewisse die Parameter verbindende Curve von den Integrationswegen niemals geschnitten werden soll. Keine dieser Curven ist absolut fest, sondern innerhalb der für Integrationswege überhaupt gestatteten Grenzen an und für sich verschiebbar. Bewegen sich nun alle diese Curven mit \(\xi\), so kann man den Integrationswegen oder der Curve \(\mu\xi\) jederzeit solche Ausweichungen ertheilen, dass ein Schnitt nicht erfolgt. Und es ist nicht einmal nöthig, diese Ausweichungen continuirlich erfolgen zu lassen; alle diese Curven können in jedem Moment nach Bedürfniss gelegt werden, wie man will, so lange jede für sich nur aus einer Lage in die andere ohne Ueberschreitung der Unstetigkeits- oder Verzweigungspuncte geführt werden kann. Die Integrationscurven können selbst so geändert werden, dass man beide Unstetigkeitspuncte zugleich überschreitet; und die Parametercurve kann so verlegt werden, dass dabei zwei zusammengehörige Grenzpuncte gleichzeitig überschritten werden. Denn es ändert den Werth des Integrals dritter Gattung nicht, wenn man einen Integrationsweg durch einen andern so ersetzt, dass die zwischen den Parametern verlaufende Curve ganz zwischen beiden liegt, und ebenso wird dieser Werth nicht geändert, wenn man die Parametercurve so durch eine andere ersetzt, dass zwischen beiden der ganze Integrationsweg enthalten ist.

Da hier die Transcendente \(T_{\xi\mu}(c)\) nicht blos bis auf Perioden, sondern absolut bestimmt ist, so ist auch \(\frac12T_{\xi\mu}(c)\) vollständig bestimmt. Wir können ferner auf diese einfache Werthereihe zunächst die Transcendente \(T_{\xi\eta}(c)\) zurückführen, wenn wir sie durch die Gleichung
\[
(1)\ldots\quad T_{\xi\eta}(c)=T_{\xi\mu}(c)-T_{\eta\mu}(c)
\]
definiren. Wenn \(\xi\) und \(\eta\) auf bestimmten Wegen aus \(\lambda\) entstanden sind, so liegt der Werth von \(T_{\xi\eta}(c)\) aus dieser Formel vollständig fest.

Endlich kann man jede andere Transcendente \(T_{\xi\eta}(x)\) aus dieser Transcendente \(T_{\xi\eta}(c)\) und aus der allgemeinen Transcendente \(T_{\xi\eta}\binom{x}{c}\) zusammensetzen, wenn man als Definition die oben als zulässig nachgewiesene Gleichung auffasst:
\[
(2)\ldots\quad T_{\xi\eta}(x)=2T_{\xi\eta}\binom{x}{c}+T_{\xi\eta}(c).
\]
Denken wir uns in dieser Formel, um die in der Transcendente \(T_{\xi\eta}\binom{x}{c}\) liegenden Voraussetzungen zu fixiren, die \(x\) in gewisser Weise den \(c\) zugeordnet und aus diesen auf bestimmten Wegen entstanden, so muss während der Beschreibung dieser Wege nur die Parametercurve \(\xi\eta\) (welche aus den ursprünglichen Curven \(\xi\mu\) und \(\eta\mu\) zusammengesetzt ist) in solcher Weise ausweichen, dass schliesslich die von den \(x\) nach den \(c\) führenden Wege sich mit der Endlage dieser Curve nicht schneiden, wodurch denn schliesslich \(T_{\xi\eta}\binom{x}{c}\), also die rechte Seite der obigen Gleichung überhaupt eine ganz bestimmte Bedeutung erhält, mithin auch \(T_{\xi\eta}(x)\) vollkommen fixirt ist. Verschiebt man insbesondere die Parametercurve so, dass auch die Linien \(0c\) nicht von ihr geschnitten werden, so kann man die Zuordnung der \(x\) zu den \(c\) beliebig abändern, wenn man nur wieder voraussetzt, dass die neuen Wege von den \(c\) zu den \(x\) sich aus den früheren und aus den Linien \(0c\) zusammensetzen, wodurch der Werth von \(T_{\xi\eta}\binom{x}{c}\) nicht geändert wird.

Bezeichnen wir jetzt die den \(x\) entsprechenden untern Grenzpuncte durch \(y\), so entsprechen auch umgekehrt den \(y\) als obere Grenzpuncte die \(x\) als untere, und wir erhalten zwei Transcendente
\[
T_{\xi\eta}(x),\qquad T_{\xi\eta}(y),
\]
welche dieselben Grenzpuncte haben, deren Summe also ein Periodicitätsmodul eines Integrals dritter Gattung sein muss. Wir wollen untersuchen, unter welchen Voraussetzungen insbesondere diese Summe gleich Null ist.

Zu diesem Zwecke denken wir uns jetzt die Function \(T_{\lambda\mu}(c)\) selbst aus einer andern continuirlich abgeleitet. Gehen wir nämlich von einer Berührungscurve \((n-2)^{\mathrm{ter}}\) Ordnung aus, welche durch die Doppel- und Rückkehrpuncte von \(f=0\), so wie durch die \(n-2\) andern Schnittpuncte der Geraden \(\lambda\mu\) geht, und die Curve \(f=0\) ausserdem in \(p\) Puncten \(\varepsilon^{(1)},\varepsilon^{(2)},\ldots,\varepsilon^{(p)}\) berührt. Denken wir uns jetzt eine dieser Curve sehr nahe Curve gelegt, welche noch durch die Doppel- und Rückkehrpuncte, durch die \(n-2\) andern Puncte der Geraden \(\lambda\mu\) und durch Puncte \(\xi^{(1)},\xi^{(2)}\ldots\xi^{(p)}\) geht, welche den \(\varepsilon\) sehr nahe liegen, so wird dieselbe die Curve \(f=0\) noch in \(p\) Punkten \(\eta^{(1)},\eta^{(2)},\ldots,\eta^{(p)}\) schneiden, welche den Punkten \(\varepsilon\) sehr nahe liegen, und mit den \(\xi\) paarweise durch sehr kleine Wege verbunden sind. Betrachtet man jetzt die \(\varepsilon\) an Stelle der obigen Grundpuncte \(c\), die \(\eta\) an Stelle der \(b\), so ist \(T_{\lambda\mu}(\xi)\) eine unendlich kleine Grösse; aber auch \(T_{\lambda\mu}(\eta)\), wenn es aus \(T_{\lambda\mu}(\xi)\) durch die unendlich kleinen Bewegungen der \(\xi\) nach den \(\eta\) entsteht, kann nur unendlich klein sein, und die \(\eta\) können auch nur sehr kleine Wege beschrieben haben. Daher kann die Summe \(T_{\lambda\mu}(\xi)+T_{\lambda\mu}(\eta)\) kein Aggregat der endlich bleibenden Periodicitätsmoduln sein, sondern muss verschwinden.

Gehen wir nun von den \(\varepsilon\) zu den \(c\), und lassen zugleich die \(\eta\) zu den \(b\) gehen, so dass auf diese Weise die zu letzteren führenden Wege fixirt werden, so ändert sich die Summe
\[
T_{\lambda\mu}(\xi)+T_{\lambda\mu}(\eta)
\]
dabei stetig, und es kann also die Summe nie aufhören Null zu sein. Und ebenso ist es, wenn die \(c,b\) in die \(x,y\), und \(\lambda\) in \(\xi\) oder \(\eta\) übergeht u. s. w. Unter dieser Voraussetzung ist also die Summe
\[
T_{\xi\eta}(x)+T_{\xi\eta}(y)
\]
immer gleich Null.

Wir wollen noch folgende Bemerkung hinzufügen. Wenn \(\lambda\) bewegt wird, so geht die Berührungscurve der \(\varepsilon\) immer stetig in andere Berührungscurven über. Bewegen wir insbesondere den Punct \(\lambda\) auf der zwischen ihm und \(\mu\) gegebenen Verbindungslinie gegen \(\mu\) hin, so erhalten wir endlich, wenn beide sich einander unendlich genähert haben, eine Berührungscurve, welche durch die \(n-2\) übrigen Puncte der in \(\mu\) an \(f=0\) gezogenen Tangente geht, und in \(p\) Puncten \(\varepsilon\) berührt. Von dieser Curve werden wir später Gebrauch zu machen haben.

\subsubsection*{\S\,48. Theorie der adjungirten Transcendenten.}

Eine genuine Zerlegung der speciellen Transcendente \(T_{\xi\eta}(x)\) in die Differenz zweier gleichartigen Functionen, welche nur von je \(p+1\) Puncten abhängen, gründet sich auf die Theorie der \emph{adjungirten Transcendenten}.

Durch je \(p-1\) der Puncte \(x\) und durch die Doppel- und Rückkehrpuncte von \(f=0\) ist eine Curve \((n-3)^{\mathrm{ter}}\) Ordnung bestimmt, welche \(f=0\) noch in \(p-1\) andern Puncten schneidet. Setzen wir in \(T_{\xi\mu}(x)\), wie es in \S\,47 definirt wurde, diese an Stelle von \(p-1\) der Puncte \(x\), und vertauschen den übrigbleibenden Punct \(x\) mit \(\xi\), so entsteht eine Reihe von \(p\) verschiedenen Transcendenten, welche wir der gegebenen adjungirt nennen. Auf den absolut constanten Punct \(\mu\) nehmen wir hiebei keine Rücksicht.

Bezeichnen wir die mit
\[
x^{(1)},\ x^{(2)},\ldots,x^{(i-1)},\ x^{(i+1)},\ldots,x^{(p)}
\]
auf derselben Curve gelegenen \(p-1\) Puncte durch
\[
x^{1i},\ x^{2i},\ldots,x^{i-1,i},\ x^{i+1,i},\ldots,x^{pi},
\]
so haben wir folgendes System adjungirter Transcendenten: \footnote{Es mag als interessant hier bemerkt werden, dass die untern Grenzpuncte aller adjungirten Transcendenten die nämlichen sind, wie man leicht aus dem Abel'schen Theorem geometrisch nachweist.}
\[
\begin{aligned}
(1)\ldots\quad T^{(1)}&=T_{x^{(1)}\mu}(\xi,x^{21}\ldots x^{p1}),\\
T^{(2)}&=T_{x^{(2)}\mu}(x^{12},\xi\ldots x^{p2}),\\[-0.35em]
&\cdots\\[-0.35em]
T^{(p)}&=T_{x^{(p)}\mu}(x^{1p},x^{2p}\ldots \xi).
\end{aligned}
\]
Um diese Transcendenten vollkommen genau im Sinne des \S\,47 zu definiren, müssen wir die Wege kennen, durch welche hier die \(x^{(i)}\) mit \(\mu\), sowie die Puncte \(\xi,x^{ik}\) mit den constanten Puncten \(c\) verbunden sind. Um dies ausführen zu können, verbinden wir den Punct \(0\) (\S\,47) durch eine beliebig gewählte Linie auch noch mit dem Puncte \(\lambda\). Die bestimmte Definition von \(T_{\xi\eta}(x)\) erfordert, dass die Entstehung von \(\xi\) aus \(\lambda\) gegeben sei. Wir nehmen nun an, dass in den adjungirten Transcendenten \(x^{(i)}\) mit \(\lambda\) verbunden werde durch eine Curve, die aus dem Wege \(\lambda\xi0c^{(i)}x^{(i)}\) durch erlaubte Verschiebungen entsteht, und dass ebenso \(\xi\) mit \(c^{(i)}\) durch den Weg \(\xi0c^{(i)}\) oder einen aus ihm durch erlaubte Verschiebungen entstandenen verbunden werde.

Zur Bestimmung der Wege, auf denen die Puncte \(x^{ik}\) entstanden sein sollen, führt nun folgende Bemerkung. Wenn man die zu den \(x^{(i)}\) gehörigen Puncte \(y^{(i)}\) auf dem in \S\,47 angegebenen Wege entstehen lässt, so sind sie durch bestimmte Wege mit \(0\) verbunden. Wählt man nun an Stelle von \(x^{(1)}\) einen Punct, welcher dem Puncte \(\xi\) beliebig nahe ist, und macht übrigens denselben Process durch, so dass die übrigen Puncte sich von den \(c\) zu den übrigen \(x\), dieser eine aber von \(c^{(1)}\) gegen \(\xi\) auf dem durch \(0\) gehenden vorgeschriebenen Wege bewegt, so fällt schliesslich einer der Puncte \(y\) mit \(\mu\) zusammen, da die durch die \(n-2\) übrigen Puncte von \(\xi\mu\) zur Bestimmung der \(y\) gelegte Curve \((n-2)^{\mathrm{ter}}\) Ordnung sich in eine Curve \((n-3)^{\mathrm{ter}}\) Ordnung und in die Gerade \(\xi\mu\) selbst auflöst. Die übrigen Puncte \(y\) sind dann in beliebiger Folge die Puncte \(x^{21},x^{31}\ldots x^{p1}\), und zwar sind sie mit den \(c\) durch bestimmte Curven verbunden. Diese sollen bei der ersten adjungirten Transcendente, und die entsprechend gebildeten bei den übrigen, zur Entstehung der obern Grenzpuncte aus den \(c\) benutzt werden.

Man sieht hieraus, dass mit Hinzufügung der einen Linie \(0\xi\) alle adjungirten Transcendenten durch die gegebene Transcendente vollständig und eindeutig bestimmt sind.

Die \(p\) adjungirten Transcendenten wollen wir nun mit Hülfe des Abel'schen Theorems so transformiren, dass die in ihnen vorkommenden Grenzpuncte mit den in \(T\) benutzten möglichst übereinstimmen.

Zu diesem Zwecke stellen wir zunächst ein diesen \(p\) Transcendenten ähnliches System von Transcendenten auf, in denen die Puncte \(x\) durch beliebige constante Puncte \(\alpha^{(1)},\alpha^{(2)}\ldots\alpha^{(p)}\) ersetzt sind, in denen aber zugleich der oben als Argument auftretende Punct \(\xi\) beziehungsweise durch \(\alpha^{(1)},\alpha^{(2)}\ldots\alpha^{(p)}\) ersetzt ist. Dieses System (welches keineswegs ein System adjungirter Transcendenten ist) sei, bei analoger Bezeichnung:
\[
\begin{aligned}
(2)\ldots\quad X^{(1)}&=T_{x^{(1)}\mu}(\alpha^{(1)},\alpha^{21}\ldots\alpha^{p1}),\\
X^{(2)}&=T_{x^{(2)}\mu}(\alpha^{12},\alpha^{(2)}\ldots\alpha^{p2}),\\[-0.35em]
&\cdots\\[-0.35em]
X^{(p)}&=T_{x^{(p)}\mu}(\alpha^{1p},\alpha^{2p}\ldots\alpha^{(p)}).
\end{aligned}
\]
Es hat die Eigenschaft, dass jede der Functionen \(X\) nur von \emph{einem} variabeln Punct \(x\) abhängt, welcher in dem betreffenden \(X\) sowohl als Parameter, wie auch in den untern Grenzen auftritt. Die Puncte \(\alpha^{(1)},\alpha^{(2)}\ldots\alpha^{(p)}\) werden dabei durch beliebige, die Linien \(x^{(i)}\mu\) nicht schneidende Linien mit den Puncten \(c\) verbunden gedacht; die \(\alpha^{ik}\) aber und die zu denselben führenden Wege werden aus ihnen ebenso abgeleitet wie oben die \(x^{ik}\) aus den \(x^{(i)}\).

Da in den Transcendenten \(T^{(1)},T^{(2)}\ldots T^{(p)}\) und \(X^{(1)},X^{(2)}\ldots X^{(p)}\) beziehungsweise dieselben Parameter auftreten, so bildet stets die halbe Differenz entsprechender eine allgemeine Transcendente \(T\), und zwar wird:
\begingroup\footnotesize
\[
\left\{
\begin{aligned}
\frac12(T^{(1)}-X^{(1)})
&=T_{x^{(1)}\mu}\begin{pmatrix}\xi,x^{21}\ldots x^{p1}\\ \alpha^{(1)},\alpha^{21}\ldots\alpha^{p1}\end{pmatrix} \\
&=\int_{\alpha^{(1)}}^{\xi}d\Pi_{x^{(1)}\mu}
+\int_{\alpha^{21}}^{x^{21}}d\Pi_{x^{(1)}\mu}\cdots
+\int_{\alpha^{p1}}^{x^{p1}}d\Pi_{x^{(1)}\mu},\\[0.35em]
\frac12(T^{(2)}-X^{(2)})
&=T_{x^{(2)}\mu}\begin{pmatrix}x^{12},\xi\ldots x^{p2}\\ \alpha^{12},\alpha^{(2)}\ldots\alpha^{p2}\end{pmatrix} \\
&=\int_{\alpha^{12}}^{x^{12}}d\Pi_{x^{(2)}\mu}
+\int_{\alpha^{(2)}}^{\xi}d\Pi_{x^{(2)}\mu}\cdots
+\int_{\alpha^{p2}}^{x^{p2}}d\Pi_{x^{(2)}\mu},\\[-0.2em]
&\cdots\\[-0.2em]
\frac12(T^{(p)}-X^{(p)})
&=T_{x^{(p)}\mu}\begin{pmatrix}x^{1p},x^{2p}\ldots\xi\\ \alpha^{1p},\alpha^{2p}\ldots\alpha^{(p)}\end{pmatrix} \\
&=\int_{\alpha^{1p}}^{x^{1p}}d\Pi_{x^{(p)}\mu}
+\int_{\alpha^{2p}}^{x^{2p}}d\Pi_{x^{(p)}\mu}\cdots
+\int_{\alpha^{(p)}}^{\xi}d\Pi_{x^{(p)}\mu}.
\end{aligned}
\right.
\tag{3}
\]
\endgroup
Nach dem Abel'schen Theorem ist nun
\begingroup\footnotesize
\[
\left\{
\begin{aligned}
&\int_{\alpha^{1i}}^{x^{1i}}d\Pi_{x^{(i)}\mu}
+\int_{\alpha^{2i}}^{x^{2i}}d\Pi_{x^{(i)}\mu}\cdots
+\int_{\alpha^{i-1,i}}^{x^{i-1,i}}d\Pi_{x^{(i)}\mu} \\
&\quad +\int_{\alpha^{i+1,i}}^{x^{i+1,i}}d\Pi_{x^{(i)}\mu}\cdots
+\int_{\alpha^{pi}}^{x^{pi}}d\Pi_{x^{(i)}\mu} \\
&= -\left\{
\int_{\alpha^{(1)}}^{x^{(1)}}d\Pi_{x^{(i)}\mu}
+\int_{\alpha^{(2)}}^{x^{(2)}}d\Pi_{x^{(i)}\mu}\cdots
+\int_{\alpha^{(i-1)}}^{x^{(i-1)}}d\Pi_{x^{(i)}\mu}
\right.\\
&\qquad\left.
+\int_{\alpha^{(i+1)}}^{x^{(i+1)}}d\Pi_{x^{(i)}\mu}\cdots
+\int_{\alpha^{(p)}}^{x^{(p)}}d\Pi_{x^{(i)}\mu}\right\}
+\log\frac{\varphi(\mu)\psi(x^{(i)})}{\psi(\mu)\varphi(x^{(i)})},
\end{aligned}
\right.
\tag{4}
\]
\endgroup
wo \(\varphi=0\) die Curve \((n-3)^{\mathrm{ter}}\) Ordnung ist, welche durch \(x^{(1)},x^{(2)}\ldots x^{(i-1)},x^{(i+1)}\ldots x^{(p)}\) gelegt ist und die Puncte \(x^{1i},x^{2i}\ldots x^{i-1,i},x^{i+1,i}\ldots x^{pi}\) bestimmt, \(\psi=0\) aber die Curve \((n-3)^{\mathrm{ter}}\) Ordnung, welche durch \(\alpha^{(1)},\alpha^{(2)}\ldots\alpha^{(i-1)},\alpha^{(i+1)}\ldots\alpha^{(p)}\) gelegt ist und die Puncte \(\alpha^{1i},\alpha^{2i}\ldots\alpha^{i-1,i},\alpha^{i+1,i}\ldots\alpha^{pi}\) bestimmt.

Bezeichnet man nun insbesondere durch
\[
\varphi(y,y^{(1)},y^{(2)}\ldots y^{(p-1)})=0
\]
die Gleichung einer Curve \((n-3)^{\mathrm{ter}}\) Ordnung, welche durch die Doppel- und Rückkehrpuncte und durch \(y^{(1)},y^{(2)}\ldots y^{(p-1)}\) geht, während \(y\) als variabler Punct betrachtet wird. Der Ausdruck \(\varphi\) ist eine Determinante, welche durch Vertauschung irgend zweier der Puncte \(y\) nur das Zeichen ändert, und man kann obige Gleichung daher ebensowohl als Gleichung einer durch \(y^{(1)},y^{(2)}\ldots y^{(i-1)},y^{(i+1)}\ldots y^{(p-1)}\) gelegten Curve betrachten, in welcher \(y^{(i)}\) der variable Punct ist.

Mit Benutzung dieser Bezeichnungsweise wird die unter dem Logarithmuszeichen befindliche Function in (4):
\begingroup\scriptsize
\[
\frac{\varphi(x^{(1)},x^{(2)}\ldots x^{(i-1)},\mu,x^{(i+1)},\ldots,x^{(p)})\,
\varphi(\alpha^{(1)},\alpha^{(2)}\ldots\alpha^{(i-1)},x^{(i)},\alpha^{(i+1)},\ldots,\alpha^{(p)})}
{\varphi(\alpha^{(1)},\alpha^{(2)}\ldots\alpha^{(i-1)},\mu,\alpha^{(i+1)},\ldots,\alpha^{(p)})\,
\varphi(x^{(1)},x^{(2)}\ldots x^{(i-1)},x^{(i)},x^{(i+1)},\ldots,x^{(p)})}.
\]
\endgroup
Daher kann man die Gleichungen (3) durch folgende ersetzen:
\begingroup\scriptsize
\[
\begin{aligned}
\frac12(T^{(1)}-X^{(1)})={}&\int_{\alpha^{(1)}}^{\xi}d\Pi_{x^{(1)}\mu}
-\int_{\alpha^{(2)}}^{x^{(2)}}d\Pi_{x^{(1)}\mu}\cdots
-\int_{\alpha^{(p)}}^{x^{(p)}}d\Pi_{x^{(1)}\mu} \\
&+\log\frac{\varphi(\mu,x^{(2)}\ldots x^{(p)})\,\varphi(x^{(1)},\alpha^{(2)}\ldots\alpha^{(p)})}
{\varphi(\mu,\alpha^{(2)}\ldots\alpha^{(p)})\,\varphi(x^{(1)},x^{(2)}\ldots x^{(p)})},\\[0.4em]
\frac12(T^{(2)}-X^{(2)})={}&-\int_{\alpha^{(1)}}^{x^{(1)}}d\Pi_{x^{(2)}\mu}
+\int_{\alpha^{(2)}}^{\xi}d\Pi_{x^{(2)}\mu}\cdots
-\int_{\alpha^{(p)}}^{x^{(p)}}d\Pi_{x^{(2)}\mu} \\
&+\log\frac{\varphi(x^{(1)},\mu\ldots x^{(p)})\,\varphi(\alpha^{(1)},x^{(2)}\ldots\alpha^{(p)})}
{\varphi(\alpha^{(1)},\mu\ldots\alpha^{(p)})\,\varphi(x^{(1)},x^{(2)}\ldots x^{(p)})},\\[-0.2em]
&\cdots\\[-0.2em]
\frac12(T^{(p)}-X^{(p)})={}&-\int_{\alpha^{(1)}}^{x^{(1)}}d\Pi_{x^{(p)}\mu}
-\int_{\alpha^{(2)}}^{x^{(2)}}d\Pi_{x^{(p)}\mu}\cdots
+\int_{\alpha^{(p)}}^{\xi}d\Pi_{x^{(p)}\mu} \\
&+\log\frac{\varphi(x^{(1)},x^{(2)}\ldots\mu)\,\varphi(\alpha^{(1)},\alpha^{(2)}\ldots x^{(p)})}
{\varphi(\alpha^{(1)},\alpha^{(2)}\ldots\mu)\,\varphi(x^{(1)},x^{(2)}\ldots x^{(p)})}.
\end{aligned}
\]
\endgroup
Fügt man diesen Gleichungen die Definition von \(T\) hinzu:
\[
\frac12T=\int_{\alpha^{(1)}}^{x^{(1)}}d\Pi_{\xi\mu}
+\int_{\alpha^{(2)}}^{x^{(2)}}d\Pi_{\xi\mu}\cdots
+\int_{\alpha^{(p)}}^{x^{(p)}}d\Pi_{\xi\mu}+\Xi,
\]
wo \(\Xi\) die den Transcendenten \(X^{(1)},X^{(2)}\ldots X^{(p)}\) analog gebildete Function
\[
\Xi=T_{\xi\mu}(\alpha^{(1)},\alpha^{(2)},\ldots,\alpha^{(p)})
\]
bedeutet, so sieht man, dass in dem ganzen System von Transcendenten nur die Grenzen \(\alpha,x,\xi\) bei Integralen dritter Gattung auftreten, abgesehen von den Functionen \(X^{(1)},X^{(2)}\ldots X^{(p)}\), welche nur je von einem der variablen Puncte abhängig sind.

Untersuchen wir nun diejenigen Terme in der Differenz \(\frac12(T-T^{(i)})\), welche zugleich von \(\xi\) und von \(x^{(i)}\) abhängen. Diese Terme können nur aus der Differenz
\[
\int_{\alpha^{(i)}}^{x^{(i)}}d\Pi_{\xi\mu}-\int_{\alpha^{(i)}}^{\xi}d\Pi_{x^{(i)}\mu}
\]
entspringen; aber diese Differenz wird gleich
\begingroup\footnotesize
\[
\int_{\alpha^{(i)}}^{\delta}d\Pi_{\xi\varepsilon}
+\int_{\delta}^{x^{(i)}}d\Pi_{\xi\varepsilon}
+\int_{\alpha^{(i)}}^{x^{(i)}}d\Pi_{\varepsilon\mu}
-\int_{\alpha^{(i)}}^{\varepsilon}d\Pi_{x^{(i)}\delta}
-\int_{\varepsilon}^{\xi}d\Pi_{x^{(i)}\delta}
-\int_{\alpha^{(i)}}^{\xi}d\Pi_{\delta\mu},
\]
\endgroup
wo die von beiden Puncten \(\xi,x^{(i)}\) zugleich abhängigen Terme sich aufheben. Solche Terme existiren also überhaupt nicht, und es ist daher
\[
\frac{\partial}{\partial x^{(i)}}\left(\frac{\partial T}{\partial\xi}\right)
=\frac{\partial}{\partial\xi}\left(\frac{\partial T^{(i)}}{\partial x^{(i)}}\right).
\]
Aber ganz ebenso enthält die Differenz
\[
\frac12(T^{(i)}-T^{(k)})
\]
keinen von \(x^{(i)}\) und \(x^{(k)}\) zugleich abhängigen Term. Daher ist auch:
\[
\frac{\partial}{\partial x^{(i)}}\left(\frac{\partial T^{(k)}}{\partial x^{(k)}}\right)
=\frac{\partial}{\partial x^{(k)}}\left(\frac{\partial T^{(i)}}{\partial x^{(i)}}\right).
\]
Diese Gleichungen zeigen nun, dass folgendes merkwürdige und für alle folgenden Untersuchungen fundamentale Theorem stattfindet:

Die Differentialquotienten der adjungirten Transcendenten
\[
\frac12T,\quad \frac12T^{(1)},\quad \frac12T^{(2)},\ldots,\frac12T^{(p)}
\]
nach den ihnen entsprechenden Variabeln
\[
\xi,\quad x^{(1)},\quad x^{(2)},\ldots,x^{(p)}
\]
sind die Coefficienten eines vollständigen Differentials:
\[
dU=\frac12\left(\frac{\partial T}{\partial\xi}d\xi
+\frac{\partial T^{(1)}}{\partial x^{(1)}}dx^{(1)}\cdots
+\frac{\partial T^{(p)}}{\partial x^{(p)}}dx^{(p)}\right).
\]
Es verdient übrigens bemerkt zu werden, dass diese Differentialquotienten den Punct \(\mu\) in keiner Weise mehr enthalten.

\subsubsection*{\S\,49. Die Function \(U\). Die Transcendente \(T_{\xi\eta}\) durch sie ausgedrückt.}

Wir wollen die Function \(U\), welche durch obige Gleichung bis auf eine Constante bestimmt ist, vollständig durch die Bestimmung definiren, dass \(U\) verschwinden soll, wenn \(\xi,x^{(1)},x^{(2)}\ldots x^{(p)}\) mit einem beliebig gegebenen constanten Punctsystem \(\zeta,\gamma^{(1)},\gamma^{(2)}\ldots\gamma^{(p)}\) zusammenfallen.

Denken wir uns also auch das Punctsystem \(\zeta,\gamma\), wie oben das Punctsystem \(\xi,x\) mit den \(\lambda,c\) durch bestimmte gegebene Curven verbunden.

Indem wir die Differentialquotienten der \(T\) dann in den Ausdruck für \(U\) einführen und zwischen den Werthsystemen \(\zeta,\gamma^{(1)},\gamma^{(2)}\ldots\gamma^{(p)}\) und \(\xi,x^{(1)},x^{(2)}\ldots x^{(p)}\) auf den durch die graphische Darstellung vorgeschriebenen, oder doch durch erlaubte Verschiebungen aus diesen abgeleiteten Wegen integriren, erhalten wir für \(U\) folgenden völlig bestimmten Ausdruck:
\begingroup\scriptsize
\[
\begin{aligned}
U={}&\sum_{i=1}^{i=p}\int_{\zeta,\gamma^{(i)}}^{\xi,x^{(i)}}
\left(\frac{\partial\Pi_{x^{(i)}\alpha^{(i)}}(\xi)}{\partial\xi}d\xi
+\frac{\partial\Pi_{\xi\alpha^{(i)}}(x^{(i)})}{\partial x^{(i)}}dx^{(i)}\right)\\
&-\frac12\sum\sum_{i\le k}\int_{\gamma^{(i)},\gamma^{(k)}}^{x^{(i)},x^{(k)}}
\left(\frac{\partial\Pi_{x^{(i)}\alpha^{(i)}}(x^{(k)})}{\partial x^{(k)}}dx^{(k)}
+\frac{\partial\Pi_{x^{(k)}\alpha^{(k)}}(x^{(i)})}{\partial x^{(i)}}dx^{(i)}\right)\\
&+\frac12\int_{\zeta}^{\xi}\frac{\partial\Xi}{\partial\xi}d\xi
+\frac12\sum_{i=1}^{i=p}\int_{\gamma^{(i)}}^{x^{(i)}}\frac{\partial X^{(i)}}{\partial x^{(i)}}dx^{(i)}\\
&-\log\frac{\varphi(x^{(1)},x^{(2)}\ldots x^{(p)})\,\varphi(\gamma^{(1)},\alpha^{(2)}\ldots\alpha^{(p)})\,\varphi(\alpha^{(1)},\gamma^{(2)}\ldots\alpha^{(p)})\cdots\varphi(\alpha^{(1)},\alpha^{(2)}\ldots\gamma^{(p)})}
{\varphi(\gamma^{(1)},\gamma^{(2)}\ldots\gamma^{(p)})\,\varphi(x^{(1)},\alpha^{(2)}\ldots\alpha^{(p)})\,\varphi(\alpha^{(1)},x^{(2)}\ldots\alpha^{(p)})\cdots\varphi(\alpha^{(1)},\alpha^{(2)}\ldots x^{(p)})}.
\end{aligned}
\]
\endgroup
Die zweite Summe ist durch 2 dividirt, damit \(i\) und \(k\) unabhängig von einander alle Werthepaare, die gleichen ausgeschlossen, durchlaufen können, wobei denn immer zwei gleiche Glieder auftreten. Die ersten Theile dieses Ausdrucks führt man leicht auf Integrale dritter Gattung zurück. Denn es wird nach bekannten Regeln
\begingroup\footnotesize
\[
\begin{aligned}
&\int_{\zeta,\gamma^{(i)}}^{\xi,x^{(i)}}
\left(\frac{\partial\Pi_{x^{(i)}\alpha^{(i)}}(\xi)}{\partial\xi}d\xi
+\frac{\partial\Pi_{\xi\alpha^{(i)}}(x^{(i)})}{\partial x^{(i)}}dx^{(i)}\right)\\
&\quad=\int_{\zeta}^{\xi}\frac{\partial\Pi_{x^{(i)}\alpha^{(i)}}(\xi)}{\partial\xi}d\xi
+\int_{\gamma^{(i)}}^{x^{(i)}}\frac{\partial\Pi_{\xi\alpha^{(i)}}(x^{(i)})}{\partial x^{(i)}}dx^{(i)},
\end{aligned}
\]
\endgroup
oder nach der früher immer angewandten Bezeichnung:
\[
=\int_{\zeta}^{\xi}d\Pi_{x^{(i)}\alpha^{(i)}}+
\int_{\gamma^{(i)}}^{x^{(i)}}d\Pi_{\xi\alpha^{(i)}}.
\]
Ganz ebenso hat man
\begingroup\footnotesize
\[
\begin{aligned}
&\int_{\gamma^{(i)},\gamma^{(k)}}^{x^{(i)},x^{(k)}}
\left(\frac{\partial\Pi_{x^{(i)}\alpha^{(i)}}(x^{(k)})}{\partial x^{(k)}}dx^{(k)}
+\frac{\partial\Pi_{x^{(k)}\alpha^{(k)}}(x^{(i)})}{\partial x^{(i)}}dx^{(i)}\right)\\
&\quad=\int_{\gamma^{(k)}}^{x^{(k)}}d\Pi_{x^{(i)}\alpha^{(i)}}+
\int_{\gamma^{(i)}}^{x^{(i)}}d\Pi_{\gamma^{(k)}\alpha^{(k)}}.
\end{aligned}
\]
\endgroup
Es ist ferner nach dem am Schlusse des \S\,46 angeführten Satze:
\[
\int_{\zeta}^{\xi}\frac{\partial\Xi}{\partial\xi}d\xi
=T_{\xi\zeta}(\alpha^{(1)},\alpha^{(2)}\ldots\alpha^{(p)}),
\]
\[
\int_{\gamma^{(i)}}^{x^{(i)}}\frac{\partial X^{(i)}}{\partial x^{(i)}}dx^{(i)}
=T_{x^{(i)}\gamma^{(i)}}(\alpha^{1i},\alpha^{2i}\ldots\alpha^{i-1,i},\alpha^{(i)},\alpha^{i+1,i}\ldots\alpha^{pi}).
\]
Daher wird endlich der Ausdruck von \(U\) folgender:
\begingroup\scriptsize
\[
\begin{aligned}
U={}&\sum_{i=1}^{i=p}\int_{\zeta}^{\xi}d\Pi_{x^{(i)}\alpha^{(i)}}
+\sum_{i=1}^{i=p}\int_{\gamma^{(i)}}^{x^{(i)}}d\Pi_{\xi\alpha^{(i)}}\\
&-\frac12\sum\sum_{i\le k}\int_{\gamma^{(k)}}^{x^{(k)}}d\Pi_{x^{(i)}\alpha^{(i)}}
-\frac12\sum\sum_{i\le k}\int_{\gamma^{(i)}}^{x^{(i)}}d\Pi_{\gamma^{(k)}\alpha^{(k)}}\\
&+\frac12T_{\xi\zeta}(\alpha^{(1)},\alpha^{(2)}\ldots\alpha^{(p)})
+\frac12\sum_{i=1}^{i=p}T_{x^{(i)}\gamma^{(i)}}(\alpha^{1i},\alpha^{2i}\ldots\alpha^{i-1,i},\alpha^{(i)},\alpha^{i+1,i}\ldots\alpha^{pi})\\
&-\log\frac{\varphi(x^{(1)},x^{(2)}\ldots x^{(p)})\,\varphi(\gamma^{(1)},\alpha^{(2)}\ldots\alpha^{(p)})\,\varphi(\alpha^{(1)},\gamma^{(2)}\ldots\alpha^{(p)})\cdots\varphi(\alpha^{(1)},\alpha^{(2)}\ldots\gamma^{(p)})}
{\varphi(\gamma^{(1)},\gamma^{(2)}\ldots\gamma^{(p)})\,\varphi(x^{(1)},\alpha^{(2)}\ldots\alpha^{(p)})\,\varphi(\alpha^{(1)},x^{(2)}\ldots\alpha^{(p)})\cdots\varphi(\alpha^{(1)},\alpha^{(2)}\ldots x^{(p)})}.
\end{aligned}
\]
\endgroup
Dass in diesem Ausdrucke die \(\xi,\gamma\) nur auf eine additive Constante von Einfluss sind, folgt aus der Bildungsweise des Ausdrucks. Aber es ist bemerkenswerth, dass die constanten Puncte \(\alpha^{(i)}\), also auch die von diesen abhängigen constanten Puncte \(\alpha^{ik}\) auf den Werth von \(U\) gar keinen Einfluss haben, also nur formal darin vorkommen. In

der That, bilden wir eine zweite Function \(U'\), welche sich von der vorigen dadurch unterscheidet, dass die \(\alpha\) durch andere Puncte \(\beta\) ersetzt wird, so findet man \(U-U'=0\). Dieses folgt auch unmittelbar daraus, dass schon die conjugirten Transcendenten \(T,T^{(1)},\ldots,T^{(p)}\) ihrer Entstehung nach von den \(\alpha\) unabhängig waren, daher auch \(dU\) und mithin \(U\) selbst.

Da \(\mu\) ganz aus dem Ausdrucke von \(U\) verschwunden ist, so sieht man, dass \(U\) die Form haben muss:
\[
U=\Phi(\xi,x^{(1)},x^{(2)},\ldots,x^{(p)})-\Phi(\zeta,\gamma^{(1)},\gamma^{(2)},\ldots,\gamma^{(p)}),
\]
d. h. dass \(U\) bis auf eine additive Constante ausschliesslich eine Function der \(p+1\) Puncte \(\xi,x^{(1)},x^{(2)},\ldots x^{(p)}\) ist. Bezeichnen wir sie insofern durch
\[
U(x^{(1)},x^{(2)},\ldots,x^{(p)};\xi).
\]
Wir werden uns dieser Function jetzt bedienen, um die specielle und die allgemeine Transcendente \(T\) durch sie darzustellen.

Es ist nämlich die Differenz
\[
U(x^{(1)},x^{(2)}\ldots x^{(p)};\xi)-U(x^{(1)},x^{(2)}\ldots x^{(p)};\mu)
\]
nichts anderes als das Integral von \(dU\), genommen zwischen den Werthsystemen:
\[
x^{(1)},x^{(2)}\ldots\mu\quad \text{bis}\quad x^{(1)},x^{(2)}\ldots\xi.
\]
Und zwar wollen wir \(\mu\xi\) durch die bereits oben erwähnte Curve verbinden, und längs dieser integriren. Daher ist mit Rücksicht auf den Ausdruck von \(dU\):
\[
\begin{aligned}
&U(x^{(1)},x^{(2)}\ldots x^{(p)};\xi)-U(x^{(1)},x^{(2)}\ldots x^{(p)};\mu)\\
&\quad=\frac12\int_{\mu}^{\xi}\frac{\partial T}{\partial\xi}\,d\xi\\
&\quad=\frac12T_{\xi\mu}(x^{(1)},x^{(2)}\ldots x^{(p)})
-\frac12T_{\mu\mu}(x^{(1)},x^{(2)}\ldots x^{(p)}),
\end{aligned}
\]
wo im ersten Gliede rechts jene Curve die Parameterpuncte verbindet, im zweiten Gliede aber die Verbindungslinie unendlich klein wird. Daher verschwindet die so definirte Transcendente \(T_{\mu\mu}\) und man hat die Formel:
\[
\frac12T_{\xi\mu}(x^{(1)},x^{(2)}\ldots x^{(p)})=U(x^{(1)},x^{(2)}\ldots x^{(p)};\xi)-U(x^{(1)},x^{(2)}\ldots x^{(p)};\mu);
\]
daher auch, nach der Definition von \(T_{\xi\eta}(x)\):
\[
\frac12T_{\xi\eta}(x^{(1)},x^{(2)}\ldots x^{(p)})=U(x^{(1)},x^{(2)}\ldots x^{(p)};\xi)-U(x^{(1)},x^{(2)}\ldots x^{(p)};\eta).
\]
Durch die Function \(U\) ist also die genuine Zerlegung von \(\frac12T_{\xi\eta}(x)\) in zwei Functionen mit je \(p+1\) Argumenten geleistet. Man hat sofort auch für die allgemeine Transcendente \(T_{\xi\eta}\binom{x}{z}\) den Ausdruck:
\[
T_{\xi\eta}\binom{x^{(1)},x^{(2)}\ldots x^{(p)}}{z^{(1)},z^{(2)}\ldots z^{(p)}}
=U(x^{(1)},\ldots;\xi)-U(x^{(1)},\ldots;\eta)-U(z^{(1)},\ldots;\xi)+U(z^{(1)},\ldots;\eta).
\]

\subsubsection*{\S\,50. Unstetigkeiten der Function \(U\). Die Function \(V=e^{-U}\).}

Wir wollen jetzt die Unstetigkeiten von \(U\) untersuchen. Diese Function ist aus Integralen dritter Gattung, aus Transcendenten \(\frac12T_{\xi\eta}(x)\) und aus Logarithmen algebraischer Functionen zusammengesetzt. Was die Unstetigkeiten der Transcendente \(\frac12T_{\xi\eta}(x)\) betrifft, so genügt es, die von \(T_{\xi\mu}(c)\) zu kennen, wo die \(c,\mu\) constante Puncte sind, deren erstere mit den Doppel- und Rückkehrpuncten nicht auf einer Curve \((n-3)^{\mathrm{ter}}\) Ordnung liegen. Denn man kann nach \S\,45 die Transcendente \(T_{\xi\eta}(x)\) immer durch Integrale dritter Gattung mit unabhängigen Grenzen und durch \(T_{\xi\mu}(c)\) ausdrücken; die Unstetigkeiten der letzteren Function aber sind in \S\,47 dargestellt worden.

Betrachten wir den Ausdruck \(U\), wie er in dem vorigen Paragraph explicite gegeben vorliegt, so müssen wir zunächst beachten, dass die \(\alpha\) in Wirklichkeit den Werth von \(U\) gar nicht afficiren, \(\xi\) und die \(\gamma\) nur eine demselben hinzugefügte additive Constante. Deswegen braucht man auf diejenigen Fälle keine Rücksicht zu nehmen, in welchen ein Glied von \(U\) unstetig wird, indem \(\xi\) oder eines der \(x\) sich einem dieser Puncte nähert; es existirt dann nothwendig immer ein anderes Glied, welches diese Unstetigkeit aufhebt. Daher bleiben folgende Fälle übrig:

1) Der Punct \(\xi\) fällt mit einem der Puncte \(x\), also etwa mit \(x^{(i)}\) zusammen. In solchem Falle verhält sich \(U\) wie das Glied
\[
\int_{\zeta}^{\xi}d\Pi_{x^{(i)}\alpha^{(i)}},
\]
also wie \(-\log\Sigma\pm\xi_1x_2^{(i)}\alpha_3^{(i)}\).

2) Ein Punct \(x^{(i)}\) fällt mit einem Puncte \(x^{(k)}\) zusammen. Dann verhält sich \(U\) wie
\[
-\int_{\gamma^{(k)}}^{x^{(k)}}d\Pi_{x^{(i)}\alpha^{(i)}}-\log\varphi(x^{(1)},x^{(2)}\ldots x^{(p)}),
\]
also wie \(\log\Sigma\pm x_1^{(k)}x_2^{(i)}\alpha_3^{(i)}-\log\varphi(x^{(1)},x^{(2)}\ldots x^{(p)})\); was als Logarithmus des Quotienten zweier einfach verschwindenden Functionen endlich bleibt.

3) Ebenso ist es, wenn etwa \(x^{(i)}\) mit einem der Doppel- oder Rückkehrpuncte von \(f=0\) zusammenfällt. In diesem Falle verhält sich \(U\) wie
\[
-\log\frac{\varphi(x^{(1)},x^{(2)}\ldots x^{(p)})}{\varphi(\alpha^{(1)},\ldots x^{(i)}\ldots\alpha^{(p)})},
\]
bleibt also endlich.

4) Dagegen wird \(U\) wirklich noch unendlich gross, wenn die \(p\) Puncte \(x\) auf einer Curve \((n-3)^{\mathrm{ter}}\) Ordnung
\[
\varphi(x^{(1)},x^{(2)}\ldots x^{(p)})=0
\]
liegen, welche durch die Doppel- und Rückkehrpuncte geht, und zwar verhält sich dabei \(U\) wie
\[
-\log\varphi(x^{(1)},x^{(2)}\ldots x^{(p)}).
\]
Man kann also folgenden Satz aussprechen:

Die Function \(U(x^{(1)},x^{(2)}\ldots x^{(p)};\xi)\) wird nur unendlich, wenn \(\xi\) mit einem der \(x\) zusammenfällt, oder die letztern auf einer durch die Doppel- und Rückkehrpuncte von \(f=0\) gelegten Curve \((n-3)^{\mathrm{ter}}\) Ordnung liegen; und zwar verhält sich in allen diesen Fällen \(U\) wie der negative Logarithmus einer verschwindenden Grösse.

Es liegt nun nahe, statt der Function \(U\) selbst die Function
\[
V(x^{(1)},x^{(2)}\ldots x^{(p)};\xi)=e^{-U(x^{(1)},x^{(2)}\ldots x^{(p)};\xi)}
\]
zu untersuchen. Die Function enthält die Constanten \(\zeta,\gamma\) nur in einem Factor. Ihre wesentlichste Eigenschaft ist diese, welche aus dem Schlusssatz des Vorigen sich unmittelbar ergiebt:

Die Function \(V=e^{-U}\) ist eine Function der Puncte \(x^{(1)},x^{(2)}\ldots x^{(p)},\xi\), welche niemals (ausser etwa in Folge unendlich oft wiederholter Integrationswege) unendlich gross wird; und welche nur verschwindet, wenn \(\xi\) mit einem der Puncte \(x\) zusammenfällt, oder wenn die Puncte \(x\) mit den Doppel- und Rückkehrpuncten auf einer Curve \((n-3)^{\mathrm{ter}}\) Ordnung liegen.

Mit Hülfe dieser Function erhält man für die specielle Transcendente \(T_{\xi\eta}(x)\) den Ausdruck
\[
\frac12T_{\xi\eta}(x^{(1)},x^{(2)}\ldots x^{(p)})
=\log\frac{V(x^{(1)},x^{(2)}\ldots x^{(p)};\eta)}{V(x^{(1)},x^{(2)}\ldots x^{(p)};\xi)};
\]
für die allgemeine Transcendente \(T_{\xi\eta}\binom{x}{z}\) findet sich:
\[
T_{\xi\eta}\binom{x^{(1)},x^{(2)}\ldots x^{(p)}}{z^{(1)},z^{(2)}\ldots z^{(p)}}
=\log\frac{V(x^{(1)},x^{(2)}\ldots x^{(p)};\eta)V(z^{(1)},z^{(2)}\ldots z^{(p)};\xi)}{V(z^{(1)},z^{(2)}\ldots z^{(p)};\eta)V(x^{(1)},x^{(2)}\ldots x^{(p)};\xi)}.
\]
Und endlich findet man für die Grundformel des Umkehrungsproblems bei gehörig gewählten Integrationswegen:
\begingroup\small
\[
\prod_{i=1}^{i=p}\frac{\left(\frac{\varphi}{\psi}\right)_{x^{(i)}}-\lambda}{\left(\frac{\varphi}{\psi}\right)_{c^{(i)}}-\lambda}
=\prod_{\xi\eta}\frac{V(x^{(1)},x^{(2)}\ldots x^{(p)};\xi)V(c^{(1)},c^{(2)}\ldots c^{(p)};\eta)}{V(x^{(1)},x^{(2)}\ldots x^{(p)};\eta)V(c^{(1)},c^{(2)}\ldots c^{(p)};\xi)},
\]
\endgroup
wodurch alles auf die Untersuchung der überall endlichen Function \(V\) zurückgeführt ist.

\subsubsection*{\S\,51. Eigenschaften der Function \(V=e^{-U}\).}

Die Differenz zweier Functionen \(U\), welche sich nur durch die Werthe von \(\xi\) unterscheiden, führte auf die specielle Transcendente \(\frac12T_{\xi\eta}(x)\). Ganz ebenso erhält man die Reihe der adjungirten Transcendenten, indem man die Differenz zweier Functionen \(U\) betrachtet, welche sich nur durch einen der Puncte \(x\) von einander unterscheiden. In der That erhält man aus dem Differentialausdrucke von \(U\):
\[
\begin{aligned}
&U(x^{(1)},x^{(2)}\ldots x^{(p)};\xi)-U(y^{(1)},x^{(2)}\ldots x^{(p)};\xi)\nobreak\hspace{0.25em}\nobreak=\frac12\int_{y^{(1)}}^{x^{(1)}}\frac{\partial T^{(1)}}{\partial x^{(1)}}\,dx^{(1)}\\
&\quad=\frac12T_{x^{(1)}\eta}(\xi,x^{21},x^{31}\ldots x^{p1})-\frac12T_{y^{(1)}\eta}(\xi,x^{21},x^{31}\ldots x^{p1})\\
&\quad=\frac12T_{x^{(1)}y^{(1)}}(\xi,x^{21},x^{31}\ldots x^{p1}).
\end{aligned}
\]
Diese Transcendente ist so zu definiren, dass sie die Differenz der beiden adjungirten Transcendenten
\[
\frac12T_{x^{(1)}\mu}(\xi,x^{21}\ldots)-\frac12T_{y^{(1)}\mu}(\xi,x^{21}\ldots)
\]
wird. Man muss also \(y^{(1)}\) durch eine Curve mit dem oben angeführten Curvensystem verbinden, und erhält dann beide Transcendenten aus \(\frac12T_{\lambda\mu}(c)\) dem Obigen zufolge eindeutig definirt.

Drückt man nun diese adjungirte Transcendente direct durch Functionen \(U\) aus, so erhält man die Relation:
\[
\tag{1}\begin{aligned}
&U(x^{(1)},x^{(2)}\ldots x^{(p)};\xi)-U(y^{(1)},x^{(2)}\ldots x^{(p)};\xi)\\
&\quad=U(\xi,x^{21},x^{31}\ldots x^{p1};x^{(1)})-U(\xi,x^{21},x^{31}\ldots x^{p1};y^{(1)}),
\end{aligned}
\]
oder wenn wir die Function \(V\) einführen:
\[
\tag{2}\frac{V(x^{(1)},x^{(2)}\ldots x^{(p)};\xi)}{V(y^{(1)},x^{(2)}\ldots x^{(p)};\xi)}
=\frac{V(\xi,x^{21},x^{31}\ldots x^{p1};x^{(1)})}{V(\xi,x^{21},x^{31}\ldots x^{p1};y^{(1)})}.
\]
Diese Formel lehrt den Quotienten zweier Functionen \(V\), welche sich durch ein Argument unterscheiden, auf den Quotienten zweier \(V\) zurückführen, welche sich durch einen Parameter unterscheiden, und umgekehrt.

Der Ausdruck \(\frac12T_{x^{(1)}y^{(1)}}(\xi,x^{21},x^{31}\ldots x^{p1})\) lässt sich aber unter Umständen noch auf eine andre Art auf Functionen \(U\) zurückführen; nämlich dann, wenn \(y^{(1)}\) mit den \(x\) auf einer Curve \((n-2)^{\mathrm{ter}}\) Ordnung liegt, welche durch die Doppel- und Rückkehrpuncte von \(f=0\) und überdies durch die übrigen \(n-2\) Schnittpuncte einer Geraden \(\xi\eta\) geht.

Bezeichnen wir die weitern Schnittpuncte dieser Curve mit \(f=0\) durch \(y^{(2)},y^{(3)}\ldots y^{(p)}\), und sind ferner \(y^{21},y^{31}\ldots y^{p1}\) die Puncte von \(f=0\), welche mit den Doppel- und Rückkehrpuncten und mit \(y^{(2)},y^{(3)}\ldots y^{(p)}\) auf einer Curve \((n-3)^{\mathrm{ter}}\) Ordnung liegen, so bilden die folgenden Curven:

\begin{center}
\begin{tabular}{@{}p{5.15in}@{}}
die Gerade \(\xi\eta\),\\
die Gerade \(x^{(1)}y^{(1)}\),\\
die Curve \((n-3)^{\mathrm{ter}}\) Ordnung durch \(x^{(2)},x^{(3)},\ldots,x^{21},x^{31}\ldots\) und die Doppel- und Rückkehrpuncte,\\
die Curve \((n-3)^{\mathrm{ter}}\) Ordnung durch \(y^{(2)},y^{(3)},\ldots,y^{31},y^{32}\ldots\) und die Doppel- und Rückkehrpuncte,
\end{tabular}
\end{center}
zusammen eine uneigentliche Curve \(2n-4\)ter Ordnung, welche die Doppel- und Rückkehrpuncte von \(f=0\) selbst zu Doppelpuncten hat. Aber von den Schnittpuncten dieser Curve mit \(f=0\) liegen der Annahme nach die \(n-2\) übrigen Puncte der Geraden \(\xi\eta\), die \(x\) und die \(y\) auf einer durch die Doppel- und Rückkehrpuncte gehenden Curve \((n-2)^{\mathrm{ter}}\) Ordnung. Dasselbe muss also mit dem Rest der Fall sein; es liegen also die übrigen Puncte der Geraden \(x^{(1)},y^{(1)}\) auf einer solchen Curve mit \(\xi,\eta,x^{21},x^{31}\ldots,y^{21},y^{31}\ldots\). Nach den oben über die Bildung der Transcendente \(\frac12T_{\xi\eta}\) gemachten Voraussetzungen hat man daher die Gleichung:
\[
T_{x^{(1)}y^{(1)}}(\xi,x^{21},x^{31}\ldots)+T_{x^{(1)}y^{(1)}}(\eta,y^{21},y^{31}\ldots)=0.
\]
Für die Functionen \(U\) giebt dies die Gleichung
\[
\begin{aligned}
&U(\xi,x^{21},x^{31}\ldots x^{p1};x^{(1)})-U(\xi,x^{21},x^{31}\ldots x^{p1};y^{(1)})\\
&\quad=-U(\eta,y^{21},y^{31}\ldots y^{p1};x^{(1)})+U(\eta,y^{21},y^{31}\ldots y^{p1};y^{(1)}),
\end{aligned}
\]
oder mit Anwendung der Gleichung (1):
\[
\begin{aligned}
&U(x^{(1)},x^{(2)}\ldots x^{(p)};\xi)-U(y^{(1)},x^{(2)}\ldots x^{(p)};\xi)\\
&
=-U(x^{(1)},y^{(2)}\ldots y^{(p)};\eta)+U(y^{(1)},y^{(2)}\ldots y^{(p)};\eta),
\end{aligned}
\]
oder endlich, wenn wir die Function \(V\) einführen:
\[
\tag{3}\frac{V(x^{(1)},x^{(2)}\ldots x^{(p)};\xi)}{V(y^{(1)},x^{(2)}\ldots x^{(p)};\xi)}
=\frac{V(y^{(1)},y^{(2)}\ldots y^{(p)};\eta)}{V(x^{(1)},y^{(2)}\ldots y^{(p)};\eta)}.
\]
Schreiben wir dies in der Form
\[
\frac{V(x^{(1)},x^{(2)}\ldots x^{(p)};\xi)}{V(y^{(1)},y^{(2)}\ldots y^{(p)};\eta)}
=\frac{V(y^{(1)},x^{(2)}\ldots x^{(p)};\xi)}{V(x^{(1)},y^{(2)}\ldots y^{(p)};\eta)},
\]
so sieht man, dass der Ausdruck links durch Vertauschung eines Paares \(x^{(1)},y^{(1)}\) sich nicht ändert; und zwar ist dabei nur vorausgesetzt, dass die im Zähler und Nenner als Argumente befindlichen Puncte mit den übrigen Puncten der Geraden \(\xi\eta\) auf einer durch die Doppel- und Rückkehrpuncte gehenden Curve \((n-2)^{\mathrm{ter}}\) Ordnung liegen. Da diese Bedingung rechts noch ebenso erfüllt ist, so können wir fortfahren, weitere Paare \(x,y\) zu vertauschen, ohne dass der Werth des Quotienten sich ändert, und wir gelangen also endlich zu der Gleichung:
\[
\tag{4}\frac{V(x^{(1)},x^{(2)}\ldots x^{(p)};\xi)}{V(y^{(1)},y^{(2)}\ldots y^{(p)};\eta)}
=\frac{V(y^{(1)},y^{(2)}\ldots y^{(p)};\xi)}{V(x^{(1)},x^{(2)}\ldots x^{(p)};\eta)}.
\]
Betrachten wir nun zweitens die Summe
\[
T_{\xi\eta}(x^{(1)}x^{(2)}\ldots x^{(p)})+T_{\xi\eta}(y^{(1)}y^{(2)}\ldots y^{(p)}),
\]
so zeigen abermals die Voraussetzungen des \S\,47, dass diese Summe verschwindet, und man hat daher auch, wenn man die so erhaltene Gleichung auf die Functionen \(V\) überträgt:
\[
\tag{5}\frac{V(x^{(1)},x^{(2)}\ldots x^{(p)};\xi)}{V(x^{(1)},x^{(2)}\ldots x^{(p)};\eta)}
=\frac{V(y^{(1)},y^{(2)}\ldots y^{(p)};\eta)}{V(y^{(1)},y^{(2)}\ldots y^{(p)};\xi)}.
\]
Die Combination der Gleichungen (4) (5) liefert nun sofort die Relation:
\[
V(x^{(1)},x^{(2)}\ldots x^{(p)};\xi)=\pm V(y^{(1)},y^{(2)}\ldots y^{(p)};\eta).
\]
Welches Zeichen hier zu wählen sei, hängt von der Art und Weise ab, in welcher die Function \(V\) definirt ist. Es genügt, dass der Continuität wegen entweder immer das obere oder immer das untere Zeichen zu wählen ist. Daher haben wir den folgenden Satz:

I. Wenn man die Function \(V\) für \(x^{(1)},x^{(2)}\ldots x^{(p)},\xi\) und für \(y^{(1)},y^{(2)}\ldots y^{(p)},\eta\) bildet, wo die \(x,y\) auf einer Curve \((n-2)^{\mathrm{ter}}\) Ordnung liegen, welche durch die Doppel- und Rückkehrpuncte, sowie durch die \(n-2\) übrigen Schnittpuncte von \(f=0\) mit der Geraden \(\xi\eta\) hindurchgeht, so ergiebt sich bei passender Wahl der Integrationswege beide Male derselbe oder der entgegengesetzte Werth.

Lassen wir nun die Puncte \(y,\eta\) fest und bewegen die \(x,\xi\), so bleibt der Werth von \(V\) dem Vorigen nach ungeändert, und man hat den folgenden Satz:

II. Lassen wir einen Strahl sich um einen Punct der Curve \(f=0\) drehen, und legen in jeder Lage desselben eine Curve \((n-2)^{\mathrm{ter}}\) Ordnung durch \(n-2\) seiner andern Schnittpuncte, durch die Doppel- und Rückkehrpuncte und durch \(p\) feste Puncte von \(f=0\), so schneidet dieselbe noch in \(p\) Puncten \(x\), und die Function \(V\), für diese \(p\) Puncte gebildet, und mit einem Parameter, welcher dem letzten Schnittpunct des Strahls entspricht, hat immer denselben Werth.

Ein specieller Fall des Satzes I. verdient besonders hervorgehoben zu werden. Wenn nämlich \(\eta\) mit einem der Puncte \(x\), etwa mit \(x^{(1)}\) zusammenfällt, so geht die Curve \((n-2)^{\mathrm{ter}}\) Ordnung in die Gerade \(\xi x^{(1)}\) und in eine Curve \((n-3)^{\mathrm{ter}}\) Ordnung über, welche durch \(x^{(2)},x^{(3)}\ldots x^{(p)}\) und die Doppel- und Rückkehrpuncte geht und dadurch bestimmt ist, und deren übrige Schnittpuncte mit \(f=0\) die Puncte \(x^{21},x^{31}\ldots x^{p1}\) sind. Man hat also die \(y\) durch \(\xi,x^{21},x^{31}\ldots x^{p1}\) zu ersetzen, und findet so den Satz:

III. Die Function \(V\) ändert sich nicht, oder doch nur ihr Zeichen, wenn man den Parameter mit einem ihrer Argumente vertauscht, und die übrigen Argumente durch diejenigen Puncte ersetzt, welche mit denselben auf einer durch die Doppel- und Rückkehrpuncte gehenden Curve \((n-3)^{\mathrm{ter}}\) Ordnung liegen.

Das Theorem II. ist äquivalent mit der Aufstellung einer partiellen Differentialgleichung, welcher die Function \(V\) genügt. Ist nämlich wieder \(\eta\) der feste Punct des Strahls, und sind \(\xi,\xi^{(1)},\xi^{(2)}\ldots\xi^{(n-2)}\) seine übrigen Schnittpuncte; liegen endlich mit den \(\xi\) und festen Puncten die \(x\) auf einer Curve \((n-2)^{\mathrm{ter}}\) Ordnung, so folgen aus dem Abel'schen Theorem für die entsprechenden Integrale erster Gattung die Differentialgleichungen:
\[
\frac{\partial u_i(\xi)}{\partial\xi}\,d\xi+
\sum_{h=1}^{h=n-2}\frac{\partial u_i(\xi^{(h)})}{\partial\xi^{(h)}}\,d\xi^{(h)}=0,
\]
\[
\sum_{h=1}^{h=p}\frac{\partial u_i(x^{(h)})}{\partial x^{(h)}}\,dx^{(h)}+
\sum_{h=1}^{h=n-2}\frac{\partial u_i(\xi^{(h)})}{\partial\xi^{(h)}}\,d\xi^{(h)}=0,
\]
oder nach Elimination der \(\xi\):
\[
\tag{6}\sum_{h=1}^{h=p}\frac{\partial u_i(x^{(h)})}{\partial x^{(h)}}\,dx^{(h)}-\frac{\partial u_i(\xi)}{\partial\xi}\,d\xi=0.
\quad (i=1,2\ldots p).
\]
Diese \(p\) Differentialgleichungen bilden ein System mit \(p+1\) Variabeln, deren vollständige Integrale durch die \(p\) Gleichungen:
\[
\tag{7}\sum_{h=1}^{h=p}\int_{c^{(h)}}^{x^{(h)}}du_i-\int_{\mu}^{\xi}du_i=\mathrm{Const.}
\quad (i=1,2\ldots p)
\]
gegeben sind. Zugleich aber ist nach Theorem II.
\[
dV=\frac{\partial V}{\partial\xi}\,d\xi+
\sum_{h=1}^{h=p}\frac{\partial V}{\partial x^{(h)}}\,dx^{(h)}=0.
\]
Combinirt man dies mit den Differentialgleichungen (6), so ergiebt sich nach Elimination der \(dx\) die partielle Differentialgleichung, welcher \(V\) genügt:
\[
\det\begin{pmatrix}
\dfrac{\partial V}{\partial x^{(1)}} & \dfrac{\partial V}{\partial x^{(2)}} & \cdots & \dfrac{\partial V}{\partial\xi}\\[0.7em]
\dfrac{\partial u_1(x^{(1)})}{\partial x^{(1)}} & \dfrac{\partial u_1(x^{(2)})}{\partial x^{(2)}} & \cdots & -\dfrac{\partial u_1(\xi)}{\partial\xi}\\
\vdots & \vdots & & \vdots\\
\dfrac{\partial u_p(x^{(1)})}{\partial x^{(1)}} & \dfrac{\partial u_p(x^{(2)})}{\partial x^{(2)}} & \cdots & -\dfrac{\partial u_p(\xi)}{\partial\xi}
\end{pmatrix}=0.
\]
Indem man diese integrirt, findet man \(V\) als Function der Integrale des Systems (6), also der linken Theile von (7). Und wir haben also den Satz:

Die Function \(V\) ist in Wirklichkeit nur von \(p\) Grössen abhängig, nämlich von den \(p\) Summen
\[
\sum_{h=1}^{h=p}\int_{c^{(h)}}^{x^{(h)}}du_i-\int_{\mu}^{\xi}du_i,
\quad (i=1,2\ldots p),
\]
oder im Sinne des Umkehrproblems, sie ist eine Function der \(p\) Grössen
\[
v_i-\int_{\mu}^{\xi}du_i.
\]
Die Bestimmung dieser Function ist die Lösung des Umkehrproblems. Wir werden diese in dem folgenden Abschnitt behandeln.

\section*{Achter Abschnitt.}
\section*{Monodromie.}

\subsubsection*{\S\,52. Differential einer Abel'schen Function. Werthsysteme der \(v\), für welche das System der \(x\) unbestimmt ist.}

Wir haben im Vorigen gesehen, dass das Umkehrproblem darin besteht, symmetrische Functionen der Coordinaten von \(p\) Puncten als Functionen der Integralsummen \(v_1,v_2\ldots v_p\) zu betrachten. Den Character dieser Functionen wollen wir jetzt näher untersuchen. Wir werden nachweisen, dass dieselben monodrome Functionen ihrer Argumente sind.

Wir können den Zusammenhang der \(x\) mit den \(v\) dahin aussprechen, dass zwischen diesen beiden Systemen von Grössen die Differentialgleichungen bestehen:
\[
\tag{1}\left\{\begin{aligned}
dv_1&=\frac{\partial u_1(x^{(1)})}{\partial x^{(1)}}dx^{(1)}+\frac{\partial u_1(x^{(2)})}{\partial x^{(2)}}dx^{(2)}\cdots+\frac{\partial u_1(x^{(p)})}{\partial x^{(p)}}dx^{(p)},\\
dv_2&=\frac{\partial u_2(x^{(1)})}{\partial x^{(1)}}dx^{(1)}+\frac{\partial u_2(x^{(2)})}{\partial x^{(2)}}dx^{(2)}\cdots+\frac{\partial u_2(x^{(p)})}{\partial x^{(p)}}dx^{(p)},\\[-0.35em]
&\hspace{4em}\cdots\qquad\cdots\qquad\cdots\\[-0.35em]
dv_p&=\frac{\partial u_p(x^{(1)})}{\partial x^{(1)}}dx^{(1)}+\frac{\partial u_p(x^{(2)})}{\partial x^{(2)}}dx^{(2)}\cdots+\frac{\partial u_p(x^{(p)})}{\partial x^{(p)}}dx^{(p)},
\end{aligned}\right.
\]
während zugleich gewisse Anfangswerthe \(x^{(1)}=c^{(1)},x^{(2)}=c^{(2)}\ldots x^{(p)}=c^{(p)}\), den Anfangswerthen \(v_1=0,v_2=0\ldots v_p=0\) entsprechend, nebst den zugehörigen \(s_x\) gegeben sind, so dass, während die \(v\) von den Werthen 0 aus wachsen, die \(x\) von den \(c\) sich allmälig entfernen. Lösen wir die Gleichungen (1) nach den \(dx\) auf, so erhalten wir zunächst die ersten Differentialquotienten der \(x\) nach den \(u\), sodann aber durch fortgesetzte Differentiation auch die höhern Differentialquotienten als Functionen der \(x\) ausgedrückt. Setzen wir in diesen Ausdrücken für die \(x\) die Werthe der \(c\), so ergeben sich die Coefficienten für eine nach der Taylor'schen Reihe ausgeführte Entwicklung nach aufsteigenden Potenzen der \(v\), welche immer innerhalb gewisser Grenzen convergirt und die Functionen \(x\) wirklich darstellt, wenn nur, was immer möglich ist, die \(c\) so gewählt sind, dass alle Coefficienten der Reihe endlich bleiben; und zwar (vgl. Briot und Bouquet, Théorie des ff. ell.) so lange convergirt, als die Differentialquotienten der \(x\) nach den \(u\) synektisch bleiben. Man kann aber auch über ein solches Gebiet hinaus eine Entwicklung aufstellen, wobei dann nur eines der bereits gefundenen Werthsysteme der \(x\) den Ausgang bildet; und man gelangt auf diese Art für jedes Werthsystem der \(v\) zu einer convergenten Entwicklung, welche freilich für verschiedene Werthsysteme verschieden sein kann, und bei denen immer diejenigen Werthsysteme vermieden werden, für welche die Coefficienten der aufgelösten Gleichungen (1) nicht mehr endlich sind.

Der Nenner, welcher bei der Auflösung der Gleichungen (1) auftritt, verschwindet, wenn zwei Puncte \(x\) im Laufe der Entwicklung zusammenfallen. Daher hört die Entwicklung der \(x\) nach dem Taylor'schen Satze auf zu convergiren, sobald ein Werthsystem erreicht ist, für welches zwei Puncte \(x\) coincidiren. Dies wird vermieden, wenn wir eine symmetrische Function der \(x\) entwickeln. Eine solche Function
\[
\frac{\varphi(x^{(1)},x^{(2)}\ldots x^{(p)})}{\psi(x^{(1)},x^{(2)}\ldots x^{(p)})}
\]
denken wir uns, um die Ideen zu fixiren, als eine algebraische Function, welche zugleich eine rationale der \(x\) und der \(s_x\) ist, und welche sich also nicht ändert, wenn man zwei der \(x\) und gleichzeitig zwei der \(s_x\) mit einander vertauscht. Es ist dann nämlich
\[
\tag{2}\frac{\partial\frac{\varphi}{\psi}}{\partial v_i}
=\frac{\displaystyle\sum \pm\frac{\partial u_1(x^{(1)})}{\partial x^{(1)}}\frac{\partial u_2(x^{(2)})}{\partial x^{(2)}}\cdots\frac{\partial\frac{\varphi}{\psi}}{\partial x^{(i)}}\cdots\frac{\partial u_p(x^{(p)})}{\partial x^{(p)}}}{\displaystyle\sum \pm\frac{\partial u_1(x^{(1)})}{\partial x^{(1)}}\frac{\partial u_2(x^{(2)})}{\partial x^{(2)}}\cdots\frac{\partial u_i(x^{(i)})}{\partial x^{(i)}}\cdots\frac{\partial u_p(x^{(p)})}{\partial x^{(p)}}},
\]
der Quotient zweier nach den \(x\) genommenen algebraischen Functionaldeterminanten: deren jede im Allgemeinen bei der Coincidenz zweier \(x\) nur einfach verschwindet, und deren Quotient daher endlich bleibt.

Die Functionaldeterminanten in (2) sind übrigens einer bemerkenswerthen Umgestaltung fähig. Man hat nämlich
\[
\frac{\partial u_k(x)}{\partial x}=\frac{\Theta_k(x)}{f'(s_x)},
\]
wo \(\Theta_k=0\) die Gleichung einer Curve \((n-3)^{\mathrm{ter}}\) Ordnung ist, welche durch die Doppel- und Rückkehrpuncte der Curve \(f=0\) hindurchgeht. Setzt man nun auch
\[
\frac{\partial\frac{\varphi}{\psi}}{\partial x^{(i)}}=\frac{N_i}{f'(s_x^{(i)})},
\]
so geht die Formel (2) über in
\[
\tag{3}\frac{\partial\frac{\varphi}{\psi}}{\partial v_i}
=\frac{\sum\pm\Theta_1(x^{(1)})\Theta_2(x^{(2)})\cdots N_i\cdots\Theta_p(x^{(p)})}{\sum\pm\Theta_1(x^{(1)})\Theta_2(x^{(2)})\cdots\Theta_i(x^{(i)})\cdots\Theta_p(x^{(p)})},
\]
oder wenn man die Determinante im Nenner jetzt durch \(H\), und ihre Unterdeterminanten durch \(H_i^{(k)}\) bezeichnet, wird
\[
\tag{4}\frac{\partial\frac{\varphi}{\psi}}{\partial v_i}=\frac{H_i^{(1)}N_1+H_i^{(2)}N_2\cdots+H_i^{(p)}N_p}{H}.
\]
Der Zuwachs der Function \(\frac{\varphi}{\psi}\) ist also jederzeit durch die Formel
\[
\tag{5}d\frac{\varphi}{\psi}=\frac{\sum_i\sum_k H_i^{(k)}N_k\,dv_i}{H}
\]
bestimmt, wo die Coefficienten rechts für jedes System der \(x\) im Allgemeinen vollkommen bestimmte Werthe annehmen.

Eine Unterbrechung der Monodromie kann also nur für solche Werthsysteme eintreten, für welche einer der Coefficienten von \(d\varphi\) unendlich gross oder unbestimmt wird. Dies kann nur geschehen, wenn entweder \(H=0\), oder wenn der Nenner des Ausdrucks
\[
N_k=f'(s_x^{(k)})\cdot\frac{\psi\frac{\partial\varphi}{\partial x^{(k)}}-\varphi\frac{\partial\psi}{\partial x^{(k)}}}{\psi^2},
\]
also \(\psi\) verschwindet. Aber in dem letzten Falle braucht man nur die reciproke Function \(\frac{\psi}{\varphi}\) zu betrachten, um zu erkennen, dass die Monodromie nicht gestört wird. Denn man hat
\[
d\frac{\psi}{\varphi}=\frac{\sum_i\sum_k H_i^{(k)}P_k\,dv_i}{H},
\]
wo
\[
P_k=f'(s_x^{(k)})\cdot\frac{\varphi\frac{\partial\psi}{\partial x^{(k)}}-\psi\frac{\partial\varphi}{\partial x^{(k)}}}{\varphi^2};
\]
und das Differential von \(\frac{\psi}{\varphi}\) wird also für \(\psi=0\) keineswegs mehr unendlich, sondern die einzelnen Coefficienten nehmen bestimmte endliche Werthe an. Somit wird für diese Werthsysteme \(\frac{\varphi}{\psi}\) nur unendlich, ohne dass von einer Verzweigung die Rede ist.

Es bleibt also nur das Verhalten der Function für diejenigen Werthsysteme der \(v\) zu untersuchen, denen Puncte \(x\) entsprechen, welche mit den Doppel- und Rückkehrpuncten von \(f=0\) auf einer Curve \((n-3)^{\mathrm{ter}}\) Ordnung liegen. Denn in der That ist dieses der Sinn der Gleichung \(H=0\). Der Anblick der Determinante, welche durch \(H\) bezeichnet wurde, lehrt, dass \(H\) für die Coordinaten jedes der \(p\) Puncte von der \((n-3)^{\mathrm{ten}}\) Ordnung ist; dass es verschwindet, wenn eines der \(x\) mit einem andern oder mit einem Doppel- oder Rückkehrpunct zusammenfällt. Sieht man daher \(p-1\) Puncte \(x\) als constant an, den \(p^{\mathrm{ten}}\) als variabel, so ist \(H=0\) die Gleichung einer oben bezeichneten Curve \((n-3)^{\mathrm{ter}}\) Ordnung, oder \(H=0\) ist die Bedingung dafür, dass alle \(p\) Puncte sich auf einer solchen Curve befinden.

Hier tritt nun zuvörderst der bemerkenswerthe Umstand ein, dass diese Bedingung, obschon nur eine Bedingung für die \(x\), doch zwei Bedingungen für die \(v\) involvirt. Dies lehrt das Abel'sche Theorem. Bezeichnen wir nämlich durch \(\alpha^{(1)},\alpha^{(2)},\ldots,\alpha^{(2p-2)}\) irgend \(2p-2\) Puncte von \(f=0\), welche mit den Doppel- und Rückkehrpuncten auf einer Curve \((n-3)^{\mathrm{ter}}\) Ordnung liegen, und sind ferner \(x^{(p+1)},x^{(p+2)}\ldots x^{(2p-2)}\) diejenigen Puncte, in welchen in solchem Fall die durch die übrigen \(x\) bestimmte Curve \((n-3)^{\mathrm{ter}}\) Ordnung unsre Curve \(f=0\) noch schneidet. Dann kann man die \(\alpha,x\) einander so zuordnen und die Integrationswege so wählen, dass
\[
\sum_{i=1}^{i=2p-2}\int_{\alpha^{(i)}}^{x^{(i)}}du_k=0\qquad (k=1,2\ldots p).
\]
Und man kann also durch passende Wahl der Integrationswege auch die Gleichungen herstellen:
\[
\sum_{i=1}^{i=2p-2}\int_{\alpha^{(i)}}^{x^{(p+i)}}du_k=-v_k+\sum_{i=1}^{i=p}\int_{c^{(i)}}^{\alpha^{(i)}}du_k\qquad(k=1,2\ldots p).
\]
Da nun rechts neben \(-v_k\) eine Constante steht, und links die untern Grenzen ebenfalls constant sind, so hängen die \(p\) Grössen \(v_k\) nur noch von den \(p-2\) Puncten \(x^{(p+1)},x^{(p+2)},\ldots,x^{(2p-2)}\) ab, sind also Functionen von \(p-2\) Variabeln, und es müssen daher zwei Gleichungen zwischen den \(v\) stattfinden, damit zwischen den \(x\) die eine Bedingung eintrete, dass sie auf einer durch die Doppel- und Rückkehrpuncte gelegten Curve \((n-3)^{\mathrm{ter}}\) Ordnung liegen.

Ebendies tritt in anderer Weise durch folgende Betrachtung hervor. Nehmen wir an, wir hätten ein System der \(x\) gefunden, für welches \(H=0\), und wir kennten das entsprechende System der \(v\). Aus den Gleichungen
\[
dv_k=\frac{\partial u_k(x^{(1)})}{\partial x^{(1)}}dx^{(1)}+\frac{\partial u_k(x^{(2)})}{\partial x^{(2)}}dx^{(2)}\cdots+\frac{\partial u_k(x^{(p)})}{\partial x^{(p)}}dx^{(p)}
\]
folgt dann durch Multiplication mit den Unterdeterminanten \(H_i^{(k)}\) und Addition:
\[
H_i^{(1)}dv_1+H_i^{(2)}dv_2\cdots+H_i^{(p)}dv_p=0\qquad(i=1,2\ldots p).
\]
Diese Gleichung stellt indessen in Wahrheit nicht \(p\) verschiedene Gleichungen dar, sondern nur eine einzige; denn da \(H\) verschwindet, so kann man immer solche Reihen von Grössen \(a,b\) bestimmen, dass die Unterdeterminanten die Werthe annehmen:
\[
H_i^{(k)}=a_i b_k,
\]
wodurch dann das obige Gleichungssystem sich auf die eine Gleichung reducirt:
\[
\tag{6} b_1dv_1+b_2dv_2\cdots+b_pdv_p=0.
\]
Man sieht also, dass einem System der \(v\), welches den oben bemerkten zwei Bedingungsgleichungen genügt, unendlich viele Systeme der \(x\) genügen müssen, dergestalt, dass die Wahl eines dieser Systeme der \(x\), durch welche die Grössen \(b\) bestimmte Werthe erhalten, eine Bedingung (6) für die Zuwächse involvirt, welche die \(v\) erhalten sollen.

Das Auftreten solcher Werthsysteme ist in der That eine allgemeine Eigenschaft von Functionen mit mehreren Veränderlichen. Setzt man eine solche Function gleich Null, so erhält man eine Bedingung zwischen den Argumenten; setzt man sie gleich unendlich, so ergiebt sich eine zweite; das Zusammentreffen beider Bedingungen liefert Werthsysteme, für welche die Function unbestimmt ist, wo also einem System von Argumenten unendlich viele Functionswerthe entsprechen.

Sind \(v_1,v_2,\ldots,v_p\) die Argumente, und ist \(f=\frac{\varphi}{\psi}\) die zu untersuchende Function, so geben die Gleichungen \(\varphi=0,\psi=0\) zusammen die Werthe, für welche \(f\) unbestimmt wird. Indem man aber in \(\varphi,\psi\) für die \(v\) solche Werthe \(v+\delta v\) einführt, welche von denjenigen, die \(\varphi,\psi\) verschwinden machen, nur unendlich wenig verschieden sind, ergiebt sich
\[
f=\frac{\sum\frac{\partial\varphi}{\partial v_i}\delta v_i}{\sum\frac{\partial\psi}{\partial v_i}\delta v_i},
\]
so dass also jedem bestimmten unter den unendlich vielen dem System der \(v\) zugehörigen Functionswerthen eine lineare Relation der Zuwächse \(\delta v\) entspricht. Nennen wir das ganze Grössensystem, welches die \(p\) Grössen \(v\) unabhängig von einander durchlaufen können, einen Ort \(p^{\mathrm{ter}}\) Dimension, so bilden diejenigen Werthsysteme, für welche \(f\) unbestimmt wird, einen Ort \((p-2)^{\mathrm{ter}}\) Dimension. Man sieht daraus, dass die Monodromie einer Function dadurch nicht nothwendig aufgehoben wird, dass sie unendlich viele Werthe für Werthsysteme annimmt, welche einen Ort \((p-2)^{\mathrm{ter}}\) Dimension bilden: dass also eine übrigens vollkommen eindeutig bestimmte Function in einem Ort \((p-2)^{\mathrm{ter}}\) Dimension unendlich vieldeutig werden kann, ohne dass man berechtigt wäre, derselben die Eindeutigkeit abzusprechen. So ist schon die Function \(z=\frac{x}{y}\) für \(x=0,y=0\) vollkommen unbestimmt, obschon sie völlig monodrom ist. Aber die verschiedenen Werthe von \(z\), welche für \(x=0,y=0\) eintreten, ordnen sich so, dass jedem speciellen Werthe eine bestimmte Fortschreitung
\[
\delta x=z\delta y
\]
entspricht, dass man also zu einem ganz bestimmten Werthe von \(z\) gelangt, wenn man \(x\) und \(y\) in einer bestimmten Weise sich dem gemeinsamen Werth Null nähern lässt.

Man muss daher den Begriff der Monodromie bei einer Function von mehreren Variabeln mit der nöthigen Vorsicht definiren. Man darf nicht sagen, dass eine monodrome Function eine solche sei, welche für alle Werthsysteme der Variabeln nur einen Werth besitze; sondern es genügt im Allgemeinen schon, dass die Function überall eindeutig sei, mit Ausnahme eines Orts \((p-2)^{\mathrm{ter}}\) Dimension, in welchem sie jeden Werth annehmen kann.

Dass die symmetrische Function \(\frac{\varphi}{\psi}\) in unserm Falle wirklich die Form \(\frac00\) annimmt, geht daraus hervor, dass sie nach dem Frühern als Quotient von Functionen \(V\) darstellbar ist, welche sämmtlich verschwinden, sobald die \(x\) auf einer der gedachten Curven \((n-3)^{\mathrm{ter}}\) Ordnung liegen.

Begnügen wir uns daher für den Augenblick damit, nachgewiesen zu haben, dass \(\frac{\varphi}{\psi}\) ausserhalb dieses Orts \((p-2)^{\mathrm{ter}}\) Dimension eine monodrome Function der \(v\) ist, und gehen wir nun zur Untersuchung der Function \(V\) selbst über, aus welcher wir auch die Function \(\frac{\varphi}{\psi}\) an diesen Ausnahmestellen zu characterisiren lernen werden.

\subsubsection*{\S\,53. Beweis, dass die Function \(V\) eine synektische Function der \(v\) ist.}

Die Function \(V(x^{(1)},x^{(2)}\ldots x^{(p)};\xi)\) hat, wie oben nachgewiesen worden, die Eigenschaft, immer endlich zu bleiben. Die Differentialquotienten ihres Logarithmus nach den \(x\) (und deswegen auch nach den \(v\)) sind keineswegs algebraische Functionen; denn da bei den in \(-U\) auftretenden Integralen dritter Gattung Parameter und Argument veränderlich sind, so setzen sich die Differentialquotienten eines solchen Integrals aus algebraischen Theilen und aus den Periodicitätsmoduln \(S^{(k)}\) der unvollständig normirten Integrale dritter Gattung zusammen. Aber alle diese Glieder haben die Eigenschaft, zum Theil für gegebene Puncte \(x\) ganz bestimmte Werthe anzunehmen, anderntheils wenigstens völlig bestimmt zu sein, wenn die variabeln Puncte \(x\), wie dies bei dem Umkehrproblem der Fall ist, auf bestimmten Wegen aus gegebenen constanten Puncten entstanden sind. Sie haben ferner die Eigenschaft, sich dem Logarithmus einer einfach verschwindenden Function, ohne weitern Factor, zu nähern, wo \(\log V\) unendlich gross wird. Daher ist die Function \(V\) eine monodrome Function der \(v\) in demselben Umfange, in welchem dieses den \(x\) selbst zukommt. Es kann daher nur in folgenden Fällen eine Störung der Monodromie von \(V\) möglich erscheinen:

\begin{center}
\begin{tabular}{@{}p{5.15in}@{}}
1) Wenn zwei Puncte \(x\) coincidiren;\\
2) Wenn die Puncte \(x\) mit den Doppel- und Rückkehrpuncten von \(f=0\) auf einer Curve \((n-3)^{\mathrm{ter}}\) Ordnung liegen.
\end{tabular}
\end{center}

Wenn wir nun bei Betrachtung des Falles 1), wie nach \S\,51 gestattet ist, die Function \(V\) durch die mit ihr identische Function
\[
\varepsilon V(y^{(1)},y^{(2)}\ldots y^{(p)};\eta)\qquad(\varepsilon=\pm1)
\]
ersetzen, wo die \(y,x\) mit den übrigen Puncten der Geraden \(\xi\eta\), so wie mit den Doppel- und Rückkehrpuncten auf einer Curve \((n-2)^{\mathrm{ter}}\) Ordnung liegen, so entsprechen die \(y\) einem andern Umkehrproblem, bei welchem die \(v\) bis auf additive Constanten durch die \(-v\) ersetzt sind. Die \(y\) haben keine den Fällen 1) oder 2) unterworfene Lage; sie sind an der betreffenden Stelle monodrome Functionen der \(v\), mithin auch \(V\) selbst, und der Fall 1) begründet also eine Unterbrechung der Monodromie nicht.

Dasselbe Hülfsmittel dient auch zur Behandlung des zweiten Falles. In diesem Falle sind die \(x\) in der That keine bestimmten Functionen der \(v\) mehr, aber die \(y\) sind es noch; die Curve \((n-2)^{\mathrm{ter}}\) Ordnung zerfällt in die Gerade \(\xi\eta\) und in eine Curve \((n-3)^{\mathrm{ter}}\) Ordnung, welche die \(x\) und \(p-2\) andere Puncte enthält. Diese, zusammen mit den Puncten \(\xi,\eta\) selbst, bilden das System der \(y\); es fällt also in der zweiten Form von \(V\) ein Argument mit dem Parameter zusammen. In diesem Falle verhält sich, wie wir wissen, \(\log V\) wie der Logarithmus einer einfach verschwindenden algebraischen Function, und die Function bleibt in der Nähe eines solchen Werthsystems monodrom. Daher kann auch der Fall 2) eine Störung der Monodromie nicht begründen. Und so hat man den Satz:

\textbf{Die Function \(V(x^{(1)},x^{(2)}\ldots x^{(p)};\xi)\) ist eine stets endliche und monodrome, also eine synektische Function der \(v\).}

Hiemit ist das Umkehrproblem auf die Betrachtung synektischer Functionen zurückgeführt. Es wird sich nun darum handeln, die wirkliche Darstellung dieser Functionen zu geben.

Schon am Ende des \(7^{\mathrm{ten}}\) Abschnittes wurde nachgewiesen, dass die von \(p+1\) Puncten abhängige Function \(V\) diese nur in den \(p\) Verbindungen
\[
\sum_{i=1}^{i=p}\int_{c^{(i)}}^{x^{(i)}}du_h-\int_{\mu}^{\xi}du_h
\]
enthält. Wir wollen sie deswegen als Function von den \(p\) Variabeln
\[
w_h=\sum_{i=1}^{i=p}\int_{c^{(i)}}^{x^{(i)}}du_h-\int_{\mu}^{\xi}du_h+K_h
\]
betrachten, wo die \(K_h\) sogleich näher zu bestimmende Constanten bedeuten sollen.

Ersetzen wir nun die \(x\) durch die Puncte \(y\), \(\xi\) durch \(\eta\), so bleibt \(V\) bis auf das Vorzeichen ungeändert, d. h. es ist bis auf das Vorzeichen dieselbe Function der Grössen
\[
w'_h=\sum_{i=1}^{i=p}\int_{c^{(i)}}^{y^{(i)}}du_h-\int_{\mu}^{\eta}du_h+K_h.
\]
Und zwar besteht zwischen \(w,w'\) die Beziehung, dass \(w_h+w'_h\) eine Constante ist, wie das Abel'sche Theorem sofort lehrt. Wir können daher die Constanten \(K_h\) immer so bestimmen, dass diese Summe Null wird. Bei solcher Bestimmung werden die \(w'\) den \(w\) gleich und entgegengesetzt, und man hat also den Satz:

\textbf{Die Function \(V\), als Function der \(w\) betrachtet, ändert sich nicht oder doch nur ihr Zeichen, wenn man die \(w\) sämmtlich gleichzeitig ihr Vorzeichen ändern lässt.}

Die Constanten \(K\) bestimmen sich aus den Gleichungen
\[
2K_h+\sum_{i=1}^{i=p}\int_{c^{(i)}}^{x^{(i)}}du_h+\sum_{i=1}^{i=p}\int_{c^{(i)}}^{y^{(i)}}du_h-\int_{\mu}^{\xi}du_h-\int_{\mu}^{\eta}du_h=0,
\]
wenn man darin für \(\xi,\eta\) irgend zwei beliebig gewählte Puncte, für die \(x,y\) irgend \(2p\) constante Puncte setzt, welche mit den übrigen Puncten der die erstern beiden verbindenden Gerade und den Doppel- und Rückkehrpuncten von \(f=0\) auf einer Curve \((n-2)^{\mathrm{ter}}\) Ordnung liegen. Ueber diesen Gegenstand wird weiter unten ausführlicher zu handeln sein.

\subsubsection*{\S\,54. Periodicitätsmoduln von \(-U=\log V\).}

Untersuchen wir jetzt, in welcher Weise die Function \(V\), welche nur von den \(w\) abhängt, sich ändert, wenn die \(w\) um die allgemeinsten Periodicitätsmoduln von Integralen erster Gattung zunehmen, wenn also \(w_i\) in
\[
w_i'=w_i+2m_i\pi\sqrt{-1}+a_{1i}q_1+a_{2i}q_2\cdots+a_{pi}q_p
\]
übergeht, wo die \(m,q\) beliebige positive oder negative ganze Zahlen sind. Die Function, welche aus \(V\) entsteht, hängt nur von den \(w_i'\) ab; aber diese selbst unterscheiden sich von den \(w\) nur durch Constante, daher ist auch die aus \(V\) entstandene Function \(V'\) nur eine Function der \(w\). Ebendies ist sonach der Fall mit der Function
\[
\psi(w_1,w_2\ldots w_p)=\log V'-\log V,
\]
welche den gleichzeitigen Zuwachs des \(\log V\) oder der Function \(-U\) angiebt. Diese Function \(\psi\) wollen wir jetzt bestimmen.

Gehen wir von der Formel aus (\S\,50):
\[
\tag{1}T_{\xi\eta}\binom{x^{(1)},x^{(2)}\ldots x^{(p)}}{z^{(1)},z^{(2)}\ldots z^{(p)}}
=\log\frac{V(x^{(1)},x^{(2)}\ldots x^{(p)};\eta)V(z^{(1)},z^{(2)}\ldots z^{(p)};\xi)}{V(x^{(1)},x^{(2)}\ldots x^{(p)};\xi)V(z^{(1)},z^{(2)}\ldots z^{(p)};\eta)},
\]
in welcher vier Functionen \(V\), also auch ihnen entsprechend vier Systeme der \(w\) auftreten, nämlich
\[
\begin{aligned}
w_h&=\sum\int_{c^{(i)}}^{x^{(i)}}du_h-\int_{\mu}^{\xi}du_h+K_h,\\
w_h^{\circ}&=\sum\int_{c^{(i)}}^{z^{(i)}}du_h-\int_{\mu}^{\xi}du_h+K_h,\\
W_h&=\sum\int_{c^{(i)}}^{x^{(i)}}du_h-\int_{\mu}^{\eta}du_h+K_h,\\
W_h^{\circ}&=\sum\int_{c^{(i)}}^{z^{(i)}}du_h-\int_{\mu}^{\eta}du_h+K_h.
\end{aligned}
\]
Lassen wir nun einen Punct \(x\) einen solchen Periodenweg durchlaufen, durch dessen Zufügung \(w_h\) in \(w'_h\) übergeht, so geht zugleich \(W_h\) in \(W'_h\) über, und \(T_{\xi\eta}\binom{x}{z}\) wächst um (vgl. \S\,41)
\[
\int_{\eta}^{\xi}dI=\int_{\eta}^{\xi}(q_1du_1+q_2du_2\cdots+q_pdu_p),
\]
während die linke Seite der Gleichung (1) um
\[
\psi(W_1,W_2\ldots)-\psi(w_1,w_2\ldots)
\]
zunimmt. Man hat daher die Gleichung:
\[
\tag{2}\psi(W_1,W_2\ldots)-\psi(w_1,w_2\ldots)=\int_{\eta}^{\xi}dI.
\]
Lässt man zweitens \(\eta\) denselben Periodenweg, aber in entgegengesetzter Richtung beschreiben, so gehen die \(W_h\) abermals in die \(W'_h\), zugleich aber die \(W_h^{\circ}\) in \(W_h^{\circ\prime}\) über; die rechte Seite von (1) also wächst um
\[
\psi(W_1,W_2,\ldots)-\psi(W_1^{\circ},W_2^{\circ},\ldots),
\]
während die linke um
\[
\sum_{i=1}^{i=p}\int_{z^{(i)}}^{x^{(i)}}dI
\]
zunimmt. Daher hat man zweitens die Gleichung
\[
\tag{3}\psi(W_1,W_2\ldots)-\psi(W_1^{\circ},W_2^{\circ}\ldots)=\sum_{i=1}^{i=p}\int_{z^{(i)}}^{x^{(i)}}dI.
\]
Zieht man nun (2) von (3) ab, so findet sich:
\[
\psi(w_1,w_2\ldots)-\psi(W_1^{\circ},W_2^{\circ}\ldots)=\sum_{i=1}^{i=p}\int_{z^{(i)}}^{x^{(i)}}dI-\int_{\eta}^{\xi}dI.
\]
Setzt man an Stelle der \(z,\eta\) hier die constanten Puncte \(c,\mu\), so wird \(\psi(W_1^{\circ},W_2^{\circ}\ldots)\) eine constante Grösse, und man kann den Ausdruck für \(\psi(w_1,w_2\ldots)\) in die Form schreiben:
\[
\tag{4}\psi(w_1,w_2\ldots)=q_1w_1+q_2w_2\cdots+q_pw_p+C.
\]
Die Constante \(C\), welche nur noch von den Zahlen \(m,q\) abhängen kann, bleibt näher zu bestimmen.

Wenn wir mit Hülfe der Gleichung (4) die Function \(V'\) durch \(V\) ausdrücken, so haben wir
\[
\tag{5}V'=e^{w_1q_1+w_2q_2\cdots+w_pq_p+C}\cdot V.
\]
Da \(V'\) dieselbe Function von den \(w'\), wie \(V\) von den \(w\) ist, so geht \(V'\) in \(\varepsilon V\) über \((\varepsilon=\pm1)\), wenn man an Stelle der Grössen \(w'_h\) die Grössen
\[
-\{w_h+2m_h\pi\sqrt{-1}+q_1a_{1h}+q_2a_{2h}\cdots+q_pa_{ph}\}
\]
setzt; denn hierdurch verwandeln sich die \(w'\) in die \(-w\), und also \(V'\) in \(\varepsilon V\), da diese Function durch Aenderung der Vorzeichen aller \(w\) bis auf ihr Vorzeichen ungeändert bleibt. Zugleich aber geht aus demselben Grunde \(V\) in \(\varepsilon V'\) über, und man hat daher aus (5) die zweite Gleichung:
\[
\tag{6}V=e^{-\sum q_hw_h-\sum\sum q_hq_ka_{hk}+C}\cdot V';
\]
also, wenn man (5) mit (6) multiplicirt, und die gleichen Factoren auf beiden Seiten auslässt:
\[
1=e^{-\sum\sum q_hq_ka_{hk}+2C}.
\]
Mithin wird
\[
\tag{7}C=\frac12\sum\sum q_hq_ka_{hk}+L\pi\sqrt{-1},
\]
wo \(L\) eine ganze Zahl ist. Es bleibt endlich übrig, diese ganze Zahl zu bestimmen.

Zu diesem Zwecke denken wir uns dem obigen Zuwachs der \(w\) einen zweiten Zuwachs
\[
2m'_h\pi\sqrt{-1}+q'_1a_{1h}+q'_2a_{2h}\cdots+q'_pa_{ph}
\]
hinzugefügt. Die Function
\[
V'=e^{\sum w_hq_h+\frac12\sum\sum q_hq_ka_{hk}+L\pi\sqrt{-1}}\cdot V
\]
verwandelt sich hierdurch in
\[
e^{\sum w_h(q_h+q'_h)+\frac12\sum(q_h+q'_h)(q_k+q'_k)a_{hk}+(L+L')\pi\sqrt{-1}}\cdot V,
\]
wo \(L'\) von den Zahlen \(q',m'\) ebenso abhängt wie \(L\) von den Zahlen \(q,m\). Dasselbe Resultat muss man erhalten, wenn man in \(V\) direct die \(w_h\) um die Grössen
\[
2(m_h+m'_h)\pi\sqrt{-1}+a_{1h}(q_1+q'_1)+a_{2h}(q_2+q'_2)\cdots
\]
vermehrt. Hierdurch aber geht \(V\) über in
\[
e^{\sum w_h(q_h+q'_h)+\frac12\sum(q_h+q'_h)(q_k+q'_k)a_{hk}+L''\pi\sqrt{-1}}\cdot V,
\]
wo \(L''\) von den Summen \(m+m',q+q'\) abhängt, wie \(L,L'\) von den einzelnen Zahlenreihen.

Die Vergleichung giebt also (da von geraden Vielfachen von \(\pi\sqrt{-1}\) abgesehen werden kann):
\[
L+L'=L''.
\]
Wir können daher überhaupt die Zahl \(L\), welche einer Summe von beliebig vielen Zahlenreihen entspricht, gleich der Summe aller Zahlen \(L\) setzen, welche den einzelnen Zahlenreihen entsprechen. Gehen wir nun von denjenigen Zahlenreihen aus, in denen immer nur eine Zahl \(m\), oder nur eine Zahl \(q\) von Null verschieden und zwar gleich 1 ist. Aus diesen können wir die Zahlenreihen \(m_1,m_2,\ldots,q_1,q_2\ldots\) zusammensetzen, indem wir die erste \(m_1\)fach, die zweite \(m_2\)fach etc. wiederholen. Wenn wir also die Werthe von \(L\), welche den ersten \(2p\) Zahlenreihen entsprechen, durch
\[
\varkappa_1,\varkappa_2\ldots\varkappa_p,\lambda_1,\lambda_2\ldots\lambda_p
\]
bezeichnen (und zwar brauchen wir diesen Zahlen nur die Werthe 0 oder 1 beizulegen, wie dies auch bei den \(L\) der Fall ist), so haben wir für \(L\) den Ausdruck:
\[
L=\varkappa_1m_1+\varkappa_2m_2\cdots+\varkappa_pm_p+\lambda_1q_1+\lambda_2q_2\cdots+\lambda_pq_p.
\]
Und so haben wir endlich den Satz:

\textbf{Wenn die \(w_h\) um die Ausdrücke}
\[
2m_h\pi\sqrt{-1}+\sum_k a_{hk}q_k
\]
\textbf{wachsen, so geht die Function \(V\) über in}
\[
V\cdot e^{\sum w_hq_h+\frac12\sum a_{hk}q_hq_k+(\sum\varkappa_hm_h+\sum\lambda_hq_h)\pi\sqrt{-1}},
\]
\textbf{wo die \(\varkappa,\lambda\) ganze Zahlen sind, welchen man die Werthe 0 oder 1 beilegen kann.}

Wir haben also die Bestimmung des Periodicitätsmoduls von \(\log V\) so weit geführt, dass nur noch die \(2p\) ganzen Zahlen \(\varkappa,\lambda\) unbestimmt bleiben. Diese sind durch bestimmte Integrale definirt; sie sind die Periodicitätsmoduln von \(\log V\), vermindert um die bereits anderweitig gefundenen Bestandtheile desselben.

\subsubsection*{\S\,55. Reihenentwicklung.}

Setzen wir in der Formel, welche die Aenderung von \(V\) bei einer Aenderung der \(w\) um Periodicitätsmoduln angiebt, an Stelle der \(w\) die um halbe Periodicitätsmoduln vermehrten Grössen
\[
\tag{1}\omega_h=w_h-\lambda_h\pi\sqrt{-1}-\frac12(\varkappa_1a_{1h}+\varkappa_2a_{2h}\cdots+\varkappa_pa_{ph})
\]
ein. Indem wir die Gleichung
\[
V'=V\cdot e^{\sum w_hq_h+\frac12\sum q_hq_ka_{hk}+\sum(\varkappa_hm_h+\lambda_hq_h)\pi\sqrt{-1}}
\]
mit
\[
e^{-\frac12\sum\varkappa_h\omega_h-\sum\varkappa_hm_h\pi\sqrt{-1}-\frac12\sum\varkappa_hq_ka_{hk}}
\]
multipliciren, finden wir die Gleichung
\[
\begin{aligned}
&e^{-\frac12\sum\varkappa_h(\omega_h+2m_h\pi\sqrt{-1}+a_{1h}q_1+a_{2h}q_2\cdots)}\cdot V'\\
&\qquad=e^{-\frac12\sum\varkappa_h\omega_h}\cdot V\cdot e^{\sum q_h\omega_h+\frac12\sum q_hq_ka_{hk}}.
\end{aligned}
\]
Die Function
\[
e^{-\frac12\sum\varkappa_h\omega_h}\cdot V
\]
ändert sich also nicht, wenn man die \(\omega\) um Vielfache von \(2\pi\sqrt{-1}\) vermehrt, und ändert sich nur um den Factor
\[
e^{\sum q_h\omega_h+\frac12\sum q_hq_ka_{hk}},
\]
wenn die \(\omega_h\) um die Ausdrücke \(q_1a_{1h}+q_2a_{2h}\cdots\) vermehrt werden.

Aus dem ersten Umstande folgt, dass diese Function in eine stets convergirende Reihe nach den positiven und negativen Potenzen der \(p\) Grössen \(e^{\omega_h}\) entwickelt werden kann, dass man also setzen darf:
\[
e^{-\frac12\sum\varkappa_h\omega_h}\cdot V=\sum A_{r_1r_2\ldots r_p}e^{r_1\omega_1+r_2\omega_2\cdots+r_p\omega_p},
\]
wo die Summe rechts über alle Werthe der \(r\) von \(-\infty\) bis \(+\infty\) ausgedehnt werden muss.

Die Anwendung des zweiten Theiles von obigem Satze führt sodann auf die Formel:
\[
\sum A_{r_1r_2\ldots r_p}e^{\sum r_h\omega_h+\sum r_hq_ka_{hk}}
=e^{\sum q_h\omega_h+\frac12\sum q_hq_ka_{hk}}\sum A_{r_1r_2\ldots r_p}e^{\sum r_h\omega_h}.
\]
Setzt man hier links \(r_1+q_1,r_2+q_2\ldots\) für \(r_1,r_2\ldots\), was, da die Summe für die \(r\) von \(-\infty\) bis \(+\infty\) zu nehmen ist, nichts ändert, so erhält man
\[
\begin{aligned}
&\sum A_{r_1+q_1,r_2+q_2\ldots}e^{\sum(r_h+q_h)\omega_h+\sum r_hq_ka_{hk}+\sum q_hq_ka_{hk}}\\
&\qquad=e^{\frac12\sum q_hq_ka_{hk}}\sum A_{r_1r_2\ldots}e^{\sum(r_h+q_h)\omega_h}.
\end{aligned}
\]
Wenn man daher die Coefficienten gleicher Exponentialgrössen vergleicht, so findet man:
\[
A_{r_1+q_1,r_2+q_2\ldots}=e^{-\sum r_hq_ka_{hk}-\frac12\sum q_hq_ka_{hk}}\cdot A_{r_1r_2\ldots},
\]
oder wenn man alle \(r\) gleich Null setzt:
\[
A_{q_1q_2\ldots q_p}=e^{-\frac12\sum q_hq_ka_{hk}}A_{000\ldots}.
\]
Durch die obigen Periodicitätsbedingungen sind also alle Coefficienten der Reihenentwicklung vollständig bestimmt, bis auf einen allen gemeinsamen Factor, und man hat:
\[
e^{-\frac12\sum\varkappa_h\omega_h}\cdot V=A_{00\ldots}\sum e^{-\frac12\sum r_hr_ka_{hk}+\sum r_h\omega_h}.
\]
Bezeichnen wir jetzt durch \(\Theta(\omega_1,\omega_2\ldots\omega_p)\) die Function:
\[
\tag{2}\Theta(\omega_1,\omega_2\ldots\omega_p)=\sum e^{-\frac12\sum r_hr_ka_{hk}+\sum r_h\omega_h},
\]
so wird
\[
V=A_{00\ldots}\cdot e^{\frac12\sum\varkappa_h\omega_h}\cdot\Theta(\omega_1,\omega_2\ldots\omega_p).
\]
Die Constante \(A_{00\ldots}\) ist leicht zu bestimmen. Denn ihrer Entstehung nach wurde die Function \(U\) gleich Null, also \(V=1\), wenn die \(x,\xi\) in beliebige Anfangspuncte \(\gamma,\xi\) übergingen. Bezeichnet man also durch \(\omega_h^{\circ}\) dasjenige, was bei diesem Uebergange aus \(\omega_h\) wird, so ist
\[
\tag{3}V=e^{\frac12\sum\varkappa_h(\omega_h-\omega_h^{\circ})}\cdot\frac{\Theta(\omega_1,\omega_2\ldots\omega_p)}{\Theta(\omega_1^{\circ},\omega_2^{\circ}\ldots\omega_p^{\circ})}.
\]
Diese Reihe für \(V\) ist nach dem Frühern stets convergent. Die Function \(\Theta\) muss also in Folge der bei diesen Entwicklungen auftretenden Werthe der \(a_{hk}\) ebenfalls die Eigenschaft haben, für alle endlichen Werthe der \(\omega\) zu convergiren.

\subsubsection*{\S\,56. Die Argumente der Function \(\Theta\).}

Gehen wir jetzt zur näheren Betrachtung derjenigen Grössen \(\omega\) über, welche die Argumente der Function \(\Theta\) bilden. Diese entstanden durch Abziehen von halben Perioden aus den \(w\), so dass wir schreiben können
\[
\tag{1}\omega_h=w_h-\frac12P_h^{\varkappa,\lambda}.
\]


wenn
\[
P_h^{\varkappa,\lambda}=2\lambda_h\pi\sqrt{-1}+\varkappa_1a_{1h}+\varkappa_2a_{2h}\cdots+\varkappa_pa_{ph}
\]
gesetzt ist. Die Argumente \(w\) aber entstanden aus den \(v\) mit Hülfe der Gleichung
\[
w_h=v_h-\int_{\mu}^{\xi}du_h+K_h
=\sum_{i=1}^{i=p}\int_{c^{(i)}}^{x^{(i)}}du_h-\int_{\mu}^{\xi}du_h+K_h.
\]
Die Constanten \(K_h\) endlich sind dadurch gegeben, dass
\[
2K=\int_{\mu}^{\xi}du_h+\int_{\mu}^{\eta}du_h
-\sum_{i=1}^{i=p}\int_{c^{(i)}}^{x^{(i)}}du_h
-\sum_{i=1}^{i=p}\int_{c^{(i)}}^{y^{(i)}}du_h,
\]
wo die \(x,y\) mit den Doppel- und Rückkehrpuncten und mit den übrigen Puncten der Geraden \(\xi\eta\) auf einer Curve \((n-2)^{\text{ter}}\) Ordnung liegen sollten. Lassen wir nun diese Curve in die in \(\S\,47\) erwähnte Berührungscurve übergehen, bei welcher \(\xi,\eta\) in einen Punct \(\mu\), die \(x\) mit den \(y\) paarweise in Puncte \(\varepsilon\) zusammenfielen; Puncte, die mit den Puncten \(\xi,\eta,c,\mu,x,y\) durch Wege zusammenhängen, welche in die hier auftretenden Integrationswege überführbar sind. Alsdann haben wir
\[
K_h=-\sum_{i=1}^{i=p}\int_{c^{(i)}}^{\varepsilon^{(i)}}du_h,
\]
und wir können also schreiben:
\[
\tag{2}w_h=\sum_{i=1}^{i=p}\int_{\varepsilon^{(i)}}^{x^{(i)}}du_h-\int_{\mu}^{\xi}du_h.
\]

Halten wir den Punct \(\mu\) fest, so ist die Aufgabe, die Puncte \(\varepsilon\) zu bestimmen, keineswegs eindeutig lösbar. Die Puncte \(\varepsilon\) sind die Berührungspuncte einer Curve \((n-2)^{\text{ter}}\) Ordnung, welche durch die Doppel- und Rückkehrpuncte geht, und durch die übrigen Schnittpuncte der im Puncte \(\mu\) an \(f=0\) gelegten Tangente. Wir wollen untersuchen, wie viel solcher Berührungscurven es giebt. Die \(\varepsilon\) gehörten einer beliebig gewählten Berührungscurve an; die Berührungspuncte einer andern dieser Curven bezeichnen wir als die Puncte \(\alpha\). Dann bestehen nach dem Abel'schen Theorem die Gleichungen
\[
2\sum_{i=1}^{i=p}\int_{\varepsilon^{(i)}}^{\alpha^{(i)}}du_h=0,
\]
wo, wenn die Wege von den \(\varepsilon\) zu den \(\alpha\) beliebig gedacht werden, rechts ein beliebiger Periodicitätsmodul eines Integrals erster Gattung hinzugefügt werden kann. Man kann also die \(\alpha\) definiren aus der Gleichung:
\[
\tag{3}\sum_{i=1}^{i=p}\int_{\varepsilon^{(i)}}^{\alpha^{(i)}}du_h=\frac12P_h^{\varkappa,\lambda}.
\]
Jedem System von Zahlen \(\varkappa,\lambda\) entspricht ein System der \(\alpha\), also eine Berührungscurve, und auch nur eine einzige. Unter diesen verschiedenen Systemen ist eines das System der \(\varepsilon\) selbst, nämlich dasjenige, für welches alle \(\varkappa,\lambda\) verschwinden. So oft nun zwei halbe Perioden \(\frac12P_h^{\varkappa,\lambda}\) sich nicht um ganze Perioden unterscheiden, müssen auch die entsprechenden Systeme der \(\alpha\) von einander verschieden sein; es giebt also so viel Berührungscurven, als Zahlensysteme \(\varkappa,\lambda\) aus den Zahlen 0 und 1 zusammengesetzt werden können, also \(2^{2p}\). Eine dieser Curven ist die ursprünglich betrachtete; eine zweite ist diejenige, deren Zahlen \(\varkappa,\lambda\) mit den in den \(\omega\) auftretenden übereinstimmen. Die Berührungspuncte der letztern wollen wir jetzt vorzugsweise als die Puncte \(\alpha\) bezeichnen.

Aus den Gleichungen (1), (2), (3) zusammen ergeben sich nun sofort die Werthe der \(\omega\) in der Form:
\[
\omega_h=\sum_{i=1}^{i=p}\int_{\alpha^{(i)}}^{x^{(i)}}du_h-\int_{\mu}^{\xi}du_h.
\]
Und man hat den Satz:

Die Argumente \(\omega\) der Function \(\Theta\) sind die von den \(v\) um Constante verschiedenen Integralsummen
\[
\omega_h=\sum_{i=1}^{i=p}\int_{\alpha^{(i)}}^{x^{(i)}}du_h-\int_{\mu}^{\xi}du_h,
\]
wo die \(\alpha\) die Berührungspuncte einer bestimmten unter den \(2^{2p}\) Curven \((n-2)^{\text{ter}}\) Ordnung sind, welche durch die Doppel- und Rückkehrpuncte gehen, so wie durch die \(n-2\) andern Schnittpuncte der in \(\mu\) an \(f=0\) gezogenen Tangente, und welche ausserdem die Curve \(f=0\) in \(p\) Puncten berühren.

Es ist zu bemerken, dass die \(2^{2p}\) Curven, welche verschiedenen Puncten \(\mu\) entsprechen, einander in bestimmter Weise zugeordnet sind, so dass, wenn man \(\mu\) bewegt und zugleich die \(\alpha\) sich continuirlich ändern lässt, der Werth von \(\omega\) keine Modification erfährt. Man braucht zu diesem Zwecke bei einem andern Puncte \(\nu\) statt der \(\alpha\) nur diejenigen Puncte \(\beta\) zu wählen, für welche
\[
\sum_{i=1}^{i=p}\int_{\alpha^{(i)}}^{\beta^{(i)}}du_h-\int_{\mu}^{\nu}du_h=0.
\]
Wählt man ein andres System der \(\beta\), so tritt rechts die Hälfte eines Periodicitätsmoduls hinzu.

\subsubsection*{\S\,57. Lösung des Umkehrproblems durch \(\Theta\)functionen.}

Wir wollen jetzt untersuchen, wie sich die Lösung des Umkehrproblems mit Hülfe der synektischen Function \(\Theta\) darstellt. Aus der Gleichung
\[
V(x^{(1)},x^{(2)}\ldots x^{(p)};\xi)
=e^{\frac12\sum\varkappa_h(\omega_h-\omega_h^\circ)}
\cdot\frac{\Theta(\omega_1,\omega_2,\ldots)}
{\Theta(\omega_1^\circ,\omega_2^\circ,\ldots)}
\]
ergiebt sich zunächst für die specielle Transcendente \(\frac12T_{\xi\eta}(x)\) der Ausdruck
\[
\begin{aligned}
\frac12T_{\xi\eta}(x)
&=\log\frac{V(x^{(1)},x^{(2)}\ldots x^{(p)};\eta)}
{V(x^{(1)},x^{(2)}\ldots x^{(p)};\xi)}\\
&=\frac12\sum\varkappa_h\int_{\eta}^{\xi}du_h
+\log\frac{\Theta\!\left(\sum_{i=1}^{i=p}\int_{\alpha^{(i)}}^{x^{(i)}}du_h-\int_{\mu}^{\eta}du_h\right)}
{\Theta\!\left(\sum_{i=1}^{i=p}\int_{\alpha^{(i)}}^{x^{(i)}}du_h-\int_{\mu}^{\xi}du_h\right)},
\end{aligned}
\]
wo bei den \(\Theta\)functionen immer nur ein Argument, aber vollständig ausgeschrieben, gesetzt ist. Sodann aber findet man für die allgemeine Transcendente die Formel:
\[
\tag{1}
\frac12T_{\xi\eta}\binom{x}{z}
=\frac12T_{\xi\eta}(x)-\frac12T_{\xi\eta}(z)
=\log\frac{
\Theta\!\left(\sum_{i=1}^{i=p}\int_{\alpha^{(i)}}^{x^{(i)}}du_h-\int_{\mu}^{\eta}du_h\right)
\Theta\!\left(\sum_{i=1}^{i=p}\int_{\alpha^{(i)}}^{z^{(i)}}du_h-\int_{\mu}^{\xi}du_h\right)}
{\Theta\!\left(\sum_{i=1}^{i=p}\int_{\alpha^{(i)}}^{x^{(i)}}du_h-\int_{\mu}^{\xi}du_h\right)
\Theta\!\left(\sum_{i=1}^{i=p}\int_{\alpha^{(i)}}^{z^{(i)}}du_h-\int_{\mu}^{\eta}du_h\right)}.
\]
Und für das Umkehrproblem hat man die Formel:
\[
\tag{2}
\prod_{i=1}^{i=p}
\frac{\left(\frac{\varphi}{\psi}\right)_{x^{(i)}}-\lambda}
{\left(\frac{\varphi}{\psi}\right)_{c^{(i)}}-\lambda}
=
\prod_{\xi\eta}
\frac{
\Theta\!\left(\sum_{i=1}^{i=p}\int_{\alpha^{(i)}}^{x^{(i)}}du_h-\int_{\mu}^{\xi}du_h\right)
\Theta\!\left(\sum_{i=1}^{i=p}\int_{\alpha^{(i)}}^{c^{(i)}}du_h-\int_{\mu}^{\eta}du_h\right)}
{\Theta\!\left(\sum_{i=1}^{i=p}\int_{\alpha^{(i)}}^{x^{(i)}}du_h-\int_{\mu}^{\eta}du_h\right)
\Theta\!\left(\sum_{i=1}^{i=p}\int_{\alpha^{(i)}}^{c^{(i)}}du_h-\int_{\mu}^{\xi}du_h\right)}.
\]
Man sieht, dass die beiden letzten, dem Umkehrprobleme wesentlich zugehörigen Formeln von den Zahlen \(\varkappa\) vollkommen frei geworden sind, welche von der zufälligen Anordnung der Transcendente \(\frac12T_{\xi\eta}(x)\) und der Function \(V\) herrührten. Aber diese Gleichungen führen naturgemäss darauf, das Umkehrproblem, welches im sechsten Abschnitte auseinandergesetzt wurde, auf das speciellere Umkehrproblem zurückzuführen, in welchem die \(c\) nicht beliebig gegebene Puncte, sondern gerade das Punktsystem der \(\alpha\) darstellen, welches einem gegebenen Puncte \(\mu\) zugehört, während wir letzteren als beliebig und numerisch gegeben betrachten können.

Wir haben dann nur, um zu dem frühern Umkehrproblem zurückzukehren, die neuen Werthe der \(v\) um constante Grössen zu vermehren.

Setzen wir also
\[
\sum_{i=1}^{i=p}\int_{\alpha^{(i)}}^{x^{(i)}}du_h=v_h\qquad(h=1,2\ldots p),
\]
so ist jetzt die Lösung des Umkehrproblems in der Formel gegeben:
\[
\tag{3}
\prod_{i=1}^{i=p}\frac{\left(\frac{\varphi}{\psi}\right)_{x^{(i)}}-\lambda}
{\left(\frac{\varphi}{\psi}\right)_{\alpha^{(i)}}-\lambda}
=
\prod_{\xi\eta}
\frac{\Theta\!\left(v_h-\int_{\mu}^{\xi}du_h\right)\Theta\!\left(\int_{\mu}^{\eta}du_h\right)}
{\Theta\!\left(v_h-\int_{\mu}^{\eta}du_h\right)\Theta\!\left(\int_{\mu}^{\xi}du_h\right)}.
\]
Wir erhalten zugleich für die Transcendente \(T_{\xi\eta}\) den Ausdruck:
\[
T_{\xi\eta}\binom{x}{\alpha}
=\log
\frac{\Theta\!\left(v_h-\int_{\mu}^{\eta}du_h\right)\Theta\!\left(\int_{\mu}^{\xi}du_h\right)}
{\Theta\!\left(v_h-\int_{\mu}^{\xi}du_h\right)\Theta\!\left(\int_{\mu}^{\eta}du_h\right)}.
\]
Wir können diese Transcendente jetzt als Function der \(v\) durch \(T_{\xi\eta}(v)\) bezeichnen, so dass man die Definitionsgleichung hat:
\[
\tag{4}
T_{\xi\eta}(v)=\log
\frac{\Theta\!\left(v_h-\int_{\mu}^{\eta}du_h\right)\Theta\!\left(\int_{\mu}^{\xi}du_h\right)}
{\Theta\!\left(v_h-\int_{\mu}^{\xi}du_h\right)\Theta\!\left(\int_{\mu}^{\eta}du_h\right)}.
\]

Aus den Eigenschaften der Function \(\Theta\) erkennt man sofort, dass wirklich die durch die Gleichungen (3), (4) dargestellten Functionen diejenigen Periodicitätseigenschaften besitzen, welche denselben oben (im 6\textsuperscript{ten} Abschnitte) beigelegt wurden. Aus der Eigenschaft, dass, wenn
\[
v'_h=v_h+2m_h\pi\sqrt{-1}+a_{1h}q_1+a_{2h}q_2\cdots+a_{ph}q_h
\qquad(h=1,2\ldots p)
\]
gesetzt wird,
\[
\Theta(v')=e^{\sum q_hv_h+\frac12\sum q_hq_ka_{hk}}\Theta(v),
\]
folgt erstens, dass die Functionen (3) sich in keiner Weise ändern, wenn die darin auftretenden Integrale um Periodicitätsmoduln vermehrt werden, dass dieselben also wirklich \(2p\)fach periodische Functionen sind. Für die Function \(T_{\xi\eta}(v)\) aber erhält man die Gleichung:
\[
T_{\xi\eta}(v')=T_{\xi\eta}(v)+\sum q_h\int_{\eta}^{\xi}du_h,
\]
was die im 6\textsuperscript{ten} Abschnitte angegebene Periodicitätseigenschaft ist.

\subsubsection*{\S\,58. Ein Integral dritter Gattung und das Doppelverhältniss \(\left(\frac{\varphi}{\psi}\right)_{\xi}:\left(\frac{\varphi}{\psi}\right)_{\eta}\) durch \(\Theta\)functionen dargestellt.}

Aus den Formeln (1), (2) kann man noch einige andere interessante Gleichungen ableiten. Lässt man in (1) die \(x,z\) paarweise zusammenfallen, so dass nur ein einziges \(x\) von dem entsprechenden \(z\) verschieden ist, etwa \(x^{(1)}\) von \(z^{(1)}\). Setzen wir dann
\[
\int_{\alpha^{(1)}}^{z^{(1)}}du_h+\int_{\alpha^{(2)}}^{x^{(2)}}du_h+\int_{\alpha^{(3)}}^{x^{(3)}}du_h\cdots+\int_{\alpha^{(p)}}^{x^{(p)}}du_h=m_h,
\]
so findet man ein einziges Integral dritter Gattung durch \(\Theta\)functionen ausgedrückt mit Hülfe der Formel:
\[
\tag{5}
\int_{z^{(1)}}^{x^{(1)}}d\Pi_{\xi\eta}
=\log\frac{
\Theta\!\left(\int_{\eta}^{x^{(1)}}du_h+m_h\right)
\Theta\!\left(\int_{\xi}^{z^{(1)}}du_h+m_h\right)}
{\Theta\!\left(\int_{\xi}^{x^{(1)}}du_h+m_h\right)
\Theta\!\left(\int_{\eta}^{z^{(1)}}du_h+m_h\right)}.
\]
Und ebenso erhält man aus der Formel (2), wenn man darin die \(c\) durch \(z\), \(p-1\) von diesen aber durch die \(x\) ersetzt, die Quotienten zweier algebraischer Functionen durch \(\Theta\)functionen ausgedrückt:
\[
\tag{6}
\frac{\left(\frac{\varphi}{\psi}\right)_{x^{(1)}}}
{\left(\frac{\varphi}{\psi}\right)_{z^{(1)}}}
=
\prod_{\xi\eta}
\frac{
\Theta\!\left(\int_{\xi}^{x^{(1)}}du_h+m_h\right)
\Theta\!\left(\int_{\eta}^{z^{(1)}}du_h+m_h\right)}
{\Theta\!\left(\int_{\eta}^{x^{(1)}}du_h+m_h\right)
\Theta\!\left(\int_{\xi}^{z^{(1)}}du_h+m_h\right)},
\]
wobei die Puncte \(\xi\) die Schnittpuncte von \(\varphi=0\) mit \(f=0\), die Puncte \(\eta\) die Schnittpuncte von \(\psi=0\) mit \(f=0\) bedeuten.

Es ist bemerkenswerth dass diese Ausdrücke nothwendig von den \(p-1\) Puncten \(x^{(2)},x^{(3)}\ldots x^{(p)}\) unabhängig sind, welche in den \(p\) Grössen \(m\) auftreten. Es folgt daraus, dass diese \(p\) Grössen selbst nur in einer Combination die rechten Theile der obigen Gleichungen beeinflussen können, da nur eine Combination der \(m\) existiren kann, welche von den \(x^{(2)},x^{(3)}\ldots x^{(p)}\) unabhängig ist.

Wenn man in der rechten Seite der Gleichung (6) die Ausdrücke (5) substituirt, so wird man zum Abelschen Theorem zurückgeführt. Man kann umgekehrt von diesem ausgehend Gleichungen entwickeln, welche der Gleichung (6) analog sind, aber rechts eine kleinere Anzahl von Factoren enthalten.

Betrachten wir zu diesem Zwecke die Gleichung des Abel'schen Theorems:
\[
\log\left(\frac{\varphi}{\psi}\right)_{\xi}-\log\left(\frac{\varphi}{\psi}\right)_{\eta}
=\sum\int_{x^{(i)}}^{y^{(i)}}d\Pi_{\xi\eta},
\]
wo jetzt die \(y\) und \(x\) Paare von Schnittpuncten der Curven \(\psi=0\) und \(\varphi=0\) mit \(f=0\) bedeuten. Die Zahl der \(x\) und der \(y\) sei \(r\). Fügen wir rechts soviel Integrale mit zusammenfallenden obern und untern Grenzen hinzu, dass die Anzahl aller Integrale durch \(p\) theilbar wird (\(q\cdot p\)), so können wir nach Belieben immer \(p\) Integrale zu einer Transcendente \(T\) vereinigen, und wir erhalten die Gleichung:
\[
\log\left(\frac{\varphi}{\psi}\right)_{\xi}-\log\left(\frac{\varphi}{\psi}\right)_{\eta}
=\sum_{i=0}^{i=q-1}
T_{\xi\eta}\binom{y^{(ip+1)},y^{(ip+2)}\ldots y^{(i+1)p}}
{x^{(ip+1)},x^{(ip+2)}\ldots x^{(i+1)p}}.
\]
Setzen wir also:
\[
v_h^{(i)}=\int_{\alpha^{(1)}}^{x^{(ip+1)}}du_h+\int_{\alpha^{(2)}}^{x^{(ip+2)}}du_h\cdots+\int_{\alpha^{(p)}}^{x^{(i+1)p}}du_h,
\]
\[
V_h^{(i)}=\int_{\alpha^{(1)}}^{y^{(ip+1)}}du_h+\int_{\alpha^{(2)}}^{y^{(ip+2)}}du_h\cdots+\int_{\alpha^{(p)}}^{y^{(i+1)p}}du_h,
\]
so erhalten wir mit Benutzung der Formel (1), indem wir von den Logarithmen zu den Zahlen übergehen:
\[
\frac{\left(\frac{\varphi}{\psi}\right)_{\xi}}
{\left(\frac{\varphi}{\psi}\right)_{\eta}}
=
\prod_{i=0}^{i=q-1}
\frac{
\Theta\!\left(v_h^{(i)}-\int_{\mu}^{\xi}du_h\right)
\Theta\!\left(V_h^{(i)}-\int_{\mu}^{\eta}du_h\right)}
{\Theta\!\left(v_h^{(i)}-\int_{\mu}^{\eta}du_h\right)
\Theta\!\left(V_h^{(i)}-\int_{\mu}^{\xi}du_h\right)}.
\]
Die Bildung des Ausdrucks rechts kann auf unendlich viele verschiedene Arten geschehen, indem die hinzugefügten Puncte \(y^{(r+1)},x^{(r+1)},\ldots,y^{q p},x^{q p}\) beliebig sind, und ausserdem noch die Vereinigung der \(q p\) Integrale in Gruppen zu \(p\) auf mannigfaltige Weise abgeändert werden kann.

\subsubsection*{\S\,59. Quotienten alternirender Functionen durch \(\Theta\)functionen ausgedrückt.}

Die im vorigen \(\S\) aus dem Abelschen Theorem entwickelte Gleichung kann auch dazu benutzt werden, im Sinne des Umkehrproblems gewisse symmetrische Functionen von \(p\) Puncten \(x\), welche nicht mehr die oben angenommene einfache Form haben, direct durch \(\Theta\)functionen auszudrücken.

Es seien zur Bestimmung des Schnittpunctsystems einer durch Doppel- und Rückkehrpuncte gehenden Curve \(m^{\text{ter}}\) Ordnung \((m>n-3)\) mit der Curve \(f=0\) soviel Schnittpuncte gegeben, dass die Hinzufügung von \(p\) weiteren Schnittpuncten die übrigen bestimmen und also das ganze Schnittpunctsystem festlegen würde. Die Gleichung einer solchen Curve setzt sich aus den Gleichungen von \(p+1\) derselben linear zusammen, so dass
\[
\varphi=k_1\varphi_1+k_2\varphi_2\cdots+k_{p+1}\varphi_{p+1}=0
\]
die Gleichung einer Curve dieser Art wird, wenn die \(k\) beliebige Constanten, \(\varphi_1=0,\varphi_2=0\ldots\varphi_{p+1}=0\) aber specielle Curven sind, welche denselben Bedingungen genügen, und zwischen deren linken Theilen keine lineare Relation besteht. Soll nun die Curve \(\varphi=0\) durch \(p\) variable Puncte \(x^{(1)},x^{(2)}\ldots x^{(p)}\) gehen, so werden die Factoren \(k\) durch lineare Gleichungen bestimmt, und man erhält, wenn \(X\) den laufenden Punct bezeichnet:
\[
\tag{1}
\varphi(X)=
\begin{vmatrix}
\varphi_1(X) & \varphi_2(X) & \cdots & \varphi_{p+1}(X)\\
\varphi_1(x^{(1)}) & \varphi_2(x^{(1)}) & \cdots & \varphi_{p+1}(x^{(1)})\\
\varphi_1(x^{(2)}) & \varphi_2(x^{(2)}) & \cdots & \varphi_{p+1}(x^{(2)})\\
\vdots & \vdots & & \vdots\\
\varphi_1(x^{(p)}) & \varphi_2(x^{(p)}) & \cdots & \varphi_{p+1}(x^{(p)})
\end{vmatrix}.
\]
Diesen Ausdruck setzen wir in der letzten Formel des vorigen \(\S\) an Stelle von \(\varphi\), an Stelle von \(\psi\) aber dieselbe Function, in welcher nur an die Stelle der \(x\) constante Puncte \(a\) gesetzt sind. Jede dieser Curven schneidet \(f=0\), ausser in den gemeinsamen Schnittpuncten, erstlich beziehungsweise in den Puncten \(x,a\), sodann aber noch in je \(p\) andern Puncten \(X,A\), welche mit den \(x,a\) algebraisch verbunden sind, und zwischen denen ausserdem die \(p\) transcendenten Gleichungen bestehen:
\[
\tag{2}
\sum_{i=1}^{i=p}\int_{a^{(i)}}^{x^{(i)}}du_h+
\sum_{i=1}^{i=p}\int_{A^{(i)}}^{X^{(i)}}du_h=0
\qquad(h=1,2\ldots p).
\]
Die Gruppen der \(x\) des vorigen \(\S\) sind hier durch die Gruppe der \(x\) und die der \(X\) vertreten, ebenso die Gruppen der \(y\) durch die Gruppe der \(a\) und die der \(A\). Setzt man also
\[
\sum_{i=1}^{i=p}\int_{a^{(i)}}^{x^{(i)}}du_h=v_h,\qquad
\sum_{i=1}^{i=p}\int_{a^{(i)}}^{a^{(i)}}du_h=v_h^\circ,\qquad
\sum_{i=1}^{i=p}\int_{a^{(i)}}^{A^{(i)}}du_h=v'_h,
\]
so ergiebt sich aus (2):
\[
\sum_{i=1}^{i=p}\int_{a^{(i)}}^{X^{(i)}}du_h=v_h^\circ+v'_h-v_h,
\]
und man erhält daher aus der letzten Gleichung des vorigen \(\S\):
\[
\tag{3}
\begin{aligned}
&\frac{\Sigma\pm\varphi_1(\xi)\varphi_2(x^{(1)})\cdots\varphi_{p+1}(x^{(p)})}
{\Sigma\pm\varphi_1(\eta)\varphi_2(x^{(1)})\cdots\varphi_{p+1}(x^{(p)})}
\cdot
\frac{\Sigma\pm\varphi_1(\eta)\varphi_2(a^{(1)})\cdots\varphi_{p+1}(a^{(p)})}
{\Sigma\pm\varphi_1(\xi)\varphi_2(a^{(1)})\cdots\varphi_{p+1}(a^{(p)})}\\
&\quad=
\frac{
\Theta\!\left(v_h-\int_{\mu}^{\xi}du_h\right)
\Theta\!\left(v_h^\circ+v'_h-v_h-\int_{\mu}^{\xi}du_h\right)
\Theta\!\left(v_h^\circ-\int_{\mu}^{\eta}du_h\right)
\Theta\!\left(v'_h-\int_{\mu}^{\eta}du_h\right)}
{\Theta\!\left(v_h-\int_{\mu}^{\eta}du_h\right)
\Theta\!\left(v_h^\circ+v'_h-v_h-\int_{\mu}^{\eta}du_h\right)
\Theta\!\left(v_h^\circ-\int_{\mu}^{\xi}du_h\right)
\Theta\!\left(v'_h-\int_{\mu}^{\xi}du_h\right)}.
\end{aligned}
\]
Schreiben wir diese Gleichung in der Form:
\[
\tag{4}
\frac{\Sigma\pm\varphi_1(\xi)\varphi_2(x^{(1)})\cdots\varphi_{p+1}(x^{(p)})}
{\Sigma\pm\varphi_1(\xi)\varphi_2(a^{(1)})\cdots\varphi_{p+1}(a^{(p)})}
=
M\cdot
\frac{
\Theta\!\left(v_h-\int_{\mu}^{\xi}du_h\right)
\Theta\!\left(v_h^\circ+v'_h-v_h-\int_{\mu}^{\xi}du_h\right)}
{\Theta\!\left(v_h^\circ-\int_{\mu}^{\xi}du_h\right)
\Theta\!\left(v'_h-\int_{\mu}^{\xi}du_h\right)},
\]
so ist \(M\) offenbar von \(\xi\) unabhängig, da es sich nach der obigen Gleichung nicht ändert, wenn man \(\xi\) durch \(\eta\) ersetzt. Bezeichnen wir nun für den Augenblick durch \(R_1,R_2\ldots R_{p+1},R_1^\circ,R_2^\circ\ldots R_{p+1}^\circ\) die aus den \(\varphi_k(x^{(i)})\) und \(\varphi_k(a^{(i)})\) zusammengesetzten Unterdeterminanten, so dass
\[
\varphi=R_1\varphi_1(\xi)+R_2\varphi_2(\xi)\cdots+R_{p+1}\varphi_{p+1}(\xi),
\]
\[
\psi=R_1^\circ\varphi_1(\xi)+R_2^\circ\varphi_2(\xi)\cdots+R_{p+1}^\circ\varphi_{p+1}(\xi),
\]
so kann man die Gleichung (4) als lineare Gleichung mit den Unbekannten \(R_1,R_2\ldots R_{p+1},M\) ansehen. Solcher Gleichungen bilden wir uns \(p+1\) verschiedene, indem wir der Reihe nach \(p+1\) verschiedene constante Puncte \(\xi^{(1)},\xi^{(2)}\ldots\xi^{(p+1)}\) an Stelle von \(\xi\) einführen. Ist dann \(\Delta_i^{(h)}\) die Unterdeterminante der Determinante
\[
\Sigma\pm\varphi_1(\xi^{(1)})\varphi_2(\xi^{(2)})\cdots\varphi_{p+1}(\xi^{(p+1)})
\]
nach dem Elemente \(\varphi_i(\xi^{(h)})\), so ergiebt sich aus der Auflösung jenes Systems die Gleichung:
\[
mR_i=
\sum_{h=1}^{h=p+1}
\left[R_1^\circ\varphi_1(\xi^{(h)})+R_2^\circ\varphi_2(\xi^{(h)})\cdots+
R_{p+1}^\circ\varphi_{p+1}(\xi^{(h)})\right]
\]
\[
\cdot\Delta_i^{(h)}
\frac{
\Theta\!\left(v-\int_{\mu}^{\xi^{(h)}}du\right)
\Theta\!\left(v^\circ+v'-v-\int_{\mu}^{\xi^{(h)}}du\right)}
{\Theta\!\left(v^\circ-\int_{\mu}^{\xi^{(h)}}du\right)
\Theta\!\left(v'-\int_{\mu}^{\xi^{(h)}}du\right)},
\]
wo \(m\) ein unbestimmter Factor ist. Man hat also hiedurch die Verhältnisse der \(p+1\) alternirenden Functionen \(R_1,R_2\ldots R_{p+1}\) durch \(\Theta\) Functionen ausgedrückt, deren Argumente sich von den Variabeln des Umkehrproblems nur um Constante unterscheiden.
% --- Continuation: source PDF pp. 228--243, printed pp. 206--221. ---
\clearpage

\begin{center}
{\bfseries Neunter Abschnitt.}\par\vspace{0.6em}
{\Large\bfseries Die Function \(\Theta\).}
\end{center}

\subsection*{\S\,60. Verschwinden der Function \(\Theta\). Unbestimmter Fall des Umkehrproblems.}

Die Function \(\Theta\), welche oben aus der Function \(V\) abgeleitet wurde, besitzt die Eigenschaften derselben; sie wird für endliche Werthe der \(\omega\) nicht unendlich, und sie verschwindet, wenn die \(\omega\) bestimmte Bedingungen erfüllen, aber in keinem andern Falle.

Setzen wir wie oben
\[
\tag{1}
\omega_h=\sum_{i=1}^{i=p}\int_{\alpha^{(i)}}^{x^{(i)}}du_h-\int_{\mu}^{\xi}du_h .
\]
Sind die \(\omega\), die Argumente von \(\Theta\), gegeben, so ist in diesen Gleichungen \(\xi\) noch willkürlich; man kann es mit \(\mu\) zusammenfallen, oder umgekehrt \(\mu\) gegen den variabeln Punct \(\xi\) hinrücken lassen, wodurch das letzte Integral zum Verschwinden gebracht wird. Hat man über \(\mu\), \(\xi\), also auch über die \(\alpha\) verfügt, so sind die \(x\) im Allgemeinen völlig und eindeutig bestimmt.

Die Function \(\Theta\), wie \(V\), verschwindet nun immer in den folgenden beiden Fällen:
\[
\begin{array}{ll}
1) & \text{wenn ein Punct }x\text{ mit }\xi\text{ zusammenfällt;}\\
2) & \text{wenn die Puncte }x\text{ auf einer durch die Doppel- und Rückkehrpuncte gelegten}\\
   & \text{Curve }(n-3)^{\mathrm{ter}}\text{ Ordnung liegen;}
\end{array}
\]
und zwar kann \(\Theta\) in keinem andern Falle verschwinden.

In beiden Fällen enthalten die \(\omega\) dann nur noch \(p-1\) Variabeln; denn im ersten Fall hebt sich \(\xi\) gegen das mit ihm zusammenfallende \(x\) direct auf; im zweiten Falle kann man die Integralsummen, welche die Puncte \(x^{(1)},x^{(2)}\ldots x^{(p)}\) zu obern Grenzen haben, durch andere ersetzen, welche die mit jenen auf einer Curve \((n-3)^{\mathrm{ter}}\) Ordnung gelegenen Puncte \(x^{(p+1)},x^{(p+2)}\ldots x^{(2p-2)}\) zu obern Grenzpuncten haben, sodass diese mit \(\xi\) zusammen als Variable übrig bleiben. In beiden Fällen findet also zwischen den \(\omega\) eine Bedingung statt.

Dass übrigens die beiden Fälle auf einander zurückkommen und also nicht wesentlich verschieden sind, ist schon oben benutzt worden. Denn da \(\Theta(-\omega)=\Theta(\omega)\), so ist ebenso wie die Gleichung (1) auch die folgende Ausdrucksweise berechtigt:
\[
\tag{2}
-\omega_h=\sum_{i=1}^{i=p}\int_{\alpha^{(i)}}^{y^{(i)}}du_h-\int_{\mu}^{\eta}du_h .
\]
Aus der Gleichung
\[
\sum_{i=1}^{i=p}\left(\int_{\alpha^{(i)}}^{x^{(i)}}du_h+
\int_{\alpha^{(i)}}^{y^{(i)}}du_h\right)-
\left(\int_{\mu}^{\xi}du_h+\int_{\mu}^{\eta}du_h\right)=0
\]
folgt nun, dass die \(x,y\) die Schnittpuncte einer Curve \((n-2)^{\mathrm{ter}}\) Ordnung mit \(f=0\) sind, welche durch die Doppel- und Rückkehrpuncte geht, so wie durch die \(p-2\) Puncte, in welchen die Gerade \(\xi\eta\) noch trifft. Liegen nun die \(x\) auf einer Curve \((n-3)^{\mathrm{ter}}\) Ordnung, so muss diese mit der Geraden \(\xi\eta\) zusammen die Stelle der Curve \((n-2)^{\mathrm{ter}}\) Ordnung einnehmen, es muss also ein Punct \(y\) in \(\xi\), ein anderer in \(\eta\) fallen. Während also bei den \(x\) der Fall (2) eintritt, tritt bei den \(y\) der Fall (1) ein, und umgekehrt, wenn bei den \(y\) der Fall (2) eintritt, erfüllen die \(x\) die Bedingungen des Falls (1).

Bei den oben gemachten Anwendungen des Abel'schen Satzes, bei welchen der Ausdruck
\[
\frac{\left(\dfrac{\varphi}{\psi}\right)_{\xi}}
     {\left(\dfrac{\varphi}{\psi}\right)_{\eta}}
\]
als Quotient von \(\Theta\)producten dargestellt wurde, können wir den Fall, wo einige der \(\Theta\) verschwinden könnten, immer umgehen. Denn wir verwandelten den Logarithmus des obigen algebraischen Ausdrucks in eine Summe von Integralen dritter Gattung; diesen können wir andere mit zusammenfallenden obern und untern Grenzen immer so zuordnen, dass sich Gruppen von je \(p\) Integralen bilden lassen, deren entsprechende \(\Theta\)quotienten nicht unbestimmt werden.

Dagegen tritt allerdings bei dem Umkehrproblem der wirklich unbestimmte Fall ein, dass die \(x\) unbestimmt werden, sobald die Integralsummen erster Gattung einer gewissen Bedingung unterliegen; derjenigen nämlich, welche aussagt, dass die \(x\) auf einer Curve \((n-3)^{\mathrm{ter}}\) Ordnung liegen, welche durch Doppel- und Rückkehrpuncte geht. Dies zeigt sich in den Formeln des Umkehrproblems dadurch an, dass die sämmtlichen \(\Theta\)functionen in Zähler und Nenner gleichzeitig verschwinden.

Man kann aber für diesen Fall das gegebene Umkehrproblem immer leicht auf ein anderes zurückführen, welches bestimmt ist. Denn obwohl die \(x\) selbst hier nothwendig unbestimmt werden, sind es doch keineswegs die \(p-2\) Puncte \(y\), in welchen \(f=0\) von der durch die \(x\) gehenden Curve \((n-3)^{\mathrm{ter}}\) Ordnung noch geschnitten wird. Fügen wir diesen \(p-2\) unbekannten Puncten \(y\) zwei andere beliebig auf \(f=0\) gewählte bekannte Puncte \(y^{(p-1)},y^{(p)}\) hinzu, welche nur mit den erstern nicht auf einer Curve \((n-3)^{\mathrm{ter}}\) Ordnung liegen. Alsdann ist
\[
\sum_{i=1}^{i=p}\int_{\alpha^{(i)}}^{x^{(i)}}du_h=v_h,
\]
aber auch
\[
\tag{3}
\sum_{i=1}^{i=p}\int_{\alpha^{(i)}}^{y^{(i)}}du_h=k_h-v_h,
\]
wo \(k_h\) die aus dem Abel'schen Theorem zu bestimmende Constante
\[
k_h=\sum_{i=1}^{i=p}\left(\int_{\alpha^{(i)}}^{x^{(i)}}du_h+
\int_{\alpha^{(i)}}^{y^{(i)}}du_h\right)
\]
bedeutet. Indem man die Gleichungen (3) zur Grundlage eines neuen Umkehrproblems macht, erhält man (vgl. \S\,57. 3.)
\[
\prod_{i=1}^{i=p}
\frac{\left(\dfrac{\varphi}{\psi}\right)_{y^{(i)}}-\lambda}
     {\left(\dfrac{\varphi}{\psi}\right)_{\alpha^{(i)}}-\lambda}
=
\prod_{\xi\eta}
\frac{\Theta\left(k_h-v_h-\int_{\mu}^{\xi}du_h\right)\Theta\left(\int_{\mu}^{\eta}du_h\right)}
     {\Theta\left(k_h-v_h-\int_{\mu}^{\eta}du_h\right)\Theta\left(\int_{\mu}^{\xi}du_h\right)},
\]
wo nun in dem Product rechts zwei Factoren, welche den gegebenen \(y\) entsprechen, völlig bekannt sind. Indem man also diese Factoren absondert, und der Zahl \(\lambda\) verschiedene Werthe beilegt, kann man eine Gleichung \((p-2)^{\mathrm{ten}}\) Grades bilden, deren Wurzeln die unbekannten Werthe von \(\left(\dfrac{\varphi}{\psi}\right)_{y^{(i)}}\) sind.

Man kann endlich eine directe Methode verlangen, den Quotienten zweier verschwindenden \(\Theta\)functionen auf seinen wahren Werth zurückzuführen. Da oben der Fall (2) immer auf den Fall (1) zurückgeführt werden konnte, so beschränken wir uns darauf, anzunehmen, dass beide \(\Theta\)functionen, in denen wir den willkürlichen Punct \(\xi\) als gemeinsam voraussetzen dürfen, durch Uebereinstimmung dieses Punctes mit einem der Grenzpuncte verschwinden. Wir betrachten also den Ausdruck
\[
L=\lim_{y=\xi}
\frac{\Theta\left(\sum_{i=1}^{i=p-1}\int_{\alpha^{(i)}}^{x^{(i)}}du_h+
\int_{\alpha^{(p)}}^{y}du_h-\int_{\mu}^{\xi}du_h\right)}
{\Theta\left(\sum_{i=1}^{i=p-1}\int_{\alpha^{(i)}}^{z^{(i)}}du_h+
\int_{\alpha^{(p)}}^{y}du_h-\int_{\mu}^{\xi}du_h\right)}.
\]
Multipliciren wir denselben mit einem ähnlichen nicht unbestimmten \(\Theta\)quotienten, so haben wir nach \S\,57. 1.
\[
\log L+
\log\frac{\Theta\left(\sum_{i=1}^{i=p-1}\int_{\alpha^{(i)}}^{z^{(i)}}du_h+
\int_{\alpha^{(p)}}^{y}du_h-\int_{\mu}^{\eta}du_h\right)}
{\Theta\left(\sum_{i=1}^{i=p-1}\int_{\alpha^{(i)}}^{x^{(i)}}du_h+
\int_{\alpha^{(p)}}^{y}du_h-\int_{\mu}^{\eta}du_h\right)}
= -\sum_{i=1}^{i=p-1}\int_{z^{(i)}}^{x^{(i)}}d\Pi_{\xi\eta}.
\]
Dieser Ausdruck ist, wie die rechte Seite lehrt, von der Lage des Punctes \(y\) abhängig; für die linke Seite können wir also den Ausdruck setzen, welcher entsteht, wenn wir \(y\) durch einen beliebig auf \(f=0\) gewählten Punct \(t\) ersetzen. Und demnach erhalten wir für \(L\) den Quotienten zweier Producte von je drei \(\Theta\)functionen; die Argumente der Functionen sind
\[
\begin{array}{c@{\quad}c}
\text{im Zähler:} & \text{im Nenner:}\\[0.4em]
\displaystyle \sum_{i=1}^{i=p-1}\int_{\alpha^{(i)}}^{x^{(i)}}du_h+\int_{\alpha^{(p)}}^{t}du_h-\int_{\mu}^{\xi}du_h
&
\displaystyle \sum_{i=1}^{i=p-1}\int_{\alpha^{(i)}}^{z^{(i)}}du_h+\int_{\alpha^{(p)}}^{t}du_h-\int_{\mu}^{\eta}du_h
\\[1.0em]
\displaystyle \sum_{i=1}^{i=p-1}\int_{\alpha^{(i)}}^{z^{(i)}}du_h+\int_{\alpha^{(p)}}^{t}du_h-\int_{\mu}^{\eta}du_h
&
\displaystyle \sum_{i=1}^{i=p-1}\int_{\alpha^{(i)}}^{z^{(i)}}du_h+\int_{\alpha^{(p)}}^{t}du_h-\int_{\mu}^{\xi}du_h
\\[1.0em]
\displaystyle \sum_{i=1}^{i=p-1}\int_{\alpha^{(i)}}^{x^{(i)}}du_h+\int_{\alpha^{(p)}}^{y}du_h-\int_{\mu}^{\eta}du_h
&
\displaystyle \sum_{i=1}^{i=p-1}\int_{\alpha^{(i)}}^{x^{(i)}}du_h+\int_{\alpha^{(p)}}^{y}du_h-\int_{\mu}^{\eta}du_h .
\end{array}
\]
Der unbestimmte \(\Theta\)quotient ist also auf den Quotienten nicht verschwindender Producte zurückgeführt, und zwar enthält der neue Ausdruck die beiden willkürlich eingeführten Puncte \(t,\eta\).

Wenn nicht nur eines, sondern mehrere Punctepaare \(x,z\) mit dem Puncte \(\xi\) coincidiren, so ist der vorliegende Ausdruck abermals unbestimmt; aber durch dasselbe Verfahren kann man die Anzahl der Punctepaare, welche die Unbestimmtheit herbeiführen, fortgesetzt vermindern, und gelangt schliesslich zu bestimmten Ausdrücken. Man hat, wenn \(k\) Punctepaare \(x,z\) mit \(\xi\) zusammenfallen, dieses Verfahren \(k\) mal hinter einander anzuwenden. Wir können hieraus vermuthen, dass in diesem Falle die Function \(\Theta\) nicht einfach, sondern \(k\)fach verschwindet. Dies wird in der That durch den Anblick von \(\log\Theta\) bestätigt, wie dieser Logarithmus in der Function \(-U\) gegeben vorliegt. Wenn in dieser (\S\,49) \(x^{(1)},x^{(2)}\ldots x^{(k)}\) sich dem gemeinsamen Endpuncte \(\xi\) nähern, so wird der unstetige Theil von \(-U\) nach den Eigenschaften des Integrals dritter Gattung:
\[
-\sum_{i=1}^{i=k}\log(\xi-x^{(i)})+
\sum_{i<h}\log(x^{(i)}-x^{(h)})-
\log\varphi(x^{(1)},x^{(2)}\ldots x^{(p)}).
\]
Aber die letzten beiden Terme bleiben zusammen endlich, auch wenn die \(x\) sich einander nähern, daher nähert sich \(U\) dem Werthe
\[
-\log\left[(\xi-x^{(1)})(\xi-x^{(2)})\cdots(\xi-x^{(k)})\right],
\]
oder es nähert sich \(\Theta\) dem Ausdruck
\[
\operatorname{Const.}\,(\xi-x^{(1)})(\xi-x^{(2)})\cdots(\xi-x^{(k)}),
\]
d. h. es verschwindet \(k\)fach, und mit ihm verschwinden seine \(k-1\) ersten Differentialquotienten nach \(\xi\).

Aus dem oben angeführten Grenzwerth von \(-U\) ersieht man übrigens, dass die Function \(\Theta\) auch dadurch in höherer Ordnung verschwindet, dass nicht nur einige der \(x\) in \(\xi\) fallen, sondern dass auch die Puncte \(x\) während dieser Annäherung fortfahren, auf einer Curve \((n-3)^{\mathrm{ter}}\) Ordnung zu liegen. Zu der Ordnung des verschwindenden Products \(\xi-x^{(1)}\cdot\xi-x^{(2)}\ldots\) ist dann die Ordnung des verschwindenden Ausdrucks
\[
\frac{\varphi}{\Pi(x^{(i)}-x^{(h)})}
\]
zu addiren, welche im Allgemeinen die erste ist.

\subsection*{\S\,61. Besondere Fälle, in welchen die Function \(\Theta\) verschwindet.}

Wir können umgekehrt einen \(\Theta\)quotienten, welcher aus irgend welchen Gründen die Form \(\frac{0}{0}\) annimmt, immer auf solche \(\Theta\)quotienten zurückführen, welche nur deswegen unbestimmt werden, weil die obern Grenzpuncte \((x,z)\) auf Curven \((n-3)^{\mathrm{ter}}\) Ordnung liegen, welche durch die Doppel- und Rückkehrpuncte gehen. Hiebei aber können einige besondere Umstände eintreten, welche eine genauere Erörterung verdienen.

Im Allgemeinen bestimmen \(mn-p\) Schnittpuncte einer durch Doppel- und Rückkehrpuncte von \(f=0\) gehenden Curve \(m^{\mathrm{ter}}\) Ordnung die \(p\) übrigen. Nur wenn \(m=n-3\), oder kleiner ist, tritt eine Ausnahme ein; wir können uns sogar auf die Betrachtung von \(m=n-3\) beschränken, indem wir bei kleinern Werthen von \(m\) durch Zufügung irrelevanter Curven uneigentliche Curven \((n-3)^{\mathrm{ter}}\) Ordnung erzeugen.

Nimmt man irgend \(p-1\) Puncte \(x^{(1)},x^{(2)}\ldots x^{(p-1)}\) beliebig auf der Curve \(f=0\) an, und lässt durch sie, sowie durch Doppel- und Rückkehrpuncte eine Curve \((n-3)^{\mathrm{ter}}\) Ordnung gehen, so ist diese bestimmt, und giebt \(p-1\) weitere Schnittpuncte auf \(f=0\). Wenn man \(x^{(p-1)}\) bewegt, so werden sich im Allgemeinen auch alle \(p-1\) weitern Schnittpuncte mit bewegen. Aber es kann geschehen, dass bei specieller Lage der Puncte \(x^{(1)},x^{(2)}\ldots x^{(p-2)}\) die Bewegung von \(x^{(p-1)}\) auf einen der \(p-1\) übrigen Schnittpuncte keinen Einfluss hat, dass also jene \(p-2\) ersten Puncte ohne Zufügung des \((p-1)^{\mathrm{ten}}\) Punctes schon einen weitern Schnittpunct \(x^{(p)}\) bestimmen, so dass eine besondere Lage von \(x^{(p-1)}\) dann nur noch die \(p-2\) fehlenden Schnittpuncte fixirt. Da hierin nur eine Bedingung liegt, so können wir die Puncte \(x^{(1)},x^{(2)}\ldots x^{(p-3)}\) beliebig annehmen, und die Frage stellen, wie \(x^{(p-2)}\) zu bestimmen ist, damit alle durch \(x^{(1)},x^{(2)}\ldots x^{(p-2)}\) gehenden Curven der gedachten Art auch noch durch einen weitern Punct \(x^{(p)}\) gehen. Legen wir nun irgend drei solcher Curven durch \(x^{(1)},x^{(2)}\ldots x^{(p-3)}\), und sind ihre Gleichungen \(\varphi_1=0,\varphi_2=0,\varphi_3=0\), so ist, wenn nur zwischen den Ausdrücken \(\varphi_1,\varphi_2,\varphi_3\) keine lineare Relation besteht, jede vierte Curve darstellbar durch die Form
\[
\alpha_1\varphi_1+\alpha_2\varphi_2+\alpha_3\varphi_3=0,
\]
wo die \(\alpha\) veränderliche Parameter bedeuten. Dieselben werden auf ein Verhältniss reducirt, wenn man diejenigen Curven betrachtet, welche durch den gesuchten Punct \(x^{(p-2)}\) gehen, für welche also
\[
\tag{1}
\alpha_1\varphi_1(x^{(p-2)})+\alpha_2\varphi_2(x^{(p-2)})+\alpha_3\varphi_3(x^{(p-2)})=0.
\]
Es wurde aber verlangt, dass alle diese Curven durch einen bestimmten andern Punct \(x^{(p)}\) gehen sollten, d. h. dass aus der Gleichung (1) immer die andere folgen solle:
\[
\tag{2}
\alpha_1\varphi_1(x^{(p)})+\alpha_2\varphi_2(x^{(p)})+\alpha_3\varphi_3(x^{(p)})=0.
\]
Dies ist nur möglich, wenn die drei Gleichungen bestehen:
\[
\tag{3}
\begin{cases}
\varphi_1(x^{(p)})=\lambda\varphi_1(x^{(p-2)}),\\
\varphi_2(x^{(p)})=\lambda\varphi_2(x^{(p-2)}),\\
\varphi_3(x^{(p)})=\lambda\varphi_3(x^{(p-2)}).
\end{cases}
\]
Aus diesen Gleichungen bestimmen sich \(\lambda\) und die Puncte \(x^{(p)},x^{(p-2)}\), welche gesucht wurden.

Die Gleichungen (3) zeigen eine Reciprocität zwischen den Puncten \(x^{(p)},x^{(p-2)}\). Es giebt also eine Anzahl von Punctepaaren, welche die Eigenschaft haben, mit gegebenen \(p-3\) Puncten zusammen die gemeinsamen Schnittpuncte eines Systems von Curven \((n-3)^{\mathrm{ter}}\) Ordnung mit \(f=0\) zu bilden, also mit diesen zusammen eine Curve \((n-3)^{\mathrm{ter}}\) Ordnung ausnahmsweise nicht zu bestimmen.

Die Anzahl solcher Punctepaare, welche einem System von \(p-3\) gegebenen Puncten zugehören, ist sehr leicht mit Hülfe der Untersuchungen zu bestimmen, welche im dritten Abschnitte geführt sind. Die Gleichungen nämlich
\[
\tag{4}
\begin{cases}
\mu y_1=\varphi_1(x_1,x_2,x_3),\\
\mu y_2=\varphi_2(x_1,x_2,x_3),\\
\mu y_3=\varphi_3(x_1,x_2,x_3)
\end{cases}
\]
sind die Gleichungen für eine eindeutige Transformation \((n-3)^{\mathrm{ter}}\) Ordnung. Die Curven \(\varphi=0\) haben sämmtliche Doppel- und Rückkehrpuncte gemein, so wie die \(p-3\) willkürlich gewählten Puncte \(x^{(1)},x^{(2)}\ldots x^{(p-3)}\). Die Punctepaare \(x^{(p)},x^{(p-2)}\) sind also keine andern, als diejenigen, welche bei dieser eindeutigen Substitution zu neuen Doppel- oder Rückkehrpuncten vereinigt werden. Daher ist auch die Zahl solcher Punctepaare leicht bestimmt; sie ist die Zahl der Doppel- und Rückkehrpuncte der neuen Curve. Da nun (vgl. \S\,18) die Ordnung der neuen Curve \(p+1\) wird, so ist die Anzahl ihrer Doppel- und Rückkehrpuncte
\[
d^{(1)}+r^{(1)}=\frac{p\cdot p-3}{2}.
\]
\[
\text{Es giebt also }\frac{p\cdot p-3}{2}\text{ Punctepaare, welche mit }p-3
\]
gegebenen Puncten zusammen Systeme von solchen \(p-1\) Puncten liefern, welche eine durch die Doppel- und Rückkehrpuncte von \(f=0\) gehende Curve \((n-3)^{\mathrm{ter}}\) Ordnung nicht bestimmen.

Ebenso wie hier zu \(p-3\) gegebenen Puncten zwei andere gesucht werden, welche mit jenen zusammen nicht eine, sondern unendlich viele Curven \((n-3)^{\mathrm{ter}}\) Ordnung bestimmen, die sich durch die Werthe eines willkürlichen Parameters unterscheiden -- ebenso kann man die Frage behandeln, unter welchen Umständen schon einer geringern Zahl gegebener Puncte eine Zahl fester Puncte zugeordnet ist, welche zusammen ein Curvensystem \((n-3)^{\mathrm{ter}}\) Ordnung mit einer grössern Anzahl von Parametern bestimmen. Zunächst können wir fragen, wie \(p-3\) Puncte \(x^{(1)},x^{(2)}\ldots x^{(p-3)}\) liegen müssen, damit jede durch sie (und die Doppel- und Rückkehrpuncte) gelegte Curve \((n-3)^{\mathrm{ter}}\) Ordnung noch durch einen gemeinsamen Punct gehe. Zu diesem Zwecke denken wir uns den Puncten \(x^{(1)},x^{(2)}\ldots x^{(p-3)}\) zwei bewegliche Puncte \(x^{(p-2)},x^{(p-1)}\) hinzugefügt, und suchen die Bedingung auf, unter welcher bei Bewegung dieser beiden neuen Puncte einer der übrigen Schnittpuncte der Curve mit \(f=0\), \(x^{(p)}\), fest bleibt. Diese Forderung involvirt zwei Bedingungen. Sehen wir also nur \(x^{(1)},x^{(2)}\ldots x^{(p-5)}\) als willkürlich gegeben an, und fragen nun, wie \(x^{(p-4)}\) und \(x^{(p-3)}\) liegen müssen, damit alle durch sie gelegten Curven noch durch einen gemeinsamen Punct \(x^{(p)}\) gehen. Da hier vier Puncte weniger gegeben sind, als zur Festlegung einer Curve erforderlich sind, so kann man fünf Curven \(\varphi_1=0,\varphi_2=0\ldots\varphi_5=0\) legen, aus denen eine beliebige durch jene \(p-5\) und die Doppel- und Rückkehrpuncte gehende Curve \((n-3)^{\mathrm{ter}}\) Ordnung sich zusammensetzt. Die allgemeine Curve des durch diese Puncte bestimmten Systems hat also die Form:
\[
\alpha_1\varphi_1+\alpha_2\varphi_2\cdots+\alpha_5\varphi_5=0.
\]
Die Forderung ist, dass jede durch \(x^{(p-4)}\) und \(x^{(p-3)}\) gehende Curve des Systems auch durch \(x^{(p)}\) gehe; dass also aus den Gleichungen
\[
\begin{aligned}
&\alpha_1\varphi_1(x^{(p-4)})+\alpha_2\varphi_2(x^{(p-4)})\cdots+\alpha_5\varphi_5(x^{(p-4)})=0,\\
&\alpha_1\varphi_1(x^{(p-3)})+\alpha_2\varphi_2(x^{(p-3)})\cdots+\alpha_5\varphi_5(x^{(p-3)})=0
\end{aligned}
\]
von selbst die Gleichung folge:
\[
\alpha_1\varphi_1(x^{(p)})+\alpha_2\varphi_2(x^{(p)})\cdots+\alpha_5\varphi_5(x^{(p)})=0.
\]
Hierzu ist nöthig und hinreichend, dass die Functions-werthe der letzten Gleichung sich aus denen der ersten beiden linear zusammensetzen, dass also die Gleichungen bestehen:
\[
\varphi_i(x^{(p)})=\lambda\varphi_i(x^{(p-4)})+\mu\varphi_i(x^{(p-3)}),
\qquad i=1,2\ldots 5.
\]
Dies sind fünf Gleichungen, welche \(\lambda,\mu\) und die drei Puncte \(x^{(p-4)},x^{(p-3)},x^{(p)}\) bestimmen. Die Bestimmung der Anzahl ihrer Lösungen erfordert eine Erweiterung der im dritten Abschnitt ausgeführten Betrachtungen.

Allgemein kann man fragen, wann alle durch \(x^{(1)},x^{(2)}\ldots x^{(k)}\) gehenden Curven der gedachten Art noch durch einen gemeinsamen Punct \(x^{(p)}\) gehen. Um dies zu untersuchen, fügt man den fest gedachten Puncten \(x^{(1)},x^{(2)}\ldots x^{(k)}\) bewegliche Puncte \(x^{(k+1)},x^{(k+2)}\ldots x^{(p-1)}\) hinzu, \(p-k-1\) an der Zahl, und untersucht, wann die Bewegung derselben einen weitern Schnittpunct \(x^{(p)}\) der durch sie gelegten Curven mit \(f=0\) unverändert lässt. Dies involvirt \(p-k-1\) Bedingungen; wir sehen also zunächst nur \(x^{(1)},x^{(2)}\ldots x^{(2k-p+1)}\) Puncte als willkürlich gegeben an, und bestimmen die übrigen \(x^{(2k-p+2)},x^{(2k-p+3)}\ldots x^{(k)}\) so, dass diese Bedingungen befriedigt werden. Durch jene \(2k-p+1\) Puncte kann man nun \(2p-2k-1\) specielle Curven \((n-3)^{\mathrm{ter}}\) Ordnung \(\varphi_1=0,\varphi_2=0\ldots,\varphi_{2p-2k-1}=0\) legen, aus denen sich die allgemeinste Curve in der Form
\[
\alpha_1\varphi_1+\alpha_2\varphi_2\cdots+\alpha_{2p-2k-1}\varphi_{2p-2k-1}=0
\]
zusammensetzt. Wenn nun eine Curve dieses Systems durch die Puncte \(x^{(2k-p+2)},x^{(2k-p+3)}\ldots x^{(k)}\) geht, so soll sie von selbst durch einen weitern Punct \(x^{(p)}\) gehen, d. h. aus den Gleichungen
\[
\begin{aligned}
&\alpha_1\varphi_1(x^{(2k-p+2)})+\alpha_2\varphi_2(x^{(2k-p+2)})\cdots+
\alpha_{2p-2k-1}\varphi_{2p-2k-1}(x^{(2k-p+2)})=0,\\
&\alpha_1\varphi_1(x^{(2k-p+3)})+\alpha_2\varphi_2(x^{(2k-p+3)})\cdots+
\alpha_{2p-2k-1}\varphi_{2p-2k-1}(x^{(2k-p+3)})=0,\\
&\hfill\cdots\hfill,\\
&\alpha_1\varphi_1(x^{(k)})+\alpha_2\varphi_2(x^{(k)})\cdots+
\alpha_{2p-2k-1}\varphi_{2p-2k-1}(x^{(k)})=0
\end{aligned}
\]
soll die weitere Gleichung
\[
\alpha_1\varphi_1(x^{(p)})+\alpha_2\varphi_2(x^{(p)})\cdots+
\alpha_{2p-2k-1}\varphi_{2p-2k-1}(x^{(p)})=0
\]
folgen. Man muss daher die Gleichungen haben:
\[
\varphi_i(x^{(p)})=\lambda\varphi_i(x^{(2k-p+2)})+
\mu\varphi_i(x^{(2k-p+3)})\cdots+
\sigma\varphi_i(x^{(k)}),
\]
\[
i=1,2\ldots 2p-2k-1,
\]
welche gerade hinreichen, um die \(\lambda,\mu\ldots\sigma\), und ausserdem die Puncte \(x^{(2k-p+2)},x^{(2k-p+3)}\ldots x^{(k)}\) zu bestimmen.

Wenn die obern Grenzpuncte \(x^{(1)},x^{(2)}\ldots x^{(p)}\) in der \(\Theta\)function solche sind, für welche einer der hier behandelten speciellen Fälle eintritt, so verschwindet die \(\Theta\)function immer, und zwar, dem Fall (2) des \S\,60 entsprechend, unabhängig von \(\xi\). Aber sie verschwindet auch unabhängig von denjenigen Puncten \(x\), deren Lage im Vorigen unbestimmt blieb; also, wenn jede durch \(x^{(1)},x^{(2)}\ldots x^{(p-2)}\) gelegte Curve \((n-3)^{\mathrm{ter}}\) Ordnung auch durch \(x^{(p)}\) geht, unabhängig von der speciellen Lage des Punctes \(x^{(p-1)}\); wenn jede durch \(x^{(1)},x^{(2)}\ldots x^{(p-3)}\) gelegte Curve durch \(x^{(p)}\) geht, so verschwindet sie unabhängig von \(x^{(p-1)}\) und \(x^{(p-2)}\); allgemein: wenn jede durch \(x^{(1)},x^{(2)}\ldots x^{(k)}\) gelegte Curve durch \(x^{(p)}\) geht, so verschwindet sie unabhängig von der Lage der Puncte \(x^{(k+1)},x^{(k+2)}\ldots x^{(p-1)}\). Die Function \(U\) nähert sich dann der Function
\[
-\log\varphi(x^{(1)},x^{(2)}\ldots x^{(p)}),
\]
also \(\Theta\), bis auf eine Constante, dem Werthe
\[
\varphi(x^{(1)},x^{(2)}\ldots x^{(p)});
\]
und dieser Ausdruck verschwindet, aber so, dass dabei nicht alle Puncte \(x\) zur Mitwirkung kommen, sondern dass eine Reihe von Unterdeterminanten der Determinante \(\varphi\) verschwinden, welche nur einige der Puncte \(x\) enthalten.

Es ist übrigens leicht, die oben betrachteten besondern Curven \((n-3)^{\mathrm{ter}}\) Ordnung in den einfachsten Fällen genauer zu verfolgen. So bei einer Curve \(5^{\mathrm{ter}}\) Ordnung mit zwei Doppelpunkten, welche das allgemeinste Beispiel für \(p=4\) ist. Wählt man auf der Curve irgendwie einen Punct, so kann man zwei Punctepaare angeben, welche mit jenem und den Doppelpunkten Grundpuncte eines Systems von Curven zweiter Ordnung werden; und zwar sind es die Punctepaare, in welchen die Verbindungslinie des gegebenen Puncts mit einem der Doppelpunkte die gegebene Curve noch schneidet; diese Geraden mit beliebigen andern Geraden combinirt, bilden die beiden Systeme von Kegelschnitten.

\subsection*{\S\,62. Die Umkehrung des Abel'schen Satzes.}

Wir kommen jetzt zu denjenigen Betrachtungen, welche umgekehrt Combinationen von \(\Theta\)functionen durch algebraische Functionen ausdrücken lehren, und welche in Folge dessen zugleich algebraische Relationen zwischen \(\Theta\)functionen ergeben.

Das Abel'sche Theorem, wie es am Ende des \S\,58 ausgedrückt ist, lehrt das Doppelverhältniss der vier Functionswerthe
\[
\varphi_\xi,\quad \varphi_\eta,\quad \psi_\xi,\quad \psi_\eta
\]
durch \(\Theta\)functionen darstellen mit Hülfe der Formel:
\[
\tag{1}
\frac{\left(\dfrac{\varphi}{\psi}\right)_{\xi}}
     {\left(\dfrac{\varphi}{\psi}\right)_{\eta}}
=
\prod_{i=0}^{i=q-1}
\frac{\Theta\left(v_h^{(i)}-\int_{\mu}^{\xi}du_h\right)
      \Theta\left(V_h^{(i)}-\int_{\mu}^{\eta}du_h\right)}
     {\Theta\left(v_h^{(i)}-\int_{\mu}^{\eta}du_h\right)
      \Theta\left(V_h^{(i)}-\int_{\mu}^{\xi}du_h\right)},
\]
wo
\[
\tag{2}
\begin{aligned}
v_h^{(i)}&=\int_{\alpha^{(1)}}^{x^{(ip+1)}}du_h+
\int_{\alpha^{(2)}}^{x^{(ip+2)}}du_h\cdots+
\int_{\alpha^{(p)}}^{x^{(i-1)p}}du_h,\\
V_h^{(i)}&=\int_{\alpha^{(1)}}^{y^{(ip+1)}}du_h+
\int_{\alpha^{(2)}}^{y^{(ip+2)}}du_h\cdots+
\int_{\alpha^{(p)}}^{y^{(i-1)p}}du_h.
\end{aligned}
\]
Die \(q p\) Puncte \(x\) einerseits, und die \(q p\) Puncte \(y\) andrerseits haben dabei keine weitern Eigenschaften, als dass, abgesehen von Punctepaaren \(x,y\), welche zusammenfielen und ganz willkürlich sein konnten, die Puncte \(y\) Schnittpuncte von \(\psi=0\) mit \(f=0\), die Puncte \(x\) Schnittpuncte von \(\varphi=0\) mit \(f=0\) sind, während alle übrigen Schnittpuncte der Curven \(\varphi=0\), \(f=0\) zugleich den Curven \(\psi=0\), \(f=0\) angehören; dass ferner die \(y\) allmälig in die \(x\) übergehen, wenn man die Curve \(\psi=0\) der Curve \(\varphi=0\) nähert. Wenn wir nun noch insbesondere annehmen, dass die Curven \(\varphi=0\), \(\psi=0\) durch die Doppel- und Rückkehrpuncte von \(f=0\) gehen sollen, so wird dieser Zusammenhang vermöge des Abel'schen Theorems durch die \(p\) Gleichungen
\[
\tag{3}
 v_h^0+v_h'\cdots+v_h^{(q-1)}=V_h^0+V_h'\cdots+V_h^{(q-1)}
 \qquad(h=1,2\ldots p)
\]
vollständig ausgedrückt, und die Gleichung (1) besteht also immer, wenn die Gleichung (3) nur erfüllt ist. Man kann deswegen den Satz des Abel'schen Theorems umkehren und folgendermassen als Satz über \(\Theta\)producte aussprechen:
\[
\begin{gathered}
\text{Das Product von }q\text{ Doppelverhältnissen}\quad
\prod_{i=0}^{i=q-1}
\frac{\dfrac{\Theta\left(v_h^{(i)}-\int_{\mu}^{\xi}du_h\right)}
              {\Theta\left(V_h^{(i)}-\int_{\mu}^{\xi}du_h\right)}}
     {\dfrac{\Theta\left(v_h^{(i)}-\int_{\mu}^{\eta}du_h\right)}
              {\Theta\left(V_h^{(i)}-\int_{\mu}^{\eta}du_h\right)}}
\end{gathered}
\]
in welchen nur die übrigens willkürlichen Systeme der \(v,V\) verbunden sein sollen durch die Gleichungen (3), ist eine rationale Function der Coordinaten von \(\xi,\eta\) und von den \(q p\) Punctepaaren, welche aus den Gleichungen (2) bestimmt sind; und zwar ist es gleich
\[
\frac{\left(\dfrac{\varphi}{\psi}\right)_{\xi}}
     {\left(\dfrac{\varphi}{\psi}\right)_{\eta}},
\]
wo die Curve \(\varphi=0\) für diejenigen \(x\), die Curve \(\psi=0\) für diejenigen \(y\) verschwindet, welche nicht zusammenfallenden Punctepaaren angehören, und wo beide Curven durch die Doppel- und Rückkehrpuncte von \(f=0\) gehen.

Aber es ist zu bemerken, dass dieser Satz nichts aussagt, wenn \(q=1\), wenn also nur ein einziges Doppelverhältniss von \(\Theta\) Functionen vorhanden ist. Denn in diesem Fall sagen die Gleichungen (3) entweder aus, dass alle \(x\) und \(y\) paarweise zusammenfallen, oder dass die \(x\) sowie die \(y\) Curven \((n-3)^{\mathrm{ter}}\) Ordnung angehören, wobei denn sämmtliche \(\Theta\)functionen verschwinden.

Man kann den vorliegenden Satz auch so auffassen, dass er das Product
\[
\prod_{i=0}^{i=q-1}
\frac{\Theta\left(v_h^{(i)}-\int_{\mu}^{\xi}du_h\right)}
     {\Theta\left(V_h^{(i)}-\int_{\mu}^{\xi}du_h\right)}
\]
auf eine algebraische Function zurückführen lehrt, und auf einen \(\Theta\)quotienten gleicher Art, in welchem nur \(\xi\) durch einen beliebigen andern Punct \(\eta\) ersetzt ist. In diesem Sinne giebt unser Satz den des Hrn. Riemann (Crelles Journal Bd. 54. p. 153), nach welchem jener \(\Theta\)quotient, als Function von \(\xi\) allein betrachtet, eine algebraische Function ist.

\subsection*{\S\,63. Producte von \(\Theta\)functionen auf andre und auf algebraische Functionen zurückgeführt.}

Setzen wir in dem oben ausgesprochenen Satze die Ausdrücke
\[
V_h^{(i)}-\int_{\mu}^{\xi}du_h
\]
mit verschiedenen obern Indices \(i\) einander gleich und gleich \(w_h\), und setzen wir ferner
\[
v_h^{(i)}=V_h^{(i)}+m_h^{(i)},
\]
wo wir unter den Grössen \(m_h^{(i)}\) Constante verstehen, so erhalten wir einen \(\Theta\)quotienten, welchen wir durch
\[
\tag{1}
\frac{\Theta(w+m^{(1)}),\ \Theta(w+m^{(2)})\cdots\Theta(w+m^{(q)})}{\Theta^{q}(w)}
\]
bezeichnen können. Die variabel gedachten Argumente \(w\) sind hier völlig willkürlich, denn die Gleichungen (3) des vorigen Paragraphen gehen hier über in:
\[
\tag{2}
m_h^{(1)}+m_h^{(2)}+\cdots+m_h^{(q)}=0
\qquad(h=1,2\ldots p).
\]
Unter Voraussetzung dieser zwischen den Constanten bestehenden Gleichungen kann man nach dem Vorigen den Ausdruck (1) zurückführen auf einen ähnlichen Ausdruck und auf eine algebraische Function. Setzt man nämlich
\[
\tag{3}
w_h=\int_{\alpha^{(1)}}^{x^{(1)}}du_h+
\int_{\alpha^{(2)}}^{x^{(2)}}du_h\cdots+
\int_{\alpha^{(p)}}^{x^{(p)}}du_h-
\int_{\mu}^{\xi}du_h,
\]
so lässt sich der Quotient (1) auf einen ähnlichen zurückführen, in welchem nur \(\xi\) durch \(\eta\) ersetzt ist, und auf eine rationale Function der Coordinaten von \(x^{(1)},x^{(2)}\ldots x^{(p)},\xi,\eta\). Zunächst zwar würde der Satz des vorigen Paragraphen verlangen, dass man auch die Puncte \(y\) einführt, welche durch die Gleichungen
\[
\tag{4}
w_h+m_h^{(i)}=\int_{\alpha^{(1)}}^{y^{(i,1)}}du_h+
\int_{\alpha^{(2)}}^{y^{(i,2)}}du_h\cdots+
\int_{\alpha^{(p)}}^{y^{(i,p)}}du_h-
\int_{\mu}^{\xi}du_h
\]
definirt sind. Wenn man dann die Curve \(\varphi=0\) durch sämmtliche Puncte \(y\), die Curve \(\psi=0\) durch sämmtliche Puncte \(x\) legt, und zwar so, dass letztere in jedem der Puncte \(x\) \(q\) Schnittpuncte mit \(f=0\) gemein hat (\(q\) punctig berührt), die übrigen Schnittpuncte von \(\varphi=0\), \(f=0\) aber mit den übrigen von \(\psi=0\), \(f=0\) zusammenfallen, und dass endlich beide Curven durch die Doppel- und Rückkehrpuncte gehen, so ist
\[
\frac{\left(\dfrac{\varphi}{\psi}\right)_{\xi}}
     {\left(\dfrac{\varphi}{\psi}\right)_{\eta}}
\]
die gesuchte algebraische Function.

Inzwischen kann man in dem vorliegenden Falle die Puncte \(y\) immer auf constante Puncte und auf die Puncte \(x\) zurückführen. Bezeichnen wir nämlich durch \(c^{(1)},c^{(2)}\ldots c^{(p)}\) beliebige auf \(f=0\) gelegene Puncte, und bestimmen wir dann \(q p\) Puncte \(b^{(i,k)}\) so, dass
\[
\tag{5}
m_h^{(i)}=\int_{c^{(1)}}^{b^{(i,1)}}du_h+
\int_{c^{(2)}}^{b^{(i,2)}}du_h\cdots+
\int_{c^{(p)}}^{b^{(i,p)}}du_h,
\qquad (i=1,2\ldots q,\ h=1,2\ldots p).
\]
Die Gleichungen (2) gehen dann in das Abel'sche Theorem über; legt man eine Curve, welche in jedem der Puncte \(c\) \(q\) punctig berührt, und welche durch die Doppel- und Rückkehrpuncte hindurchgeht, so schneidet diese noch in andern Puncten; legt man durch sie, durch die Doppel- und Rückkehrpuncte und durch die \(q-1\) ersten Systeme der \(b\) eine Curve gleich hoher Ordnung, so schneidet dieselbe noch in dem \(q^{\mathrm{ten}}\) Systeme der \(b\).

Zugleich aber erhält man aus (3), (4), (5):
\[
\int_{x^{(1)}}^{y^{(i,1)}}du_h+\int_{x^{(2)}}^{y^{(i,2)}}du_h\cdots+
\int_{x^{(p)}}^{y^{(i,p)}}du_h
=
\int_{c^{(1)}}^{b^{(i,1)}}du_h+\int_{c^{(2)}}^{b^{(i,2)}}du_h\cdots+
\int_{c^{(p)}}^{b^{(i,p)}}du_h.
\]
Dies ist abermals das Abel'sche Theorem. Legt man eine Curve durch die \(x\), die Doppel- und Rückkehrpuncte und durch ein System der \(b\), so schneidet diese noch in andern Puncten; legt man durch diese, die Doppel- und Rückkehrpuncte und die Puncte \(c\) eine neue Curve derselben Ordnung, so schneidet sie \(f=0\) in dem entsprechenden Systeme der \(y\). Man bestimmt diese also durch die \(x\) selbst und durch die constanten Puncte \(c,b\). Wir können demnach die Coefficienten von \(\varphi\) als rationale Functionen der Coordinaten von \(x^{(1)},x^{(2)}\ldots x^{(p)}\) darstellen, wie weiter unten genauer entwickelt werden soll.

Sprechen wir sonach den Satz aus:

\medskip
\noindent I. \emph{Der Ausdruck}
\[
\frac{\Theta(w+m^{(1)}),\ \Theta(w+m^{(2)})\cdots\Theta(w+m^{(q)})}{\Theta^{q}(w)},
\]
in welchem
\[
m_h^{(1)}+m_h^{(2)}+\cdots+m_h^{(q)}=0\qquad(h=1,2\ldots p),
\]
ist, wenn
\[
w_h=\int_{\alpha^{(1)}}^{x^{(1)}}du_h+
\int_{\alpha^{(2)}}^{x^{(2)}}du_h\cdots+
\int_{\alpha^{(p)}}^{x^{(p)}}du_h-
\int_{\mu}^{\xi}du_h
\]
gesetzt wird, gleich einer rationalen Function der Coordinaten von \(x^{(1)},x^{(2)}\ldots x^{(p)},\xi,\eta\), multiplicirt mit dem \(\Theta\)quotienten, welcher aus dem gegebenen entsteht, wenn darin \(\xi\) durch \(\eta\) ersetzt wird.

Dieser Satz findet seine nothwendige Ergänzung in folgendem Satze:

\medskip
\noindent II. \emph{Der obige Ausdruck ist auch gleich einer rationalen Function der Coordinaten von \(x^{(1)},x^{(2)}\ldots x^{(p)},\xi,a\), multiplicirt mit dem \(\Theta\)quotienten, welcher aus dem gegebenen entsteht, wenn man einen der Puncte \(x\) durch den beliebig auf \(f=0\) gewählten Punct \(a\) ersetzt.}

Man führt leicht diesen zweiten Satz auf den ersten zurück. Dann schreiben wir nur in dem vorgelegten Ausdruck die \(\Theta\)functionen sämmtlich mit negativem Argument, betrachten wir also den Ausdruck
\[
\frac{\Theta(-w-m^{(1)}),\ \Theta(-w-m^{(2)})\cdots\Theta(-w-m^{(q)})}{\Theta^{q}(-w)},
\]
so erfüllen die Grössen \(-m\) noch immer dieselben Bedingungen wie die \(m\) selbst; setzt man also
\[
-w_h=\int_{\alpha^{(1)}}^{z^{(1)}}du_h+
\int_{\alpha^{(2)}}^{z^{(2)}}du_h\cdots+
\int_{\alpha^{(p)}}^{z^{(p)}}du_h-
\int_{\mu}^{\xi}du_h,
\]
so wird nach Satz I. der gegebene Ausdruck gleich einem ähnlichen, in welchem nur \(\xi\) durch einen andern Punct, \(a\), ersetzt ist, multiplicirt mit einer rationalen Function der Coordinaten von \(z^{(1)},z^{(2)}\ldots z^{(p)},\xi,a\). Die \(z,\xi\) hängen aber mit den \(x,\xi\) durch die Gleichung
\[
\sum_{i=1}^{i=p}\left(\int_{\alpha^{(i)}}^{z^{(i)}}du_h+
\int_{\alpha^{(i)}}^{x^{(i)}}du_h\right)-
\int_{\mu}^{\xi}du_h-\int_{\mu}^{\xi}du_h=0
\]
zusammen, aus welcher hervorgeht, dass die \(x,z\) die Schnittpuncte von \(f=0\) mit einer Curve \((n-2)^{\mathrm{ter}}\) Ordnung sind, welche durch die Doppel- und Rückkehrpuncte und durch die \(n-2\) übrigen Puncte der Geraden \(\xi\eta\) geht. Lässt man nun \(\xi\) mit einem der \(x\) zusammenfallen, so geht eines der \(z\) in \(\xi\) über, und die übrigen Puncte \(z\) liegen mit den übrigen \(x\) und mit den Doppel- und Rückkehrpuncten auf einer Curve \((n-3)^{\mathrm{ter}}\) Ordnung. Man sieht, dass diese von dem in \(a\) überzuführenden \(x\) unabhängig werden; daher ist auch der \(\Theta\)quotient, auf welchen der gegebene zurückkommt, von diesem ganz unabhängig, enthält vielmehr an dessen Stelle immer den Punct \(a\); die algebraische Function aber wird eine rationale Function der Coordinaten dieses Puncts \(x\), der Puncte \(a,\xi\) und der mit den übrigen Puncten auf jener Curve \((n-3)^{\mathrm{ter}}\) Ordnung gelegenen Puncte. Da endlich letztere in \(\varphi,\psi\) nur symmetrisch vorkommen, so lassen sich diese Functionen auch durch die Coordinaten der jene Puncte eindeutig bestimmenden übrigen Puncte \(x\) rational darstellen, und die algebraische Function wird also eine rationale Function der Coordinaten von \(x^{(1)},x^{(2)}\ldots x^{(p)},\xi,a\), wie es verlangt wurde.

% --- Continuation: source PDF pp. 244--251, printed pp. 222--229. ---
\clearpage

\subsection*{\S\,64. Ueber die Darstellung der im Vorigen enthaltenen algebraischen Ausdrücke.}

Um die im Vorigen gegebenen Beweise zu ergänzen, muss man zeigen, dass wirklich die Curven \(\varphi,\psi\) rational von den Coordinaten der gegebenen Puncte abhängen. Es kommt dies auf den Beweis des folgenden algebraischen Satzes zurück:

\begin{quote}
Wenn eine Curve \(\Psi=0\) durch gegebene Puncte \(x\) auf \(f=0\) gehen soll, und eine zweite Curve \(\Phi=0\) durch die übrigen Schnittpuncte von \(\Psi=0\) mit \(f=0\), so involvirt dies für die Coefficienten von \(\Phi=0\) eine Reihe von linearen Bedingungen, welche die Coordinaten der Puncte \(x\) nur in rationaler Weise enthalten.
\end{quote}

Man kann übrigens \(\Phi=0\) durch andre für seine Coefficienten lineare Bedingungen vollständig fixiren, und dann durch den Rest des Schnittpunctsystems eine neue Curve \(\Phi_1=0\) legen; auch für diese erhält man dann Bedingungen, welche für die Coefficienten der Curve linear sind, und die \(x\) rational enthalten u.~s.~w. Die ganze hierin enthaltene Reihe von Sätzen lässt sich zusammenfassen, indem man einen Satz von allgemeinerer Fassung nachweist.

Seien nämlich, wie in Abschnitt I, \(x\) und \(s_x\) die Parameter zweier Strahlbüschel. Indem wir die Gleichungen zweier Curven \(f=0\), \(\Psi=0\) in \(x\) und \(s_x\) ausdrücken, und sodann \(s_x\) eliminiren, erhalten wir eine Gleichung
\[
X=0,
\]
welche die Parameter \(x\) derjenigen Strahlen des ersten Büschels liefert, welche nach den Schnittpuncten von \(f=0\), \(\Psi=0\) gerichtet sind. Zugleich aber findet man im Verlauf des Eliminationsprocesses \(s_x\) durch eine rationale Function von \(x\)
\[
s_x=S(x)
\]
ausgedrückt. Diese könnte die Form \(\frac{0}{0}\) nur annehmen, wenn zwei Schnittpuncte auf demselben Strahl lägen; was aber durch Verlegung des Scheitels immer zu vermeiden ist.

Wenn nun ein Theil der Schnittpuncte einzeln oder als Ganzes durch irgend welche Bestimmungen direct gegeben ist, so kann man immer einen Theil des Büschels \(X=0\) direct angeben. Dieser Theil sei \(X^{(1)}=0\). Der Rest
\[
X^{(2)}=0=\frac{X}{X^{(1)}}
\]
entsteht dann durch blosse Division und ist also von den Coefficienten von \(X^{(1)}\) nur in rationaler Weise abhängig.

Soll nun eine Curve \(\Phi=0\) durch die dem Büschel \(X^{(2)}=0\) entsprechenden Schnittpuncte gehen, so verwandle man die Gleichung \(\Phi=0\) zunächst in eine Gleichung zwischen \(x\) und \(s_x\), und setze dann für \(s_x\) seinen Werth \(S(x)\) ein. Die Gleichung \(\Phi=0\) geht dann in eine Gleichung
\[
\Omega(x)=0
\]
über, welche für die in \(X^{(2)}=0\) enthaltenen Strahlen bestehen muss, also \(X^{(2)}=0\) als Factor enthalten muss. Dies ist die Bedingung dafür, dass \(\Phi=0\) durch die zu \(X^{(2)}=0\) gehörigen Schnittpuncte von \(f=0\), \(\Psi=0\) gehen soll. Setzt man
\[
\Omega(x)=X^{(2)}\cdot X^{(3)},
\]
so erhält man durch Gleichsetzung der Coefficienten gleich hoher Potenzen von \(x\) eine Reihe von Gleichungen, welche für die unbekannten Coefficienten von \(\Phi\) und \(X^{(3)}\) linear sind, also auch für erstere allein, wenn man die letztern eliminirt. Diese Bedingungen hängen aber offenbar in rationaler Weise von den Coefficienten von \(X^{(1)}\) ab. Und so hat man den Satz, der den oben angegebenen als speciellen Fall in sich schliesst:

\begin{quote}
Es sei ein Theil der Schnittpuncte von \(f=0\), \(\Psi=0\) gegeben, so dass man die Parameter der nach ihnen gerichteten Strahlen \(x\) eines Büschels als Wurzeln einer Gleichung
\[
X^{(1)}=0
\]
darstellen, und für die nach denselben gerichteten Strahlen eines zweiten Büschels ihre Parameter
\[
s_x=S(x)
\]
rational durch \(x\) ausdrücken kann; dann führt die Bedingung, dass eine Curve \(\Phi=0\) durch die übrigen Schnittpuncte von \(f=0\), \(\Psi=0\) gehen soll, auf eine Reihe von Gleichungen, welche linear sind für die Coefficienten von \(\Phi\) und rational für die Coefficienten von \(X^{(1)}\), \(S(x)\).
\end{quote}

Eine specielle Anwendung dieses Satzes enthält schon \S\,6, wo die Bestimmung der Function \(\Omega\) auf ein ähnliches Problem führte, welches nur dadurch wesentlich vereinfacht wurde, dass \(\Psi=0\) dort eine Gerade war.

\subsection*{\S\,65. Darstellung von Quotienten von \(\Theta\)producten durch algebraische Functionen.}

Wir wollen der Kürze wegen ein Product
\[
\Theta(w+m^{(1)}),\quad \Theta(w+m^{(2)})\cdots \Theta(w+m^{(q)})
\]
ein vollständiges \(\Theta\)product nennen, sobald die Grössen \(m\) den \(p\) Gleichungen
\[
m_h^{(1)}+m_h^{(2)}+\cdots+m_h^{(q)}=0
\]
genügen. Die Potenz \(\Theta^q(w)\) ist selbst ein solches vollständiges \(\Theta\)product. Setzen wir
\[
w_h=\int_{a^{(1)}}^{x^{(1)}}du_h+
\int_{a^{(2)}}^{x^{(2)}}du_h+\cdots+
\int_{a^{(p)}}^{x^{(p)}}du_h-
\int_{\mu}^{\xi}du_h\qquad(h=1,2\ldots p),
\]
und bezeichnen wir durch \(w_h^\circ,w_h^{(1)}\ldots w_h^{(p)}\) die Systeme, welche aus diesem Systeme der Reihe nach hervorgehen, wenn zunächst \(\xi\) durch \(\eta\), sodann auch noch \(x^{(1)}\) durch \(a^{(1)}\), sodann ausserdem \(x^{(2)}\) durch \(a^{(2)}\) etc. ersetzt wird. Aus dem Quotienten
\[
F(w)=\frac{\Theta(w+m^{(1)})\,\Theta(w+m^{(2)})\cdots\Theta(w+m^{(q)})}{\Theta^q(w)}
\]
entstehen auf diese Weise folgende Functionen, in deren Argumenten die beistehenden Puncte als obere Grenzen auftreten:
\[
\begin{array}{r@{\ :\ }l}
F(w)&x^{(1)},x^{(2)}\ldots x^{(p)},\xi,\\
F(w^\circ)&x^{(1)},x^{(2)}\ldots x^{(p)},\eta,\\
F(w^{(1)})&a^{(1)},x^{(2)}\ldots x^{(p)},\eta,\\
\vdots&\vdots\\
F(w^{(p)})&a^{(1)},a^{(2)}\ldots a^{(p)},\eta.
\end{array}
\]
Nach den Theoremen I. II. des \S\,63 ist nun der Quotient je zweier auf einander folgender Functionen \(F\) eine algebraische Function ihrer Argumente, und zwar hat man
\[
\tag{1}
\begin{aligned}
\frac{F(w)}{F(w^\circ)}&=\frac{\varphi_0(\xi)\psi_0(\eta)}{\varphi_0(\eta)\psi_0(\xi)},\\
\frac{F(w^\circ)}{F(w^{(1)})}&=\frac{\varphi_1(x^{(1)})\psi_1(a^{(1)})}{\varphi_1(a^{(1)})\psi_1(x^{(1)})},\\
&\hspace{2em}\cdots\\
\frac{F(w^{(p-1)})}{F(w^{(p)})}&=\frac{\varphi_p(x^{(p)})\psi_p(a^{(p)})}{\varphi_p(a^{(p)})\psi_p(x^{(p)})}.
\end{aligned}
\]
Die Bildungsweise der Curven \(\varphi_i=0\), \(\psi_i=0\) ist in \S\,63 gegeben; bemerken wir hier nur noch, dass die Coefficienten der einzelnen Functionen rational von den Coordinaten der im Folgenden angeführten Puncte abhängen:
\[
\begin{array}{r@{,\ }l@{\quad}l}
\varphi_0&\psi_0&\text{von }x^{(1)},x^{(2)}\ldots x^{(p)},\\
\varphi_1&\psi_1&\text{von }a^{(1)},x^{(2)}\ldots x^{(p)},\\
\varphi_2&\psi_2&\text{von }a^{(1)},a^{(2)}\ldots x^{(p)},\\
\multicolumn{3}{c}{\cdots}\\
\varphi_p&\psi_p&\text{von }a^{(1)},a^{(2)}\ldots a^{(p)}.
\end{array}
\]
Multipliciren wir die Gleichungen (1) mit einander, so erhalten wir die Gleichung
\[
\tag{2}
\frac{F(w)}{F(w^{(p)})}=\frac{
\dfrac{\varphi_0(\xi)\varphi_1(x^{(1)})\cdots\varphi_p(x^{(p)})}{\psi_0(\xi)\psi_1(x^{(1)})\cdots\psi_p(x^{(p)})}}
{\dfrac{\varphi_0(\eta)\varphi_1(a^{(1)})\cdots\varphi_p(a^{(p)})}{\psi_0(\eta)\psi_1(a^{(1)})\cdots\psi_p(a^{(p)})}}.
\]
Sehen wir nun \(\eta,a^{(1)},a^{(2)}\ldots a^{(p)}\) als constante Puncte an, so erhalten wir \(F(w)\) als rationale Function der Coordinaten der Puncte \(x^{(1)},x^{(2)}\ldots\xi\). Dividiren wir endlich zwei Functionen \(F\), die sich durch verschiedene Werthe der \(m\) unterscheiden, so erhalten wir folgenden allgemeinen Satz:

\begin{quote}
Der Quotient zweier vollständigen \(\Theta\)producte von gleich viel Factoren und mit denselben Variabeln \(w\) ist eine rationale Function der Coordinaten der Puncte \(x^{(1)},x^{(2)}\ldots x^{(p)},\xi\), welche durch die Gleichungen
\[
w_h=\int_{\alpha^{(1)}}^{x^{(1)}}du_h+
\int_{\alpha^{(2)}}^{x^{(2)}}du_h+\cdots+
\int_{\alpha^{(p)}}^{x^{(p)}}du_h-
\int_{\mu}^{\xi}du_h
\]
\[(h=1,2\ldots p)
\]
verbunden sind.
\end{quote}

Dieses ist die Ausdehnung auf Abel'sche Functionen für denjenigen Satz, welchen Hr. Hermite als Grundlage für eine grosse Reihe weiterer Sätze in die Theorie der elliptischen Functionen eingeführt hat.

Die algebraische Function ist nothwendig für die Puncte \(x^{(1)},x^{(2)}\ldots x^{(p)}\) symmetrisch, obgleich sie der Form und Bildung nach es nicht ist. Man kann aber z.~B. der linken Seite der Gleichung (2) eine solche symmetrische Form geben, indem man die \(x\) permutirt, und die auf diese Weise entstehenden Ausdrücke addirt und durch ihre Anzahl dividirt.

Man kann hieraus schliessen, dass auch, wenn man
\[
w_h=\int_{\gamma^{(1)}}^{z^{(1)}}du_h+
\int_{\gamma^{(2)}}^{z^{(2)}}du_h+\cdots+
\int_{\gamma^{(s)}}^{z^{(s)}}du_h\qquad(h=1,2\ldots p)
\]
gegeben hat, diese Function noch immer eine rationale Function der Coordinaten von \(z^{(1)},z^{(2)}\ldots z^{(s)}\) ist. Denn bestimmt man die \(x\) aus den Gleichungen
\[
\sum_{i=1}^{i=s}\int_{\gamma^{(i)}}^{z^{(i)}}du_h=\sum_{i=1}^{i=p}\int_{\alpha^{(i)}}^{x^{(i)}}du_h-\int_{\mu}^{\xi}du_h,
\]
welche wieder das Abel'sche Theorem darstellen, so kann man nach \S\,64 die symmetrischen Functionen der Coordinaten der \(x\) als rationale Functionen der Coordinaten der \(z,\xi\) ausdrücken, und also auch die rechte Seite von (2) rational durch die Coordinaten der \(z\) darstellen, wobei dann \(\xi\) ganz herausgehen muss.

\subsection*{\S\,66. Relationen zwischen \(\Theta\)functionen.}

Wir können dem Vorigen nach die Verhältnisse vollständiger \(\Theta\)producte als rationale symmetrische Functionen der Coordinaten von \(p\) Puncten \(x\) darstellen. Auf den Punct \(\xi\) nehmen wir für den Augenblick keine Rücksicht. Er ist willkürlich; man kann ihn etwa mit \(\mu\) zusammenfallen und so aus den Ausdrücken der \(w\) ganz verschwinden lassen.

Bezeichnen wir durch \(P_1(w),P_2(w)\ldots P_{p+2}(w)\), \(p+2\) vollständige \(\Theta\)producte, welche sich nur durch die Werthe der \(m\) von einander unterscheiden. Wir haben dann Formeln der folgenden Art:
\[
\tag{1}
\begin{aligned}
M P_1(w)&=\Phi_1(x^{(1)},x^{(2)}\ldots x^{(p)}),\\
M P_2(w)&=\Phi_2(x^{(1)},x^{(2)}\ldots x^{(p)}),\\
&\hspace{2em}\cdots\\
M P_{p+2}(w)&=\Phi_{p+2}(x^{(1)},x^{(2)}\ldots x^{(p)}),
\end{aligned}
\]
wobei \(M\) ein gemeinschaftlicher Factor ist, dessen Natur hier nicht weiter untersucht werden soll, die \(\Phi\) aber rationale und symmetrische Functionen der Coordinaten der \(p\) Puncte \(x\) bedeuten. Verbindet man mit diesen Gleichungen die Gleichungen
\[
\tag{2}
f(x^{(1)})=0,\quad f(x^{(2)})=0\ldots f(x^{(p)})=0,
\]
oder die symmetrischen Gleichungen, welche sich aus diesen zusammensetzen lassen, so kann man aus (1), (2) die \(x\) und den Proportionalitätsfactor \(M\) eliminiren. Man erhält also den Satz:

\begin{quote}
Zwischen je \(p+2\) vollständigen \(\Theta\)producten von gleich viel Factoren besteht eine homogene algebraische Gleichung.
\end{quote}

Zugleich erhält man im Allgemeinen während des Eliminationsprocesses die symmetrischen Functionen der Coordinaten der \(x\) als homogene rationale Functionen nullter Ordnung der vorgelegten \(p+2\) \(\Theta\)producte. Man kann daher das Umkehrproblem auf eine neue Art auffassen, indem man sagt:

\begin{quote}
Die symmetrischen Functionen der Coordinaten der \(x\), mithin auch die Coefficienten einer Gleichung, deren Wurzeln die Functionswerthe \(\varphi(x^{(1)}),\varphi(x^{(2)})\ldots\varphi(x^{(p)})\) sind, lassen sich als homogene rationale Functionen von \(p+2\) im Allgemeinen beliebig zu wählenden vollständigen \(\Theta\)producten mit gleich viel Factoren darstellen, zwischen welchen selbst eine algebraische Gleichung stattfindet.
\end{quote}

Es ist bemerkenswerth, dass hier eine einzige Irrationalität in den Ausdrücken der \(x\) durch die \(p+1\) Verhältnisse der vollständigen \(\Theta\)producte auftritt, diejenige nämlich, welche durch die erwähnte Relation herbeigeführt ist. Führt man diese \(p+1\) Verhältnisse als neue algebraische Grössen ein, so hat man die symmetrischen Functionen von \(p\) beliebigen Puncten der Curve durch diese \(p+1\) neuen Variabeln algebraisch und rational ausgedrückt, während sie selbst durch eine höhere Gleichung verbunden sind.

Unter die einfachsten Relationen zwischen \(\Theta\)producten gehören diejenigen, welche sich aus den Formeln des Umkehrproblems selbst ergeben, wie sie in \S\,57 und \S\,59 des vorigen Abschnitts dargestellt sind. Bezeichnen wir durch \(\varphi-\lambda\psi=0\) Curven eines Büschels, und durch \(\xi^{(i)}\) die Schnittpuncte von \(f=0\) mit einer Curve des Büschels, deren Parameter \(\lambda^{(i)}\) ist, so haben wir nach \S\,57 die Formeln:
\[
\tag{3}
\begin{aligned}
&M\,[\varphi(x^{(1)})-\lambda^{(i)}\psi(x^{(1)})]\cdot[\varphi(x^{(2)})-\lambda^{(i)}\psi(x^{(2)})]\cdots[\varphi(x^{(p)})-\lambda^{(i)}\psi(x^{(p)})]\\
&\qquad = C^{(i)}\prod_{\xi^{(i)}}\Theta\!\biggl(v_h-\int_{\mu}^{\xi^{(i)}}du_h\biggr),
\end{aligned}
\]
wo \(M\) ein unbestimmter Factor ist, wo ferner
\[
v_h=\sum_{i=1}^{i=p}\int_{a^{(i)}}^{x^{(i)}}du_h,
\]
und
\[
C^{(i)}=\frac{[\varphi(a^{(1)})-\lambda^{(i)}\psi(a^{(1)})]\cdot[\varphi(a^{(2)})-\lambda^{(i)}\psi(a^{(2)})]\cdots[\varphi(a^{(p)})-\lambda^{(i)}\psi(a^{(p)})]}{\displaystyle\prod_{\xi^{(i)}}\Theta\!\biggl(\int_{\mu}^{\xi^{(i)}}du_h\biggr)}.
\]
Die in (3) auftretenden \(\Theta\) bilden ein System vollständiger \(\Theta\)producte, wenn man die \(v\) um geeignete Constante vermehrt, denn die Summe der Argumente mit gleichem \(h\) hat nach dem Abel'schen Theorem in jedem Producte denselben Werth. Setzt man nun in (3) für \(\lambda^{(i)}\), \(p+2\) verschiedene Werthe, so kann man die \(\varphi(x^{(1)})\), \(\psi(x^{(1)})\) eliminiren, und erhält zwischen den \(\Theta\)producten die lineare Relation:
\[
\sum_{i=1}^{i=p+2}
\frac{C^{(i)}\prod_{\xi^{(i)}}\Theta\!\bigl(v_h-\int_{\mu}^{\xi^{(i)}}du_h\bigr)}
{(\lambda^{(i)}-\lambda^{(1)})(\lambda^{(i)}-\lambda^{(2)})\cdots(\lambda^{(i)}-\lambda^{(p+2)})}=0.
\]
In ähnlicher Weise ergeben sich Relationen zwischen \(\Theta\)producten mit je zwei Factoren aus den Formeln des \S\,59. Setzt man unter Beibehaltung der dort gewählten Bezeichnungen
\[
K^{(i)}=\frac{\sum \pm \varphi_1(\xi^{(i)})\varphi_2(a^{(1)})\cdots\varphi_{p+1}(a^{(p)})}
{\Theta\!\bigl(v_h^\circ-\int_{\mu}^{\xi^{(i)}}du_h\bigr)\Theta\!\bigl(v'_h-\int_{\mu}^{\xi^{(i)}}du'_h\bigr)},
\]
so hat man die Gleichung
\[
\begin{aligned}
R_1\varphi_1(\xi^{(i)})+R_2\varphi_2(\xi^{(i)})+\cdots+R_{p+1}\varphi_{p+1}(\xi^{(i)})
= {}& M K^{(i)}\Theta\!\biggl(v_h-\int_{\mu}^{\xi^{(i)}}du_h\biggr)\\
&\cdot\Theta\!\biggl(v_h^\circ+v'_h-v_h-\int_{\mu}^{\xi^{(i)}}du_h\biggr).
\end{aligned}
\]
Daher, wenn man für \(\xi^{(i)}\), \(p+2\) verschiedene Puncte einsetzt, und dann die \(R\) eliminirt:
\[
0=\sum_{i=1}^{i=p+2}K^{(i)}\Delta^{(i)}\Theta\!\biggl(v_h-\int_{\mu}^{\xi^{(i)}}du_h\biggr)\Theta\!\biggl(v_h^\circ+v'_h-v_h-\int_{\mu}^{\xi^{(i)}}du_h\biggr),
\]
wobei \(\Delta^{(1)},\Delta^{(2)}\ldots\Delta^{(p+2)}\) die verschiedenen aus dem System der Grössen \(\varphi_k(\xi^{(i)})\) zu bildenden Determinanten sind.
% --- Continuation: source PDF pp. 252--270, printed pp. 230--248. ---
\clearpage

\section*{Zehnter Abschnitt.}
\subsection*{Die Theilung.}

\subsection*{\S\,67. Multiplication.}

Wenn eine Berührungscurve gelegt werden soll, welche die Curve \(f=0\) in \(p\) gegebenen Puncten \(m\)punctig berührt, welche durch die Doppel- und Rückkehrpuncte von \(f=0\) geht, und für deren Schnittpunctsystem ausserdem alle Puncte bis auf \(p\) gegeben sind, so führt die Bestimmung der letzten auf eine Gleichung \(p^{\text{ten}}\) Grades, wie dies im zweiten Abschnitte auseinandergesetzt ist. Denken wir uns eine feste und eine bewegliche Curve dieser Art; beide sollen sich nur durch die \(p\) Berührungspuncte und durch \(p\) andre Puncte unterscheiden; für die feste Curve sollen ausserdem beide Schaaren von Puncten mit den oben durch \(\alpha\) bezeichneten, einem gegebenen Punct \(\mu\) zugehörigen Puncten zusammenfallen, so dass die feste Curve in der That in den \(p\) Puncten \(\alpha\) \((m+1)\)punctig berührt. Die Berührungspuncte der beweglichen Curve seien \(y^{(1)},y^{(2)}\ldots y^{(p)}\), die übrigen, durch sie bestimmten Schnittpuncte \(x^{(1)},x^{(2)}\ldots x^{(p)}\). Setzen wir
\[
v_h=\sum_{i=1}^{i=p}\int_{\alpha^{(i)}}^{x^{(i)}}du_h,\qquad
w_h=\sum_{i=1}^{i=p}\int_{\alpha^{(i)}}^{y^{(i)}}du_h,\qquad
(h=1,2\ldots p),
\]
so gehen die Gleichungen des Abel'schen Theorems über in
\[
v_h+mw_h=0.\qquad(h=1,2\ldots p)
\]

Nun kann man nach dem Vorigen die Quotienten vollständiger \(\Theta\)producte gleich hoher Ordnung mit den Argumenten \(v\) rational und symmetrisch durch die Coordinaten der Puncte \(x\), sodann wieder nach dem Verfahren des \S\,64 rational und symmetrisch durch die Coordinaten der Puncte \(y\), und endlich den so transformirten Ausdruck nach dem Umkehrproblem durch vollständige \(\Theta\)producte mit den Argumenten \(w\) ausdrücken. Man erhält auf diese Weise den folgenden Satz über die Multiplication der Abel'schen Functionen:

\begin{quote}
Der Quotient zweier vollständigen \(\Theta\)producte gleich hoher Ordnung mit den Argumenten \(mv_1,mv_2\ldots mv_p\) lässt sich rational und homogen durch vollständige \(\Theta\)producte mit den Argumenten \(v_1,v_2\ldots v_p\) darstellen.
\end{quote}

Drücken wir nach \S\,66 insbesondere die Verhältnisse von \(p+1\) vollständigen \(\Theta\)producten mit den Argumenten \(mv\) durch \(p+2\) vollständige \(\Theta\)producte mit den Argumenten \(v\) aus, zwischen welchen letztern eine algebraische Gleichung besteht, so lassen sich auch umgekehrt die Verhältnisse der \(\Theta\)producte mit den Argumenten \(v\) durch jene mit den Argumenten \(mv\) ausdrücken. Diese algebraische Aufgabe führt geometrisch auf die Lösung der Berührungsaufgabe, welche wir im Folgenden behandeln werden, und welche wir mit dem Namen der allgemeinen Theilungsaufgabe belegen, da es sich immer darum handelt, von Functionen mit den Argumenten \(mv\) zu Functionen mit den Argumenten \(v\), oder von Functionen mit den Argumenten \(v\) zu solchen mit den Argumenten \(\frac{v}{m}\) zu gelangen.

\subsection*{\S\,68. Das Problem der allgemeinen Theilung.}

Seien \(c^{(1)},c^{(2)}\ldots c^{(p)}\) Puncte, in welchen man die Curve \(f=0\) durch eine andre Curve \(m\)punctig berühren lässt. Diese Curve schneide \(f=0\) in den Doppel- und Rückkehrpuncten, und in gewissen andern Puncten \(a^{(1)},a^{(2)}\ldots a^{(q)},b^{(1)},b^{(2)}\ldots b^{(r)}\). Durch die Puncte \(b\) und durch gegebene Puncte \(x^{(1)},x^{(2)}\ldots x^{(q)}\), sowie durch die Doppel- und Rückkehrpuncte sollen wir eine Curve von der Ordnung der vorigen legen, welche \(f=0\) in \(p\) Puncten \(m\)punctig berührt; man sucht die Berührungspuncte \(y^{(1)},y^{(2)}\ldots y^{(p)}\). Nach dem Abel'schen Theorem sind die \(y\) dann durch die transcendenten Gleichungen gegeben:
\[
(1)\ldots\quad
m\sum_{i=1}^{i=p}\int_{c^{(i)}}^{y^{(i)}}du_h+
\sum_{i=1}^{i=q}\int_{a^{(i)}}^{x^{(i)}}du_h=0
\qquad(h=1,2\ldots p).
\]

Verstehen wir insbesondere unter den Integralen der zweiten Summe die Werthe, welche dieselben auf bestimmten willkürlich gewählten Integrationswegen annehmen, so kann man auf der rechten Seite von (1) die Null durch das allgemeinste System von Periodicitätsmoduln \(P_h\) der Integrale erster Gattung ersetzen, und findet dann die Bestimmung der \(y\) durch die folgenden Gleichungen, welche die Form des Umkehrproblems haben:
\[
(2)\ldots\quad
\sum_{i=1}^{i=p}\int_{c^{(i)}}^{y^{(i)}}du_h
=
\frac{1}{m}\left[
P_h-\sum_{i=1}^{i=q}\int_{a^{(i)}}^{x^{(i)}}du_h
\right]\qquad(h=1,2\ldots p).
\]
Man sieht hieraus, dass bei bestimmter Wahl der Periodicitätsmoduln \(P\) das System der \(y\), also die Berührungscurve, eindeutig bestimmt ist. Zugleich aber bemerkt man, dass dieses System sich nicht ändert, wenn die \(P_h\) gleichzeitig um das \(m\)fache eines Systems von Periodicitätsmoduln vermehrt oder vermindert werden. Daher braucht man, um alle möglichen Berührungscurven zu erhalten, den in den \(P_h\) auftretenden ganzen Zahlen \(m_1,m_2\ldots m_p,q_1,q_2\ldots q_p\) nur die Werthe \(0,1\ldots m-1\) beizulegen. Man erhält also im Ganzen \(m^{2p}\) Systeme von Periodicitätsmoduln, und also auch ebensoviel Berührungscurven. Das vorgelegte Problem hat daher \(m^{2p}\) Lösungen.\footnote{Ueber die geometrischen Fragen, welche sich hieran knüpfen, vergleiche man Crelles Journal Bd. 63. p. 189 folgg.} Wir wollen jetzt sehen, wie man die Lösungen auf algebraischem Wege entwickelt.

Denken wir uns die Gleichung der gesuchten Berührungscurve, deren Ordnung die \(\nu^{\text{te}}\) sein mag, mit unbestimmten Coefficienten angeschrieben, und bezeichnen wir diese Gleichung durch
\[
\varphi=0.
\]
Diese mit \(f=0\) combinirt, liefert die \(\nu n\) Schnittpuncte beider Curven. Legen wir die Strahlen eines beliebigen Büschels \(\alpha+\lambda\beta=0\), welche nach diesen Puncten gerichtet sind, so finden wir für \(\lambda\) eine Gleichung \(\nu n^{\text{ten}}\) Grades, indem wir aus den gleichzeitig bestehenden Gleichungen
\[
\varphi=0,\qquad f=0,\qquad \alpha+\lambda\beta=0
\]
die \(x\) eliminiren.

Die resultirende Gleichung in \(\lambda\), welche durch \(L=0\) bezeichnet sein mag, enthält eine Reihe bekannter Factoren. Denn die Strahlen \(\lambda\), welche nach den gegebenen Puncten \(a,b\), und die Doppelstrahlen, welche nach den Doppel- und Rückkehrpuncten gerichtet sind, kennt man von vorn herein. Indem man die ihnen entsprechenden Factoren von \(L=0\) entfernt, bleibt eine Gleichung \(mp^{\text{ten}}\) Grades \(L'=0\) zurück, welche die übrigen Schnittpuncte liefert. Diese sollen sich zu \(m\) in \(p\) Puncten vereinigen; der Ausdruck \(L'\) muss also die \(m^{\text{te}}\) Potenz eines Ausdrucks \(\Lambda\) vom \(p^{\text{ten}}\) Grade sein, so dass
\[
(3)\ldots\quad L'=\Lambda^m,
\]
und dass schliesslich die Gleichung \(\Lambda=0\) die \(p\) Berührungspuncte liefert. Denn im Verlauf der oben angedeuteten Elimination findet man die Verhältnisse der Coordinaten der Schnittpuncte als rationale Functionen von \(\lambda\); mithin sind diese selbst gegeben, wenn die betreffenden Werthe von \(\lambda\) bekannt sind.

Denken wir uns die Gleichung \(\varphi=0\) so gebildet, dass man \((m-1)p+1\) specielle Curven \(\varphi_1=0,\varphi_2=0\ldots\nu^{\text{ter}}\) Ordnung durch die Puncte \(a,b\) und durch die Doppel- und Rückkehrpuncte legt, und sodann die Combination
\[
\varphi=\varkappa_1\varphi_1+\varkappa_2\varphi_2+\cdots+
\varkappa_{(m-1)p+1}\varphi_{(m-1)p+1}
\]
bildet, so enthält \(L\) nur noch die \((m-1)p+1\) unbestimmten Constanten \(\varkappa\), oder vielmehr die \((m-1)p\) Verhältnisse derselben; die Gleichung \(L=0\) spaltet sich von selbst in zwei Gleichungen, deren eine von den \(\varkappa\) unabhängig ist und das System der bekannten Strahlen liefert, und in die Gleichung \(L'=0\), welche die \(m^{\text{te}}\) Potenz der Gleichung \(\Lambda=0\) werden soll, und welche dies wird durch eine geeignete Bestimmung der Grössen \(\varkappa\).

Die linke Seite der Gleichung (3) enthält also auf rationale Weise die \(\varkappa\) und die Coordinaten der \(x\). Damit \(L'\) die \(m^{\text{te}}\) Potenz eines Ausdrucks \(p^{\text{ter}}\) Ordnung werde, sind \((m-1)p\) Bedingungen zu erfüllen. In der That, es enthält \(\Lambda\) noch \(p+1\) willkürliche Coefficienten; setzt man nun die entsprechenden Coefficienten von \(L'\) und \(\Lambda^m\) einander gleich, so findet man \(mp+1\) Gleichungen, in denen die \(p+1\) Coefficienten von \(\Lambda\) und die \((m-1)p+1\) Grössen \(\varkappa\) als Unbekannte auftreten. Eliminirt man erstere, so bleiben \((m-1)p\) Gleichungen übrig, welche nur noch die Verhältnisse der \(\varkappa\) enthalten und zur Bestimmung derselben dienen. Eliminiren wir aus den gedachten Gleichungen auch noch \(\varkappa_3,\varkappa_4\ldots\varkappa_{(m-1)p+1}\), so bleibt endlich eine einzige Gleichung übrig, welche nur den Quotienten von \(\varkappa_1\) und \(\varkappa_2\) als Unbekannte enthält.

Diese Gleichung zerfällt nothwendiger Weise in Factoren, je nach dem Character ihrer Wurzeln. Denn wenn auch oben gezeigt wurde, dass für jede der gesuchten Berührungscurven die Gleichung (3) identisch erfüllt sein müsse, so folgt doch noch keinesweges das Umgekehrte; vielmehr können \(m\) Strahlen auch dadurch zusammenfallen, dass die \(m\) Schnittpuncte, welche ihnen entsprechen, sich irgendwie gruppenweise auf die Puncte vertheilen, welche ein solcher Strahl mit der Curve \(f=0\) gemein hat. Aber diejenigen Curven \(\varphi=0\), für welche etwas dieser Art eintritt, hängen offenbar von der Lage des Büschelscheitels \(a=0,b=0\) ab und ändern sich mit diesem, während die gesuchten Berührungscurven von der Lage dieses Scheitels vollkommen unabhängig sind. Daher sieht man, dass die verschiedenen Gleichungen, welche für \(\varkappa_1:\varkappa_2\) erhalten werden, wenn man verschiedene Büschel \(a,b\) benutzt, in einem Theile ihrer Wurzeln übereinstimmen müssen, in einem andern aber nicht; und dass nur der erste Theil für die hier behandelte Frage zur Verwendung kommt. Man hat also nichts zu thun, als die oben bezeichnete Gleichung für zwei verschiedene Büschel \(a,b\) zu bilden, und den beiden gemeinsamen Factor aufzusuchen; dieser Factor giebt die Gleichung für \(\varkappa_1:\varkappa_2\), welche die gesuchten Berührungscurven liefert.

Die so gewonnene neue Gleichung ist nothwendig vom Grade \(m^{2p}\), wie die oben gefundene Zahl der Berührungscurven beweist. Wir wollen diese Gleichung durch \(K=0\) bezeichnen.

Man kann übrigens die Verhältnisse so einrichten, dass die Unbekannte \(\varkappa_1:\varkappa_2\) dieser Gleichung eine einfache geometrische Bedeutung gewinnt. Zu diesem Zwecke braucht man nur die Curven \(\varphi_3=0,\varphi_4=0\ldots\varphi_{(m-1)p+1}=0\) so zu wählen, dass sie in einem bestimmten der gemeinsamen Schnittpuncte \(a,b\) Doppelpunkte besitzen. Da dies zwei weitere Bedingungen involvirt, so giebt es in der That \((m-1)p-1\) solcher Curven, zwischen deren Gleichungen keine lineare Relation stattfindet. In diesem Falle verschwinden für den fraglichen Punct die Differentialquotienten von \(\varphi_3,\varphi_4\) etc., und die Gleichung der Tangente von \(\varphi=0\) in diesem Puncte wird:
\[
\begin{aligned}
(4)\ldots\quad&
\varkappa_1\left(
\frac{\partial\varphi_1}{\partial x_1}X_1+
\frac{\partial\varphi_1}{\partial x_2}X_2+
\frac{\partial\varphi_1}{\partial x_3}X_3
\right)\\
&+\varkappa_2\left(
\frac{\partial\varphi_2}{\partial x_1}X_1+
\frac{\partial\varphi_2}{\partial x_2}X_2+
\frac{\partial\varphi_2}{\partial x_3}X_3
\right)=0.
\end{aligned}
\]
Es sind also die den verschiedenen Berührungscurven entsprechenden Werthe von \(\varkappa_1:\varkappa_2\) die Parameter der in jenem Puncte an diese Curven gelegten Tangenten, welche dem durch die Gleichung (4) repräsentirten Strahlbüschel angehören.

Die Aufsuchung der Puncte \(y\) führt also auf zwei algebraische Gleichungen. Die erste, mit der Unbekannten \(\varkappa_1:\varkappa_2\), ist vom Grade \(m^{2p}\); ihre Coefficienten enthalten die Coordinaten der Puncte \(x\) in rationaler und symmetrischer Weise. Kennt man eine Wurzel dieser Gleichung, so sind \(\varkappa_3,\varkappa_4\) etc. rationale Functionen derselben und der Coordinaten der \(x\), mithin überhaupt die Gleichung der entsprechenden Berührungscurve rational durch diese Wurzel und die Coordinaten der Puncte \(x\) dargestellt. Alsdann ist zweitens die Gleichung \(\Lambda=0\) vom \(p^{\text{ten}}\) Grade zu lösen, welche die \(p\) zugehörigen Werthe von \(\lambda\) und also auch, rational durch je einen solchen Werth ausgedrückt, die Coordinaten der Berührungspuncte liefert; auch die Coefficienten dieser Gleichung enthalten die Coefficienten der Puncte \(x\) in rationaler und symmetrischer Weise.

\subsection*{\S\,69. Das Problem der speciellen Theilung. Zurückführung der allgemeinen Theilung auf die specielle.}

Man kann nun, wie Jacobi für die elliptischen (Crelles Journal Bd. 3, p. 86), und Hermite für die hyperelliptischen Functionen (ib. Bd. 32, p. 293) gezeigt haben, die Auflösung der ersten Gleichung durch Wurzelausziehen bewerkstelligen, sobald man eine specielle Theilungsaufgabe als gelöst betrachtet. Diese specielle Aufgabe entsteht, wenn man die Puncte \(x\) mit den Puncten \(a\) zusammenfallen lässt; wenn man also, von der in den Puncten \(c\) berührenden Curve ausgehend, die übrigen \(m^{2p}-1\) Berührungscurven aufsucht, welche durch die andern Schnittpuncte der erstgenannten Curve mit \(f=0\), insbesondere auch durch die Doppel- und Rückkehrpuncte hindurchgeht. Bezeichnen wir durch \(z^{(1)},z^{(2)}\ldots z^{(p)}\) die Berührungspuncte einer solchen Curve, so hat man der Gleichung (2) \S\,68 entsprechend zur Definition derselben die Gleichungen:
\[
(1)\ldots\quad
\sum_{i=1}^{i=p}\int_{c^{(i)}}^{z^{(i)}}du_h=\frac{P_h}{m}
\qquad(h=1,2\ldots p),
\]
welche, wenn \(P_h=0\), die selbstverständliche Lösung \(z^{(1)}=c^{(1)},z^{(2)}=c^{(2)}\ldots z^{(p)}=c^{(p)}\), für jedes andre System der \(P\) aber eine bestimmte Berührungscurve liefern. Wir wollen nun zunächst sehen, wie man die allgemeine Theilung auf diese specielle zurückführt; sodann aber, wie die Aufgabe der speciellen Theilung zu behandeln ist.

Bezeichnen wir für den Augenblick durch \(y\) die Berührungspuncte einer bestimmten, beliebig gewählten Berührungscurve des allgemeinen Problems, durch \(\eta\) die einer andern, welche dem Periodensystem \(P_h'\) entspricht; endlich durch \(\xi\) die Berührungspuncte einer Curve des speciellen Problems, welche dem Periodensystem \(P_h-P_h'\) zugehört. Man hat dann aus den Gleichungen des Abel'schen Theorems:
\[
\sum_{i=1}^{i=p}\left(
\int_{c^{(i)}}^{y^{(i)}}du_h-
\int_{c^{(i)}}^{\eta^{(i)}}du_h-
\int_{c^{(i)}}^{\xi^{(i)}}du_h
\right)=0,
\]
oder
\[
\sum_{i=1}^{i=p}\left(
\int_{y^{(i)}}^{\eta^{(i)}}du_h+
\int_{c^{(i)}}^{\xi^{(i)}}du_h
\right)=0.
\]
Man findet also die Puncte \(\eta\) auf folgende Art, wenn die \(y\) gegeben sind. Durch die \(y\) und \(c\) legt man eine Curve, welche zugleich die Doppel- und Rückkehrpuncte enthält; durch letztere und die übrigen Schnittpuncte der Curve mit \(f=0\) legt man eine Curve ebenso hoher Ordnung, welche auch durch die \(\xi\) geht; dieselbe schneidet \(f=0\) in den Puncten \(\eta\). Indem man der Reihe nach alle Systeme der \(\xi\) benutzt, erhält man so aus einem System \(y\) alle übrigen, aus einer beliebig gewählten Berührungscurve des allgemeinen Problems die sämmtlichen andern. Die symmetrischen Functionen der Coordinaten der \(\eta\) kann man also rational durch die symmetrischen Functionen der Coordinaten der \(y\) und durch die symmetrischen Functionen der Coordinaten der \(\xi\) darstellen.

Nun entsprach oben jedem System der \(y\), d.~h. jeder Berührungscurve des allgemeinen Problems, eine bestimmte Wurzel
\[
\varkappa=\frac{\varkappa_2}{\varkappa_1}
\]
der Gleichung \(K=0\). Die symmetrischen Functionen der Coordinaten der \(y\) liessen sich rational durch \(\varkappa\), also auch umgekehrt, und zwar auf unendlich viele Arten, \(\varkappa\) als rationale und symmetrische Functionen der Coordinaten der \(y\) darstellen. In ähnlicher Weise stehen die \(\eta\) in Verbindung mit einer zweiten Wurzel \(\varkappa'\) jener Gleichung.

Aber wenn wir in den Gleichungen der vorigen Paragraphen überall die \(a\) mit den \(x\) zusammenfallen lassen, so geht die Aufgabe der allgemeinen Theilung in die specielle über. Die Gleichung \(K=0\) verwandelt sich in eine andere Gleichung \(M=0\), von welcher wir eine Wurzel kennen; denn die Tangentenrichtung der einen Berührungscurve, der ursprünglich gegebenen, ist bekannt. Indem wir also unter \(M=0\) die Gleichung des Grades \(m^{2p}-1\) verstehen, welche nach Hinweglassung des bekannten linearen Factors übrig bleibt, entspricht jeder andern Wurzel \(\mu\) dieser Gleichung ein System der \(\xi\) dergestalt, dass die symmetrischen Functionen der Coordinaten dieser Puncte rationale Functionen von \(\mu\) sind. Wenn wir also vorhin fanden, dass die symmetrischen Functionen der \(y\) sich durch die der \(\eta\) und die der \(\xi\) rational ausdrücken lassen, so können wir dies als Eigenschaft der Gleichung \(K=0\) jetzt folgendermassen aussprechen:

\begin{quote}
Durch irgend eine Wurzel \(\varkappa\) der Gleichung \(K=0\) können wir alle übrigen Wurzeln derselben rational ausdrücken, und zwar mit Hülfe der Formel
\[
\varkappa'=\Theta(\varkappa,\mu),
\]
wo \(\Theta\) eine gewisse rationale Function ist, \(\mu\) aber eine Wurzel der Gleichung \((m^{2p}-1)^{\text{ten}}\) Grades \(M=0\), und wo man alle Wurzeln \(\varkappa'\) der Reihe nach erhält, wenn man für \(\mu\) der Reihe nach alle Wurzeln dieser Gleichung setzt.
\end{quote}

Alle Wurzeln der Gleichung \(K=0\) sind also rationale Functionen von einer derselben. Seien nun \(\varkappa'\), \(\varkappa''\) zwei Wurzeln, welche entstehen, indem man \(\mu\) durch zwei Wurzeln \(\mu'\), \(\mu''\) von \(M=0\) ersetzt, so dass also:
\[
(2)\ldots\quad
\varkappa'=\Theta(\varkappa,\mu'),\qquad \varkappa''=\Theta(\varkappa,\mu'').
\]
Man kann dann, indem man \(\varkappa'\) oder \(\varkappa''\) an Stelle von \(\varkappa\) treten lässt, andre Wurzeln von \(K=0\) bilden. Es ist aber insbesondere leicht zu sehen, dass die beiden Wurzeln
\[
\Theta(\varkappa'',\mu')=\Theta[\Theta(\varkappa,\mu''),\mu'],
\qquad
\Theta(\varkappa',\mu'')=\Theta[\Theta(\varkappa,\mu'),\mu'']
\]
identisch sind, dass also bei abwechselnder Benutzung verschiedener Wurzeln von \(M=0\) die Reihenfolge, in welcher dieselben benutzt werden, keinen Unterschied macht.

In der That nämlich entstand aus der zu dem Periodensystem \(P_h\) gehörigen Wurzel \(\varkappa\) durch Benutzung einer Wurzel \(\mu'\), welche dem Periodensystem \(P_h-P_h'=Q_h'\) angehört, eine Wurzel \(\varkappa'\), welche dem Periodensystem \(P_h'=P_h-Q_h'\) entspricht. Ebenso entsteht aus \(\varkappa\) durch Benutzung der zu \(Q_h''\) gehörigen Wurzel \(\mu''\) die Wurzel \(\varkappa''\), welche dem Periodensystem \(P_h-Q_h''\) entspricht. Setzt man in dem Ausdruck von \(\varkappa'\) an Stelle von \(\varkappa\) die Wurzel \(\varkappa''\), so ist das Periodensystem \(P_h\) durch \(P_h-Q_h''\) zu ersetzen, und man findet die zu dem Periodensystem \(P_h-Q_h'-Q_h''\) gehörige Wurzel
\[
\varkappa'''=\Theta(\varkappa'',\mu')=\Theta[\Theta(\varkappa,\mu''),\mu'].
\]
Setzt man umgekehrt an Stelle von \(\varkappa\) in dem Ausdruck von \(\varkappa''\) die Wurzel \(\varkappa'\), so muss man \(P_h\) durch \(P_h-Q_h'\) ersetzen, und man gelangt zu der demselben Periodensystem \(P_h-Q_h'-Q_h''\) angehörigen Wurzel, welche aber jetzt den Ausdruck hat:
\[
\varkappa'''=\Theta(\varkappa',\mu'')=\Theta[\Theta(\varkappa,\mu'),\mu''].
\]
Hiedurch ist der oben ausgesprochene Satz bewiesen.

Wenn man zweitens immer mit Anwendung der nämlichen Wurzel \(\mu\) die aufeinanderfolgenden Wurzeln von \(K=0\) bildet:
\[
\varkappa'=\Theta(\varkappa,\mu),\quad
\varkappa''=\Theta(\varkappa',\mu),\quad
\varkappa'''=\Theta(\varkappa'',\mu)\ldots,
\]
so findet man, dass dieser Cyclus sich nach \(m\) wiederholten Bildungen, aber nicht früher, schliesst, so dass also
\[
\varkappa^{(m)}=\varkappa.
\]
Denn bezeichnen wir wieder das zu \(\varkappa\) gehörige Periodensystem durch \(P_h\), das zu \(\mu\) gehörige durch \(Q_h\), so gehört zu \(\varkappa'\) das Periodensystem \(P_h-Q_h\), zu \(\varkappa''\) das System \(P_h-2Q_h\), endlich zu \(\varkappa^{(m)}\) das Periodensystem \(P_h-mQ_h\), dessen \(m^{\text{ter}}\) Theil sich von \(\frac{P}{m}\) nur durch ganze Periodicitätsmoduln unterscheidet, welche weggelassen werden können. Es ist daher das zu \(\varkappa^{(m)}\) gehörige Periodensystem mit dem zu \(\varkappa\) gehörigen identisch, daher auch \(\varkappa^{(m)}=\varkappa\), was zu beweisen war.

\subsection*{\S\,70. Lösung der allgemeinen Theilungsgleichung.}

Die Gleichung \(K=0\) gehört also in die Classe derjenigen, welche, wie Abel (Crelles Jornal, Bd. 4, p. 131) nachgewiesen hat, algebraisch auflösbar sind. Um die Auflösung genauer darzustellen, führen wir eine Bezeichnung der Wurzeln \(\varkappa,\mu\) ein, welche dieselben in ihrer Abhängigkeit von den Systemen der Periodicitätsmoduln darstellt. Ist
\[
P_h=2m_h\pi\sqrt{-1}+a_{1h}q_1+a_{2h}q_2\cdots+a_{ph}q_p,
\]
so bezeichnen wir die zu diesem Periodensystem gehörigen Wurzeln \(\varkappa,\mu\) durch
\[
{}^{\varkappa}_{m_1,m_2\ldots,q_1,q_2\ldots},\qquad
{}^{\mu}_{m_1,m_2\ldots,q_1,q_2\ldots}.
\]
Die Gleichung (2) geht dann über in
\[
\Theta\!\left(
\varkappa_{m_1,m_2\ldots,q_1,q_2\ldots},
\mu_{n_1,n_2\ldots,r_1,r_2\ldots}
\right)
=
\varkappa_{m_1-n_1,m_2-n_2\ldots,q_1-r_1,q_2-r_2\ldots}.
\]
Betrachten wir nun die Summe
\[
\begin{aligned}
(3)\ldots\quad S_{a_1a_2\ldots b_1b_2\ldots}
&=\sum_{m_i,q_i}
e^{2(a_1m_1+a_2m_2\ldots+b_1q_1+b_2q_2\ldots)
\frac{\pi\sqrt{-1}}{m}}\\
&\qquad\cdot\varkappa_{m_1,m_2\ldots,q_1,q_2\ldots},
\end{aligned}
\]
in welcher die \(a,b\) ganze Zahlen sind, und die \(m_i,q_i\) unabhängig von einander alle Werthe von \(0\) bis \(m-1\) annehmen sollen. Diese Summe kann man auch in der Form schreiben
\[
\begin{aligned}
(4)\ldots\quad S_{a_1a_2\ldots b_1b_2\ldots}
&=\sum_{n_i,r_i}
e^{2\left[a_1(m_1-n_1)\ldots+b_1(q_1-r_1)\ldots\right]
\frac{\pi\sqrt{-1}}{m}}\\
&\qquad\cdot\varkappa_{m_1-n_1,\ldots,q_1-r_1,\ldots}\\
&=e^{2\left[a_1m_1\ldots+b_1q_1\ldots\right]\frac{\pi\sqrt{-1}}{m}}
\sum_{n_i,r_i}
e^{-2(a_1n_1\ldots+b_1r_1\ldots)\frac{\pi\sqrt{-1}}{m}}\\
&\qquad\cdot
\Theta\!\left(
\varkappa_{m_1\ldots q_1\ldots},
\mu_{n_1\ldots r_1\ldots}
\right),
\end{aligned}
\]
wobei nun die \(n_i,r_i\) unabhängig von einander alle Werthe von \(0\) bis \(m-1\) erhalten müssen. Der letzte Ausdruck ist eine rationale Function von \(\varkappa_{m_1,m_2\ldots q_1q_2\ldots}\), und zwar ist die \(m^{\text{te}}\) Potenz derselben von den Zahlen \(m_i,q_i\) nur insofern abhängig, als sie die vor den übrigen bevorzugte Wurzel bezeichnen. Der Ausdruck \(S_{a_1a_2\ldots b_1b_2\ldots}\) ändert sich aber nicht, wenn man diese Wurzel durch eine andere ersetzt. Setzt man also der Reihe nach für \(\varkappa_{m_1,m_2\ldots q_1q_2\ldots}\) alle Wurzeln von \(K=0\), und addirt die \(m^{\text{ten}}\) Potenzen der aus dem Obigen so entstandenen Ausdrücke, so erhält man \(m^{2p}\) multiplicirt mit der \(m^{\text{ten}}\) Potenz von \(S_{a_1a_2\ldots b_1b_2\ldots}\) gleich einer symmetrischen Function dieser Wurzeln, also gleich einer rationalen Function der Coefficienten von \(K=0\). Man hat also
\[
(5)\ldots\quad
S_{a_1a_2\ldots b_1b_2\ldots}
=
\sqrt[m]{T_{a_1a_2\ldots b_1b_2\ldots}},
\]
wo \(T_{a_1a_2\ldots b_1b_2\ldots}\) eine Function ist, welche zwar die Wurzeln \(\mu\) der Gleichung \(M=0\) enthält, von den Coordinaten der Puncte \(x\) aber nur in rationaler und symmetrischer Weise abhängt.

Wenn man nun die Gleichung (5) für alle Werthsysteme der \(a,b\) bildet, in denen diese Zahlen unabhängig von einander die Werthe \(0,1,\ldots,m-1\) annehmen, so werden die so erhaltenen \(m^{2p}\) Gleichungen ebensoviel lineare Gleichungen zur Bestimmung der \(\varkappa\), und zwar erhält man durch Auflösung
\[
\begin{aligned}
m^{2p}\varkappa_{m_1,m_2\ldots q_1q_2\ldots}
&=\sum_{a_i,b_i}
e^{-2(a_1m_1+a_2m_2\ldots+b_1q_1+b_2q_2\ldots)
\frac{\pi\sqrt{-1}}{m}}\\
&\qquad\cdot
\sqrt[m]{T_{a_1a_2\ldots b_1b_2\ldots}}.
\end{aligned}
\]
Nur das Glied \(\sqrt[m]{T_{00\ldots00\ldots}}\) ist durch eine rationale Function zu ersetzen; denn \(S_{00\ldots00\ldots}\) ist nichts anderes als die Summe aller \(\varkappa\), und daher ein Coefficient der Gleichung \(K=0\) selbst.

Man kann endlich die \(m^{2p}-1\) Wurzelgrössen, welche hier vorkommen, auf nur \(2p\) zurückführen, welche dann schliesslich von einander unabhängig sind und in ihren Combinationen die verschiedenen Wurzeln \(\varkappa\) liefern. Gehen wir nämlich von den \(2p\) Ausdrücken \(S\) aus, in denen immer nur ein Index von Null verschieden ist:
\[
S_{10\ldots00\ldots},\quad
S_{01\ldots00\ldots},\quad\ldots,\quad
S_{00\ldots10\ldots},\quad
S_{00\ldots01\ldots},\quad\ldots,
\]
so ist
\[
\frac{
S_{a_1a_2\ldots b_1b_2\ldots}
}{
S_{10\ldots00\ldots}^{a_1}
S_{01\ldots00\ldots}^{a_2}\cdots
S_{00\ldots10\ldots}^{b_1}
S_{00\ldots01\ldots}^{b_2}\cdots
}
=U_{a_1a_2\ldots b_1b_2\ldots}
\]
eine rationale Function von \(\varkappa_{m_1,m_2\ldots q_1q_2\ldots}\), welche, nach (4) gebildet, sofort zeigt, dass sie von der Wahl der Wurzel \(\varkappa\) unabhängig, und also durch die Coefficienten der Gleichung \(K=0\) rational ausdrückbar ist. Mit Hülfe dieser Gleichung wird nun auch
\[
\begin{aligned}
\sqrt[m]{T_{a_1a_2\ldots b_1b_2\ldots}}
&=
U_{a_1a_2\ldots b_1b_2\ldots}\\
&\quad\cdot
\sqrt[m]{
T_{10\ldots00\ldots}^{a_1}
T_{01\ldots00\ldots}^{a_2}
\cdots
T_{00\ldots10\ldots}^{b_1}
T_{00\ldots01\ldots}^{b_2}
\cdots},
\end{aligned}
\]
so dass also in den endlichen Ausdrücken wirklich nur die \(m^{\text{ten}}\) Wurzeln aus den \(2p\) Ausdrücken
\[
T_{10\ldots00\ldots},\quad
T_{01\ldots00\ldots},\quad\ldots,\quad
T_{00\ldots10\ldots},\quad
T_{00\ldots01\ldots},\quad\ldots
\]
zurückbleiben, deren jede dann auf \(m\) verschiedene Arten gewählt werden kann, und so zu der Entstehung der \(m^{2p}\) verschiedenen Wurzeln \(\varkappa\) Veranlassung giebt.

\subsection*{\S\,71. Zurückführung der speciellen Theilungsgleichung auf den Fall, in welchem \(m\) Potenz einer Primzahl ist.}

Wir kommen jetzt zu der Aufgabe der speciellen Theilung, oder zu der Untersuchung der Gleichung \(M=0\), welche wir bei Auflösung der Gleichung \(K=0\) als gelöst voraussetzten. Auch die Wurzeln dieser Gleichung \(M=0\) stehen unter einander in merkwürdigen Beziehungen. Zunächst sieht man, dass die Periodensysteme sich aus \(2p\) Systemen immer linear und ganzzahlig zusammensetzen, sobald nur die Determinante der in diesen \(2p\) Systemen vorkommenden ganzen Zahlen gleich 1 ist. Wenn wir nun diese \(2p\) Periodensysteme durch die ihnen entsprechenden Integralsummen ersetzen, so geht die Gleichung (1) in eine Form über, in welcher die einer beliebigen speciellen Berührungscurve entsprechenden Integralsummen sich linear und ganzzahlig aus jenen \(2p\) Systemen von Integralsummen zusammensetzen. Nach dem Abel'schen Theorem lassen sich also die symmetrischen Functionen der Coordinaten von den Berührungspuncten irgend einer der Curven durch die symmetrischen Functionen der Coordinaten von den Berührungspuncten solcher \(2p\) auf mannigfache Weise herauszuwählenden Curven ausdrücken; oder es lassen sich alle Wurzeln der Gleichung \(M=0\) durch solche \(2p\) Wurzeln derselben immer rational ausdrücken, deren zugeordnete Periodensysteme die Determinante 1 haben.

Uebrigens verdient es erwähnt zu werden, dass man nur ein einziges specielles Theilungsproblem für jede Zahl \(m\) zu lösen hat. Hat man die Puncte \(z\) aus den Gleichungen
\[
\sum_{i=1}^{i=p}\int_{c^{(i)}}^{z^{(i)}}du_h=\frac{P_h}{m}
\]
bestimmt, so findet man für andre Ausgangspuncte \(\gamma\) die Punctsysteme \(\xi\), welche die Gleichungen
\[
\sum_{i=1}^{i=p}\int_{\gamma^{(i)}}^{\xi^{(i)}}du_h=\frac{P_h}{m}
\]
befriedigen, sofort durch das Abel'sche Theorem, indem man diese Gleichung von der vorigen abzieht.

Aber der Character der Gleichung \(M=0\) ist völlig abhängig von dem Character der Zahl \(m\). Sind \(\nu_1,\nu_2\ldots\) Primzahlen und
\[
m=\nu_1^{\alpha_1}\nu_2^{\alpha_2}\cdots,
\]
so kann man die Behandlung der vorgelegten Gleichung \(M=0\), welche die specielle Theilung für die Zahl \(m\) liefert, immer auf diejenigen Gleichungen zurückführen, welche die Theilung für die einzelnen Primzahlpotenzen \(\nu_1^{\alpha_1},\nu_2^{\alpha_2}\ldots\) ergeben. Denn man kann dem characteristischen Periodensystem, welches zu einer Wurzel der Gleichung \(M=0\) gehört, immer die Form geben:
\[
(1)\ldots\quad
\frac{P'_h}{m}
=
\frac{P_h^{(1)}}{\nu_1^{\alpha_1}}
+
\frac{P_h^{(2)}}{\nu_2^{\alpha_2}}
+\cdots,
\]
wobei nun die in den \(P^{(1)},P^{(2)}\) auftretenden ganzen Zahlen stets kleiner als die entsprechenden Nenner sind; und man erhält alle Periodensysteme \(\frac{P_h}{m}\), wenn man in dieser Formel, unabhängig von einander, den in \(P_h^{(1)}\) enthaltenen ganzen Zahlen alle Werthe von \(0\) bis \(\nu_1^{\alpha_1}-1\), den in \(P_h^{(2)}\) enthaltenen alle Werthe von \(0\) bis \(\nu_2^{\alpha_2}-1\) etc. beilegt. Wenn man aber in (1) die verschiedenen Periodensysteme durch die den Berührungscurven entsprechenden Integralsummen ersetzt, so ergiebt sich aus (1) eine Form des Abel'schen Theorems, welche die Berührungscurve des speciellen Problems für die Zahl \(m\) aus denjenigen für die Zahlen \(\nu_1^{\alpha_1},\nu_2^{\alpha_2}\ldots\) zusammensetzen lehrt, oder welche die Wurzeln der Gleichung \(M=0\) als rationale Functionen derjenigen Wurzeln liefert, welche aus den entsprechenden Gleichungen für \(\nu_1^{\alpha_1},\nu_2^{\alpha_2}\ldots\) erhalten werden.

Man bemerkt insbesondere, dass in der Gleichung (1) der Generalnenner immer kleiner als \(m\) wird, sobald in einem der Periodensysteme \(P_h^{(1)},P_h^{(2)}\ldots\) alle ganzen Zahlen zu Null werden oder eine Potenz des betreffenden Nenners als Factor enthalten. Solche Fälle nun, in denen also auch \(P_h\) einen Factor von \(m\) in allen seinen ganzen Zahlen enthält, geben gar nicht wirkliche Berührungscurven der gesuchten Ordnung. Denn indem man die ursprüngliche Gleichung
\[
m\sum_{i=1}^{i=p}\int_{c^{(i)}}^{z^{(i)}}du_h=P_h
\]
in solchem Falle durch diesen gemeinschaftlichen Factor \(\delta\) dividirt, findet man die Gleichung
\[
\frac{m}{\delta}\sum_{i=1}^{i=p}\int_{c^{(i)}}^{z^{(i)}}du_h=P_h^{(1)},
\]
welche lehrt, dass die den \(z\) entsprechende Curve nicht \(m\)punctig, sondern nur \(\frac{m}{\delta}\)punctig berührt, und also bei dem vorliegenden Problem nur insofern in Frage kommen kann, als sie, \(\delta\)fach gerechnet, eine uneigentliche \(m\)punctig berührende Curve darstellt. Die Anzahl dieser auszuschliessenden Fälle, für welche die entsprechenden Wurzeln \(\mu\) schon als Wurzeln niederer Theilungsgleichungen auftreten, erhält man, wenn man in einem der Periodensysteme \(P_h^{(1)},P_h^{(2)}\ldots\) den ganzen Zahlen nur Werthe beilegt, welche durch den entsprechenden Primfactor theilbar sind, während die ganzen Zahlen der übrigen Periodensysteme beliebig bleiben.

Nun ist in einem Periodensystem \(P_h^{(i)}\) die Anzahl der Zahlensysteme, deren sämmtliche Zahlen nicht durch \(\nu_i\) theilbar sind, leicht zu bestimmen. Man setzt für die erste der darin vorkommenden ganzen Zahlen irgend eine der nicht durch \(\nu_i\) theilbaren ganzen Zahlen, die kleiner als \(\nu_i^{\alpha_i}\) sind, und nimmt die \(2p-1\) andern Zahlen beliebig; dann setzt man für die erste Zahl eine durch \(\nu_i\) theilbare, für die zweite eine nicht theilbare, und nimmt die \(2p-2\) andern beliebig; sodann für die ersten beiden Zahlen solche, die durch \(\nu_i\) theilbar sind, für die dritte eine nicht theilbare, und nimmt die \(2p-3\) übrigen beliebig u.~s.~w. Auf diese Weise erhält man die folgende Anzahl von Combinationen, welche eigentlichen Berührungscurven zugehören:
\[
\begin{aligned}
&\nu_i^{\alpha_i-1}(\nu_i-1)\nu_i^{\alpha_i(2p-1)}
+\nu_i^{2(\alpha_i-1)}(\nu_i-1)\nu_i^{\alpha_i(2p-2)}\\
&\quad+\nu_i^{3(\alpha_i-1)}(\nu_i-1)\nu_i^{\alpha_i(2p-3)}
+\cdots+\nu_i^{2p(\alpha_i-1)}(\nu_i-1)\\
&\qquad=\nu_i^{2p(\alpha_i-1)}(\nu_i^{2p}-1).
\end{aligned}
\]
Und daher ist die gesammte Anzahl der in dem speciellen Theilungsproblem auftretenden eigentlichen Berührungscurven:
\[
\nu_1^{2p(\alpha_1-1)}\nu_2^{2p(\alpha_2-1)}\cdots
(\nu_1^{2p}-1)(\nu_2^{2p}-1)\cdots.
\]

\subsection*{\S\,72. Lösung der speciellen Theilungsgleichung für eine Primzahl und Potenzen einer solchen.}

Es bleibt übrig die Gleichung \(M=0\) unter der Voraussetzung zu lösen, dass \(m\) Potenz einer Primzahl ist, \(m=\nu^\alpha\). Man kann diese Aufgabe immer auf die specielle Theilung für \(m=\nu\) und auf die oben behandelte allgemeine Theilung zurückführen. Denken wir uns nämlich die specielle Theilung für \(m=\nu\) ausgeführt, d.~h. denken wir uns die symmetrischen Functionen der \(z\) so bestimmt, dass
\[
\nu\sum_{i=1}^{i=p}\int_{c^{(i)}}^{z^{(i)}}du_h=P_h,
\]
und bestimmen wir ferner durch die allgemeine Theilung, für \(m=\nu\) andere \(\alpha-1\) Punctsysteme so, dass
\[
\left\{
\begin{aligned}
\nu\sum_{i=1}^{i=p}\int_{c^{(i)}}^{z_1^{(i)}}du_h
&=
\sum_{i=1}^{i=p}\int_{c^{(i)}}^{z^{(i)}}du_h+P_h^{(1)},\\
\nu\sum_{i=1}^{i=p}\int_{c^{(i)}}^{z_2^{(i)}}du_h
&=
\sum_{i=1}^{i=p}\int_{c^{(i)}}^{z_1^{(i)}}du_h+P_h^{(2)},\\
&\ \vdots\\
\nu\sum_{i=1}^{i=p}\int_{c^{(i)}}^{z_{\alpha-1}^{(i)}}du_h
&=
\sum_{i=1}^{i=p}\int_{c^{(i)}}^{z_{\alpha-2}^{(i)}}du_h+P_h^{(\alpha-1)},
\end{aligned}
\right.
\tag{2}
\]
so ergiebt sich durch Combination aller dieser Gleichungen:
\[
\nu^\alpha\sum_{i=1}^{i=p}\int_{c^{(i)}}^{z_{\alpha-1}^{(i)}}du_h
=
P_h+\nu P_h^{(1)}+\nu^2P_h^{(2)}+\cdots+\nu^{\alpha-1}P_h^{(\alpha-1)}.
\]
Das System der \(z_{\alpha-1}\) giebt also die allgemeinste Berührungscurve des speciellen Problems für \(m=\nu^\alpha\). Die Gleichungen (2) als Gleichungen des Abel'schen Theorems lehren zugleich, dass die Wurzeln der speciellen Theilungsgleichung für \(m=\nu^\alpha\) durch wiederholtes Wurzelausziehen aus rationalen Functionen der Wurzeln der Theilungsgleichung für \(m=\nu\) gefunden werden. Diese letzte bleibt also zu behandeln.

Die specielle Theilung für eine Primzahl \(\nu\) führt auf eine Gleichung in \(\mu\) vom Grade \(\nu^{2p}-1\). Die ihren Wurzeln entsprechenden Periodensysteme werden gebildet, wenn man allein in den \(P_h\) vorkommenden ganzen Zahlen unabhängig von einander die Werthe \(0,1,\ldots,\nu-1\) beilegt, und nur den einen Fall ausschliesst, in welchem sämmtliche Zahlen den Werth 0 haben würden; dieser Fall entspricht der gegebenen Berührungscurve, von welcher man ausgeht. Bezeichnen die Grössen \(P_h\) irgend ein Periodensystem, so ist
\[
P_h,\quad 2P_h,\quad 3P_h,\ldots,\quad(\nu-1)P_h
\]
eine Gruppe von Periodensystemen, welche sich nicht ändert, wenn man statt des Systems \(P_h\) irgend ein anderes System der Gruppe zu Grunde legt. Nach der wiederholt gemachten Anwendung des Abel'schen Theorems lässt sich jede dieser Gruppe zugehörige Wurzel durch jede andre rational ausdrücken, und eine symmetrische Function \(A\) der Wurzeln dieser Gruppe ist als dieselbe rationale Function jeder der \(\nu-1\) Wurzeln der Gruppe ausdrückbar.

Die Zahl der Gruppen (welche nie ein Periodensystem oder eine Wurzel gemein haben können) ist
\[
\frac{\nu^{2p}-1}{\nu-1}.
\]
Bildet man die symmetrische Function \(A\) für alle Gruppen, und untersucht die Summe gleich hoher Potenzen dieser \(A\), multiplicirt mit \((\nu-1)\), indem man für \((\nu-1)A\) die Summe der \(\nu-1\) Darstellungen von \(A\) durch die einzelnen Wurzeln der Gruppe setzt, so findet man diese Potenzsummen gleich symmetrischen Functionen aller \(\mu\), also gleich rationalen Functionen der Coefficienten der gegebenen Gleichung. Man kann also sofort eine Gleichung vom Grade \(\frac{\nu^{2p}-1}{\nu-1}\) aufstellen, deren Wurzeln die Werthe sind, welche die symmetrische Function \(A\) für die verschiedenen Gruppen annimmt. Ist sodann \(B\) eine andre symmetrische Function der Wurzeln einer Gruppe, so kann man die Summe aller \(B\), die Summe aller \(BA\), die Summe aller \(BA^2\) etc. ebenso durch die Coefficienten der gegebenen Gleichung rational ausdrücken. Man erhält auf diese Art, wenn die \(A\) bekannt sind, für jede andre symmetrische Function der Wurzeln einer Gruppe eine lineare Bestimmung, und zwar sieht man leicht, dass jedes \(B\) bei der Auflösung dieser Gleichungen ersten Grades so umgeformt werden kann, dass es nur das eine \(A\) enthält, welches ihm entspricht.

Wenn man also die Gleichung des Grades
\[
\frac{\nu^{2p}-1}{\nu-1}
\]
aufgelöst hat, von welcher die \(A\) abhängen, so kann man durch eine Wurzel \(A\) dieser Gleichung die übrigen symmetrischen Functionen der Wurzeln der entsprechenden Gruppe darstellen, und man kann daher sofort eine Gleichung \((\nu-1)^{\text{ten}}\) Grades anschreiben, deren Wurzeln die Wurzeln dieser Gruppe sind, und deren Coefficienten nur von einer Wurzel \(A\) der erstgenannten Gleichung abhängen.

Die Gleichung für die \(A\) ist die einzige in diesen Problemen, welche algebraisch nicht lösbar ist. Denn die den einzelnen Gruppen entsprechenden Gleichungen \((\nu-1)^{\text{ten}}\) Grades kann man wieder nach Abel mit Hülfe von Wurzelgrössen auflösen. Zu diesem Zwecke sei \(g\) eine primitive Wurzel von \(\nu\); dann kann man das Periodensystem einer Gruppe darstellen durch
\[
P_h,\quad gP_h,\quad g^2P_h,\ldots,\quad g^{\nu-2}P_h.
\]
Nennen wir die entsprechenden Wurzeln der Gruppe
\[
\mu_0,\mu_1,\mu_2,\ldots,\mu_{\nu-2},
\]
so ist nach dem Abel'schen Theorem
\[
\mu_1=\Theta(\mu_0),\quad
\mu_2=\Theta(\mu_1),\ldots,\quad
\mu_{\nu-2}=\Theta(\mu_{\nu-3});
\]
aber da \(g^{\nu-1}\equiv 1\pmod{\nu}\), so kommt das der Wurzel
\[
\mu_{\nu-1}=\Theta(\mu_{\nu-2})
\]
entsprechende Periodensystem \(g^{\nu-1}P_h\) auf \(P_h\) selbst zurück und es ist also \(\mu_{\nu-1}=\mu_0\), oder
\[
\mu_0=\Theta(\mu_{\nu-2}).
\]
Die Wurzeln \(\mu_0,\mu_1,\mu_2\ldots\mu_{\nu-2}\) bilden daher, so geordnet, einen Cyclus, so dass, wenn statt \(\mu\) eine andre Wurzel der Gruppe zu Grunde gelegt wird, diese Wurzeln sich nur cyclisch vertauschen. Bildet man nun den Ausdruck
\[
F_\lambda=\sum_{i=1}^{i=\nu-2}
e^{\frac{2\lambda i\pi\sqrt{-1}}{\nu-1}}\mu_i,
\]
so kann man \(F_\lambda\) als rationale Function jeder beliebigen der \(\nu-1\) Wurzeln \(\mu\) darstellen. Setzen wir
\[
\Theta[\Theta(\mu)]=\Theta_1(\mu),\quad
\Theta[\Theta_1(\mu)]=\Theta_2(\mu)\quad\text{etc.},
\]
so hat man allgemein zwischen zwei Wurzeln \(\mu_k,\mu_h\) die Relation
\[
\mu_k=\Theta_{k-h-1}(\mu_h),
\]
wobei, wenn \(k-h-1\) negativ wird, dieser Index um \(\nu-1\) zu vermehren ist. Um also \(F_\lambda\) durch \(\mu_h\) rational auszudrücken, setzt man
\[
F_\lambda=
\sum_{i=1}^{i=\nu-2}
e^{\frac{2\lambda i\pi\sqrt{-1}}{\nu-1}}\,
\Theta_{i-h-1}(\mu_h),
\]
oder, wenn man \(i-h-1=k\) setzt:
\[
F_\lambda=
e^{\frac{2\lambda(h+1)\pi\sqrt{-1}}{\nu-1}}
\sum_{k}
e^{\frac{2\lambda k\pi\sqrt{-1}}{\nu-1}}\Theta_k(\mu_h).
\]
Da hier nun die Zahl \(h\) ausser in \(\mu_h\) nur in einer als Factor vortretenden \((\nu-1)^{\text{ten}}\) Wurzel der Einheit vorkommt, so ist \(F_\lambda^{\nu-1}\) ein Ausdruck, welcher vollständig dieselbe Gestalt hat, durch welche Wurzel \(\mu_h\) man ihn auch darstelle. Bildet man also die \(\nu-1\) Ausdrücke von \(F_\lambda^{\nu-1}\) durch die verschiedenen Wurzeln \(\mu\) und addirt alle, so erhält man \((\nu-1)F_\lambda^{\nu-1}\) gleich einer symmetrischen Function dieser Wurzeln, und also \(F_\lambda^{\nu-1}\) durch die Coefficienten unserer Gleichung \((\nu-1)^{\text{ten}}\) Grades ausgedrückt. Sei demnach
\[
F_\lambda^{\nu-1}=\Phi_\lambda;
\]
so hat man durch Wurzelziehen die \(\nu-1\) Gleichungen:
\[
F_1=\sqrt[\nu-1]{\Phi_1},\qquad
F_2=\sqrt[\nu-1]{\Phi_2},\qquad
\ldots,\qquad
F_{\nu-2}=\sqrt[\nu-1]{\Phi_{\nu-2}},
\]
während \(F_0\) von vornherein eine symmetrische Function der Wurzeln und ein Coefficient der Gleichung \((\nu-1)^{\text{ten}}\) Grades selbst ist. Man hat also \(\nu-1\) Gleichungen ersten Grades, aus denen die \(\nu-1\) Wurzeln der Gruppe sich linear bestimmen.

Zwischen den verschiedenen \((\nu-1)^{\text{ten}}\) Wurzeln, welche in den Ausdrücken der \(F\) vorkommen, bestehen noch Relationen. Denn man sieht, ganz wie oben die entsprechende Eigenschaft von \(F_\lambda^{\nu-1}\) nachgewiesen wurde, dass das Product
\[
\frac{F_\lambda}{F_1^\lambda}
\]
sich rational durch die Coefficienten der Gleichung \((\nu-1)^{\text{ten}}\) Grades ausdrücken lässt. Ist also
\[
\frac{F_\lambda}{F_1^\lambda}=U_\lambda,
\]
wo nun \(U_\lambda\) eine rationale Function der gegebenen Coefficienten bedeutet, so hat man zugleich
\[
F_\lambda=U_\lambda\left(\sqrt[\nu-1]{\Phi_1}\right)^\lambda.
\]
In den Ausdrücken der Wurzeln \(\mu\) kommen also jetzt nur noch die Potenzen einer Wurzelgrösse \(\sqrt[\nu-1]{\Phi_1}\) vor, und man erhält alle Wurzeln \(\mu\) aus einer derselben, indem man dieser Wurzelgrösse ihre verschiedenen Werthe beilegt.

\subsection*{\S\,73. Algebraisch darstellbare \(\Theta\)quotienten.}

Mit den vorliegenden Untersuchungen in genauem Zusammenhange steht die Betrachtung der einfachen algebraisch darstellbaren \(\Theta\)quotienten.



% --- Continuation: source PDF pp. 271--278, printed pp. 249--256. ---
% Source witness: Gordan_Abelsche_Functionen_Google.pdf, rendered at 300dpi.
% Continuation of §73 from previous tranche; source page 278 also begins §74.

Ist \(\frac{\Theta(w+A)}{\Theta(w)}\) ein \(\Theta\)quotient, welcher sich als algebraische Function der Coordinaten der Puncte \(x^{(1)},x^{(2)}\ldots x^{(p)},\xi\) darstellen lässt, welche mit den \(w\) durch die Gleichungen
\[
w_h=\sum_{i=1}^{i=p}\int_{\alpha^{(i)}}^{x^{(i)}}du_h-\int_{\mu}^{\xi}du_h
\qquad(h=1,2\ldots p)
\]
verbunden sind, so muss bei Veränderung der Integrationswege, d.~h. bei Vermehrung der Argumente um Periodicitätsmoduln, jener Quotient nur eine endliche Anzahl verschiedener Werthe anzunehmen fähig sein. Setzt man aber
\[
w_h'=w_h+2m_h\pi\sqrt{-1}+a_{1h}q_1+a_{2h}q_2\cdots+a_{ph}q_p
\]
an Stelle von \(w_h\) (vgl. oben p.~200), so wird
\[
\Theta(w'+A)=e^{(w_1+A_1)q_1+(w_2+A_2)q_2\cdots}\Theta(w+A),
\]
\[
\Theta(w')=e^{w_1q_1+w_2q_2\cdots}\Theta(w),
\]
also
\[
(1)\ldots\quad
\frac{\Theta(w'+A)}{\Theta(w')}
=
e^{A_1q_1+A_2q_2\cdots+A_pq_p}\cdot
\frac{\Theta(w+A)}{\Theta(w)}.
\]
Daher muss, wenn alle \(q\) bis auf eins ungeändert bleiben, der Exponentialfactor seinen ursprünglichen Werth wieder annehmen, sobald diese eine Zahl \(q_i\) um eine gewisse ganze Zahl \(r_i\) gewachsen ist. Man muss sonach haben
\[
e^{A_i(q_i+r_i)}=e^{A_iq_i},\qquad
A_i=\frac{2\pi s_i\sqrt{-1}}{r_i},
\]
wo \(s_i\) wieder eine ganze Zahl ist.

Verlangt man nur, dass \(\frac{\Theta(w+A)}{\Theta(w)}\) multiplicirt mit einem Exponentialfactor
\[
e^{B_1w_1+B_2w_2\cdots+B_pw_p}
\]
bei Vermehrung der Argumente um Periodicitätsmoduln eine endliche Reihe von Werthen besitze, so tritt an Stelle von (1) die Gleichung
\[
\frac{\Theta(w'+A)}{\Theta(w')}\cdot e^{\sum B_iw_i'}=
e^{\sum A_iq_i+2\pi\sqrt{-1}\sum m_iB_i+\sum\sum B_iq_ka_{ik}}
\cdot
\frac{\Theta(w+A)}{\Theta(w)}\cdot e^{\sum B_iw_i},
\]
und der Factor
\[
e^{\sum A_iq_i+2\pi\sqrt{-1}\sum m_iB_i+\sum\sum B_iq_ka_{ik}}
\]
muss also eine endliche Reihe von Werthen haben. Bleibt er daher ungeändert, wenn man \(q_i\) um \(r_i\), oder \(m_i\) um \(\varrho_i\) wachsen lässt, so folgt
\[
e^{2\pi\sqrt{-1}B_i\varrho_i}=1,\qquad
e^{(A_i+\sum_kB_ka_{ik})r_i}=1,
\]
also
\[
B_i=-\frac{\sigma_i}{\varrho_i},\qquad
A_i=\frac{2s_i\pi\sqrt{-1}}{r_i}+\frac{1}{\varrho_i}\sum_k a_{ik}\sigma_k,
\]
wo wieder die \(\sigma,s\) ganze Zahlen sind. Diese Ausdrücke kann man in eine Form bringen, in welcher statt der \(r_i,\varrho_i\) der grösste gemeinschaftliche Dividuus \(m\) dieser Zahlen auftritt, und man darf also setzen:
\[
B_i=-\frac{\sigma_i}{m},\qquad
A_i=\frac1m\left[2s_i\pi\sqrt{-1}+\sum_k a_{ik}\sigma_k\right].
\]
Daher ist das System der \(A\) der \(m^{\text{te}}\) Theil eines allgemeinen Systems von Periodicitätsmoduln, und man hat also den Satz (vgl. Riemann, l.~c. p.~154):

\begin{quote}
Damit
\[
\frac{\Theta(w+A)}{\Theta w}\,e^{\sum w_iB_i}
\]
algebraisch durch die \(\xi,x\) ausdrückbar sei, welche aus den Gleichungen
\[
w_h=\sum_{i=1}^{i=p}\int_{c^{(i)}}^{x^{(i)}}du_h-\int_{\mu}^{\xi}du_h
\]
bestimmt sind, ist es nöthig, dass die \(A\) \(r^{\text{te}}\) Theile eines Systems von Periodicitätsmoduln seien, und wenn
\[
A_i=\frac1m\left[2s_i\pi\sqrt{-1}+\sum_k a_{ik}\sigma_k\right],
\]
so ist
\[
B_i=-\frac{\sigma_i}{m}.
\]
\end{quote}

Der erweiterte Hermitesche Satz lehrt nun diese Darstellung folgendermassen ausführen.

Setzen wir
\[
m_h^{(1)}=m_h^{(2)}\cdots=m_h^{(m-1)}=\frac{P_h}{m},
\]
wo \(P_h\) den Periodicitätsmodul
\[
P=2s_h\pi\sqrt{-1}+\sum_k a_{hk}\sigma_k
\]
bedeutet, so muss, damit der Satz des \S\,65 Anwendung finde,
\[
m_h^{(m)}=-P_h+\frac{P_h}{m}
\]
sein, also
\[
\Theta(w+m^{(m)})=
e^{-\sum_h\left(w_h+\frac1mP_h\right)\sigma_h}
\cdot\Theta\left(w+\frac Pm\right).
\]
Der algebraisch darstellbare Quotient ist also
\[
\left[\frac{\Theta\left(w+\frac Pm\right)}{\Theta(w)}\right]^m
\cdot
e^{-\sum_h\left(w_h+\frac{P_h}{m}\right)\sigma_h}.
\]
Indem wir den constanten \(\Theta\)quotienten auf der rechten Seite des Theorems analog transformiren, erhalten wir für ihn den Ausdruck
\[
\left[\frac{\Theta\left(w^{(p)}+\frac Pm\right)}{\Theta(w^{(p)})}\right]^m
\cdot
e^{-\sum_h\left(w_h^{(p)}+\frac{P_h}{m}\right)\sigma_h},
\]
wo die Grössen
\[
w_h^{(p)}=\sum_{i=1}^{i=p}\int_{\alpha^{(i)}}^{a^{(i)}}du_h-\int_{\mu}^{\eta}du_h
\]
wegen der beliebig auf \(f=0\) zu wählenden Puncte \(a\) beliebige constante Werthe annehmen. Man kann also endlich sagen, dass der Ausdruck
\[
F=
\frac{\Theta\left(w+\frac Pm\right)}{\Theta(w)}
\cdot
\frac{\Theta(w^{(p)})}{\Theta\left(w^{(p)}+\frac Pm\right)}
\cdot
e^{\sum_h\frac{w_h^{(p)}-w_h}{m}\sigma_h}
\]
die \(m^{\text{te}}\) Wurzel einer rationalen Function der Coordinaten von \(x^{(1)},x^{(2)}\ldots x^{(p)},\xi\) (und der \(a\)) sei, wie Riemann l.~c. gezeigt hat. Eine weitere Betrachtung aber führt auf eine bemerkenswerthe Umformung dieser algebraischen Function.

Nach der Darstellung des \S\,65 wird
\[
e^{-\sum_h\frac{w_h\sigma_h}{m}}
\frac{\Theta\left(w+\frac Pm\right)}{\Theta(w)}
=
e^{-\sum_h\frac{w_h^{(p)}\sigma_h}{m}}
\frac{\Theta\left(w^{(p)}+\frac Pm\right)}{\Theta(w^{(p)})}
\sqrt[m]{
\frac{
\dfrac{\varphi_0(\xi)\varphi_1(x^{(1)})\cdots\varphi_p(x^{(p)})}
      {\psi_0(\xi)\psi_1(x^{(1)})\cdots\psi_p(x^{(p)})}
}{
\dfrac{\varphi_0(\eta)\varphi_1(a^{(1)})\cdots\varphi_p(a^{(p)})}
      {\psi_0(\eta)\psi_1(a^{(1)})\cdots\psi_p(a^{(p)})}
}
}.
\]
Eine Formel, welche sich aus den einzelnen Gleichungen zusammensetzt:
\[
(1)\ldots
\begin{cases}
e^{-\sum_h\frac{w_h\sigma_h}{m}}
\dfrac{\Theta\left(w+\frac Pm\right)}{\Theta(w)}
=
e^{-\sum_h\frac{w_h^\circ\sigma_h^\circ}{m}}
\dfrac{\Theta\left(w^\circ+\frac Pm\right)}{\Theta(w^\circ)}
\sqrt[m]{\dfrac{\dfrac{\varphi_0(\xi)}{\psi_0(\xi)}}{\dfrac{\varphi_0(\eta)}{\psi_0(\eta)}}},\\[1.0em]
e^{-\sum_h\frac{w_h^\circ\sigma_h^\circ}{m}}
\dfrac{\Theta\left(w^\circ+\frac Pm\right)}{\Theta(w)}
=
e^{-\sum_h\frac{w_h^{(1)}\sigma_h^{(1)}}{m}}
\dfrac{\Theta\left(w^{(1)}+\frac Pm\right)}{\Theta(w^{(1)})}
\sqrt[m]{\dfrac{\dfrac{\varphi_1(x^{(1)})}{\psi_1(x^{(1)})}}{\dfrac{\varphi_1(a^{(1)})}{\psi_1(a^{(1)})}}},\\
\cdots.
\end{cases}
\]

Die Definition der verschiedenen \(w\) siehe \S\,65, wo auch angegeben ist, von welchen der Puncte \(x,\xi\) die Coefficienten der verschiedenen Functionen \(\varphi,\psi\) abhängen, nämlich:
\[
\begin{array}{ll}
\varphi_0,\psi_0 & \text{von }x^{(1)},x^{(2)}\ldots x^{(p)},\\
\varphi_1,\psi_1 & \text{von }a^{(1)},x^{(2)}\ldots x^{(p)},\\
\varphi_2,\psi_2 & \text{von }a^{(1)},a^{(2)}\ldots x^{(p)},\\
\multicolumn{2}{c}{\cdots}\\
\varphi_p,\psi_p & \text{von }a^{(1)},a^{(2)}\ldots a^{(p)},
\end{array}
\]
während zugleich
\[
(2)\ldots
\begin{cases}
w_h^\circ=\displaystyle\int_{\alpha^{(1)}}^{x^{(1)}}du_h+\int_{\alpha^{(2)}}^{x^{(2)}}du_h\cdots+\int_{\alpha^{(p)}}^{x^{(p)}}du_h-\int_{\mu}^{\eta}du_h,\\[0.55em]
w_h^{(1)}=\displaystyle\int_{\alpha^{(1)}}^{a^{(1)}}du_h+\int_{\alpha^{(2)}}^{x^{(2)}}du_h\cdots+\int_{\alpha^{(p)}}^{x^{(p)}}du_h-\int_{\mu}^{\eta}du_h,\\[0.55em]
w_h^{(2)}=\displaystyle\int_{\alpha^{(1)}}^{a^{(1)}}du_h+\int_{\alpha^{(2)}}^{a^{(2)}}du_h\cdots+\int_{\alpha^{(p)}}^{x^{(p)}}du_h-\int_{\mu}^{\eta}du_h,\\
\cdots,\\
w_h^{(p)}=\displaystyle\int_{\alpha^{(1)}}^{a^{(1)}}du_h+\int_{\alpha^{(2)}}^{a^{(2)}}du_h\cdots+\int_{\alpha^{(p)}}^{a^{(p)}}du_h-\int_{\mu}^{\eta}du_h,
\end{cases}
\]
Ersetzen wir nun in den Gleichungen (1) in der ersten \(x^{(1)},x^{(2)}\ldots x^{(p)}\) durch \(a^{(1)},a^{(2)}\ldots a^{(p)}\), in der zweiten \(x^{(2)},x^{(3)}\ldots x^{(p)}\) durch \(a^{(2)},a^{(3)}\ldots a^{(p)}\), in der dritten \(x^{(3)},x^{(4)}\ldots x^{(p)}\) durch \(a^{(3)},a^{(4)}\ldots a^{(p)}\) u.~s.~w., so gehen die in den Argumenten der \(\Theta\)functionen Gleichung (1) links enthaltenen \(w_h^{(i)}\) bezüglich über in:
\[
\begin{aligned}
W_h&=\int_{\alpha^{(1)}}^{a^{(1)}}du_h+\int_{\alpha^{(2)}}^{a^{(2)}}du_h\cdots+\int_{\alpha^{(p)}}^{a^{(p)}}du_h-\int_{\mu}^{\xi}du_h
=w_h^{(p)}-\int_{\eta}^{\xi}du_h,\\
W_h^\circ&=\int_{\alpha^{(1)}}^{x^{(1)}}du_h+\int_{\alpha^{(2)}}^{a^{(2)}}du_h\cdots+\int_{\alpha^{(p)}}^{a^{(p)}}du_h-\int_{\mu}^{\eta}du_h
=w_h^{(p)}+\int_{a^{(1)}}^{x^{(1)}}du_h,\\
W_h^{(1)}&=\int_{\alpha^{(1)}}^{a^{(1)}}du_h+\int_{\alpha^{(2)}}^{x^{(2)}}du_h\cdots+\int_{\alpha^{(p)}}^{a^{(p)}}du_h-\int_{\mu}^{\eta}du_h
=w_h^{(p)}+\int_{a^{(2)}}^{x^{(2)}}du_h,\\
&\hspace{11em}\cdots;
\end{aligned}
\]
zugleich gehen \(\varphi_0,\psi_0\) in zwei Ausdrücke \(\Phi,\Psi\) über, deren Coefficienten nur noch von constanten Puncten abhängen, und \(\varphi_1,\varphi_2\ldots\varphi_p\), welche ganz gleich gebildet sind, in \(\Phi_1\), \(\psi_1,\psi_2\ldots\psi_p\) in \(\Psi_1\), welche ebenfalls in ihren Coefficienten nur von constanten Puncten abhängen. Die Gleichungen (1) verwandeln sich daher in folgende:
\[
\begin{aligned}
e^{-\sum_h\frac{W_h\sigma_h}{m}}
\frac{\Theta\left(W+\frac Pm\right)}{\Theta(W)}
&=
e^{-\sum_h\frac{w_h^{(p)}\sigma_h}{m}}
\frac{\Theta\left(w^{(p)}+\frac Pm\right)}{\Theta(w^{(p)})}
\sqrt[m]{\dfrac{\dfrac{\Phi(\xi)}{\Psi(\xi)}}{\dfrac{\Phi(\eta)}{\Psi(\eta)}}},\\
e^{-\sum_h\frac{W_h^\circ\sigma_h}{m}}
\frac{\Theta\left(W^\circ+\frac Pm\right)}{\Theta(W^\circ)}
&=
e^{-\sum_h\frac{w_h^{(p)}\sigma_h}{m}}
\frac{\Theta\left(w^{(p)}+\frac Pm\right)}{\Theta(w^{(p)})}
\sqrt[m]{\dfrac{\dfrac{\Phi_1(x^{(1)})}{\Psi_1(x^{(1)})}}{\dfrac{\Phi_1(a^{(1)})}{\Psi_1(a^{(1)})}}},\\
e^{-\sum_h\frac{W_h^{(1)}\sigma_h}{m}}
\frac{\Theta\left(W^{(1)}+\frac Pm\right)}{\Theta(W^{(1)})}
&=
e^{-\sum_h\frac{w_h^{(p)}\sigma_h}{m}}
\frac{\Theta\left(w^{(p)}+\frac Pm\right)}{\Theta(w^{(p)})}
\sqrt[m]{\dfrac{\dfrac{\Phi_1(x^{(2)})}{\Psi_1(x^{(2)})}}{\dfrac{\Phi_1(a^{(2)})}{\Psi_1(a^{(2)})}}},\\
&\hspace{9em}\cdots,
\end{aligned}
\]
bei welchen die algebraischen Functionen nur noch von je einem variabeln Puncte abhängen.

Wenn wir nun jede dieser Gleichungen in die entsprechende Gleichung (1) dividiren, so finden wir
\[
(3)\ldots
\begin{aligned}
\frac{\Theta\left(w+\frac Pm\right)\Theta(W)\Theta(w^\circ)\Theta\left(w^{(p)}+\frac Pm\right)}
{\Theta(w)\Theta\left(W+\frac Pm\right)\Theta\left(w^\circ+\frac Pm\right)\Theta(w^{(p)})}
&=
\sqrt[m]{\dfrac{\dfrac{\varphi_0(\xi)}{\psi_0(\xi)}}{\dfrac{\varphi_0(\eta)}{\psi_0(\eta)}}}
\cdot
\sqrt[m]{\dfrac{\dfrac{\Phi(\eta)}{\Psi(\eta)}}{\dfrac{\Phi(\xi)}{\Psi(\xi)}}},\\[0.7em]
\frac{\Theta\left(w^\circ+\frac Pm\right)\Theta(W^\circ)\Theta(w^{(1)})\Theta\left(w^{(p)}+\frac Pm\right)}
{\Theta(w^\circ)\Theta\left(W^\circ+\frac Pm\right)\Theta\left(w^{(1)}+\frac Pm\right)\Theta(w^{(p)})}
&=
\sqrt[m]{\dfrac{\dfrac{\varphi_1(x^{(1)})}{\psi_1(x^{(1)})}}{\dfrac{\varphi_1(a^{(1)})}{\psi_1(a^{(1)})}}}
\cdot
\sqrt[m]{\dfrac{\dfrac{\Phi_1(a^{(1)})}{\Psi_1(a^{(1)})}}{\dfrac{\Phi_1(x^{(1)})}{\Psi_1(x^{(1)})}}},\\[0.7em]
\frac{\Theta\left(w^{(1)}+\frac Pm\right)\Theta(W^{(1)})\Theta(w^{(2)})\Theta\left(w^{(p)}+\frac Pm\right)}
{\Theta(w^{(1)})\Theta\left(W^{(1)}+\frac Pm\right)\Theta\left(w^{(2)}+\frac Pm\right)\Theta(w^{(p)})}
&=
\sqrt[m]{\dfrac{\dfrac{\varphi_2(x^{(1)})}{\psi_2(x^{(1)})}}{\dfrac{\varphi_2(a^{(1)})}{\psi_2(a^{(1)})}}}
\cdot
\sqrt[m]{\dfrac{\dfrac{\Phi_1(a^{(2)})}{\Psi_1(a^{(2)})}}{\dfrac{\Phi_1(x^{(2)})}{\Psi_1(x^{(2)})}}},\\
&\hspace{9em}\cdots .
\end{aligned}
\]
Aber die Ausdrücke links kann man direct als rationale Functionen der Coordinaten der \(x,a,\xi,\eta\) darstellen, indem man auf das Theorem des \S\,62 zurückgeht. Auf den \(\Theta\)quotienten der ersten Gleichung ist jener Satz in seiner dort gegebenen Gestalt ohne Weiteres anwendbar. Der \(\Theta\)quotient wird gleich
\[
\frac{\left(\frac{U}{\overline V}\right)_\xi}{\left(\frac{U}{\overline V}\right)_\eta},
\]
wenn die Curve \(V=0\) durch die Puncte \(x\) und die aus den Gleichungen
\[
\sum_{i=1}^{i=p}\int_{\alpha^{(i)}}^{z^{(i)}}du_h
=
\sum_{i=1}^{i=p}\int_{\alpha^{(i)}}^{a^{(i)}}du_h+\frac{P_h}{m}
\]
oder
\[
(4)\ldots\quad
\sum_{i=1}^{i=p}\int_{a^{(i)}}^{z^{(i)}}du_h=\frac{P_h}{m}
\]
bestimmten Puncte, die Curve \(U=0\) durch die aus den Gleichungen
\[
\sum_{i=1}^{i=p}\int_{\alpha^{(i)}}^{y^{(i)}}du_h
=
\sum_{i=1}^{i=p}\int_{\alpha^{(i)}}^{x^{(i)}}du_h+\frac{P_h}{m}
\]
oder
\[
(5)\ldots\quad
\sum_{i=1}^{i=p}\int_{x^{(i)}}^{y^{(i)}}du_h=\frac{P_h}{m}
\]
bestimmten Puncte und durch die Puncte \(a\) geht; während alle übrigen Schnittpuncte von \(U=0\) und \(V=0\) mit \(f=0\) gemeinsam sind, die Doppel- und Rückkehrpuncte eingeschlossen.

Man bemerke, dass die Gleichung (5) keine Irrationalität in Bezug auf die \(x\) herbeiführt. Denn obgleich zunächst die \(y\) einem speciellen Theilungsprobleme angehören, bei welchem die \(x\) zu Grunde gelegt sind, so kann man doch nach dem Abel'schen Theorem, indem man statt \(\frac Pm\) etwa die Ausdrücke (4) setzt, diese Irrationalität auf eine solche zurückführen, welche einem speciellen Theilungsproblem mit constanten Grundpuncten entnommen ist.

Die andern \(\Theta\)quotienten muss man zuerst auf eine analoge Form bringen, indem man die Methoden anwendet, welche am Ende des \S\,63 benutzt sind, so dass man die Argumente der \(\Theta\)functionen durch Anwendung des Abel'schen Theorems in solche verwandelt, welche sich (in Zähler und Nenner) nur durch das negativ genommene Integral unterscheiden. Es gehen dadurch die linken Theile der übrigen Gleichungen (3) in rationale Functionen
\[
\frac{\left(\frac{U_1}{V_1}\right)_{x^{(1)}}}{\left(\frac{U_1}{V_1}\right)_{a^{(1)}}},\quad
\frac{\left(\frac{U_2}{V_2}\right)_{x^{(2)}}}{\left(\frac{U_2}{V_2}\right)_{a^{(2)}}},\ldots
\]
über. Setzt man diese in die Gleichungen (3) ein, so findet sich
\[
\begin{aligned}
\sqrt[m]{\dfrac{\dfrac{\varphi_0(\xi)}{\psi_0(\xi)}}{\dfrac{\varphi_0(\eta)}{\psi_0(\eta)}}}
&=
\frac{\left(\frac{U}{V}\right)_\xi}{\left(\frac{U}{V}\right)_\eta}
\cdot
\sqrt[m]{\dfrac{\dfrac{\Phi(\xi)}{\Psi(\xi)}}{\dfrac{\Phi(\eta)}{\Psi(\eta)}}},\\
\sqrt[m]{\dfrac{\dfrac{\varphi_1(x^{(1)})}{\psi_1(x^{(1)})}}{\dfrac{\varphi_1(a^{(1)})}{\psi_1(a^{(1)})}}}
&=
\frac{\left(\frac{U_1}{V_1}\right)_{x^{(1)}}}{\left(\frac{U_1}{V_1}\right)_{a^{(1)}}}
\cdot
\sqrt[m]{\dfrac{\dfrac{\Phi_1(x^{(1)})}{\Psi_1(x^{(1)})}}{\dfrac{\Phi_1(a^{(1)})}{\Psi_1(a^{(1)})}}},\\
\sqrt[m]{\dfrac{\dfrac{\varphi_2(x^{(2)})}{\psi_2(x^{(2)})}}{\dfrac{\varphi_2(a^{(2)})}{\psi_2(a^{(2)})}}}
&=
\frac{\left(\frac{U_2}{V_2}\right)_{x^{(2)}}}{\left(\frac{U_2}{V_2}\right)_{a^{(2)}}}
\cdot
\sqrt[m]{\dfrac{\dfrac{\Phi_1(x^{(2)})}{\Psi_1(x^{(2)})}}{\dfrac{\Phi_1(a^{(2)})}{\Psi_1(a^{(2)})}}},\\
&\hspace{9em}\cdots
\end{aligned}
\]
und man erhält also für unsern \(\Theta\)quotienten folgenden Ausdruck:
\[
\begin{aligned}
(6)\ldots\quad
e^{-\sum_h\frac{w_h\sigma_h}{m}}
\frac{\Theta\left(w+\frac Pm\right)}{\Theta(w)}
&=
e^{-\sum_h\frac{w_h^{(p)}\sigma_h}{m}}
\frac{\Theta\left(w^{(p)}+\frac Pm\right)}{\Theta(w^{(p)})}\\
&\quad{}\cdot
\frac{\left(\frac UV\right)_\xi\left(\frac{U_1}{V_1}\right)_{x^{(1)}}\cdots}
     {\left(\frac UV\right)_\eta\left(\frac{U_1}{V_1}\right)_{a^{(1)}}\cdots}
\sqrt[m]{\frac{
\dfrac{\Phi(\xi)}{\Psi(\xi)}\dfrac{\Phi_1(x^{(1)})}{\Psi_1(x^{(1)})}\cdots
}{
\dfrac{\Phi(\eta)}{\Psi(\eta)}\dfrac{\Phi_1(a^{(1)})}{\Psi_1(a^{(1)})}\cdots
}}.
\end{aligned}
\]
Der gesuchte \(\Theta\)quotient drückt sich also schliesslich als rationale Function der Coordinaten von \(\xi,x^{(1)},x^{(2)}\ldots\) aus, multiplicirt mit der \(m^{\text{ten}}\) Wurzel aus dem Producte rationaler Functionen, deren jede nur einen der variabeln Puncte enthält.

Diese Formel enthält ausser der \(m^{\text{ten}}\) Wurzel keine andern Irrationalitäten, als eine Wurzel \(\mu\) der Gleichung \(M=0\) für das specielle Theilungsproblem.

Dass der Quotient
\[
e^{-\sum_h\frac{(w_h-w_h^{(p)})\sigma_h}{m}}
\frac{\dfrac{\Theta\left(w+\frac Pm\right)}{\Theta(w)}}
     {\dfrac{\Theta\left(w^{(p)}+\frac Pm\right)}{\Theta(w^{(p)})}}
\]
sich durch die \(m^{\text{te}}\) Wurzel einer algebraischen Function ausdrückt, findet darin seine Begründung, dass jener Quotient sich um eine \(m^{\text{te}}\) Wurzel der Einheit ändert, wenn man die Zahlen \(\varrho,\sigma\) um ganze Vielfache von \(m\) ändert, wobei die algebraische Function völlig ungeändert bleibt.

\subsection*{\S\,74. Algebraische Bestimmung von \(\dfrac{\Theta\left(\frac Pm\right)}{\Theta(0)}\).}

Das in der Gleichung (6) enthaltene transcendente Glied
\[
e^{-\sum_h\frac{w_h^{(p)}\sigma_h}{m}}
\frac{\Theta\left(w^{(p)}+\frac Pm\right)}{\Theta(w^{(p)})}
\]
kann selbst auf algebraische Ausdrücke zurückgeführt, und dadurch der gegebene \(\Theta\)quotient rein algebraisch ausgedrückt werden,


% --- Source pages 279--291; printed pages 257--269. ---

indem man die bisher willkürlichen constanten Puncte \(a\) in geeigneter Weise wählt. Es war
\[
w_h^{(p)}=\sum_{i=1}^{i=p}\int_{\alpha^{(i)}}^{a^{(i)}}du_h-\int_{\mu}^{\eta}du_h.
\]
Lassen wir der Einfachheit wegen den Punct \(\eta\) in den Punct \(\mu\) fallen, und setzen sodann
\[
\sum_{i=1}^{i=p}\int_{\alpha^{(i)}}^{a^{(i)}}du_h=\frac{Q_h}{m},
\]
wo die \(Q_h\) ein System von Periodicitätsmoduln bedeuten. Es sind also die \(a\) in diesem Falle Lösung eines speciellen Theilungsproblems. Insbesondere kann man alle \(Q_h\) gleich Null setzen, und die \(a\) mit den \(\alpha\) coincidiren lassen. Dann ist \(w_h^{p}=0\), und der transcendente Factor reducirt sich auf
\[
\frac{\Theta\left(\frac Pm\right)}{\Theta(o)}.
\]
Es handelt sich jetzt darum diese Grössen algebraisch auszudrücken, was folgendermassen geschieht.

Setzen wir in der Formel (6) alle \(w_h\) gleich Null; lassen wir \(\xi,\eta\) mit dem festen Puncte \(\mu\) zusammenfallen, und also die \(x\) mit den \(\alpha\). Die \(w_h^{(p)}\) hingegen bestimmen wir durch die Gleichungen:
\[
w_h^{(p)}=\sum_{i=1}^{i=p}\int_{\alpha^{(i)}}^{a^{(i)}}du_h
=\frac{m-1}{m}P_h.
\]
Die symmetrischen Functionen der Coordinaten der \(a\) werden dann rationale Functionen der zu \(\frac{m-1}{m}P_h\) zugeordneten Wurzel der Gleichung \(M=0\). Aber diese steht mit der zu \(\frac{P_h}{m}\) selbst zugeordneten Wurzel, welche bereits benutzt wurde, in der besondern Beziehung, dass jede dieselbe rationale Function der andern ist, wie der Satz des \(\S\,64\) oder das Abelsche Theorem lehrt. Man kann also auch sagen, die symmetrischen Functionen der Coordinaten der \(a\) seien rationale Functionen der zu \(\frac{P_h}{m}\) zugeordneten Wurzel von \(M=0\).

Durch die angegebenen Werthe gehen die beiden Quotienten
\[
e^{-\sum\frac{w_h\sigma_h}{m}}
\frac{\Theta\left(w+\frac Pm\right)}{\Theta(w)},\qquad
e^{-\sum\frac{w_h^{(p)}\sigma_h}{m}}
\frac{\Theta\left(w^{(p)}+\frac Pm\right)}{\Theta(w^{(p)})}
\]
über in
\[
\frac{\Theta\left(\frac Pm\right)}{\Theta(o)},\qquad
e^{-\frac{m-1}{m^2}\sum P_h\sigma_h}\cdot
\frac{\Theta(P)}{\Theta\left(\frac{m-1}{m}P\right)}
\]
von welchen Ausdrücken der letzte mit Benutzung der für die \(\Theta\)function geltenden Periodicitätsgleichung (p. 200) auch die Form
\[
e^{\sum\frac{P_h\sigma_h}{m^2}}\cdot
\frac{\Theta(o)}{\Theta\left(\frac Pm\right)}
\]
annimmt. Die Gleichung (6) kann man also in der Gestalt schreiben:
\[
\begin{aligned}
e^{-\sum\frac{P_h\sigma_h}{m^2}}
\left[\frac{\Theta\left(\frac Pm\right)}{\Theta(o)}\right]^2
&=
\frac{\left(\frac{U_1}{V_1}\right)_{\alpha^{(1)}}
      \left(\frac{U_2}{V_2}\right)_{\alpha^{(2)}}\cdots}
     {\left(\frac{U_1}{V_1}\right)_{a^{(1)}}
      \left(\frac{U_2}{V_2}\right)_{a^{(2)}}\cdots} \\
&\quad{}\cdot
\sqrt[m]{\frac{
\dfrac{\Phi_1(\alpha^{(1)})\Phi_1(\alpha^{(2)})\cdots}
      {\Psi_1(\alpha^{(1)})\Psi_1(\alpha^{(2)})\cdots}
}{
\dfrac{\Phi_1(a^{(1)})\Phi_1(a^{(2)})\cdots}
      {\Psi_1(a^{(1)})\Psi_1(a^{(2)})\cdots}
}}.
\end{aligned}
\]
Der Ausdruck rechts ist algebraisch. Der unter dem Wurzelzeichen befindliche Ausdruck ist bereits eine symmetrische Function der Coordinaten der verschiedenen Puncte \(\alpha\) einerseits und der verschiedenen Puncte \(a\) andrerseits, also eine rationale Function der zu \(\frac Pm\) zugeordneten Wurzel der speciellen Theilungsgleichung \(M=0\). Der rationale Theil vor dem Wurzelzeichen ist unsymmetrisch gebildet, ändert sich aber durch Vertauschung der \(\alpha\) unter einander und der \(a\) unter einander nicht, und kann daher auch in symmetrische Form gebracht werden. Ist also \(\mu\) die betreffende Wurzel der Theilungsgleichung, so hat man den Satz:

Die Ausdrücke
\[
e^{-\sum\frac{P_h\sigma_h}{2m^2}}
\frac{\Theta\left(\frac Pm\right)}{\Theta(o)}
\]
lassen sich in der Form
\[
\sqrt{\,A(\mu)\sqrt[m]{B(\mu)}\,}
\]
darstellen, wo \(A,B\) rationale Functionen, und \(\mu\) die entsprechende Wurzel der Gleichung \(M=0\) des speciellen Theilungsproblems ist.

Dass dieser Ausdruck sich durch eine \((2m)^{\text{te}}\) Wurzel darstellt, wird dadurch erklärt, dass derselbe sich um eine \((2m)^{\text{te}}\) Wurzel der Einheit ändert, wenn man die Zahlen \(s,\sigma\) um ganze Vielfache von \(m\) ändert, wobei die algebraischen Ausdrücke \(A,B\) völlig ungeändert bleiben.

Den Periodicitätsmoduln \(P_h\) und \(-P_h\) werden zwei verschiedene Wurzeln von \(M=0\) zugeordnet, \(\mu,\mu'\), welche offenbar in obige Formel eingesetzt, beide dasselbe Resultat liefern. Erhebt man obige Gleichung zur Potenz \(2m\), so dass
\[
A^2(\mu)\cdot B(\mu)
=
e^{-\sum\frac{P_h\sigma_h}{m}}
\left[\frac{\Theta\left(\frac Pm\right)}{\Theta(o)}\right]^{2m},
\]
so kann man die Gleichung \(M=0\) mit Hülfe dieser Relation in eine Gleichung \(N=0\) verwandeln, deren Wurzeln die Grössen
\[
e^{-\sum\frac{P_h\sigma_h}{m}}
\left[\frac{\Theta\left(\frac Pm\right)}{\Theta(o)}\right]^{2m}
\]
sind. Hat man diese Gleichung \(N=0\) gelöst, so werden im Allgemeinen aus jeder ihrer Wurzeln zwei zusammengehörige Wurzeln \(\mu,\mu'\) von \(M=0\) gefunden, deren symmetrische Functionen sich während der Bildung der neuen Gleichung durch die neue Unbekannte rational ausdrücken. Nach dem Vorigen hat man nur die Fälle zu unterscheiden und zu behandeln, in denen \(m=2\), und in denen \(m\) eine ungerade Primzahl ist. Im ersten Falle ist die neue Gleichung vom Grade
\[
\frac{m^{2p}-1}{2},
\]
oder jede der algebraisch auflösbaren Gleichungen, welche nach Lösung der Hülfs\-gleichung vom Grade
\[
\frac{m^{2p}-1}{m-1}
\]
übrig bleiben, wird in eine Gleichung vom Grade
\[
\frac{m-1}{2}
\]
verwandelt. Der andere Fall wird im Folgenden genauer erörtert werden.

\subsection*{\S\,75. Bestimmung von \(\dfrac{\Theta\left(\frac P2\right)}{\Theta(o)}\). Curven \((n-3)^{\text{ter}}\) Ordnung, welche in \(p\) Puncten berühren, und durch die Doppel- und Rückkehrpuncte gehen.}

Der Fall von \(m=2\) ist dadurch ausgezeichnet, dass eine Reihe unter den Ausdrücken \(\Theta^{P/2}\) identisch verschwindet. Denn nach der Periodicitätseigenschaft der Function \(\Theta\) (p. 200) ist
\[
\Theta\left(\frac P2\right)
=
\Theta\left(-\frac P2\right)\cdot
e^{-\sum\frac{P_h\sigma_h}{2}+\frac12\sum\sum\sigma_h\sigma_k a_{hk}}.
\]
Da nun
\[
P_h=2s_h\pi\sqrt{-1}+\sum_k a_{hk}\sigma_k,
\]
so geht diese Gleichung über in
\[
\Theta\left(\frac P2\right)=
\Theta\left(-\frac P2\right)\cdot(-1)^{\sum s_h\sigma_h}.
\]
Da nun ausserdem
\[
\Theta\left(\frac P2\right)=\Theta\left(-\frac P2\right),
\]
so verschwindet \(\Theta\left(\frac P2\right)\) immer, sobald \(\sum s_h\sigma_h\) ungerade ist, und man hat den Satz:

Die Function \(\Theta\) verschwindet immer, wenn ihre Argumente eines der Systeme von halben Perioden
\[
\frac12P_h=s_h\pi\sqrt{-1}+\frac12\sum_k a_{hk}\sigma_k
\]
sind, für welche die Summe \(s_1\sigma_1+s_2\sigma_2+\cdots+s_p\sigma_p\) eine ungerade Zahl ist.

Von den Functionen \(\Theta\left(\frac P2\right)\) verschwinden also immer so viele, als die Gleichung
\[
s_1\sigma_1+s_2\sigma_2+\cdots+s_p\sigma_p=2H+1
\]
Lösungen hat, wenn den Zahlen \(s,\sigma\) nur die Werthe \(0\) oder \(1\) beigelegt werden sollen. Um die Anzahl dieser Lösungen zu finden, nehmen wir an, es seien \(k\) der Zahlen \(s\) gleich \(0\), die übrigen \(p-k\) gleich \(1\); die den ersten Zahlen \(s\) entsprechenden \(\sigma\) können dann \(0\) oder \(1\) sein, was \(2^k\) Combinationen giebt. Die Summe der andern \(\sigma\) aber muss ungerade sein, oder es ist immer das letzte durch die übrigen \(p-k-1\) bestimmt. Diese letzten Zahlen \(\sigma\) können daher noch auf \(2^{p-k-1}\) verschiedene Arten gewählt werden, oder es giebt im Ganzen in diesem Falle \(2^{p-1}\) Systeme aller \(\sigma\). Je nachdem man nun unter den Zahlen \(s_1,s_2\ldots s_p\) solche \(k\) herauswählt, welche gleich Null sein sollen, hat man
\[
\frac{p\cdot(p-1)\cdots(p-k+1)}{1\cdot2\cdots k}
\]
verschiedene Fälle, und daher im Ganzen
\[
\frac{p\cdot(p-1)\cdots(p-k+1)}{1\cdot2\cdots k}\cdot2^{p-1}
\]
Systeme der \(s,\sigma\) für ein gegebenes \(k\). Nun kann endlich \(k\) die Werthe \(0,1\ldots p-1\) annehmen; die Anzahl aller möglichen Lösungen ist daher
\[
2^{p-1}\left\{1+\frac p1+\frac{p(p-1)}{1\cdot2}\cdots+p\right\}
=2^{p-1}(2^p-1).
\]
Dieses ist also die Anzahl der Systeme von halben Periodicitätsmoduln, für welche die \(\Theta\)function identisch verschwindet.

Lässt man nun wieder in den allgemeinen Gleichungen
\[
w_h=\sum_{i=1}^{i=p}\int_{\alpha^{(i)}}^{x^{(i)}}du_h-\int_{\mu}^{\xi}du_h
\]
den Punct \(\xi\) mit \(\mu\) zusammenfallen, und ersetzt die \(w\) durch ein System von halben Periodicitätsmoduln, setzt man also
\[
\frac12P_h=\sum_{i=1}^{i=p}\int_{\alpha^{(i)}}^{x^{(i)}}du_h,
\]
so wird die Function \(\Theta\left(\frac P2\right)\) nur dann und immer dann gleich Null, wenn entweder ein Punct \(x^{(i)}\) mit \(\mu\) zusammenfällt, oder wenn die Puncte \(x\) auf einer Curve \((n-3)^{\text{ter}}\) Ordnung liegen, welche durch die Doppel- und Rückkehrpuncte geht. Erinnern wir uns an die Entstehung der \(\alpha\): sie waren die Berührungspuncte einer Curve \((n-2)^{\text{ter}}\) Ordnung, welche durch die Doppel- und Rückkehrpuncte von \(f=0\), sowie durch die \((n-2)\) Puncte geht, in denen die in \(\mu\) an \(f=0\) gezogene Tangente diese Curve noch schneidet; welche endlich in \(p\) Puncten die Curve \(f=0\) berührt. Durch jene festen Puncte gehen noch andre Berührungscurven gleicher Art; eine derselben hat die \(x\) zu Berührungspuncten. Wenn nun einer der Puncte \(x\) in \(\mu\) fällt, so zerfällt deren Berührungscurve in eine Berührungscurve \((n-3)^{\text{ter}}\) Ordnung, welche durch die Doppel- und Rückkehrpuncte geht, und in den \(p-1\) übrigen Puncten \(x\) berührt, und in die Tangente des Punctes \(\mu\) selbst.

Die Berührungspuncte einer solchen Curve \((n-3)^{\text{ter}}\) Ordnung seien \(\beta^{(1)},\beta^{(2)}\ldots\beta^{(p-1)}\); dann giebt es ein gewisses System \(Q_h\) von Periodicitätsmoduln, für welches
\[
(1)\ldots\quad
\frac{Q_h}{2}
=\int_{\alpha^{(1)}}^{\beta^{(1)}}du_h
+\int_{\alpha^{(2)}}^{\beta^{(2)}}du_h
+\cdots+
\int_{\alpha^{(p-1)}}^{\beta^{(p-1)}}du_h
+\int_{\alpha^{(p)}}^{\mu}du_h.
\]
Sei nun \(R_h\) ein solches Periodensystem, für welches die Gleichungen bestehen:
\[
(2)\ldots\quad
\frac{R_h}{2}
=\int_{\alpha^{(1)}}^{\gamma^{(1)}}du_h
+\int_{\alpha^{(2)}}^{\gamma^{(2)}}du_h
+\cdots+
\int_{\alpha^{(p)}}^{\gamma^{(p)}}du_h,
\]
wobei die \(\gamma\) auf einer Curve \((n-3)^{\text{ter}}\) Ordnung liegen, welche durch die Doppel- und Rückkehrpuncte geht. Diese schneidet dann noch etwa in \(\delta^{(1)},\delta^{(2)},\ldots,\delta^{(p-2)}\), und man hat nach dem Abel'schen Theorem:
\[
\begin{aligned}
(3)\ldots\quad 0
&=\int_{\beta^{(1)}}^{\gamma^{(1)}}du_h
+\int_{\beta^{(2)}}^{\gamma^{(2)}}du_h
+\cdots+
\int_{\beta^{(p-1)}}^{\gamma^{(p-1)}}du_h\\
&\quad{}+\int_{\beta^{(1)}}^{\delta^{(1)}}du_h
+\int_{\beta^{(2)}}^{\delta^{(2)}}du_h
+\cdots+
\int_{\beta^{(p-2)}}^{\delta^{(p-2)}}du_h
+\int_{\beta^{(p-1)}}^{\gamma^{(p)}}du_h.
\end{aligned}
\]
Aus einer Combination der Gleichungen (1), (2), (3) findet sich nun:
\[
\frac{Q_h-R_h}{2}
=\int_{\beta^{(1)}}^{\delta^{(1)}}du_h
+\int_{\beta^{(2)}}^{\delta^{(2)}}du_h
+\cdots+
\int_{\beta^{(p-2)}}^{\delta^{(p-2)}}du_h
+\int_{\beta^{(p-1)}}^{\mu}du_h.
\]
Da hier, wenn man mit \(2\) multiplicirt, links ein Periodicitätsmodul steht, und da die \(\beta\) Berührungspuncte einer Berührungscurve \((n-3)^{\text{ter}}\) Ordnung sind, so müssen auch \(\delta^{(1)},\delta^{(2)},\ldots,\delta^{(p-2)},\mu\) Berührungspuncte einer solchen sein. Aber solcher Berührungscurven giebt es eine endliche Anzahl; daher kann man den Punct \(\mu\) immer so wählen, dass er keiner ihrer Berührungspuncte wird, und dass also obige Gleichung, sowie die Voraussetzung, dass die \(\gamma\) auf einer Curve \((n-3)^{\text{ter}}\) Ordnung liegen, unmöglich wird. Wir können also daraus schliessen, dass so oft \(\Theta\left(\frac P2\right)\) auf die angegebene Weise identisch verschwindet, die Berührungscurve \((n-2)^{\text{ter}}\) Ordnung dadurch in die Tangente von \(\mu\) und in eine Berührungscurve \((n-3)^{\text{ter}}\) Ordnung zerfällt, dass einer der Puncte \(x\) mit dem Puncte \(\mu\) zusammenfällt.

Sobald \(s_1\sigma_1+s_2\sigma_2\cdots+s_p\sigma_p=2H+1\), giebt die allgemeine Periodicitätsgleichung der \(\Theta\)function:
\[
\Theta\left(w+\frac P2\right)
=-\Theta\left(-w+\frac P2\right)\cdot e^{\sum\sigma_h w_h},
\]
was man auch in die Form bringen kann:
\[
e^{-\frac12\sum\sigma_h w_h}\Theta\left(w+\frac P2\right)
=
-e^{\frac12\sum\sigma_h w_h}\Theta\left(-w+\frac P2\right).
\]
Die \(2^{p-1}(2^p-1)\) Functionen
\[
e^{-\frac12\sum\sigma_h w_h}\Theta\left(w+\frac P2\right),
\]
für welche \(s_1\sigma_1+s_2\sigma_2\cdots=2H+1\), haben also den Character ungerader Functionen, indem sie nur das Vorzeichen ändern, wenn man alle \(w\) gleichzeitig ihr Zeichen wechseln lässt. Umgekehrt haben die \(2^{p-1}(2^p+1)\) übrigen Functionen
\[
e^{-\frac12\sum\sigma_h w_h}\Theta\left(w+\frac P2\right)
\]
den Charakter von geraden Functionen, da sie sich nicht ändern, wenn man alle \(w\) gleichzeitig ihr Vorzeichen ändern lässt. Aus diesen letzten Functionen aber entstehen ungerade Functionen durch Differentiation nach den verschiedenen Argumenten \(w\); und wenn man also durch \(\varphi\) eine dieser Functionen bezeichnet, so verschwinden die Ausdrücke
\[
\frac{\partial\varphi}{\partial w_1},\quad
\frac{\partial\varphi}{\partial w_2},\ \ldots,\ 
\frac{\partial\varphi}{\partial w_p},
\]
wenn alle \(w\) gleichzeitig gleich Null gesetzt werden. Daher verschwinden auch, wenn man
\[
w_h+\frac{P_h}{2}
=\sum_{i=1}^{i=p}\int_{\alpha^{(i)}}^{x^{(i)}}du_h-\int_{\mu}^{\xi}du_h
\]
setzt, und \(\xi\) mit \(\mu\) zusammenfallen lässt, die Ausdrücke
\[
\frac{\partial\varphi}{\partial x^{(i)}}
=
\frac{\partial\varphi}{\partial w_1}\frac{\partial w_1}{\partial x^{(i)}}
+
\frac{\partial\varphi}{\partial w_2}\frac{\partial w_2}{\partial x^{(i)}}+\cdots.
\]
Nun ist oben bewiesen, dass, wenn eine der hier betrachteten Functionen verschwindet, indem die \(w\) sämmtlich verschwinden, einer der Puncte \(x\) in \(\mu\) falle, die andern Puncte \(x\) aber in die Berührungspuncte einer durch Doppel- und Rückkehrpuncte gehenden Curve \((n-3)^{\text{ter}}\) Ordnung übergehen. Da \(\mu\) beliebig, die Zahl der Berührungscurven aber und ihrer Berührungspuncte eine endliche ist, so kann man immer \(\mu\) von diesen Berührungspuncten verschieden voraussetzen. Wenn daher die betrachtete Function mit ihren Differentialquotienten, sie selbst also mindestens in zweiter Ordnung verschwinden soll, so kann der Grund hiervon nicht darin liegen, dass mehrere der Puncte \(x\) sich in dem Puncte \(\xi\) oder \(\mu\) vereinigen, sondern nur darin, dass, während eines der \(x\), etwa \(x^{(p)}\), in \(\mu\) fällt, die sämmtlichen Puncte \(x\) ausserdem auf einer durch Doppel- und Rückkehrpuncte gehenden Curve \((n-3)^{\text{ter}}\) Ordnung liegen.

Man schliesst hieraus, dass die Functionen unter den oben betrachteten, welche den Character von geraden Functionen haben, nicht verschwinden, sobald alle \(w\) zu Null werden. Denn in solchem Falle müsste dem System von halben Periodicitätsmoduln
\[
\frac{P_1}{2},\quad \frac{P_2}{2},\ldots,\frac{P_p}{2}
\]
eine solche Berührungscurve \((n-3)^{\text{ter}}\) Ordnung entsprechen, deren Berührungspuncte mit jedem Puncte \(\mu\), d. h. mit jedem andern Puncte der Curve auf einer durch Doppel- und Rückkehrpuncte gehenden Curve \((n-3)^{\text{ter}}\) Ordnung lägen. Diese Berührungspuncte müssten also zwei ganz verschiedene und im Allgemeinen durchaus nicht zusammentreffende Eigenschaften, welche durch ganz verschiedene algebraische Gleichungen definirt sind, gleichzeitig besitzen: die Berührungspuncte einer Berührungscurve \((n-3)^{\text{ter}}\) Ordnung zu sein, und ein System von \(p-1\) Puncten zu bilden, welche eine durch die Doppel- und Rückkehrpuncte gehende Curve \((n-3)^{\text{ter}}\) Ordnung nicht bestimmen.

Es verschwinden also im Allgemeinen nur diejenigen Functionen \(\Theta\left(\frac P2\right)\), für welche \(\sum s_i\sigma_i\) ungerade ist, und man hat zugleich den Satz:

Es giebt \(2^{p-1}(2^p-1)\) Curven \((n-3)^{\text{ter}}\) Ordnung, welche durch die Doppel- und Rückkehrpuncte gehen, und die Curve \(f=0\) in \(p-1\) Puncten berühren.

Die Curven werden aus den Gleichungen
\[
\sum_{i=1}^{i=p-1}\int_{\alpha^{(i)}}^{x^{(i)}}du_h
+\int_{\alpha^{(p)}}^{\mu}du_h
=\frac12P_h
\qquad (h=1,2\ldots p)
\]
gefunden, welche zusammen bestehen können und in den \(x\) die Berührungspuncte der Curven liefern, sobald für die in den \(P\) auftretenden ganzen Zahlen die Gleichung besteht
\[
s_1\sigma_1+s_2\sigma_2\cdots+s_p\sigma_p=2H+1.
\]
Die Lösung der besondern Classe von Umkehrproblemen, denen das vorliegende angehört, ist in \(\S\,60\) gegeben.

Man bemerkt, dass der Fall, wo alle \(P\) verschwinden, in den hier betrachteten Fällen nicht enthalten ist. Daher ist dieser Fall einer unter denjenigen, welche in folgendem Satz zusammengefasst erscheinen, der den Vorigen ergänzt:

Es giebt \(2^{p-1}(2^p+1)\) Curven \((n-2)^{\text{ter}}\) Ordnung, welche durch die Doppel- und Rückkehrpuncte, so wie durch die \(n-2\) Schnittpuncte der in einem beliebigen Puncte \(\mu\) gelegten Tangente mit \(f=0\) hindurchgehen und die Curve \(f=0\) in \(p\) Puncten berühren.

Die Berührungspuncte dieser Curven findet man aus den Gleichungen:
\[
\sum_{i=1}^{i=p}\int_{\alpha^{(i)}}^{x^{(i)}}du_h=\frac12P_h,
\qquad (h=1,2\ldots p)
\]
wenn man für die \(P_h\) alle diejenigen Systeme von Periodicitätsmoduln setzt, für welche
\[
s_1\sigma_1+s_2\sigma_2\cdots+s_p\sigma_p=2H.
\]

\subsection*{\S\,76. Reduction der Theilungsprobleme für \(m=2\).}

Wenn man unmittelbar auf algebraischem Wege eine Curve \((n-3)^{\text{ter}}\) Ordnung \(\varphi=0\) zu bestimmen sucht, welche die Curve \(f=0\) in \(p-1\) Puncten berührt und durch die Doppel- und Rückkehrpuncte geht, so wird man einen Coefficienten von \(\varphi\) gleich \(1\) setzen, und aus den aufzustellenden Bedingungen alle übrigen durch einen derselben rational ausdrücken können, während dieser selbst \((\varrho)\) von einer höhern Gleichung \(R=0\) abhängt. Der Grad dieser Gleichung ist nach dem Vorigen \(2^{p-1}(2^p-1)\). Wenn man eine Wurzel derselben kennt, so erfordert die Bestimmung der zugehörigen Berührungspuncte die Lösung einer Gleichung \((p-1)^{\text{ten}}\) Grades.

Die Frage nach den Berührungscurven \((n-2)^{\text{ter}}\) Ordnung, welche durch die Doppel- und Rückkehrpuncte, so wie durch \(n-2\) feste Puncte gehen, und \(f=0\) in \(p\) Puncten berühren, ist nach \(\S\,68\) ein Problem der allgemeinen Theilung, und führt auf eine Gleichung \(K=0\) vom Grade \(2^{2p}\). Die Auflösung derselben geschieht nach \(\S\,70\) mit Zurückführung auf das specielle Theilungsproblem, welches darin besteht, die übrigen Lösungen des obigen Berührungsproblems zu suchen, wenn eine derselben gegeben ist, und welches von einer Gleichung des Grades \(2^{2p}-1\) abhängt.

Aber in dem vorliegenden Falle treten andre Verhältnisse ein, als bei der Theilung nach einer ungeraden Primzahl, indem nicht die specielle Theilungsgleichung die einfachste ist, auf welche die Lösung des allgemeinen Theilungsproblems zurückkommt, sondern indem beide Probleme durch einen besondern Fall des allgemeinen Theilungsproblems einfacher erledigt werden. Dieses besondere Problem besteht darin, dass die \(n-2\) festen Puncte des allgemeinen Problems als die übrigen Schnittpuncte der in irgend einem Puncte \(\mu\) an \(f=0\) gelegten Tangente angenommen werden. Wie oben gezeigt, sind in diesem Falle nur \(2^{p-1}(2^p+1)\) wirkliche Berührungscurven vorhanden, während die übrigen in die Tangente selbst in Verbindung mit den verschiedenen Berührungscurven \((n-3)^{\text{ter}}\) Ordnung übergehen, welche durch die oben betrachtete Gleichung \(R=0\) gegeben werden. Bilden wir also die Gleichung \(K=0\) des allgemeinen Theilungsproblems für diesen besondern Fall, und bemerken wir, dass die Unbekannte in derselben der Parameter eines Strahlbüschels war, welcher in einem der \(n-2\) festen Puncte gelegt wurde, und zwar der Parameter desjenigen Strahls, welcher Tangente einer der gesuchten Berührungscurven wurde. Von diesen fallen hier \(2^{p-1}(2^p-1)\) mit der Tangente des Punctes \(\mu\) zusammen, welche einen Bestandtheil von ebensoviel zerfallenden Berührungscurven bildete. Die Gleichung \(K=0\) hat also für diesen besondern Fall eine bekannte \(2^{p-1}(2^p-1)\) fache Wurzel, und nach Absonderung derselben bleibt eine Gleichung \(N=0\) vom Grade \(2^{p-1}(2^p+1)\) zurück, welche aufzulösen ist, und welche die eigentlichen Berührungscurven \((n-2)^{\text{ter}}\) Ordnung liefert. Die Wurzeln dieser Gleichung vom Grade \(2^{p-1}(2^p+1)\) sind durch die Wurzeln der Gleichung \(R=0\) vom Grade \(2^{p-1}(2^p-1)\) rational ausdrückbar.

Dies beruht auf Folgendem. Bezeichnen wir der Kürze wegen Systeme von Periodicitätsmoduln von der Form
\[
2s_h\pi\sqrt{-1}+a_{h1}\sigma_1+a_{h2}\sigma_2+\cdots+a_{hp}\sigma_p,
\qquad (h=1,2\ldots p),
\]
in welchen die \(s,\sigma\) die Werthe \(0\) oder \(1\) haben sollen, durch \(G_h\) oder \(U_h\), je nachdem \(\sum s_i\sigma_i\) gerade oder ungerade ist, so bemerkt man leicht, dass man auf mannigfache Weise \(2p\) solche Systeme \(U_h\) aufstellen kann, bei denen die aus den \(4p^2\) Zahlen \(s,\sigma\) gebildete Determinante gleich \(\pm1\) ist; z. B. indem man als die verschiedenen Zahlensysteme der \(s,\sigma\) die Horizontalreihen des Schemas benutzt:
\[
\scriptsize
\begin{array}{c|ccccc|ccccc}
& s_1&s_2&s_3\cdots&s_{p-1}&s_p
& \sigma_1&\sigma_2&\sigma_3\cdots&\sigma_{p-1}&\sigma_p\\ \hline
1)   &1&0&0\cdots&0&0&1&0&0\cdots&0&0\\
2)   &0&1&0\cdots&0&0&0&1&0\cdots&0&0\\
3)   &0&0&1\cdots&0&0&0&0&1\cdots&0&0\\
     & & &\cdots& & & & &\cdots& &\\
p-1) &0&0&0\cdots&1&0&0&0&0\cdots&1&0\\
p)   &0&0&0\cdots&0&1&0&0&0\cdots&0&1\\
p+1) &1&0&0\cdots&0&0&1&1&0\cdots&0&0\\
p+2) &0&1&0\cdots&0&0&0&1&1\cdots&0&0\\
p+3) &0&0&1\cdots&0&0&0&0&1\cdots&0&0\\
     & & &\cdots& & & & &\cdots& &\\
2p-1)&0&0&0\cdots&1&0&0&0&0\cdots&1&1\\
2p)  &0&0&0\cdots&0&1&1&0&0\cdots&0&1
\end{array}
\]
In dreien der vier Quadrate, in welche dies Schema zerfällt, sind nur die Diagonalelemente von \(0\) verschieden, im letzten auch noch die rechts neben diesen stehenden. Betrachtet man das Schema als Determinante und zieht immer die \(i^{\text{te}}\) Reihe von der \((p+i)^{\text{ten}}\) ab, so bleibt eine Determinante übrig, welche offenbar gleich \(\pm1\) ist.

Hieraus folgt, dass man durch solche \(2p\) Systeme von \(U_h\) die ursprünglichen Systeme von Periodicitätsmoduln
\[
\left(
\begin{array}{cccc|cccc}
2\pi\sqrt{-1}&0&\cdots&0&a_{11}&a_{12}&\cdots&a_{1p}\\
0&2\pi\sqrt{-1}&\cdots&0&a_{21}&a_{22}&\cdots&a_{2p}\\
\vdots&\vdots&\ddots&\vdots&\vdots&\vdots&\ddots&\vdots\\
0&0&\cdots&2\pi\sqrt{-1}&a_{p1}&a_{p2}&\cdots&a_{pp}
\end{array}
\right)
\]
linear und ganzzahlig ausdrücken kann; und man kann also auch die Systeme der \(G_h\), so wie die übrigen Systeme der \(U_h\) durch jene \(2p\) Systeme der \(U_h\) linear und ganzzahlig darstellen.

Sei daher irgend ein System der \(G_h\) mit jenen \(2p\) Systemen der \(U_h\) verbunden durch die Gleichungen
\[
G_h=\tau^{(1)}U_h^{(1)}+\tau^{(2)}U_h^{(2)}
+\cdots+\tau^{(2p)}U_h^{(2p)};
\qquad (h=1,2\ldots p)
\]
wo die \(\tau\) ganze Zahlen sind. Dividiren wir diese Gleichung durch \(2\), und ersetzen die halben Perioden immer durch die Integralsummen, denen sie in dem Problem der speciellen Theilung gleich werden. An Stelle von \(\frac12G_h\) tritt dann die Summe
\[
\sum_{i=1}^{i=p}\int_{\alpha^{(i)}}^{\beta^{(i)}}du_h,
\]
in welcher die \(\alpha,\beta\) Berührungspuncte zweier Curven \((n-2)^{\text{ter}}\) Ordnung sind, deren eine (die der \(\beta\)) von dem gewählten Systeme der \(G_h\) abhängt; die \(\frac12U_h\) aber sind zu ersetzen durch Summen der Form
\[
\sum_{i=1}^{i=p-1}\int_{\alpha^{(i)}}^{\gamma^{(i)}}du_h
+\int_{\alpha^{(p)}}^{\mu}du_h,
\]
in welche die \(\gamma\) Berührungspuncte einer Curve \((n-3)^{\text{ter}}\) Ordnung sind. Entsprechen nun den \(\alpha,\beta\) die Wurzeln \(\nu,\nu'\) von \(N=0\), den verschiedenen Systemen der \(\gamma\) die Wurzeln \(\varrho^{(1)},\varrho^{(2)}\ldots\varrho^{(2p)}\) von \(R=0\), so ist nach der mehrfach gemachten Anwendung des Abel'schen Theorems \(\nu'\) eine rationale Function von \(\nu,\varrho^{(1)},\varrho^{(2)}\ldots\varrho^{(2p)}\). Aber endlich ist für jedes nicht unter den obigen \(2p\) Systemen enthaltene System der \(U_h\):
\[
U_h=v^{(1)}U_h^{(1)}+v^{(2)}U_h^{(2)}
+\cdots+v^{(2p)}U_h^{(2p)},\qquad (h=1,2\ldots p)
\]
wo die \(v\) ganze Zahlen sind. Man kann leicht zeigen, dass die Summe der Zahlen \(v\) immer gleich der Summe der Zahlen \(s\) ist, welche in dem ursprünglichen Ausdruck des Systems der \(U_h\) auftreten. Man kann daher immer ein solches System \(U_h\) aufstellen, für welches die Summe der \(v\) gleich \(2\) ist. In Bezug auf die im obigen Schema gegebenen \(U^{(1)},U^{(2)}\ldots U^{(2p)}\) ist diese Bedingung z. B. durch das System der \(U_h\) erfüllt, welches durch die Zahlen
\[
1\quad1\quad0\cdots0\quad0\ .\quad1\quad0\quad0\cdots0\quad0
\]
charakterisirt ist, und für welches
\[
U_h=2U_h^{(1)}+U_h^{(2)}-U_h^{(p+1)}.
\]
Dividirt man diese Gleichung durch \(2\) und trägt für die halben Periodicitätsmoduln wieder die Integralsummen ein, in deren obern Grenzen Systeme der \(\gamma\) mit \(\mu\), in deren untern Grenzen nur die Puncte \(\alpha\) auftreten, so kann man in der entstehenden Gleichung die Grenzen offenbar so zusammenziehen, dass die untern Grenzen \(\alpha\) nur noch in einer Integralsumme erscheinen, welche überdies mit \(1\) multiplicirt ist. Das Abel'sche Theorem, auf die so gebildete Gleichung angewandt, lehrt dann, dass die Wurzel \(\nu\) von \(N=0\), welche der Curve der \(\alpha\) entspricht, rational ausdrückbar ist durch die Wurzeln \(\varrho,\varrho^{(1)},\varrho^{(2)},\varrho^{(p+1)}\), welche den benutzten Berührungscurven \((n-3)^{\text{ter}}\) Ordnung zugehören.

Auf diese Weise ist \(\nu\), und damit auch jede andre Wurzel \(\nu'\) von \(N=0\) rational durch Wurzeln von \(R=0\) ausgedrückt, was zu beweisen war. Uebrigens erleiden diese Betrachtungen natürlich eine Ausnahme für \(p=0\) und \(p=1\).

Wie nun endlich die specielle Theilung für \(m=2\) auf die soeben behandelten Aufgaben zurückgeführt wird, liegt auf der Hand. Wenn man die der Gleichung
\[
\sum_{i=1}^{i=p}\int_{c^{(i)}}^{x^{(i)}}du_h=\frac{P_h}{2}
\qquad (h=1,2\ldots p)
\]
entsprechende Lösung des speciellen Theilungsproblems sucht, so ersetzt man \(\frac{P_h}{2}\) durch die Integralsumme des entsprechenden Falles jener Aufgabe, und findet dann die symmetrischen Functionen der \(x\), also auch eine Wurzel der speciellen Theilungsgleichung, durch die Wurzeln von \(R=0\) und \(N=0\) rational ausgedrückt. Es verdient bemerkt zu werden, dass hienach die Lösungen des speciellen Theilungsproblems sich ebenfalls in zwei Classen sondern, je nachdem \(P_h\) der Classe der \(U_h\) oder der Classe der \(G_h\) angehört, deren erstere \(2^{p-1}(2^p-1)\), deren letztere \(2^{p-1}(2^p+1)-1\) Periodensysteme umfasst. Die Unterscheidung der beiden Classen von Lösungen auf geometrischem Wege erfolgt sofort durch die geometrische Deutung für den Fall des Abel'schen Theorems, welcher zur algebraischen Lösung benutzt wurde.

% --- Source pages 292--299; printed pages 270--277. ---

\section*{Elfter Abschnitt.\\Das erweiterte Umkehrproblem.}

\subsection*{\S\,77. Zurückführung des erweiterten Umkehrproblems auf eine von nur \(p+1\) Puncten abhängige Function \(U^{(q)}\).}

Wir haben im sechsten Abschnitt eine Lösung des erweiterten Umkehrproblems ausführlich dargestellt, durch welche dasselbe auf das Jacobi'sche zurückgeführt wurde. Aber wenn es sich darum handelt, jenes Problem auf \(\Theta\)functionen zurückzuführen, so ist es zweckmässiger, den andern Weg einzuschlagen, welcher in \(\S\,44\) kurz angedeutet ist. Dieser Weg behandelt das erweiterte Umkehrproblem vollkommen nach der Analogie des Jacobi'schen, indem er dasselbe auf die Transcendente
\[
T_{\xi\eta}^{(q)}
\begin{pmatrix}
x^{(1)},&x^{(2)}&\ldots\\
c^{(1)},&c^{(2)}&\ldots
\end{pmatrix}
=\sum_{i=1}^{i=p+q}\int_{c^{(i)}}^{x^{(i)}}d\Pi_{\xi\eta}
\]
zurückführt. Da diese ganz ebenso wie oben die Transcendente \(T_{\xi\eta}\) behandelt werden kann, so wird es genügen, die Analogieen hervorzuheben, und die eintretenden Besonderheiten kurz darzustellen.

Zunächst aber wollen wir, um die in diesen Erweiterungen auftretenden Sätze einfacher aussprechen zu können, das erweiterte Umkehrproblem auf eine speciellere Fassung zurückführen, welche demselben ohne Beeinträchtigung der Allgemeinheit immer gegeben werden kann. Durch eine eindeutige Transformation nämlich können wir immer die gegebene Curve \(f=0\) in eine andre verwandeln, bei welcher gegebene Punctepaare von \(f=0\) sich in Doppelpuncte vereinigen. Dabei werden noch andre Doppelpuncte auftreten, welche diesen nothwendig conjugirt sind, auf die es hier aber zunächst nicht ankommt. Wir denken uns diese eindeutige Transformation so eingerichtet, dass die in dem erweiterten Umkehrproblem des sechsten Abschnitts auftretenden Punctepaare \(\xi^{(1)},\eta^{(1)};\ \xi^{(2)},\eta^{(2)},\ldots;\ \xi^{(q)},\eta^{(q)}\) sich in Doppelpuncte vereinigen. Mit andern Worten, wir wollen annehmen, dass die in dem erweiterten Umkehrproblem gegebenen Integrale dritter Gattung nur in \emph{Doppelpuncten} der gegebenen Curve unstetig werden; einige dieser Puncte können auch in Rückkehrpuncte übergehen, was später berücksichtigt werden soll. Nehmen wir also zunächst an, die \(q\) Grundintegrale dritter Gattung seien in \(q\) Doppelpuncten von \(f=0\) unstetig; während ausserdem \(\delta\) Doppelpuncte und \(r\) Rückkehrpuncte vorhanden sein mögen, so dass also
\[
p=\frac{(n-1)\cdot(n-2)}{2}-q-\delta-r.
\]
Da die Grundintegrale dritter Gattung jetzt äusserlich die Form von Integralen erster Gattung haben, so wollen wir sie durch
\[
u_{p+1},\ u_{p+2}\ldots u_{p+q}
\]
bezeichnen, und erhalten also das Umkehrproblem in der Form:
\[
v_h=\sum_{i=1}^{i=p+q}\int_{c^{(i)}}^{x^{(i)}}du_h,
\qquad (h=1,2\ldots p+q).
\]
Wendet man nun nach der Analogie des \(\S\,42\) das Abel'sche Theorem auf einen Curvenbüschel \(\varphi-\lambda\psi=0\) an, welches durch die \(\delta+r\) Ausnahmspuncte geht, so findet man die Lösung des erweiterten Umkehrproblems, wie a.\ a.\ O. die des Jacobi'schen, unter der Form:
\[
\prod_{i=1}^{i=p+q}
\frac{\left(\dfrac{\varphi}{\psi}\right)_{x^{(i)}}-\lambda}
     {\left(\dfrac{\varphi}{\psi}\right)_{c^{(i)}}-\lambda}
=
e^{-\sum T_{\xi\eta}^{(q)}
\begin{pmatrix}
x^{(1)},x^{(2)}\ldots x^{(p+q)}\\
c^{(1)},c^{(2)}\ldots c^{(p+q)}
\end{pmatrix}},
\]
wobei die Summe sich über alle Punctepaare \(\xi,\eta\) erstreckt, und wo die \(\xi\) die Schnittpuncte von \(f=0\) mit \(\varphi-\lambda\psi=0\), die \(\eta\) aber die Schnittpuncte von \(f=0\) mit \(\psi=0\) bedeuten.

Bei der Anwendung des Abel'schen Theorems hat man hier immer Curven zu benutzen, welche durch die \(\delta+r\) übrigen Ausnahmspuncte, nicht aber durch diejenigen gehen, welche die Unstetigkeitspuncte der im Umkehrproblem gegebenen Integrale dritter Gattung sind. In dem Schnittpunctsystem einer solchen Curve sind, wenn ihre Ordnung die \((n-3)^{\mathrm{te}}\) übersteigt, immer \(p+q\) Puncte durch die übrigen bestimmt (\(\S\,9\)), dagegen nur \(p+q-1\), wenn die Curve von der \((n-3)^{\mathrm{ten}}\) Ordnung ist. Legt man insbesondere eine Curve \((n-2)^{\mathrm{ter}}\) Ordnung durch die \(\delta+r\) Ausnahmspuncte, und durch die \(n-2\) Puncte, in welchen die Curve \(f=0\) von einer Geraden \(\xi\eta\) noch geschnitten wird, und durch \(p+q\) gegebene Puncte \(x\), so schneidet sie \(f=0\) noch in \(p+q\) andern Puncten \(y\). Eine Transcendente \(T_{\xi\eta}^{(q)}\), welche die \(x\) zu obern und die zugehörigen \(y\) zu untern Grenzen hat, wollen wir als specielle Transcendente \(T_{\xi\eta}^{(q)}(x^{(1)},x^{(2)}\ldots x^{(p+q)})\) bezeichnen. Man kann dann, genau wie in \(\S\,45\), die allgemeine Transcendente \(T_{\xi\eta}^{(q)}\) durch zwei specielle ausdrücken mit Hülfe der Formel:
\[
T_{\xi\eta}^{(q)}
\begin{pmatrix}
x^{(1)},x^{(2)}\ldots\\
c^{(1)},c^{(2)}\ldots
\end{pmatrix}
=\frac12\left[
T_{\xi\eta}^{(q)}(x^{(1)},x^{(2)}\ldots)
-
T_{\xi\eta}^{(q)}(c^{(1)},c^{(2)}\ldots)
\right].
\]
Man kann nun auch die Theorie der adjungirten Transcendenten auf diesen Fall ausdehnen. Legt man durch die \(\delta+r\) Ausnahmspuncte und durch \(p+q-1\) der \(x\) eine Curve \((n-3)^{\mathrm{ter}}\) Ordnung, so schneidet sie noch in \(p+q-1\) Puncten, welche je nach dem ausgelassenen Puncte \(x\) wie folgt bezeichnet werden sollen:
\[
\begin{array}{c@{\qquad}l}
x^{(1)}&x^{21},\ x^{31}\ldots\\
x^{(2)}&x^{12},\ x^{32}\ldots\\
x^{(3)}&x^{13},\ x^{23}\ldots\\
&\ldots
\end{array}
\]
Bildet man dann die adjungirten Transcendenten
\[
T^{(1)}=T_{x^{(1)}\eta}(\xi,x^{21},x^{31}\ldots),
\qquad
T^{(2)}=T_{x^{(2)}\eta}(x^{12},\xi,x^{32}\ldots),
\]
und bezeichnet ausserdem durch
\[
T^0=T_{\xi\eta}^{(q)}(x^{(1)},x^{(2)},x^{(3)}\ldots)
\]
die ursprüngliche Transcendente, so ist wieder
\[
dU^{(q)}=\frac12\left[
\frac{\partial T^0}{\partial\xi}\,d\xi
+\frac{\partial T^{(1)}}{\partial x^{(1)}}\,dx^{(1)}
+\frac{\partial T^{(2)}}{\partial x^{(2)}}\,dx^{(2)}
+\cdots+
\frac{\partial T^{(p+q)}}{\partial x^{(p+q)}}\,dx^{(p+q)}
\right]
\]
ein vollständiges Differential. Indem man constante Hülfspuncte \(\alpha\), und sodann constante untere Grenzen \(\zeta,\gamma\) einführt, erhält man genau wie im siebenten Abschnitt einen Ausdruck für die Function \(U^{(q)}\) selbst, nämlich:
\[
\begin{aligned}
U^{(q)}(x^{(1)},x^{(2)},\ldots;\xi)
={}&
\sum_{i=1}^{i=p+q}\int_{\zeta}^{\xi}d\Pi_{x^{(i)}\alpha^{(i)}}
+\sum_{i=1}^{i=p+q}\int_{\gamma^{(i)}}^{x^{(i)}}d\Pi_{\xi\alpha^{(i)}}\\
&-\frac12\sum_{i>k}\sum\int_{\gamma^{(k)}}^{x^{(k)}}d\Pi_{x^{(i)}\alpha^{(i)}}
-\frac12\sum_{i>k}\sum\int_{\gamma^{(i)}}^{x^{(i)}}d\Pi_{\gamma^{(k)}\alpha^{(k)}}\\
&+\frac12 T_{\xi\zeta}^{(q)}(\alpha^{(1)},\alpha^{(2)}\ldots\alpha^{(p+q)})\\
&+\frac12\sum_{i=1}^{i=p+q}
T_{x^{(i)}\gamma^{(i)}}^{(q)}
(\alpha^{1i},\alpha^{2i}\ldots\alpha^{i-1,i},\alpha^{(i)},\alpha^{i+1,i}\ldots\alpha^{p+q,i})\\
&+\log\frac{%
\begin{gathered}\scriptstyle
\varphi(x^{(1)},x^{(2)}\ldots x^{(p+q)})\,
\varphi(\gamma^{(1)},\alpha^{(2)}\ldots\alpha^{(p+q)})\\[-0.15em]
\scriptstyle{}\cdot
\varphi(\alpha^{(1)},\gamma^{(2)}\ldots\alpha^{(p+q)})\cdots
\varphi(\alpha^{(1)},\alpha^{(2)}\ldots\gamma^{(p+q)})
\end{gathered}}{%
\begin{gathered}\scriptstyle
\varphi(\gamma^{(1)},\gamma^{(2)}\ldots\gamma^{(p+q)})\,
\varphi(x^{(1)},\alpha^{(2)}\ldots\alpha^{(p+q)})\\[-0.15em]
\scriptstyle{}\cdot
\varphi(\alpha^{(1)},x^{(2)}\ldots\alpha^{(p+q)})\cdots
\varphi(\alpha^{(1)},\alpha^{(2)}\ldots x^{(p+q)})
\end{gathered}}.
\end{aligned}
\]
Die Function \(\varphi\) giebt dabei, gleich Null gesetzt, die Bedingung dafür an, dass die als Argumente ihr beigesetzten Puncte mit den \(\delta+r\) Ausnahmspuncten auf einer Curve \((n-3)^{\mathrm{ter}}\) Ordnung liegen.

Man beweist nun, wie in \(\S\,51\), die folgenden beiden Sätze:

I. Wenn man die Function \(U^{(q)}\) für \(x^{(1)},x^{(2)}\ldots x^{(p+q)}, \xi\) und für \(y^{(1)},y^{(2)}\ldots y^{(p+q)}, \eta\) bildet, wo die \(x,y\) auf einer Curve \((n-2)^{\mathrm{ter}}\) Ordnung liegen, welche durch die \(\delta+r\) abgesonderten Ausnahmspuncte, sowie durch die \(n-2\) übrigen Schnittpuncte von \(f=0\) mit der Geraden \(\xi\eta\) hindurchgeht, so ergiebt sich bei passender Wahl der Integrationswege beide Male derselbe oder der entgegengesetzte Werth.

II. Lassen wir einen Strahl sich um einen Punct der Curve \(f=0\) drehen, und legen in jeder Lage desselben eine Curve \((n-2)^{\mathrm{ter}}\) Ordnung durch \(n-2\) seiner andern Schnittpuncte, durch die \(\delta+r\) Ausnahmspuncte von \(f=0\) und durch \(p+q\) feste Puncte von \(f=0\), so schneidet dieselbe noch in \(p+q\) Puncten \(x\), und die Function \(U^{(q)}\), für diese \(p+q\) Puncte gebildet, und mit einem Parameter, welcher dem letzten Schnittpunct des Strahls entspricht, hat immer denselben Werth.

Der zweite Satz versieht wieder die Stelle einer partiellen Differentialgleichung, welcher die Function \(U^{(q)}\) genügt; er sagt aus, dass diese Function nur von den \(p+q\) Summen
\[
\sum_{h=1}^{h=p+q}\int_{c^{(h)}}^{x^{(h)}}du_k-\int_{\mu}^{\xi}du_k,
\qquad (k=1,2\ldots p+q)
\]
oder von den \(p+q\) Argumenten
\[
v_k-\int_{\mu}^{\xi}du_k
\]
abhänge, wo \(\mu\) ein beliebiger constanter Punct sein kann.

\subsection*{\S\,78. Unstetigkeiten von \(U^{(q)}\). Die synektische Function \(V^{(q)}=e^{-U^{(q)}}\).}

Bis hieher ist die Analogie mit den Untersuchungen des siebenten Abschnitts so vollständig, dass man die Resultate demselben ohne Weiteres entnehmen kann. Aber dies ändert sich in etwas, wenn wir zur Untersuchung der Unstetigkeiten von \(U^{(q)}\) übergehen. Diese Function wird erstlich, wie die Function \(U\) des siebenten Abschnitts, unendlich gross,

\[
\begin{array}{ll}
1)&\text{wenn }\xi\text{ mit einem der Puncte }x\text{ zusammenfällt;}\\
2)&\text{wenn die Puncte }x\text{ auf einer Curve }(n-3)^{\mathrm{ter}}\text{ Ordnung liegen,}\\
&\text{welche durch die }\delta+r\text{ übrigen Ausnahmspuncte geht,}
\end{array}
\]
und nähert sich in beiden Fällen dem negativen Logarithmus einer einfach verschwindenden algebraischen Function. Aber die Function \(U^{(q)}\) wird noch in einem dritten Falle unendlich gross, nämlich, wenn einer der Puncte \(\xi,x\) sich einem der \(q\) Doppelpuncte nähert, welche in dem erweiterten Umkehrproblem ausgezeichnet sind. Um zu ermitteln, welchem Werthe in diesem Fall die Function \(U^{(q)}\) sich nähert, nehmen wir etwa an, dass \(x^{(1)}\) einem jener Doppelpuncte unendlich nahe rückt. Dann wird in \(U^{(q)}\) der Term
\[
\frac12 T_{x^{(1)}\gamma^{(1)}}^{(q)}(\alpha^{(1)},\alpha^{21},\alpha^{31}\ldots\alpha^{p+q,1})
\]
unendlich gross. Denn die Curve \((n-2)^{\mathrm{ter}}\) Ordnung, durch welche man die untern Grenzpuncte construirt, geht durch die \(n-2\) Puncte, in welchen \(f=0\) von der Geraden \(x^{(1)}\gamma^{(1)}\) noch geschnitten wird. Liegt nun \(x^{(1)}\) in der Nähe eines Doppelpuncts auf dem einen Zweige von \(f=0\), so schneidet diese Gerade in der Nähe desselben noch einmal, und zwar den andern Zweig von \(f=0\); daher muss die Curve \((n-2)^{\mathrm{ter}}\) Ordnung durch einen dem Doppelpuncte unendlich nahen Punct der Geraden \(x^{(1)}\gamma^{(1)}\) gehen, und schneidet daher in der Nähe des Doppelpuncts noch ein zweites Mal, und zwar wieder auf dem ersten Zweige von \(f=0\). Es muss also ein weiterer Schnittpunct, d.\ h. ein unterer Grenzpunct, dem Doppelpunct unendlich nahe rücken. Sei dieser Punct \(\beta^{(1)}\); dann wird der unstetige Theil der obigen Transcendente das Glied
\[
\frac12\int_{\beta^{(1)}}^{\alpha^{(1)}}d\Pi_{x^{(1)}\gamma^{(1)}},
\]
in welche das Argument \(\beta^{(1)}\) und der Parameter \(x^{(1)}\) einander unendlich nahe rücken, und zwar auf demselben der beiden Zweige von \(f=0\), welche sich in dem fraglichen Doppelpuncte durchschneiden.

Wie dieses Integral sich in der Grenze verhält, sehen wir aus der Formel (5) des \(\S\,58\), welche ein einzelnes Integral dritter Gattung durch \(\Theta\)functionen ausdrücken lehrt. Nach derselben verhält sich obiges Integral bei der Annäherung von \(\beta^{(1)}\) an \(x^{(1)}\), wie der Logarithmus einer \(\Theta\)function, welche durch Zusammenfallen des Parameters \(x^{(1)}\) mit dem Argument \(\beta^{(1)}\) verschwindet, also wie \(\log(\beta^{(1)}-x^{(1)})\). Es nähert sich also \(T^{(q)}\) dem Werthe \(\frac12\log(\beta^{(1)}-x^{(1)})\).

Inzwischen wird durch Annäherung von \(x^{(1)}\) an einen der \(q\) ausgezeichneten Doppelpuncte die betreffende Grösse \(v_{p+i}\) ebenfalls unendlich gross wegen des Integrals
\[
\int_{c^{(1)}}^{x^{(1)}}du_{p+i},
\]
in welchem das Argument \(x^{(1)}\) sich einem der Parameter nähert; und zwar demjenigen der beiden im Doppelpuncte vereinigt gedachten Puncte, auf dessen Zweig \(x^{(1)}\) sich bewegt. Bezeichnen wir daher durch \(\delta\) diesen Doppelpunct, so verhält sich das obige Integral, mithin auch \(v_i\) in der Grenze wie \(\pm\log(\delta-x^{(1)})\); und zwar ist das obere oder das untere Zeichen zu wählen, je nachdem \(x^{(1)}\) sich auf dem einen oder auf dem andern Zweige bewegt, je nachdem also \(x^{(1)}\) dem einen oder dem andern Parameter sich nähert.

Bei der erwähnten Bewegung der Puncte \(\beta\) und des Punctes \(x^{(1)}\) behalten nach dem Abel'schen Theorem die oben als Argumente einer Function \(U^{(q)}\) erwähnten \(p+q\) Ausdrücke
\[
\sum_{h=1}^{h=p+q}\int_{c^{(h)}}^{\beta^{(h)}}du_k-\int_{\mu}^{x^{(1)}}du_k,
\qquad (k=1,2\ldots p+q)
\]
constante endliche Werthe. Aber wenn \(x^{(1)}\) und mit ihm eines der \(\beta\) gegen den Doppelpunct \(\delta\) rückt, so werden in dem Ausdruck
\[
\sum_{h=1}^{h=p+q}\int_{c^{(h)}}^{\beta^{(h)}}du_{p+i}-\int_{\mu}^{x^{(1)}}du_{p+1}
\]
die beiden Integrale
\[
\int_{c^{(1)}}^{\beta^{(1)}}du_{p+i}-\int_{\mu}^{x^{(1)}}du_{p+i}
\]
unendlich gross, und nähern sich den Ausdrücken
\[
\pm\left[\log(\beta^{(1)}-\delta)-\log(x^{(1)}-\delta)\right].
\]
Daher bleibt die Grösse
\[
\log\frac{\beta^{(1)}-\delta}{x^{(1)}-\delta}
\]
endlich, oder die Summe
\[
T_{x^{(1)}\gamma^{(1)}}^{(q)}(\alpha^{(1)},\alpha^{21}\ldots)\pm v_{p+i}
\]
behält einen endlichen Werth. Man kann daher sagen, die Function \(U^{(q)}\) nähere sich dem unendlich grossen Werthe \(\pm\frac12v_{p+i}\), wenn einer der Puncte \(x\) sich dem \(i^{\mathrm{ten}}\) der \(q\) ausgezeichneten Doppelpuncte nähert.

Bemerken wir hienach, dass die Function \(U^{(q)}\) nur unendlich werden kann, wenn eine der Grössen \(v\) selbst einen unendlich grossen Werth annimmt, oder in den oben in (1), (2) angeführten Fällen, in denen sie sich dem negativen Logarithmus einer einfach verschwindenden Grösse nähert, so schliesst man, wie in \(\S\,53\), dass die Function
\[
V^{(q)}=e^{-U^{(q)}}
\]
eine synektische Function der \(p+q\) Argumente
\[
\sum_{h=1}^{h=p+q}\int_{c^{(h)}}^{x^{(h)}}du_k-\int_{\mu}^{\xi}du_k,
\qquad (k=1,2\ldots p+q)
\]
ist. Und diese Function verschwindet nur, und immer dann, wenn entweder \(\xi\) mit einem der \(x\) zusammenfällt, oder wenn die Puncte \(x\) auf einer Curve \((n-3)^{\mathrm{ter}}\) Ordnung liegen, welche durch die \(\delta+r\) abgesonderten Ausnahmspuncte hindurchgeht.

Der oben unter I. aufgeführte Satz giebt nun noch die Eigenschaft der Function \(V^{(q)}\) an (vgl. \(\S\,53\)), dass sie, als Function der \(p+q\) Argumente
\[
w_h=\sum_{i=1}^{i=p+q}\int_{c^{(i)}}^{x^{(i)}}du_h-\int_{\mu}^{\xi}du_h+K_h
\]
betrachtet, bei geeigneter Wahl der Constanten \(K\) bis auf das Vorzeichen ungeändert bleibt, wenn man alle \(w\) gleichzeitig das Zeichen wechseln lässt. Und zwar muss man die \(K\) so annehmen, dass, wenn die \(y\) solche \(p+q\) Puncte bezeichnen, welche mit den \(x\) und mit den \(n-2\) übrigen Schnittpuncten einer Geraden \(\xi\eta\), sowie mit den \(\delta+r\) Ausnahmspuncten auf einer Curve \((n-2)^{\mathrm{ter}}\) Ordnung liegen:
\[
(1)\quad
2K_h+
\sum_{i=1}^{i=p+q}\int_{c^{(i)}}^{x^{(i)}}du_h+
\sum_{i=1}^{i=p+q}\int_{c^{(i)}}^{y^{(i)}}du_h
-\int_{\mu}^{\xi}du_h-\int_{\mu}^{\eta}du_h=0,
\]
wodurch sich nach dem Abel'schen Theorem für die \(K_h\) constante Werthe ergeben.

\subsection*{\S\,79. Reihenentwicklung der Function \(V^{(q)}\) aus ihren Periodicitätseigenschaften.}

Wir bestimmen nun den Ausdruck der synektischen Function \(V^{(q)}\) vollständig durch ihre Periodicitätsmoduln. Bilden wir zunächst die Transcendenten
\[
\frac12T_{\xi\eta}^{(q)}(x^{(1)},x^{(2)}\ldots x^{(p+q)})
=
\log
\frac{V^{(q)}(x^{(1)},x^{(2)}\ldots x^{(p+q)};\eta)}
     {V^{(q)}(x^{(1)},x^{(2)}\ldots x^{(p+q)};\xi)},
\]
\[
T_{\xi\eta}^{(q)}
\begin{pmatrix}
x^{(1)},x^{(2)}\ldots x^{(p+q)}\\
z^{(1)},z^{(2)}\ldots z^{(p+q)}
\end{pmatrix}
=
\log
\frac{
V^{(q)}(x^{(1)},x^{(2)}\ldots;\eta)
V^{(q)}(z^{(1)},z^{(2)}\ldots;\xi)
}{
V^{(q)}(x^{(1)},x^{(2)}\ldots;\xi)
V^{(q)}(z^{(1)},z^{(2)}\ldots;\eta)
}.
\]
Aus der letzten Formel erhalten wir genau wie in \(\S\,54\) die Veränderung, welche \(V^{(q)}\) erfährt, wenn die \(w\), von denen \(V^{(q)}(x^{(1)},x^{(2)}\ldots;\xi)\) abhängt, um Periodicitätsmoduln wachsen. Die allgemeinste Veränderung der Grössen \(w_1,w_2\ldots w_p\) besteht darin, dass \(w_i\) in
\[
w_i'=w_i+2m_i\pi\sqrt{-1}+a_{1i}q_1+a_{2i}q_2+\cdots+a_{pi}q_p,
\qquad (i=1,2\ldots p)
\]

% --- Source pages 300--308; printed pages 278--286. ---

übergeht. Dann ändern sich die übrigen \(w\) um die Periodicitätsmoduln, welche den in ihnen auftretenden Integralen dritter Gattung zukommen, und man kann ausserdem noch von einander unabhängige Periodicitätsmoduln hinzufügen, welche ganze Vielfache von \(2\pi\sqrt{-1}\) sind. Bezeichnen wir durch \(u_{ki}\) die Grösse, um welche das Integral \(u_{p+i}\) zunimmt, wenn die \(w_1,w_2\ldots w_p\) um die Periodicitätsmoduln \(a_{1k},a_{2k}\ldots a_{pk}\) wachsen, so werden die allgemeinsten Werthe, welche die Grössen \(w_{p+1},w_{p+2}\ldots\) zugleich mit dem Uebergange von \(w_1,w_2\ldots w_p\) in \(w'_1,w'_2\ldots w'_p\) annehmen können, folgende:
\[
w'_{p+i}=w_{p+i}+2n_i\pi\sqrt{-1}+u_{1i}q_1+u_{2i}q_2+\cdots+u_{pi}q_p,
\qquad (i=1,2\ldots q).
\]
Bezüglich der \(u_{ki}\) bemerken wir, dass wir für jedes Integral \(u_{p+i}\) in der imaginären Ebene die im Doppelpunct vereinigten Parameterpuncte verbunden denken müssen durch eine Curve, welche von dem den Doppelpunct repräsentirenden Puncte ausgeht und zu ihm zurückkehrt, so aber, dass der Anfang dieser Curve dem einen durch jenen Doppelpunct entsprechenden Zweige von \(f=0\) entspricht, das Ende dem andern. Das Integral \(u_{ki}\) ist dann nach \(\S\,34\), (6) das Integral \(u_k\) über diese Curve geführt, wobei von vornherein immer festgesetzt sein muss, welchen Zweig von \(f=0\) man dem ersten, und welchen man dem zweiten Parameter entsprechen lässt, mit welchem also dieser Integrationsweg beginnt und mit welchem er endigt.

Während also die \(p+q\) Grössen \(w\) sich in die \(p+q\) Grössen \(w'\) verwandeln, geht \(V^{(q)}\) über in
\[
(2)\quad
V'^{(q)}=e^{w_1q_1+w_2q_2+\cdots+w_pq_p+C}\cdot V^{(q)},
\]
wobei \(C\) eine Constante ist. Und genau wie a. a. O. findet man für \(C\) den Ausdruck
\[
C=\frac12\sum\sum q_hq_ka_{hk}+L\pi\sqrt{-1},
\]
wo \(L\) eine ganze Zahl ist.

Bilden wir \(L\) für das hier benutzte Zahlensystem der \(m,n,q\), und bezeichnen durch \(L'\), \(L''\) die Zahlen, in welche \(L\) übergeht, wenn die \(m,n,q\) beziehungsweise durch andre Zahlen \(m',n',q'\), oder durch die Summen \(m+m',n+n',q+q'\) ersetzt werden, so hat man wie a. a. O.
\[
L+L'=L'',
\]
und schliesst daraus
\[
\begin{aligned}
L={}&\varkappa_1m_1+\varkappa_2m_2+\cdots+\varkappa_pm_p
+\lambda_1q_1+\lambda_2q_2+\cdots+\lambda_pq_p\\
&+\nu_1n_1+\nu_2n_2+\cdots+\nu_qn_q,
\end{aligned}
\]
wo die \(\varkappa,\lambda,\nu\) bestimmte ganze Zahlen bedeuten, welche nur von der ursprünglichen Definition der Transcendente \(\frac12T_{\xi\eta}^{(q)}(x^{(1)},x^{(2)}\ldots)\) abhängen, und denen es genügt die Werthe 0 oder 1 beizulegen.

Führen wir nun statt der \(w\) die Argumente
\[
\begin{aligned}
\omega_h&=w_h-\lambda_h\pi\sqrt{-1}-\frac12(\varkappa_1a_{1h}+\varkappa_2a_{2h}+\cdots+\varkappa_pa_{ph}),
&& (h=1,2\ldots p),\\
\omega_{p+h}&=w_{p+h}-\nu_h\pi\sqrt{-1}-\frac12(\varkappa_1u_{1h}+\varkappa_2u_{2h}+\cdots+\varkappa_pu_{ph}),
&& (h=1,2\ldots q),
\end{aligned}
\]
ein, welche sich von jenen um halbe Periodicitätsmoduln unterscheiden, so kann man der Gleichung (2) oder der aus ihr entstandenen Gleichung
\[
(3)\quad
V'^{(q)}=
 e^{\sum n_hq_h+\frac12\sum\sum q_hq_ka_{hk}+
 \pi\sqrt{-1}(\sum\varkappa_hm_h+\sum\lambda_hq_h+\sum\nu_hn_h)}\cdot V^{(q)}
\]
nunmehr die Form geben:
\[
\begin{aligned}
&e^{-\frac12(\varkappa_1\omega'_1+\varkappa_2\omega'_2+\cdots+\nu_1\omega'_{p+1}+\nu_2\omega'_{p+2}+\cdots)}\,V'^{(q)}\\
&=e^{-\frac12(\varkappa_1\omega_1+\varkappa_2\omega_2+\cdots+\nu_1\omega_{p+1}+\nu_2\omega_{p+2}+\cdots)}\,V^{(q)}\\
&\qquad\cdot e^{q_1\omega_1+q_2\omega_2+\cdots+\frac12\sum\sum q_hq_ka_{hk}-\frac12\sum\sum q_h\nu_ku_{hk}},
\end{aligned}
\]
wobei die \(\omega'\) aus den \(\omega\) durch dieselben Zusätze von Periodicitätsmoduln hervorgegangen gedacht werden, durch welche die \(w'\) aus den \(w\) entstanden. Die Function
\[
e^{-\frac12(\varkappa_1\omega_1+\varkappa_2\omega_2+\cdots+\nu_1\omega_{p+1}+\nu_2\omega_{p+2}+\cdots)}\,V^{(q)}
\]
ändert sich also nicht, wenn die \(\omega\) unabhängig von einander um Vielfache von \(2\pi\sqrt{-1}\) wachsen, und sie ändert sich um den Factor
\[
e^{q_1\omega_1+q_2\omega_2+\cdots+q_p\omega_p+\frac12\sum\sum q_hq_ka_{hk}-\frac12\sum\sum q_h\nu_ku_{hk}},
\]
wenn die \(\omega\) beziehungsweise um die Systeme
\[
\begin{array}{rcl}
\omega_i&:&a_{1i}q_1+a_{2i}q_2+\cdots+a_{pi}q_p,\\
\omega_{p+i}&:&u_{1i}q_1+u_{2i}q_2+\cdots+u_{pi}q_p
\end{array}
\]
zunehmen.

Wegen der ersten Eigenschaft kann man die obige Function in eine stets convergente Reihe von der Form
\[
\sum_{r,s} A_{r_1r_2\ldots r_ps_1s_2\ldots s_q}
 e^{r_1\omega_1+r_2\omega_2+\cdots+r_p\omega_p+s_1\omega_{p+1}+s_2\omega_{p+2}+\cdots+s_q\omega_{p+q}}
\]
entwickeln, wobei die Zahlen \(r,s\) unabhängig von einander die sämmtlichen ganzzahligen Werthe von \(-\infty\) bis \(+\infty\) erhalten sollen. Die letzte Eigenschaft giebt dann die identische Relation:
\[
\begin{aligned}
&\sum_{r,s} A_{r_1r_2\ldots s_1s_2\ldots}
 e^{\sum r_i\omega_i+\sum s_i\omega_{p+i}+\sum\sum q_kr_ia_{ki}+\sum\sum q_ks_iu_{ki}}\\
&=e^{q_1\omega_1+q_2\omega_2+\cdots+q_p\omega_p+\frac12\sum\sum q_hq_ka_{hk}-\frac12\sum\sum q_h\nu_ku_{hk}}
\sum_{r,s} A_{r_1r_2\ldots s_1s_2\ldots}e^{\sum r_i\omega_i+\sum s_i\omega_{p+i}}.
\end{aligned}
\]
Setzt man links \(r_1+q_1,r_2+q_2\ldots\) an Stelle von \(r_1,r_2\ldots\) und vergleicht dann, so erhält man die Relation zwischen den Coefficienten:
\[
\begin{aligned}
&A_{r_1+q_1,r_2+q_2,\ldots,s_1,s_2,\ldots}
 e^{\sum\sum q_kr_ia_{ki}+\sum\sum q_ks_iu_{ki}+\sum\sum q_kq_ia_{ki}}\\
&=A_{r_1r_2\ldots s_1s_2\ldots}
 e^{\frac12\sum\sum q_hq_ka_{hk}-\frac12\sum\sum q_h\nu_ku_{hk}},
\end{aligned}
\]
oder, wenn man alle \(r\) verschwinden lässt und \(A_{s_1s_2\ldots}\) für \(A_{00\ldots s_1s_2\ldots}\) schreibt,
\[
A_{q_1q_2\ldots s_1s_2\ldots}=A_{s_1s_2\ldots}
 e^{-\frac12\sum\sum q_hq_ka_{hk}-\sum\sum q_ks_iu_{ki}-\frac12\sum\sum q_h\nu_ku_{hk}}.
\]
Die zu bestimmende Reihe verwandelt sich somit in:
\[
\begin{aligned}
&\sum_s\left\{\sum_r e^{-\frac12\sum\sum r_hr_ka_{hk}+\sum r_h\omega_h-
\sum\sum r_h\left(s_i+\frac12\nu_i\right)u_{hi}}\right\}
 A_{s_1s_2\ldots}e^{\sum s_i\omega_{p+i}}\\
&=\sum_s A_{s_1s_2\ldots s_q}e^{\sum s_i\omega_{p+i}}
 \Theta\left[\omega_h-\left(s_1+\frac{\nu_1}{2}\right)u_{h1}-\left(s_2+\frac{\nu_2}{2}\right)u_{h2}\cdots\right],
\end{aligned}
\]
wo bei der \(\Theta\)function nur ein Argument statt aller gesetzt ist. Für die Function \(V^{(q)}\) also hat man den Ausdruck:
\[
\begin{aligned}
V={}&e^{\frac12(\sum\varkappa_h\omega_h+\sum\nu_h\omega_{p+h})}\,\cdot
\sum_s A_{s_1s_2\ldots s_q}e^{\sum s_i\omega_{p+i}}\\
&\cdot\Theta\left[\omega_h-\left(s_1+\frac{\nu_1}{2}\right)u_{h1}-\left(s_2+\frac{\nu_2}{2}\right)u_{h2}\cdots\right].
\end{aligned}
\]
Dieser Ausdruck vereinfacht sich sehr wesentlich, wenn man sich daran erinnert (vgl. \(\S\,78\)), dass bei einem unendlich grossen Werthe von \(v_{p+i}\), also auch von \(w_{p+i}\), die Function sich verhält wie \(e^{\pm\frac12w_{p+i}}\). Daraus folgt, dass sämmtliche Zahlen \(\nu\) gleich 1 gesetzt werden müssen, und dass dann in der Summe nur diejenigen Glieder beibehalten werden dürfen, in denen die \(s\) die Werthe 0 oder \(-1\) annehmen, damit in der Entwicklung von \(V^{(q)}\) keine höhere Potenz von \(e^{\pm\frac12w_{p+i}}\) als die erste vorkommt. Bezeichnet man dann \(2s_i+\nu_i\) wieder durch \(s_i\), so wird der Ausdruck von \(V^{(q)}\):
\[
\begin{aligned}
(1)\quad V^{(q)}={}&e^{\frac12\sum\varkappa_h\omega_h}
\sum_{s_1,s_2,\ldots=\pm1} A_{s_1s_2\ldots s_q}
 e^{\frac12\sum s_i\omega_{p+i}}\\
&\cdot\Theta\left(\omega_h-\frac{s_1u_{h1}+s_2u_{h2}+\cdots}{2}\right),
\end{aligned}
\]
wo die Summe jetzt noch aus \(2^q\) Termen besteht, deren Coefficienten \(A_{s_1s_2\ldots}\) zu bestimmen bleiben.

\subsection*{\S\,80. Bestimmung der Coefficienten \(A\) der Function \(V^{(q)}\).}

Ehe wir zu der Bestimmung der Coefficienten \(A\) übergehen, wollen wir die Argumente der in \(V\) auftretenden \(\Theta\)functionen genauer betrachten. Die Grössen \(\omega\) können wir in der Form
\[
\omega_h=\sum_{i=1}^{i=p+q}\int_{\beta^{(i)}}^{x^{(i)}}du_h-\int_{\mu}^{\xi}du_h,
\qquad (h=1,2\ldots p+q)
\]
darstellen, wenn die untern Grenzpunkte \(\beta\) gehörig bestimmt werden. Man beweist leicht, dass die \(\beta\) die Berührungspuncte einer Curve \((n-2)^{\mathrm{ter}}\) Ordnung sind, welche durch die \(\delta+r\) abgesonderten Doppel- und Rückkehrpuncte und durch die \(n-2\) übrigen Schnittpuncte der in dem Puncte \(\mu\) an \(f=0\) gezogenen Tangente geht, und welche in \(p+q\) Puncten berührt. In der That gehen in der Gleichung (1) die \(y\) in die \(\beta\) über, wenn man die \(\beta\) an die Stelle der \(x\) setzt und \(\xi,\eta\) mit \(\mu\) zusammenfallen lässt. Man hat dann aus (1)
\[
K_h+\sum_{i=1}^{i=p+q}\int_{c^{(i)}}^{\beta^{(i)}}du_h
\]
gleich einem halben Periodicitätsmodul, und zwar ändert dieser sich je nach dem gewählten Systeme der \(\beta\). Daher unterscheidet
\[
w_h=v_h+K_h-\int_{\mu}^{\xi}du_h
\]
sich von dem Ausdruck
\[
\sum_{i=1}^{i=p+q}\int_{\beta^{(i)}}^{x^{(i)}}du_h-\int_{\mu}^{\xi}du_h
\]
nur durch die Hälfte eines Systems von Periodicitätsmoduln, und dasselbe gilt von den \(\omega_h\), wie auch das System der \(\beta\) gewählt sei; man kann dieses daher endlich so wählen, dass obiger Ausdruck mit \(\omega_h\) genau übereinstimmt.

Denken wir uns jetzt eine Curve \((n-2)^{\mathrm{ter}}\) Ordnung so bewegt, dass sie stets durch die \((n-2)\) übrigen Schnittpuncte der Tangente in \(\mu\), so wie durch die \(\delta+r\) abgesonderten Ausnahmspuncte geht. Die Ausgangscurve sei dieselbe, welche in den Puncten \(\beta^{(1)},\beta^{(2)}\ldots\beta^{(p+q)}\) berührt; die Endcurve sei die im 8\textsuperscript{ten} Abschnitt benutzte Curve, welche durch sämmtliche Doppel- und Rückkehrpuncte geht und in den \(p\) Puncten \(\alpha\) berührt, welche als untere Grenzpuncte in dem Jacobi'schen Umkehrproblem erschienen. Bildet man für den Uebergang der einen Curve in die andere die transcendenten Gleichungen des Abel'schen Theorems, so werden die \(q\) Gleichungen illusorisch, welche den \(q\) ausgezeichneten Doppelpuncten entsprechen, für die Integrale erster Gattung aber erhält man Gleichungen der Form:
\[
\begin{aligned}
(1)\quad
2\sum_{i=1}^{i=p}\int_{\beta^{(i)}}^{\alpha^{(i)}}du_h
+\sum_{i=1}^{i=q}\left(
\int_{\beta^{(p+i)}}^{\varepsilon^{(+i)}}du_h
+\int_{\beta^{(p+i)}}^{\varepsilon^{(-i)}}du_h\right)=0,\\
(h=1,2\ldots p),
\end{aligned}
\]
wenn \(\varepsilon^{(+i)},\varepsilon^{(-i)}\) die beiden in einem der bevorzugten Doppelpuncte vereinigten Puncte von \(f=0\) bedeuten, welche je nach dem Zweige der Curve unterschieden werden, auf welchem man sich dem Doppelpunct nähert. In Folge dieser Gleichung kann man den Grössen \(\omega_1,\omega_2\ldots\omega_p\) die Form geben:
\[
\begin{aligned}
\omega_h={}&\sum_{i=1}^{i=p}\int_{\alpha^{(i)}}^{x^{(i)}}du_h
+\frac12\sum_{i=1}^{i=q}\left(\int_{\varepsilon^{(+i)}}^{x^{(p+i)}}du_h+
\int_{\varepsilon^{(-i)}}^{x^{(p+i)}}du_h\right)-\int_{\mu}^{\xi}du_h,\\
&\hfill (h=1,2\ldots p),
\end{aligned}
\]
oder auch
\[
\omega_h=\sum_{i=1}^{i=p}\int_{\alpha^{(i)}}^{x^{(i)}}du_h
+\sum_{i=1}^{i=q}\int_{\varepsilon^{(is_i)}}^{x^{(p+i)}}du_h
+\frac12\sum_{i=1}^{i=q}\int_{\varepsilon^{(-is_i)}}^{\varepsilon^{(is_i)}}du_h
-\int_{\mu}^{\xi}du_h,
\]
wobei die Zahlen \(s_1,s_2\ldots s_q\) wie oben die Bedeutung \(\pm1\) haben.

Nun können wir die Integrale \(u_{hi}\), welche oben nur bis auf Periodicitätsmoduln bestimmt waren, offenbar so wählen, dass sie mit denjenigen Integralen übereinstimmen, welche in der letzten Form auftreten, bei demselben Doppelpunct beginnend und endigend, und welche durch den oben ausgeführten Curvenübergang genau bestimmt sind. Dann wird also, wenn \(\varepsilon^{(i)}\) dem ersten, \(\varepsilon^{(-i)}\) dem zweiten Parameter des Integrals \(u_{p+i}\) entsprechend gedacht wird,
\[
\int_{\varepsilon^{(-is_i)}}^{\varepsilon^{(is_i)}}du_h=s_i\cdot u_{hi},
\]
und man kann die obige Gleichung für \(\omega_h\) in die Form bringen:
\[
\omega_h-\frac12(s_1u_{h1}+s_2u_{h2}+\cdots)
=\sum_{i=1}^{i=p}\int_{\alpha^{(i)}}^{x^{(i)}}du_h
+\sum_{i=1}^{i=q}\int_{\varepsilon^{(is_i)}}^{x^{(p+i)}}du_h-\int_{\mu}^{\xi}du_h.
\]
Man hat hier die Argumente der verschiedenen in \(V^{(q)}\) auftretenden \(\Theta\)functionen vor sich. Sie haben die Form der \(\omega_h\) selbst; nur sind die untern Grenzen hier durch die im ursprünglichen Umkehrproblem auftretenden Puncte \(\alpha\) ersetzt, und durch die bevorzugten Doppelpuncte selbst, bei welchen man immer die passende Fortschreitungsrichtung zu wählen hat.

Die Function \(V^{(q)}\) verschwindet nach \(\S\,78\), wenn einer der Punkte \(x\) mit \(\xi\) zusammenfällt. Lässt man \(x^{(p+i)}\) mit \(\xi\) zusammenfallen, und \(x^{(p+1)},x^{(p+2)}\ldots x^{(p+i-1)},x^{(p+i+1)}\ldots x^{(p+q)}\) mit \(\varepsilon^{(1\cdot s_1)},\varepsilon^{(2\cdot s_2)}\ldots\varepsilon^{((i-1)s_{i-1})},\varepsilon^{((i+1)s_{i+1})}\ldots\varepsilon^{(qs_q)}\), so dass also diese letzten \(x\) sich auf vorbestimmten Zweigen \(q-1\) der bevorzugten Doppelpuncte nähern, so verschwindet \(V^{(q)}\), während zugleich die Grössen \(\omega_{p+1},\omega_{p+2}\ldots\omega_{p+i-1},\omega_{p+i+1},\ldots\omega_{p+q}\) unendlich grosse Werthe erhalten, und zwar wird, da \(\varepsilon^{(k)}\) den ersten, \(\varepsilon^{(-k)}\) den zweiten Parameter darstellt, \(e^{s_k\omega_k}\) unendlich gross, sobald einer der Puncte \(x\) sich dem Puncte \(\varepsilon^{(ks_k)}\) nähert. Beschränkt man also die Function \(V^{(q)}\) auf diejenigen Glieder, gegen welche alle übrigen unendlich klein werden, so erhält man die Gleichung:
\[
\begin{aligned}
(2)\quad 0={}&A_{s_1s_2\ldots s_{i-1},1,s_{i+1}\ldots s_q}
 e^{\frac12\omega_{p+i}}
 \Theta\left(\sum_{k=1}^{k=p}\int_{\alpha^{(k)}}^{x^{(k)}}du_h+
 \int_{\varepsilon^{(i)}}^{\mu}du_h\right)\\
&+A_{s_1s_2\ldots s_{i-1},-1,s_{i+1}\ldots s_q}
 e^{-\frac12\omega_{p+i}}
 \Theta\left(\sum_{k=1}^{k=p}\int_{\alpha^{(k)}}^{x^{(k)}}du_h+
 \int_{\varepsilon^{(-i)}}^{\mu}du_h\right),
\end{aligned}
\]
welche für alle beliebigen Puncte \(x^{(1)},x^{(2)}\ldots x^{(p)}\) bestehen muss, sobald man zugleich der Grösse \(\omega_{p+i}\) ihren Werth giebt:
\[
\omega_{p+i}=\sum_{k=1}^{k=p}\int_{\beta^{(k)}}^{x^{(k)}}du_{p+i}
+\sum_{k=1}^{k=i-1}\int_{\beta^{(p+k)}}^{\varepsilon^{(ks_k)}}du_{p+i}
+\int_{\beta^{(p+i)}}^{\mu}du_{p+i}
+\sum_{k=i+1}^{k=q}\int_{\beta^{(p+k)}}^{\varepsilon^{(ks_k)}}du_{p+i}.
\]
Nun können wir insbesondere den Puncten \(x\) solche Lagen \(\gamma\) geben, dass
\[
(3)\quad
2\sum_{k=1}^{k=p}\int_{\alpha^{(k)}}^{\gamma^{(k)}}du_h+
\int_{\varepsilon^{(i)}}^{\mu}du_h+
\int_{\varepsilon^{(-i)}}^{\mu}du_h=0,
\qquad (h=1,2\ldots p).
\]
In diesem Falle werden die beiden \(\Theta\)functionen in der Gleichung (2) einander gleich, und man erhält:
\[
(4)\quad
A_{s_1s_2\ldots s_{i-1},-1,s_{i+1}\ldots s_q}
=-e^{\omega_{p+i}}A_{s_1s_2\ldots s_{i-1},1,s_{i+1}\ldots s_q},
\]
wo nun
\[
\omega_{p+i}=\sum_{k=1}^{k=p}\int_{\beta^{(k)}}^{\gamma^{(k)}}du_{p+i}
+\sum_{k=1}^{k=i-1}\int_{\beta^{(p+k)}}^{\varepsilon^{(ks_k)}}du_{p+i}
+\int_{\beta^{(p+i)}}^{\mu}du_{p+i}
+\sum_{k=i+1}^{k=q}\int_{\beta^{(p+k)}}^{\varepsilon^{(ks_k)}}du_{p+i}.
\tag{5}
\]
Die Gleichung (3) definirt die \(\gamma\) als die Berührungspuncte einer durch alle Doppel- und Rückkehrpuncte gehenden Curve \((n-2)^{\mathrm{ter}}\) Ordnung. Denken wir uns die Tangente in \(\mu\) und einen beliebigen Strahl durch den Doppelpunct \(\varepsilon^{(i)},\varepsilon^{(-i)}\) gezogen. Wir können dann die Integralsumme
\[
\int_{\varepsilon^{(i)}}^{\mu}du_h+
\int_{\varepsilon^{(-i)}}^{\mu}du_h
\]
durch die negative Summe der \(n-2\) Integrale ersetzen, deren obere Grenzen die Schnittpuncte der Tangente und deren untere die Schnittpuncte des Strahls sind. Trägt man diese Integrale in die Gleichung (3) ein, so sagt sie nach dem Abel'schen Theorem aus, dass, wie die \(\alpha\) Berührungspuncte einer durch alle Doppel- und Rückkehrpuncte und die \(n-2\) Schnittpuncte der Tangente gelegten Curve \((n-2)^{\mathrm{ter}}\) Ordnung sind, ebenso die \(\gamma\) die Berührungspuncte einer solchen Curve werden, welche durch alle Doppel- und Rückkehrpuncte und die \(n-2\) Schnittpuncte des Strahls gehen; aber letztere wird von dem Strahl demnach in \(n+2\) Puncten geschnitten, und muss also diesen Strahl ganz enthalten. Demnach sind die \(\gamma\) die Berührungspuncte einer Curve \((n-3)^{\mathrm{ter}}\) Ordnung, welche durch alle Doppel- und Rückkehrpuncte, einen ausgenommen, hindurchgeht, und \(f=0\) in \(p\) Puncten berührt.

Da diese Puncte \(\gamma\) dem Doppelpuncte \(\varepsilon^{(i)},\varepsilon^{(-i)}\) zugeordnet sind, so wollen wir sie durch
\[
\gamma^{i,1},\gamma^{i,2},\ldots,\gamma^{i,p}
\]
bezeichnen, und ferner setzen:
\[
\begin{aligned}
a_i={}&\sum_{k=1}^{k=p}\int_{\beta^{(k)}}^{\gamma^{i,k}}du_{p+i}
+\int_{\beta^{(p+i)}}^{\mu}du_{p+i}
+\frac12\sum_{k=1}^{k=i-1}\left(\int_{\beta^{(p+k)}}^{\varepsilon^{(k)}}du_{p+i}+\int_{\beta^{(p+k)}}^{\varepsilon^{(-k)}}du_{p+i}\right)\\
&+\frac12\sum_{k=i+1}^{k=q}\left(\int_{\beta^{(p+k)}}^{\varepsilon^{(k)}}du_{p+i}+\int_{\beta^{(p+k)}}^{\varepsilon^{(-k)}}du_{p+i}\right),\\
b_{ik}=b_{ki}={}&\int_{\varepsilon^{(-i)}}^{\varepsilon^{(i)}}du_{p+k}=\int_{\varepsilon^{(-k)}}^{\varepsilon^{(k)}}du_{p+i}.
\end{aligned}
\]
Die verschiedenen aus (4) abzuleitenden Gleichungen geben dann offenbar
\[
A_{s_1s_2\ldots s_q}=(-1)^{\frac{q-s_1-s_2-\cdots-s_q}{2}}
 e^{-\frac12\sum_i s_ia_i-\frac14\sum_{i\ge k}\sum s_is_kb_{ik}}\cdot C,
\]
und man erhält demnach
\[
V^{(q)}=e^{\frac12\sum\varkappa_h\omega_h}\cdot C\cdot \Theta^{(q)}(\omega_1,\omega_2\ldots\omega_{p+q}),
\]
wenn man unter \(\Theta^{(q)}\) die Function versteht:
\[
\begin{aligned}
(6)\quad \Theta^{(q)}(\omega)={}&
\sum_{s_1s_2\ldots s_q=\pm1}
(-1)^{\frac{-s_1-s_2-\cdots-s_q}{2}}
 e^{\frac12\sum_i s_i(\omega_{p+i}-a_i)-\frac14\sum\sum s_is_kb_{ik}}\\
&\cdot\Theta\left(\omega_h-\frac{s_1u_{h1}+s_2u_{h2}+\cdots}{2}\right).
\end{aligned}
\]
Die in den Exponentialgrössen auftretenden Argumente haben hier eine ganz ähnliche Form, wie die Argumente der \(\Theta\)function; denn es ist
\[
\begin{aligned}
\omega_{p+i}-a_i-\sum_{k\ge i}s_kb_{ik}
={}&\sum_{k=1}^{k=p}\int_{\gamma^{i,k}}^{x^{(k)}}du_{p+i}
+\sum_{k=1}^{k=i-1}\int_{\varepsilon^{(ks_k)}}^{x^{(p+k)}}du_{p+i}
+\int_{\xi}^{x^{(p+i)}}du_{p+i}\\
&+\sum_{k=i+1}^{k=q}\int_{\varepsilon^{(ks_k)}}^{x^{(p+k)}}du_{p+i}.
\end{aligned}
\]
Die Constante \(C\) bestimmt sich dadurch, dass \(U^{(q)}=-\log V^{(q)}\) verschwinden muss für \(\xi=\zeta\), \(x^{(i)}=\gamma^{(i)}\). Setzt man demnach \(\omega_1^\circ,\omega_2^\circ\ldots\omega_{p+q}^\circ\) für die Werthe, welche die \(w\) annehmen, wenn \(\xi\) in \(\zeta\) und die \(x\) in die \(\gamma\) fallen, so hat man
\[
V^{(q)}=e^{\frac12\sum\varkappa_h(\omega_h-\omega_h^\circ)}\cdot
\frac{\Theta^{(q)}(\omega)}{\Theta^{(q)}(\omega^\circ)}.
\]
Mit Hülfe der Transcendente \(\Theta^{(q)}\) erfolgt die Lösung des erweiterten Umkehrproblems wie die des Jacobi'schen. Man hat, wenn man wieder immer nur ein Argument der Function \(\Theta^{(q)}\) schreibt:
\[
\begin{aligned}
\frac12T_{\xi\eta}^{(q)}(x^{(1)},x^{(2)}\ldots x^{(p+q)})
&=\log\frac{V^{(q)}(x^{(1)}x^{(2)}\ldots x^{(p+q)};\eta)}{V^{(q)}(x^{(1)}x^{(2)}\ldots x^{(p+q)};\xi)}\\
&=\frac12\sum\varkappa_h\int_{\eta}^{\xi}du_h+
\log\frac{\Theta^{(q)}\left(\sum_{i=1}^{i=p+q}\int_{\beta^{(i)}}^{x^{(i)}}du_h-\int_{\mu}^{\eta}du_h\right)}{\Theta^{(q)}\left(\sum_{i=1}^{i=p+q}\int_{\beta^{(i)}}^{x^{(i)}}du_h-\int_{\mu}^{\xi}du_h\right)}.
\end{aligned}
\]
Ferner
\[
\begin{aligned}
&T_{\xi\eta}^{(q)}
\begin{pmatrix}
x^{(1)},x^{(2)}\ldots x^{(p+q)}\\
z^{(1)},z^{(2)}\ldots z^{(p+q)}
\end{pmatrix}\\
&=\log\frac{
\Theta^{(q)}\left(\sum_{i=1}^{i=p+q}\int_{\beta^{(i)}}^{x^{(i)}}du_h-\int_{\mu}^{\eta}du_h\right)
\Theta^{(q)}\left(\sum_{i=1}^{i=p+q}\int_{\beta^{(i)}}^{z^{(i)}}du_h-\int_{\mu}^{\xi}du_h\right)
}{
\Theta^{(q)}\left(\sum_{i=1}^{i=p+q}\int_{\beta^{(i)}}^{z^{(i)}}du_h-\int_{\mu}^{\eta}du_h\right)
\Theta^{(q)}\left(\sum_{i=1}^{i=p+q}\int_{\beta^{(i)}}^{x^{(i)}}du_h-\int_{\mu}^{\xi}du_h\right)
}.
\end{aligned}
\]
Sind \(\varphi=0\), \(\psi=0\) die Gleichungen zweier Curven gleich hoher Ordnung, und bedeuten die \(\xi\) die Schnittpuncte von \(f=0\) mit \(\varphi-\lambda\psi=0\), die \(\eta\) aber die Schnittpuncte von \(f=0\) mit \(\psi=0\), so findet man die symmetrischen Functionen der Werthe, welche \(\frac{\varphi}{\psi}\) für die Puncte \(x^{(1)},x^{(2)}\ldots x^{(p+q)}\) annimmt, mit Hülfe der Formel:
\[
\prod_{i=1}^{i=p+q}
\frac{\left(\dfrac{\varphi}{\psi}\right)_{x^{(i)}}-\lambda}{\left(\dfrac{\varphi}{\psi}\right)_{\beta^{(i)}}-\lambda}
=\prod_{\xi\eta}
\frac{
\Theta^{(q)}\left(v_h-\int_{\mu}^{\xi}du_h\right)
\Theta^{(q)}\left(\int_{\mu}^{\eta}du_h\right)
}{
\Theta^{(q)}\left(v_h-\int_{\mu}^{\eta}du_h\right)
\Theta^{(q)}\left(\int_{\mu}^{\xi}du_h\right)
},
\]
wobei die \(v\) die Integralsummen bedeuten:
\[
v_h=\sum_{i=1}^{i=p+q}\int_{\beta^{(i)}}^{x^{(i)}}du_h,
\qquad (h=1,2\ldots p+q).
\]

% --- Source pages 309--321; printed pages 287--299. ---

\subsection*{\S\,81. Die Function \(\Theta^{(q)}\) ist bei passender Anordnung eine Function ihrer Argumente, welche mit \(q\) gerade oder ungerade ist.}

Wir wollen jetzt untersuchen, was aus der Function \(\Theta^{(q)}(\omega)\) wird, wenn ihre Argumente sämmtlich das Zeichen wechseln. Setzt man in der Formel für \(\Theta^{(q)}\) die \(-\omega\) an Stelle der \(\omega\), und zugleich \(-s_1,-s_2\ldots\) an Stelle von \(s_1,s_2\ldots\) (welches letzteres durchaus nichts ändert), so hat man:
\[
\begin{aligned}
\Theta^{(q)}(-\omega)={}&
\sum_{s_1,s_2\ldots s_q=\pm1}(-1)^{\frac{s_1+s_2+\cdots+s_q}{2}}
\exp\!\left
\{\frac12\sum_i s_i(\omega_{p+i}+a_i)-\frac14\sum_{i>k}\sum s_i s_k b_{ik}\right\}\\
&\cdot \Theta\!\left(\omega_h-\frac{s_1u_{h1}+s_2u_{h2}+\cdots}{2}\right).
\end{aligned}
\]
Die entsprechenden Glieder von \(\Theta^{(q)}(\omega)\) und \(\Theta^{(q)}(-\omega)\) unterscheiden sich also nur durch die Coefficienten
\[
(-1)^{-\frac{s_1+s_2+\cdots}{2}}e^{-\frac12\sum_i s_ia_i},\qquad
(-1)^{\frac{s_1+s_2+\cdots}{2}}e^{\frac12\sum_i s_ia_i},
\]
deren Quotient den Werth hat:
\[
e^{-\sum_i s_i a_i}\cdot(-1)^q.
\]
Betrachten wir den im Exponenten von \(e\) stehenden Theil genauer. Derselbe wird mit Einsetzung der Werthe der \(a\):
\[
\begin{aligned}
-\frac12\sum_i s_i\Bigg[
&2\sum_{k=1}^{k=p}\int_{\beta^{(k)}}^{\gamma^{i,k}}du_{p+i}
+\sum_{k=1}^{k=i-1}\left(\int_{\beta^{(p+k)}}^{\varepsilon^{(k)}}du_{p+i}
+\int_{\beta^{(p+k)}}^{\varepsilon^{(-k)}}du_{p+i}\right)\\
&+2\int_{\beta^{(p+i)}}^{\mu}du_{p+i}
+\sum_{k=i+1}^{k=q}\left(\int_{\beta^{(p+k)}}^{\varepsilon^{(k)}}du_{p+i}
+\int_{\beta^{(p+k)}}^{\varepsilon^{(-k)}}du_{p+i}\right)
\Bigg].
\end{aligned}
\]
Ergänzen wir hier in dem eingeklammerten Theile Integrale mit zusammenfallenden obern und untern Grenzen, welche die Puncte darstellen, in denen \(f=0\) von der Tangente in \(\mu\) noch geschnitten wird, so bilden die untern Grenzen das vollständige Schnittpunctsystem der Curve \((n-2)^{\mathrm{ter}}\) Ordnung, durch welche die \(\beta\) definirt wurden, die obern aber das Aggregat der Schnittpuncte der Tangente in \(\mu\) mit denen der Curve \((n-3)^{\mathrm{ter}}\) Ordnung, welche die \(\gamma\) definirt. Nach dem Abel'schen Theorem ist also der eingeklammerte Ausdruck gleich einem Periodicitätsmodul des Integrals \(u_{p+i}\), vermehrt um den Logarithmus einer algebraischen Function. Der letztere verschwindet, weil die Coordinaten der Parameterpuncte von \(u_{p+i}\) coincidiren, und doch beide Curven nicht durch den betreffenden Doppelpunct gehen, also die linken Theile ihrer Gleichungen für \(\varepsilon^{(i)}\) und \(\varepsilon^{(-i)}\) gleiche und nicht verschwindende Werthe annehmen. Daher bleibt nur der Periodicitätsmodul zu bestimmen. Dieser besteht aus einem Vielfachen von \(2\pi\sqrt{-1}\), welcher dem Integrale dritter Gattung \(u_{p+i}\) als solchem zukommt, und aus einem weitern Theile, der entsprechende Glieder auch bei den Integralen erster Gattung erfordert, und daher mit diesen Gliedern verschwindet. Diesen zweiten Theil kann man daher bei den Integralen erster Gattung untersuchen, also bei den Integralen:
\[
\begin{aligned}
&2\sum_{k=1}^{k=p}\int_{\beta^{(k)}}^{\gamma^{i,k}}du_h
+\sum_{k=1}^{k=i-1}\left(\int_{\beta^{(p+k)}}^{\varepsilon^{(k)}}du_h
+\int_{\beta^{(p+k)}}^{\varepsilon^{(-k)}}du_h\right)
+2\int_{\beta^{(p+i)}}^{\mu}du_h\\
&\hspace{3em}
+\sum_{k=i+1}^{k=q}\left(\int_{\beta^{(p+k)}}^{\varepsilon^{(k)}}du_h
+\int_{\beta^{(p+k)}}^{\varepsilon^{(-k)}}du_h\right),
\qquad (h=1,2\ldots p).
\end{aligned}
\]
Vergleicht man nun die Gleichungen \(\S\,80\) (1) und (3), welche sich auf Integrale erster Gattung und auf dieselben Integrationswege beziehen, so findet man sofort, dass diese Summe verschwindet. Daraus folgt also, dass auch die Periodicitätsmoduln der Integrale dritter Gattung nur aus einem Terme \(2\varrho_i\pi\sqrt{-1}\) bestehen können, und dass also der zu untersuchende Quotient den Werth hat:
\[
(-1)^{q+\sum\varrho_i s_i},
\]
oder da alle \(s_i\) gleich \(\pm1\) sind:
\[
(-1)^{q+\sum\varrho_i},
\]
so dass endlich
\[
\Theta^{(q)}(-\omega)=(-1)^{q+\sum\varrho_i}\cdot\Theta^{(q)}(\omega).
\]
Was nun die Zahlen \(\varrho\) betrifft, so liegt in diesen eine wirkliche Willkürlichkeit. Denken wir uns nämlich das System der \(\beta\) in der Weise bestimmt, wie dies auf p. 281 angegeben wurde. Dann kann man andre Systeme \(\beta'\), andere Berührungscurven finden, welche sich zu den ursprünglichen so verhalten, dass
\[
\sum_{k=1}^{k=p+q}\int_{\beta^{(k)}}^{\beta'^{(k)}}du_h=0,
\qquad (h=1,2\ldots p),
\]
\[
\sum_{k=1}^{k=p+q}\int_{\beta^{(k)}}^{\beta'^{(k)}}du_{p+h}=\sigma_h\pi\sqrt{-1},
\qquad (h=1,2\ldots q),
\]
und welche sich nach dem Abel'schen Theorem als Berührungspuncte solcher Curven erweisen, wie man sofort durch Multiplication der vorstehenden Gleichungen mit 2 sieht. Führt man nun die \(\beta'\) statt der \(\beta\) ein, so ändern sich die \(\omega_1,\omega_2\ldots\omega_p\), also die Argumente der in \(\Theta^{(q)}\) auftretenden \(\Theta\)functionen nicht, wohl aber ändert sich jede der Grössen \(\omega_{p+i},a_{i,s_1,s_2\ldots}\) um \(\sigma_i\pi\sqrt{-1}\). Man sieht, dass hiedurch die Function \(\Theta^{(q)}\) nur insofern geändert wird, als die Variabeln \(\omega_{p+i}\) eine veränderte Bedeutung erhalten. Aber zugleich ändert sich der oben untersuchte Quotient um
\[
-\sum s_i\sigma_i,
\]
oder um \(\sum\sigma_i\). Wir wollen nun im Folgenden immer statt der \(\beta\) dasjenige System der \(\beta'\) gesetzt denken, für welches \(\sigma_i=\varrho_i\) wird \((i=1,2\ldots q)\). Dann hat man
\[
\Theta^{(q)}(-\omega)=(-1)^q\Theta^{(q)}(\omega),
\]
und die so definirte Function \(\Theta^{(q)}\) ist also eine gerade oder ungerade Function ihrer Argumente, jenachdem \(q\) gerade oder ungerade ist.

Für diese so definirte Function ist immer
\[
a_i=0,
\qquad (i=1,2\ldots q)
\]
und man hat also endlich für \(\Theta^{(q)}\) den Ausdruck:
\[
\begin{aligned}
\Theta^{(q)}(\omega)={}&
\sum_{s_1s_2\ldots s_q=\pm1}
\exp\!\left\{\frac12\sum_i s_i\omega_{p+i}
-\frac14\sum_{i>k}\sum s_i s_k b_{ik}
-\frac{\pi\sqrt{-1}}{2}\sum_i s_i\right\}\\
&\cdot\Theta\!\left(\omega_h-\frac{s_1u_{h1}+s_2u_{h2}+\cdots}{2}\right),
\end{aligned}
\]
während
\[
\omega_h=\sum_{i=1}^{i=p+q}\int_{c^{(i)}}^{x^{(i)}}du_h-\int_{\mu}^{\xi}du_h,
\qquad (h=1,2\ldots p+q).
\]

\subsection*{\S\,82. Die Theilung für das erweiterte Umkehrproblem.}

Wir wollen auf die Erweiterungen derjenigen Sätze nicht eingehen, welche den Inhalt des neunten Abschnitts ausmachen, und welche leicht auf die Aufgaben des erweiterten Umkehrproblems und auf die Function \(\Theta^{(q)}\) ausgedehnt werden. Dagegen scheint es angemessen, der \emph{Theilung} einige Worte zu widmen, und zwar der speciellen, da die Erweiterung der allgemeinen Theilungsaufgabe ganz wie die im neunten Abschnitt behandelte allgemeine Theilung gelöst wird. Die Aufgabe der allgemeinen Theilung wird hier sein, aus den Gleichungen
\[
m\sum_{i=1}^{i=p+q}\int_{c^{(i)}}^{x^{(i)}}du_h
+\sum_{i=1}^{i=r}\int_{a^{(i)}}^{y^{(i)}}du_h=0,
\qquad (h=1,2\ldots p+q)
\]
die \(x\) zu bestimmen, oder durch gegebene Puncte \(y\) eine Curve zu legen, welche durch die \(\delta+r\) abgesonderten Ausnahmspuncte geht, welche in \(p+q\) Puncten \(m\)punctig berührt, und welche ihre übrigen Schnittpuncte mit \(f=0\) mit einer ähnlichen Berührungscurve gemein hat, die in den Puncten \(c\) berührt und durch die Puncte \(a\) geht. Die Aufgabe lässt \(m^{2p+q}\) Lösungen zu; der Exponent von \(m\) giebt die Anzahl der von einander unabhängigen Systeme von Periodicitätsmoduln des erweiterten Umkehrproblems an. Man findet die Wurzeln der jenen Lösungen entsprechenden algebraischen Gleichung wie in \(\S\,70\) auf algebraische Weise unter Voraussetzung der Wurzeln einer speciellen Theilungsgleichung.

Die letztere, welche aus der Gleichung
\[
(1)\quad m\sum_{i=1}^{i=p+q}\int_{c^{(i)}}^{x^{(i)}}du_h=P_h,
\qquad (h=1,2\ldots p+q),
\]
entspringt, hat ausser der selbstverständlichen Lösung \(x^{(1)}=c^{(1)}\), \(x^{(2)}=c^{(2)}\) etc., deren noch \(m^{2p+q}-1\). Wir werden die entsprechende algebraische Gleichung, also die specielle Theilungsgleichung des erweiterten Umkehrproblems, auf die entsprechende des Jacobischen zurückführen.

Sei ein Problem der speciellen Theilung für das Jacobi'sche Umkehrproblem durch die Gleichung gegeben:
\[
(2)\quad m\sum_{i=1}^{i=p}\int_{\gamma^{(i)}}^{z^{(i)}}du_h=P_h,
\qquad (h=1,2\ldots p).
\]
Vergleichen wir nun zwei Lösungen der Probleme (1) und (2) mit einander, für welche beziehungsweise:
\[
(3)\quad
\left\{\begin{aligned}
\sum_{i=1}^{i=p+q}\int_{c^{(i)}}^{x^{(i)}}du_h
&=\frac{2\lambda_h\pi\sqrt{-1}}{m}
+\frac1m(\varkappa_1a_{1h}+\varkappa_2a_{2h}+\cdots),
&& (h=1,2\ldots p),\\
\sum_{i=1}^{i=p+q}\int_{c^{(i)}}^{x^{(i)}}du_{p+h}
&=\frac{2\nu_h\pi\sqrt{-1}}{m}
+\frac1m(\varkappa_1u_{1h}+\varkappa_2u_{2h}+\cdots),
&& (h=1,2\ldots q),
\end{aligned}\right.
\]
und
\[
(4)\quad
\sum_{i=1}^{i=p}\int_{\gamma^{(i)}}^{z^{(i)}}du_h
=\frac{2\lambda_h\pi\sqrt{-1}}{m}
+\frac1m(\varkappa_1a_{1h}+\varkappa_2a_{2h}+\cdots),
\qquad (h=1,2\ldots p).
\]
Aus der Combination der ersten Gleichung (3) mit (4) ergiebt sich:
\[
(5)\quad
\sum_{i=1}^{i=p+q}\int_{c^{(i)}}^{x^{(i)}}du_h
+\sum_{i=1}^{i=p}\int_{z^{(i)}}^{\gamma^{(i)}}du_h=0,
\qquad (h=1,2\ldots p).
\]
Setzt man also die Puncte \(z\) als bekannt voraus, oder doch die symmetrischen Functionen ihrer Coordinaten, d. h. eine Wurzel der speciellen Theilungsgleichung für das Jacobi'sche Umkehrproblem, so kann man eine Curve \(\psi=0\) durch sie, durch die \(c\) und durch sämmtliche Doppel- und Rückkehrpuncte legen; diese schneidet in weitern Puncten \(b\); durch sie so wie durch die \(\gamma\) und die Doppel- und Rückkehrpuncte geht dann eine zweite Curve \(\varphi=0\) gleich hoher Ordnung wie die vorige, deren übrige Schnittpuncte mit \(f=0\) die \(x\) sind. Aber diese Curve, oder selbst ihr Schnittpunctsystem, sind hiedurch noch nicht völlig bestimmt; vielmehr bleiben zunächst noch \(q\) der \(x\) oder ebensoviel durch \(f=0\) nicht zerstörbare Coefficienten von \(\varphi=0\) unbestimmt.

Denken wir uns nun zunächst die Parameterpuncte der Integrale dritter Gattung \(u_{p+h}\) von den \(q\) bevorzugten Doppelpuncten wenig verschieden, und in der Nähe jedes der Doppelpuncte auf den verschiedenen Zweigen von \(f=0\) in \(\varepsilon^{(h)}\), \(\varepsilon^{(-h)}\) gelegen. Alsdann hat man wegen der Gleichung (5) nach dem Abel'schen Satz auch:
\[
(6)\quad
\sum_{i=1}^{i=p+q}\int_{c^{(i)}}^{x^{(i)}}du_{p+h}
+\sum_{i=1}^{i=p}\int_{z^{(i)}}^{\gamma^{(i)}}du_{p+h}
=-\log
\frac{\dfrac{\varphi(\varepsilon^{(h)})}{\psi(\varepsilon^{(h)})}}
{\dfrac{\varphi(\varepsilon^{(-h)})}{\psi(\varepsilon^{(-h)})}}.
\]
Lässt man nun die Puncte \(\varepsilon^{(h)}\) und \(\varepsilon^{(-h)}\) sich dem betreffenden Doppelpunct mehr und mehr nähern, so verschwinden zugleich \(\varphi(\varepsilon^{(h)})\), \(\psi(\varepsilon^{(h)})\), \(\varphi(\varepsilon^{(-h)})\), \(\psi(\varepsilon^{(-h)})\). Ist also
\[
\begin{array}{lll}
\varepsilon_1^{(h)}=\vartheta_1^{(h)}+\delta\xi_1^{(h)},&
\varepsilon_2^{(h)}=\vartheta_2^{(h)}+\delta\xi_2^{(h)},&
\varepsilon_3^{(h)}=\vartheta_3^{(h)}+\delta\xi_3^{(h)},\\
\varepsilon_1^{(-h)}=\vartheta_1^{(h)}+\delta\eta_1^{(h)},&
\varepsilon_2^{(-h)}=\vartheta_2^{(h)}+\delta\eta_2^{(h)},&
\varepsilon_3^{(-h)}=\vartheta_3^{(h)}+\delta\eta_3^{(h)},
\end{array}
\]
wo die \(\delta\xi\), \(\delta\eta\) die unendlich kleinen Incremente der Coordinaten eines Puncts bedeuten, der sich von dem Doppelpunct aus auf dem einen oder andern Zweige von \(f=0\) bewegt, so hat man
\[
\lim
\frac{\dfrac{\varphi(\varepsilon^{(h)})}{\psi(\varepsilon^{(h)})}}
{\dfrac{\varphi(\varepsilon^{(-h)})}{\psi(\varepsilon^{(-h)})}}
=
\frac{\dfrac{\varphi'(\vartheta_1^{(h)})\delta\xi_1^{(h)}+\varphi'(\vartheta_2^{(h)})\delta\xi_2^{(h)}+\varphi'(\vartheta_3^{(h)})\delta\xi_3^{(h)}}
{\psi'(\vartheta_1^{(h)})\delta\xi_1^{(h)}+\psi'(\vartheta_2^{(h)})\delta\xi_2^{(h)}+\psi'(\vartheta_3^{(h)})\delta\xi_3^{(h)}}}
{\dfrac{\varphi'(\vartheta_1^{(h)})\delta\eta_1^{(h)}+\varphi'(\vartheta_2^{(h)})\delta\eta_2^{(h)}+\varphi'(\vartheta_3^{(h)})\delta\eta_3^{(h)}}
{\psi'(\vartheta_1^{(h)})\delta\eta_1^{(h)}+\psi'(\vartheta_2^{(h)})\delta\eta_2^{(h)}+\psi'(\vartheta_3^{(h)})\delta\eta_3^{(h)}}}.
\]
Von den drei \(d\xi\) so wie von den drei \(d\eta\) kann man eines nach Belieben gleich Null setzen; das Verhältniss der andern findet man dann aus der quadratischen Gleichung, welche die niedrigsten Terme von \(f=0\) in der Nähe des Doppelpuncts liefern:
\[
\sum_i\sum_k
\frac{\partial^2f(\vartheta^{(h)})}{\partial\vartheta_i^{(h)}\partial\vartheta_k^{(h)}}
\,\delta\xi_i^{(h)}\delta\xi_k^{(h)}=0,
\]
und welche auch befriedigt werden muss, wenn man die \(d\xi\) durch die \(d\eta\) ersetzt.

Bemerken wir ferner, dass wir die Gleichung (4) als eine Form des Abel'schen Theorems ansehen können, welche aus der Gleichung
\[
m\sum_{i=1}^{i=p}\int_{\gamma^{(i)}}^{z^{(i)}}du_h
=\lambda\int_t^{\delta}du_h+\varkappa_1\int_{t^{(1)}}^{\delta^{(1)}}du_h
+\varkappa_2\int_{t^{(2)}}^{\delta^{(2)}}du_h+\cdots+\varkappa_p\int_{t^{(p)}}^{\delta^{(p)}}du_h,
\qquad (h=1,2\ldots p)
\]
entspringt, indem man die variabeln Puncte \(t\) Periodenwege durchlaufen lässt, bis sie mit den Puncten \(\delta\) zusammenfallen, und dabei die auf der rechten Seite von (5) befindlichen Periodicitätsmoduln erzeugen. Legt man also eine Curve \(\Phi=0\), welche durch Doppel- und Rückkehrpuncte geht, in den Puncten \(\gamma\) \(m\)punctig, in den Puncten \(\delta,\delta^{(1)},\delta^{(2)}\ldots\) aber beziehungsweise \(\lambda\)punctig, \(\varkappa_1\)punctig etc. berührt; und legt man durch ihre übrigen Schnittpuncte mit \(f=0\), so wie durch die Doppel- und Rückkehrpuncte eine andre Curve \(\Psi=0\) gleich hoher Ordnung, welche in den Puncten \(t,t^{(1)},t^{(2)}\ldots\) beziehungsweise \(\lambda\)punctig, \(\varkappa_1\)punctig etc. berührt; denkt man sich ferner die zweite Curve aus der ersten durch continuirliche Bewegung so entstanden, dass sie in keinem Augenblicke ihre Berührungseigenschaften verliert, so wird die zweite Curve in den Puncten \(z\) berühren, und diese werden die betreffende Lösung der speciellen Theilungsaufgabe des Jacobi'schen Umkehrproblems bilden, wenn die \(t\) auf den angezeigten Periodenwegen zu den \(\delta\) zurückgekehrt sind. Daher hat man auch für irgend ein Integral dritter Gattung zunächst:
\[
m\sum_{i=1}^{i=p}\int_{\gamma^{(i)}}^{z^{(i)}}d\Pi_{\xi\eta}
+\lambda\int_{\delta}^{t}d\Pi_{\xi\eta}
+\varkappa_1\int_{\delta^{(1)}}^{t^{(1)}}d\Pi_{\xi\eta}+\cdots
=-\log\frac{\dfrac{\Phi(\xi)}{\Psi(\xi)}}{\dfrac{\Phi(\eta)}{\Psi(\eta)}}.
\]
Wenn aber die \(t\) ihre Periodenwege durchlaufen haben, und wenn zugleich das Integral \(\Pi_{\xi\eta}\) in das Integral \(u_{p+h}\) übergeht, so geht diese Gleichung ähnlich wie (5), in die Form über:
\[
\begin{aligned}
\sum_{i=1}^{i=p}\int_{\gamma^{(i)}}^{z^{(i)}}du_{p+h}
={}&\frac1m(\varkappa_1u_{1h}+\varkappa_2u_{2h}+\cdots)\\
&-\frac1m\log
\frac{\dfrac{\Phi'(\vartheta_1^{(h)})\delta\xi_1^{(h)}+\Phi'(\vartheta_2^{(h)})\delta\xi_2^{(h)}+\cdots}
{\Psi'(\vartheta_1^{(h)})\delta\xi_1^{(h)}+\Psi'(\vartheta_2^{(h)})\delta\xi_2^{(h)}+\cdots}}
{\dfrac{\Phi'(\vartheta_1^{(h)})\delta\eta_1^{(h)}+\Phi'(\vartheta_2^{(h)})\delta\eta_2^{(h)}+\cdots}
{\Psi'(\vartheta_1^{(h)})\delta\eta_1^{(h)}+\Psi'(\vartheta_2^{(h)})\delta\eta_2^{(h)}+\cdots}},\\
&\hspace{13em}(h=1,2\ldots q).
\end{aligned}
\]
Die Combination dieser Gleichung mit der zweiten Gleichung (3) und mit (6) liefert nun sofort, dass der Ausdruck
\[
\log\frac{\dfrac{\varphi'(\vartheta_1^{(h)})\delta\xi_1^{(h)}+\cdots}
{\psi'(\vartheta_1^{(h)})\delta\xi_1^{(h)}+\cdots}}
{\dfrac{\varphi'(\vartheta_1^{(h)})\delta\eta_1^{(h)}+\cdots}
{\psi'(\vartheta_1^{(h)})\delta\eta_1^{(h)}+\cdots}}
-\frac1m
\log\frac{\dfrac{\Phi'(\vartheta_1^{(h)})\delta\xi_1^{(h)}+\cdots}
{\Psi'(\vartheta_1^{(h)})\delta\xi_1^{(h)}+\cdots}}
{\dfrac{\Phi'(\vartheta_1^{(h)})\delta\eta_1^{(h)}+\cdots}
{\Psi'(\vartheta_1^{(h)})\delta\eta_1^{(h)}+\cdots}}
\]
gleich dem Logarithmus einer \(m^{\mathrm{ten}}\) Wurzel der Einheit ist, oder dass:
\[
(7)\quad
\frac{\dfrac{\varphi'(\vartheta_1^{(h)})\delta\xi_1^{(h)}+\cdots}
{\psi'(\vartheta_1^{(h)})\delta\xi_1^{(h)}+\cdots}}
{\dfrac{\varphi'(\vartheta_1^{(h)})\delta\eta_1^{(h)}+\cdots}
{\psi'(\vartheta_1^{(h)})\delta\eta_1^{(h)}+\cdots}}
=
\sqrt[m]{
\frac{\dfrac{\Phi'(\vartheta_1^{(h)})\delta\xi_1^{(h)}+\cdots}
{\Psi'(\vartheta_1^{(h)})\delta\xi_1^{(h)}+\cdots}}
{\dfrac{\Phi'(\vartheta_1^{(h)})\delta\eta_1^{(h)}+\cdots}
{\Psi'(\vartheta_1^{(h)})\delta\eta_1^{(h)}+\cdots}}
},
\qquad (h=1,2\ldots q).
\]
In dieser Gleichung ist, wenn die betreffende Lösung des speciellen Theilungsproblems für das Jacobi'sche Umkehrproblem bekannt ist, alles gegeben, bis auf die Coefficienten der Function \(\varphi\); und man hat daher aus vorliegender Formel \(q\) weitere lineare Gleichungen, welche dazu dienen, diese Coefficienten vollständig zu bestimmen. Man erhält dieselben ausgedrückt durch rationale Functionen und \(m^{\mathrm{te}}\) Wurzeln rationaler Functionen einer Wurzel des als bekannt vorausgesetzten Theilungsproblems, und kann also auch eine Wurzel des erweiterten Umkehrproblems auf solche Weise durch jene ausdrücken. Zu jeder Wurzel des Jacobi'schen Theilungsproblems gehören also \(m^q\) Wurzeln des erweiterten, welche man erhält, indem man den \(q\) verschiedenen \(m^{\mathrm{ten}}\) Wurzeln unabhängig von einander alle möglichen Werthe beilegt. Die \(m^q\) Curven einer solchen Gruppe stehen in der oben erwähnten Beziehung zu einander. Wenn man durch alle Doppel- und Rückkehrpuncte, durch die \(c\) und die \(\gamma\) eine Curve \(\psi=0\) legt, so liegt jedes einer Curve der Gruppe angehörige System der \(x\) auf einer Curve \(\varphi=0\), welche mit \(\psi=0\) von gleicher Ordnung ist, mit den Doppel- und Rückkehrpuncten, den Puncten \(z\) und den übrigen Schnittpuncten von \(\psi=0\) mit \(f=0\).

Unter den Gruppen ist aber eine, welche der identischen Wurzel der Jacobi'schen Theilungsgleichung zugeordnet ist, für welche die \(z\) mit den \(\gamma\) zusammenfallen. Für diese wird die Function, aus welcher in (7) die \(m^{\mathrm{te}}\) Wurzel gezogen wurde, der Einheit gleich; und eine der Curven \(\varphi\) fällt mit \(\psi\) zusammen, diejenige nämlich, für welche die \(m^{\mathrm{te}}\) Wurzel selbst gleich 1 gesetzt wird. Man bemerkt, dass die übrigen \(m^q-1\) Curven dieser Gruppe gefunden werden ohne Hinzunahme der Theilungsgleichung für das Jacobi'sche Umkehrproblem, dass sie also keine Irrationalitäten erfordern als die \(m^{\mathrm{ten}}\) Wurzeln der Einheit.

\subsection*{\S\,83. Theilung für \(m=2\).}

Was nun insbesondre die specielle Theilung des erweiterten Umkehrproblems für \(m=2\) betrifft, so bemerken wir, dass sich aus dem auf p. 279 gegebenen Theorem für \(\Theta^{(q)}\) die Periodicitätsformel ergiebt:
\[
\Theta^{(q)}(\omega+P)=
\exp\!\left\{\sum q_i\omega_i+\frac12\sum\sum q_hq_ka_{hk}
+\pi\sqrt{-1}\sum n_i\right\}\Theta^{(q)}(\omega),
\]
wenn nämlich
\[
\begin{array}{ll}
P_h=2m_h\pi\sqrt{-1}+q_1a_{1h}+q_2a_{2h}+\cdots,&(h=1,2\ldots p),\\
P_{p+h}=2n_h\pi\sqrt{-1}+q_1u_{1h}+q_2u_{2h}+\cdots,&(h=1,2\ldots q).
\end{array}
\]
Setzen wir nun \(\omega_h=-\frac{P_h}{2}\), so haben wir
\[
\Theta^{(q)}\!\left(\frac P2\right)
=e^{\pi\sqrt{-1}\left(\sum n_i-\sum m_iq_i\right)}
\Theta^{(q)}\!\left(-\frac P2\right),
\]
oder, mit Rücksicht auf die Gleichung
\[
\Theta^{(q)}(-\omega)=(-1)^q\Theta^{(q)}(\omega),
\]
auch:
\[
0=\Theta^{(q)}\!\left(\frac P2\right)
\left(1-(-1)^{\sum m_iq_i+\sum n_i+q}\right).
\]
Der Ausdruck \(\Theta^{(q)}\!\left(\frac P2\right)\) verschwindet also immer, wenn
\[
(1)\quad \sum m_iq_i+\sum n_i+q=2H+1,
\]
und man kann wie in \(\S\,75\) zeigen, dass er in jedem andern Falle im Allgemeinen nicht verschwindet. In der Formel (1) kann man offenbar alle \(m_i,q_i,n_i\) willkürlich wählen, bis auf eines der letztern, welches durch die übrigen bestimmt ist. Man sieht also, dass unter den Ausdrücken \(\Theta^{(q)}\!\left(\frac P2\right)\) eine Anzahl von \(2^{2p+q-1}\) verschwindet, und dass eine ebenso grosse Anzahl nicht verschwindet; dass also auch unter den \(2^{2p+q}\) Functionen \(\Theta^{(q)}\!\left(\omega+\frac P2\right)\) ebenso viele gerade als ungerade sind.

Man zeigt wie im zehnten Abschnitt, dass die nicht verschwindenden Functionen \(\Theta^{(q)}\!\left(\frac P2\right)\) den wirklichen Berührungscurven entsprechen, welche durch die \(\delta+r\) abgesonderten Doppel- und Rückkehrpuncte und durch die \(n-2\) übrigen Schnittpuncte der in dem Puncte \(\mu\) gezogenen Tangente gehen, und ausserdem in \(p+q\) Puncten die gegebene Curve berühren. Die andern Functionen \(\Theta^{(q)}\!\left(\frac P2\right)\) dagegen entsprechen solchen Berührungscurven, welche in die Tangente in \(\mu\) selbst und in Curven \((n-3)^{\mathrm{ter}}\) Ordnung zerfallen. Man schliesst daraus weiter, dass es \(2^{2p+q-1}\) Curven \((n-3)^{\mathrm{ter}}\) Ordnung giebt, welche durch \(\delta+r=d+r-q\) Doppel- und Rückkehrpuncte der Curve gehen und die Curve \(f=0\) in \(p+q\) Puncten berühren. Zu der Auffindung dieser Curven ist aber nach dem Vorigen die Auflösung keiner andern Gleichung erforderlich als derjenigen, welche die durch alle Doppel- und Rückkehrpuncte gelegten, in \(p\) Puncten berührenden Curven \((n-3)^{\mathrm{ter}}\) Ordnung liefert.

\subsection*{\S\,84. Einfluss von Rückkehrpuncten.}

Es bleibt endlich übrig, derjenigen Modificationen Erwähnung zu thun, welche die vorstehenden Untersuchungen erfahren, sobald einige der \(q\) bevorzugten Doppelpuncte in Rückkehrpuncte übergehen. Setzen wir im Vorigen \(q+r\) an die Stelle von \(r\), und denken wir uns in der Gleichung \(f=0\) \(r\) variable Parameter, welche sich unabhängig von einander der Grenze Null nähern, während dabei zugleich die Unstetigkeitspuncte der \(r\) letzten Integrale \(u_{p+q+1},u_{p+q+2},\ldots,u_{p+q+r}\) aus Doppelpuncten zu Rückkehrpuncten werden. Die \(r\) variabeln Parameter seien \(\tau_1,\tau_2\ldots\tau_r\); wir können sie uns geometrisch etwa so definirt denken, dass man in den betreffenden Doppelpuncten der Curve \(f=0\) Strahlbüschel legt, und die Differenz der Parameter derjenigen Strahlen in jedem Büschel durch \(\tau_i\) bezeichnet, welche die Tangenten des Doppelpuncts sind. Während die \(\tau\) sich der Null, und also die Doppelpuncte Rückkehrpuncten nähern, werden jene \(r\) letzten Integrale durch Annäherung der Parameterpuncte unendlich klein, und man kann setzen:
\[
u_{p+q+i}=\tau_iU_i,
\qquad (i=1,2\ldots r),
\]
wo nun die \(U\) endliche Integrale zweiter Gattung sind. In den Gleichungen des Umkehrproblems hat man dann die \(r\) letzten Gleichungen:
\[
\sum_{h=1}^{h=p+q+r}\int_{c^{(h)}}^{x^{(h)}}du_{p+q+i}=v_{p+q+i},
\qquad (i=1,2\ldots r)
\]
zu ersetzen durch
\[
(1)\quad \sum_{h=1}^{h=p+q+r}\int_{c^{(h)}}^{x^{(h)}}dU_i=V_i,
\]
indem
\[
v_{p+q+i}=\tau_iV_i
\]
wird, und die \(V_i\) jetzt endliche Variable bedeuten.

Die Integrale (1) besitzen als Integrale zweiter Gattung keinen Periodicitätsmodul \(2\pi\sqrt{-1}\) mehr, dagegen mag \(V_i\) in
\[
V_i'=V_i+U_{1i}q_1+U_{2i}q_2+\cdots+U_{pi}q_p,
\qquad (i=1,2\ldots r)
\]
übergehen, während die Integralsummen erster und dritter Gattung sich in
\[
\begin{array}{ll}
v_i'=v_i+2m_i\pi\sqrt{-1}+a_{1i}q_1+a_{2i}q_2+\cdots+a_{pi}q_p,&(i=1,2\ldots p),\\
v_{p+i}'=v_{p+i}+2n_i\pi\sqrt{-1}+u_{1i}q_1+u_{2i}q_2+\cdots+u_{pi}q_p,&(i=1,2\ldots q)
\end{array}
\]
verwandeln.

Den Transcendenten \(U^{(q)}\), \(V^{(q)}\) entsprechen bei diesem Grenzübergange andere Transcendenten \(U^{(q,r)}\), \(V^{(q,r)}\), deren Form man aus \(\S\,77\) sofort abschreiben kann, und welche sich bezüglich ihrer Stetigkeit und ihres Verschwindens ganz ähnlich wie jene Functionen verhalten.

Anders ist es mit der Function \(\Theta^{(q,r)}\), welche aus \(\Theta^{(q+r)}\) durch den Grenzübergang hervorgeht. Man sieht nämlich leicht, dass die Function \(\Theta^{(q)}\) verschwindet, sobald irgend einer der \(q\) in ihr auftretenden Doppelpuncte sich in einen Rückkehrpunct verwandelt. Geht z. B. der \(q^{\mathrm{te}}\) Doppelpunct in einen solchen über, so vereinigen sich die bisher als getrennt zu betrachtenden Puncte \(\varepsilon^{(q)},\varepsilon^{(-q)}\) in zwei unendlich nahe Puncte der Curve; es verschwindet \(\omega_{p+q}\) nebst den \(u_{hq}\) und den \(a_q,b_{hq}\); in \(\Theta^{(q)}\) werden die mit
\[
(-1)^{\frac{-s_1-s_2-\cdots-1}{2}}
\quad\text{und}\quad
(-1)^{\frac{-s_1-s_2-\cdots+1}{2}}
\]
multiplicirten Glieder einander gleich, und zerstören sich also, da diese Potenzen von \(\sqrt{-1}\) immer entgegengesetzte Werthe haben.

Man sieht hieraus, dass, indem \(r\) Doppelpuncte von \(\Theta^{(q+r)}\) in Rückkehrpuncte übergehen, diese Function sich der Grenze nähert
\[
\lim \Theta^{(q+r)}=\tau_1\tau_2\ldots\tau_r\Theta^{(q,r)},
\]
wo nun erst \(\Theta^{(q,r)}\) eine endliche Function ist. Da überall nur Quotienten der \(\Theta^{(q)}\) vorkommen, so erhält man die Formeln des neuen Umkehrproblems, indem man überall diese Function an Stelle von \(\Theta^{(q+r)}\) benutzt, eine Function, welche nach dem Vorigen sowohl synektisch sein muss, als auch die Eigenschaft behält, nur zu verschwinden unter den bei \(\Theta^{(q)}\) angegebenen Umständen.

Was nun die wirkliche Bildung der Function \(\Theta^{(q,r)}\) betrifft, so setzt man:
\[
\begin{array}{lll}
\omega_{p+q+i}=\tau_iW_i,&(i=1,2\ldots r),\\
b_{h,q+i}=\tau_iB_{hi},&(i=1,2\ldots r,\ h=1,2\ldots q),\\
b_{q+k,q+i}=\tau_i\tau_kC_{ki},&(i,k=1,2\ldots r),\\
u_{h,q+i}=\tau_iU_{hi},&(i=1,2\ldots r,\ h=1,2\ldots p),
\end{array}
\]
und findet dann
\[
\begin{aligned}
&\Theta^{(q,r)}(w_1,w_2\ldots w_{p+q},W_1,W_2\ldots W_r)\\
&=\frac{\partial^q}{\partial\tau_1\partial\tau_2\cdots}
\Bigg\{\sum_{s_1s_2\ldots s_r=\pm1}
\exp\!\left(\frac12\sum_{i=1}^{i=r}s_i\tau_iW_{p+i}
-\frac{\pi\sqrt{-1}}2\sum_{i=1}^{i=r}s_i
-\frac14\sum_{i\ge k}\sum s_is_k\tau_i\tau_kC_{ik}\right)\\
&\qquad\cdot
\Theta^{(q)}\!\left(\omega_h-\frac{\tau_1s_1U_{h1}+\tau_2s_2U_{h2}+\cdots}{2},\
\omega_{p+h}-\frac{\tau_1s_1B_{h1}+\tau_2s_2B_{h2}+\cdots}{2}\right)\Bigg\},
\end{aligned}
\]
wo rechts in \(\Theta^{(q)}\) nur zwei Argumente an Stelle von \(p+q\) angegeben sind, und wo nach Ausführung der Differentiation sämmtliche \(\tau\) gleich Null zu setzen sind.

Die Aufgabe der Theilung wird bei dem Uebergange von Doppelpuncten in Rückkehrpuncte dadurch modificirt, dass den Integralen zweiter Gattung keine eigenthümlichen Periodicitätsmoduln mehr zukommen. Daher hat die Aufgabe der allgemeinen Theilung für das vorliegende Umkehrproblem mit \(p+q+r\) Variabeln doch nur \(m^{2p+q}\) Lösungen, die der speciellen Theilung deren \(m^{2p+q}-1\).

Dagegen bleibt sowohl die Eigenschaft eine gerade oder ungerade Function zu sein, je nachdem \(q+r\) gerade oder ungerade ist, der Function \(\Theta^{(q,r)}\) erhalten, als auch die Eigenschaft, durch Vermehrung der Variabeln um ein System \(P_h\) von Periodicitätsmoduln den Factor
\[
\exp\!\left\{\sum_{i=1}^{i=p}q_i\omega_i
+\frac12\sum_{h,k=1}^{h,k=p}q_hq_ka_{hk}
+\pi\sqrt{-1}\sum_{i=1}^{i=q}n_i\right\}
\]
zu erwerben. Daher hat man auch, wie in \(\S\,83\), \(\Theta^{(q,r)}\left(\frac P2\right)=0\), sobald
\[
\sum_{i=1}^{i=p}m_iq_i+\sum_{i=1}^{i=q}n_i+q+r=2H+1.
\]
Es giebt also hier wie dort \(2^{2p+q-1}\) eigentliche Berührungscurven \((n-2)^{\mathrm{ter}}\) Ordnung, welche durch die ausser den \(q+r\) bevorzugten noch vorhandenen Ausnahmspuncte und durch die Schnittpuncte der Tangente von \(\mu\) gehen, und ausserdem in \(p+q+r\) Puncten berühren; ebenso gross ist die Zahl von Curven \((n-3)^{\mathrm{ter}}\) Ordnung, welche durch jene Ausnahmspuncte gehen und in \(p+q+r-1\) Puncten berühren. Davon ausgenommen ist nur der Fall \(q=0\), der Fall also, in welchem sämmtliche bevorzugte Doppelpuncte zu Rückkehrpunkten werden. In diesem Falle kann man nicht mehr die \(m_i,q_i\) beliebig wählen, sondern man hat für sie die Gleichung
\[
\sum_{i=1}^{i=p}m_iq_i+r=2H+1,
\]
damit das entsprechende System von halben Periodicitätsmoduln auf eine Berührungscurve \((n-3)^{\mathrm{ter}}\) Ordnung führe. Daher giebt es in diesem Fall \(2^{p-1}(2^p-1)\) Berührungscurven \((n-3)^{\mathrm{ter}}\) Ordnung, welche durch die Doppel- und Rückkehrpuncte, \(r\) der letzteren ausgenommen, gehen, und in \(p+r-1\) Puncten berühren, sobald \(r\) gerade ist, aber \(2^{p-1}(2^p+1)\) solcher Curven, wenn \(r\) ungerade; und gerade umgekehrt verhält es sich mit den eigentlichen Berührungscurven \((n-2)^{\mathrm{ter}}\) Ordnung, welche durch die eben bezeichneten Doppel- und Rückkehrpuncte so wie durch die \(n-2\) Schnittpuncte der Tangente eines Punctes \(\mu\) gehen, und welche in \(p+r\) Puncten berühren.


% --- Source pages 322--331; printed pages 300--309. ---

\section*{Zwölfter Abschnitt.\\Die unendlich vielen Formen der Function \(\Theta\).}

\subsection*{\S\,85. Kanonische Determinanten.}

Im vierten Abschnitte wurde gezeigt, wie man zu einem Systeme normaler Periodicitätsmoduln gelangen könne. Dies geschah dadurch, dass man mittelst der linearen ganzzahligen Relationen
\[
\begin{aligned}
(1)\quad
M^{(h)}&=\alpha_1^{(h)}(a_1)+\alpha_2^{(h)}(a_2)+\cdots+\alpha_{2p}^{(h)}(a_{2p}),\\
N^{(h)}&=\alpha_1^{(h)}(b_1)+\alpha_2^{(h)}(b_2)+\cdots+\alpha_{2p}^{(h)}(b_{2p})
\qquad (h=1,2\ldots2p)
\end{aligned}
\]
die identische Transformationsgleichung
\[
(2)\quad
\sum\sum c_{ik}(a_i)(b_k)
=\sum_{h=1}^{h=p}\left(M^{(2h-1)}N^{(2h)}-N^{(2h-1)}M^{(2h)}\right)
\]
herstellte, wobei die \(c\) ganze Zahlen bedeuteten, \(c_{ik}=-c_{ki}\), und die Determinante der \(c\) gleich 1 war. Sind dann die \((a),(b)\) Periodicitätsmoduln von irgend zwei Integralen erster Gattung, so waren die Ausdrücke (1) normale Periodicitätsmoduln derselben, wie sie zur Aufstellung der im Vorigen ausgeführten Theorie nothwendig waren.

Die Transformationsformel (2) ist auf unendlich viele Arten erfüllbar. Nehmen wir an, sie sei auf eine Weise erfüllt; es sei also ein System der \(\alpha,M,N\) gegeben; suchen wir die Bedingungen auf, denen alle andern Lösungen des Transformationsproblems zu genügen haben.

Setzen wir also, um ein neues System normaler Periodicitätsmoduln zu erhalten:
\[
\begin{aligned}
(3)\quad
R_k&=\beta_{1,k}M^{(1)}+\beta_{2,k}M^{(2)}+\cdots+\beta_{2p,k}M^{(2p)},\\
S_k&=\beta_{1,k}N^{(1)}+\beta_{2,k}N^{(2)}+\cdots+\beta_{2p,k}N^{(2p)}
\qquad (k=1,2\ldots2p).
\end{aligned}
\]
Die ganzen Zahlen \(\beta\) sind dann so zu bestimmen, dass identisch
\[
(4)\quad
\sum_{k=1}^{k=p}(R_{2k-1}S_{2k}-S_{2k-1}R_{2k})
=\sum_{h=1}^{h=p}\left(M^{(2h-1)}N^{(2h)}-N^{(2h-1)}M^{(2h)}\right),
\]
oder, wenn man links die Grössen (3) durch ihre Werthe ersetzt, dass:
\[
\begin{aligned}
&\sum_{k=1}^{k=p}\sum_{s,r=1}^{s,r=2p}
(\beta_{s,2k-1}\beta_{r,2k}-\beta_{r,2k-1}\beta_{s,2k})M^{(s)}N^{(r)}\\
&\hspace{2em}=\sum_{h=1}^{h=p}\left(M^{(2h-1)}N^{(2h)}-N^{(2h-1)}M^{(2h)}\right).
\end{aligned}
\]
Die Vergleichung der Coefficienten zeigt also, dass die ganzen Zahlen \(\beta\) die folgenden Bedingungen befriedigen müssen, in denen wir, um keine doppelt zu erhalten, \(s<r\) voraussetzen können:
\[
(5)\quad
\sum_{k=1}^{k=p}\left(\beta_{s,2k-1}\beta_{r,2k}-\beta_{r,2k-1}\beta_{s,2k}\right)=0,
\]
ausgenommen, wenn \(s\) eine ungerade Zahl, und \(r\) die darauf folgende gerade Zahl ist, wo denn
\[
(6)\quad
\sum_{k=1}^{k=p}\left(\beta_{2h-1,2k-1}\beta_{2h,2k}-\beta_{2h,2k-1}\beta_{2h-1,2k}\right)=1.
\]
Determinanten der Art, wie sie hier aus den ganzen Zahlen \(\beta\) zusammengesetzt werden\footnote{Solche Determinanten hat zuerst Hermite in den Comptes Rendus von 1855 untersucht, und zwar bei Gelegenheit allgemeiner Transformationen der Function \(\Theta\) für \(p=2\).}, wollen wir \emph{kanonische} nennen. Indem wir also als Elemente einer kanonischen Determinante solche ganze Zahlen definiren, welche den Gleichungen (5), (6) genügen, können wir den in dieser Definition enthaltenen Satz aufstellen:

\begin{center}
\textbf{Der Uebergang von einem System normaler Periodicitätsmoduln zu einem andern geschieht mit Hülfe einer kanonischen Transformationsdeterminante, und jede solche ergiebt einen der verlangten Uebergänge.}
\end{center}

Die kanonische Determinante hat den Werth \(\pm1\); denn ihr Quadrat, multiplicirt mit der Determinante
\[
\left|
\begin{array}{rrrrrrr}
0&1&0&0&\cdots&0&0\\
-1&0&0&0&\cdots&0&0\\
0&0&0&1&\cdots&0&0\\
0&0&-1&0&\cdots&0&0\\
\multicolumn{7}{c}{\cdots}\\
0&0&0&0&\cdots&0&1\\
0&0&0&0&\cdots&-1&0
\end{array}
\right|=1
\]
der lineo-linearen Function, welche die rechte Seite von (4) bildet, giebt die mit der vorigen ganz identische Determinante der lineo-linearen Function auf der linken Seite von (4); das Quadrat der kanonischen Determinante ist also 1. Bezeichnen wir den Werth der kanonischen Determinante selbst durch \(\varepsilon=\pm1\), ferner durch \(\gamma_{hk}\) die dem Elemente \(\beta_{hk}\) zugeordnete Unterdeterminante. Die Auflösung der Gleichungen (3) giebt dann:
\[
\begin{aligned}
(7)\quad
M^{(h)}&=\varepsilon[\gamma_{h1}R_1+\gamma_{h2}R_2+\cdots+\gamma_{h,2p}R_{2p}],\\
N^{(h)}&=\varepsilon[\gamma_{h1}S_1+\gamma_{h2}S_2+\cdots+\gamma_{h,2p}S_{2p}].
\end{aligned}
\]
Ersetzt man nun auf der linken Seite der Gleichung (4) die \(R\) durch ihre Werthe aus (3), auf der rechten Seite die \(N\) durch ihre Werthe aus (7), so findet man die Gleichung:
\[
\begin{aligned}
&\sum_{k=1}^{k=p}\sum_{r=1}^{r=2p}(\beta_{r,2k-1}S_{2k}-\beta_{r,2k}S_{2k-1})M^{(r)}\\
&\hspace{2em}=\varepsilon\sum_{h=1}^{h=p}\sum_{s=1}^{s=2p}
(\gamma_{2h,s}M^{(2h-1)}-\gamma_{2h-1,s}M^{(2h)})S_s,
\end{aligned}
\]
und durch Vergleichung der beiderseitigen Coefficienten erhält man für die Unterdeterminanten \(\gamma\) die folgenden Werthe:
\[
\begin{array}{rlrl}
(8)\quad
\gamma_{2h,2k}&=\varepsilon\beta_{2h-1,2k-1},&
\gamma_{2h,2k-1}&=-\varepsilon\beta_{2h-1,2k},\\
\gamma_{2h-1,2k}&=-\varepsilon\beta_{2h,2k-1},&
\gamma_{2h-1,2k-1}&=\varepsilon\beta_{2h,2k}.
\end{array}
\]
Daher kann man die Gleichungen (7), die Auflösungen von (4), auch in der Form schreiben:
\[
\begin{aligned}
(9)\quad
M^{(2h-1)}&=\beta_{2h,2}R_1-\beta_{2h,1}R_2+\beta_{2h,4}R_3-\beta_{2h,3}R_4+\cdots,\\
M^{(2h)}&=-\beta_{2h-1,2}R_1+\beta_{2h-1,1}R_2-\beta_{2h-1,4}R_3+\beta_{2h-1,3}R_4+\cdots,\\
N^{(2h-1)}&=\beta_{2h,2}S_1-\beta_{2h,1}S_2+\beta_{2h,4}S_3-\beta_{2h,3}S_4+\cdots,\\
N^{(2h)}&=-\beta_{2h-1,2}S_1+\beta_{2h-1,1}S_2-\beta_{2h-1,4}S_3+\beta_{2h-1,3}S_4+\cdots.
\end{aligned}
\]
Der Definition nach ist jede Transformationsdeterminante, welche zwei Systeme normaler Periodicitätsmoduln verbindet, eine kanonische; unter anderm also auch die Determinante der Gleichungen (7) oder (9). Sie kommt auch der Determinante der \(\gamma\) selbst zu; denn da die Gleichungen (5), (6) erfüllt sind, wenn man die \(\beta\) durch die \(\varepsilon\cdot\gamma\) ersetzt, so sind sie auch erfüllt, wenn man die \(\gamma\) allein an Stelle der \(\beta\) treten lässt.

Drückt man in den Gleichungen (3) die \(M,N\) durch irgend zwei andre Systeme normaler Periodicitätsmoduln aus:
\[
\begin{aligned}
M^{(h)}&=\delta_{1,h}V_1+\delta_{2,h}V_2+\cdots+\delta_{2p,h}V_{2p},\\
N^{(h)}&=\delta_{1,h}W_1+\delta_{2,h}W_2+\cdots+\delta_{2p,h}W_{2p},
\end{aligned}
\]
so gehen die Formeln (3) über in die folgenden:
\[
\begin{aligned}
R_k&=\vartheta_{1,k}V_1+\vartheta_{2,k}V_2+\cdots+\vartheta_{2p,k}V_{2p},\\
S_k&=\vartheta_{1,k}W_1+\vartheta_{2,k}W_2+\cdots+\vartheta_{2p,k}W_{2p},
\end{aligned}
\]
wo
\[
\vartheta_{i,k}=\beta_{1,k}\delta_{i,1}+\beta_{2,k}\delta_{i,2}+\cdots+\beta_{2p,k}\delta_{i,2p}.
\]
Die Determinanten der \(\delta\) und der \(\vartheta\) sind der Definition nach kanonisch. Daher hat man den Satz:

\begin{center}
\textbf{Durch Multiplication zweier kanonischen Determinanten}
\end{center}
\[
\left|\begin{array}{cccc}
\beta_{11}&\beta_{21}&\cdots&\beta_{2p,1}\\
\beta_{12}&\beta_{22}&\cdots&\beta_{2p,2}\\
\multicolumn{4}{c}{\cdots}\\
\beta_{1,2p}&\beta_{2,2p}&\cdots&\beta_{2p,2p}
\end{array}\right|,
\qquad
\left|\begin{array}{cccc}
\delta_{11}&\delta_{21}&\cdots&\delta_{2p,1}\\
\delta_{12}&\delta_{22}&\cdots&\delta_{2p,2}\\
\multicolumn{4}{c}{\cdots}\\
\delta_{1,2p}&\delta_{2,2p}&\cdots&\delta_{2p,2p}
\end{array}\right|
\]
entsteht, indem bei der einen Horizontal- und Verticalreihen vertauscht werden, die ebenfalls kanonische Determinante
\[
\left|\begin{array}{cccc}
\vartheta_{11}&\vartheta_{21}&\cdots&\vartheta_{2p,1}\\
\vartheta_{12}&\vartheta_{22}&\cdots&\vartheta_{2p,2}\\
\multicolumn{4}{c}{\cdots}\\
\vartheta_{1,2p}&\vartheta_{2,2p}&\cdots&\vartheta_{2p,2p}
\end{array}\right|,
\]
in welcher
\[
\vartheta_{ik}=\beta_{1,k}\delta_{i,1}+\beta_{2,k}\delta_{i,2}+\cdots+\beta_{2p,k}\delta_{i,2p}.
\]
Da es übrigens für das Bestehen dieses Satzes gleichgültig sein muss, bei welcher der beiden Determinanten Horizontal- und Verticalreihen vertauscht werden, und da die beiden entstehenden Determinanten sich dann selbst nur durch Vertauschung der Horizontal- und Verticalreihen unterscheiden, so folgt der weitere Satz:

\begin{center}
\textbf{Eine kanonische Determinante verliert diese Eigenschaft nicht, wenn man in ihr die Horizontalreihen mit den Verticalreihen vertauscht.}
\end{center}

Der obige Satz über die Multiplication bleibt also auch noch richtig, wenn man den Zusatz über die Reihenvertauschung auslässt; aber dann geben die Elemente der resultirenden Determinante nicht mehr die Coefficienten der Substitution an, welche sich aus den Substitutionen zusammensetzt, die den einzelnen mit einander multiplicirten Determinanten entsprechen.

Die Frage nach der allgemeinsten kanonischen Determinante oder dem allgemeinsten System normaler Periodicitätsmoduln kann nun so aufgefasst werden, dass man diejenigen einfachsten kanonischen Determinanten aufsucht, aus denen durch Multiplication alle andern sich ableiten lassen. Die entsprechenden linearen Substitutionen führen dann von gewissen ursprünglichsten Systemen normaler Periodicitätsmoduln zu allen möglichen.

Statt aber aus möglichst einfachen Determinanten die gegebene zusammenzusetzen, kann man, wenn nur, wie angenommen werden soll, die den einfachsten adjungirten kanonischen Determinanten wieder zu den einfachsten gehören, auch so verfahren, dass man die gegebene Determinante durch Multiplication mit einfachen auf eine einfache zurückführt. Die letzte mit den adjungirten der andern multiplicirt, ist dann gleich der gegebenen, also diese in einfache Factoren aufgelöst. Es entstehen also die beiden Fragen: erstlich, welche kanonische Determinanten wir als \emph{einfache} betrachten wollen; zweitens, wie man durch Multiplication mit einfachen eine gegebene selbst auf eine einfache zurückführt.

\subsection*{\S\,86. Reduction der kanonischen Determinanten auf einfache.}

Es genügt, als \emph{einfachste} kanonische Determinanten folgende aufzuführen:

I. Die Diagonalglieder sind 1, bis auf zwei, \(\beta_{2h-1,2h-1}\) und \(\beta_{2h,2h}\), welche Null sind. Dafür sind \(\beta_{2h,2h-1}=\pm1\), \(\beta_{2h-1,2h}=\mp1\); alle übrigen Elemente sind Null. Solcher Determinanten giebt es, abgesehen von dem Zeichen \(\pm\), \(p\). Für \(p=2\) sind sie:
\[
\left|\begin{array}{rrrr}0&\pm1&0&0\\ \mp1&0&0&0\\ 0&0&1&0\\ 0&0&0&1\end{array}\right|,
\qquad
\left|\begin{array}{rrrr}1&0&0&0\\ 0&1&0&0\\ 0&0&0&\pm1\\ 0&0&\mp1&0\end{array}\right|.
\]

II. Die Diagonalglieder sind 1, ein Element \(\beta_{2h-1,2h}=\pm1\), alle übrigen Glieder Null. Solcher Determinanten giebt es, abgesehen von dem Zeichen \(\pm\), ebenfalls \(p\). Für \(p=2\):
\[
\left|\begin{array}{rrrr}1&0&0&0\\ \pm1&1&0&0\\ 0&0&1&0\\ 0&0&0&1\end{array}\right|,
\qquad
\left|\begin{array}{rrrr}1&0&0&0\\ 0&1&0&0\\ 0&0&1&0\\ 0&0&\pm1&1\end{array}\right|.
\]

III. Die Diagonalglieder sind 1, ein Element \(\beta_{2h-1,2k-1}\ (k\gtrless h)=\pm1\), \(\beta_{2k,2h}=\mp1\), die übrigen Glieder Null. Abgesehen von dem Zeichen \(\pm\) giebt es \(p(p-1)\) solcher Determinanten. Für \(p=2\):
\[
\left|\begin{array}{rrrr}1&0&0&0\\ 0&1&0&\mp1\\ \pm1&0&1&0\\ 0&0&0&1\end{array}\right|,
\qquad
\left|\begin{array}{rrrr}1&0&\pm1&0\\ 0&1&0&0\\ 0&0&1&0\\ 0&\mp1&0&1\end{array}\right|.
\]
Dass diese Determinanten sämmtlich kanonisch sind, ergiebt sich sogleich aus der Betrachtung der zu erfüllenden Bedingungsgleichungen. Aber auch die zweite Bedingung ist erfüllt, dass die aus den Unterdeterminanten gebildeten Determinanten wieder von derselben Art sind wie die ursprünglichen. In der That, wenn man für irgend eine der obigen Determinanten die adjungirte bildet, so kommt immer die ursprüngliche heraus, nur dass \(\pm\) in \(\mp\), \(\mp\) in \(\pm\) verwandelt ist.

Sehen wir nun, welche Wirkung die Multiplication einer gegebenen kanonischen Determinante
\[
\left|\begin{array}{cccc}
\beta_{11}&\beta_{21}&\cdots&\beta_{2p,1}\\
\beta_{12}&\beta_{22}&\cdots&\beta_{2p,2}\\
\multicolumn{4}{c}{\cdots}\\
\beta_{1,2p}&\beta_{2,2p}&\cdots&\beta_{2p,2p}
\end{array}\right|
\]
mit einer der unter I, II, III aufgeführten einfachsten kanonischen Determinanten auf erstere äussert. Damit die der resultirenden Determinante entsprechende Substitution eine Combination der beiden Substitutionen sei, welche der gegebenen und der einfachen Determinante entsprechen, vertauschen wir vor der Multiplication bei der letztern die Horizontalreihen mit den Verticalreihen.

Der Effect der Multiplication der Determinante der \(\beta\) mit einer einfachen Determinante I. ist offenbar (vgl. \(\S\,86\)), dass zwei Verticalreihen, die \((2h-1)^{\mathrm{te}}\) und die \((2h)^{\mathrm{te}}\) vertauscht werden, und dass dabei eine derselben das Vorzeichen ändert.

Der Effect der Multiplication mit einer einfachen Determinante II. ist kein anderer, als dass die \((2h)^{\mathrm{te}}\) Verticalreihe der \((2h-1)^{\mathrm{ten}}\) hinzugefügt oder von derselben abgezogen wird.

Der Effect endlich der Multiplication mit einer Determinante III. besteht darin, dass die \((2h-1)^{\mathrm{te}}\) Verticalreihe um die \((2k-1)^{\mathrm{te}}\) vermehrt oder vermindert, und zugleich umgekehrt die \((2k)^{\mathrm{te}}\) um die \((2h)^{\mathrm{te}}\) vermindert oder vermehrt wird.

Wir können also mit Hülfe der Multiplication mit einfachen kanonischen Determinanten immer folgende Operationen an der gegebenen Determinante vornehmen:

I. Eine gerade mit der vorhergehenden ungeraden Verticalreihe vertauschen und dabei das Vorzeichen der einen ändern.

II. Eine gerade Verticalreihe von der vorhergehenden ungeraden abziehen oder zu derselben addiren.

III. Eine ungerade Verticalreihe von einer andern ungeraden abziehen oder zu ihr addiren, und zugleich die auf die letztere folgende gerade Reihe zu der auf erstere folgenden geraden Reihe addiren oder von derselben abziehen.

Wir wollen nun sehen, ob wir durch Anwendung dieser Processe die gegebene Determinante in eine einfache, oder in die identische Substitutionsdeterminante verwandeln können, bei welcher alle nicht in der Diagonale stehenden Glieder durch Nullen ersetzt sind.

Zunächst können wir den obigen drei Processen einen erlaubten vierten hinzufügen. Wenden wir auf die Verticalreihen \((2h-1)\), \((2h)\) den Process I. an, so können wir an die \((2h-1)^{\mathrm{te}}\) Stelle die Reihe \(-(2h)\), an die \((2h)^{\mathrm{te}}\) die Reihe \((2h-1)\) setzen. Mit Hülfe des Processes II. geht nun die an \((2h-1)^{\mathrm{ter}}\) Stelle befindliche Reihe \(-(2h)\) in das Aggregat der Reihe \(\pm(2h-1)-(2h)\) über. Wendet man nun den Process I. nochmals an, so haben wir an \((2h-1)^{\mathrm{ter}}\) Stelle die unveränderte Reihe \((2h-1)\), an \((2h)^{\mathrm{ter}}\) Stelle aber die Differenz der Reihen \((2h)\mp(2h-1)\). Man hat also gefunden, dass man auch folgenden Process immer vornehmen kann:

IV. Eine ungerade Verticalreihe von der darauf folgenden geraden abziehen oder zu derselben addiren.

Durch wechselweise Anwendung der Processe II. und IV. kann man offenbar in irgend einer, z. B. in der ersten Horizontalreihe alle geraden Glieder zerstören. Wenn nun alle geraden Glieder dieser Reihe Null sind, so reducirt sich die Wirkung des Processes III., so weit sie diese Horizontalreihe betrifft, auf das Addiren und Abziehen der übriggebliebenen ungeraden Glieder. Man kann also diese abermals durch einander zerstören, so dass schliesslich ein einziges Glied übrig bleibt, und dieses kann man zum ersten Gliede der Reihe machen. Da die Determinante den Werth \(\pm1\) hat, so ist auch dieses Glied nothwendig \(\pm1\).

Man behandelt jetzt ähnlich die zweite Horizontalreihe, indem man nur das erste Glied derselben stets ausschliesst. Man zerstört alle übrigen ungeraden Glieder durch die Processe II., IV., und sodann die geraden durch III. bis auf das zweite, welches dann ein Factor der Determinante wird und daher den Werth \(\pm1\) hat. Durch dieses zerstört man endlich nach II. das erste Glied.

Nun lehrt die Gleichung (5) \(\S\,85\), dass, wenn in zwei auf einander folgenden Horizontalreihen \((2h-1)\), \((2h)\) alle Elemente bis auf die Diagonalglieder verschwinden, bei den entsprechenden Verticalgliedern dasselbe stattfinden muss. Da die Natur der Determinante sich durch Vertauschung von Horizontal- und Verticalreihen sich nicht ändert, so ist es hinreichend, den umgekehrten Satz zu beweisen. Setzen wir in (5) \(r=2h-1\) oder \(r=2h\), und nehmen an, dass alle \(\beta_{2h-1,i}\), \(\beta_{2h,i}\) bis auf \(\beta_{2h-1,2h-1}\), \(\beta_{2h,2h}\) verschwinden, so giebt die Gleichung (5) sofort \(\beta_{s,2h}=0\), \(\beta_{s,2h-1}=0\), was zu beweisen war.

Die reducirte Determinante enthält also in den ersten Horizontal- und Verticalreihen nur noch die Diagonalglieder. Mit den übrigen Reihen der Determinante kann man also operiren, als ob die ersten beiden Reihen gar nicht vorhanden wären. Man verwandelt dann nach dem obigen Verfahren auch die dritte und vierte Horizontal- und Verticalreihe in solche, welche nur Diagonalglieder enthalten, sodann die fünfte und sechste u. s. w., bis überhaupt nur noch Diagonalglieder übrig sind.

Man kann also jede kanonische Determinante mit Hülfe der Determinanten I., II., III. zusammensetzen, daher auch jede kanonische Substitution aus den \(2p(p+1)\) einfachen Substitutionen, welche jenen \(2p(p+1)\) einfachen Determinanten entsprechen\footnote{Die Idee, ganzzahlige Substitutionsdeterminanten, welche den Werth 1 haben, aus einfachen Determinanten zusammenzusetzen, in denen nur die Ziffern 1 und 0 auftreten, gehört, so viel wir wissen, Hrn. Kronecker an.}.

\subsection*{\S\,87. Transformation der Integrale erster und dritter Gattung bei Veränderung des Systems von Normalperioden.}

Wir wollen jetzt untersuchen, wie die Normalintegrale erster und dritter Gattung sich gestalten, wenn man ein verändertes System normaler Periodicitätsmoduln zu Grunde legt. Zu diesem Zwecke gehen wir zunächst auf die Formeln des \(\S\,30\) zurück. Setzen wir in diesen an Stelle der \(H\) die Normalintegrale erster Gattung \(u_1,u_2,u_3\), welche einem gegebenen Periodensystem entsprechen, und bezeichnen wir durch \(v_1,v_2,v_3\) die neuen Normalintegrale erster Gattung, deren System normaler Periodicitätsmoduln mit dem erstern durch die kanonische Determinante
\[
\left|\begin{array}{cccc}
\beta_{11}&\beta_{21}&\cdots&\beta_{2p,1}\\
\beta_{12}&\beta_{22}&\cdots&\beta_{2p,2}\\
\multicolumn{4}{c}{\cdots}\\
\beta_{1,2p}&\beta_{2,2p}&\cdots&\beta_{2p,2p}
\end{array}\right|
\]
verbunden ist. Nimmt man als ursprüngliches System von normalen Periodicitätsmoduln das folgende:
\[
\begin{array}{cccccc}
2\pi\sqrt{-1}&a_{11}&0&a_{12}&0&a_{13}\cdots\\
0&a_{21}&2\pi\sqrt{-1}&a_{22}&0&a_{23}\cdots\\
0&a_{31}&0&a_{32}&2\pi\sqrt{-1}&a_{33}\cdots\\
\multicolumn{6}{c}{\cdots}
\end{array}
\]
so wird das neue System von Periodicitätsmoduln:
\[
\begin{array}{cccc}
u_1^{(1)}&u_1^{(2)}&\cdots&u_1^{(2p)}\\
u_2^{(1)}&u_2^{(2)}&\cdots&u_2^{(2p)}\\
\multicolumn{4}{c}{\cdots}\\
u_p^{(1)}&u_p^{(2)}&\cdots&u_p^{(2p)}
\end{array}
\]
durch die Formel gegeben:
\[
\left\{
\begin{aligned}
(1)\quad u_1^{(h)}&=\beta_{1,h}\cdot2\pi\sqrt{-1}+a_{11}\beta_{2,h}+a_{21}\beta_{4,h}+a_{31}\beta_{6,h}+\cdots,\\
u_2^{(h)}&=\beta_{3,h}\cdot2\pi\sqrt{-1}+a_{12}\beta_{2,h}+a_{22}\beta_{4,h}+a_{32}\beta_{6,h}+\cdots,\\
u_3^{(h)}&=\beta_{5,h}\cdot2\pi\sqrt{-1}+a_{13}\beta_{2,h}+a_{23}\beta_{4,h}+a_{33}\beta_{6,h}+\cdots,\\
&\hspace{5em}\cdots\qquad (h=1,2\ldots2p).
\end{aligned}
\right.
\]
Setzt man nun
\[
\begin{aligned}
(2)\quad
v_1&=m_{11}u_1+m_{12}u_2+\cdots+m_{1p}u_p,\\
v_2&=m_{21}u_1+m_{22}u_2+\cdots+m_{2p}u_p,\\
&\hspace{5em}\cdots\\
v_p&=m_{p1}u_1+m_{p2}u_2+\cdots+m_{pp}u_p,
\end{aligned}
\]
so bestimmen sich die \(m\) nach \(\S\,30\) aus der Formel
\[
(3)\quad
m_{hk}=2\pi\sqrt{-1}\,\frac{\partial\log\Delta}{\partial u_k^{(2h-1)}},
\]
wo
\[
(4)\quad
\Delta=\Sigma\pm u_1^{(1)}u_2^{(3)}u_3^{(5)}\cdots u_p^{(2p-1)}.
\]
In Folge dieser Bestimmung wird das System der Periodicitätsmoduln der neuen Normalintegrale \(v\) das folgende:
\[
\begin{array}{cccccc}
2\pi\sqrt{-1}&b_{11}&0&b_{12}&0&b_{13}\cdots\\
0&b_{21}&2\pi\sqrt{-1}&b_{22}&0&b_{23}\cdots\\
0&b_{31}&0&b_{32}&2\pi\sqrt{-1}&b_{33}\cdots\\
\multicolumn{6}{c}{\cdots,}
\end{array}
\]
wo
\[
(5)\quad
b_{hk}=m_{h1}u_1^{(2k)}+m_{h2}u_2^{(2k)}+\cdots
=m_{k1}u_1^{(2h)}+m_{k2}u_2^{(2h)}+\cdots.
\]
Nachdem die neuen Integrale erster Gattung auf diese Weise normirt worden, ist es leicht, die Integrale dritter Gattung entsprechend zu normiren. Aus dem Integrale \(S_{\xi\eta}\) entstand das Normalintegral \(\Pi_{\xi\eta}\), indem man diejenigen Periodicitätsmoduln durch Hinzufügen von Integralen erster Gattung verschwinden machte, welche den Periodicitätsmoduln \(2\pi\sqrt{-1}\) entsprechen.

% --- Source pages 332--342; printed pages 310--320. ---

Daher entsteht aus dem für die \(u\) normirten Integrale \(\int_z^x d\Pi_{\xi\eta}\) das für die \(v\) normirte Integral \(\int_z^x dP_{\xi\eta}\) mit Hülfe der Formel:
\[
\int_z^x dP_{\xi\eta}
=
\int_z^x d\Pi_{\xi\eta}
-\frac{1}{2\pi\sqrt{-1}}
 \int_z^x(K_1\,dv_1+K_2\,dv_2+\cdots+K_p\,dv_p),
\]
indem wir diejenigen Periodicitätsmoduln zum Verschwinden bringen, welche den Periodicitätsmoduln \(2\pi\sqrt{-1}\) der \(v\) entsprechen, d.\ h. den Periodicitätsmoduln
\[
u_k^{(1)},\ u_k^{(3)},\ldots,u_k^{(2p-1)}
\]
der \(u\). Man muss daher \(K_h\) gleich dem Periodicitätsmodul von \(\Pi\) setzen, welcher zu \(u_k^{(2h-1)}\) gehört, oder man muss setzen:
\[
K_h=\int_{\eta}^{\xi}(\beta_{2,2h-1}du_1+\beta_{4,2h-1}du_2+\cdots+\beta_{2p,2h-1}du_p),
\]
so dass man für \(P_{\xi\eta}\) die Definitionsgleichung erhält:
\[
(6)\quad
\int_z^x dP_{\xi\eta}
=
\int_z^x d\Pi_{\xi\eta}
-\frac{1}{2\pi\sqrt{-1}}
\sum_{h=1}^{h=p}\sum_{k=1}^{k=p}
\beta_{2k,2h-1}\int_{\eta}^{\xi}du_k\int_z^x dv_h.
\]
Man kann den letzten Theil dieses Ausdrucks, welcher noch die \(u,v\) gemischt enthält, mit Vortheil durch die letztern allein darstellen, und sieht dann auch, dass wirklich, wie es nach dem Satze über Vertauschung von Parameter und Argument nothwendig ist, derselbe ungeändert bleibt, wenn man \(x,z\) mit \(\xi,\eta\) vertauscht. Zu diesem Zwecke stellen wir zunächst die Auflösungen der Gleichungen (2) her, welche sehr einfach sind. Da nämlich die \(u_k\) sich auf die \(u_k^{(2p-1)}\) reduciren, wenn alle \(v\) verschwinden, bis auf \(v_p\), und wenn dieses letztere gleich \(2\pi\sqrt{-1}\) wird, so hat man nothwendig:
\[
(7)\quad
u_k=\frac{1}{2\pi\sqrt{-1}}
\left(u_k^{(1)}v_1+u_k^{(3)}v_2+\cdots+u_k^{(2p-1)}v_p\right).
\]
Daher erhält man aus (6) mit Benutzung dieser Formel:
\[
(8)\quad
\int_z^x dP_{\xi\eta}
=
\int_z^x d\Pi_{\xi\eta}
+\frac{1}{4\pi^2}
\sum_{h=1}^{h=p}\sum_{m=1}^{m=p}
c_{hm}\int_z^x dv_h\int_{\eta}^{\xi}dv_m,
\]
wobei
\[
(9)\quad
c_{hm}=\sum_{k=1}^{k=p}\beta_{2k,2h-1}u_k^{(2m-1)}.
\]
Vertauscht man nun in (8) Parameter und Argument, so findet man die letzte Summe nur insofern geändert, als \(c_{mh}\) an Stelle von \(c_{hm}\) tritt. Man muss also haben
\[
c_{hm}=c_{mh},
\]
oder, wenn man die Werthe der \(u_k^{(2m-1)}\) aus (1) einträgt:
\[
\begin{aligned}
0={}&2\pi\sqrt{-1}\sum_{k=1}^{k=p}
(\beta_{2k,2h-1}\beta_{2k-1,2m-1}
-\beta_{2k,2m-1}\beta_{2k-1,2h-1})\\
&+\sum_{s=1}^{s=p}\sum_{k=1}^{k=p}a_{sk}
(\beta_{2k,2h-1}\beta_{2s,2m-1}
-\beta_{2k,2m-1}\beta_{2s,2h-1}).
\end{aligned}
\]
Diese Gleichung ist in der That erfüllt. Denn da offenbar \(a_{sk}\) und \(a_{ks}\) in gleiche und entgegengesetzte Zahlen multiplicirt sind, so heben sich die Terme der letzten Summe paarweise auf; der Coefficient von \(2\pi\sqrt{-1}\) aber verschwindet wegen der Eigenschaften einer kanonischen Determinante.

\subsection*{\S\,88. Transformation der \(\Theta\)function mittelst kanonischer Substitutionen.}

Summirt man die Gleichung (8) über \(p\) Punctepaare \(x,z\), so gelangt man zu der neuen Form der Transcendente \(T\); und indem man diese weiter behandelt wie es in Abschnitt VII. und VIII. geschah, gelangt man zu den neuen Formen der Functionen \(U,V,\Theta\). Da es schliesslich nur auf die letzte ankommt, so wollen wir von der Gleichung (8) sofort auf diese übergehen. Bezeichnen wir die Transcendente \(\Theta\) jetzt, wo auch die Constanten im Exponenten verändert werden, durch \(\Theta(\omega,a)\), die neue Transcendente aber durch \(\Theta(w,b)\).

Aus der Formel
\[
\sum_{i=1}^{i=p}\int_{z^{(i)}}^{x^{(i)}}d\Pi_{\xi\eta}
=
-\log
\frac{
\Theta\!\left(\sum\int_{\alpha^{(i)}}^{x^{(i)}}du_h-\int_{\mu}^{\xi}du_h,a\right)
\Theta\!\left(\sum\int_{\alpha^{(i)}}^{z^{(i)}}du_h-\int_{\mu}^{\eta}du_h,a\right)
}{
\Theta\!\left(\sum\int_{\alpha^{(i)}}^{x^{(i)}}du_h-\int_{\mu}^{\eta}du_h,a\right)
\Theta\!\left(\sum\int_{\alpha^{(i)}}^{z^{(i)}}du_h-\int_{\mu}^{\xi}du_h,a\right)
}
\]
folgt sofort die analoge
\[
\sum_{i=1}^{i=p}\int_{z^{(i)}}^{x^{(i)}}dP_{\xi\eta}
=
-\log
\frac{
\Theta\!\left(\sum\int_{\gamma^{(i)}}^{x^{(i)}}dv_h-\int_{\mu}^{\xi}dv_h,b\right)
\Theta\!\left(\sum\int_{\gamma^{(i)}}^{z^{(i)}}dv_h-\int_{\mu}^{\eta}dv_h,b\right)
}{
\Theta\!\left(\sum\int_{\gamma^{(i)}}^{x^{(i)}}dv_h-\int_{\mu}^{\eta}dv_h,b\right)
\Theta\!\left(\sum\int_{\gamma^{(i)}}^{z^{(i)}}dv_h-\int_{\mu}^{\xi}dv_h,b\right)
}.
\]
In dieser letztern Formel sind die untern Grenzen \(\alpha\) der erstern durch neue Grenzen \(\gamma\) ersetzt; denn man kann von vorn herein nicht entscheiden, ob nicht die Berührungscurve der \(\alpha\) hier durch eine andere ersetzt werden muss, deren Berührungspuncte wir als Puncte \(\gamma\) bezeichnen.

Indem wir nun die Gleichung (8) über die \(p\) Punctepaare \(x,z\) summiren, und für die Integralsummen obige Werthe einsetzen, erhalten wir eine Relation zwischen den verschiedenen Arten von \(\Theta\), nämlich:
\[
\resizebox{0.98\linewidth}{!}{$
\begin{aligned}
\log
\frac{
\frac{
\Theta\!\left(\sum\int_{\gamma^{(i)}}^{x^{(i)}}dv_h-\int_{\mu}^{\xi}dv_h,b\right)
\Theta\!\left(\sum\int_{\gamma^{(i)}}^{z^{(i)}}dv_h-\int_{\mu}^{\eta}dv_h,b\right)}
{\Theta\!\left(\sum\int_{\alpha^{(i)}}^{x^{(i)}}du_h-\int_{\mu}^{\xi}du_h,a\right)
\Theta\!\left(\sum\int_{\alpha^{(i)}}^{z^{(i)}}du_h-\int_{\mu}^{\eta}du_h,a\right)}
}{
\frac{
\Theta\!\left(\sum\int_{\gamma^{(i)}}^{x^{(i)}}dv_h-\int_{\mu}^{\eta}dv_h,b\right)
\Theta\!\left(\sum\int_{\gamma^{(i)}}^{z^{(i)}}dv_h-\int_{\mu}^{\xi}dv_h,b\right)}
{\Theta\!\left(\sum\int_{\alpha^{(i)}}^{x^{(i)}}du_h-\int_{\mu}^{\eta}du_h,a\right)
\Theta\!\left(\sum\int_{\alpha^{(i)}}^{z^{(i)}}du_h-\int_{\mu}^{\xi}du_h,a\right)}
}
={}&-\frac{1}{4\pi^2}
\sum_{h=1}^{h=p}\sum_{m=1}^{m=p}
c_{hm}
\left(\sum\int_{z^{(i)}}^{x^{(i)}}dv_h\right)
\int_{\eta}^{\xi}dv_m .
\end{aligned}
$}
\]
Nun ist offenbar identisch
\[
\begin{aligned}
&\sum\sum c_{hm}\left(\sum\int_{z^{(i)}}^{x^{(i)}}dv_h\right)\int_{\eta}^{\xi}dv_m\\
={}&\frac12\sum\sum c_{hm}
\left(\sum\int_{\gamma^{(i)}}^{x^{(i)}}dv_h-\int_{\mu}^{\eta}dv_h\right)
\left(\sum\int_{\gamma^{(i)}}^{x^{(i)}}dv_m-\int_{\mu}^{\eta}dv_m\right)\\
&+\frac12\sum\sum c_{hm}
\left(\sum\int_{\gamma^{(i)}}^{z^{(i)}}dv_h-\int_{\mu}^{\xi}dv_h\right)
\left(\sum\int_{\gamma^{(i)}}^{z^{(i)}}dv_m-\int_{\mu}^{\xi}dv_m\right)\\
&-\frac12\sum\sum c_{hm}
\left(\sum\int_{\gamma^{(i)}}^{x^{(i)}}dv_h-\int_{\mu}^{\xi}dv_h\right)
\left(\sum\int_{\gamma^{(i)}}^{x^{(i)}}dv_m-\int_{\mu}^{\xi}dv_m\right)\\
&-\frac12\sum\sum c_{hm}
\left(\sum\int_{\gamma^{(i)}}^{z^{(i)}}dv_h-\int_{\mu}^{\eta}dv_h\right)
\left(\sum\int_{\gamma^{(i)}}^{z^{(i)}}dv_m-\int_{\mu}^{\eta}dv_m\right).
\end{aligned}
\]
Setzt man also der Kürze wegen:
\[
\resizebox{0.98\linewidth}{!}{$
\begin{aligned}
&\log
\frac{
\Theta\!\left(\sum\int_{\gamma^{(i)}}^{x^{(i)}}dv_h-\int_{\mu}^{\xi}dv_h,b\right)}
{\Theta\!\left(\sum\int_{\alpha^{(i)}}^{x^{(i)}}du_h-\int_{\mu}^{\xi}du_h,a\right)}
-\frac{1}{8\pi^2}\sum\sum c_{hm}
\left(\sum\int_{\gamma^{(i)}}^{x^{(i)}}dv_h-\int_{\mu}^{\xi}dv_h\right)
\left(\sum\int_{\gamma^{(i)}}^{x^{(i)}}dv_m-\int_{\mu}^{\xi}dv_m\right)\\
&\hspace{8em}=
\varphi\!\left(\sum\int_{\alpha^{(i)}}^{x^{(i)}}du-\int_{\mu}^{\xi}du\right),
\end{aligned}
$}
\]
so hat man zwischen den vier von nur je \(p\) Argumenten abhängigen Functionen \(\varphi\) die Gleichung:
\[
\begin{aligned}
&\varphi\!\left(\sum\int_{\alpha^{(i)}}^{x^{(i)}}du-\int_{\mu}^{\xi}du\right)
+\varphi\!\left(\sum\int_{\alpha^{(i)}}^{z^{(i)}}du-\int_{\mu}^{\eta}du\right)\\
&-\varphi\!\left(\sum\int_{\alpha^{(i)}}^{x^{(i)}}du-\int_{\mu}^{\eta}du\right)
-\varphi\!\left(\sum\int_{\alpha^{(i)}}^{z^{(i)}}du-\int_{\mu}^{\xi}du\right)=0.
\end{aligned}
\]
Diese Gleichung ist offenbar (wie man übrigens durch Differentiation nach den Argumenten leicht beweist) nur erfüllbar, wenn die Function \(\varphi\) eine \emph{lineare Function} ihrer Argumente ist. Setzen wir also
\[
\left.
\begin{aligned}
\sum\int_{\alpha^{(i)}}^{x^{(i)}}du_h-\int_{\mu}^{\xi}du_h&=\omega_h,\\
\sum\int_{\gamma^{(i)}}^{x^{(i)}}dv_h-\int_{\mu}^{\xi}dv_h&=w_h,
\end{aligned}
\right\}
\qquad (h=1,2\ldots p),
\]
so hat man
\[
\varphi(\omega)=\log C+c_1\omega_1+c_2\omega_2+\cdots+c_p\omega_p,
\]
wobei die \(c,C\) weiter zu bestimmende Constante bedeuten, und zwischen den beiden Arten von \(\Theta\)functionen besteht die Gleichung:
\[
(10)\quad
e^{-\frac{1}{8\pi^2}\sum\sum c_{hm}w_hw_m}
\Theta(w,b)=C\cdot e^{\sum c_h\omega_h}\cdot\Theta(\omega,a).
\]
Zwischen den \(\omega,w\) bestehen, abgesehen von additiven halben Periodicitätsmoduln, welche, um die Grenzen gleich zu machen, hinzugefügt werden müssen, dieselben Relationen, welche die \(u\) mit den \(v\) verbanden, nämlich:
\[
(11)\quad
\omega_k-\frac12P_k
=
\frac{1}{2\pi\sqrt{-1}}
\left(u_k^{(1)}w_1+u_k^{(3)}w_2+\cdots+u_k^{(2p-1)}w_p\right);
\]
die halben Periodicitätsmoduln aber dienen zur Definition der \(\gamma\) aus den \(\alpha\), indem man hat:
\[
(12)\quad
\sum\int_{\alpha^{(i)}}^{\gamma^{(i)}}du_k
=\frac12P_k
=\varrho_k\pi\sqrt{-1}
+\frac12(a_{1k}\sigma_1+a_{2k}\sigma_2+\cdots+a_{pk}\sigma_p).
\]
In der Relation (10) bleiben nun noch die Constanten \(c,C\) zu bestimmen, und die ganzen Zahlen \(\varrho,\sigma\).

Um diese Bestimmungen auszuführen, lassen wir sämmtliche \(w_h\) in \(-w_h\) übergehen, wodurch denn \(\omega_k\) in
\[
\frac12P_k-\left(\omega_k-\frac12P_k\right)=P_k-\omega_k
\]
übergeht. Durch diesen Process bleibt die linke Seite der Gleichung (10) ungeändert; die rechte also muss ebenfalls ungeändert bleiben, und man hat daher:
\[
\begin{aligned}
e^{\sum c_h\omega_h}\Theta(\omega,a)
&=e^{\sum c_h(P_h-\omega_h)}\Theta(P-\omega,a)\\
&=e^{\sum c_h(P_h-\omega_h)-\sum\sigma_h\omega_h+\frac12\sum\sum\sigma_h\sigma_ka_{hk}}\Theta(\omega,a).
\end{aligned}
\]
Die Vergleichung der Exponenten giebt also zunächst
\[
c_h=-\frac12\sigma_h;
\]
sodann aber
\[
1=(-1)^{\sum\sigma_h\varrho_h}.
\]
Die Bestimmung der \(c\) ist also auf die der \(\sigma\) zurückgeführt; die letzte Gleichung lehrt, dass \(\sum\sigma_h\varrho_h\) stets eine gerade Zahl ist, was für die Berührungscurven den Satz giebt:

\begin{center}
\textbf{Bei der Ueberführung von \(\Theta\)functionen in andere durch kanonische Substitutionen erscheinen als untere Grenzpuncte stets Berührungspuncte eigentlicher Berührungscurven \((n-2)^{\mathrm{ter}}\) Ordnung (\(\S\,75\)).}
\end{center}

\subsection*{\S\,89. Berührungscurven, welche in den untern Grenzen auftreten.}

Die Bestimmung der \(\varrho,\sigma\) selbst findet man durch Untersuchung der Periodicitätseigenschaften beider Seiten von (10). Wenn wir das System der \(w_h\) um die Periodicitätsmoduln
\[
Q_h=2m_h\pi\sqrt{-1}+b_{1h}q_1+b_{2h}q_2+\cdots
\]
wachsen lassen, so wachsen die \(\omega_h\) um die entsprechenden Periodicitätsmoduln
\[
\begin{aligned}
R_h&=m_1u_h^{(1)}+m_2u_h^{(3)}+m_3u_h^{(5)}+\cdots
+q_1u_h^{(2)}+q_2u_h^{(4)}+\cdots\\
&=2\mu_h\pi\sqrt{-1}+a_{1h}\nu_1+a_{2h}\nu_2+\cdots,
\end{aligned}
\]
wobei die Zahlen \(\mu,\nu\) mit Rücksicht auf die Gleichungen (1) die Werthe erhalten:
\[
\begin{aligned}
\mu_h={}&m_1\beta_{2h-1,1}+m_2\beta_{2h-1,3}+m_3\beta_{2h-1,5}+\cdots\\
&+q_1\beta_{2h-1,2}+q_2\beta_{2h-1,4}+q_3\beta_{2h-1,6}+\cdots,\\
\nu_h={}&m_1\beta_{2h,1}+m_2\beta_{2h,3}+m_3\beta_{2h,5}+\cdots\\
&+q_1\beta_{2h,2}+q_2\beta_{2h,4}+q_3\beta_{2h,6}+\cdots .
\end{aligned}
\tag{13}
\]
Setzen wir nun in (10) gleichzeitig \(w_h+Q_h\) an Stelle von \(w_h\), und \(\omega_h+R_h\) an Stelle von \(\omega_h\), und vergleichen wir die Aenderungen, welche beide Seiten der Gleichung (10) dadurch erfahren, so erhalten wir:
\[
\resizebox{0.98\linewidth}{!}{$
\begin{aligned}
e^{-\frac{1}{4\pi^2}\sum\sum c_{hk}w_hQ_k
-\frac{1}{8\pi^2}\sum\sum c_{hk}Q_hQ_k
+\sum w_hq_h+\frac12\sum\sum q_hq_kb_{hk}}
=
e^{-\frac12\sum\sigma_hR_h+\sum\nu_k\omega_k+\frac12\sum\sum\nu_h\nu_ka_{hk}},
\end{aligned}
$}
\]
oder wenn wir \(\omega_h\) aus (11) durch die \(w\) ausdrücken:
\[
=e^{-\frac12\sum\sigma_hR_h+\frac12\sum\nu_hP_h
+\frac{1}{2\pi\sqrt{-1}}\sum\sum\nu_kw_hu_k^{(2h-1)}
+\frac12\sum\sum\nu_h\nu_ka_{hk}}.
\]
Wir setzen in dieser Gleichung \(w_h=-\frac12Q_h\), und finden
\[
\begin{aligned}
(-1)^{\sum q_hm_h}
&=(-1)^{\sum(\nu_h\varrho_h-\sigma_h\mu_h)}
e^{-\frac{1}{4\pi\sqrt{-1}}\sum\sum\nu_kQ_hu_k^{(2h-1)}
+\frac12\sum\sum\nu_h\nu_ka_{hk}}.
\end{aligned}
\]
Bemerken wir nun, dass, wenn wie in (7) die \(v_h\) um die Periodicitätsmoduln \(Q_h\) wachsen, \(u_k\) um den entsprechenden Periodicitätsmodul \(R_k\) zunehmen muss. Man hat also
\[
R_k=\frac{1}{2\pi\sqrt{-1}}\sum_h Q_hu_k^{(2h-1)},
\]
und die obige Gleichung geht über in:
\[
\begin{aligned}
(-1)^{\sum q_hm_h}
&=(-1)^{\sum(\nu_h\varrho_h-\sigma_h\mu_h)}
e^{-\frac12\sum\nu_kR_k+\frac12\sum\sum\nu_h\nu_ka_{hk}}\\
&=(-1)^{\sum(\nu_h\varrho_h-\sigma_h\mu_h-\nu_h\mu_h)}.
\end{aligned}
\]
Man muss also, und zwar für alle ganzzahligen Werthsysteme der \(m,q\), die Gleichung erfüllt haben:
\[
(14)\quad
\sum q_hm_h-\sum\nu_h\mu_h+\sum\nu_h\varrho_h-\sum\sigma_h\mu_h
\equiv0 \pmod{2}.
\]
Setzen wir in dieser Gleichung
\[
m_k=\beta_{2h,2k},\qquad q_k=-\beta_{2h,2k-1}\qquad (k=1,2\ldots p),
\]
so lehren die Gleichungen (13), dass wegen der Eigenschaften der Elemente einer kanonischen Determinante alle \(\mu\) und \(\nu\) verschwinden, bis auf \(\mu_h\), welches gleich 1 wird. Die Gleichung (14) liefert also
\[
(15)\quad
\sigma_h\equiv\sum_{k=1}^{k=p}\beta_{2h,2k}\beta_{2h,2k-1}\pmod{2}.
\]
Setzt man dagegen
\[
m_k=-\beta_{2h-1,2k},\qquad q_k=\beta_{2h-1,2k-1}\qquad (k=1,2\ldots p),
\]
so verschwinden nach (13) alle \(\mu,\nu\) bis auf \(\nu_h\), welches gleich 1 wird, und (14) ergiebt:
\[
(16)\quad
\varrho_h\equiv\sum_{k=1}^{k=p}\beta_{2h-1,2k}\beta_{2h-1,2k-1}\pmod{2}.
\]
Die Zahlen \(\varrho,\sigma\) sind hienach so weit bestimmt, als die Natur der Sache es überhaupt zulässt; denn wenn man sie um gerade Zahlen vermehrt, ändert die rechte Seite der Gleichung (10) sich nur um eine Constante; die Formel bleibt also ungeändert, und nur die Constante \(C\) erhält eine andere Bedeutung. In der That, ändern wir, indem wir die \(w\) uns unverändert denken, \(\varrho_h,\sigma_h\) um \(2r_h,2s_h\), so erhält \(\omega_k\) den Zuwachs
\[
M_k=2r_k\pi\sqrt{-1}+a_{1k}s_1+a_{2k}s_2+\cdots .
\]
Daher verwandelt sich
\[
e^{-\frac12\sum\sigma_h\omega_h}\Theta(\omega,a)
\]
in
\[
e^{-\frac12\sum(\sigma_h+2s_h)(\omega_h+M_h)}
e^{\sum s_h\omega_h+\frac12\sum\sum s_hs_ka_{hk}}\Theta(\omega,a),
\]
was wirklich von dem vorhergehenden Ausdruck nur um eine Constante verschieden ist.

Der Zusammenhang der beiden, aus verschiedenen Systemen normaler Periodicitätsmoduln entspringenden \(\Theta\)functionen ist also schliesslich durch die Gleichung ausgedrückt:
\[
(17)\quad
e^{-\frac{1}{8\pi^2}\sum\sum c_{hm}w_hw_m}\Theta(w,b)
=
C\,e^{-\frac12\sum\sigma_h\omega_h}\Theta(\omega,a),
\]
wobei die \(w,\omega\) durch die Gleichungen
\[
\omega_k-\frac12P_k=
\frac{1}{2\pi\sqrt{-1}}
\left(u_k^{(1)}w_1+u_k^{(3)}w_2+\cdots+u_k^{(2p-1)}w_p\right),
\]
\[
P_k=\varrho_k\pi\sqrt{-1}+\frac12(a_{1k}\sigma_1+a_{2k}\sigma_2+\cdots+a_{pk}\sigma_p)
\]
verbunden sind, und die \(\varrho,\sigma\) durch die Gleichungen (15), (16) vermittelst der Elemente der kanonischen Determinante definirt werden. Man kann voraussetzen, um sie vollständig zu bestimmen, dass sie nur die Werthe 0 oder 1 annehmen sollen.

Es bleibt übrig zu zeigen, dass die Zahlen \(\sigma,\varrho\), wie sie aus (15), (16) gegeben sind, wirklich unabhängig von einander gerade oder ungerade sein können, mit der einzigen Einschränkung
\[
\sum\sigma_h\varrho_h\equiv0\pmod{2};
\]
mit andern Worten, zu zeigen, dass bei verschiedener Wahl der normalen Periodicitätsmoduln in den untern Grenzen sämmtliche eigentlichen Berührungscurven \(\alpha,\gamma\) auftreten können.

Dies geschieht einfach, indem man kanonische Determinanten aufstellt, welche wirklich alle Berührungscurven liefern. Betrachten wir zunächst die kanonische Determinante:
\[
\left|
\begin{array}{cccccc}
a_1&b_1&0&0&0&0\cdots\\
c_1&d_1&0&0&0&0\cdots\\
0&0&a_2&b_2&0&0\cdots\\
0&0&c_2&d_2&0&0\cdots\\
0&0&0&0&a_3&b_3\cdots\\
0&0&0&0&c_3&d_3\cdots\\
&&\cdots
\end{array}
\right|.
\]
Die Bedingungen, unter welchen sie kanonisch wird, sind:
\[
a_1d_1-b_1c_1=\pm1,\qquad
a_2d_2-b_2c_2=\pm1,\qquad
a_3d_3-b_3c_3=\pm1\quad\text{etc.}
\]
Zugleich hat man:
\[
\varrho_1=a_1c_1,\quad\varrho_2=a_2c_2,\quad\varrho_3=a_3c_3\cdots,
\]
\[
\sigma_1=b_1d_1,\quad\sigma_2=b_2d_2,\quad\sigma_3=b_3d_3\cdots .
\]
Nun sieht man sofort, dass bei gehöriger Bestimmung der \(a,b,c,d\) die \(\varrho,\sigma\) unabhängig von einander gerade oder ungerade sein können; nur die Beschränkung tritt ein, dass keines der Paare \(\varrho_i,\sigma_i\) aus zwei \emph{ungeraden} Zahlen bestehen kann. Dieses letztere kann überhaupt der Gleichung \(\sum\sigma_h\varrho_h\equiv0\pmod{2}\) wegen nur bei einer geraden Anzahl von Paaren eintreten. Sollen \(\sigma_i,\varrho_i\) und \(\sigma_k,\varrho_k\) zwei solche Paare werden, so braucht man nur die Elemente
\[
\begin{array}{ccccc}
a_i&b_i&\cdots&0&0\\
c_i&d_i&\cdots&0&0\\
&&\cdots\\
0&0&\cdots&a_k&b_k\\
0&0&\cdots&c_k&d_k
\end{array}
\]
der obigen Determinante zu ersetzen durch:
\[
\begin{array}{ccccc}
1&1&\cdots&0&-1\\
0&1&\cdots&0&-1\\
&&\cdots\\
1&0&\cdots&1&0\\
1&0&\cdots&1&1
\end{array}.
\]
So oft dies geschieht, hat man zwei Paare \(\sigma,\varrho\) vor sich, welche aus zwei ungeraden Zahlen bestehen; die übrigen Paare aber setzen sich, wie vorher, unabhängig von einander aus geraden und ungeraden, oder aus zwei geraden Zahlen zusammen.

So sind wirklich Determinanten gebildet, welche alle verlangten Fälle liefern, und damit den obigen Satz beweisen. Zugleich sieht man, dass alle kanonischen Determinanten sich in \(2^{p-1}(2^p+1)\) Gruppen theilen, je nachdem die betreffenden Systeme der \(\varrho,\sigma\) auf eine oder die andere Berührungscurve führen.

\subsection*{\S\,90. Reduction der in der Transformationsformel für \(\Theta\) auftretenden Constante auf eine numerische.}

Wir kommen jetzt dazu, den Werth der Constanten \(C\) in der Formel
\[
(1)\quad
e^{-\frac{1}{8\pi^2}\sum\sum c_{hm}w_hw_m+\frac12\sum\sigma_h\omega_h}
\Theta(w,b)
=C\cdot\Theta(\omega,a)
\]
aufzusuchen. Zu diesem Ende bemerken wir zunächst, dass die Function
\[
\Theta(\omega,a)=\sum e^{-\frac12\sum\sum a_{ik}m_im_k+\sum m_i\omega_i}
\]
offenbar den folgenden partiellen Differentialgleichungen genügt\footnote{In diesen Gleichungen werden die \(a_{ik}\) als von einander unabhängig betrachtet. Aber in der That gilt auch die Formel (1), wie man leicht sieht, für den ganz allgemeinen Fall, in welchem die \(a_{ik}\) beliebige Grössen sind, nur mit der Eigenschaft begabt, die \(\Theta\)-Reihe convergiren zu lassen. Man hat, um dies zu beweisen, nur die synektische Function links durch ihre Periodicitätseigenschaften zu definiren, und kann sie dann in eine Reihe entwickeln, welche mit der rechten Seite von (1) übereinstimmt.}:
\[
(2)\quad
\frac{\partial\Theta(\omega,a)}{\partial a_{ii}}
=-\frac12\frac{\partial^2\Theta(\omega,a)}{\partial\omega_i^2},
\qquad
\frac{\partial\Theta(\omega,a)}{\partial a_{ik}}
=-\frac{\partial^2\Theta(\omega,a)}{\partial\omega_i\partial\omega_k}.
\]
Dividiren wir in (1) auf beiden Seiten mit \(C\), so muss auch die linke Seite der Gleichung diesen partiellen Differentialgleichungen genügen; wobei denn die \(c,C\) nur Functionen der \(a_{ik}\) sind, die \(w\) aber sowohl von den \(a_{ik}\) als von den \(\omega_h\) abhängen. Und da die linke Seite der Gleichung (1) eine Summe verschiedenartiger Exponentialgrössen ist, so muss jedes Glied für sich jenen Differentialgleichungen genügen. Betrachten wir insbesondere dasjenige Glied, welches von dem Gliede 1 der Function \(\Theta(w,b)\) herrührt; es ist \(\frac{e^\psi}{C}\), wenn man der Kürze wegen
\[
\psi=-\frac{1}{8\pi^2}\sum\sum c_{hm}w_hw_m+\frac12\sum\sigma_h\omega_h
\]
setzt. Die Function \(\frac{e^\psi}{C}\) muss nun den Gleichungen (2) genügen, oder der einen totalen Differentialgleichung, in welche man jene zusammenfassen kann:
\[
\sum\frac{\partial\frac{e^\psi}{C}}{\partial a_{ii}}\,da_{ii}
+\sum\frac{\partial\frac{e^\psi}{C}}{\partial a_{ik}}\,da_{ik}
=
-\frac12\sum\frac{\partial^2\frac{e^\psi}{C}}{\partial\omega_i^2}\,da_{ii}
-\sum\frac{\partial^2\frac{e^\psi}{C}}{\partial\omega_i\partial\omega_k}\,da_{ik},
\]
wobei in den Summen jede Combination \(ik\) nur einfach, nicht mit dem Factor 2 behaftet, anzusetzen ist. Man kann dafür schreiben:
\[
\begin{aligned}
-\frac{dC}{C}
+\sum\frac{\partial\psi}{\partial a_{ii}}\,da_{ii}
+\sum\frac{\partial\psi}{\partial a_{ik}}\,da_{ik}
={}&-\frac12\sum\frac{\partial^2\psi}{\partial\omega_i^2}\,da_{ii}
-\sum\frac{\partial^2\psi}{\partial\omega_i\partial\omega_k}\,da_{ik}\\
&-\frac12\sum\left(\frac{\partial\psi}{\partial\omega_i}\right)^2da_{ii}
-\sum\frac{\partial\psi}{\partial\omega_i}\frac{\partial\psi}{\partial\omega_k}da_{ik},
\end{aligned}
\]
oder
\[
\begin{aligned}
(3)\quad d\log C={}&
\sum\frac{\partial\psi}{\partial a_{ii}}\,da_{ii}
+\sum\frac{\partial\psi}{\partial a_{ik}}\,da_{ik}
+\frac12\sum\frac{\partial^2\psi}{\partial\omega_i^2}\,da_{ii}\\
&+\sum\frac{\partial^2\psi}{\partial\omega_i\partial\omega_k}\,da_{ik}
+\frac12\sum\left(\frac{\partial\psi}{\partial\omega_i}\right)^2da_{ii}
+\sum\frac{\partial\psi}{\partial\omega_i}\frac{\partial\psi}{\partial\omega_k}da_{ik}.
\end{aligned}
\]
Bildet man nun die Differentialquotienten von \(\psi\) und setzt sodann für die \(\omega_h\) diejenigen constanten Werthe ein, für welche die \(w_h\) verschwinden, so findet man:
\[
(4)\quad
\frac{\partial\psi}{\partial a_{ii}}=0,\qquad
\frac{\partial\psi}{\partial a_{ik}}=0,\qquad
\frac{\partial\psi}{\partial\omega_i}=\frac12\sigma_i,
\]
\[
\frac{\partial^2\psi}{\partial\omega_i\partial\omega_k}
=
-\frac{1}{4\pi^2}\sum\sum c_{hm}
\frac{\partial w_h}{\partial\omega_i}
\frac{\partial w_m}{\partial\omega_k}.
\]
Man hat ferner, indem man den Werth
\[
c_{hm}=\sum_r\beta_{2r,2h-1}u_r^{(2m-1)}
\]
einträgt:
\[
\frac{\partial^2\psi}{\partial\omega_i\partial\omega_k}
=
\sum_h\sum_m\sum_r
\beta_{2r,2h-1}u_r^{(2m-1)}
\frac{\partial}{\partial\omega_i}\frac{w_h}{2\pi\sqrt{-1}}\,
\frac{\partial}{\partial\omega_k}\frac{w_m}{2\pi\sqrt{-1}},
\]
oder mit Rücksicht auf (11) \(\S\,88\):
\[
=
\sum_h\sum_r\beta_{2r,2h-1}
\frac{\partial}{\partial\omega_i}\frac{w_h}{2\pi\sqrt{-1}}\,
\frac{\partial\omega_r}{\partial\omega_k}
=
\sum_h\beta_{2k,2h-1}
\frac{\partial}{\partial\omega_i}\frac{w_h}{2\pi\sqrt{-1}}.
\]
Führen wir endlich die Function \(\Delta\) des \(\S\,87\) ein, so wird
\[
\frac{\partial}{\partial\omega_i}\frac{w_h}{2\pi\sqrt{-1}}
=
\frac{\partial\log\Delta}{\partial u_i^{(2h-1)}},
\]
und da auch
\[
\beta_{2k,2h-1}=\frac{\partial u_i^{(2h-1)}}{\partial a_{ik}},
\]
so wird schliesslich
\[
(5)\quad
\frac{\partial^2\psi}{\partial\omega_i\partial\omega_k}
=
\frac{\partial\log\Delta}{\partial a_{ik}},
\]
wenn man den letzten Differentialquotienten so versteht, dass nach \(a_{ik}\) differenzirt wird, soweit dasselbe in den \(u\) mit unterm Index \(i\) vorkommt. Tragen wir nun die Werthe (4), (5) in (3) ein, so ergiebt sich:
\[
d\log C=\frac12d\log\Delta+\frac18d\!\left(\sum\sum\sigma_i\sigma_ka_{ik}\right),
\]
oder
\[
C=\gamma\cdot\sqrt{\Delta}\cdot e^{\frac18\sum\sum\sigma_i\sigma_ka_{ik}},
\]
wo nun die Constante \(\gamma\) weder von den \(\omega\), noch von den \(a_{ik}\) abhängt, sondern eine rein numerische sein muss. Die Formel für den Zusammenhang der beiden Functionen \(\Theta\) aber verwandelt sich in folgende:
\[
e^{-\frac{1}{8\pi^2}\sum\sum c_{hm}w_hw_m}
\Theta(w,b)
=
\gamma\sqrt{\Delta}\,
e^{-\frac12\sum\sigma_h\omega_h+\frac18\sum\sum\sigma_i\sigma_ka_{ik}}
\Theta(\omega,a).
\]


% --- Source pages 343--355; printed pages 321--333. ---

\subsection*{\S\,91. Die numerische Constante bei Zusammensetzung von Transformationen.}

Die Grösse \(\gamma\) ist eine rein numerische Constante, welche von den Eigenschaften der Zahlen \(\beta\) abhängt. Um sie zu bestimmen, müssen wir auf die Zusammensetzung der kanonischen Determinante aus einfachen zurückgehen. Denken wir uns zwei auf einander folgende Substitutionen von kanonischer Determinante zusammengesetzt, wie dies in \(\S\,85\) ausgeführt wurde. Die eine Substitution, bei welcher wir die Elemente der kanonischen Determinanten durch \(\beta\), und die betreffenden Combinationen derselben durch \(\varrho,\sigma\) bezeichnen, führe von den Variabeln \(\omega\) zu den Variabeln \(v\) mit Hülfe der Gleichungen:
\[
(1)\quad
\omega_k-\frac12P_k
=\frac{1}{2\pi\sqrt{-1}}\left(u_k^{(1)}v_1+u_k^{(3)}v_2+\cdots+u_k^{(2p-1)}v_p\right).
\]
Die zweite Substitution führe von den \(v\) zu neuen Variabeln \(W\). Die Elemente der dabei benutzten kanonischen Determinante bezeichnen wir durch \(\beta'\), ebenso die betreffenden Combinationen derselben durch \(\varrho',\sigma'\); die fundamentalen Periodicitätsmoduln, welche sich hier aus den \(b_{ik}\) und \(2\pi\sqrt{-1}\) ganzzahlig zusammensetzen und welche den \(u_k^{(i)}\) entsprechen, seien durch \(v_k^{(i)}\) dargestellt. Endlich mögen die Perioden \(P\) hier durch Perioden \(P'\) ersetzt werden, welche sich ebenfalls aus den \(b_{ik}\) und aus \(2\pi\sqrt{-1}\) ganzzahlig zusammensetzen. Die Formel für den Zusammenhang der Variabeln ist dann
\[
(2)\quad
v_k-\frac12P_k'
=\frac{1}{2\pi\sqrt{-1}}\left(v_k^{(1)}W_1+v_k^{(3)}W_2+\cdots+v_k^{(2p-1)}W_p\right).
\]
Endlich setzt sich aus diesen beiden eine Substitution zusammen, welche von den \(\omega\) zu den \(W\) führt:
\[
(3)\quad
\omega_k-\frac12P_k''
=\frac{1}{2\pi\sqrt{-1}}\left(w_k^{(1)}W_1+w_k^{(3)}W_2+\cdots+w_k^{(2p-1)}W_p\right).
\]
In dieser Substitution bezeichnen wir die den \(\beta,\varrho,\sigma\) entsprechenden Grössen durch \(\beta'',\varrho'',\sigma''\); die Perioden \(w_k^{(h)},P_k''\) setzen sich hier wieder aus den \(a_{ik}\) und \(2\pi\sqrt{-1}\) zusammen. Setzt man die Werthe der \(v_k\) aus (2) in (1) ein, und vergleicht die entstehende Formel mit (3), so ergeben sich folgende Relationen:
\[
(4)\quad
\left\{\begin{aligned}
2\pi\sqrt{-1}\,w_k^{(1)}&=u_k^{(1)}v_1^{(1)}+u_k^{(3)}v_2^{(1)}+\cdots+u_k^{(2p-1)}v_p^{(1)},\\
2\pi\sqrt{-1}\,w_k^{(3)}&=u_k^{(1)}v_1^{(3)}+u_k^{(3)}v_2^{(3)}+\cdots+u_k^{(2p-1)}v_p^{(3)},\\
&\hspace{4em}\cdots\\
2\pi\sqrt{-1}\,w_k^{(2p-1)}&=u_k^{(1)}v_1^{(2p-1)}+u_k^{(3)}v_2^{(2p-1)}+\cdots+u_k^{(2p-1)}v_p^{(2p-1)}.
\end{aligned}\right.
\]
\[
(5)\quad
2\pi\sqrt{-1}\,(P_k''-P_k)
=u_k^{(1)}P_1'+u_k^{(3)}P_2'+\cdots+u_k^{(2p-1)}P_{2p}'.
\]
Die ersten \(p\) von diesen Gleichungen liefern nichts anderes als einige von den Relationen
\[
(6)\quad
\beta_{k,r}''=\beta_{k1}\beta_{1r}'+\beta_{k2}\beta_{2r}'+\cdots+\beta_{k,2p}\beta_{2p,r}',
\]
vermittelst deren die Elemente der zusammengesetzten kanonischen Determinante gebildet werden. Die letzte Gleichung aber lehrt, dass die Zahlen \(\varrho'',\sigma''\) nicht blos, wie die \(\varrho,\sigma,\varrho',\sigma'\), aus den Elementen \(\beta,\beta'\) bis auf gerade Zahlen bestimmt werden, sondern dass sie sich aus diesen und den \(\varrho,\sigma,\varrho',\sigma'\) auf völlig bestimmte Weise zusammensetzen. Setzt man nämlich für die \(u_k^{(h)},P_k,P_k',P_k''\) ihre Werthe
\[
(7)\quad
u_k^{(h)}=\beta_{2k-1,h}\,2\pi\sqrt{-1}+\sum_m\beta_{2m,h}a_{mk}
\]
\[
(8)\quad
\left\{\begin{aligned}
P_k&=\varrho_k\cdot2\pi\sqrt{-1}+\sigma_1a_{1k}+\sigma_2a_{2k}+\cdots+\sigma_pa_{pk},\\
P_k''&=\varrho_k''\cdot2\pi\sqrt{-1}+\sigma_1''a_{1k}+\sigma_2''a_{2k}+\cdots+\sigma_p''a_{pk},\\
P_k'&=\varrho_k'\cdot2\pi\sqrt{-1}+\sigma_1'b_{1k}+\sigma_2'b_{2k}+\cdots+\sigma_p'b_{pk},
\end{aligned}\right.
\]
ein, so ergiebt sich:
\[
(\varrho_k''-\varrho_k)2\pi\sqrt{-1}+\sum_h(\sigma_h''-\sigma_h)a_{hk}
=\sum_h u_k^{(2h-1)}\varrho_h'
+\frac{1}{2\pi\sqrt{-1}}\sum_h\sum_i u_k^{(2h-1)}b_{hi}\sigma_i'.
\]
Inzwischen wissen wir, dass dem Zuwachs der \(v_k\) um die Periodicitätsmoduln \(b_{ki}\) ein Zuwachs der \(\omega_k\) um \(u_k^{(2i)}\) entspricht, so dass also
\[
(9)\quad
u_k^{(2i)}=\frac{1}{2\pi\sqrt{-1}}\sum_hu_k^{(2h-1)}b_{hi}.
\]
Demnach geht die obige Formel über in:
\[
(\varrho_k''-\varrho_k)2\pi\sqrt{-1}+\sum_h(\sigma_h''-\sigma_h)a_{hk}
=\sum_hu_k^{(2h-1)}\varrho_h'+\sum_hu_k^{(2h)}\sigma_h'.
\]
oder man erhält zur Bestimmung der \(\varrho'',\sigma''\) die Formeln:
\[
(10)\quad
\left\{\begin{aligned}
\varrho_k''-\varrho_k&=\sum_h(\varrho_h'\beta_{2k-1,2h-1}+\sigma_h'\beta_{2k-1,2h}),\\
\sigma_k''-\sigma_k&=\sum_h(\varrho_h'\beta_{2k,2h-1}+\sigma_h'\beta_{2k,2h}).
\end{aligned}\right.
\]
Die drei Formeln für die \(\Theta\)functionen erhalten wir aus der frühern Formel, wenn wir bei der zweiten und dritten Substitution die \(c,a,b,\Delta\) beziehungsweise durch \(c',b,A,\Delta'\) und durch \(c'',a,A,\Delta''\) ersetzen, so dass wir erhalten:
\[
\begin{aligned}
e^{-\frac{1}{8\pi^2}\sum\sum c_{hm}v_hv_m}\Theta(v,b)
&=\gamma\sqrt{\Delta}\,e^{-\frac12\sum\sigma_h\omega_h+\frac18\sum\sum\sigma_i\sigma_ka_{ik}}\Theta(\omega,a),\\
e^{-\frac{1}{8\pi^2}\sum\sum c_{hm}'W_hW_m}\Theta(W,A)
&=\gamma'\sqrt{\Delta'}\,e^{-\frac12\sum\sigma_h'v_h+\frac18\sum\sum\sigma_i'\sigma_k'b_{ik}}\Theta(v,b),\\
e^{-\frac{1}{8\pi^2}\sum\sum c_{hm}''W_hW_m}\Theta(W,A)
&=\gamma''\sqrt{\Delta''}\,e^{-\frac12\sum\sigma_h''\omega_h+\frac18\sum\sum\sigma_i''\sigma_k''a_{ik}}\Theta(\omega,a).
\end{aligned}
\]
Wir dividiren nun das Product der ersten beiden Gleichungen durch die dritte, und setzen zugleich die \(W\) gleich Null, so dass \(\omega_h=\frac12P_h'',\ v_h=\frac12P_h'\) wird. Wir finden dann:
\[
\begin{aligned}
(11)\quad
e^{-\frac{1}{32\pi^2}\sum\sum c_{hm}P_h'P_m'}
={}&\frac{\gamma\gamma'}{\gamma''}\cdot\frac{\sqrt\Delta\sqrt{\Delta'}}{\sqrt{\Delta''}}\\
&\cdot e^{\frac14\sum[(\sigma_h''-\sigma_h)P_h''-\sigma_h'P_h']
-\frac18\sum\sum[(\sigma_i''\sigma_k''-\sigma_i\sigma_k)a_{ik}-\sigma_i'\sigma_k'b_{ik}]}.
\end{aligned}
\]
Da nun \(\Delta\) die Determinante der \(u_k^{(2h-1)}\), \(\Delta'\) die der \(v_k^{(2h-1)}\) und \(\Delta''\) die der \(w_k^{(2h-1)}\) ist, so hat man nach (4)
\[
(12)\quad
\Delta\Delta'=\Delta''\cdot(2\pi\sqrt{-1})^{\frac p4}.
\]
Indem wir nun \((2\pi\sqrt{-1})^{\frac p2}\) in einer ein für alle Mal willkürlich fixirten Bedeutung nehmen, wollen wir das Vorzeichen von \(\sqrt{\Delta''}\) immer so bestimmt denken, dass auch
\[
(13)\quad
\sqrt\Delta\cdot\sqrt{\Delta'}=\sqrt{\Delta''}\cdot(2\pi\sqrt{-1})^{\frac p2}.
\]
Wir wollen ferner in (11), da es sich nur um Bestimmung des rein numerischen Verhältnisses \(\frac{\gamma\gamma'}{\gamma''}\) handelt, die \(a_{ik}\) gleich Null setzen. Sofort reducirt sich dann \(P_h\) auf \(\varrho_h\cdot2\pi\sqrt{-1}\), \(P_h''\) auf \(\varrho_h''\cdot2\pi\sqrt{-1}\), und die \(b_{ik}\) bestimmen sich aus den Gleichungen, in welche die Gleichungen (9) jetzt übergehen:
\[
(14)\quad
\beta_{2k-1,2i}\cdot2\pi\sqrt{-1}=\sum_h\beta_{2k-1,2h-1}b_{ih}.
\]
Wenn man ferner die Werthe der \(c_{hm}\) einführt, so hat man:
\[
\begin{aligned}
(15)\quad
\sum\sum c_{hm}P_h'P_m'
&=2\pi\sqrt{-1}\sum\sum\sum\beta_{2k,2h-1}u_k^{(2m-1)}P_m'P_h'\\
&=2\pi\sqrt{-1}\sum\sum\beta_{2k,2h-1}(P_k''-P_k)P_h',
\end{aligned}
\]
was sich, weil die \(a_{ik}\) verschwinden sollten, auf
\[
\begin{aligned}
&(2\pi\sqrt{-1})^2\sum\sum\beta_{2k,2h-1}(\varrho_k''-\varrho_k)P_h'\\
&=(2\pi\sqrt{-1})^3\sum\sum\beta_{2k,2h-1}(\varrho_k''-\varrho_k)\varrho_h'\\
&\quad+(2\pi\sqrt{-1})^2\sum\sum\sum\beta_{2k,2h-1}(\varrho_k''-\varrho_k)\sigma_m'b_{hm}
\end{aligned}
\]
reducirt. Setzen wir im zweiten Theile dieses Ausdrucks für \(\varrho_k''-\varrho_k\) seinen Werth aus (10), so wird
\[
\begin{aligned}
(16)\quad
\sum_k\sum_h\beta_{2k,2h-1}(\varrho_k''-\varrho_k)b_{hm}
={}&\sum_k\sum_h\sum_i\beta_{2k,2h-1}\bigl(\varrho_i'\beta_{2k-1,2i-1}\nobreak+\sigma_i'\beta_{2k-1,2i}\bigr)b_{hm}.
\end{aligned}
\]
Nun ist nach den zwischen den Elementen der kanonischen Determinante bestehenden Gleichungen:
\[
(17)\quad
\left\{\begin{aligned}
\sum_k\beta_{2k,2h-1}\beta_{2k-1,2i-1}
&=\sum_k\beta_{2k-1,2h-1}\beta_{2k,2i-1},\\
\sum_k\beta_{2k,2h-1}\beta_{2k-1,2i}
&=\sum_k\beta_{2k-1,2h-1}\beta_{2k,2i},
\end{aligned}\right.
\]
und nur wenn \(h=i\), ist die letztere Gleichung zu ersetzen durch
\[
(18)\quad
\sum_k\beta_{2k,2i-1}\beta_{2k-1,2i}
=-1+\sum_k\beta_{2k-1,2i-1}\beta_{2k,2i}.
\]
Der Ausdruck (16) verwandelt sich daher in:
\[
\sum_k\sum_h\sum_i\beta_{2k-1,2h-1}(\varrho_i'\beta_{2k,2i-1}+\sigma_i'\beta_{2k,2i})b_{hm}
-
\sum_h\sigma_h'b_{hm},
\]
und wenn man nach \(h\) summirt, so erhält man nach (14):
\[
2\pi\sqrt{-1}\sum_k\sum_i\sum_m\beta_{2k-1,2m}(\varrho_i'\beta_{2k,2i-1}+\sigma_i'\beta_{2k,2i})\sigma_m'
-
\sum_h\sum_m\sigma_h'\sigma_m'b_{hm}.
\]
Der Ausdruck (15) wird also
\[
\begin{aligned}
&(2\pi\sqrt{-1})^3\sum_k\sum_h\beta_{2k,2h-1}(\varrho_k''-\varrho_k)\varrho_h'\\
&\quad+(2\pi\sqrt{-1})^3\sum_k\sum_i\sum_m\beta_{2k-1,2m}(\varrho_i'\beta_{2k,2i-1}+\sigma_i'\beta_{2k,2i})\sigma_m'\\
&\quad-(2\pi\sqrt{-1})^2\sum_h\sum_m\sigma_h'\sigma_m'b_{hm}.
\end{aligned}
\]
Fasst man jetzt alles zusammen, so ergibt sich endlich aus (11) die Formel:
\[
\begin{aligned}
(19)\quad
\frac{\gamma\gamma'}{\gamma''}
={}&\frac{1}{(2\pi\sqrt{-1})^{p/2}}
\exp\!\left(\frac{2\pi\sqrt{-1}}{8}\sum_k\sum_h\beta_{2k,2h-1}(\varrho_k''-\varrho_k)\varrho_h'\right)\\
&\cdot\exp\!\left(\frac{2\pi\sqrt{-1}}{8}\sum_k\sum_m\sum_i\beta_{2k-1,2m}\sigma_m'(\varrho_i'\beta_{2k,2i-1}+\sigma_i'\beta_{2k,2i})\right)\\
&\cdot\exp\!\left(-\frac{2\pi\sqrt{-1}}{4}\bigl[(\sigma_h''-\sigma_h)\varrho_h''-\sigma_h'\varrho_h'\bigr]\right).
\end{aligned}
\]
Abgesehen also von dem Factor \((2\pi\sqrt{-1})^{p/2}\), welchen man auch als integrirenden Theil jedes Factors \(\gamma\) betrachten kann, unterscheidet sich \(\gamma''\) von dem Producte \(\gamma\gamma'\) nur durch eine achte Wurzel der Einheit. Setzen wir
\[
\gamma=\delta(2\pi\sqrt{-1})^{-p/2},\qquad
\gamma'=\delta'(2\pi\sqrt{-1})^{-p/2},\qquad
\gamma''=\delta''(2\pi\sqrt{-1})^{-p/2},
\]
so sind die \(\delta\) verbunden durch die Gleichung:
\[
\begin{aligned}
(20)\quad
\frac{\delta\delta'}{\delta''}
={}&\exp\!\left(\frac{2\pi\sqrt{-1}}{8}\sum\sum\beta_{2k,2h-1}(\varrho_k''-\varrho_k)\varrho_h'\right)\\
&\cdot\exp\!\left(
\frac{2\pi\sqrt{-1}}{8}\sum\sum\sum\beta_{2k-1,2m}\sigma_m'(\varrho_i'\beta_{2k,2i-1}+\sigma_i'\beta_{2k,2i})
\right.\\
&\qquad\left.
-\frac{2\pi\sqrt{-1}}{4}\bigl[(\sigma_h''-\sigma_h)\varrho_h''-\sigma_h'\varrho_h'\bigr]\right),
\end{aligned}
\]
und die Formel für die \(\Theta\) wird:
\[
\begin{aligned}
(21)\quad
e^{-\frac{1}{8\pi^2}\sum\sum c_{hm}v_hv_m}\Theta(v,b)
={}&\delta\sqrt{\frac{\Delta}{(2\pi\sqrt{-1})^p}}\,
 e^{-\frac12\sum\sigma_h\omega_h+\frac18\sum\sum\sigma_i\sigma_ka_{ik}}\Theta(\omega,a).
\end{aligned}
\]
Die Formel (20) lässt sich übrigens noch wesentlich einfacher darstellen. Man hat nach (10):
\[
\sum\sum\beta_{2k,2h-1}(\varrho_k''-\varrho_k)\varrho_h'
=\sum\sum\sum\beta_{2k,2h-1}\varrho_i'(\beta_{2k-1,2i-1}\varrho_i'+\beta_{2k-1,2i}\sigma_i').
\]
Wendet man nun auf die nach \(k\) zu nehmenden Summen die Formeln (17) an, so wird hieraus:
\[
\sum\sum\sum\beta_{2k-1,2h-1}\varrho_h'(\beta_{2k,2i-1}\varrho_i'+\beta_{2k,2i}\sigma_i'),
\]
und der in (20) mit \(\frac{2\pi\sqrt{-1}}8\) multiplicirte Ausdruck ist daher:
\[
\begin{aligned}
&\sum_k\left[\sum_h(\beta_{2k-1,2h-1}\varrho_h'+\beta_{2k-1,2h}\sigma_h')\right]
\left[\sum_i(\beta_{2k,2i-1}\varrho_i'+\beta_{2k,2i}\sigma_i')\right]\\
&=\sum_k(\varrho_k''-\varrho_k)(\sigma_k''-\sigma_k),
\end{aligned}
\]
so dass die Formel (20) in die Gestalt übergeht:
\[
(22)\quad
\frac{\delta\delta'}{\delta''}
=\exp\!\left(\frac{2\pi\sqrt{-1}}8\sum_k[2\sigma_k'\varrho_k'-(\varrho_k''+\varrho_k)(\sigma_k''-\sigma_k)]\right).
\]
Man vereinfacht endlich diese Formel noch, indem man
\[
\delta=\varepsilon e^{\frac{2\pi\sqrt{-1}}4\sum_k\varrho_k\sigma_k},
\qquad
\delta'=\varepsilon' e^{\frac{2\pi\sqrt{-1}}4\sum_k\varrho_k'\sigma_k'},
\qquad
\delta''=\varepsilon'' e^{\frac{2\pi\sqrt{-1}}4\sum_k\varrho_k''\sigma_k''}
\]
setzt. Die Formel für die \(\Theta\)functionen wird dann:
\[
\begin{aligned}
(23)\quad
e^{-\frac{1}{8\pi^2}\sum\sum c_{hm}v_hv_m}\Theta(v,b)
={}&\varepsilon\sqrt{\frac{\Delta}{(2\pi\sqrt{-1})^p}}
 e^{\begin{subarray}{l}\frac{2\pi\sqrt{-1}}4\sum\varrho_k\sigma_k-\frac12\sum\sigma_h\omega_h\\[0.1em]
 +\frac18\sum\sum\sigma_i\sigma_ka_{ik}\end{subarray}}\Theta(\omega,a),
\end{aligned}
\]
und die Formel für den Zusammenhang der Zahlen \(\varepsilon\) ist:
\[
(24)\quad
\frac{\varepsilon\varepsilon'}{\varepsilon''}
=\exp\!\left(\frac{2\pi\sqrt{-1}}8\sum_k(\varrho_k''-\varrho_k)(\sigma_k''+\sigma_k)\right).
\]

\subsection*{\S\,92. Bestimmung der numerischen Constante für die einfachen Transformationen.}

Es bleibt jetzt nur noch übrig, die allgemeine Formel zwischen den verschiedenen \(\Theta\) für den Fall aufzustellen, in welchem die kanonische Determinante eine \emph{einfache} ist, und für diesen Fall die Coefficienten \(\varepsilon\) wirklich zu bestimmen. Da man jede Substitution aus einfachen zusammensetzen kann, so giebt die Formel (24) des Vorigen das Mittel, die \(\varepsilon\) für jede Substitution daraus zu bestimmen. Betrachten wir der Reihe nach die verschiedenen einfachen Substitutionen, wobei wir der Einfachheit der Darstellung wegen die nicht verschwindenden Elemente der einfachen Determinante an bestimmte Stellen setzen. Es ist sofort ersichtlich, was sich ändert, wenn man die \(\pm1\) an andre Stellen setzt.

\begin{center}I.\end{center}

Determinante der Substitution:
\[
\left|\begin{array}{ccccc}
0&\pm1&0&0&\cdots\\
\mp1&0&0&0&\cdots\\
0&0&1&0&\cdots\\
0&0&0&1&\cdots\\
\cdots&\cdots&\cdots&\cdots&
\end{array}\right|.
\]
Alle \(\varrho,\sigma\) verschwinden. Die \(v\) werden durch die Formeln gegeben:
\[
\begin{gathered}
v_1=\pm\frac{2\pi\sqrt{-1}}{a_{11}}\omega_1,\qquad
v_2=\frac{a_{11}\omega_2-a_{21}\omega_1}{a_{11}},\qquad
v_3=\frac{a_{11}\omega_3-a_{31}\omega_1}{a_{11}},\\
\cdots\qquad
v_p=\frac{a_{11}\omega_p-a_{p1}\omega_1}{a_{11}}.
\end{gathered}
\]
Ferner wird:
\[
\begin{gathered}
b_{11}=\frac{4\pi^2}{a_{11}},
\quad b_{12}=\pm2\pi\sqrt{-1}\frac{a_{12}}{a_{11}},
\quad b_{13}=\pm2\pi\sqrt{-1}\frac{a_{13}}{a_{11}},\ldots,
\quad b_{1p}=\pm2\pi\sqrt{-1}\frac{a_{1p}}{a_{11}},\\
b_{ik}=\frac{a_{ik}a_{11}-a_{1i}a_{1k}}{a_{11}},\qquad (i,k=2,3\ldots p),\\
\Delta=\pm(2\pi\sqrt{-1})^{p-1}a_{11},
\qquad c_{11}=a_{11},\quad\hbox{alle andern }c_{ik}\hbox{ verschwinden.}
\end{gathered}
\]
Die Formel für die Transformation der \(\Theta\)function wird also:
\[
e^{-\frac{a_{11}}{8\pi^2}v_1^2}\Theta(v,b)
=\varepsilon\sqrt{\frac{\pm a_{11}}{2\pi\sqrt{-1}}}\Theta(\omega,a).
\]
Die Bestimmung von \(\varepsilon\) erfolgt dadurch, dass man die Formel direct ableitet. Es ist nämlich
\[
\Theta(\omega,a)=\sum e^{-\frac12\sum\sum a_{ik}m_im_k+\sum m_i\omega_i}
\]
\[
=\sum e^{-\frac12\sum{}'\sum{}'a_{ik}m_im_k+\sum{}'m_i\omega_i}
\Theta(\omega_1-a_{12}m_2-a_{13}m_3\ldots,a_{11}),
\]
wenn das Zeichen \(!\) andeutet, dass die Summation sich auf die Indices \(2,3\ldots p\) beschränken soll, und wenn die \(\Theta\)function unter dem Summenzeichen die elliptische ist, für welche die Definition gilt:
\[
\Theta(\omega_1,a_{11})=\sum e^{-\frac12a_{11}m_1^2+m_1\omega_1}.
\]
Nach Jacobi aber ist dann zugleich
\[
\Theta(\omega_1,a_{11})=\sqrt{\frac{2\pi}{a_{11}}}
 e^{\frac{\omega_1^2}{2a_{11}}}
\Theta\left(\pm\frac{2\pi\sqrt{-1}}{a_{11}}\omega_1,\frac{4\pi^2}{a_{11}}\right).
\]
Was das Zeichen der Wurzel angeht, so kann man \(\omega_1=0\), \(a_{11}=2\pi\) setzen, und sieht, dass in diesem Falle \(\sqrt{\frac{2\pi}{a_{11}}}\) in 1 übergehen muss. Da nun der reelle Theil von \(a_{11}\) stets positiv sein muss, damit die \(\Theta\)function überhaupt einen Sinn habe, so schliesst man, dass auch \(\sqrt{\frac{2\pi}{a_{11}}}\), dessen reeller Theil demnach ebenfalls nie Null werden kann, immer so zu wählen ist, dass sein reeller Theil, wie in dem speciellen Fall von 1, positiv ist. Indem man nun statt \(\omega_1\) in obiger Formel
\[
\omega_1-m_2a_{21}-m_3a_{31}\ldots
\]
setzt und dann den transformirten Werth in \(\Theta(\omega,a)\) einführt, erhält man:
\[
\Theta(\omega,a)=\sqrt{\frac{2\pi}{a_{11}}}e^{\frac{\omega_1^2}{2a_{11}}}\sum e^{-\frac12\sum\sum b_{ik}m_im_k+\sum m_iv_i},
\]
wo die \(v,b\) genau mit den oben gegebenen Definitionen übereinstimmen, und wo auch
\[
e^{\frac{\omega_1^2}{2a_{11}}}=e^{-\frac{a_{11}}{8\pi^2}v_1^2},
\]
so dass die Uebereinstimmung der Formel mit der aus unserer allgemeinen Theorie abgeleiteten vollständig wird, wenn man
\[
\varepsilon\sqrt{\frac{\pm a_{11}}{2\pi\sqrt{-1}}}=\sqrt{\frac{a_{11}}{2\pi}}
\]
setzt, und die Wurzel \(\sqrt{\frac{a_{11}}{2\pi}}\) so nimmt, dass der reelle Theil positiv wird.

\begin{center}II.\end{center}

Determinante der Substitution:
\[
\left|\begin{array}{ccccc}
1&0&0&0&\cdots\\
\pm1&1&0&0&\cdots\\
0&0&1&0&\cdots\\
0&0&0&1&\cdots\\
\cdots&\cdots&\cdots&\cdots&
\end{array}\right|.
\]
Die \(\varrho,\sigma\) verschwinden bis auf \(\varrho_1\), welches congruent \(\pm1\) wird, und also gleich 1 gesetzt werden kann. Man hat daher
\[
P_1=\pm2\pi\sqrt{-1},\qquad P_2=0,
\qquad P_3=0\ldots
\]
Die Variabeln hängen durch die Formeln zusammen:
\[
v_1=\omega_1\mp\pi\sqrt{-1},\qquad
v_2=\omega_2,
\qquad v_3=\omega_3\ldots
\]
Ausserdem hat man:
\[
b_{11}=a_{11}\pm2\pi\sqrt{-1},\quad
\hbox{in allen andern Fällen }b_{ik}=a_{ik},
\]
\[
c_{ik}=0\quad(\hbox{für alle Werthe der }i,k),
\qquad\Delta=(2\pi\sqrt{-1})^p.
\]
Die Formel für die Transformation der \(\Theta\)function wird also:
\[
\Theta(v,b)=\varepsilon\Theta(\omega,a),
\]
oder:
\[
\sum e^{-\frac12\sum\sum a_{ik}m_im_k-m_1^2\pi\sqrt{-1}+\sum m_i\omega_i\mp m_1\pi\sqrt{-1}}
=\varepsilon\sum e^{-\frac12\sum\sum a_{ik}m_im_k+\sum m_i\omega_i}.
\]
Da \(e^{-m_1(m_1\pm1)\pi\sqrt{-1}}=1\), so stimmen die \(\Theta\)functionen beiderseits völlig überein, und man hat also
\[
\varepsilon=1.
\]

\begin{center}III.\end{center}

Determinante der Substitution:
\[
\left|\begin{array}{cccccc}
1&0&\pm1&0&0&\cdots\\
0&1&0&0&0&\cdots\\
0&0&1&0&0&\cdots\\
0&\mp1&0&1&0&\cdots\\
0&0&0&0&1&\cdots\\
\cdots&\cdots&\cdots&\cdots&\cdots&
\end{array}\right|.
\]
Alle \(\varrho,\sigma\) verschwinden. Man hat ferner
\[
v_i=\omega_i,
\]
mit alleiniger Ausnahme von \(v_2\), für welches die Formel gilt:
\[
v_2=\omega_2\mp\omega_1.
\]
Alle \(c_{ik}\) verschwinden. Die \(b_{ik}\) erhält man aus den Formeln:
\[
\begin{gathered}
b_{11}=a_{11},\qquad b_{12}=a_{12}\mp a_{11},\qquad b_{22}=a_{22}\mp a_{12}+a_{11},\\
b_{2k}=a_{2k}\mp a_{1k}\qquad(k=3,4\ldots),
\end{gathered}
\]
in allen andern Fällen ist \(b_{ik}=a_{ik}\). Endlich \(\Delta=(2\pi\sqrt{-1})^p\). Die Transformationsformel für \(\Theta\) wird also:
\[
\Theta(v,b)=\varepsilon\Theta(\omega,a),
\]
oder:
\[
\begin{aligned}
&\sum\sum e^{-\frac12\sum\sum a_{ik}m_im_k\pm(a_{11}m_1+a_{12}m_2+a_{13}m_3+\cdots)m_2}
 e^{-\frac12a_{11}m_2^2+\sum m_i\omega_i\mp m_2\omega_1}\\
&\qquad=\varepsilon\sum\sum e^{-\frac12\sum\sum a_{ik}m_im_k+\sum m_i\omega_i}.
\end{aligned}
\]
Hier stimmen wieder beide Summen völlig überein, wenn nur in der Summe links \(m_1\pm m_2\) an Stelle von \(m_1\) gesetzt wird, was ohne Aenderung des Werthes der Summe geschehen kann. Man hat also abermals
\[
\varepsilon=1.
\]
Wir wollen noch angeben, welche Gestalt die Reductionsformel für \(\frac{\varepsilon\varepsilon'}{\varepsilon''}\) annimmt, wenn bei der Zusammensetzung der Substitutionen eine beliebige \((\beta)\) mit einer der einfachen \((\beta')\) zusammengesetzt wird. Wir haben dann die Mittel in der Hand, für irgend eine kanonische Substitution die Zahl \(\varepsilon\) zu bestimmen, sobald man die Zerlegung der kanonischen Determinante in einfache gefunden, und also umgekehrt jene aus diesen zusammensetzt. Die Formeln werden:

Zusammensetzung einer beliebigen Determinante mit I.:
\[
\varepsilon''=\varepsilon\varepsilon',
\]
wo \(\varepsilon'\) durch die Formel
\[
\varepsilon'\sqrt{\frac{\pm a_{11}}{2\pi\sqrt{-1}}}=\sqrt{\frac{a_{11}}{2\pi}}
\]
gegeben ist.

Zusammensetzung einer beliebigen Determinante mit II.:
\[
\varepsilon''=\varepsilon\cdot
\exp\!\left(\mp\frac{2\pi\sqrt{-1}}8\sum_k\beta_{2k-1,1}(\beta_{2k,1}+2\sigma_k)\right).
\]
Zusammensetzung einer beliebigen Determinante mit III.
\[
\varepsilon''=\varepsilon.
\]
Wir fügen diesen Untersuchungen nur noch einige Bemerkungen hinzu, welche sich auf die Bestimmung der Coefficienten \(a_{ik}\), so wie der untern Grenzpuncte \(\alpha\) bei irgend einer nach \(\Theta\)functionen anzustellenden Entwicklung beziehen.

\subsection*{\S\,93. Berechnung der \(a_{ik}\) für die verschiedenen Arten von \(\Theta\)functionen. Directe Aufstellung der Reihenentwicklungen.}

Wir haben in den Entwicklungen der frühern Abschnitte die untern Grenzpuncte \(\alpha\), d. h. die bevorzugte Berührungscurve \((n-2)^{\mathrm{ter}}\) Ordnung nirgend definirt. Die Betrachtungen der vorigen Paragraphen lehrten nun, dass in der That alle \(2^{p-1}(2^p+1)\) eigentlichen Berührungscurven \((n-2)^{\mathrm{ter}}\) Ordnung die untern Grenzpuncte liefern können, und zwar richtet sich das Auftreten einer oder der andern Berührungscurve nach dem Systeme von fundamentalen Periodicitätsmoduln, welche man zu Grunde legt.

Statt daher diese zuerst zu bestimmen, kann man davon ausgehen, eine beliebig gewählte Berührungscurve \((n-2)^{\mathrm{ter}}\) Ordnung zu wählen, deren Berührungspuncte in den untern Grenzpuncten der Argumente der \(\Theta\)functionen erscheinen sollen.

Wir haben oben (\(\S\,75\)) gezeigt, wie man die \(2^{p-1}(2^p+1)-1\) verschiedenen, nicht verschwindenden Functionen
\[
(1)\quad
e^{-\frac12\sum P_h\sigma_h}\left(\frac{\Theta\left(\frac P2\right)}{\Theta(0)}\right)^4
\]
rational und algebraisch durch die Wurzeln der Gleichung ausdrücken kann, von welchen die Auffindung jener Berührungscurve abhängt. Wenn man die bevorzugte Berührungscurve kennt, welche wir hier nach Belieben wählen, so kann man jeden der Ausdrücke (1) durch die beiden Wurzeln jener Gleichung ausdrücken, welche der bevorzugten Curve und irgend einer andern angehören. Zwar weiss man a priori nicht, wie die verschiedenen Ausdrücke (1) sich diesen verschiedenen Wurzeln zuordnen. Aber wenn man die symmetrischen Functionen der Ausdrücke (1) bildet, so erhält man vollständig bekannte Ausdrücke, welche nur noch von jener willkürlich bevorzugten Wurzel abhängen.

Diese so gebildeten Gleichungen kann man dazu benutzen, um die Grössen \(e^{a_{ik}}\) durch algebraische bekannte Functionen annähernd zu berechnen. In der That gehört hiezu nichts, als dass man die \(\Theta\)functionen, die in den Gleichungen (1) vorkommen, zunächst auf ihre ersten Glieder beschränkt, und sich dadurch Gleichungen bildet, welche die Auflösung direct gestatten und als erste Annäherungsgleichungen hinreichen. Indem man die Annäherung dann fortsetzt, erhält man die Grössen \(e^{a_{ik}}\) mit beliebiger Genauigkeit.

Indem man also von einer willkürlich gewählten Berührungscurve ausgeht, findet man \(a_{ik}\), welche einer bestimmten Classe von \(\Theta\)functionen angehören, und man findet also die Mittel, um eine der Entwicklungen nach \(\Theta\)functionen anzustellen. Setzt man diese gefundenen Werthe der \(a_{ik}\) in die Ausdrücke (1) ein, so ergiebt sich aus der Vergleichung derselben mit den entsprechenden algebraischen Ausdrücken, welche der verschiedenen Berührungscurven den verschiedenen Ausdrücken (1) zugeordnet sind.

Die hier angedeutete Methode stimmt genau mit derjenigen überein, welche in der Theorie der elliptischen Functionen sich ergibt, um die Grösse \(q\) durch \(\sqrt{k}\) und \(\sqrt{k'}\) aus denjenigen Formeln auszudrücken, welche auf p. 184 der Fundamente gegeben sind. Sie giebt aber überhaupt eine directe Anweisung, die Lösung des Umkehrproblems auf eine nicht vollkommen unübersehbare Anzahl von Operationen zurückzuführen. Man sieht, dass es für die Herstellung der Endformeln \emph{nicht} nothwendig ist, die Lage der Verzweigungspuncte und ihre Zuordnung zu den einem Puncte \(O\) entsprechenden Wurzeln \(s_x\) vermittelst willkürlich angenommener Schleifensysteme zu bestimmen, die Periodicitätsmoduln durch bestimmte Integrale auszurechnen u. s. w.

Die Aufstellung der Endformeln verlangt erstlich die Kenntniss der \(a_{ik}\), von welcher soeben die Rede war, sodann aber die Kenntniss derjenigen Combinationen von Integralen erster Gattung, welche wir die zugehörigen Normalintegrale genannt haben. Sind \(I_1,I_2\ldots I_p\) irgend welche von einander unabhängige Integrale erster Gattung, und ist das Umkehrproblem in der Form gegeben:
\[
(2)\quad
\begin{aligned}
V_1&=\sum_{i=1}^{i=p}\int_{\alpha^{(i)}}^{x^{(i)}}dI_1,\\
V_2&=\sum_{i=1}^{i=p}\int_{\alpha^{(i)}}^{x^{(i)}}dI_2,\\[-0.25em]
&\quad\cdots\\[-0.25em]
V_p&=\sum_{i=1}^{i=p}\int_{\alpha^{(i)}}^{x^{(i)}}dI_p,
\end{aligned}
\]
so kommt es darauf an, diejenigen Combinationen
\[
u_k=\beta_{1k}I_1+\beta_{2k}I_2\ldots+\beta_{pk}I_p
\]
zu finden, welche die Normalintegrale erster Gattung sind, oder die Combinationen
\[
v_k=\beta_{1k}V_1+\beta_{2k}V_2\ldots+\beta_{pk}V_p,
\]
welche als Argumente der \(\Theta\)functionen auftreten. Die Coefficienten \(\beta\) sind zu bestimmen.

Nun konnten wir nach den Betrachtungen des zehnten Abschnitts die Function
\[
\left[\frac{\Theta\left(v+\frac P2\right)}{\Theta(v)}\right]^4
\]
als genau bestimmte algebraische Function der \(x\) darstellen. Differenziren wir die durch diese Darstellung gegebene Gleichung, so erhalten wir eine Gleichung der Form:
\[
A_1dv_1+A_2dv_2\ldots+A_pdv_p=B^{(1)}dx^{(1)}+B^{(2)}dx^{(2)}\ldots+B^{(p)}dx^{(p)},
\]
oder auch, wenn wir die Differentiale der \(x\) aus (2) durch die \(dV\) ausdrücken:
\[
=C_1dV_1+C_2dV_2\ldots+C_pdV_p.
\]
Die \(A,C\) enthalten sowohl die \(v\) als die \(x\); aber indem man die \(x\) den \(\alpha\) gleichsetzt, verschwinden alle \(v\), und die Coefficienten \(A,C\) werden hiemit völlig bekannte constante Grössen. Wir erhalten also lineare Gleichungen mit constanten Coefficienten zwischen den \(dv\) und \(dV\). Bilden wir \(p\) solcher Gleichungen, so können wir also die einen Differentiale linear durch die andern ausdrücken, und die erhaltenen Coefficienten sind die Grössen \(\beta\), welche gesucht wurden.

So ist es möglich, von vornherein mit beliebiger Annäherung zur Aufstellung der Reihenentwicklungen zu gelangen. Aber wir sind hiemit an die Grenze derjenigen Untersuchungen gelangt, welche mit Hülfe der in diesem Werke zu entwickelnden Methoden anzustellen sind. Die genauere Ausführung der angedeuteten Betrachtungen verlangt Methoden, welche den hier entwickelten fern liegen, und müssen ihrem Umfange nach einer selbständigen Darstellung vorbehalten bleiben.
\end{document}
